% Options for packages loaded elsewhere
\PassOptionsToPackage{unicode}{hyperref}
\PassOptionsToPackage{hyphens}{url}
\PassOptionsToPackage{dvipsnames,svgnames,x11names}{xcolor}
%
\documentclass[
  11pt,
  oneside]{book}
\usepackage{amsmath,amssymb}
\usepackage{iftex}
\ifPDFTeX
  \usepackage[T1]{fontenc}
  \usepackage[utf8]{inputenc}
  \usepackage{textcomp} % provide euro and other symbols
\else % if luatex or xetex
  \usepackage{mathtools}
\usepackage{unicode-math} % this also loads fontspec
  \defaultfontfeatures{Scale=MatchLowercase}
  \defaultfontfeatures[\rmfamily]{Ligatures=TeX,Scale=1}
\fi
\usepackage{lmodern}
\ifPDFTeX\else
  % xetex/luatex font selection
\fi
% Use upquote if available, for straight quotes in verbatim environments
\IfFileExists{upquote.sty}{\usepackage{upquote}}{}
\IfFileExists{microtype.sty}{% use microtype if available
  \usepackage[]{microtype}
  \UseMicrotypeSet[protrusion]{basicmath} % disable protrusion for tt fonts
}{}
\makeatletter
\@ifundefined{KOMAClassName}{% if non-KOMA class
  \IfFileExists{parskip.sty}{%
    \usepackage{parskip}
  }{% else
    \setlength{\parindent}{0pt}
    \setlength{\parskip}{6pt plus 2pt minus 1pt}}
}{% if KOMA class
  \KOMAoptions{parskip=half}}
\makeatother
\usepackage{xcolor}
\usepackage[a4paper,top=23mm,bottom=23mm,left=23mm,right=23mm]{geometry}
\usepackage{longtable,booktabs,array}
\usepackage{calc} % for calculating minipage widths
% Correct order of tables after \paragraph or \subparagraph
\usepackage{etoolbox}
\makeatletter
\patchcmd\longtable{\par}{\if@noskipsec\mbox{}\fi\par}{}{}
\makeatother
% Allow footnotes in longtable head/foot
\IfFileExists{footnotehyper.sty}{\usepackage{footnotehyper}}{\usepackage{footnote}}
\makesavenoteenv{longtable}
\setlength{\emergencystretch}{3em} % prevent overfull lines
\providecommand{\tightlist}{%
  \setlength{\itemsep}{0pt}\setlength{\parskip}{0pt}}
\setcounter{secnumdepth}{-\maxdimen} % remove section numbering
\usepackage{fontspec}
\IfFileExists{UEOT_CJK.ttf}{
  \setmainfont{UEOT_CJK.ttf}[BoldFont=UEOT_CJK.ttf,BoldFeatures={FakeBold=2},ItalicFont=UEOT_CJK.ttf,ItalicFeatures={FakeSlant=.15}]
  \setsansfont{UEOT_CJK.ttf}[BoldFont=UEOT_CJK.ttf,BoldFeatures={FakeBold=2}]
  \setmonofont{UEOT_CJK.ttf}
}{
  \setmainfont{Noto Sans CJK SC}
  \setsansfont{Noto Sans CJK SC}
  \setmonofont{Noto Sans CJK SC}
}
\XeTeXlinebreaklocale "zh"
\XeTeXlinebreakskip=0pt plus 1pt
\usepackage{mathtools}
\usepackage{fvextra}
\usepackage{etoolbox}
\usepackage{newunicodechar}
\newunicodechar{ℏ}{\ensuremath{\hbar}}
\newunicodechar{∎}{\ensuremath{\square}}
\usepackage{fancyhdr}
\pagestyle{fancy}
\fancyhf{}
\fancyhead[L]{\small UEOT PHYSICS}
\fancyhead[R]{\small 物理学 · 第三版}
\fancyfoot[C]{\small\thepage}
\setlength{\headheight}{16pt}
\renewcommand{\headrulewidth}{0.3pt}
\setcounter{tocdepth}{0}
\setlength{\parindent}{0pt}
\setlength{\parskip}{5pt plus 1pt}
\setlength{\emergencystretch}{8em}
\setlength{\tabcolsep}{3pt}
\renewcommand{\arraystretch}{1.16}
\AtBeginEnvironment{longtable}{\small}
\allowdisplaybreaks[2]
\renewcommand{\contentsname}{目录}
\makeatletter
\renewcommand{\maketitle}{
  \begin{titlepage}
  \thispagestyle{empty}
  \vspace*{30mm}
  {\color{teal}\fontsize{46}{52}\selectfont UEOT\par}
  \vspace{5mm}
  {\fontsize{32}{40}\selectfont\bfseries 物理学\par}
  \vspace{11mm}
  {\Large 从对象数学到量子物理\par}
  \vspace{4mm}
  {\large 理论推导、形式证明与研究边界\par}
  \vspace{10mm}
  {\color{teal}\rule{\textwidth}{1pt}\par}
  \vspace{7mm}
  {\large 第三版\quad 3.4 正式整合版\par}
  \vspace{6mm}
  {钱德桂（Degui Qian）\par}
  \vspace{5mm}
  {63章\quad 6个附录\quad UEOT-QM Kernel 1.8\par}
  \vfill
  {对象与协议\quad 量子操作与场论\quad 证明与反例\par}
  \vspace{7mm}
  {2026年10月20日\par}
  \end{titlepage}
}
\makeatother
\ifLuaTeX
  \usepackage{selnolig}  % disable illegal ligatures
\fi
\IfFileExists{bookmark.sty}{\usepackage{bookmark}}{\usepackage{hyperref}}
\IfFileExists{xurl.sty}{\usepackage{xurl}}{} % add URL line breaks if available
\urlstyle{same}
\hypersetup{
  pdftitle={UEOT物理学},
  pdfauthor={钱德桂（Degui Qian）},
  colorlinks=true,
  linkcolor={teal},
  filecolor={Maroon},
  citecolor={Blue},
  urlcolor={teal},
  pdfcreator={LaTeX via pandoc}}

\title{UEOT物理学}
\usepackage{etoolbox}
\makeatletter
\providecommand{\subtitle}[1]{% add subtitle to \maketitle
  \apptocmd{\@title}{\par {\large #1 \par}}{}{}
}
\makeatother
\subtitle{从对象数学到量子物理：理论推导、形式证明与研究边界}
\author{钱德桂（Degui Qian）}
\date{2026-10-20}

\begin{document}
\frontmatter
\maketitle

{
\hypersetup{linkcolor=}
\setcounter{tocdepth}{0}
\tableofcontents
}
\mainmatter
\hypertarget{ch-0}{%
\chapter{0. 文档定位、论域与证明责任}\label{ch-0}}

本书采用与 UEOT Core
完整数学主文相同的方式：一份连续、自足、带稳定编号的 Markdown
主文件，包含定义、推导、标准定理的适用条件、反例、应用、接口和可执行验证。它不是历史版本补丁集合，不要求先寻找旧版才能补齐关键公式。

理论起点仍是原始 UEOT
论文提出的对象持续、结构化开放、相位一致、预测与资源约束问题。后续量子分支提供其具体物理实现。本文接入通用数学
Core，但\textbf{不把通用对象数学当作已足以唯一产生量子物理的公理}。

本书整合的量子主线包括有限维对象与测量、P2
量子选择、连续变量、自旋与相同粒子、多体与开放系统、自由相对论场、相互作用与重整化、QED，以及标准模型候选和三代研究。这里``完整''指这一声明架构的完整主文；不指完成所有量子物理专门学科、严格构造四维相互作用
QFT、量子引力或终极统一理论。

\hypertarget{ux5185ux5bb9ux72b6ux6001ux4e0eux8bc1ux660eux8d23ux4efbux5fc5ux987bux5206ux5f00}{%
\section{0.1
内容状态与证明责任必须分开}\label{ux5185ux5bb9ux72b6ux6001ux4e0eux8bc1ux660eux8d23ux4efbux5fc5ux987bux5206ux5f00}}

\begin{longtable}[]{@{}
  >{\raggedright\arraybackslash}p{(\columnwidth - 4\tabcolsep) * \real{0.1000}}
  >{\raggedright\arraybackslash}p{(\columnwidth - 4\tabcolsep) * \real{0.3600}}
  >{\raggedright\arraybackslash}p{(\columnwidth - 4\tabcolsep) * \real{0.5400}}@{}}
\toprule\noalign{}
\begin{minipage}[b]{\linewidth}\raggedright
标记
\end{minipage} & \begin{minipage}[b]{\linewidth}\raggedright
含义
\end{minipage} & \begin{minipage}[b]{\linewidth}\raggedright
证明责任
\end{minipage} \\
\midrule\noalign{}
\endhead
\bottomrule\noalign{}
\endlastfoot
D & 定义、约定、协议 & 给出类型、支撑、单位 \\
P & 本文给出证明或显式推导 & 条件性数学结论；不自动表示原创 \\
K & 标准数学/物理结果 & 明确论域；重定理不伪装成已自行完整证明 \\
H & 量子或粒子领域附加原则 & 不由 Core 自动推出；须有独立理由 \\
S & UEOT 组织解释和跨模块连接 & 不改变原物理预测时不称新自然定律 \\
E & 实验或数值证据 & 区分真实数据、合成数据、代数检查 \\
N / O & 反例、不可能性 / 开放问题 & 不进入已证正向依赖链 \\
\end{longtable}

本书的 P 编号在正文有证明或推导；采用的重型 K
定理列入附录。数值检查不是机器形式化证明，也不是自然界实证。

另设一条独立的\textbf{验证轴}：正文证明、机器内核证明、精确有限计算、浮点数值检查、合成数据检验、实际实验数据。一个条目可以同时是``附加物理假设下的
P''和``已经通过 Lean
内核的条件定理''；内核通过不会删除其物理假设。反过来，一项实验支持的是指定模型、装置、协议和误差范围，也不会把有限支持提升为无条件数学定理。正文的
P/K/H 分类标记论证角色；验证轴回答目前怎样核对它。另由第48章的
dependency\_\allowbreak class
字段明确Core/桥/扇区/标准物理/有效模型/可观测后果/开放猜想角色，三轴分别保留。{[}S2，§0.3、§28.2；S3，§1--2；S4，§2、§8--9{]}

\hypertarget{ux539fux59cbux601dux60f3ux7684ux4e25ux683cux843dux70b9}{%
\section{0.2
原始思想的严格落点}\label{ux539fux59cbux601dux60f3ux7684ux4e25ux683cux843dux70b9}}

\begin{longtable}[]{@{}
  >{\raggedright\arraybackslash}p{(\columnwidth - 4\tabcolsep) * \real{0.2500}}
  >{\raggedright\arraybackslash}p{(\columnwidth - 4\tabcolsep) * \real{0.3400}}
  >{\raggedright\arraybackslash}p{(\columnwidth - 4\tabcolsep) * \real{0.4100}}@{}}
\toprule\noalign{}
\begin{minipage}[b]{\linewidth}\raggedright
原始思想
\end{minipage} & \begin{minipage}[b]{\linewidth}\raggedright
有效的量子实现
\end{minipage} & \begin{minipage}[b]{\linewidth}\raggedright
不允许的跳跃
\end{minipage} \\
\midrule\noalign{}
\endhead
\bottomrule\noalign{}
\endlastfoot
对象是可持续关系结构 & 协议未来等价、有效因子、历史运输 &
一个名称或波函数标签自动定义对象 \\
相位闭合 & 射影表示、holonomy、边界条件；规范异常一致性 &
闭轨积分直接推出 CCR \\
Π/Φ & Fisher--作用量的条件模型、量子信息成本、控制目标 &
任意两个函数唯一选择 QM \\
父对象与子接口 & 张量组合、纯化、instrument dilation &
一般父对象必为纯量子父系统 \\
观测与分辨率 & POVM、instrument、公共交换代数、过程粗化 &
粗化本身导致客观坍缩 \\
GOD/GOA & 指定量子控制任务的最优控制和长期结构 &
所有量子态自然优化同一目标 \\
粒子``最小闭合'' & 不可约表示、无异常的物理实现包 & Betti
数等于世代数 \\
跨尺度统一 & 类型明确的 Core/过程/QFT 桥接 & 同名 RG 就是同一变换 \\
\end{longtable}

\hypertarget{ux4e3bux6587ux4f7fux7528ux65b9ux5f0f}{%
\section{0.3 主文使用方式}\label{ux4e3bux6587ux4f7fux7528ux65b9ux5f0f}}

第1--9章给基础和测量，第 10--15 章给连续变量与对称性，第 16--20
章给多体、环境和控制，第 21--28 章给 QFT/RG/QED，第 29--32
章给粒子结构及研究边界，第 33--35
章给接口、应用和装配。附录内嵌验证代码与带版本身份的实际回执；全书新增模块见0.5节。

统一的逻辑方向是：\textbf{声明物理实现 → 定义可实施协议 → 计算响应 →
建立对象与预测证书 → 在附加任务下求控制 →
检验领域假设}。不能把希望得到的结论反向写入``对象必须如此''的定义，再称为推导。

\hypertarget{ux5b8cux6574ux6027ux5ba1ux67e5ux4e09ux4e2aux4e0dux540cux7684ux5b8cux6210}{%
\section{0.4
完整性审查：三个不同的``完成''}\label{ux5b8cux6574ux6027ux5ba1ux67e5ux4e09ux4e2aux4e0dux540cux7684ux5b8cux6210}}

本版的结论是：UEOT
已有可装配的量子物理研究架构；在明确量子实现和领域假设后，许多模块可以完整定义、推导和计算。它尚未仅从通用
Core
唯一导出量子公理、自然界的规范群与全部参数，也没有严格构造四维相互作用场论和量子引力。因此本书完成的是声明范围内的统一理论主文与可运行基准，不能把文档完整性当作终极物理已完成。

逐模块用四个问题审查：定义和证明是否闭合？物理输入是否独立声明？输出是否可计算？候选是否有区分性检验？不使用没有客观分母的``完成百分比''。

\begin{longtable}[]{@{}
  >{\raggedright\arraybackslash}p{(\columnwidth - 4\tabcolsep) * \real{0.2500}}
  >{\raggedright\arraybackslash}p{(\columnwidth - 4\tabcolsep) * \real{0.3400}}
  >{\raggedright\arraybackslash}p{(\columnwidth - 4\tabcolsep) * \real{0.4100}}@{}}
\toprule\noalign{}
\begin{minipage}[b]{\linewidth}\raggedright
模块
\end{minipage} & \begin{minipage}[b]{\linewidth}\raggedright
已达到的状态
\end{minipage} & \begin{minipage}[b]{\linewidth}\raggedright
仍缺的关键条件或工作
\end{minipage} \\
\midrule\noalign{}
\endhead
\bottomrule\noalign{}
\endlastfoot
Core 到量子选择 & 给出经典反模型；列明 OPT 背景与选择原则 &
重建重定理仍为外部依赖；原则的物理必然性未证 \\
态、效应、通道、仪器 & 有限维定义、构造证明、组合、误差闭合 &
现实制备测量的可用性与校准是实验输入 \\
对象、记录、预测 & 未来等价、公共记录、SRN 和有记忆接口 &
对象划分依赖允许协议，不能无条件唯一化 \\
测量和实际化 & 相干扩张、条件态、具体 SSE 的吸收与 Born 条件定理 & Core
未选择该随机方程；本体解释仍不可由任意系统通道区分 \\
连续变量与全局相位 & 自伴域、Weyl 表示、能量截断、环上 holonomy &
物理边界与作用尺度为输入；不能由闭轨整数反推 CCR \\
自旋、统计与多体 & 对称表示、Fock、有限模型与热环境基准 &
相对论自旋统计为外部定理；复杂多体无通用有效求解 \\
开放系统、控制、热力学 & GKSL/Davies 条件、过程记忆、控制梯度、自由能 &
强耦合和无限系统需逐模型闭合 \\
拓扑与稳定性 & 本版给 Berry 丛和 SSH 边态的完整受控例 &
尚无一般相互作用拓扑分类或粒子世代选择 \\
自由场和微扰 QFT & 场、传播、因果、散射、一圈重整化与 QED 基准 &
高阶、红外与非微扰连续极限并未全部完成 \\
格点规范理论 & 本版给有限纯规范测度、规范不变性、连续展开与可穷举模型 &
四维连续极限、QCD 质量隙、手征费米子与实时误差仍需专门研究 \\
标准模型候选 & 规范表示、局域异常、Higgs、Yukawa、中微子接口 &
群、全局形式、耦合、质量与混合参数未唯一推出 \\
三代与味结构 & 低能条件选择、高能双候选、秩一 FCNC 和式与下界 & HGP
来源、排除五色的独立理由、完整谱和现实约束 \\
新物理识别 & 三分支 CP 界、四能级见证、能量与退相干联测 &
尚无本次新增真实实验；排除噪声子类不等于发现新自然定律 \\
引力与宇宙学 & 第47章提供受控起点和必要接口 &
未构造量子引力，未导出暗物质、暗能量和宇宙初态 \\
\end{longtable}

``基于 UEOT''在本书有严格含义：物理模型提供干预概率，Core
组织未来等价、载体、持续性、资源与控制，领域桥接说明何时可以合成。若删除
UEOT 解释后预测完全等同标准 QM，这部分属于 S/K
层的组织与重建；只有额外假设产生不同且可识别的概率约束，才进入新物理候选
H 的检验。

\hypertarget{ux7edfux4e00ux4e3bux6587ux7684ux88c5ux914dux4e0eux672cux7248ux589eux91cf}{%
\section{0.5
统一主文的装配与本版增量}\label{ux7edfux4e00ux4e3bux6587ux7684ux88c5ux914dux4e0eux672cux7248ux589eux91cf}}

第1--35章给完整基础、计算与接口；第36--42章把三代高能审计、味约束、相干见证、实际化证明和研究协议接入主线；第43--46章补全全局相位、拓扑、格点和相关试次统计；第47--48章给研究验收合同与总著维护。第49--56章展开有限内核、误差、随机化、长期平均、条件生存、量子选择分离和一般有限量子Doob构造；第57--59章给完整有限复纯化、保留记录与记忆的真实协议及指定非交换模型条件动力学；第60--62章补上全效应分离与任意联合态复局部层析、完整凸纯性及有类型的纯化酉唯一性、共同因果控制器逐动作记录预算。附录保留符号、来源、全部稳定编号、复现代码和完整命名审计。

阅读第31章三代条件定理时应同时使用第36章的高低能区分；第3章的鞅条件结论由第40章在特定
SSE 内补足吸收证明；第7--8章的
SRN/层析由第38--39章发展成噪声类见证；第26章的跨尺度接口由第45章给出有限规范模型的实际载体。所有这些连接都在一个文档内。

v2.0 附录保存其当时执行的量子主线、前沿、全局相位、SSH
边态、有限格点和相关试次检查；这些是带版本的历史回执。后续内核、精确证书和实际实验的状态，由对应版本的源清单、执行日志和审计回执给出。历史通过不替代当前版本复核，新的局部证明也不反向覆盖未纳入形式化的旧章节。有限数值检查不升级外部重定理的证明状态，不替代全称证明，也不构成实际实验数据。

\hypertarget{ux4eceux7814ux7a76ux613fux671bux5230ux53efux5426ux8bc1ux95eeux9898}{%
\section{0.6
从研究愿望到可否证问题}\label{ux4eceux7814ux7a76ux613fux671bux5230ux53efux5426ux8bc1ux95eeux9898}}

每个候选理论写成四元组
\((\mathfrak M,\Theta,\mathcal U,\mathcal P)\)：模型类、参数域、允许协议、响应规律。给定记录
\(D\)，所谓``解释''必须先确定哪些参数由训练数据固定，哪些量是留出预测。若对任意
\(D\) 都可任意更换模型、环境或协议，模型类不具区分力。

对两候选定义指定协议族上的分离度 \[
\Delta_{\mathcal U}(M_0,M_1)
=\sup_{u\in\mathcal U}\operatorname{TV}(p_0(\cdot|u),p_1(\cdot|u)).
\] 这里 \(\operatorname{TV}(P,Q)=\sup_B|P(B)-Q(B)|\)，有限结果空间中为
\(\tfrac12\sum_y|p_y-q_y|\)。它比较同一协议、同一记录空间上的概率律。
若两边都可调参数，应研究两个响应集合的距离及交集，而不是只比较各自最喜欢的一个实例。第41章的环境等价障碍是集合包含问题；第39章的见证则提供一个明确分离超平面。第47章所有研究合同遵守这一点。

\hypertarget{ux6838ux5fc3ux54f2ux5b66ux5728ux53d8ux5316ux4e2dux8bc6ux522bux5bf9ux8c61}{%
\section{0.7
核心哲学：在变化中识别对象}\label{ux6838ux5fc3ux54f2ux5b66ux5728ux53d8ux5316ux4e2dux8bc6ux522bux5bf9ux8c61}}

UEOT
的出发点是一个可以精确提出的问题：当组成材料、微观状态和环境输入不断变化时，为什么我们仍能追踪一个对象，并据此预测其下一步响应？原始论文把答案集中在持续关系、反馈、开放性和层级组成上。这些思想决定了本书优先研究什么，却还没有独自指定一个
Hamiltonian、一套量子仪器或一个粒子表示谱。{[}S1，§1.3，PDF
第2--4页；§2.2--2.4，第4--8页{]}

本书采用 Core
已完成的数学重述：先固定过程、可实施协议和候选自身的未来，识别共同响应所需的信息；再分别研究这种信息能否递归更新、由何种物理装置实现、怎样经受干预，以及怎样沿历史运输。这里的``关系身份''有可检查的内容：改变坐标后，相同实验仍产生相同概率；保留恰当状态后，完整历史中被删除的信息不再影响声明的未来。对象的存在证书可以强，也可以弱；主张恢复、自维持、遗传或最优控制时，再增加对应机制。{[}S2，§0.1、§3--4、§12.1{]}

这使原始``Being 是持续的
Becoming''具有具体意义：持续的是经验证的响应与实现关系，状态本身可以改变。例如移动波包不必固定在某个位置；受噪声扰动并经恢复的逻辑比特不必始终占据同一个物理态；相同输入输出响应也可能由不同硬件实现。预测身份、物理载体和历史实例因此相互连接，却不是一个概念。

原论文曾把一般稳定存在写成非零一阶同调 \(H_1\ne0\)
的结果。这不能作为本书的普遍定理。取 \(\dot x=-x\)，解为
\(x(t)=x(0)e^{-t}\)，原点全局渐近稳定，而状态空间 \(\mathbb R\) 与吸引子
\(\{0\}\) 都可缩，\(H_1=0\)。反方向，在圆柱 \(\mathbb R\times S^1\) 上取
\(\dot u=u,\ \dot\theta=1\)，状态空间有非平凡一阶同调，\(u=0\)
给闭合轨道，但横向扰动按 \(e^t\)
放大。拓扑环、因果依赖图中的反馈边和动力学稳定性需要不同的定义及连接证明。{[}S1，§2.2，第4页，印刷行147--160；S2，§9.4、§12.2{]}

同样，作用量驻值不等于系统沿时间单调增加某种利益。谐振子的作用量满足变分驻值，其相轨道可以周期运动；它不会因为这个事实就收敛到一个唯一的最大收益状态。RG
的尺度参数也不同于物理时间。要把某一尺度流解释成动力学、把某一漂移解释成最优行为，必须分别提供过程桥和评价桥。{[}S1，§2.3，第6--7页；S2，§18--23、§28.3--4{]}

因此，UEOT
在物理学中的统一性首先体现在共同的识别、实现、运输和证书语言。标准量子模型提供实际概率规律；UEOT
让我们明确哪些状态区别与某项任务有关，怎样证明压缩后仍保留所有允许的未来，以及何时需要扩大记忆、操作库或对象边界。能够产生超出标准模型的概率约束时，新增部分成为待检验的物理假设。

\hypertarget{ux6765ux6e90ux5c42ux7ea7ux4e0eux7a33ux5b9aux6587ux672c}{%
\section{0.8
来源层级与稳定文本}\label{ux6765ux6e90ux5c42ux7ea7ux4e0eux7a33ux5b9aux6587ux672c}}

本书对理论含义采用明确的来源次序：Core v3
的冻结定义、前提和证明决定通用数学语义；量子章节补充专属物理结构；可执行代码与
Lean
声明说明哪些具体结论已得到何种核对；原始论文保留研究动机与尚未闭合的猜想。原始哲学措辞不覆盖后来修正的数学定义。{[}S2，§0、附录E；S4，§1--2{]}

本组基础章节使用以下稳定源号，读者可以直接定位原始责任：

\begin{longtable}[]{@{}
  >{\raggedright\arraybackslash}p{(\columnwidth - 4\tabcolsep) * \real{0.2500}}
  >{\raggedright\arraybackslash}p{(\columnwidth - 4\tabcolsep) * \real{0.3400}}
  >{\raggedright\arraybackslash}p{(\columnwidth - 4\tabcolsep) * \real{0.4100}}@{}}
\toprule\noalign{}
\begin{minipage}[b]{\linewidth}\raggedright
源号
\end{minipage} & \begin{minipage}[b]{\linewidth}\raggedright
原始材料
\end{minipage} & \begin{minipage}[b]{\linewidth}\raggedright
本书采用的角色
\end{minipage} \\
\midrule\noalign{}
\endhead
\bottomrule\noalign{}
\endlastfoot
S1 & \texttt{UEOT\_\allowbreak v3\_\allowbreak top.\allowbreak pdf}，43页，用户提供本 &
原始对象、双驱动与物理猜想的动机；引文使用 PDF 实际页码及印刷行 \\
S2 & \texttt{UEOT\_\allowbreak Core\_\allowbreak Mathematics\_\allowbreak v3.\allowbreak 0\_\allowbreak Complete.\allowbreak md}，2026-09-05 &
冻结通用定义、稳定 P-ID、数学证明和适配条件 \\
S3 & \texttt{CORE\_\allowbreak V3\_\allowbreak OPEN\_\allowbreak PORT\_\allowbreak MATRIX.\allowbreak md}，2026-10-07 &
七端口的精确缺口、工作包与历史源---Lean 映射 \\
S4 & \texttt{CORE\_\allowbreak V3\_\allowbreak POST\_\allowbreak CORE\_\allowbreak SCIENTIFIC\_\allowbreak CLOSURE.\allowbreak md}，2026-10-07 &
核心理解、四生成器边界和科学闭合的独立审查 \\
S5 & \texttt{THEORY\_\allowbreak COMPLETION\_\allowbreak ROADMAP\_\allowbreak V2.\allowbreak md}，2026-10-07 &
后续研究的依赖与验收规则；其 PLANNED 项不是完成回执 \\
Q2 & \texttt{UEOT\_\allowbreak Physics\_\allowbreak Unified\_\allowbreak v2.\allowbreak 0.\allowbreak md}，2026-09-06 &
本书第0--48章既有连续推导的直接继承文本 \\
\end{longtable}

S3--S5 中的日期、分支、阶段和 CI
数字是各自记录时的事实，不自动成为本书当前版本的运行状态。正式修订保留旧编号和有用证明，在准确位置增补前提或更正公式；开放猜想不会通过更换版本号转成定理。

\hypertarget{ux672cux6b63ux5f0fux6574ux5408ux7248ux7684ux7ed3ux6784ux4e0eux9605ux8bfbux8defux5f84}{%
\section{0.9
本正式整合版的结构与阅读路径}\label{ux672cux6b63ux5f0fux6574ux5408ux7248ux7684ux7ed3ux6784ux4e0eux9605ux8bfbux8defux5f84}}

本版承接《UEOT PHYSICS
统一主文2.0》的完整49章，不缩减既有量子基础、QFT/QED、世代、实际化、拓扑、格点和实验统计推导。第0--13章明确原始思想与严格Core、量子选择背景及测量/预测接口；第14--32、36--46章保留高级物理及反模型；第33--35、47--48章更新实际接口、误差、路线与总著维护；第49--50章展开既有有限量子内核和稳定性证明，第51--53章补齐本轮实际混合、Cesàro
与 Core 条件生存桥。

理论作者为钱德桂（Degui
Qian）；原始思想及项目成果的身份由附录C和源清单保留。本版是在该理论工作上进行的持续整合、推导补全与协作审查。标准物理定理仍明确归属于其原始文献，机器验证和实验身份分别记录。

从理论核心入手，可依次阅读第0--5、49--62章，再按所需物理域进入后续章节。核查三代问题时必须把第31、36--37与49.7一起阅读。附录A给符号和外部定理责任，附录B保留稳定结果注册，附录C记录来源，附录D包含全部继承验证代码及3.0版实际执行回执；3.1--3.4版仅对未变程序核对身份。附录E给编译维护合同，附录F列出本版实际形式化范围和全部命名公理请求。

\hypertarget{ch-1}{%
\chapter{1. UEOT Core 到量子物理的桥接}\label{ch-1}}

\hypertarget{ux64cdux4f5cux5386ux53f2ux4e0eux91cfux5b50ux54cdux5e94}{%
\section{1.1
操作历史与量子响应}\label{ux64cdux4f5cux5386ux53f2ux4e0eux91cfux5b50ux54cdux5e94}}

协议 \(u\) 是一串可实施制备、变换、测量与允许的条件反馈。记录历史为
\(h\)。对固定未来目标 \(Y^+\)，定义规范响应 \[
C(h)=\bigl(P(Y^+\in\cdot\mid h,u)\bigr)_{u\in\mathcal U}.
\]
这里的``未来''是已声明的记录目标，可以包含一段完整实验路径；它不表示全宇宙的所有未来事实。一般概率表述取标准
Borel 历史和记录空间、同一历史参考律
\(\mu\)、有限或可数协议族，并选择可测响应核
\(K_u(h,\cdot)\)。对于之后才实施的协议，历史参考律保持相同；若干预改变了历史生成机制，须另给共同结构模型及支撑说明。{[}S2，§1、§2.3--4、§4.1{]}

若统计量 \(M\) 对每个协议充分，即存在核 \(L_u\) 使
\(K_u(H,\cdot)=L_u(M(H),\cdot)\)
几乎处处，则取全部协议及各未来空间可数决定事件族的共同概率一集合，得到
\[
C(H)=\bigl(L_u(M(H),\cdot)\bigr)_{u\in\mathcal U}=\varphi(M(H)).
\] 右侧为可测函数。反过来，每个 \(K_u\) 是 \(C\) 的坐标，因此 \(C\)
本身充分。这证明的是给定协议、目标和参考律下的最小预测因子，而不是预设经典或量子状态空间，也不唯一指定物理载体。{[}S2，P-PRED-01{]}

可数性不能省略。若 \(H\) 均匀分布于 \([0,1]\)，协议 \(u\in[0,1]\)
的核版本仅在 \(h=u\)
处被改动，每个版本仍分别几乎处处相同，但全部例外点的并是整个区间。逐协议的模零等式因此不保证不可数协议的共同点态等式。点态给定的结构核可以另走点态路线；可数稠密协议加共同连续版本也可以扩展，但必须显式具备这些条件。{[}S2，N-PRED-01{]}

量子实现提供具体概率生成器，例如 \[
p(a_1,\ldots,a_T)=
\operatorname{Tr}\bigl[
\mathcal M_{a_T}^{(T)}\circ\cdots\circ
\mathcal M_{a_1}^{(1)}(\rho_0)\bigr].
\] 自适应协议使后续 instrument 依赖已见记录，不能依赖未来结果。

在离散记录中，若协议族对动作前缀及正概率观察后的允许延续封闭，则当前响应还决定更新后的响应：
\[
P(B\mid h,a,o)=\frac{P(o,B\mid h,a)}{P(o\mid h,a)},\qquad P(o\mid h,a)>0.
\]
分子和分母都是当前规范响应所保存的量。对零概率分支，条件态不由此比值决定；数学上选取任意延拓不产生经验支持。这个分母条件也将贯穿量子仪器、有限坐标和后选择误差。{[}S2，P-PRED-03{]}

\hypertarget{ux901aux7528ux6838ux5fc3ux4e0dux9009ux62e9ux91cfux5b50ux6247ux533a}{%
\section{1.2
通用核心不选择量子扇区}\label{ux901aux7528ux6838ux5fc3ux4e0dux9009ux62e9ux91cfux5b50ux6247ux533a}}

\protect\hypertarget{p-qbr-01}{}{}

\textbf{P-QBR-01：经典实现反模型。}
通用预测因子、状态闭合、访问载体、持续性和控制结构可以在有限经典 Markov
系统上同时实现，因此这些结构不能单独逻辑蕴含非交换观测、复杂 Hilbert
扇区或量子纯化。

\textbf{证明。}
取有限状态受控链，将历史按全部未来响应分组，以状态读出和经典概率核实现。访问充分族、其极小元与击中集均可定义；有吸引复归类时有持续性，有限折扣任务有最优策略。其状态空间仍为概率单纯形，所有经典效应可共同读取。混合经典态不能是纯经典联合态的边缘，因为纯联合态是确定点。故量子特有结论不是通用前提的逻辑后果。\(\square\)

这不否定 UEOT
的统一组织价值，只界定了它必须通过领域原则才能承担的物理内容。

\hypertarget{ux8868ux793aux89c4ux8303ux4e0eux53efux89c2ux6d4bux7b49ux4ef7}{%
\section{1.3
表示、规范与可观测等价}\label{ux8868ux793aux89c4ux8303ux4e0eux53efux89c2ux6d4bux7b49ux4ef7}}

表示变换 \(T\)
必须同步运输制备、通道、效应和复合规则。若全部允许协议概率保持，则两个实现属于该协议的操作等价类。可逆坐标变换不自动是物理规范对称；物理规范轨道是否代表冗余，由具体领域模型决定。

固定目标和协议后，物理 gauge
不可见性可使响应下降到商空间。但``所有实验当前未区分''比``所有物理实验永远不能区分''弱；有限检验只产生有限证书。

\hypertarget{ux56dbux8ba1ux6570ux751fux6210ux5668ux5982ux4f55ux8fdbux5165ux7269ux7406ux63a8ux5bfc}{%
\section{1.4
四计数生成器如何进入物理推导}\label{ux56dbux8ba1ux6570ux751fux6210ux5668ux5982ux4f55ux8fdbux5165ux7269ux7406ux63a8ux5bfc}}

S4 §6.2 对四生成器的正式范围给出精确账本：106 项中 11 项标为
generated、89 项保留为 retained\_\allowbreak adapter、6 项保留为
retained\_\allowbreak boundary。四者的非冗余性指注册推导系统内的
\texttt{nonredundant\_\allowbreak under\_\allowbreak declared\_\allowbreak derivation\_\allowbreak system}。它不是对所有数学证明语言的绝对独立性定理，更不是从四条物理公理无条件推出106项结论。下面说明四种数学机制在量子应用中实际承担什么。

\textbf{M-QD：商与下降。} 给映射 \(q:X\to C\) 和响应 \(F:X\to Y\)，若 \[
q(x)=q(x')\Longrightarrow F(x)=F(x'),
\] 则在 \(q(X)\) 上定义 \(\bar F(q(x))=F(x)\)，其值不依赖代表元，而且
\(F=\bar F\circ q\)。有限集合中这就完成构造；可测空间中还需要对应的可测因子化条件。Core
预测最小性、结构化内部---环境分解利用这一模式。量子中的对应问题是：全部允许记录的
Born
响应，是否只依赖所保留的有限坐标？它不会在证明下降的同时选择复数标量域或
Born 配对。{[}S2，P-PRED-01、P-INT-02；S3，附录106项映射{]}

\textbf{M-TC：有限定量运输。} 若真实演化 \(x_{t+1}\) 与参考映射
\(F_t(x_t)\) 的缺陷不超过 \(\varepsilon_t\)，而 \(F_t\) 的 Lipschitz
常数为 \(L_t\)，则插入 \(F_t(x_t)\) 得 \[
e_{t+1}\le L_te_t+\varepsilon_t,
\] 归纳展开为 \[
e_n\le\left(\prod_{j<n}L_j\right)e_0+
\sum_{k<n}\varepsilon_k\prod_{j=k+1}^{n-1}L_j.
\]
若选定操作距离下的量子通道收缩，\(L_t\le1\)，误差便受初始误差与各步缺陷之和控制。但``通道在此距离下收缩''必须由量子正性、归一性和效应拉回证明；它不是从运输模板自动产生的假设。后选择分母、记忆辅助系统和协议支撑也不能省略。{[}S2，P-DYN-03--04、P-ID-01--02；本书§4.4{]}

\textbf{M-PE：正特征结构。} 在本原非负次随机矩阵 \(Q\) 的例子中，正左右
Perron 数据满足 \[
qQ=rq,\quad Qh=rh,\quad q\mathbf1=1,\quad qh=1.
\] 于是 \(qQ^n=r^nq\)，存活条件化后仍为 \(q\)；并且 \[
Q^\uparrow_{ij}=\frac{Q_{ij}h_j}{rh_i},\qquad
\sum_jQ^\uparrow_{ij}=1,\qquad
\sum_iq_i h_iQ^\uparrow_{ij}=q_jh_j.
\] 正性、谱数据和归一化共同产生准平稳分布与条件过程。量子 CP
映射也作用于正锥，因而可以研究相应的正算子谱问题；然而不可约性、本原性、存在正特征算子、谱隙及条件化类型必须针对该映射重新建立。经典矩阵的逐项非负和量子映射的完全正不是同一个输入类型。第53章已在对角测量制备模型中证明这座实际桥，直接消费
Core
谱定理；一般非交换谱义务仍须独立处理。{[}S2，§10.2，P-QSD-02；S4，§6.2{]}

\textbf{M-OI：占据极限与不变闭合。} 对有限量子系统的固定 CPTP 通道
\(\mathcal T\)，考虑时刻态的 Cesàro 平均 \[
\bar\rho_N=\frac1N\sum_{t=0}^{N-1}\mathcal T^t(\rho_0).
\] 有限维密度矩阵集合紧，故存在收敛子列；而 \[
\mathcal T(\bar\rho_N)-\bar\rho_N=
\frac{\mathcal T^N(\rho_0)-\rho_0}{N},\qquad
\|\mathcal T(\bar\rho_N)-\bar\rho_N\|_1\le\frac2N.
\]
由线性映射的连续性，每个子列极限都是不变态。这说明不变结构怎样从占据平均得到；它没有证明唯一性、逐时刻混合，或者每条随机轨迹的经验极限。无限维情形的紧性也不能直接沿用。量子酉周期运动可以具有这样的平均极限，同时永远不收敛到单个态。{[}S2，P-PER-02、P-GOA-01；这里的有限量子版本是所声明
CPTP 模型中的直接推导{]}

第52章已把实际密度集合紧性、Cesàro 矩阵残差、全部效应的 \(1/N\)
界以及平稳子极限存在写成 Lean 定理。该机器证明使用 Euclidean 算符范数的
\(2/N\)
界保证矩阵残差趋零；本节的迹范数公式保留解析责任，不因两者常数相同就被自动计入形式化。

四种机制解释了内核为什么可以复用。实际量子层仍须负责复矩阵、半正定性、仪器完备性、真实迹配对和实验接口。新增一个可复用证明，不会相应增加一条自然界基本定律；保留一个专属适配，也不是核心失败。

\hypertarget{ux5bf9ux8c61ux5b9eux73b0ux4e0e-ux3c9-ux7684ux91cfux5b50ux8bedux4e49}{%
\section{1.5 对象、实现与 Ω
的量子语义}\label{ux5bf9ux8c61ux5b9eux73b0ux4e0e-ux3c9-ux7684ux91cfux5b50ux8bedux4e49}}

量子对象证书首先记录候选、其自身未来、允许操作及环境接口；再记录预测与递归闭合、干预条件、历史运输和误差。恢复、纠错、自维持或功能目标是按主张添加的模块。一个纯态标签、一个拓扑数或一个有限矩阵均不能替代这份证书。{[}S2，D-OBJ、D-OBJ-STRONG{]}

给定有限访问接口集
\(V\)，若增加可访问接口只扩大允许实验，则充分集族向上闭合。设极小充分集族为
\(\mathcal M_0\)。仅删除访问权限 \(D\) 后仍能预测目标，当且仅当 \[
\exists M\in\mathcal M_0:\ M\cap D=\varnothing.
\] 否定它便得到``击中每个最小实现''的条件，极小击中集就是访问
blocker。量子实现中必须说明访问是经典记录、局部测量、联合测量还是相干控制；更强访问不能凭空视为可实施。{[}S2，§13.1--2{]}

物理删除另作定义。固定证书和时域 \(T\)，令 \(\chi_{O,T}(D)=1\)
表示移除机制集 \(D\) 后仍通过证书。若失败向上单调，则 \[
\mathcal C_{\Omega,T}=\min_{\subseteq}\{D:\chi_{O,T}(D)=0\},\qquad
\chi_{O,T}(D)=0\iff\exists C\in\mathcal C_{\Omega,T},\ C\subseteq D.
\] 正向在有限失败子集中取极小元，反向由单调性。这里的
\(\Omega_T=(\mathcal M,\mathcal C_{\Omega,T})\)
是给定机制和证书的失效超图。它与原论文几何环的直觉有联系，但没有把非零同调设成普遍存在门槛。删除抑制器若会恢复系统，失败不再单调，须保留完整干预响应。{[}S2，P-OMG-01{]}

将访问 blocker 解释成物理失效还须满足 frame
条件：剩余读出及动力学保持、删除语义与集合对应、恢复和供能等机制已纳入，并且剩余访问充分性确实等价于目标证书。切断一个供电端口可能同时破坏原本仍``可访问''的存储器，正是不能跳过这一桥的例子。{[}S2，§13.3{]}

\hypertarget{godgoa-ux4e0e-ux3c0ux3c6-ux7684ux9644ux52a0ux4efbux52a1}{%
\section{1.6 GOD、GOA 与 Π/Φ
的附加任务}\label{godgoa-ux4e0e-ux3c0ux3c6-ux7684ux9644ux52a0ux4efbux52a1}}

给定量子制备与仪器，并没有同时给定收益。评价桥至少包含允许策略、奖励或风险泛函、时域、资源预算和评价主体/层级。在同一量子存储装置上，``保持未知态''``尽快读出标签''``估计外场''``冷却到基态''可以要求不同动作；对象与器件都没有改变。{[}S2，D-BRIDGE、N-TEL-01；Q2，§20{]}

最简单的逻辑反例是同一过程可选 \(a,b\)，评价
\(J_a=\mathbf1_{\{A_0=a\}}\) 与 \(J_b=\mathbf1_{\{A_0=b\}}\)
给相反最优动作。故本书使用的正式依赖是 \[
(\text{量子过程},\text{对象证书},\text{评价桥})
\longrightarrow V^*\longrightarrow\mathcal A^*\longrightarrow\text{已定类型的长期结构}.
\] 例如有限动作、折扣 \(0<\beta<1\)、有界奖励时，条件态上的 Bellman
算子为 \[
(TV)(\rho)=\max_a\left\{r(\rho,a)+\beta\sum_o p(o\mid\rho,a)
V(\rho_{a,o})\right\}.
\] 对零概率分支可指定任意占位态，其加权项为零。由于各 \(p_o\ge0\)
且和为一，\(\|TV-TW\|_\infty\le\beta\|V-W\|_\infty\)。有限动作的可测性及相应状态空间条件满足时，固定点和最优选择由标准折扣控制构造得到。GOD
首先是达到最大值的动作集合，不保证单值，更不必表现为欧氏空间的梯度。{[}S2，§19；Q2，§20.2{]}

GOA 则要指明讨论的是不变态、Cesàro 极限、轨道
\(\omega\)-极限集、QSD，还是非自治跟踪对象。即使闭环有唯一不变概率，也未必逐时刻混合：两状态确定翻转在唯一均匀不变律之外持续振荡。量子中酉共轭同样允许周期或准周期运动。因此``沿最优方向必趋向唯一最优固定点''不是通用结论。{[}S1，§2.3.2，第6--7页；S2，§21；S4，§4{]}

若使用 \(J=\Pi-\lambda\Phi-C\)，应先给出 \(\Pi,\Phi,C\)
的量纲、符号和独立测量。任意 \(\chi\) 产生 \[
\Pi'=\Pi+\lambda\chi,\quad \Phi'=\Phi+\chi,
\qquad \Pi'-\lambda\Phi'-C=J.
\]
所以相同最优行为不能唯一识别物理轴和信息轴。量子控制能耗、制备时间、记忆需求、预测
KL 或贝叶斯风险可以组成具体桥；Fisher
信息、熵、信息缺口和信息收益也不能只因都叫``信息''就同号相加。{[}S2，§23.1--3、§23.6{]}

若另提出所有环境共享二维驱动 \(z(B)\in\mathbb R^2\)，响应
\(m_e(B)=g_e(z(B))\)，链式法则仅在同一预算点给 \[
\begin{bmatrix}Dm_1(B)\\\vdots\\Dm_E(B)\end{bmatrix}
=\begin{bmatrix}Dg_1(z(B))\\\vdots\\Dg_E(z(B))\end{bmatrix}Dz(B),
\qquad\operatorname{rank}\le2.
\]
这是可反驳的必要条件，尚不是由量子形式主义推出的普适二维律。不同环境分别秩二，或不同预算点的导数拼接成低秩，都不能替代共同同点检验。标量奖励也不强迫二维响应：\(J(a,B)=-\|a-B\|^2/2\)
的最优解 \(a=B\)
可以具有任意维数。{[}S2，P-DDH-04；S3，C5工作包；S4，§5{]}

\hypertarget{ux4e03ux79d1ux5b66ux7aefux53e3ux4e0eux91cfux5b50ux5185ux6838ux7684ux5206ux5de5}{%
\section{1.7
七科学端口与量子内核的分工}\label{ux4e03ux79d1ux5b66ux7aefux53e3ux4e0eux91cfux5b50ux5185ux6838ux7684ux5206ux5de5}}

七端口是开放研究问题的组织方式，不是新增公理或必须依次完成的七个阶段。在量子分支中，它们可以落实为下表的具体责任。{[}S2，§31.2；S3，§3--10；S5，§2--7{]}

\begin{longtable}[]{@{}
  >{\raggedright\arraybackslash}p{(\columnwidth - 4\tabcolsep) * \real{0.2500}}
  >{\raggedright\arraybackslash}p{(\columnwidth - 4\tabcolsep) * \real{0.3400}}
  >{\raggedright\arraybackslash}p{(\columnwidth - 4\tabcolsep) * \real{0.4100}}@{}}
\toprule\noalign{}
\begin{minipage}[b]{\linewidth}\raggedright
端口
\end{minipage} & \begin{minipage}[b]{\linewidth}\raggedright
与实际 QM 的准确连接
\end{minipage} & \begin{minipage}[b]{\linewidth}\raggedright
何种结果可以据此声称
\end{minipage} \\
\midrule\noalign{}
\endhead
\bottomrule\noalign{}
\endlastfoot
C1：存在与一般状态 &
无限维正常态、可测条件版本、隐藏记忆、非紧极限；有限维证明逐项声明扩展条件
& 指定数学模型中的存在、递归或极限定理 \\
C2：对象与协议发现 &
从允许仪器和记录中寻找充分响应空间，分开候选搜索与独立验证 &
给定操作库中的证书；不等同识别全部未知设备 \\
C3：实现与毁损 & 访问权限、联合量子读出、纠错/恢复机制、物理故障图 &
声明操作语义下的极小实现或失效族 \\
C4：身份与长期运输 & 编码变更、漂移器件、时间相关通道、稳态或亚稳态运输
& 固定距离、支撑和时间范围中的运输界 \\
C5：共同低维机制 & 预先锚定能量/信息预算，检查同一点跨环境响应秩 &
拒绝或支持指定局部低维假设，不能唯一命名驱动 \\
C6：目的桥识别 &
独立确定控制任务或实际学习/选择机制，预测未拟合预算下动作 &
指定机制对行为的支持；不从态稳定性读出意图 \\
C7：统计与复现 &
冻结协议、误差、停止规则，保存原始记录、源版本和拒绝标准 &
分开的数学证明、模型检查、实验支持或假设拒绝 \\
\end{longtable}

有限精确矩阵可以完全封闭某个计算实例；其 Lean
证明还可推出该实例在任意有限记录长度上的数学性质。但设备是否实现该矩阵、噪声是否落入被声明模型、训练与检验是否独立，仍属于
C2/C7 的实物与统计责任。S5 中标为 PLANNED
的历史工作包也不能因出现了一个内核就整体改成完成。{[}S5，状态说明；本书§5.5及后续内核证据章{]}

\hypertarget{ch-2}{%
\chapter{2. 量子选择：何时选出复数量子理论}\label{ch-2}}

\hypertarget{ux56faux5b9aux64cdux4f5cux7406ux8bbaux80ccux666f}{%
\section{2.1
固定操作理论背景}\label{ux56faux5b9aux64cdux4f5cux7406ux8bbaux80ccux666f}}

采用有限维因果操作概率理论。量子重构的背景包含系统与有限结果测试、平凡概率系统、顺序和并行复合及其相容律、非空可归一状态、确定效应、粗化和条件测试、有限维且闭的归一态集合、非平凡可区分系统，以及复合系统交换次序的等价。变换相等由\textbf{包含任意辅助系统}的完整操作统计定义；这比只比较单系统输入输出更强。完整背景固定为下述原始论文
Appendix A.1--A.2 的定义和12项条件。

不在背景中预先假定 Hilbert 空间、纯化、局部层析、测量无约束或全部 CPTP
通道。若一开始把通道相等定义为局部输入输出相同，那么再提出 CIT
就成为定义重述，失去了原重构定理所要求的内容。

量子选择必须在这类完整操作理论中陈述，不能只给一个状态锥就声称同时重建了所有物理变换。

\hypertarget{ux4e24ux4e2aux72ecux7acbux9886ux57dfux539fux5219}{%
\section{2.2
两个独立领域原则}\label{ux4e24ux4e2aux72ecux7acbux9886ux57dfux539fux5219}}

\textbf{H-QSEL-1：局部等价 / CIT。} 对系统 \(A,B\) 的通道
\(\Lambda,\mathcal E\)，若每个输入态 \(\rho\) 和每个输出效应 \(f\) 满足
\(f(\Lambda\rho)=f(\mathcal E\rho)\)，则
\(\Lambda=\mathcal E\)，包括所有辅助系统中的操作等价。在上述固定背景中，它等价于复合态的局部可区分性。

\textbf{H-QSEL-2：等价系统纯化 ESP。} 对每个系统
\(A\)，存在一个与其操作等价的纯化系统 \(\widetilde A\)，每个归一态都是
\(A\widetilde A\) 中某纯态的边缘；同一态在此纯化系统中的任意两种纯化只差
\(\widetilde A\) 上的可逆通道。

\textbf{K-QSEL-01。} 在该固定背景与上述两个原则下，Nakahira 的 Theorem 1
刻画有限维复数操作量子理论：对应系统的全部密度矩阵、POVM、Born
配对、Hilbert 张量积和 CPTP
通道。本文调用此项外部重构，未将其完整证明收入 UEOT Lean
内核。\href{https://arxiv.org/abs/2605.23217}{原始预印本}；\href{https://arxiv.org/html/2605.23217v1}{v1
正文及 Appendix A--D}。原文 Appendix A.3
还明确区分：对理论中存在的系统恢复量子结构，不等于证明每个自然数维数的物理系统都必须存在。

为使正文后续自足，后续直接明列并使用这些量子结构，而不把读者理解量子计算规则的责任转给该外部论文。该调用仅服务于``为什么这组选择原则足够''的来源说明。

\hypertarget{ux72ecux7acbux6027ux4e0eux91cfux5b50ux6ee1ux8db3ux6027}{%
\section{2.3
独立性与量子满足性}\label{ux72ecux7acbux6027ux4e0eux91cfux5b50ux6ee1ux8db3ux6027}}

\protect\hypertarget{p-qsel-01}{}{}

\textbf{P-QSEL-01：经典与实量子删除见证。}
有限经典概率理论满足局部层析但不满足非平凡混合态纯化。实量子理论有等价系统纯化，却不满足一般局部层析。

\textbf{证明。}
经典联合分布由各局部确定事件的乘积概率恢复；纯联合分布是一个确定点，其边缘不能混合。实对称密度矩阵可经实谱分解纯化；两个纯化的矩阵因子只差纯化空间正交变换。对两个
rebit， \[
\rho_\pm=\tfrac14(I\otimes I\pm Y\otimes Y).
\] \(Y\otimes Y\) 是实对称矩阵，\(\rho_\pm\) 半正定且不同。单 rebit
实对称效应由 \(I,X,Z\) 张成，与 \(Y\)
正交，故任意产品效应对两态给相同概率。\(\square\)

这个反例可以逐项核对。取
\(Y=\left(\begin{smallmatrix}0&-i\\i&0\end{smallmatrix}\right)\)，虽然
\(Y\) 本身不是实量子可观测量，\(Y\otimes Y\) 却实对称，平方为 \(I\)
且迹为零，其特征值为 \(+1,-1\) 各两次。因此 \(\rho_\pm\) 的谱都是
\(\{1/2,1/2,0,0\}\)。对任意实对称
\(E,F\)，\(\operatorname{Tr}(EY)=\operatorname{Tr}(FY)=0\)，所以 \[
\operatorname{Tr}[(E\otimes F)(\rho_+-\rho_-)]=
\tfrac12\operatorname{Tr}(EY)\operatorname{Tr}(FY)=0.
\] 联合效应 \((I+Y\otimes Y)/2\)
则分别给概率1和0，直接显示局部效应未分离全部联合态。这不是把一个禁用的单
rebit
效应偷放入实验，而是在允许的实联合态空间中展示局部层析缺口。{[}Q2，P-QSEL-01，补充显式计算{]}

复数 QM 的局部层析由 Hermitian 基的张量积张成全部复合 Hermitian
空间得到；纯化由谱分解 \(\rho=\sum_i p_i|i\rangle\langle i|\) 对应
\(|\Psi\rangle=\sum_i\sqrt{p_i}|i\rangle|i\rangle\) 得到。纯化唯一性来自
Schmidt 分解和等谱因子分解。

两个原则的相对不可删减性不等于``QM
在任何语言中恰需两个公理''。将它们合并成一个合取句，或移入背景，会改变公理条数而不改变逻辑内容。历史
CRP/CIT 的纯化、局部层析数学属于既有理论；UEOT
提供父系统与接口的解释位置。

\hypertarget{ux91cdux6784ux7ed3ux8bbaux7684ux4f7fux7528ux8fb9ux754c}{%
\section{2.4
重构结论的使用边界}\label{ux91cdux6784ux7ed3ux8bbaux7684ux4f7fux7528ux8fb9ux754c}}

本书可以按两条明确路线使用后续公式。重构路线接受固定 OPT 背景、CIT、ESP
并调用 K-QSEL-01；直接建模路线则把有限维复 Hilbert、Born
概率、张量复合和 CP
仪器作为量子领域输入。两者都能支持后续数学，但不能把后者的输入倒写成前者已经从
Core 推出的结论。

可数或有限实验会约束设备的响应集合，不能单凭有限通过证明 CIT 或 ESP
的所有系统、态和变换量词。一次成功的纯化制备也不确定唯一的实际环境本体。若研究有限精度的近似纯化或近似层析，需要另给稳定重构定理；严格重构不能仅将等号换成小数误差后照搬。

\hypertarget{ch-3}{%
\chapter{3. Fisher--作用量、复振幅与 Born 的条件推导}\label{ch-3}}

\hypertarget{ux8fdeux7eedux7edfux8ba1ux4f5cux7528ux91cf}{%
\section{3.1
连续统计作用量}\label{ux8fdeux7eedux7edfux8ba1ux4f5cux7528ux91cf}}

在无节点的光滑区域，取 \(\rho>0\)、实相位作用
\(S\)，满足所需衰减或边界条件。给定质量 \(m>0\)、势 \(V\) 和作用尺度
\(\hbar>0\)，定义 \[
\mathcal A[\rho,S]=\int dt\,dx\left[
\rho\left(\partial_tS+\frac{|\nabla S|^2}{2m}+V\right)
+\frac{\hbar^2}{8m}\frac{|\nabla\rho|^2}{\rho}\right].
\] 这一局部、二阶、特定 Fisher 系数的作用量是
\textbf{H}，不是由任意对象闭合自动选出。采用电磁耦合时还须独立指定
\(\nabla S-qA\)、\(q\phi\) 和 gauge 变换。

\protect\hypertarget{p-qact-01}{}{}

\textbf{P-QACT-01：Madelung--Schrödinger 条件等价。}
对上述作用量做紧支撑变分，得到 \[
\partial_t\rho+\nabla\cdot(\rho\nabla S/m)=0,
\] \[
\partial_tS+\frac{|\nabla S|^2}{2m}+V
-\frac{\hbar^2}{2m}\frac{\Delta\sqrt\rho}{\sqrt\rho}=0.
\] 令 \(\psi=\sqrt\rho\,e^{iS/\hbar}\)，两式等价于 \[
i\hbar\partial_t\psi=
-\frac{\hbar^2}{2m}\Delta\psi+V\psi
\] 在该光滑无节点区域的实部与虚部。

\textbf{证明。} 写 \(S\mapsto S+\epsilon\eta\)，其中 \(\eta\) 紧支撑，则
\[
\left.\frac{d\mathcal A}{d\epsilon}\right|_0
=\int\rho\left(\partial_t\eta+\frac{\nabla S}{m}\cdot\nabla\eta\right)\,dt\,dx
=-\int\eta\left[\partial_t\rho+\nabla\cdot(\rho\nabla S/m)\right]dt\,dx.
\] 任意 \(\eta\) 给连续方程。再对 \(\rho\) 变分；对小参数仍保持局部
\(\rho>0\)。Fisher 项的 Euler 导数为 \[
\frac{\hbar^2}{8m}
\left(-2\frac{\Delta\rho}{\rho}
+\frac{|\nabla\rho|^2}{\rho^2}\right)
=-\frac{\hbar^2}{2m}\frac{\Delta\sqrt\rho}{\sqrt\rho}.
\] 令 \(R=\sqrt\rho\)，直接计算 \[
\partial_t\psi=e^{iS/\hbar}(\partial_tR+iR\partial_tS/\hbar),
\] \[
\Delta\psi=e^{iS/\hbar}\left[\Delta R+
\frac{2i}{\hbar}\nabla R\cdot\nabla S+
\frac{iR}{\hbar}\Delta S-\frac{R}{\hbar^2}|\nabla S|^2\right].
\] 代入 Schrödinger 方程，实部除以 \(R>0\) 给 Hamilton--Jacobi
式；虚部乘 \(2R\)
给连续方程。每一步可反向进行，故是所声明区域内的等价。\(\square\)

这里 \(|\psi|^2=\rho\)
是变量映射的一部分。它说明给定统计作用量与波函数动力学如何一致，不单独证明所有可观测量的
POVM Born
规则，也不处理节点、涡旋和全局多值相位。全局波函数须另满足适当线丛/边界条件；闭轨相位整数只限制相位积分。

系数和单位也属于结论的前提：\(S\) 与 \(\hbar\)
都有作用量单位，\(|\nabla S|^2/(2m)\) 为能量，Fisher 项系数
\(\hbar^2/(8m)\) 才使两方程线性合成为所写 Schrödinger
方程。若改变此系数、改变概率密度的含义或增加其他梯度项，应重新推导结果。原论文提出的
Π/Φ
动机没有唯一固定这些选择。{[}S1，§4.2.1，第30--32页；S2，§23；Q2，P-QACT-01{]}

\hypertarget{fisher-ux4e0bux754cux4f55ux65f6ux6210ux4e3aux4f4dux7f6eux52a8ux91cfux4e0dux786eux5b9aux6027}{%
\section{3.2 Fisher
下界何时成为位置---动量不确定性}\label{fisher-ux4e0bux754cux4f55ux65f6ux6210ux4e3aux4f4dux7f6eux52a8ux91cfux4e0dux786eux5b9aux6027}}

设一维归一密度 \(\rho\) 有有限方差和位置 Fisher 信息
\(I_x=\int(\rho')^2/\rho\,dx\)，且边界项消失。

\protect\hypertarget{p-qact-02}{}{}

\textbf{P-QACT-02。} \[
\operatorname{Var}(X)I_x\ge1.
\] 若另使用上述波函数和量子动量 \(P=-i\hbar\partial_x\)，满足所需
Sobolev 定义域，则 \[
\operatorname{Var}(P)=\operatorname{Var}_\rho(S')
+\frac{\hbar^2}{4}I_x,
\quad \Delta X\Delta P\ge\hbar/2.
\]

\textbf{证明。} 分部积分给 \(\int(x-\bar x)\rho'=-1\)，因此 \[
1=\left|\int(x-\bar x)\sqrt\rho\frac{\rho'}{\sqrt\rho}\,dx\right|^2
\le\left[\int(x-\bar x)^2\rho\,dx\right]
\left[\int\frac{(\rho')^2}{\rho}\,dx\right].
\] 对 \(\psi=\sqrt\rho e^{iS/\hbar}\)， \[
P\psi=e^{iS/\hbar}\left(\sqrt\rho\,S'-i\hbar\frac{\rho'}{2\sqrt\rho}\right).
\] 边界项消失使
\(\langle P\rangle=E_\rho S'\)；模平方中实部和虚部不产生交叉项，故
\(\|P\psi\|^2=E_\rho(S')^2+\hbar^2I_x/4\)。作为自伴动量谱的二阶矩，这在
\(\psi\in D(P)\) 且有关积分有限时已有意义，不必先假定 \(P^2\psi\)
存在。相减得到方差分解，再用第一不等式即得
\(\Delta X\Delta P\ge\hbar/2\)。\(\square\)

因此估计论与量子不确定性有明确连接，但动量算子和振幅相位结构是额外桥梁。仅知道有限信息容量不能直接推出
CCR。

一般概率论中可以独立规定 \(X\) 和另一个实随机变量 \(P\)
的分布，使两方差都任意小；Fisher
信息本身不会指定它们是正则共轭变量。上述动量方差恒等式正是把统计不等式变成物理不确定性的桥，不能用``信息有限''替代它。{[}S1，§4.2.1，第30页，印刷行1468--1487{]}

\hypertarget{ux4e8cux6b21ux578bux4e0eux4e24ux76f8ux4f4dux68c0ux9a8cux7684ux4e25ux683cux542bux4e49}{%
\section{3.3
二次型与``两相位检验''的严格含义}\label{ux4e8cux6b21ux578bux4e0eux4e24ux76f8ux4f4dux68c0ux9a8cux7684ux4e25ux683cux542bux4e49}}

\textbf{K-QBORN-01。} 在复向量空间上，连续非负函数 \(F\) 若满足
\(F(\lambda x)=|\lambda|^2F(x)\) 及全部 \(x,y\)
的平行四边形恒等式，则由复极化公式得到唯一 Hermitian 半正定形式
\(B\)，使 \(F(x)=B(x,x)\)。有限维即
\(F(x)=x^\dagger Ax,A\succeq0\)。这是二次型/内积的标准刻画，不是有限个数据点自动证明的全局结论。

按内积对第二变量线性的约定， \[
B(x,y)=\frac14\{F(x+y)-F(x-y)
-iF(x+iy)+iF(x-iy)\}.
\] 若对每个事件 \(a\) 都有这样的二次型
\(F_a(x)=x^\dagger E_ax\)，而对所有向量 \(x\) 都满足
\(\sum_aF_a(x)=\|x\|^2\)，则 \[
x^\dagger\left(\sum_aE_a-I\right)x=0\quad\forall x.
\] 极化唯一性给 \(\sum_aE_a=I\)；各 \(E_a\succeq0\) 又推出
\(E_a\preceq I\)。此时确实得到一套
POVM。若经典随机混合按凸线性运输，\(\rho=\sum_ip_i|\psi_i\rangle\langle\psi_i|\)
的概率为
\(\sum_ip_iF_a(\psi_i)=\operatorname{Tr}(E_a\rho)\)。这段推导显式依赖二次性、复向量结构、总概率及混合规则。

对于一套有限实验，``两个相位''通常指用实/虚干涉坐标恢复一个已声明二次模型，而不是只测两个数就排除任意非二次理论。

有限测试集合外可以改变响应函数而保持全部已测值不变。因此原两角计划的有效结论必须携带模型族、量词和误差；不得从有限通过升级为无条件
Born 唯一性。

\hypertarget{ux65e0ux504fux5206ux652fux6743ux91cdux4e0eux5b9eux9645ux7ed3ux679c}{%
\section{3.4
无偏分支权重与实际结果}\label{ux65e0ux504fux5206ux652fux6743ux91cdux4e0eux5b9eux9645ux7ed3ux679c}}

\protect\hypertarget{p-qborn-02}{}{}

\textbf{P-QBORN-02：无偏最终分支的概率。} 若分支权重 \(p_a(t)\in[0,1]\)
是相对同一过滤族的鞅，且几乎处处趋向某个 one-hot 向量
\(1_{\{A=a\}}\)，则 \[
P(A=a\mid\mathcal F_0)=p_a(0).
\]

\textbf{证明。} 对任意 \(B\in\mathcal F_0\)，鞅性给
\(E[\mathbf1_Bp_a(t)]=E[\mathbf1_Bp_a(0)]\)。有界性允许支配收敛，得到
\(E[\mathbf1_B\mathbf1_{\{A=a\}}]=E[\mathbf1_Bp_a(0)]\)，这就是条件概率恒等式。固定初始制备使
\(p_a(0)\) 为常数时，取期望得到原式 \(P(A=a)=p_a(0)\)；随机初始制备则为
\(P(A=a)=Ep_a(0)\)。\(\square\)

若初权重来自 \(|c_a|^2\)，便得到该 actualization 模型的 Born
结果概率。但``无偏鞅''``最终单分支''``具体物理机制''都是需要独立说明的前提；它们尚未由公共记录客观性证明。

\hypertarget{ch-4}{%
\chapter{4. 有限维量子操作：态、效应、通道与仪器}\label{ch-4}}

\hypertarget{ux57faux7840ux7ea6ux5b9a}{%
\section{4.1 基础约定}\label{ux57faux7840ux7ea6ux5b9a}}

有限维复 Hilbert 空间
\(\mathcal H\cong\mathbb C^d\)，内积对第二变量线性。态为
\(\rho\succeq0,\operatorname{Tr}\rho=1\)，效应为
\(0\preceq E\preceq I\)，概率 \[
p(E|\rho)=\operatorname{Tr}(\rho E).
\] 这个迹确为实数且在 \([0,1]\) 内。因为
\(\rho^{1/2}E\rho^{1/2}\succeq0\)，其迹非负，迹循环给
\(\operatorname{Tr}(\rho E)\ge0\)；把 \(E\) 换成 \(I-E\) 得
\(1-\operatorname{Tr}(\rho E)\ge0\)。这里使用的是已声明的量子状态、效应与配对，不能把单凭公式写出的任意复矩阵称为概率模型。

POVM 是一族 \(\{E_a\}\) 满足
\(\sum_aE_a=I\)。投影测量是互斥正交投影的特例；一般 POVM 不指定测后态。

复合系统使用 \(\mathcal H_A\otimes\mathcal H_B\)，局部态由
\(\operatorname{Tr}(\rho_AE)=\operatorname{Tr}[\rho_{AB}(E\otimes I)]\)
唯一确定。分离态是
\(\sum_i p_i\rho_i^A\otimes\rho_i^B\)，不具此形式者纠缠。

\hypertarget{ux5b8cux5168ux6b63choi-ux4e0e-kraus}{%
\section{4.2 完全正、Choi 与
Kraus}\label{ux5b8cux5168ux6b63choi-ux4e0e-kraus}}

线性映射
\(\mathcal E:\mathcal B(\mathcal H_A)\to\mathcal B(\mathcal H_B)\)
完全正，若所有辅助维数 \(n\) 上 \(\operatorname{id}_n\otimes\mathcal E\)
保持正性。采用未归一 Choi 约定 \[
J(\mathcal E)=\sum_{ij}|i\rangle\langle j|\otimes
\mathcal E(|i\rangle\langle j|),\qquad
\mathcal E(\rho)=\operatorname{Tr}_A[(\rho^\top\otimes I)J].
\]

\protect\hypertarget{p-qop-01}{}{}

\textbf{P-QOP-01：有限 Choi--Kraus 等价。} \[
\mathcal E\ {\rm CP}\iff J(\mathcal E)\succeq0
\iff\mathcal E(\rho)=\sum_\mu K_\mu\rho K_\mu^\dagger.
\] 保迹等价于 \(\sum_\mu K_\mu^\dagger K_\mu=I_A\)，亦等价于
\(\operatorname{Tr}_B J=I_A\)。

\textbf{证明。} CP 作用于未归一最大纠缠投影给正 Choi。正 Choi
的谱分解写作 \(\sum_\mu|K_\mu\rangle\rangle\langle\langle K_\mu|\)，其中
\(|K\rangle\rangle=\sum_i|i\rangle\otimes K|i\rangle\)；逐矩阵元代入重构式得
Kraus 表达。Kraus
表达在任意辅助空间仍为正项之和。保迹条件由迹循环及任意输入成立得到。\(\square\)

具体地，正 Choi 分解可先写
\(J=\sum_\mu\lambda_\mu|v_\mu\rangle\langle v_\mu|\)，将
\(\sqrt{\lambda_\mu}v_\mu\) 按输入索引和输出索引重排成 \(K_\mu\)。随后
\[
\operatorname{Tr}\mathcal E(X)
=\operatorname{Tr}\left[X\sum_\mu K_\mu^\dagger K_\mu\right]
\] 对所有 \(X\) 等于 \(\operatorname{Tr}X\)，当且仅当括号中的算子为
\(I\)。若 Choi
约定改用归一最大纠缠向量或交换输入/输出顺序，上述偏迹和重构式的系数或转置位置必须同步改变。

仅正映射不够：转置保持单系统密度矩阵正性，但其对 Bell
态一侧作用得到具有负特征值的部分转置。因而若允许任意纠缠辅助系统，CP
是物理兼容条件；其依据来自所声明的复合扇区。

\hypertarget{ux4eeaux5668ux548cux5f52ux4e00ux5316}{%
\section{4.3 仪器和归一化}\label{ux4eeaux5668ux548cux5f52ux4e00ux5316}}

Instrument \(\{\mathcal M_a\}\) 的每支 CP、迹非增，总和 CPTP。 \[
p_a=\operatorname{Tr}\mathcal M_a(\rho),\quad
E_a=\mathcal M_a^\dagger(I),\quad
\rho_a=\mathcal M_a(\rho)/p_a\quad(p_a>0).
\] \(p_a=0\) 时不能把除零状态当作已认证后验。Lüders 更新
\(\mathcal M_a(\rho)=P_a\rho P_a\) 是特定 instrument，不由相同 POVM
自动决定。

有限 Kraus 实现写作
\(\mathcal M_o(X)=\sum_kK_{o,k}XK_{o,k}^\dagger\)。全仪器的实际完备条件是
\(\sum_o\sum_kK_{o,k}^\dagger K_{o,k}=I\)，因此
\(E_o=\sum_kK_{o,k}^\dagger K_{o,k}\)。每支求和的索引数可以不同；统一数组存储时加零
Kraus 项不改变通道，但不能凭借漏掉一支或漏掉某个项来满足完备性。

\protect\hypertarget{p-qop-02}{}{}

\textbf{P-QOP-02：无信号。} 局部 CPTP \(\mathcal E_A\) 不改变未条件化的
\(B\) 态。

\textbf{证明。} 对任意 \(B\) 效应 \(F\)，对偶单元性给
\((\mathcal E_A^\dagger\otimes\operatorname{id})(I\otimes F)=I\otimes F\)，故全部局部概率不变。筛选某个测量结果可改变条件态，但接收者若未获得结果标签，看到的仍是总和。\(\square\)

\hypertarget{ux8dddux79bbux548cux8befux5dee}{%
\section{4.4 距离和误差}\label{ux8dddux79bbux548cux8befux5dee}}

迹距离 \(D(\rho,\sigma)=\frac12\|\rho-\sigma\|_1\)。通道的钻石范数采用
\[
\|\Phi\|_\diamond=\sup_{\|X\|_1=1}
\|(\operatorname{id}_d\otimes\Phi)(X)\|_1,
\] 有限输入维 \(d\) 足够；通道区分距离取其一半。

\protect\hypertarget{p-qop-03}{}{}

\textbf{P-QOP-03：操作误差与传播。} 任意效应概率差不超过迹距离，CPTP
不增加态迹距离；若每步通道钻石距离至多
\(\epsilon_t\)，使用相同因果量子控制及相同初态，最终可区分概率误差不超过
\(\sum_t\epsilon_t\)。

\textbf{证明。} 令 \(\Delta=\rho-\sigma=\Delta_+-\Delta_-\)
为正负谱分解。因为 \(\operatorname{Tr}\Delta=0\)，两部分迹相同且为
\(\|\Delta\|_1/2\)。对 \(0\preceq E\preceq I\)， \[
-\operatorname{Tr}\Delta_-\le\operatorname{Tr}(E\Delta)
\le\operatorname{Tr}\Delta_+.
\] 取 \(E\) 为正谱支撑投影时上界达到，故 \[
\sup_{0\preceq E\preceq I}|\operatorname{Tr}[E(\rho-\sigma)]|
=\tfrac12\|\rho-\sigma\|_1.
\] 对 CPTP \(\mathcal T\)，对偶映射正且单元，因此
\(0\preceq\mathcal T^*(E)\preceq I\)。由迹对偶，输出效应的概率差等于输入拉回效应的概率差，再取上确界即得收缩。

对于多步比较，将已有经典记录、控制器和允许的量子记忆都纳入电路；两边使用同一因果控制规则。逐个替换第
\(t\) 个通道形成一串中间电路。每次替换的半钻石范数界为
\(\epsilon_t\)，其前后相同 CPTP
处理及最终效应不会增加区分距离。三角不等式给
\(\min(1,\sum_t\epsilon_t)\) 的最终概率界。\(\square\)

若比较测量仪器，应把结果标签保留在通道中：\(\rho\mapsto\sum_o|o\rangle\langle o|\otimes\mathcal M_o(\rho)\)。仅约束遗忘标签后的总通道，不保证条件反馈电路接近。有限协议族的最大响应差也可以定义操作伪距离，但一般只不大于完整效应上确界；两者不能共用``已经层析完备''的默认前提。

后选择会放大误差。若两未归一后验 \(\sigma,\tau\succeq0\) 满足
\(\|\sigma-\tau\|_1\le\delta\)，且
\(p=\operatorname{Tr}\sigma,q=\operatorname{Tr}\tau\ge p_{\min}>0\)，则
\[
D(\sigma/p,\tau/q)\le\delta/p_{\min}.
\] 证明是加减 \(\tau/p\)，利用
\(|p-q|\le\delta\)。稀有结果不能沿用未条件化的同一误差半径。

本节的迹距离、钻石范数和 Choi
等价是完整有限维数学论述。具体机器内核覆盖哪些命题，应以逐个 Lean
声明及其审计回执为准；一个模型的精确数值检查不会自动形式化这里使用的全部谱论和范数理论。{[}Q2，§4；S4，§2、§9{]}

\hypertarget{ch-5}{%
\chapter{5. Quantum FFIPS 与量子对象实现}\label{ch-5}}

\hypertarget{ux6001ux4f5cux4e3aux5b8cux6574ux672aux6765ux9884ux6d4bux8868ux793a}{%
\section{5.1
态作为完整未来预测表示}\label{ux6001ux4f5cux4e3aux5b8cux6574ux672aux6765ux9884ux6d4bux8868ux793a}}

\protect\hypertarget{p-qff-01}{}{}

\textbf{P-QFF-01：有限量子 FFIPS。} 假定当前 \(\rho_h\)
对全部未来过程充分，且未来协议包含一套分离 Hermitian 算子的效应，则 \[
h\sim h'\iff\rho_h=\rho_{h'}.
\]

\textbf{证明。}
态相同经同一未来仪器序列有相同概率。反向，若全部允许未来概率相同，分离效应族对
\(\Delta=\rho_h-\rho_{h'}\) 的配对均为零，由分离性得
\(\Delta=0\)。在允许全部效应的特例中，可以直接取非零迹零 Hermitian
\(\Delta\) 的正谱支撑投影 \(P_+\)，其 \(\operatorname{Tr}(P_+\Delta)>0\)
给出区分实验。有限分离族无需真的包含每个
\(P_+\)；其张成性已足够。\(\square\)

若实际协议只测 \(Z\)，则 \(|+\rangle,|-\rangle\)
的当前响应相同；若未来允许 Hadamard 后测
\(Z\)，两者立即可分。因此规范态相对于协议，而不是相对于分析者预先喜欢的坐标。

环境记忆或初始系统---环境关联存在时，\(\rho_h\) 未必充分；第 18
章给明确因果断开反例。此时应使用过程响应而非把约化态硬称为全未来状态。

\hypertarget{ux91cfux5b50ux8bfbux51faux4e0eux901aux7528ux8f7dux4f53ux7684ux8fdeux63a5}{%
\section{5.2
量子读出与通用载体的连接}\label{ux91cfux5b50ux8bfbux51faux4e0eux901aux7528ux8f7dux4f53ux7684ux8fdeux63a5}}

固定有限物理读出块 \(V\) 与目标响应 \(C\)。访问集合 \(S\subseteq V\)
提供可实施量子/经典接口；充分性要求能恢复指定全部未来响应。极小充分族和访问
blocker 仍可按 Core
定义，但量子接口应说明是否允许联合测量、辅助系统和相干访问。

逻辑上删除一个输出端口、物理上擦除子系统、测量该系统并遗忘结果，是三个不同操作。它们可能分别保持、限制或破坏相干性，不能共用一条未注明语义的毁损结论。

\hypertarget{ux91cfux5b50ux52a8ux529bux5b66ux5546}{%
\section{5.3 量子动力学商}\label{ux91cfux5b50ux52a8ux529bux5b66ux5546}}

给逻辑与物理系统的编码 \(\mathcal E:M\to X\)、解码
\(\mathcal D:X\to M\)，均 CPTP，并在声明逻辑态类上满足
\(\mathcal D\mathcal E=\operatorname{id}\)。对每个允许动作 \(a\) 和结果
\(o\)，给物理、逻辑仪器分支 \(\mathcal M_{a,o}\) 与
\(\overline{\mathcal M}_{a,o}\)。要求逐分支交换： \[
\mathcal M_{a,o}\mathcal E=\mathcal E\overline{\mathcal M}_{a,o}.
\]
它应在完整逻辑算子空间上成立，或至少在包含全部可达分支输出的闭合类上成立。若最终还测试效应
\(F\)，其逻辑效应须为
\(\bar F=\mathcal E^*(F)\)，或在所声明输入类上具有相同配对。

\protect\hypertarget{p-qff-02}{}{}

\textbf{P-QFF-02：量子商的路径保持。}
在上述逐分支交换条件下，若物理初态为
\(\mathcal E(\bar\rho)\)，则任意允许有限记录的未归一分支态由编码精确对应，因而记录概率相同；对应终端效应的联合概率也相同。

\textbf{证明。} 对记录
\(w=((a_1,o_1),\ldots,(a_t,o_t))\)，连续代入逐分支等式，得到 \[
\mathcal M_{a_t,o_t}\cdots\mathcal M_{a_1,o_1}\mathcal E(\bar\rho)
=\mathcal E\bigl(\overline{\mathcal M}_{a_t,o_t}\cdots
\overline{\mathcal M}_{a_1,o_1}(\bar\rho)\bigr).
\] 编码保迹，故两边迹相同；加入终端效应后，用对偶式
\(\operatorname{Tr}[F\mathcal E(X)]=\operatorname{Tr}[\bar F X]\)。自适应协议在一条已见记录路径上选择相同动作，逐分支等式仍成立；对有限叶集求和就得到任意叶事件的概率相同。零概率分支保留零未归一算子，无需除法。\(\square\)

逐分支要求不能削弱成``总通道相同且单次 POVM 相同''。在一个比特上，令 \[
\mathcal M_0(X)=\tfrac12\operatorname{Tr}(X)|0\rangle\langle0|,
\quad\mathcal M_1(X)=\tfrac12\operatorname{Tr}(X)|1\rangle\langle1|,
\] 而
\(\mathcal N_0=\mathcal M_1,\ \mathcal N_1=\mathcal M_0\)。两仪器都
CP，所有分支效应均为 \(I/2\)，总通道同为
\(X\mapsto\operatorname{Tr}(X)I/2\)。然而紧接一次 \(Z\)
读出时，\(\mathcal M\) 的两次标签总是相同，\(\mathcal N\)
的两次标签总是相反。这个精确反例说明一步边缘和无条件演化不会独自保留多时刻预测。{[}Q2，§5.3，修订前提；S2，§7.5、§28.3--4{]}

在编码像上成立的商不覆盖任意物理初态；泄漏、编码误差和辅助系统完整性须另计。近似交换的误差应采用受限或完整钻石范数，明确其输入类。

\hypertarget{ux5bf9ux8c61ux6301ux7eedux4e0eux91cfux5b50ux5b58ux5728}{%
\section{5.4
对象持续与量子存在}\label{ux5bf9ux8c61ux6301ux7eedux4e0eux91cfux5b50ux5b58ux5728}}

可用的持续结构包括不变子空间、噪声无扰子系统、稳态、亚稳态、可纠错逻辑态、拓扑保护扇区或有效准粒子。每一种都有自己的机制条件。一般波函数不因``不呈现经典闭环''而失去对象资格；一般稳定谱隙也不自动意味着拓扑保护。

第 20
章在额外任务下赋予量子对象控制价值。这里的对象证书本身不指定唯一收益函数。

\hypertarget{ux4eceux6709ux9650ux95edux5408ux68c0ux67e5ux5230ux5168ux90e8ux6709ux9650ux672aux6765}{%
\section{5.5
从有限闭合检查到全部有限未来}\label{ux4eceux6709ux9650ux95edux5408ux68c0ux67e5ux5230ux5168ux90e8ux6709ux9650ux672aux6765}}

有限量子系统不必为每个未来长度重新枚举一次证明。关键是证明由真实仪器生成的效应空间闭合，而不是把一个有限长度的测试表称为``完整未来''。固定有限操作库，令事件为
\(e=(a,o)\)，其 Heisenberg 拉回为 \[
\mathcal M_e^*(E)=\sum_kK_{e,k}^\dagger E K_{e,k}.
\] 对按发生顺序排列的记录定义 \[
E_{\varnothing}=I,\qquad E_{e::w}=\mathcal M_e^*(E_w).
\] 迹循环给
\(\operatorname{Tr}[E\mathcal M_e(X)]=\operatorname{Tr}[\mathcal M_e^*(E)X]\)，于是对词长归纳得到
\[
p(w\mid\rho)=\operatorname{Tr}(E_w\rho).
\] 定义实线性空间
\(\mathcal O=\operatorname{span}_{\mathbb R}\{E_w\}\)。所有 \(E_w\) 都是
Hermitian，因此在 \(d\) 维量子系统中
\(\dim\mathcal O\le d^2\)。这个维数上界来自 Hermitian
矩阵的实坐标数，并不意味着全部态可由任意少量测量恢复。

给一份有限候选列表 \(B_1,\ldots,B_r\)，需要检查三类有限条件：其中一个
\(B_{j_0}=I\)；每个 \(B_j\) 都是真实词效应
\(E_{w_j}\)；对每个允许事件和每个候选元素，存在实系数 \[
\mathcal M_e^*(B_j)=\sum_iT_{e,ji}B_i.
\]
前两类保证候选空间没有放入与真实操作库无关的自由算子；第三类保证一步拉回闭合。任意
\(E_w\)
属于候选实线性空间可由词长归纳证明：空词使用单位元素；前缀步骤使用拉回线性性和已检查的每个局部等式。因此候选空间恰等于
\(\mathcal O\)。

定义未归一坐标 \(x_j(X)=\operatorname{Tr}(B_jX)\)，则 \[
x_j(\mathcal M_e(X))=\sum_iT_{e,ji}x_i(X),\qquad
x(\mathcal M_e(X))=T_ex(X).
\] 对于 \(w=(e_1,\ldots,e_t)\)，先执行 \(e_1\)，所以 \[
p(w\mid\rho)=\bigl[T_{e_t}\cdots T_{e_1}x(\rho)\bigr]_{j_0}.
\]
两个态的有限坐标相同，当且仅当所有有限词概率相同：正向由闭合和递推，反向因为每个
\(B_j\)
自己就是一个词效应。候选列表可以有线性冗余；这不损害充分性，却意味着不能仅凭该证书声称坐标已最小或独立。

相同结果也覆盖完整有限自适应树。每片叶有一条按历史选择动作的真实记录词，其概率由同一初始坐标计算；各叶概率相加得到记录事件。正概率叶的条件坐标再除以叶概率，零概率叶保持未归一语义。这里没有把后验闭合或全词等价作为输入，而是从局部拉回恒等式推出它们。

UEOT-QM 的 \texttt{Observables.\allowbreak lean}、\allowbreak \texttt{Finite\allowbreak Coordinates.\allowbreak lean} 与
\texttt{Adaptive.\allowbreak lean} 实现这条真实
Kraus---迹配对链。\texttt{Exact.\allowbreak lean}
把一对有理实部/虚部矩阵可靠解释为复矩阵，证明加法、乘法、伴随、实有理缩放和有限求和被保持；\texttt{Exact\allowbreak Certificate.\allowbreak lean}
进一步检查非负权重、全部仪器完备性、空词、每个基矩阵的真实词见证和局部闭合，再构造普通复数量子仪器及
\texttt{Finite\allowbreak Observable\allowbreak Presentation}。它检查的有限数据可以由
\texttt{decide\ +kernel}
证明为真，随后用通用可靠性定理提升到任意有限记录。{[}源接口：\texttt{UEOTQM.\allowbreak record\allowbreak Effect}、\allowbreak \texttt{observable\allowbreak Space}、\allowbreak \texttt{response\_\allowbreak eq\_\allowbreak iff\_\allowbreak pairing\_\allowbreak observable\allowbreak Space}；\texttt{Exact\allowbreak Certificate.\allowbreak Conditions}、\allowbreak \texttt{exact\allowbreak Record\allowbreak Effect\_\allowbreak to\allowbreak Matrix}、\allowbreak \texttt{coordinates\_\allowbreak eq\_\allowbreak iff\_\allowbreak response\_\allowbreak eq}、\allowbreak \texttt{adaptive\_\allowbreak leaf\allowbreak Probability\_\allowbreak eq}{]}

这些数学保证对应被绑定到源代码的那个有限模型。它们不自动证明外部 JSON
解析器、搜索算法和所有导出器实现正确，也不证明输入数值由实际设备实现；更没有从有限量子响应空间推出粒子世代、质量谱或完整
QFT。源绑定和负控解决的是``究竟编译了哪份数据''，实验与领域桥解决的是``这份数据怎样对应物理''。

\hypertarget{ch-6}{%
\chapter{6. 测量：相干父系统、公共商与条件预测态}\label{ch-6}}

\hypertarget{ux4efbux610fux6709ux9650ux4eeaux5668ux7684ux76f8ux5e72ux5b8cux6210}{%
\section{6.1
任意有限仪器的相干完成}\label{ux4efbux610fux6709ux9650ux4eeaux5668ux7684ux76f8ux5e72ux5b8cux6210}}

令
\(\mathcal M_a(\rho)=\sum_\mu A_{a\mu}\rho A_{a\mu}^\dagger\)，总和保迹。定义
\[
V|\psi\rangle=\sum_{a,\mu}A_{a\mu}|\psi\rangle\otimes|a,\mu\rangle,\qquad
Q_a=\sum_\mu|a,\mu\rangle\langle a,\mu|.
\]

\protect\hypertarget{p-qmeas-01}{}{}

\textbf{P-QMEAS-01：父系统完成。} \(V^\dagger V=I\)，且 \[
\mathcal M_a(\rho)=
\operatorname{Tr}_R[(I\otimes Q_a)V\rho V^\dagger(I\otimes Q_a)].
\] 有限扩维后 \(V\) 可以由酉操作和固定辅助初态实现。

\textbf{证明。} 辅助基正交使
\(V^\dagger V=\sum A_{a\mu}^\dagger A_{a\mu}=I\)。投影选择标签
\(a\)，偏迹只留下相同 \(\mu\)
项，得到仪器。将等距映射的像基补成全空间正交基，给酉延拓。\(\square\)

纯输入对应纯联合父态。父态是否在实验中可逆保持，取决于环境访问和噪声，不能只凭数学
dilation 宣称实际可完成逆操作。

\hypertarget{ux6761ux4ef6ux5316ux5c31ux662fux7ed3ux679cux6269ux5c55ux540eux7684ux9884ux6d4bux6001}{%
\section{6.2
条件化就是结果扩展后的预测态}\label{ux6761ux4ef6ux5316ux5c31ux662fux7ed3ux679cux6269ux5c55ux540eux7684ux9884ux6d4bux6001}}

\protect\hypertarget{p-qmeas-02}{}{}

\textbf{P-QMEAS-02。} 在结果 \(a\) 已发生且 \(p_a>0\)
后，任意未来量子效应 \(F\) 的条件概率为 \[
P(F|h,a)=\operatorname{Tr}(F\rho_a).
\] 若未来协议分离，\(\rho_a\) 是该结果扩展历史的唯一量子 FFIPS。

\textbf{证明。} 联合概率为
\(\operatorname{Tr}[F\mathcal M_a(\rho)]\)，除以 \(p_a\)；唯一性调用
P-QFF-01。\(\square\)

所以标准测量链严格表达为相干父完成、记录代数访问、结果条件化。这与标准
QM 的 instrument 预测一致，是 UEOT 的组织连接，不是另加坍缩动力学。

\hypertarget{ux516cux5171ux8bb0ux5f55ux770bux4e0dux89c1ux4ec0ux4e48}{%
\section{6.3
公共记录看不见什么}\label{ux516cux5171ux8bb0ux5f55ux770bux4e0dux89c1ux4ec0ux4e48}}

设 \(\Delta_R(\Omega)=\sum_a(I\otimes Q_a)\Omega(I\otimes Q_a)\)。

\protect\hypertarget{p-qmeas-03}{}{}

\textbf{P-QMEAS-03：公共盲性。} 若 \([B,I\otimes Q_a]=0\) 对全部 \(a\)
成立，则 \[
\operatorname{Tr}(B\Omega)=\operatorname{Tr}[B\Delta_R(\Omega)].
\]

\textbf{证明。} \(B\)
是记录块对角，展开投影分解并用迹循环，所有不同标签交叉项为零。\(\square\)

因此相干父态与其去相干态对公共记录具有相同统计。冗余记录、稳定指针或观察者一致，不能单独判定全局相干是否基本消失。区分必须加入跨记录块的相干测试，例如完整擦除/反演。

这里的盲性适用于所有记录块对角效应。块对角算子代数
\(\bigoplus_a\mathcal B(\mathcal H_a)\)
在分支内部维数大于一时仍可非交换，不能将它整体称为经典代数。下一节只取由标签投影张成的
\(\mathcal Z=\operatorname{span}\{P_a\}\)；这个较小代数才自动交换。公共标签经典性和分支内部保留量子信息可以同时成立。

\hypertarget{ux4ea4ux6362ux4ee3ux6570ux4e0aux7684ux7ecfux5178ux8fc7ux7a0b}{%
\section{6.4
交换代数上的经典过程}\label{ux4ea4ux6362ux4ee3ux6570ux4e0aux7684ux7ecfux5178ux8fc7ux7a0b}}

公共投影 \(P_a\) 张成交换代数 \(\mathcal Z\)。若
\(\mathcal T^\dagger(\mathcal Z)\subseteq\mathcal Z\)，则 \[
\mathcal T^\dagger(P_b)=\sum_a M_{ab}P_a,\quad
M_{ab}\ge0,\quad\sum_bM_{ab}=1.
\] 于是公共概率行向量按 \(p'=pM\) 演化。

\protect\hypertarget{p-qmeas-04}{}{}

\textbf{P-QMEAS-04。} 上述 \(M\) 是经典随机矩阵。

\textbf{证明。}
正性给非负展开系数，单元性给每行和一，迹配对给概率更新。\(\square\)

这将量子公共经典性与 Core
的经典预测商直接接通，不要求整个父系统变成经典概率混合。

当要求多时刻经典记录过程时，还需相应测量实施或各记录分支对该代数保持适当闭合。仅有一个无条件
\(\mathcal T^*\) 的不变代数，直接保证的是它反复作用下公共概率的 Markov
更新；涉及其他中间仪器或反馈，须按§5.3逐分支核对。

\hypertarget{ux5ba2ux89c2ux8bb0ux5f55ux8bc1ux4e66-orc}{%
\section{6.5 客观记录证书
ORC}\label{ux5ba2ux89c2ux8bb0ux5f55ux8bc1ux4e66-orc}}

预先指定标签、误差 \(\delta\)、持续窗 \(T\)、记录数量 \(m\)
和置信度。证书至少要求：公共代数交换；\(m\)
个不交片段分别能识别标签；记录有持续与历史运输；到已定义失败类有正裕量；全部统计条件仍有置信余量。

ORC
是操作定义。观察到环境片段互信息平台是相关证据，但不能自动代替片段可独立访问、不可扰动读取和完整失败模型。它也不蕴含下面额外
actualization 通道的强度。

\hypertarget{ch-7}{%
\chapter{7. 强记录无扰动力学的完整分类}\label{ch-7}}

后续接口：第38章给循环相位界，第39章给四能级噪声类见证，第46章给相关试次的检验。

\hypertarget{srn-ux5b9aux4e49ux4e0e-schur-ux5f62ux5f0f}{%
\section{7.1 SRN 定义与 Schur
形式}\label{srn-ux5b9aux4e49ux4e0e-schur-ux5f62ux5f0f}}

令有限维
\(\mathcal H=\bigoplus_{a=1}^m\mathcal H_a\)，每个分支非零，\(\sum_aP_a=I\)。CPTP
通道满足
\textbf{SRN}，若它逐点固定每一分支内部的全部密度矩阵，而不只是分支总概率。零维分支应先删除，避免为不可观测块赋予虚假的独立参数。

\protect\hypertarget{p-qsrn-01}{}{}

\textbf{P-QSRN-01：SRN 完备分类。} SRN 当且仅当 \[
\mathcal A_K(X)=\sum_{a,b}K_{ab}P_aXP_b,\qquad
K\succeq0,\quad K_{aa}=1.
\]

\textbf{证明。} 取 Kraus \(L_\mu\)。分支内每个纯态投影被固定，正项
\(L_\mu|\psi\rangle\langle\psi|L_\mu^\dagger\) 的和秩一，故每个
\(L_\mu|\psi\rangle\) 都平行于
\(|\psi\rangle\)。对同一分支两独立向量及其和使用线性性，比例系数必须相同，即
\(L_\mu P_a=c_{\mu a}P_a\)。所以
\(K_{ab}=\sum_\mu c_{\mu a}\overline{c_{\mu b}}\) 为 Gram
矩阵，保迹给对角一。反向，从任意这样的 Gram 分解取
\(L_\mu=\sum_ac_{\mu a}P_a\)，即得 CPTP 和分支内逐点固定。\(\square\)

此算子结构与 Schur 通道、无扰代数及 genuinely incoherent operations
的既有理论一致。本文重证它来保证测量扩展的参数化完备，不称为 UEOT
首创算子定理。

\hypertarget{ux590dux5408ux76f2ux6027ux4e0eux4e0dux786eux5b9aux6027}{%
\section{7.2
复合、盲性与不确定性}\label{ux590dux5408ux76f2ux6027ux4e0eux4e0dux786eux5b9aux6027}}

\protect\hypertarget{p-qsrn-02}{}{}

\textbf{P-QSRN-02。} \[
\mathcal A_K\mathcal A_L=\mathcal A_{K\odot L}.
\] 对全部块对角效应，\(\mathcal A_K\) 与恒等通道统计相同。

\textbf{证明。} 每个 \((a,b)\) 块先后乘以
\(L_{ab},K_{ab}\)；公共效应只读取 \(a=b\) 的块，而其系数一。\(\square\)

全一矩阵 \(J\) 表示恒等通道，单位矩阵 \(I\)
表示完全分支去相干；两者有同样公共记录。因此仅由 ORC 无法推出唯一
\(K=F(\mathrm{ORC})\)。如果预先给定标准完整相干动力学，则已经选了相应通道，不能再称为由客观性独立推出。

\hypertarget{ux8fdeux7eedux534aux7fa4ux7684ux5168ux5c40ux751fux6210ux5143}{%
\section{7.3
连续半群的全局生成元}\label{ux8fdeux7eedux534aux7fa4ux7684ux5168ux5c40ux751fux6210ux5143}}

连续 \(K(t+s)=K(t)\odot K(s),K(0)=J\) 逐项给
\(K_{ab}(t)=e^{t\lambda_{ab}}\)。有限时零值会通过开方递推与零点连续性矛盾；非零连续乘法函数取局部对数，得到指数。

\protect\hypertarget{p-qsrn-03}{}{}

\textbf{P-QSRN-03：生成元锥的必要充分条件。} 全部 \(t\ge0\) 的上述映射为
CPTP，当且仅当 \[
\Lambda=\Lambda^\dagger,\quad\lambda_{aa}=0,\quad
z^\dagger\Lambda z\ge0\quad(\mathbf1^\dagger z=0).
\]

\textbf{证明。} 必要性：\(K(t)=J+t\Lambda+o(t)\)，在零和子空间消去
\(J\)，由正性取右导数。充分性可直接构造，不必将 Schoenberg
定理当黑箱。固定锚 \(r\)，令 \[
G_{ab}=\lambda_{ab}-\lambda_{ar}-\lambda_{rb}.
\] 对任意向量 \(x\)，\(x^\dagger Gx\) 等于 \(z^\dagger\Lambda z\)，其中
\(z=x-(\sum_ax_a)e_r\)，故 \(G\succeq0\)。取
\(\theta_a=\operatorname{Im}\lambda_{ar}\)，则 \[
\lambda_{ab}=G_{ab}-\tfrac12G_{aa}-\tfrac12G_{bb}
+i(\theta_a-\theta_b).
\] Schur 乘积保持正性，所以
\(\exp_\odot(tG)=J+tG+t^2G\odot G/2+\cdots\succeq0\)。 再左右乘对角矩阵
\(D_{aa}=e^{-tG_{aa}/2+it\theta_a}\) 与其共轭，恰得到
\(K(t)\)。对角为一，故由 P-QSRN-01 为 CPTP；半群律逐项成立。\(\square\)

这是对全局相容性的显式证明。仅要求每对衰减率非负不够，例如 \[
\Lambda=\begin{pmatrix}0&-1&-1\\-1&0&-10\\-1&-10&0\end{pmatrix}
\] 对 \(z=(2,-1,-1)\) 有 \(z^\top\Lambda z=-12\)，因此小时间通道不物理。

实对称情形令 \(D_{ab}=-\lambda_{ab}\)，上述 Gram 构造可给
\(D_{ab}=\|v_a-v_b\|^2\)。一般复生成元还包含 Gram
的复几何，不应只留下任意相位差。

具体地，若 \(G_{ab}=\langle u_a,u_b\rangle\) 为实 Gram 表示，则
\(D_{ab}=\|u_a-u_b\|^2/2\)；取 \(v_a=u_a/\sqrt2\)
得到所写距离形式。复情形的虚部循环相位受同一个正 Gram
结构约束，并非所有成对非负衰减率加上任意相位都能产生 CP
半群。上述分类针对固定分支分解、有限维、连续且时间齐次的半群；记录随时间改变或存在环境记忆时，应重新写模型。

\hypertarget{ch-8}{%
\chapter{8. QDRR：记录写入---擦除反演与有限数据}\label{ch-8}}

\hypertarget{ux5b8cux6574ux53cdux6f14ux7684ux63a5ux53e3}{%
\section{8.1
完整反演的接口}\label{ux5b8cux6574ux53cdux6f14ux7684ux63a5ux53e3}}

区分两种类型：\(U_W\) 是在系统与记录寄存器上的完整酉操作，寄存器以固定
\(|0\rangle_R\) 初始化；真正的写入等距映射为 \[
W:\mathcal H\to\mathcal H\otimes\mathcal R,\qquad
W|\psi\rangle=U_W(|\psi\rangle\otimes|0\rangle_R).
\] 其像中的分支投影为 \(\widetilde P_a=WP_aW^\dagger\)；它们之和是
\(WW^\dagger\)，不是自动等于整个扩展空间的单位算子。中间通道在这些正交像分支上为
SRN
\(K\)，不把输入泄漏到像外；若将通道定义于整个空间，须显式补全正交余空间的行动。

\protect\hypertarget{p-qqdr-01}{}{}

\textbf{P-QQDR-01：相干拉回。} \[
W^\dagger\widetilde{\mathcal A}_K(WXW^\dagger)W=\mathcal A_K(X).
\]

相应完整物理逆操作满足 \[
U_W^\dagger\widetilde{\mathcal A}_K
\bigl[U_W(X\otimes|0\rangle\langle0|)U_W^\dagger\bigr]U_W
=\mathcal A_K(X)\otimes|0\rangle\langle0|.
\]

\textbf{证明。} 将 \(X\) 展成 \(\sum_{a,b}P_aXP_b\)，由
\(W^\dagger W=I\)， \[
\widetilde P_a(WXW^\dagger)\widetilde P_b=WP_aXP_bW^\dagger.
\] 中间 SRN 将该块乘 \(K_{ab}\)，故其输出为
\(W\mathcal A_K(X)W^\dagger\)。左右乘 \(W^\dagger,W\) 给第一式；再使用
\(U_W^\dagger W|\psi\rangle=|\psi\rangle\otimes|0\rangle\)
给第二式。\(\square\)

反演必须覆盖实际保留 which-branch
信息的全部相关环境。只逆转部分寄存器的残余退相干，不能直接归因于基本坍缩。

\hypertarget{ux5168-srn-ux7c7bux7684ux5c42ux6790}{%
\section{8.2 全 SRN
类的层析}\label{ux5168-srn-ux7c7bux7684ux5c42ux6790}}

准备 \((|a\rangle+e^{i\phi}|b\rangle)/\sqrt2\)，测
\(X_{ab}=|a\rangle\langle b|+|b\rangle\langle a|\)。

\protect\hypertarget{p-qqdr-02}{}{}

\textbf{P-QQDR-02。} \[
\langle X_{ab}\rangle_\phi
=\operatorname{Re}(e^{-i\phi}K_{ab}).
\] \(\phi=0,\pi/2\) 分别给实部、虚部。全部 \(a<b\) 共 \(m(m-1)\)
个实坐标恢复整个 SRN 通道。

\textbf{证明。} 输入 \((a,b)\) 元为 \(e^{-i\phi}/2\)，输出乘以
\(K_{ab}\)；两个共轭项相加即得。对角已知为一，故恢复唯一。\(\square\)

这是对单次有限维 CPTP-SRN
类的完备性，不覆盖人口转移、泄漏、分支内部扰动或任意历史依赖的新物理。

\hypertarget{ux6709ux9650ux91c7ux6837ux4e0eux7269ux7406ux6295ux5f71}{%
\section{8.3
有限采样与物理投影}\label{ux6709ux9650ux91c7ux6837ux4e0eux7269ux7406ux6295ux5f71}}

设 \(m\ge2\)，每设置取得 \(n\ge1\) 个独立、具有该设置固定期望的
\([-1,1]\)
输出样本，\(L=m(m-1)\)，\(0<\alpha<1\)。并合界不要求不同设置之间独立；每个设置内部使用
Hoeffding 时的抽样假设必须满足。\(m=1\) 时没有非对角参数，不能把 \(L=0\)
代入下述对数公式。

\protect\hypertarget{p-qqdr-03}{}{}

\textbf{P-QQDR-03：同时误差。} 以至少 \(1-\alpha\) 的概率， \[
\max_\ell|\widehat\mu_\ell-\mu_\ell|
\le\epsilon_n=\sqrt{\frac{2\log(2L/\alpha)}n},
\] 并有 \[
\|\widehat K-K\|_{\rm op}\le\|\widehat K-K\|_F
\le\epsilon_n\sqrt{2m(m-1)}.
\]

\textbf{证明。} Hoeffding 对宽度二的变量给 \(2e^{-n\epsilon^2/2}\)，对
\(L\) 个设置并合。每个复非对角误差平方至多 \(2\epsilon^2\)，Hermitian
对称再计两次。\(\square\)

取 \(m=3,\epsilon=.12,\alpha=.05\) 需每设置 \(n\ge762\)。不能使用
\([0,1]\) 变量的系数而未换算测量输出。

经验矩阵未必半正定。最近相关矩阵投影可提高物理可用性；投影残差应作为模型一致性诊断保留，不能删除后只汇报``投影后通过''。欧氏凸投影到含真实
\(K\) 的集合不增加 Frobenius 真误差，但它不能把错误 SRN
模型变成经验证的真实模型。

\hypertarget{ux786cux4ef6ux4e0eux65b0ux7269ux7406ux7684ux53efux8bc6ux522bux6027}{%
\section{8.4
硬件与新物理的可识别性}\label{ux786cux4ef6ux4e0eux65b0ux7269ux7406ux7684ux53efux8bc6ux522bux6027}}

只有硬件噪声与候选 actualization 都是独立 SRN 时，才有
\(K_{\rm obs}=K_{\rm hw}\odot K_{\rm act}\)。 逐项相除在 \(K_{\rm hw}\)
接近零时不稳定；一般噪声不可这样分解。

实验至少需要匹配的无记录、立即反演、冗余打乱、物理比特置换、随机交错和独立校准。原量子
Darwinism
数据通常只测试公共记录；缺少相干反演和匹配校准时，不能事后把记录冗余解释为基本坍缩证据。本文没有执行新的硬件实验。

\hypertarget{ch-9}{%
\chapter{9. 量子信息、纠错与计量接口}\label{ch-9}}

\hypertarget{ux7ea0ux7f20ux4e0eux4fe1ux606f}{%
\section{9.1 纠缠与信息}\label{ux7ea0ux7f20ux4e0eux4fe1ux606f}}

von Neumann 熵 \(S(\rho)=-\operatorname{Tr}\rho\log\rho\)，量子相对熵 \[
D(\rho\|\sigma)=\operatorname{Tr}\rho(\log\rho-\log\sigma)
\] 在支撑不包含时为无穷。\textbf{K-QINFO-01}：CPTP
下相对熵数据处理、强次可加性及 Holevo 可访问信息界采用标准量子信息定理。

纯父态 Schmidt 系数 \(p_i\) 给
\(S(\rho_A)=S(\rho_B)\)。局部混合可由纠缠产生，但不是每种局部混合都证明其实际纯化环境已经被实验确定。

量子互信息 \(I(A:B)=S(A)+S(B)-S(AB)\)
衡量总相关。纠缠、共同经典记录、可访问互信息和父对象切割不可约性回答不同问题。

\hypertarget{ux65e0ux514bux9686}{%
\section{9.2 无克隆}\label{ux65e0ux514bux9686}}

\protect\hypertarget{p-qinfo-01}{}{}

\textbf{P-QINFO-01。} 不存在对所有未知纯态实施
\(|\psi\rangle|0\rangle\mapsto|\psi\rangle|\psi\rangle\) 的同一酉操作。

\textbf{证明。} 对两态内积 \(c\)，酉保持内积要求 \(c=c^2\)，与任意
\(0<|c|<1\) 不相容。若允许最终环境依赖输入，纯态复制要求联合输出为
\(|\psi\rangle|\psi\rangle|e_\psi\rangle\)，则内积保持给
\(c=c^2\langle e_\psi|e_\varphi\rangle\)。取模和 Cauchy--Schwarz 得
\(|c|\le|c|^2\)，仍与 \(0<|c|<1\) 矛盾。一般确定 CPTP
复制可通过等距扩张归于同一论证。\(\square\)

正交经典记录可被冗余复制，与无克隆不矛盾。ORC
复制的是可区分公共标签，不是任意未知量子态。

\hypertarget{ux7ea0ux9519ux6761ux4ef6ux4e0eux5bf9ux8c61ux7ef4ux6301}{%
\section{9.3
纠错条件与对象维持}\label{ux7ea0ux9519ux6761ux4ef6ux4e0eux5bf9ux8c61ux7ef4ux6301}}

\textbf{K-QEC-01：Knill--Laflamme。} 编码投影 \(P\) 上，一组噪声 Kraus
\(E_a\) 存在共同完美恢复，当且仅当 \[
PE_a^\dagger E_bP=c_{ab}P.
\] 这是完整辅助系统相干性要求，不只是两个基态分别可恢复。

其充分性的构造是对正矩阵 \(c\)
对角化，得到在代码上互相正交的错误像；测量错误综合征后逆转各等距错误，不读取逻辑信息。必要性来自可恢复通道的互补环境不能区分逻辑态。严格通用证明作为标准
K 调用。

在 UEOT
中，代码子空间是候选实现，恢复动作是维持协议，逻辑未来是预测目标。码距和阈值取决于噪声、几何和可实现门，不由通用
Ω 直接数值确定。

\hypertarget{ux91cfux5b50-fisher-ux4e0eux8ba1ux91cf}{%
\section{9.4 量子 Fisher
与计量}\label{ux91cfux5b50-fisher-ux4e0eux8ba1ux91cf}}

参数态 \(\rho_\theta\) 的 SLD 满足
\(\partial_\theta\rho=(L\rho+\rho L)/2\)，量子 Fisher 为
\(F_Q=\operatorname{Tr}\rho L^2\)。给定 POVM 的经典 Fisher 不超过
\(F_Q\)；局部无偏估计满足相应 Cramér--Rao 界，需支持和正则性条件。

对纯酉族 \(|\psi_\theta\rangle=e^{-i\theta G}|\psi\rangle\)：

\protect\hypertarget{p-qinfo-02}{}{}

\textbf{P-QINFO-02。} \[
F_Q=4(\langle G^2\rangle-\langle G\rangle^2).
\]

\textbf{证明。} 纯态切向投影公式
\(F_Q=4(\langle\dot\psi|\dot\psi\rangle- |\langle\psi|\dot\psi\rangle|^2)\)，代入
\(|\dot\psi\rangle=-iG|\psi\rangle\)。该公式也可直接由 SLD
在态支撑及其正交切向子空间上求解得到。\(\square\)

计量最优性不是所有任务的 GOD。多参数最优测量未必相容，单参数 Fisher
界不能逐项同时声称可达到。

\hypertarget{ch-10}{%
\chapter{10. 无限维量子理论：定义域与正常态}\label{ch-10}}

\hypertarget{ux65e0ux9650-hilbert-ux7ef4ux4e0dux7b49ux4e8eux65e0ux9650ux81eaux7531ux5ea6}{%
\section{10.1 无限 Hilbert
维不等于无限自由度}\label{ux65e0ux9650-hilbert-ux7ef4ux4e0dux7b49ux4e8eux65e0ux9650ux81eaux7531ux5ea6}}

一个谐振子的 \(\mathcal H=L^2(\mathbb R)\)
已无限维，但只有一对正则变量。量子场具有无限多个模式，不能直接使用有限自由度的
Stone--von Neumann 唯一性。有限维操作重建也不自动唯一指定无限维扩展。

采用可分复 Hilbert 空间。\textbf{H-QINF}
限定正常操作扇区：\(\mathcal B(\mathcal H)\)
上的正常态由正迹类、迹一密度算子表示；Schrödinger
图景的操作作用于迹类，对偶图景采用正常的相应 CP
映射。一般算子代数还可有奇异态，不将它们悄悄当作密度矩阵。此处的无限维结构是新增论域，不由有限维重构和有限
\texttt{Fin\ d} 矩阵证书直接提供。

有界效应仍满足 Born 配对；POVM
是弱可数可加的正算子测度。连续位置的概率是 \(\int_B|\psi(x)|^2dx\)，广义
\(|x\rangle\) 不是 Hilbert
中单位向量，单点零概率不能按离散后验直接归一化。

\hypertarget{ux81eaux4f34ux6027ux662fux7269ux7406ux8f93ux5165ux7684ux4e00ux90e8ux5206}{%
\section{10.2
自伴性是物理输入的一部分}\label{ux81eaux4f34ux6027ux662fux7269ux7406ux8f93ux5165ux7684ux4e00ux90e8ux5206}}

无界算子由作用表达式及稠密定义域共同指定。对称 \(A\subseteq A^\dagger\)
不等于自伴 \(A=A^\dagger\)。

\textbf{K-QINF-01：谱定理。} 自伴 \(A\) 对应唯一 PVM \(E_A\)， \[
A=\int\lambda\,dE_A(\lambda),\quad
D(A)=\{\psi:\int\lambda^2d\langle\psi,E_A(\lambda)\psi\rangle<\infty\}.
\] 测量概率始终先经 PVM；均值、方差还需相应谱矩有限。

在 \(L^2(\mathbb R)\)，位置定义域为 \(x\psi\in L^2\)，动量为
\(-i\hbar\partial_x\) 配 Sobolev 域
\(H^1(\mathbb R)\)。区间上同一导数表达式配
\(\psi(L)=e^{i\vartheta}\psi(0)\)
可给不同自伴扩展和动量谱；边界相位属于模型，不由符号 \(-i\partial_x\)
决定。

\textbf{K-QINF-02。} 缺陷指标 \(n_\pm=\dim\ker(A^\dagger\mp i)\)
相等时，闭对称算子具有自伴扩展；半有界情形可用 Friedrichs
扩展。Kato--Rellich 要求扰动对已自伴 \(A\)
相对有界且相对界小于一；闭半有界二次型也可通过表示定理生成自伴
Hamiltonian。这些是选择合法 Hamiltonian 的标准接口。

\hypertarget{stone-ux4e0eux5f3a-schruxf6dinger-ux65b9ux7a0b}{%
\section{10.3 Stone 与强 Schrödinger
方程}\label{stone-ux4e0eux5f3a-schruxf6dinger-ux65b9ux7a0b}}

\textbf{K-QINF-03：Stone。} 强连续一参数酉群唯一写为
\(U(t)=e^{-itA}\)，\(A\) 自伴。定义能量 \(H=\hbar A\) 后，对于
\(\psi_0\in D(H)\)， \[
i\hbar\partial_t\psi_t=H\psi_t.
\] 域外初态仍有连续酉轨迹，但未必有强导数。Stone 给时间生成元，不给
\(\hbar\) 数值。

\protect\hypertarget{p-qinf-01}{}{}

\textbf{P-QINF-01：正常态的迹范数连续性。} 若 \(U(t)\) 强连续，任意迹类
\(\rho\) 的 \(U(t)\rho U(t)^\dagger\) 在迹范数连续。

\textbf{证明。} 先对有限秩展开，使用
\(\||x\rangle\langle y|\|_1=\|x\|\|y\|\)；任意迹类用有限秩逼近，酉共轭保持迹范数，再用三角不等式。\(\square\)

von Neumann 方程 \(\dot\rho=-i[H,\rho]/\hbar\)
只在迹类共轭群生成元的定义域内是迹范数强方程。不能对任意迹类态把两个未定义乘积写成普通矩阵式。

\hypertarget{ux65e0ux9650ux7ef4-ffips-ux4e0eux80fdux91cfux622aux65ad}{%
\section{10.4 无限维 FFIPS
与能量截断}\label{ux65e0ux9650ux7ef4-ffips-ux4e0eux80fdux91cfux622aux65ad}}

可数正交基的 \(|i\rangle\)、\((|i\rangle+|j\rangle)/\sqrt2\)、
\((|i\rangle+i|j\rangle)/\sqrt2\)
投影期望确定全部密度矩阵元，故在分离协议下正常态仍唯一代表充分未来。可数可识别不等于有限样本可高效学习。

\protect\hypertarget{p-qinf-02}{}{}

\textbf{P-QINF-02：能量截断误差。} 设 \(H\ge0\)、截断投影 \(\Pi\) 与
\(H\) 对易，且截断外 \(H\ge E_{\rm cut}(I-\Pi)\)。若
\(\operatorname{Tr}\rho H\le E\)，则 \[
\epsilon=\operatorname{Tr}[(I-\Pi)\rho]\le E/E_{\rm cut}.
\] 当 \(p=1-\epsilon>0\)，\(\rho_\Pi=\Pi\rho\Pi/p\) 满足 \[
D(\rho,\rho_\Pi)\le\sqrt{\epsilon}.
\]

\textbf{证明。} 第一式由算子序积分。取 \(\rho\) 的纯化
\(|\Psi\rangle\)，将系统投影并归一化，纯化重叠为
\(\sqrt p\)，纯态迹距离为 \(\sqrt{1-p}\)。偏迹收缩给结论。\(\square\)

若用该式获得有限维可计算模型，还需 \(\Pi\)
有限秩；仅有能量上界不自动给有限态数。多步截断须控制各步的能量尾与动力学缺陷。

\hypertarget{ch-11}{%
\chapter{11. 相空间、Weyl--CCR 与作用尺度}\label{ch-11}}

\hypertarget{ux4eceux5c04ux5f71ux4f5cux7528ux5f00ux59cb}{%
\section{11.1
从射影作用开始}\label{ux4eceux5c04ux5f71ux4f5cux7528ux5f00ux59cb}}

经典相空间 \(V=\mathbb R^{2n}\)，定义
\(\sigma((q,p),(q',p'))=q\cdot p'-p\cdot q'\)。 射影酉作用满足 \[
W(z)W(z')=m(z,z')W(z+z').
\] 结合律要求 2-cocycle 方程。重相位 \(W'(z)=f(z)W(z)\) 改变 \(m\) 一个
coboundary，不改变物理射线作用。

\textbf{K-QCCR-01。} 连续有限维实向量群的相应 multiplier
类可由实交替双线性形式表示。若另给完整线性辛协变，则该形式为
\(\kappa\sigma\)。这是表示理论和运动学条件，不是作用积分量子化的后果。

为固定系数，本章在相应重相位后取代表 \(m(z,z')=\exp[iB(z,z')/2]\)。\(B\)
的指数须无量纲；当 \(q,p\) 使用物理位置、动量单位时，\(\kappa\)
具有作用量倒数的单位。若更换 multiplier 的符号或二分之一约定，下文的
\(\kappa\) 定义也同步更换。

\protect\hypertarget{p-qccr-01}{}{}

\textbf{P-QCCR-01：辛协变的单尺度。} 若交替形式 \(B\)
被全部辛变换保持，则 \(B=\kappa\sigma\)。

\textbf{证明。} 写 \(B(u,v)=\sigma(Au,v)\)，协变使 \(A\)
与全部辛变换对易。Darboux 分块中，与
\(\operatorname{diag}(\lambda I,\lambda^{-1}I)\) 对易杀去非对角块；与
\(\operatorname{diag}(R,R^{-T})\) 对易使对角块各为标量；与辛剪切
\(\left(\begin{smallmatrix}I&C\\0&I\end{smallmatrix}\right)\)、\(C=C^\top\)
对易使两标量相等。\(\square\)

\(\kappa=0\) 仍允许。因此射影性和辛协变本身不强迫量子非交换性。

\hypertarget{uca-ux7684ux6761ux4ef6ux975eux5e73ux51e1ux6027}{%
\section{11.2 UCA
的条件非平凡性}\label{uca-ux7684ux6761ux4ef6ux975eux5e73ux51e1ux6027}}

\textbf{H-UCA。} 相对全部 canonical 位移，不存在非平凡不变正交投影；且
canonical 因子维数大于一。这是不可约性的投影表述，只施加于轨道 canonical
因子，不排除自旋等额外内部重数。

\protect\hypertarget{p-qccr-02}{}{}

\textbf{P-QCCR-02。} 在复 Hilbert 表示中，UCA 和维数大于一排除平凡
multiplier 类。

\textbf{证明。} 平凡类可重相位成普通 Abelian
酉表示，所有酉算子对易。若全为标量，则任意非平凡投影不变；若存在非标量酉算子，其非平凡谱投影与整个对易族对易。两种情况均违背
UCA。\(\square\)

结合辛协变得 \(\kappa\ne0\)，按本章 multiplier 约定选择方向使
\(\kappa=-1/\hbar\)，其中 \(\hbar=1/|\kappa|>0\)。这把直接的非零 cocycle
假设替换为特定不可约物理扇区假设，并未由一般对象性无条件证明 UCA
或作用尺度。

\hypertarget{weyl-ux5173ux7cfbux4e0eux4e25ux683c-ccr}{%
\section{11.3 Weyl 关系与严格
CCR}\label{weyl-ux5173ux7cfbux4e0eux4e25ux683c-ccr}}

采用 \[
W(q,p)=\exp[i(p\cdot Q-q\cdot P)/\hbar],\qquad
W(z)W(z')=e^{-i\sigma(z,z')/(2\hbar)}W(z+z').
\] 在 Schrödinger 表示中，可用有界酉算子直接定义 \[
[W(q,p)\psi](x)=e^{ip\cdot(x-q/2)/\hbar}\psi(x-q).
\] 逐次作用即可验证 Weyl 相位。

\textbf{K-QCCR-02：Stone--von Neumann。} 有限 canonical
自由度、非退化辛形式、非零固定中心特征、正则不可约 Weyl 表示，酉等价于
Schrödinger
表示。去除不可约性得到重数因子；去除正则性或进入无限自由度场论不再有此唯一结论。

\protect\hypertarget{p-qccr-03}{}{}

\textbf{P-QCCR-03。} 在 Schwartz 共同不变稠密域， \[
[Q_j,P_k]\psi=i\hbar\delta_{jk}\psi.
\]

\textbf{证明。} 对
\(Q_j\psi=x_j\psi\)、\(P_k\psi=-i\hbar\partial_k\psi\)
用乘积求导。两个复合均在该域定义。\(\square\)

有限维不可能精确满足
\([Q,P]=i\hbar I\)：取迹左侧零、右侧非零。截断谐振子必须报告边界修正，不得把有限矩阵误当严格
CCR 实现。

\hypertarget{ux4e0dux786eux5b9aux6027ux4e0eux95edux8f68ux754cux9650}{%
\section{11.4
不确定性与闭轨界限}\label{ux4e0dux786eux5b9aux6027ux4e0eux95edux8f68ux754cux9650}}

对归一 \(\psi\) 在相关复合算子域中，Cauchy--Schwarz 给 \[
\Delta A\,\Delta B\ge\tfrac12|\langle[A,B]\rangle|.
\] 证明是对 \((A-\langle A\rangle)\psi,(B-\langle B\rangle)\psi\)
的内积取虚部；更完整的实部项给 Robertson--Schrödinger 界。

经典周期轨道也能满足选定作用积分为
\(nh\)，而其函数仍按普通乘法交换。因此 \(\oint p\,dq=nh\) 不能推出算子
CCR。WKB 的闭轨量子化还可能带 Maslov 修正，不是一切谱的精确普适公式。

\hypertarget{ux4e0dux540cux6247ux533aux7684ux5171ux540c-hbar}{%
\section{\texorpdfstring{11.5 不同扇区的共同
\(\hbar\)}{11.5 不同扇区的共同 \textbackslash hbar}}\label{ux4e0dux540cux6247ux533aux7684ux5171ux540c-hbar}}

两独立相空间可有
\(B=\kappa_1\sigma_1\oplus\kappa_2\sigma_2\)，复合一致性不强迫相等。若增加在联合相空间上混合两扇区的完整辛协变，P-QCCR-01
给共同
\(\kappa\)，从而统一尺度。自然界全局可混合性及数值定标仍是领域事实，不能由纯
tensor composition 推出。

\hypertarget{ch-12}{%
\chapter{12. 标准非相对论计算模型}\label{ch-12}}

\hypertarget{ux81eaux7531ux6ce2ux5305ux4e0eux5e72ux6d89}{%
\section{12.1
自由波包与干涉}\label{ux81eaux7531ux6ce2ux5305ux4e0eux5e72ux6d89}}

自由 Hamiltonian \(H=P^2/(2m)\) 在适当二阶 Sobolev 域自伴。Fourier
分量相位为 \(e^{-ip^2t/(2m\hbar)}\)。初始最小不确定 Gaussian 的位置方差
\[
\sigma_x^2(t)=\sigma_0^2+\frac{\hbar^2t^2}{4m^2\sigma_0^2},
\quad \langle X_t\rangle=x_0+p_0t/m.
\] 可由 \(X_t=X_0+tP_0/m\)、初始协方差零及
\(\operatorname{Var}P=\hbar^2/(4\sigma_0^2)\) 直接推导。

双路径末态振幅为 \(\psi_1+e^{i\phi}\psi_2\)，概率包含
\(2\operatorname{Re}(e^{i\phi}\psi_2\overline{\psi_1})\)。
若路径分别关联环境 \(|E_1\rangle,|E_2\rangle\)，交叉项乘以重叠
\(\langle E_1|E_2\rangle\)。它是偏迹后的可见度系数，不自动是基本
actualization 率。

\hypertarget{ux8c10ux632fux5b50}{%
\section{12.2 谐振子}\label{ux8c10ux632fux5b50}}

\[
H=\frac{P^2}{2m}+\frac{m\omega^2Q^2}{2},\quad
a=\sqrt{\frac{m\omega}{2\hbar}}Q+\frac{iP}{\sqrt{2m\hbar\omega}}.
\] 共同域上 \([a,a^\dagger]=I\)，
\(H=\hbar\omega(a^\dagger a+1/2)\)。基态方程 \(a\psi_0=0\) 给
Gaussian；逐次升算子给 \[
E_n=\hbar\omega(n+1/2),\quad
a|n\rangle=\sqrt n|n-1\rangle.
\] 完整谱和基完备性可由标准自伴谐振子谱定理保证。

截断到 \(0,\ldots,N-1\) 时， \[
[a_N,a_N^\dagger]=I-N|N-1\rangle\langle N-1|.
\] 该式由两对角乘积逐项计算，是有限数值边界的精确诊断。

\hypertarget{ux52bfux9631ux4e0eux52bfux5792}{%
\section{12.3 势阱与势垒}\label{ux52bfux9631ux4e0eux52bfux5792}}

区间 \((0,L)\) 的 Dirichlet 无限阱： \[
\psi_n=\sqrt{2/L}\sin(n\pi x/L),\quad
E_n=\frac{\hbar^2\pi^2n^2}{2mL^2}.
\] 这是解二阶本征方程并施加两个零边界得到的谱。

有限方阱 \(V=-V_0\) 于 \(|x|<a\)，束缚能 \(-V_0<E<0\)，令
\(k=\sqrt{2m(E+V_0)}/\hbar,\kappa=\sqrt{-2mE}/\hbar\)。连续匹配波函数和导数给
\[
k\tan(ka)=\kappa\quad\text{偶态},\qquad
-k\cot(ka)=\kappa\quad\text{奇态}.
\] 须按三角函数连续分支寻根，不能跨极点误报本征值。

矩形势垒高度 \(V_0>E\)、宽 \(a\)： \[
T(E)=\left[1+
\frac{V_0^2\sinh^2(\kappa a)}{4E(V_0-E)}
\right]^{-1},\quad
\kappa=\sqrt{2m(V_0-E)}/\hbar.
\]
推导为三段平面/指数波匹配并取流比；左右渐近势不同还要加入速度比。概率流
\(j=(\hbar/m)\operatorname{Im}\bar\psi\nabla\psi\) 给透射反射守恒检验。

对于左右势都为零的势垒，令 \(k=\sqrt{2mE}/\hbar\)，在 \(0<x<a\)
内，波函数与导数的传播矩阵是 \[
\begin{pmatrix}\psi(a)\\\psi'(a)\end{pmatrix}
=\begin{pmatrix}\cosh\kappa a&\sinh(\kappa a)/\kappa\\
\kappa\sinh(\kappa a)&\cosh\kappa a\end{pmatrix}
\begin{pmatrix}\psi(0)\\\psi'(0)\end{pmatrix}.
\] 代入左侧 \(e^{ikx}+re^{-ikx}\)、右侧 \(te^{ikx}\) 的边界值并消去
\(r\)，得到 \[
|t|^{-2}=\cosh^2(\kappa a)+
\frac{(\kappa^2-k^2)^2}{4\kappa^2k^2}\sinh^2(\kappa a)
=1+\frac{(\kappa^2+k^2)^2}{4\kappa^2k^2}\sinh^2(\kappa a).
\] 使用 \(\kappa^2+k^2=2mV_0/\hbar^2\) 即得所写
\(T=|t|^2\)。此计算的能量、宽度和质量均来自已指定模型。

\hypertarget{coulomb-ux4e0eux539fux5b50}{%
\section{12.4 Coulomb 与原子}\label{coulomb-ux4e0eux539fux5b50}}

给定 Coulomb 势与约化质量 \(\mu\)， \[
H=-\frac{\hbar^2}{2\mu}\Delta-\frac{Ze^2}{4\pi\epsilon_0r}.
\] 分离球谐函数并令径向解在原点正规、无穷远可积，合流超几何级数终止给 \[
E_n=-\frac{\mu Z^2e^4}{2(4\pi\epsilon_0)^2\hbar^2n^2}.
\] \(a_0=4\pi\epsilon_0\hbar^2/(\mu Ze^2)\)，基态
\(\psi_{100}=(\pi a_0^3)^{-1/2}e^{-r/a_0}\)。
这是采用电磁势、质量、空间维数和量子扇区后的标准计算，不是从 Ω
数出氢谱常数。

令
\(g=Ze^2/(4\pi\epsilon_0)\)、\(\psi=u(r)Y_{\ell m}/r\)，束缚态径向方程为
\[
u''+\left[-\kappa^2+\frac{2\mu g}{\hbar^2r}
-\frac{\ell(\ell+1)}{r^2}\right]u=0,
\qquad \kappa=\frac{\sqrt{-2\mu E}}{\hbar}.
\] 取 \(s=2\kappa r\)、\(u=e^{-s/2}s^{\ell+1}v(s)\)，代入得到 \[
s v''+(2\ell+2-s)v'+\left(\frac{\mu g}{\hbar^2\kappa}-\ell-1\right)v=0.
\] 在通常的正规 Coulomb
自伴实现中，原点正规且无穷远平方可积要求相应级数终止：\(\mu g/(\hbar^2\kappa)=n_r+\ell+1=n\)。消去
\(\kappa\) 给
\(E_n=-\mu g^2/(2\hbar^2n^2)\)。它说明离散能量来自具体微分算子和全局边界条件；局部相位闭合的直觉不能代替本征问题。

多电子问题还需反对称性、电子相互作用与近似误差。Hartree--Fock、组态相互作用、密度泛函等是可接入方法，不由本章一体解自动获得精度保证。

\hypertarget{ch-13}{%
\chapter{13. 有限模 Gaussian 扇区}\label{ch-13}}

设 \(R=(Q_1,P_1,\ldots,Q_n,P_n)\)，
\([R_j,R_k]=i\hbar\Omega_{jk}\)。均值 \(d=\langle R\rangle\)，对称协方差
\(V_{jk}=\langle\{R_j-d_j,R_k-d_k\}\rangle/2\)。

\protect\hypertarget{p-qgauss-01}{}{}

\textbf{P-QGAUSS-01：协方差必要条件。} \[
V+\frac{i\hbar}{2}\Omega\succeq0.
\]

\textbf{证明。} 对任意复向量 \(z\)，算子 \(A=\sum_jz_j(R_j-d_j)\) 满足
\(\langle A^\dagger A\rangle\ge0\)；展开为相应二次型。\(\square\)

\textbf{K-QGAUSS-01。} 正定协方差有 Williamson 形式

\[
V=S\operatorname{diag}(\nu_1,\nu_1,\ldots,\nu_n,\nu_n)S^\top,
\]

其中 \(S\) 为辛矩阵。Gaussian 物理态要求 \(\nu_j\ge\hbar/2\)；纯
Gaussian 各 \(\nu_j=\hbar/2\)。非 Gaussian 态不由前二阶矩确定。

Gaussian 通道采用 \[
d'=Xd+d_0,\quad V'=XVX^\top+Y,
\] 其中 \(X,d_0\) 为实数据、\(Y\) 实对称，输入输出模数及相应辛形式固定。
其 CP 条件为 \[
Y+\frac{i\hbar}{2}(\Omega-X\Omega X^\top)\succeq0.
\]
这一标准条件比输出一个输入态时协方差正性更强；完整辅助模一致性是关键。

二次 Hamiltonian \(H=R^\top GR/2+g^\top R\) 导致
\(\dot d=\Omega(Gd+g)\)、\(\dot V=\Omega GV+VG^\top\Omega^\top\)，在所声明单位中由
Heisenberg 方程得到。阻尼和噪声再加入相应漂移/扩散项，须满足 CP。

这里 \(G\) 为实对称矩阵，二次项按对称次序理解，\(g\) 为实向量。由
\([R_j,R_k]=i\hbar\Omega_{jk}\)，逐项计算
\(\dot R=(i/\hbar)[H,R]=\Omega(GR+g)\)；对均值取期望得第一式。中心化算子
\(\delta R=R-d\) 满足
\(\dot{\delta R}=A\delta R\)，\(A=\Omega G\)，因此对反对易协方差求导得
\(\dot V=AV+VA^\top\)。这说明所写方程依赖已建立的 CCR 和二次
Hamiltonian，而不是从``信息几何''单独产生。

物理协方差条件可以再展开为 \[
\left\langle\delta R_j\delta R_k\right\rangle
=V_{jk}+\tfrac12\langle[R_j,R_k]\rangle
=V_{jk}+\frac{i\hbar}{2}\Omega_{jk}.
\] 于是
\(\langle A^\dagger A\rangle=z^\dagger(V+i\hbar\Omega/2)z\ge0\)，明确了符号与二分之一因子。Gaussian
通道的 CP 条件则是对所有输入及所有辅助模的要求，属于标准 Gaussian
通道定理；检查一个输出协方差只证明该输入的二阶矩没有违反必要条件。数值截断还须接上§10.4的尾界，不能以有限矩阵近似的
CCR 代替真实无限维模。

\hypertarget{ch-14}{%
\chapter{14. 对称性、自旋与非相对论粒子}\label{ch-14}}

本章及第 15--32、36--46
章以标准量子扇区和各章明列的领域输入为起点。文中的纸面命题、采用的外部重定理、精确算术回归、浮点计算和
Lean 定理按不同证据登记。有限 Kraus
内核使其中若干操作合同可以机器复查；它不自动证明本章开始使用的无限维表示、自伴域、场论或热力学极限。具体覆盖以形式化命题登记和冻结回执为准。

\hypertarget{wigner-ux4e0eux751fux6210ux5143}{%
\section{14.1 Wigner 与生成元}\label{wigner-ux4e0eux751fux6210ux5143}}

\textbf{K-QSYM-01：Wigner。}
保持全部射线转移概率的双射由酉或反酉算符实现，唯一到相位。这是针对已经给定复
Hilbert
射线空间的标准定理。它不从一般状态集合、对象名称或有限个角度的测试重建全部量子理论。

若一个群的射线作用适当连续，且选定单位元附近的连续实现，连通单位元分支由酉射影表示实现；射影乘子仍须处理，不能据此断言整个群已有单值普通酉表示。适用时可转到覆盖群或中心扩张。离散时间反演可以反酉。

无界生成元不是任意可交换的矩阵。对每个强连续酉单参数子群，Stone
定理给自伴生成元；对有限维 Lie
群的强连续酉表示，可在稠密光滑向量空间上定义共同的无穷小表示并讨论 Lie
括号。以下含 \(P,J,K,H\)
的对易式首先在这样的共同不变域上使用，随后再处理相应算符的闭包与域；形式对易式并不单独证明全局群表示存在。

有限维可逆 CPTP 动力学若其逆也 CPTP，则为酉共轭。对连续参数微分得到 \[
\dot\rho=-i[H,\rho]/\hbar.
\] 这个标准量子序同构结果不包含反酉转置通道，也不决定具体
\(H\)。例如复矩阵转置虽然保持单系统正性，但其 Choi
算符是交换算符，反对称子空间有负本征值，故它不是完全正映射。反酉是对称操作对射线的实现类别，不能因此列入任意系统上可实施的
CPTP 控制库。

\hypertarget{ux65cbux8f6cux4e0eux81eaux65cb}{%
\section{14.2 旋转与自旋}\label{ux65cbux8f6cux4e0eux81eaux65cb}}

角动量代数 \[
[J_i,J_j]=i\hbar\epsilon_{ijk}J_k
\] 的有限维不可约酉表示由 \(s=0,\tfrac12,1,\ldots\) 标记，并满足 \[
J^2=\hbar^2s(s+1),\qquad J_z=\hbar m,
\qquad m=-s,-s+1,\ldots,s.
\] 令 \(J_\pm=J_x\pm iJ_y\)，由代数得到 \[
\|J_\pm|s,m\rangle\|^2
=\hbar^2[s(s+1)-m(m\pm1)].
\] 范数非负、最高与最低权的终止及相邻权相差一，要求
\(2s\in\mathbb N\)。这是离散表示标签的来源。轨道单值标量球谐对应整数
\(\ell\)；内禀自旋属于另外声明的内部表示。

半整数自旋的向量在 \(2\pi\) 转动下变号，而射线不变；\(SO(3)\)
的相应射影表示由 \(SU(2)\)
覆盖处理。空间中的半扭转图形不替代表示论证明。仅给 \(SU(2)\)
表示也不决定粒子的质量、磁矩大小或磁矩中的 \(g\) 因子；这些仍需
Hamiltonian、物质模型或实验输入。

\hypertarget{galilei-ux4e2dux5fc3ux6269ux5f20ux4e0eux6709ux8d28ux91cfux7c92ux5b50}{%
\section{14.3 Galilei
中心扩张与有质量粒子}\label{galilei-ux4e2dux5fc3ux6269ux5f20ux4e0eux6709ux8d28ux91cfux7c92ux5b50}}

采用生成元 \(P_i,J_i,K_i,H\)，并在共同光滑域上写 \[
[K_i,P_j]=i\hbar\delta_{ij}M,\qquad
[K_i,H]=i\hbar P_i,\qquad [P_i,H]=0,
\] 其中 \(M\) 与群表示的全部生成元对易。不可约性使中心算符取标量
\(M=mI\)。本节明确选择 \(m>0\) 的自由 massive 分支；零质量或其他奇异
Galilei 表示不由下面除以 \(m\) 的推导覆盖。

定义 \[
Q_i=K_i/m,\qquad
\mathbf S=\mathbf J-\mathbf Q\times\mathbf P,
\qquad C_E=H-\frac{P^2}{2m}.
\] 第一式给轨道 CCR。由完整的 Galilei 对易关系，\(\mathbf S\) 与
\(\mathbf Q,\mathbf P\) 对易，并自身满足角动量代数。又 \[
[K_i,P^2/(2m)]=i\hbar P_i,
\] 所以 \(C_E\) 与 boosts
对易；它也与空间平移、旋转和时间平移对易。对所声明不可约表示应用 Schur
论证，得 \(C_E=E_{\rm int}I\)。在正则轨道表示与有限自旋因子下，因而有 \[
\mathcal H_1=L^2(\mathbb R^3)\otimes\mathbb C^{2s+1},\qquad
H=P^2/(2m)+E_{\rm int}.
\] 这里分类标签是
\((m,s,E_{\rm int})\)，不是从拓扑绕数算出的质量谱。整体粒子不可约性属于中心扩张后的对称群；Weyl
位移只作用于轨道因子，故其在含自旋的整个空间上可约并不矛盾。

中心扩张还有可直接检查的后果：boost 与空间平移的两种次序相差与
\(m\mathbf v\cdot\mathbf a/\hbar\)
成比例的相位，具体正负号由指数约定决定。历史基准分别作用两条算符序列检验该相位，而不是只把同一目标式代入两次。它验证已采用的表示，不能反向决定
\(m\) 或 \(\hbar\) 的数值。

\hypertarget{ux65f6ux95f4ux53cdux6f14ux4e0eux6807ux51c6ux57faux51c6}{%
\section{14.4
时间反演与标准基准}\label{ux65f6ux95f4ux53cdux6f14ux4e0eux6807ux51c6ux57faux51c6}}

时间反演满足 \(Q\to Q,P\to-P\)。若它是酉，实现后的 CCR 左侧变号而 \(i\)
不变，矛盾；反酉使 \(i\to-i\)，相容。在标准自旋表示
\(T=e^{-i\pi J_y/\hbar}K\)，\(T^2=(-1)^{2s}\)。若 \(T^2=-1\) 且
Hamiltonian 时间反演不变，\(\psi,T\psi\) 正交给 Kramers 简并。

Stern--Gerlach 自旋半态沿 \(\mathbf n\) 的两效应
\((I\pm\mathbf n\cdot\boldsymbol\sigma)/2\)，Bloch 态向量 \(\mathbf r\)
给 \(p_\pm=(1\pm\mathbf r\cdot\mathbf n)/2\)。 Rabi Hamiltonian
\(H=\hbar(\Delta Z+\Omega X)/2\) 给从计算基态到另一基态的概率 \[
P_{0\to1}(t)=\frac{\Omega^2}{\Omega^2+\Delta^2}
\sin^2\left(\frac{\sqrt{\Omega^2+\Delta^2}\,t}{2}\right).
\] 由 \((\Delta Z+\Omega X)^2=(\Delta^2+\Omega^2)I\) 展开矩阵指数得到。

CHSH 对局部确定 \(\pm1\) 值有绝对上界二；量子 Bell 态和适当 Pauli
方向达到 \(2\sqrt2\)。算子平方
\(\mathcal B^2=4I-[A_0,A_1]\otimes[B_0,B_1]\) 给范数上界
\(2\sqrt2\)。实验违反依赖测量选择、探测和无信号等实际条件，不由一次模拟替代。

\hypertarget{ch-15}{%
\chapter{15. 相同粒子、Fock 与超选}\label{ch-15}}

\hypertarget{ux4e0dux53efux533aux5206ux6027ux4e0dux5355ux72ecux9009ux62e9-bosefermi}{%
\section{15.1 不可区分性不单独选择
Bose/Fermi}\label{ux4e0dux53efux533aux5206ux6027ux4e0dux5355ux72ecux9009ux62e9-bosefermi}}

标号张量积是辅助表示。若可观测代数与 \(S_N\)
的置换表示对易，可将载体分解为 \[
\mathcal h^{\otimes N}\simeq\bigoplus_\lambda V_\lambda\otimes M_\lambda,
\qquad U_\pi=\bigoplus_\lambda U_\lambda(\pi)\otimes I.
\]
置换不变可观测量在这些因子上具有交换代数的形式。对密度矩阵本身施加置换不变性时，Schur
论证给 \[
\rho=\bigoplus_\lambda\frac{I_{V_\lambda}}{\dim V_\lambda}\otimes\tau_\lambda,
\qquad \tau_\lambda\succeq0,\qquad\sum_\lambda\operatorname{Tr}\tau_\lambda=1.
\] \(N\ge3\) 时 \(S_N\)
不只具有平凡和符号两种不可约表示。因此``交换标签不改变可观测量''不能单独排除其他表示扇区，也不能独自强迫状态只落在
Bose 或 Fermi 子空间。

标准 Bose/Fermi 选择采用额外的物理扇区条件 \[
P_\pm=\frac1{N!}\sum_{\pi\in S_N}
(\operatorname{sgn}\pi)^{(1\mp1)/2}U_\pi.
\] 群乘法与符号乘法给
\(P_\pm^2=P_\pm=P_\pm^\dagger\)。被采用的态支撑应明确写入其中一个子空间。历史文献还研究对所有辅助扩展的更强置换不变性、量子控制的置换等选择条件；那些定理的量词和许可操作不同于普通置换不变性，也不同于纯化仅在辅助可逆变换意义下唯一。不能把它们压成一个无需前提的``不可区分原则''。

三维非相对论理论中，自旋与统计仍是独立标签：形式上可以构造自旋 \(1/2\)
的 Bose 多体扇区或自旋 \(0\) 的 Fermi
多体扇区。这些是对非相对论推导能力的反例，不是自然界存在这些基本粒子的主张。相对论自旋---统计联系需要局域性、正能量、正度量及相应域条件。二维交换配置的拓扑以编织群取代普通置换群，还允许任意子；本章的三维对称化公式不能排除它们。

\hypertarget{fock-ux6784ux9020}{%
\section{15.2 Fock 构造}\label{fock-ux6784ux9020}}

\[
\mathcal F_\pm(\mathcal h)=\bigoplus_{N=0}^\infty
P_\pm\mathcal h^{\otimes N}.
\] Bose 模满足 \([a_i,a_j^\dagger]=\delta_{ij}\)，Fermi 模满足
\(\{c_i,c_j^\dagger\}=\delta_{ij}\)、\(\{c_i,c_j\}=0\)。 Fermi 的
\((c_i^\dagger)^2=0\) 给单模式占据零或一；Bose 占据任意非负整数。

Bose 产生湮灭算子无界，应先在有限粒子向量共同域定义。单模 Fermi
算子有界。数算子 \(N=\sum_i a_i^\dagger a_i\)，一体算子的二次量子化
\(d\Gamma(h)=\sum_{ij}h_{ij}a_i^\dagger a_j\)。 这组织多体 Hilbert
空间，不单独构成相对论场论。

两正交轨道 \(u,v\) 的双粒子波函数 \[
\Psi_\pm(x_1,x_2)=
[u(x_1)v(x_2)\pm v(x_1)u(x_2)]/\sqrt2
\] 给交换干涉项。轨道不正交时归一化须改为
\(\sqrt{2(1\pm|\langle u,v\rangle|^2)}\)。

\hypertarget{ux8d85ux9009ux7684ux64cdux4f5cux7c7bux578b}{%
\section{15.3
超选的操作类型}\label{ux8d85ux9009ux7684ux64cdux4f5cux7c7bux578b}}

超选是相对于可观测代数和允许操作定义的。若 \(\mathcal A\)
的中心投影给出互斥扇区，且全部允许操作不能访问其间相干，便得到该实现中的超选结构。必须区分中心
\(Z(\mathcal A)\) 与交换子代数 \(\mathcal A'\)：一般 \(\mathcal A'\)
中的投影可能只选择一个不可观测重数因子，不自动定义新的基本电荷。代数量子场论中的超选还可表现为不等价表示，不能被某个预先固定
Hilbert 空间的中心投影列表穷尽。

扩大参考系、允许操作或边界自由度可能改变操作超选。Bargmann
质量超选的推论依赖采用的 Galilei
表示与参考框架；它不说明任何涉及不同质量的形式叠加都在所有扩展理论中被禁止。Fock
构造自身也不强迫普适粒子数超选。只有在固定同一种质量 \(m\)
的粒子、采用相应质量超选并使总质量为 \(mN\)
等条件下，才得到与粒子数相连的结论。

有限模式 Fock 空间和无限模式场表示还需区分：有限个 canonical DOF 的
Stone--von Neumann 唯一性不能直接控制无限模 CCR
的全部表示。量子场中的真空、热态和其他相的表示问题在第 21 章另立合同。

\hypertarget{ch-16}{%
\chapter{16. 多体结构、局域性与相}\label{ch-16}}

\hypertarget{ux5c40ux57df-hamiltonian-ux4e0eux65e0ux9650ux4f53ux79ef}{%
\section{16.1 局域 Hamiltonian
与无限体积}\label{ux5c40ux57df-hamiltonian-ux4e0eux65e0ux9650ux4f53ux79ef}}

有限格点 \(\Lambda\) 的空间 \(\bigotimes_{x\in\Lambda}\mathcal H_x\)，
\[
H_\Lambda=\sum_{X\subseteq\Lambda}\Phi(X),
\] 其中 \(\Phi(X)\) 自伴且只作用于
\(X\)。有限维局域空间和有限范围/足够快衰减的相互作用是重要的可控类。

\textbf{K-QMB-01：Lieb--Robinson 接口。}
在有界局域项、受控格点增长与相互作用衰减条件下，分离区域的局域算子满足形如
\[
\|[A_X(t),B_Y]\|
\le C\|A_X\|\|B_Y\||X|
e^{-\mu(d(X,Y)-v|t|)}.
\] 常数由具体相互作用范数决定；这不是相对论光锥，也不无条件适用于无界
Bose 模或任意长程相互作用。

无限体积通常使用准局域 \(C^*\)
代数和状态泛函，不应假定整个无限系统仍有同一个迹一 Gibbs
矩阵。有限体积极限与长时极限、外场趋零极限可能不交换。

\hypertarget{ising-ux4e0e-heisenberg}{%
\section{16.2 Ising 与 Heisenberg}\label{ising-ux4e0e-heisenberg}}

横场 Ising： \[
H=-J\sum_iZ_iZ_{i+1}-h\sum_iX_i.
\] 一维均匀无限链在标准约定下，经 Jordan--Wigner、Fourier 和 Bogoliubov
变换的准粒子色散为 \[
\epsilon(k)=2\sqrt{J^2+h^2-2Jh\cos k}.
\] 热力学极限在 \(|h|=|J|\)
处闭隙。有限周期链的费米子边界条件依赖奇偶扇区；不能将单一动量网格用于全部有限态。

两自旋 Heisenberg \(H=(J/4)(X_1X_2+Y_1Y_2+Z_1Z_2)\)： 总自旋分解给
singlet 能量 \(-3J/4\)、triplet 能量 \(J/4\)。这是
\(\mathbf S_1\cdot\mathbf S_2=(\mathbf S_{\rm tot}^2-\mathbf S_1^2-\mathbf S_2^2)/2\)
的直接结果。

Bose--Hubbard 模型 \[
H=-t\sum_{\langle ij\rangle}(a_i^\dagger a_j+\mathrm{h.c.})
+\frac U2\sum_i n_i(n_i-1)-\mu\sum_i n_i
\] 还需占据截断误差。Fermi hopping 必须保持 CAR
的符号；把各格点独立矩阵简单相乘会丢失反对易结构，应使用固定排序的
Jordan--Wigner 或显式 Fock 基。

\hypertarget{ux76f8ux5173ux7ea0ux7f20ux4e0eux76f8ux5206ux7c7b}{%
\section{16.3
相关、纠缠与相分类}\label{ux76f8ux5173ux7ea0ux7f20ux4e0eux76f8ux5206ux7c7b}}

关联长度、谱隙、序参量、纠缠熵和拓扑指标分别回答不同问题。有限系统的能级交叉或局部涨落峰不是自动的热力学相变证明。自发对称破缺通常按
\[
\lim_{h\to0^\pm}\lim_{|\Lambda|\to\infty}
\]
选择无限体积态；交换次序可恢复对称态。有限尺寸外推因此要同时说明边界、破缺场、初态及所用观测量，不能以一张有限链曲线替代极限定理。

MPS 表示 \[
|\psi\rangle=\sum_{i_1\cdots i_L}
 v_L^\dagger A_1^{i_1}\cdots A_L^{i_L}v_R
 |i_1\cdots i_L\rangle,
\] 键维限制各切割的 Schmidt 秩。对单个切割，若归一化纯态丢弃的 Schmidt
权重为 \(\delta<1\)，归一化截断态与原态的保真振幅为
\(\sqrt{1-\delta}\)，所以纯态迹距离恰为
\(\sqrt\delta\)。这一单切割结论不能直接当作整条链、全部时刻的统一误差；顺次压缩、时间离散、Trotter、有限体积和初始近似仍需组合。

一维、有限局域维数、受控局域相互作用与统一谱隙等条件可支持面积律及 MPS
近似结果。近似存在、找到该近似的算法复杂度、任意高纠缠实时态能否低键维表示是三种不同命题。不能仅因某低能基态可用
MPS，就宣布一般多体动力学高效。

数值基准也应有区分力。单粒子 hopping 不涉及交换两费米子的符号；验证
Fermi 后端须使用真正多粒子 CAR 或显式反对称 Fock 基，例如固定 \(N=2\)
的链，并对照独立的二次量子化构造。

\hypertarget{ux591aux4f53ux5bf9ux8c61ux8bc1ux4e66}{%
\section{16.4 多体对象证书}\label{ux591aux4f53ux5bf9ux8c61ux8bc1ux4e66}}

对象候选可以是局域激发、逻辑子空间、序参量有效态或空间区域。必须先冻结目标未来、初态/环境类、允许干预、持续窗和资源尺度，再声明各证书的距离与估计方法。一个可继承的登记形式为
\[
(\epsilon_{\rm pred},\epsilon_{\rm mem},\epsilon_{\rm leak},
\Pi_{\rm pers},R_\delta,C_{\rm glob}).
\]
这些分别记录指定协议中的预测误差、已选记忆模型后的剩余误差、泄漏、持续性、可访问记录冗余及全局相干等信息；没有跨平台通用的未定义数值。\(\epsilon_{\rm mem}\)
衡量遗漏了多少记忆，不要求真实物理过程本身 Markov。

共同相关不证明父对象因果不可约。若从多体 Hamiltonian
推有效对象动力学，必须给投影、可达输入类和记忆控制；能带图或相关图本身不完成
Core 接口。一个响应充分的泛对象也未必是具有内部修复和资源维持能力的强
\(\Omega\) 对象，后者要另加合法破坏、修复和持续性合同。

同样，环境片段的 Holevo 信息
plateau、总互信息、获得给定记录精度所需的最小片段大小和可独立访问的冗余片段数不能互相替代。强客观性需要相应的操作条件；多个片段可携带同一公共记录，而联合父系统仍保留不可由公共读出看见的全局相干。第
6--8 章与第 41 章的不识别性界在多体情形继续适用。

\hypertarget{ch-17}{%
\chapter{17. 开放系统：GKSL、热浴与轨迹}\label{ch-17}}

\hypertarget{ux5faeux89c2ux7ea6ux5316ux4e0eux534aux7fa4ux4e0dux662fux540cux4e00ux5047ux8bbe}{%
\section{17.1
微观约化与半群不是同一假设}\label{ux5faeux89c2ux7ea6ux5316ux4e0eux534aux7fa4ux4e0dux662fux540cux4e00ux5047ux8bbe}}

初态为 \(\rho_S\otimes\sigma_E\) 且共同酉演化时， \[
\Lambda_t(\rho_S)=\operatorname{Tr}_E[
U_t(\rho_S\otimes\sigma_E)U_t^\dagger]
\] 是 CPTP。但一般不满足
\(\Lambda_{t+s}=\Lambda_t\Lambda_s\)，因为环境保留关联。若改变输入时初始关联也变化，未必存在对全部系统态的同一
CP assignment。

\textbf{K-QOPEN-01：有限维 GKSL。} 连续 CPTP 半群的生成元恰为 \[
\mathcal L(\rho)=-i[H,\rho]/\hbar+
\sum_\mu\left(L_\mu\rho L_\mu^\dagger-
\tfrac12\{L_\mu^\dagger L_\mu,\rho\}\right).
\] 这刻画半群物理性，不推出特定 \(L_\mu\) 的微观来源。

\hypertarget{ux4e24ux4e2aux663eux5f0fux901aux9053}{%
\section{17.2 两个显式通道}\label{ux4e24ux4e2aux663eux5f0fux901aux9053}}

Amplitude damping 参数 \(\eta=e^{-\gamma t}\)： \[
K_0=|0\rangle\langle0|+\sqrt\eta|1\rangle\langle1|,
\quad K_1=\sqrt{1-\eta}|0\rangle\langle1|.
\] 激发人口变为 \(\eta\rho_{11}\)，相干变为 \(\sqrt\eta\rho_{01}\)。

纯退相干取 \(L=\sqrt{\gamma/2}Z\)，则
\(\dot\rho=(\gamma/2)(Z\rho Z-\rho)\)，对角不变，非对角乘
\(e^{-\gamma t}\)。这也是两分支 SRN 的标准噪声，不因具有 SRN
形式就成为新坍缩证据。

\hypertarget{davies-ux63a5ux53e3}{%
\section{17.3 Davies 接口}\label{davies-ux63a5ux53e3}}

GKSL 是有限维连续 CPTP
半群的结构刻画；要从一个具体浴推导该半群，必须给微观条件。设 \[
H_S=\sum_\epsilon\epsilon\Pi_\epsilon,\qquad
H_{
\rm int}=\lambda\sum_\alpha A_\alpha\otimes B_\alpha,
\] 并声明浴的平稳态、相关衰减/混合、耦合算符的条件，以及弱耦合
\(\lambda\to0\)、宏观时间按 \(\lambda^{-2}\)
缩放的极限。有限但长的时间、耦合强度和误差范围另需估计，不能把极限生成元写成所有耦合强度下的精确方程。

定义 Bohr 频率分量 \[
A_\alpha(\omega)=
\sum_{\epsilon'-\epsilon=\hbar\omega}
\Pi_\epsilon A_\alpha\Pi_{\epsilon'},
\qquad [H_S,A_\alpha(\omega)]=-\hbar\omega A_\alpha(\omega).
\] 在相容的 Fourier 约定下，浴谱矩阵 \(\gamma_{\alpha\beta}(\omega)\)
半正定，使每个频率块的耗散项 \[
\sum_{\alpha\beta}\gamma_{\alpha\beta}(\omega)
\left[A_\beta(\omega)\rho A_\alpha(\omega)^\dagger
-\tfrac12\{A_\alpha(\omega)^\dagger A_\beta(\omega),\rho\}\right]
\] 具有 GKSL 形式。热 KMS 关系为 \[
\gamma_{\alpha\beta}(-\omega)
=e^{-\beta_T\hbar\omega}\gamma_{\beta\alpha}(\omega),
\qquad \beta_T=(k_BT)^{-1}.
\] 该标准弱耦合极限中的 Lamb shift 与 \(H_S\)
对易；简并或等频跃迁须整体处理，不能先任意删去同一 Bohr 块中的交叉项。

利用
\(A_\alpha(\omega)\rho_T=e^{-\beta_T\hbar\omega}\rho_T A_\alpha(\omega)\)
与正负频率的热关系，可以逐块验证 Gibbs 态驻定。驻定只意味着
\(\mathcal L(\rho_T)=0\)；唯一性、混合速率和任意初态热化还需不可约性/遍历性等条件。额外守恒量或无噪声子空间可留下多个驻定态。

对于 gap \(\Delta>0\) 的二能级系统，热 Bose 占据 \[
n_T=(e^{\beta_T\Delta}-1)^{-1}
\] 给
\(\gamma_\downarrow=\kappa(n_T+1)\)、\(\gamma_\uparrow=\kappa n_T\)，二者比值为
\(e^{-\beta_T\Delta}\)。这把速率连到声明的浴模型，不是把预先期望的平衡态反填进任意速率。

\hypertarget{ux91cfux5b50ux8f68ux8ff9ux4e0eux8ba1ux7b97}{%
\section{17.4
量子轨迹与计算}\label{ux91cfux5b50ux8f68ux8ff9ux4e0eux8ba1ux7b97}}

在短步 \(dt\)，跳跃概率
\(dp_\mu=dt\,\langle L_\mu^\dagger L_\mu\rangle\)，跳后态归一化为
\(L_\mu|\psi\rangle\)；未跳跃演化由
\(H_{\rm eff}=H-i\hbar\sum L_\mu^\dagger L_\mu/2\)
给出。对结果平均并展开到 \(dt\) 得 GKSL。

不同测量环境方式给不同
unraveling，却可产生同一无条件主方程。因此随机轨迹的数学存在不是某种本体坍缩被实验选择的证明。

计算中检验迹、Hermiticity、正性、步长收敛与稳态残差
\(\|\mathcal L(\rho_*)\|\)。小残差并不独自给离唯一稳态的距离，需逆生成元或混合间隙控制。

\hypertarget{ch-18}{%
\chapter{18. 非 Markov 过程与记忆对象}\label{ch-18}}

\hypertarget{ux4e3aux4ec0ux4e48ux5355ux65f6ux523bux901aux9053ux4e0dux591f}{%
\section{18.1
为什么单时刻通道不够}\label{ux4e3aux4ec0ux4e48ux5355ux65f6ux523bux901aux9053ux4e0dux591f}}

CP-divisibility 要求：对 \(t\ge s\)，存在 CPTP \(V_{t:s}\) 使
\(\Lambda_t=V_{t:s}\Lambda_s\)。它约束单时刻映射族；这些映射通常是在固定初始系统---环境制备下测得的。干预会改变中间的系统---环境关联，因此单时刻可分解性不能自动确定全部多时刻干预的统计。

过程张量将一串 CP interventions
映为末态或记录概率。相对于明确的输入/输出腿排序、Choi
归一化和转置约定，它由正 Choi
算符与逐级因果迹约束刻画；消去最后输出腿得到``相应输入单位算子 × 前一级
comb''。对历史依赖的
instrument，保留每个结果，最后再取所需的完整记录事件。把输出腿当输入腿或漏掉转置，会得到不同的配对，并非一个无害记号差。

\textbf{K-QMEM-01。} 在完整 instrument 协议下，操作 Markov
性要求当前因果断开之后，未来条件统计不再依赖更早的操作和结果；在相同
Choi 约定下，Markov 过程对应初态与逐段通道 Choi
的时间乘积。这是多时刻操作判据，采用其原始定理条件，不能用
CP-divisibility 替换。\href{https://arxiv.org/abs/1512.00589}{原始研究}

``完整''是一个量词：若只准许一小类被动读出，相同数据可能掩盖环境记忆；某些扩大的主动干预才能区分历史。有限
Kraus
内核对固定时齐操作库的全部有限词与完整有限自适应树所证明的充分性，也不能自动扩展到新增的环境接入、漂移、非时齐库或未声明的过程插槽。

\hypertarget{ux4e00ux4e2aux5b8cux5168ux53efux6267ux884cux7684ux8bb0ux5fc6ux53cdux4f8b}{%
\section{18.2
一个完全可执行的记忆反例}\label{ux4e00ux4e2aux5b8cux5168ux53efux6267ux884cux7684ux8bb0ux5fc6ux53cdux4f8b}}

系统 \(S\)、环境 \(E\) 均为 qubit。历史输入
\(S=|b\rangle\)，环境为零。第一步 \(\mathrm{CNOT}_{S\to E}\) 写入
\(b\)；将系统丢弃并重新制备为 \(|0\rangle\)；第二步
\(\mathrm{CNOT}_{E\to S}\)
回写环境记忆。两历史在断开时都有同一系统态，未来却分别为
\(|0\rangle,|1\rangle\)。

\protect\hypertarget{p-qmem-01}{}{}

\textbf{P-QMEM-01：遗漏记忆的不可消除误差。} 两历史的真未来态距离为
\(d\)，任何只用其共同当前态的预测器，最坏迹距离误差至少 \(d/2\)。

\textbf{证明。} 预测器给同一 \(\tau\)，而
\(d\le D(\sigma_0,\tau)+D(\tau,\sigma_1)\)。\(\square\)

本例 \(d=1\)，下界 \(1/2\)。它是标准量子过程对错误 Markov
压缩的反例，不是标准 QM 的违反。

\hypertarget{ux8bb0ux5fc6ux6838ux548cux6a21ux62dfux9009ux62e9}{%
\section{18.3
记忆核和模拟选择}\label{ux8bb0ux5fc6ux6838ux548cux6a21ux62dfux9009ux62e9}}

设联合线性演化 \(\dot R=\mathcal LR\)，选投影超算符 \(\mathcal P\) 与
\(\mathcal Q=1-\mathcal P\)。在算符与半群定义合法的范围内，先解 \[
\mathcal QR(t)=e^{t\mathcal Q\mathcal L}\mathcal QR(0)
+\int_0^t e^{(t-s)\mathcal Q\mathcal L}
\mathcal Q\mathcal L\mathcal PR(s)\,ds.
\] 代回 \(\mathcal P\) 方程，得到 Nakajima--Zwanzig 恒等式 \[
\frac d{dt}\mathcal PR(t)=\mathcal P\mathcal L\mathcal PR(t)
+\int_0^t\mathcal P\mathcal L e^{(t-s)\mathcal Q\mathcal L}
\mathcal Q\mathcal L\mathcal PR(s)\,ds
+\mathcal P\mathcal L e^{t\mathcal Q\mathcal L}\mathcal QR(0).
\] 在具体系统投影下，它可写成 \[
\dot\rho_S(t)=\mathcal L_0\rho_S(t)
+\int_0^t\mathcal K(t-s)\rho_S(s)\,ds+\mathcal I(t).
\] 核与非齐次项依赖所选投影和初始相关，不是一个普适、唯一的``非 Markov
数值''。原始恒等式也不保证任意截断的近似核保持
CP；这必须通过可用的结构条件或真实协议概率单独检验。

系统与单个环境的激发交换给 \(G(t)=\cos gt\)，激发区分度
\(D(t)=\cos^2gt\) 可复苏。时间局部率 \(2g\tan gt\)
在映射不可逆时奇异；不能跨零点把它当普通光滑半群生成元。区分度复苏是该模型的一个证据，不能代替对全部可允许干预的过程判据。

HEOM、pseudomode、链映射与张量网络可作为模拟后端。其误差来自浴谱逼近、辅助层截断、时间积分和初态假设，须换算到共同的目标响应后再比较。过程张量提供操作比较语言，不保证永远比其他算法便宜。

\hypertarget{ch-19}{%
\chapter{19. 热平衡、响应与熵产生}\label{ch-19}}

\hypertarget{gibbs-ux7684ux53d8ux5206ux542bux4e49}{%
\section{19.1 Gibbs
的变分含义}\label{gibbs-ux7684ux53d8ux5206ux542bux4e49}}

有限系统 \(\rho_T=e^{-\beta_TH}/Z\)，
\(F_T(\rho)=\operatorname{Tr}\rho H-\beta_T^{-1}S(\rho)\)。

\protect\hypertarget{p-qtherm-01}{}{}

\textbf{P-QTHERM-01：自由能与相对熵。} \[
F_T(\rho)-F_T(\rho_T)
=\beta_T^{-1}D(\rho\|\rho_T)\ge0.
\]

\textbf{证明。} 代入
\(\log\rho_T=-\beta_TH-\log Z\)，整理即得。相对熵非负给最小性。\(\square\)

因此给定
Hamiltonian、温度和熵定义时，平衡具有精确资源---信息变分表达。它不证明孤立酉系统必从任意初态趋向
Gibbs；酉演化保持整体谱与熵。

\hypertarget{kms-ux4e0eux8c31ux8be6ux7ec6ux5e73ux8861}{%
\section{19.2 KMS
与谱详细平衡}\label{kms-ux4e0eux8c31ux8be6ux7ec6ux5e73ux8861}}

定义 \(\tau_t(B)=e^{iHt/\hbar}Be^{-iHt/\hbar}\)。平衡 KMS 边界为 \[
F(t)=\langle A\tau_t(B)\rangle_T,\quad
F(t+i\beta_T\hbar)=\langle\tau_t(B)A\rangle_T.
\]

\protect\hypertarget{p-qtherm-02}{}{}

\textbf{P-QTHERM-02：有限 Gibbs 的 KMS。} 上述边界成立。

\textbf{证明。} 解析延拓使
\(\tau_{t+i\beta_T\hbar}(B)= e^{-\beta_TH}\tau_t(B)e^{\beta_TH}\)； 代入
Gibbs 迹并循环移动因子得到另一边界。有限维确保解析。\(\square\)

对 Hermitian \(B\)，定义
\(S_B(\omega)=\int dt\,e^{i\omega t}\langle B(t)B(0)\rangle_T\)。
Lehmann 展开给 \[
S_B(\omega)=2\pi\sum_{nm}p_n|B_{nm}|^2
\delta(\omega-(E_m-E_n)/\hbar),
\quad S_B(-\omega)=e^{-\beta_T\hbar\omega}S_B(\omega).
\] 无限系统 KMS 应在 \(C^*\)
动力学和解析元上定义，不能要求不存在的全局有限配分函数。

\hypertarget{kubo-ux4e0eux6da8ux843dux8017ux6563}{%
\section{19.3 Kubo
与涨落耗散}\label{kubo-ux4e0eux6da8ux843dux8017ux6563}}

扰动 \(H(t)=H_0-f(t)B\)，一阶相互作用绘景展开给 \[
\delta\langle A(t)\rangle=
\int_{-\infty}^tds\,f(s)\chi_{AB}(t-s),\quad
\chi_{AB}(t)=\frac{i}{\hbar}\theta(t)\langle[A(t),B]\rangle_T.
\] 由平衡谱展开， \[
\chi_{BB}''(\omega)
=\frac{1-e^{-\beta_T\hbar\omega}}{2\hbar}S_B(\omega)
\] 在上述 Fourier 约定下成立。响应符号依赖扰动写法；不可混用 \(+fB\) 和
\(-fB\)。

\hypertarget{ux9a7bux5b9aux5e73ux8861ux5316ux4e0eux70edux5316}{%
\section{19.4
驻定、平衡化与热化}\label{ux9a7bux5b9aux5e73ux8861ux5316ux4e0eux70edux5316}}

有限闭系统的长时间平均保留完整的简并能量块： \[
\overline\rho=\lim_{T\to\infty}\frac1T\int_0^T
 e^{-iHt/\hbar}\rho_0e^{iHt/\hbar}\,dt
=\sum_E P_E\rho_0P_E.
\] 在能量基底展开后，\(E\ne E'\) 的平均系数为
\(T^{-1}\int_0^T e^{-i(E-E')t/\hbar}dt\to0\)；同一能量块中的系数恒为一。因此简并块内部的相干不能任意删除，选择一组能量本征向量并只保留其对角元会给错误结果。

时间平均存在也不意味着逐时态收敛。例如 \(H=\hbar\omega Z/2\)、初态
\(|+x\rangle\)
时，\(\langle X\rangle_t=\cos\omega t\)；它的时间平均为零，但晚时振荡始终存在。有限谱还允许复现。要得到``多数晚时局域观测接近平衡''或某种均方偏差界，须增加有效维数、能隙结构、观测类和初态条件。

ETH 是针对特定非可积系统局部矩阵元结构的假设，不是所有 Hamiltonian
的定理。可积性和额外守恒量可能要求广义 Gibbs 或其他稳态描述。有限系统
Gibbs 态的变分存在性、开放热浴中的驻定、动力学平衡化及 Gibbs
热化应分别登记。

\hypertarget{spohn-ux4e0eux7269ux7406ux71b5ux4ea7ux751f}{%
\section{19.5 Spohn
与物理熵产生}\label{spohn-ux4e0eux7269ux7406ux71b5ux4ea7ux751f}}

\protect\hypertarget{p-qtherm-03}{}{}

\textbf{P-QTHERM-03。} CPTP 半群 \(\Lambda_t\) 若有满秩驻定态
\(\sigma\)，则 \(D(\Lambda_t\rho\|\sigma)\) 非增；在可微满秩轨道上 \[
-\operatorname{Tr}[\mathcal L(\rho)(\log\rho-\log\sigma)]\ge0.
\]

\textbf{证明。} 相对熵数据处理给
\(D(\Lambda_{t+s}\rho\|\Lambda_s\sigma) \le D(\Lambda_t\rho\|\sigma)\)；因
\(\sigma\) 驻定，取右导数。\(\square\)

若 \(\sigma=\rho_T\)、Hamiltonian 固定且热功划分有效，变成
\(\dot S-\beta_T\dot Q\ge0\)。任意参考态、强耦合或非 Markov
约化下，不能自动把相对熵下降叫总热力学熵产生。

\hypertarget{ch-20}{%
\chapter{20. 量子 GOD/GOA：任务、控制与长期结构}\label{ch-20}}

\hypertarget{ux4eceux5bf9ux8c61ux5230ux76eeux7684ux5fc5ux987bux589eux52a0ux4efbux52a1}{%
\section{20.1
从对象到目的必须增加任务}\label{ux4eceux5bf9ux8c61ux5230ux76eeux7684ux5fc5ux987bux589eux52a0ux4efbux52a1}}

固定允许动作、资源、测量反馈和奖励
\(r(\rho,a)\)。例如目标态保真度、逻辑存活、能耗或信息获取有不同偏好。一个稳定量子对象不自动最大化全部指标。

无观测反馈时状态按 CPTP \(\rho'=\mathcal T_a(\rho)\)
更新；有观测时每动作给 instrument，后验 \(\rho_{a,o}\) 和概率
\(p(o|\rho,a)\)。控制状态还可含经典预算、设备模式和环境信念。

\hypertarget{bellman-ux548cux91cfux5b50ux76eeux6807ux65b9ux5411}{%
\section{20.2 Bellman
和量子目标方向}\label{bellman-ux548cux91cfux5b50ux76eeux6807ux65b9ux5411}}

有限动作、有界奖励、折扣 \(0<\beta_C<1\)，定义 \[
(TV)(\rho)=\max_a\left[r(\rho,a)+
\beta_C\sum_op(o|\rho,a)V(\rho_{a,o})\right].
\]

\protect\hypertarget{p-qctl-01}{}{}

\textbf{P-QCTL-01。} \(T\) 在有界函数空间为 \(\beta_C\)
压缩，故有唯一有界固定点；当所需可测性成立时有限动作可选可测最优策略。

\textbf{证明。} 每动作条件期望的函数差至多
\(\|V-W\|_\infty\)，取最大保持该界。Banach
不动点定理给固定点；有限动作按固定顺序打破并列，给可测选择。迭代或折扣
telescoping 给最优性。\(\square\)

GOD 是最优动作集合；若动作对应 Hamiltonian 方向或 Lindblad
生成元，再映为相应动力学集合，不强造一个欧氏向量。未来动作集合、仪器和奖励变更时，应重算价值。

\hypertarget{ux6709ux9650ux65f6ux91cfux5b50ux63a7ux5236ux68afux5ea6}{%
\section{20.3
有限时量子控制梯度}\label{ux6709ux9650ux65f6ux91cfux5b50ux63a7ux5236ux68afux5ea6}}

闭系统 \(H(t)=H_0+\sum_k u_k(t)H_k\)，目标
\(J=\langle\psi_T|O|\psi_T\rangle-\int c(u)dt\)。 定义伴随态
\(|\chi_t\rangle=U(T,t)^\dagger O|\psi_T\rangle\)，Duhamel 变分给 \[
\frac{\delta J}{\delta u_k(t)}
=\frac2\hbar\operatorname{Im}\langle\chi_t|H_k|\psi_t\rangle
-\partial_{u_k}c.
\] 由 \(\delta\psi_T=-(i/\hbar)U(T,t)H_k\psi_t\,\delta u_kdt\)
和两共轭项求和可直接检查符号。离散脉冲梯度还要说明矩阵指数导数近似；非凸控制不保证找到全局
GOD。

\hypertarget{goa-ux548cux7a33ux5b9aux6027}{%
\section{20.4 GOA 和稳定性}\label{goa-ux548cux7a33ux5b9aux6027}}

固定最优反馈后，GOA
可以是稳态集合、极限环、噪声无扰子空间、占据分布或条件存活结构。它不是普适唯一纯态。

有限时间近似通道给第 4
章误差合成；若一个闭环通道在迹距离上有统一收缩系数
\(\alpha<1\)，则有唯一固定态，扰动通道的固定态距离至多
\(\epsilon/(1-\alpha)\)。证明与 Core 相同：插入一个共同通道输出得
\(d\le\alpha d+\epsilon\)。无收缩、近简并或分岔时，不应从很小单步误差推出长期身份稳定。

第52章证明无需最优控制或严格收缩前提的有限通道平稳态存在，并给出周期轨道与多平稳态反例；第53章则先指定存活结果，再证明
QSD 生存律和 Doob
的不同不变分布。这些成果为持续结构提供实际矩阵证书，仍不产生一个未经声明的评价函数。

\hypertarget{ch-21}{%
\chapter{21. 相对论量子场：论域与新增前提}\label{ch-21}}

从本章起，除特别注明，采用 \(\hbar=c=1\)、度规
\(\eta=\operatorname{diag}(1,-1,-1,-1)\)、\(p x=p^0x^0-\mathbf p\cdot\mathbf x\)。质量、动量和能量同量纲，长度和时间为能量的倒数。连续变量章节中的有限自由度唯一性不能直接迁入无限自由度场论。

\hypertarget{ux573aux8bbaux5b9eux73b0ux7684ux7269ux7406ux8f93ux5165}{%
\section{21.1
场论实现的物理输入}\label{ux573aux8bbaux5b9eux73b0ux7684ux7269ux7406ux8f93ux5165}}

\textbf{H-QFT-1。} 声明闵可夫斯基背景、强连续的 Poincaré
酉实现、能动量联合谱位于前向闭光锥、局域性、正定物理态空间，以及所用真空的存在、生成性、正则性和必要的唯一性条件。相互作用模型还要给场内容、作用量、调节与重整化条件。这些是进入该场论扇区的物理和数学输入，不是一般对象闭合的后果。谱条件本身不保证真空唯一。

场 \(\phi(x)\) 通常是算符值分布。对 Schwartz 测试函数等指定函数类定义 \[
\phi(f)=\int d^4x\,f(x)\phi(x),
\]
并在共同不变稠密域上讨论可闭性、乘积与矩阵元的分布性质。点处场、同点场乘积和时序乘积不能被默认当成处处有界矩阵；后两者还可能需要独立的延拓或重整化。

对有界区域 \(\mathcal O\)，局域可观测代数 \(\mathcal A(\mathcal O)\)
满足包含单调性、协变性和类空交换性。自由 Bose 场可从有界 Weyl
算符生成局域代数；费米场具有分次局域性，物理偶可观测量满足通常的类空交换性。一个物理对象接口可记录区域、代数、限制态与允许局域
instruments： \[
(\mathcal O,\mathcal A(\mathcal O),\omega|_{\mathcal A(\mathcal O)},\mathfrak I_{\mathcal O}).
\] 其中再附时间/能量分辨率和近似误差，不预先强迫尖锐区域拥有独立的
Hilbert 张量因子。

\hypertarget{ux65e0ux9650ux81eaux7531ux5ea6ux8868ux793aux4e0eux5c40ux57dfux4ee3ux6570}{%
\section{21.2
无限自由度、表示与局域代数}\label{ux65e0ux9650ux81eaux7531ux5ea6ux8868ux793aux4e0eux5c40ux57dfux4ee3ux6570}}

有限 canonical DOF 的正则不可约 CCR 表示唯一性不延伸到无限模。以 Bose
Fock 表示为例，对适当的线性正则 Bogoliubov 变换，是否可在原 Fock
空间由酉算符实现，受其产生---湮灭混合部分 \(\beta\) 的 Hilbert--Schmidt
条件控制： \[
\|\beta\|_{\rm HS}^2<\infty.
\] 对独立模式的相容对角变换，\(\beta_n=1/n\) 的平方和收敛，而
\(\beta_n=1/\sqrt n\)
发散。这是同样逐模合法的变换在无限系统中可能具有不同实现状态的明确诊断；不是把每个模式截断后再悄悄交换无限极限。真空、热态及不同相因此可能需要不等价表示。

在通常连续相对论模型的标准条件下，区域局域代数具有 Type III
等不同于有限矩阵代数的结构。即使整个 Fock 空间上的
\(\mathcal B(\mathcal H)\) 是 Type I，也不能据此把每个局域代数写为
\(\mathcal B(\mathcal H_{\mathcal O})\otimes I\)。尖锐区域限制态不自动有相对于普通局域迹的密度矩阵，普通约化纠缠熵也不自动有限。

因此，``迹掉区域外''必须说明其实现：可以使用格点或能量截断，也可在附加
split/nuclearity 等条件下，对有间隔的内外区域引入中间 Type I
因子。这样的 split 选择有条件且通常不唯一。一般始终可讨论代数态限制 \[
\omega\longmapsto\omega|_{\mathcal A(\mathcal O)},
\]
并在适用框架内使用相对熵或模理论量，而非宣称一个不存在的有限局域配分函数或普通局域迹。

Reeh--Schlieder
的局域循环性结论需要谱、真空与适当生成性等定理条件。循环性不允许远端自由控制：若类空两域的
\(B\) 与全部局域 Kraus \(K_\mu\) 交换，且
\(\sum_\mu K_\mu^\dagger K_\mu=I\)，则 \[
\sum_\mu K_\mu^\dagger B K_\mu=B.
\]
故非选择局域操作不改变另一域的期望；依赖具体测量结果的条件态必须同时计入概率与经典结果通信。

Haag
定理在其真空、协变性及场条件下限制了自由场与相互作用场之间天真的全局酉等价。它不证明相互作用不可能，而是要求区别有限调节模型、微扰相互作用图像与移除调节后的严格场论。Core
的响应和限制语言可以接收这些实现；有限维密度、熵、diamond
范数或矩阵算法若要延伸，仍需明确拓扑、能量约束和极限合同。

\hypertarget{ux4e09ux79cdux95edux5408ux72b6ux6001}{%
\section{21.3 三种闭合状态}\label{ux4e09ux79cdux95edux5408ux72b6ux6001}}

自由场可显式构造；带调节的有限模型可直接定义并计算；四维相互作用的微扰级数只按阶给出形式/渐近预测。这三种状态不能合称``严格非微扰完成''。本书没有证明四维
Yang--Mills 质量隙、完整 QED 非微扰存在性或量子引力。

\hypertarget{ch-22}{%
\chapter{22. 自由标量场：模式、传播与因果}\label{ch-22}}

\hypertarget{kleingordon-ux573a}{%
\section{22.1 Klein--Gordon 场}\label{kleingordon-ux573a}}

取 \[
\mathcal L=\tfrac12\partial_\mu\phi\,\partial^\mu\phi-\tfrac12m^2\phi^2,\qquad
(\Box+m^2)\phi=0,\qquad E_{\mathbf p}=\sqrt{\mathbf p^2+m^2}.
\] 定义 \[
[a(\mathbf p),a^\dagger(\mathbf q)]=(2\pi)^3\,2E_{\mathbf p}\,
\delta^3(\mathbf p-\mathbf q),\quad [a,a]=0,
\] \[
\phi(x)=\int\frac{d^3p}{(2\pi)^3\,2E_{\mathbf p}}
\left[a(\mathbf p)e^{-ipx}+a^\dagger(\mathbf p)e^{ipx}\right].
\] 真空满足 \(a(\mathbf p)|0\rangle=0\)。通常序 Hamiltonian 为 \[
:H: =\int\frac{d^3p}{(2\pi)^3\,2E_{\mathbf p}}\,
E_{\mathbf p}a^\dagger(\mathbf p)a(\mathbf p).
\] 这选择了平直背景中的正频分解；一般弯曲时空未必有唯一全局真空。

\protect\hypertarget{p-qkg-01}{}{}

\textbf{P-QKG-01：正则关系与传播子。} 上述模式在分布意义下给出 \[
[\phi(t,\mathbf x),\pi(t,\mathbf y)]=i\delta^3(\mathbf x-\mathbf y),
\quad \pi=\partial_t\phi.
\] 证明是将两组模式代入，使用唯一非零的 \(a,a^\dagger\)
交换项；时间导数中的 \(E_{\mathbf p}\) 抵消测度分母，两项各给一半
Fourier delta。其余等时交换子因动量反演而为零。\(\square\)

\hypertarget{ux4e09ux79cdux6838ux4e0dux6df7ux6dc6}{%
\section{22.2 三种核不混淆}\label{ux4e09ux79cdux6838ux4e0dux6df7ux6dc6}}

\[
\Delta^+(x)=\langle0|\phi(x)\phi(0)|0\rangle
=\int\frac{d^3p}{(2\pi)^3\,2E_{\mathbf p}}e^{-ipx},
\] \[
[\phi(x),\phi(y)]=\Delta^+(x-y)-\Delta^+(y-x)=i\Delta(x-y),
\] \[
\Delta_F(x)=\langle0|T\phi(x)\phi(0)|0\rangle
=\int\frac{d^4p}{(2\pi)^4}\frac{i\,e^{-ipx}}{p^2-m^2+i0}.
\] 由分母乘回得 \((\Box+m^2)\Delta_F=-i\delta^4\)。因果响应使用 retarded
核，而不是把 Feynman 核解释成经典信号传播概率。

\protect\hypertarget{p-qkg-02}{}{}

\textbf{P-QKG-02：类空因果性。} 对 \(x^2<0\)，作保时向 Lorentz 变换到
\(x^0=0\)，此时交换子积分为动量反演奇项，故为零。协变性把结论带回原坐标。类空分离的适当涂抹可观测量因此交换。\(\square\)

但类空 Wightman 函数不为零；例如 \(m>0,r=|\mathbf x|\) 时 \[
\Delta^+(0,\mathbf x)=\frac{m}{4\pi^2r}K_1(mr).
\]
真空关联不等于可控超光速信号。若局域操作与另一局域效应交换，非选择操作保持另一侧期望；有选择的条件统计变化仍需要经典结果通信才能利用。

\hypertarget{ux6e90ux4e0eux63a2ux6d4bux5668ux63a5ux53e3}{%
\section{22.3
源与探测器接口}\label{ux6e90ux4e0eux63a2ux6d4bux5668ux63a5ux53e3}}

外源 \(J\) 耦合 \(\int J\phi\)
时，线性响应由交换子给出。局域两能级探测器可用
\(H_I(\tau)=g\chi(\tau)\mu(\tau)\phi(x(\tau))\)
建模，二阶跃迁概率由涂抹的 \(\Delta^+\)
积分决定。切换函数、轨迹、能隙和调节都是输入；无穷快点探测器极限可能发散。

\hypertarget{ch-23}{%
\chapter{23. Dirac 场、Maxwell 场与物理自由度}\label{ch-23}}

\hypertarget{ux65cbux91cfux4e0eux53cdux7c92ux5b50}{%
\section{23.1 旋量与反粒子}\label{ux65cbux91cfux4e0eux53cdux7c92ux5b50}}

Clifford 关系 \[
\{\gamma^\mu,\gamma^\nu\}=2\eta^{\mu\nu}I,\quad
\bar\psi=\psi^\dagger\gamma^0,\quad
\mathcal L_D=\bar\psi(i\gamma^\mu\partial_\mu-m)\psi
\] 给出 \((i\not\partial-m)\psi=0\)。动量旋量满足 \[
(\not p-m)u_s(p)=0,\quad(\not p+m)v_s(p)=0,
\] \[
\sum_su_s\bar u_s=\not p+m,\quad
\sum_sv_s\bar v_s=\not p-m
\] （对应通常协变归一化）。模式展开为 \[
\psi(x)=\sum_s\int\frac{d^3p}{(2\pi)^3\,2E_p}
\left[b_s(p)u_s(p)e^{-ipx}+d_s^\dagger(p)v_s(p)e^{ipx}\right],
\] 其中 \(b,d\) 分别服从相同归一化的
CAR。通常序能量计数粒子和反粒子的正能激发；通常序电荷给出二者相反符号。

\protect\hypertarget{p-qdir-01}{}{}

\textbf{P-QDIR-01。} 将 Dirac 算符与其共轭因子相乘，
\((i\not\partial-m)(i\not\partial+m)=-(\Box+m^2)\)，所以每个旋量分量满足
KG 方程，传播子为 \[
S_F(x)=(i\not\partial+m)\Delta_F(x),\qquad
S_F(p)=\frac{i(\not p+m)}{p^2-m^2+i0}.
\] 混合导数的反对称部分由 Clifford 关系消掉。\(\square\)

从同一模式展开还可以检验 Dirac 场的分次因果性。令 \[
D_0(x)=\Delta^+(x)-\Delta^+(-x)=i\Delta(x),
\] 并保持第 22 章的归一化与符号约定，则 \[
\{\psi_\alpha(x),\bar\psi_\beta(y)\}
=\bigl[(i\not\partial_x+m)D_0(x-y)\bigr]_{\alpha\beta}.
\] 证明是分别用 \(u\bar u=\not p+m\) 与 \(v\bar v=\not p-m\)
的自旋求和；\(i\partial\) 在正、负频率项上产生相反的动量符号。\(D_0\)
的因果支撑在取导数后仍成立，故类空分离的费米场反对易；由偶数个费米场构成的适当局域可观测量通常交换。检验有质量模型须保留
\(m\) 项，不能以无质量核的数值替身宣称已验证整个 massive 接口。

\textbf{K-QSTAT-01：自旋---统计。}
在相对论局域、正能谱、正定态空间及相应真空/域等定理条件下，整数/半整数自旋分别服从
Bose/Fermi 统计。这里采用标准重定理，不声称由 Core 独立证明。它不等于第
15
章的普通不可区分性，也不能原样搬到未经检查的二维编织统计或其他超出定理域的模型。

\hypertarget{maxwell-ux7ea6ux675fux4e0eux6a2aux5411ux6a21ux5f0f}{%
\section{23.2 Maxwell
约束与横向模式}\label{maxwell-ux7ea6ux675fux4e0eux6a2aux5411ux6a21ux5f0f}}

取 \[
\mathcal L_M=-\tfrac14F_{\mu\nu}F^{\mu\nu},\qquad
F_{\mu\nu}=\partial_\mu A_\nu-\partial_\nu A_\mu,
\qquad A_\mu\sim A_\mu+\partial_\mu\alpha.
\] \(A_0\) 无独立动能，其正则共轭为零；要求约束随时间保持产生 Gauss
约束。每个非零动量上的四个势分量经过约束与规范冗余，只剩两个物理横向偏振。它们给
Maxwell 光子的 helicity
\(\pm1\)。选择偏振向量的具体相位与局部标架是表示选择，不能把标架变化当作新增物理偏振。

在库仑规范下，横向投影为 \[
P^T_{ij}(\mathbf k)=\delta_{ij}-k_i k_j/|\mathbf k|^2.
\] 它在位置空间非局域；这不证明可观测信号超光速。物理电磁场
\(F_{\mu\nu}\) 的因果代数仍需按第 21
章处理。规范势的辅助传播子、物理场强的交换子及可制备/探测的光子态不是同一个对象。

不能无条件同时坚持非平凡 \(A_\mu\) 具有点局域、Lorentz 协变、严格
Maxwell 方程和正定 Hilbert 描述。标准 Maxwell 势的 no-go
结果解释了为何协变规范通常采用辅助不定度量空间，再用 subsidiary
condition 或 BRST
商得到物理正态空间；另一路是在物理空间直接使用横向模式或局域规范不变场。四个
\(A_\mu\) 模式都被当成正范数光子，会重复计算非物理自由度。

\hypertarget{ux5e26ux7535ux5bf9ux8c61gauss-ux5b9aux5f8bux4e0eux8fb9ux754c}{%
\section{23.3 带电对象、Gauss
定律与边界}\label{ux5e26ux7535ux5bf9ux8c61gauss-ux5b9aux5f8bux4e0eux8fb9ux754c}}

Gauss
定律把电荷与围绕区域的电通量联系起来。因此``带电粒子完全局域在一个独立
Hilbert 因子''的直觉需要修正。区域划分可能需要通量中心、边缘变量或带
dressing
的算符。不同选择会改变熵和载体分解，必须写入协议接口。规范冗余的消除与物理边界自由度的删除不是同一操作。

\hypertarget{ch-24}{%
\chapter{24. 相互作用、生成泛函与散射}\label{ch-24}}

\hypertarget{ux5e26ux8c03ux8282ux7684-lambdaphi4-ux57faux51c6}{%
\section{\texorpdfstring{24.1 带调节的 \(\lambda\phi^4\)
基准}{24.1 带调节的 \textbackslash lambda\textbackslash phi\^{}4 基准}}\label{ux5e26ux8c03ux8282ux7684-lambdaphi4-ux57faux51c6}}

必须先固定裸参数与重整化参数的记法。裸拉氏量可写成 \[
\mathcal L_0=\tfrac12(\partial\phi_0)^2
-\tfrac12m_0^2\phi_0^2-\frac{\lambda_0}{4!}\phi_0^4.
\] 它本身不再额外加一份独立反项。若在 \(d=4-\epsilon\) 的 dimensional
regularization 中改用重整化场和无量纲耦合，定义 \[
\phi_0=Z_\phi^{1/2}\phi,\qquad
Z_\phi m_0^2=m^2+\delta m^2,\qquad
Z_\phi^2\lambda_0=\mu^\epsilon(\lambda+\delta\lambda),
\qquad Z_\phi=1+\delta Z_\phi.
\] 同一个裸式于是被分解为 \[
\begin{aligned}
\mathcal L={}&\tfrac12(\partial\phi)^2-\tfrac12m^2\phi^2
-\frac{\mu^\epsilon\lambda}{4!}\phi^4+\mathcal L_{\rm ct},\\
\mathcal L_{\rm ct}={}&\tfrac12\delta Z_\phi(\partial\phi)^2
-\tfrac12\delta m^2\phi^2-\frac{\mu^\epsilon\delta\lambda}{4!}\phi^4.
\end{aligned}
\]
这是一套参数重写，不是新增两组独立质量和耦合。四维极限下二点反项顶点为
\(i(\delta Z_\phi p^2-\delta m^2)\)，四点反项顶点为
\(-i\delta\lambda\)；尚未取 \(\epsilon\to0\) 时四点顶点须保留
\(\mu^\epsilon\)。第 27 章 QED 使用 \(d=4-2\epsilon_{\rm DR}\)，其 pole
记法应另行换算，不能把两个 epsilon 直接混用。

经典实标量势的全局稳定性要求适当的正四次项；微扰方案中的裸/重整化参数关系与全局稳定性不是同一陈述。第
25 章的 Euclidean
硬截止质量基准使用其明确的减除方案；从不同调节得到的参数需要匹配后再比较物理量。

相互作用图像记号为 \[
U_I(t,t_0)=T\exp\left[-i\int_{t_0}^{t}H_I(s)\,ds\right].
\] 其第 \(n\) 阶是时间有序积分；改成全区域积分时出现
\(1/n!\)。在有限调节与适当算符条件下可以先定义动力学；对四维连续理论，Dyson
级数通常按微扰阶使用。调节移除、渐近态极限和级数求和不得凭一个指数符号交换，也不能据此避开第
21 章的 Haag 边界。

\hypertarget{wick-ux5c55ux5f00ux4e0e-feynman-ux89c4ux5219}{%
\section{24.2 Wick 展开与 Feynman
规则}\label{wick-ux5c55ux5f00ux4e0e-feynman-ux89c4ux5219}}

自由真空的时序偶数点函数等于全部成对收缩之和，奇数点函数为零（K：高斯/Wick
定理）。将 Dyson 级数与其组合，得到内部标量线 \(i/(p^2-m^2+i0)\)、顶点
\(-i\lambda\)、每个独立圈积分
\(\int d^4\ell/(2\pi)^4\)，以及来自图自同构和展开计数的对称因子。

真空归一化生成泛函为 \[
\begin{gathered}
Z[J]=\frac{\int\mathcal D\phi\,
e^{iS[\phi]+i\int J\phi}}{\int\mathcal D\phi\,e^{iS[\phi]}},
\\ W[J]=-i\log Z[J],\\
\Gamma[\varphi]=W[J]-\int J\varphi,\\
\varphi=\frac{\delta W}{\delta J}.
\end{gathered}
\] 这些符号在有限调节下或微扰意义使用；Lorentz
路径积分不是自动存在的普通概率测度。连接图由 \(W\) 生成，1PI 顶点由
\(\Gamma\) 生成。

\hypertarget{lsz-ux4e0eux622aux9762}{%
\section{24.3 LSZ 与截面}\label{lsz-ux4e0eux622aux9762}}

\textbf{K-QLSZ-01。}
若存在适用的稳定单粒子质量壳、非零极点留数、适当渐近态与谱条件，则连接相关函数去除外腿传播子、乘外腿留数因子后得到散射振幅。长程带电场的红外问题使简单带电
Fock 渐近态版本不能原样套用。

约定 \(S=I+iT\)，
\(\langle f|iT|i\rangle=i(2\pi)^4\delta^4(P_f-P_i)\mathcal M\)。
二体初态通量 \(4\sqrt{(p_1p_2)^2-m_1^2m_2^2}\)，终态测度 \[
d\Phi_n=(2\pi)^4\delta^4(P-\sum k_i)
\prod_i\frac{d^3k_i}{(2\pi)^3\,2E_i}.
\] 相同终态粒子要除以对应阶乘。质心二到二散射 \[
\frac{d\sigma}{d\Omega}=\frac1{64\pi^2s}
\frac{|\mathbf p_f|}{|\mathbf p_i|}\overline{|\mathcal M|^2}
\] 再按相同粒子与角域约定处理。自旋平均只对未制备分辨的初态进行。

\protect\hypertarget{p-qunit-01}{}{}

\textbf{P-QUNIT-01：光学定理。} \(S^\dagger S=I\) 给出
\(-i(T-T^\dagger)=T^\dagger T\)。取初态对角元并插入完整终态，
\(2\operatorname{Im}\mathcal M_{ii}\)
等于带相空间积分的中间态振幅模方总和，含统计因子。\(\square\)
它约束每个微扰阶，而不是保证任意截断振幅精确幺正。

\hypertarget{ch-25}{%
\chapter{25. 重整化：显式一圈基准与方案}\label{ch-25}}

\hypertarget{ux4e8cux70b9ux51fdux6570}{%
\section{25.1 二点函数}\label{ux4e8cux70b9ux51fdux6570}}

Euclidean 球形硬截止的径向积分是 \[
I_\Lambda(m)=\int_{|q|<\Lambda}\frac{d^4q}{(2\pi)^4}
\frac1{q^2+m^2}
=\frac1{16\pi^2}\left[\Lambda^2-m^2\log(1+\Lambda^2/m^2)\right].
\] \protect\hypertarget{p-qren-01}{}{}

\textbf{P-QREN-01。} 四维球面积 \(2\pi^2\)，令 \(u=q^2\) 后直接积分
\(u/(u+m^2)\) 得上式。四次顶点的 tadpole 组合因子为 \(1/2\)，故
\(\Sigma^{(1)}=\lambda I_\Lambda/2\)。以零动量质量条件吸收该常数，
\(m_0^2=m_R^2-\lambda I_\Lambda(m_R)/2+O(\lambda^2)\)，得
\(\Gamma_R^{(2)}(p_E)=p_E^2+m_R^2+O(\lambda^2)\)；
一圈该图无动量依赖，所以此阶波函数重整化为零。\(\square\)

二次截止项是此方案的参数关系，不是独立可观测量；改变调节并不意味着物理预测改变。

\hypertarget{ux56dbux70b9ux51fdux6570ux4e0eux9608ux503c}{%
\section{25.2
四点函数与阈值}\label{ux56dbux70b9ux51fdux6570ux4e0eux9608ux503c}}

在三个等价 Euclidean 通道的指定减除配置下，定义 \[
B_R(Q;\mu,m)=\frac1{16\pi^2}\int_0^1dx
\log\frac{m^2+x(1-x)\mu^2}{m^2+x(1-x)Q^2},
\] \[
F(Q)=\lambda-\tfrac32\lambda^2 B_R(Q;\mu,m)+O(\lambda^3).
\] 这里 \(F_{\rm tree}=\lambda\)，Minkowski 树振幅则为
\(\mathcal M_{\rm tree}=-\lambda\)。不得以符号差误判物理不一致。一般运动学必须分别使用三个通道，不能任意设
\(s=t=u\)。

\protect\hypertarget{p-qren-02}{}{}

\textbf{P-QREN-02：泡图减除。} 用
\(1/(AB)=\int_0^1dx/[xA+(1-x)B]^2\)，平移积分变量后分母为
\((\ell^2+m^2+x(1-x)Q^2)^2\)。相同调节下减去 \(\mu\)
配置，发散常数抵消，径向积分给出上述对数；三个通道各带
\(1/2\)。\(\square\)

检验点 \((\lambda,m,\mu,Q)=(0.4,1,2,5)\) 给

\[
B_R=-0.00668304090527,\qquad F=0.401603929817.
\]

这是一圈截断数值，不是精确全阶结果。

对 Minkowski 通道 \(z\)，继续为 \[
B_R(z;z_0)=\frac1{16\pi^2}\int_0^1dx
\log\frac{m^2-x(1-x)z_0-i0}{m^2-x(1-x)z-i0},
\quad z_0=-\mu^2,
\] \[
\mathcal M(s,t,u)=-\lambda+
\frac{\lambda^2}{2}\bigl[B_R(s)+B_R(t)+B_R(u)\bigr]+O(\lambda^3).
\] \protect\hypertarget{p-qren-03}{}{}

\textbf{P-QREN-03：二粒子割。} 当 \(s>4m^2\)，对数分母为负的区间长度为
\(\beta_s=\sqrt{1-4m^2/s}\)，所以 \[
\operatorname{Im}\mathcal M^{(1)}=\frac{\lambda^2\beta_s}{32\pi}.
\] 另一方面 \(\int d\Phi_2=\beta_s/(8\pi)\)，相同标量除以 \(2!\)；
光学定理右端为 \(\lambda^2\beta_s/(16\pi)\)，正好是两倍虚部。\(\square\)

\hypertarget{running-ux4e0eux4e34ux754cux70b9}{%
\section{25.3 Running 与临界点}\label{running-ux4e0eux4e34ux754cux70b9}}

\textbf{K-QREN-01。} 四维实标量的无质量/UV MS 一圈函数为 \[
\beta_\lambda=\mu\frac{d\lambda}{d\mu}
=\frac{3\lambda^2}{16\pi^2}+O(\lambda^3).
\] 正号表示向高能耦合增长。带质量动量减除方案存在阈值函数 \[
\beta_{\rm MOM}=\frac{3\lambda^2}{16\pi^2}\,
T(\mu/m)+O(\lambda^3),\quad
T(r)=\int_0^1\frac{x(1-x)r^2}{1+x(1-x)r^2}\,dx.
\] \(T(r)\sim r^2/6\) 于红外，UV 趋于1。不能把一个方案的所有尺度 beta
原样搬到另一个方案。

对连接重整化 \(n\) 点函数，若 \(\phi_0=Z_\phi^{1/2}\phi_R\)、
\(\gamma_\phi=\tfrac12\mu d\log Z_\phi/d\mu\)，则 \[
[\mu\partial_\mu+\beta_\lambda\partial_\lambda+
\beta_{m^2}\partial_{m^2}+n\gamma_\phi]G_R^{(n)}=0
\] （保持裸量不变）。1PI 顶点场异常维数项符号相反。

在 \(d=4-\epsilon\) 的 epsilon 展开中，
\(\beta_\lambda=-\epsilon\lambda+3\lambda^2/(16\pi^2)+\cdots\)，
非平凡固定点 \(\lambda_*=16\pi^2\epsilon/3+O(\epsilon^2)\)。 这是小
\(\epsilon\) 的受控级数；取 \(\epsilon=1\)
时须报告截断和重求和不确定性。

\hypertarget{ux65b9ux6848ux5339ux914dux5c55ux5f00ux53c2ux6570ux4e0eux72ecux7acb-rg-ux68c0ux9a8c}{%
\section{25.4 方案匹配、展开参数与独立 RG
检验}\label{ux65b9ux6848ux5339ux914dux5c55ux5f00ux53c2ux6570ux4e0eux72ecux7acb-rg-ux68c0ux9a8c}}

leading beta
系数的相同性有自己的合同。若两种无质量、质量无关方案只作显式尺度无关的有限重定义
\[
\lambda'=\lambda+c\lambda^2+O(\lambda^3),
\] 则链式法则给 \[
\beta_{\lambda'}=(1+2c\lambda+\cdots)\beta_\lambda
=b\lambda'^2+O(\lambda'^3)
\] 当 \(\beta_\lambda=b\lambda^2+\cdots\)。若 \(c\) 本身依赖
\(\mu/m\)，显式 \(\mu\partial_\mu c\) 已在同一阶增加项；这正是带质量 MOM
阈值不必等于全尺度 MS beta
的原因。物理比较应同时转换耦合、质量和外腿归一化，并固定匹配条件。

Wilson--Fisher 展开用 \[
g=\frac{\lambda}{16\pi^2},\qquad
\beta_g=-\epsilon g+3g^2+O(g^3)
\] 更直接显示圈展开参数，所以
\(g_*=\epsilon/3+O(\epsilon^2)\)。固定点附近还须调到临界质量面；\(\beta_g'(g_*)=\epsilon+O(\epsilon^2)\)
描述这里选定的 RG 参数方向，不能脱离 UV/IR
方向解释稳定性。\(\epsilon=1\) 的应用仍需高阶与重求和分析，不能把小
\(\epsilon\) 证明当成无误差的三维数值。

第 26 章向 IR 的薄壳积分给 \[
\int_{\rm shell}\frac{d^4q}{(2\pi)^4q^4}=\frac{ds}{8\pi^2},
\qquad\delta\lambda=-\frac32\lambda^2\int_{\rm shell}\frac{d^4q}{(2\pi)^4q^4}.
\] 它由三通道的组合因子和独立径向积分产生。第 25.2 节减除泡图对
\(\log\mu\) 的导数则从另一计算得到对应 UV
系数。在匹配的无质量方案下二者相合，方向符号相反。历史回归用独立求导/积分与阶数缩放检查这一关系；直接把一条已知
beta 公式复制到两边不构成独立验证。两系数在一阶相合也不证明完整
Wilsonian action flow 与 Callan--Symanzik 的所有非微扰数据同一。

\hypertarget{ch-26}{%
\chapter{26. Wilsonian RG、有效场论与 Core 桥接}\label{ch-26}}

\hypertarget{ux79efux6389ux5febux6a21ux5f0f}{%
\section{26.1 积掉快模式}\label{ux79efux6389ux5febux6a21ux5f0f}}

带 Euclidean 调节将场分成 \(\phi_<+\phi_>\)，定义 \[
e^{-S_k[\phi_<]}=\int\mathcal D\phi_>\,e^{-S_\Lambda[\phi_<+\phi_>]}
\] 至归一常数。所得 \(S_k\)
一般含无限多允许算符。再标度后按相关、边缘、无关算符组织；``无关''指指定固定点附近的标度，不等于任意精度下可忽略。

取向红外的 \(s=\log(\Lambda/k)\)，四维薄壳
\(\int_{\rm shell}d^4q/(2\pi)^4q^{-4}=ds/(8\pi^2)\)。
在匹配的无质量、规范化方案中
\(d\lambda/ds=-3\lambda^2/(16\pi^2)+\cdots\)，与第25章向 UV
的参数方向相反。

\hypertarget{ux53efux5b8cux5168ux6838ux5bf9ux7684ux91cdux573aux5339ux914d}{%
\section{26.2
可完全核对的重场匹配}\label{ux53efux5b8cux5168ux6838ux5bf9ux7684ux91cdux573aux5339ux914d}}

\protect\hypertarget{p-qeft-01}{}{}

\textbf{P-QEFT-01：二次 Schur 补。} 两个实场的动量二次核 \[
K(p)=\begin{pmatrix}p^2-m^2&-\kappa\\-\kappa&p^2-M^2\end{pmatrix}
\] 中消去重场，得 \[
K_{\rm eff}(p)=p^2-m^2+\frac{\kappa^2}{M^2-p^2}.
\] 证明：重场方程为 \(\chi=(p^2-M^2)^{-1}\kappa\phi\)，代回二次型即得
Schur 补。\(\square\)

若 \(|p^2|<M^2\)， \[
\frac{\kappa^2}{M^2-p^2}
=\frac{\kappa^2}{M^2}\sum_{j=0}^{n}(p^2/M^2)^j+R_n,\quad
|R_n|\le\frac{\kappa^2}{M^2}
\frac{|p^2/M^2|^{n+1}}{1-|p^2/M^2|}.
\] \(\kappa\) 的质量维数是2，不是无量纲耦合。例
\(m=0.7,M=10,\kappa=3,p^2=4\)，
精确修正0.09375；LO、NLO、NNLO分别为0.09、0.0936、0.093744。该例还需
\(m^2M^2>\kappa^2\)
保证原二次质量矩阵正定。能量接近重极点时导数展开失效。

一般 EFT 写成 \[
\mathcal L_{\rm EFT}=\sum_{d,i}\frac{C_i^{(d)}(\mu)}
{\Lambda_*^{d-4}}\mathcal O_i^{(d)}.
\]
算符基、冗余关系、重整化方案、匹配尺度、幂计数和拟合数据都是模型对象的一部分。

\hypertarget{ux540cux4e00ux91cdux5c3aux5ea6ux7684ux6563ux5c04ux5339ux914dux68c0ux9a8c}{%
\subsection{同一重尺度的散射匹配检验}\label{ux540cux4e00ux91cdux5c3aux5ea6ux7684ux6563ux5c04ux5339ux914dux68c0ux9a8c}}

二次混合中的 \(\kappa\)
具有质量维数二。另一个有区分力的匹配基准采用重场交换顶点
\(-g\chi\phi^2/2\)，其中 \(g\)
的质量维数为一；这里须声明局部稳定真空及所需的完整势，以下只计重
\(\chi\) 的树级交换贡献，不把它当作完整全局稳定势的证明。

标准顶点和传播子给 \[
\mathcal M_\chi(s,t,u)=g^2\left[\frac1{M^2-s}+\frac1{M^2-t}+\frac1{M^2-u}\right].
\] 当各 \(|z|/M^2\le r<1\)，逐通道展开为 \[
\mathcal M_\chi=
\frac{3g^2}{M^2}+\frac{g^2(s+t+u)}{M^4}
+\frac{g^2(s^2+t^2+u^2)}{M^6}+\cdots.
\] 截到每通道 \(z^n\) 后，总余项绝对值至多 \[
\frac{3g^2}{M^2}\frac{r^{n+1}}{1-r}.
\] 对质量为 \(m\) 的相同外粒子，壳上有
\(s+t+u=4m^2\)；低维局域项和导数项须在这个运动学约束下匹配，而不是把三个独立变量任意设置。例
\(M=10,g=2,m=1,s=9,t=u=-2.5\) 给交换振幅约
\(0.12200482444\)。这是完整树图与 EFT
级数的比较，独立于上一节仅修改二点核的 Schur 补。

把重场从低能活跃场表移除后，需要匹配 Wilson
系数并在相应理论中运行。一个完整理论的 MOM
耦合在阈值处平滑减弱，不等于已经自动完成这项 EFT
匹配，更不等于所有重场效应精确为零。

\hypertarget{ux4e09ux79cd-rg-ux7684ux7c7bux578b}{%
\section{26.3 三种 RG 的类型}\label{ux4e09ux79cd-rg-ux7684ux7c7bux578b}}

\begin{longtable}[]{@{}
  >{\raggedright\arraybackslash}p{(\columnwidth - 4\tabcolsep) * \real{0.2500}}
  >{\raggedright\arraybackslash}p{(\columnwidth - 4\tabcolsep) * \real{0.3400}}
  >{\raggedright\arraybackslash}p{(\columnwidth - 4\tabcolsep) * \real{0.4100}}@{}}
\toprule\noalign{}
\begin{minipage}[b]{\linewidth}\raggedright
变换
\end{minipage} & \begin{minipage}[b]{\linewidth}\raggedright
输入与输出
\end{minipage} & \begin{minipage}[b]{\linewidth}\raggedright
保持的内容
\end{minipage} \\
\midrule\noalign{}
\endhead
\bottomrule\noalign{}
\endlastfoot
Core 对象粗化 & 声明协议响应及其预测因子之间的映射 &
指定未来任务的响应或有界误差 \\
量子过程粗化 & 多时 instrument 插槽、初态与记忆过程 &
允许干预下的记录概率 \\
Wilsonian RG & 带调节的作用量、场与耦合 &
低能相关函数或匹配振幅至声明精度 \\
\end{longtable}

它们可以构造桥接，但定义域、陪域和应保持的数据不同。尤其固定点、谱隙、Markov
阶数、对象最小性是各自的概念，没有一般定理把它们无条件一一对应。

\protect\hypertarget{p-qeft-02}{}{}

\textbf{P-QEFT-02：受限的实时桥。}
给定同一个系统---环境划分、初态、联合酉演化、允许插槽干预和末端读出，直接部分迹与将环境积分为
Schwinger--Keldysh influence functional
的计算，如果均精确且相应求和/积分与极限合法，则给出相同的协议概率。

\textbf{证明。} 在有限维、有限调节实现中，完整记录 \(\mathbf o\)
的概率可先写成联合算符的线性配对：在同一个初态上，按时间次序插入每一步联合酉演化和相应的系统
CP
分支，最后对全部终态求迹。将矩阵元展开，演化振幅与其共轭分别形成闭时路径的两条腿；环境初态和末端迹提供相同的环境指标缩并。先缩并环境得到
influence
functional，最后再求系统指标，或先组成联合算符再部分迹，只是同一有限和的不同顺序。受控的连续极限要求额外收敛与换序条件，然后才能把该等式传到极限。\(\square\)

相同响应可以再按 Core 的预测等价关系取商。若近似 influence
functional、截断记忆或丢弃某些干预，须给相应概率误差，不能继续使用精确相等。Euclidean
快模式积分自身不是一般 CPTP
时间演化；初始相关、外部算符插入、规范约束和非局域记忆都可能使天真的通道解释失败。共享一种``粗化''直觉不替代这些物理接口。

\hypertarget{ch-27}{%
\chapter{27. QED：规范、BRST 与 Ward 接口}\label{ch-27}}

\hypertarget{ux4f5cux7528ux91cfux548cux5bf9ux79f0ux6027}{%
\section{27.1
作用量和对称性}\label{ux4f5cux7528ux91cfux548cux5bf9ux79f0ux6027}}

\[
D_\mu=\partial_\mu+ieA_\mu,\quad
\mathcal L=-\tfrac14F^2+\bar\psi(i\not D-m)\psi,
\quad \psi\mapsto e^{-ie\alpha}\psi,\quad A_\mu\mapsto A_\mu+\partial_\mu\alpha.
\] 相互作用为 \(-e\bar\psi\gamma^\mu A_\mu\psi\)，顶点为
\(-ie\gamma^\mu\)。 线性协变规范增加
\(-(\partial A)^2/(2\xi)\)，自由光子传播子 \[
D_{\mu\nu}(q)=\frac{-i}{q^2+i0}
\left[\eta_{\mu\nu}-(1-\xi)\frac{q_\mu q_\nu}{q^2+i0}\right].
\] Abelian Faddeev--Popov 行列式的局部非零模式不依赖 \(A\)，相应 ghost
在普通 QED 微扰中解耦；全局零模和边界仍要另行处理。

辅助场 \(B\) 下，取
\(sA=\partial c,\ s\psi=-iec\psi,\ sc=0,\ s\bar c=B,\ sB=0\)。
规范固定与 ghost 项可写为 \(s[\bar c(\partial A+\xi B/2)]\)；消去 \(B\)
得上述规范固定项。 BRST
商是在声明的量子一致性条件下构造物理扇区，不是将任意非正定态当作概率态。

\hypertarget{wardtakahashi-ux6052ux7b49ux5f0f}{%
\section{27.2 Ward--Takahashi
恒等式}\label{wardtakahashi-ux6052ux7b49ux5f0f}}

\textbf{K-QWARD-01。} 以不含整体电荷的完整顶点 \(\Gamma^\mu\)
为约定，在保持向量规范对称的正规化/重整化下， \[
q_\mu\Gamma^\mu(p+q,p)=S^{-1}(p+q)-S^{-1}(p).
\]
证明机制是对带源路径积分作无异常的局域相位变量替换，对费米源微分，然后去外腿；接触项给出右端两个逆传播子。这里明确调用标准完整恒等式，不能以树级检查替代所有圈阶证明。

树级时直接为
\(q_\mu\gamma^\mu=\not p+\not q-m-(\not p-m)\)。真空极化满足
\(q_\mu\Pi^{\mu\nu}=0\)。该恒等式约束纵向部分；它不独自确定所有有限横向顶点形状因子。

令 \(d=4-2\epsilon_{\rm DR}\)， \[
\psi_0=Z_2^{1/2}\psi,\quad A_0=Z_3^{1/2}A,\quad
e_0=\mu^{\epsilon_{\rm DR}}Z_e e,\quad
Z_1=Z_eZ_2Z_3^{1/2}.
\] Ward 恒等式给 \(Z_1=Z_2\)，从而 \(Z_e=Z_3^{-1/2}\)。

\hypertarget{ux4e00ux5708-uv-ux95edux5408ux57faux51c6}{%
\section{27.3 一圈 UV
闭合基准}\label{ux4e00ux5708-uv-ux95edux5408ux57faux51c6}}

这里冻结 Feynman 规范、一个单位电荷 Dirac 费米子及 MS-bar 的一圈 UV pole
记法： \[
a=\frac{\alpha}{4\pi\bar\epsilon},\qquad
\frac1{\bar\epsilon}=\frac1{\epsilon_{\rm DR}}-\gamma_E+\log4\pi.
\] 目标是独立计算反项的一致性，不是给任意规范下的全部有限自能和顶点。

费米自能分子先用
\(\gamma_\mu\gamma_\alpha\gamma^\mu=(2-d)\gamma_\alpha\) 化简；Feynman
参数平移后，UV 系数积分为 \[
\Sigma_{\rm UV}(p)=a\int_0^1(-2x\not p+4m)dx
=a(-\not p+4m).
\] 若逆传播子约定为 \[
\Gamma^{(2)}=\not p-m+\delta Z_2\not p
-\delta m_{\rm direct}\,m-\Sigma_{\rm loop},
\] 则分别消掉 \(\not p\) 与 \(m\) 的 UV 残差给 \[
\delta Z_2=-a,\qquad\delta m_{\rm direct}=-4a.
\] 乘法质量重整化 \(m_0=Z_m m\) 满足 \[
1+\delta m_{\rm direct}=Z_2Z_m,
\qquad\delta Z_m=\delta m_{\rm direct}-\delta Z_2=-3a.
\] 所以 direct mass coefficient 与 \(Z_m-1\)
并不相同；把它们混用会制造并不存在的 Ward 冲突。

顶点的 logarithmic UV 部分还有一条独立检查。对称 loop-momentum
张量积分给度规缩并，在 pole 的四维留数上 \[
\frac14\gamma_\alpha\gamma_\rho\gamma^\mu
\gamma^\rho\gamma^\alpha=\gamma^\mu.
\] 两次使用四维的 \(\gamma_\rho\gamma^\mu\gamma^\rho=-2\gamma^\mu\)
即得此式。\(d-4\)
的差别可影响有限部分，不能用该四维留数替代完整维数依赖顶点。它给出一圈
UV 顶点修正 \(\Lambda^\mu_{\rm UV}=a\gamma^\mu\)，故
\(\delta Z_1=-a\)。与从自能独立得到的 \(\delta Z_2\)
相比，确认此规范、此阶的 \(\delta Z_1=\delta Z_2\)。

光子 UV 系数来自 \[
\frac{2\alpha}{\pi\bar\epsilon}\int_0^1x(1-x)dx
=\frac{\alpha}{3\pi\bar\epsilon},
\] 因此 \(\delta Z_3=-\alpha/(3\pi\bar\epsilon)\)。再由完整 Ward 关系
\(Z_e=Z_3^{-1/2}\) 得对应 charge
pole。历史回归分别积分自能系数、缩并顶点 Dirac 张量和检查 photon
系数，不能把``先设 \(Z_1=Z_2\) 再比较''当成独立圈图计算。

以上不覆盖任意 \(\xi\) 的全部 finite parts、壳上 \(Z_2\)
的红外结构或多圈精密 QED。树级 Ward、此处一圈 UV checkpoint 与第 27.2
节采用的完整 Ward--Takahashi 定理是三种证据，必须分别登记。

\hypertarget{ch-28}{%
\chapter{28. QED 可计算应用：散射、running 与磁矩}\label{ch-28}}

\hypertarget{ux7535ux5b50ux6b63ux7535ux5b50ux6e6eux706d}{%
\section{28.1
电子---正电子湮灭}\label{ux7535ux5b50ux6b63ux7535ux5b50ux6e6eux706d}}

对 \(e^-e^+\to\mu^-\mu^+\)，忽略两种轻子质量、只取一光子树图， \[
\mathcal M=\frac{e^2}{s}
[\bar v(p_2)\gamma^\mu u(p_1)]
[\bar u(k_1)\gamma_\mu v(k_2)].
\] \protect\hypertarget{p-qqed-01}{}{}

\textbf{P-QQED-01。} 用自旋完备关系和 \[
\operatorname{Tr}(\not a\gamma^\mu\not b\gamma^\nu)
=4(a^\mu b^\nu+a^\nu b^\mu-\eta^{\mu\nu}a\cdot b)
\] 求迹，初态平均除以4，得 \[
\overline{|\mathcal M|^2}=\frac{2e^4}{s^2}(t^2+u^2)
=e^4(1+\cos^2\theta).
\] 代入二体相空间并积分， \[
\frac{d\sigma}{d\Omega}=\frac{\alpha^2}{4s}(1+\cos^2\theta),
\qquad \sigma=\frac{4\pi\alpha^2}{3s}.
\] 角积分为 \(16\pi/3\)，故最后系数可直接核对。\(\square\)

纵向传播子项由外腿 Dirac 方程和电流守恒消失，树截面与 \(\xi\)
无关。有限质量、Z
交换、辐射修正与探测器接受度不在这个基准内。固定复现输入
\(\alpha^{-1}=137.035999084\)、\(\sqrt s=10\) 能量单位时，
\(\sigma=2.2305875\times10^{-6}\)
逆能量平方；此处不宣称采用最新常数测量或精密实验预测。

\hypertarget{ux771fux7a7aux6781ux5316ux4e0eux53bbux8026}{%
\section{28.2
真空极化与去耦}\label{ux771fux7a7aux6781ux5316ux4e0eux53bbux8026}}

以 spacelike \(Q^2>0\) 的正 screening 函数定义 \[
\widehat\Pi(Q^2)=\frac{2\alpha}{\pi}\int_0^1dx\,x(1-x)
\log\!\left[1+\frac{Q^2}{m^2}x(1-x)\right],\quad
\alpha_{\rm eff}=\frac{\alpha}{1-\widehat\Pi}.
\] 这是特定减除与符号约定。小 \(a_Q=Q^2/m^2\) 展开给
\(\widehat\Pi=\alpha a_Q/(15\pi)+O(a_Q^2)\)； 大 \(a_Q\) 给 \[
\widehat\Pi=\frac{\alpha}{3\pi}\left(\log a_Q-\frac53\right)+O(a_Q^{-1}).
\] 前者由 \(\int x^2(1-x)^2dx=1/30\)，后者由 \(\int x(1-x)dx=1/6\)
与对数矩得到。

\protect\hypertarget{p-qqed-02}{}{}

\textbf{P-QQED-02：阈值 beta。} 对上式作 \(\log Q\) 导数，至
\(O(\alpha^2)\)： \[
\frac{d\alpha_{\rm eff}}{d\log Q}
=\frac{4\alpha^2}{\pi}\int_0^1dx\,
\frac{a_Qx^2(1-x)^2}{1+a_Qx(1-x)}.
\] 红外趋零，UV 趋 \(2\alpha^2/(3\pi)\)。多个活跃 Dirac 场乘各自
\(N_cQ_f^2\) 并求和，阈值须各按质量处理。\(\square\) 一圈 Landau
极点是外推警告，不是已经证明存在的物理能量奇点。

\hypertarget{schwinger-ux78c1ux77e9ux4e0eux7ea2ux5916ux8fb9ux754c}{%
\section{28.3 Schwinger
磁矩与红外边界}\label{schwinger-ux78c1ux77e9ux4e0eux7ea2ux5916ux8fb9ux754c}}

壳上顶点分解 \[
\bar u(p')\Gamma^\mu u(p)=\bar u(p')
\left[\gamma^\mu F_1(q^2)
+\frac{i\sigma^{\mu\nu}q_\nu}{2m}F_2(q^2)\right]u(p),
\quad a_\ell=F_2(0).
\] 这里的标准一圈来源是在外腿壳上、以 \(-ie\)
从完整顶点提出之后，按固定动量流向写出的修正积分，例如 \[
\Lambda^\mu(p',p)=-ie^2\int\frac{d^d\ell}{(2\pi)^d}
\frac{\gamma_\rho(\not p'-\not\ell+m)\gamma^\mu
(\not p-\not\ell+m)\gamma^\rho}
{[\ell^2+i0][(p'-\ell)^2-m^2+i0][(p-\ell)^2-m^2+i0]}.
\] dimensional regularization 中还须保留相应的
\(\mu^{2\epsilon_{\rm DR}}\)，需要时先引入相容的红外处理。Feynman
参数化、loop shift、Gordon 分解和 Pauli 结构投影把 \(q^2\to0\) 的有限
\(F_2\) 部分化为下式。这个解析 reduction 是标准 QED
输入；后面的二维数值积分不等于机器重新证明了这里的全部换序、正则化和极限。

\textbf{K-QQED-01。} 标准一圈顶点积分经 Feynman 参数与 Pauli
投影给出有限积分 \[
F_2(0)=\frac{\alpha}{\pi}
\int_0^1dz\int_0^{1-z}dy\,\frac{z}{1-z}
=\frac{\alpha}{2\pi}.
\] 内层积分先消去 \(1-z\)，再积分 \(z\)
即得结果。本文独立检查的是该参数积分与数值，不把它说成已用机器从完整圈图重新证明全部正则化和壳上极限。

虚光子圈与实软辐射分别可能红外发散；物理包容观测需冻结能量/角分辨率并结合实虚贡献。向量规范异常消除不意味着轴流异常不存在。QED
扇区尚未给 UEOT 提供新的独立常数预测。

\hypertarget{ch-29}{%
\chapter{29. 非 Abelian 规范场与标准模型候选}\label{ch-29}}

\hypertarget{yangmills-ux7684ux5b9aux4e49ux4e0eux8ba1ux7b97ux8d23ux4efb}{%
\section{29.1 Yang--Mills
的定义与计算责任}\label{yangmills-ux7684ux5b9aux4e49ux4e0eux8ba1ux7b97ux8d23ux4efb}}

取 Hermitian 生成元 \([T^a,T^b]=if^{abc}T^c\)，
\(\operatorname{Tr}_R(T^aT^b)=T(R)\delta^{ab}\)。 为与 QED 的
\(D=\partial+igA\) 一致，定义 \[
F_{\mu\nu}=\partial_\mu A_\nu-\partial_\nu A_\mu+ig[A_\mu,A_\nu],
\quad [D_\mu,D_\nu]=igF_{\mu\nu},
\quad\mathcal L_{\rm YM}=-\tfrac12\operatorname{Tr}(F_{\mu\nu}F^{\mu\nu})
\] （基本表示迹归一化 \(T=1/2\)）。此约定分量非线性项是
\(-g f^{abc}A_\mu^bA_\nu^c\)；换成 \(D=\partial-igA\)
时符号同步变换，不能只换其中一处。

非 Abelian 场自耦产生三、四规范玻色子顶点。ghost 不再解耦；规范固定与
BRST 变换用协变导数和 Lie 括号构造，幂零性依赖 Jacobi 恒等式。量子约束为
Slavnov--Taylor 恒等式，不是对每个顶点都套 QED 的 \(Z_1=Z_2\)。微扰 BRST
不解决全局 Gribov 问题。

\textbf{K-QYM-01。} 对 Dirac 费米子与复标量， \[
\beta(g)=-\frac{g^3}{16\pi^2}b_0+O(g^5),\quad
b_0=\frac{11}{3}C_A-\frac43\sum_{\rm Dirac}T(R_f)
-\frac13\sum_{\rm complex}T(R_s).
\] Weyl 费米子贡献是 Dirac 的一半。基本表示 QCD 六味给 \(b_0=7\)。
渐近自由不是禁闭或质量隙证明。强耦合低能预测通常需要格点、有效理论或其他非微扰方法及其误差控制。

\hypertarget{ux6807ux51c6ux6a21ux578bux662fux4e00ux79cdux5177ux4f53ux5019ux9009}{%
\section{29.2
标准模型是一种具体候选}\label{ux6807ux51c6ux6a21ux578bux662fux4e00ux79cdux5177ux4f53ux5019ux9009}}

声明规范 Lie 代数
\(\mathfrak{su}(3)_C\oplus\mathfrak{su}(2)_L\oplus\mathfrak u(1)_Y\)。
全局群还可能取与表示兼容的离散商；局域拉氏量不足以决定线算符与全部全局扇区。下表采用全部左手
Weyl 约定和 \(Q=T_3+Y\)。

\begin{longtable}[]{@{}
  >{\raggedright\arraybackslash}p{(\columnwidth - 6\tabcolsep) * \real{0.1800}}
  >{\raggedright\arraybackslash}p{(\columnwidth - 6\tabcolsep) * \real{0.3700}}
  >{\raggedleft\arraybackslash}p{(\columnwidth - 6\tabcolsep) * \real{0.1300}}
  >{\raggedleft\arraybackslash}p{(\columnwidth - 6\tabcolsep) * \real{0.3200}}@{}}
\toprule\noalign{}
\begin{minipage}[b]{\linewidth}\raggedright
场
\end{minipage} & \begin{minipage}[b]{\linewidth}\raggedright
\(SU(3)_C,SU(2)_L\)
\end{minipage} & \begin{minipage}[b]{\linewidth}\raggedleft
\(Y\)
\end{minipage} & \begin{minipage}[b]{\linewidth}\raggedleft
每代 Weyl 分量数
\end{minipage} \\
\midrule\noalign{}
\endhead
\bottomrule\noalign{}
\endlastfoot
\(Q\) & \(3,2\) & \(1/6\) & 6 \\
\(U^c\) & \(\bar3,1\) & \(-2/3\) & 3 \\
\(D^c\) & \(\bar3,1\) & \(1/3\) & 3 \\
\(L\) & \(1,2\) & \(-1/2\) & 2 \\
\(E^c\) & \(1,1\) & 1 & 1 \\
\(H\) & \(1,2\) & \(1/2\) & 复标量，非 Weyl \\
\end{longtable}

规范动能、费米动能、Higgs 动能与势、Yukawa
项组成重整化维数不超过4的普通主干；还可有拓扑 theta 项。强 CP
参数不是用``最小对象''自动置零。这个场表是领域输入，不是本书从 Core
推出的唯一粒子表。

\hypertarget{ux7cbeux786eux5f02ux5e38ux65b9ux7a0b}{%
\section{29.3 精确异常方程}\label{ux7cbeux786eux5f02ux5e38ux65b9ux7a0b}}

对一般 \(N_c\) 与相同多重态类型，局域异常要求 \[
2y_Q+y_U+y_D=0,\quad N_cy_Q+y_L=0,
\] \[
2N_cy_Q+N_cy_U+N_cy_D+2y_L+y_E=0,
\] \[
2N_cy_Q^3+N_cy_U^3+N_cy_D^3+2y_L^3+y_E^3=0.
\] 纯色立方异常由双重态的两个基本表示与两个反基本表示抵消。\(SU(2)\)
无微扰立方异常，但其 Witten 全局异常要求基本双重态总数为偶数；这里为
\(N_f(N_c+1)\)。

\protect\hypertarget{p-qsm-01}{}{}

\textbf{P-QSM-01：异常解与非唯一性。} 三个线性方程给
\(y_L=-N_cy_Q,\ y_D=-2y_Q-y_U,\ y_E=2N_cy_Q\)，立方式因式分解为 \[
6N_cy_Q[(N_c-1)y_Q-y_U][(N_c+1)y_Q+y_U].
\] 展开或代入即可验证。故除了 SM 型分支和上下型互换，还存在 \(y_Q=0\)
的退化分支，异常本身不唯一选择超荷。\(\square\)

加入普通 Yukawa 规范不变性 \(QHU^c,QH^\dagger D^c,LH^\dagger E^c\)，得到
\(y_H=N_cy_Q\)。以 \(y_H=1/2\) 归一化， \[
Y_Q=\frac1{2N_c},\quad
Y_U=-\frac{N_c+1}{2N_c},\quad
Y_D=\frac{N_c-1}{2N_c},\quad Y_L=-\frac12,\quad Y_E=1.
\] 任意满足其他条件的 \(N_c\) 都有这支局域解；并未选出3。加入 \(N^c\)
可重新开放 Abelian 混合方向，必须重解全部条件。

\hypertarget{ux5c40ux57dfux5168ux5c40ux4e0eux8eabux4efdux8fd0ux8f93}{%
\section{29.4
局域、全局与身份运输}\label{ux5c40ux57dfux5168ux5c40ux4e0eux8eabux4efdux8fd0ux8f93}}

微扰异常系数为零只是第一层。全局规范群、spin
结构、离散对称性、背景丛、边界及可能 inflow 都影响完整异常问题。

\textbf{S-QGOIT：Gauge-Orbit Identity Transport。}
若一个规范变换已被物理实现认定为表示冗余，完整物理响应须能一致地下降到其轨道商。手征费米积分的
determinant/Pfaffian line 若存在未抵消的曲率或 holonomy
障碍，孤立候选不能作为该动态规范理论的一致实现；可用合法的额外物质、反项或
inflow 完成。

这是标准异常障碍在 UEOT 身份运输框架中的位置，不是把任意 Core holonomy
等同于规范异常。未规范化全局对称性的 't Hooft
异常也不等于动态规范冗余不一致。数学来源可参见
\href{https://arxiv.org/abs/1607.01873}{Yonekura 的 determinant-line
论述} 和 \href{https://arxiv.org/abs/1808.00009}{Dai--Freed
粒子物理分析}。

\hypertarget{ch-30}{%
\chapter{30. Higgs、Yukawa、味与中微子}\label{ch-30}}

\hypertarget{ux7535ux5f31ux7834ux7f3aux7684ux663eux5f0fux8f7dux4f53}{%
\section{30.1
电弱破缺的显式载体}\label{ux7535ux5f31ux7834ux7f3aux7684ux663eux5f0fux8f7dux4f53}}

取单个复双重态 \[
\begin{gathered}
V(H)=-\mu_H^2H^\dagger H+\lambda_H(H^\dagger H)^2,
\\ \lambda_H>0,\\ \mu_H^2>0,\\
H=\frac1{\sqrt2}\binom{0}{v+h},\\ v^2=\mu_H^2/\lambda_H.
\end{gathered}
\] 这是特定真空邻域的幺正规范参数化，不是全局规范固定存在性的证明。

\protect\hypertarget{p-qhig-01}{}{}

\textbf{P-QHIG-01：质量矩阵。} 展开 \((D_\mu H)^\dagger D^\mu H\) 得 \[
m_W^2=\frac{g^2v^2}{4},\qquad
M_{\rm neutral}^2=\frac{v^2}{4}
\begin{pmatrix}g^2&-gg'\\-gg'&g'^2\end{pmatrix}
\] （中性场基可取与该约定一致的 \(W^3,B\)）。
其行列式为零，另一个本征值为 \((g^2+g'^2)v^2/4\)。定义 \[
\tan\theta_W=g'/g,\quad
A=\sin\theta_W W^3+\cos\theta_W B,\quad
Z=\cos\theta_W W^3-\sin\theta_W B,
\] 则 \(m_A=0,m_Z=v\sqrt{g^2+g'^2}/2\)，
\(e=g\sin\theta_W=g'\cos\theta_W\)。 势的二阶展开给
\(m_h^2=2\lambda_Hv^2\)。 树级
\(\rho=m_W^2/(m_Z^2\cos^2\theta_W)=1\)。\(\square\)

多个适当双重态也可能保留树级 \(\rho=1\)，故此关系不唯一选择单
Higgs；``拓扑平凡标量''也不决定双重态、势或真空尺度。

\hypertarget{ux5473ux7684ux81eaux7531ux53c2ux6570ux4e0e-cp}{%
\section{30.2 味的自由参数与
CP}\label{ux5473ux7684ux81eaux7531ux53c2ux6570ux4e0e-cp}}

四分量记法 \[
\mathcal L_Y=-\bar Q_LY_dd_RH-\bar Q_LY_uu_R\widetilde H
-\bar L_LY_ee_RH+\mathrm{h.c.},\quad \widetilde H=i\sigma_2H^*.
\] 破缺后 \(M_f=vY_f/\sqrt2\)。奇异值分解给 \[
U_{fL}^\dagger M_fU_{fR}=D_f\ge0,\qquad
V_{\rm CKM}=U_{uL}^\dagger U_{dL}.
\] 中性味对角性、带电流混合与 GIM
抵消由该结构计算。退化质量会扩大未固定的基变换群。

\protect\hypertarget{p-qflav-01}{}{}

\textbf{P-QFLAV-01：一般非退化 \(N\) 代夸克计数。} 两个复 Yukawa 有
\(4N^2\) 个实参数；\(U(N)_Q\times U(N)_u\times U(N)_d\) 基变换除去
\(3N^2-1\) 个，因为共同重子相位不改 Yukawa。剩余 \(N^2+1\)，分解为 \[
2N\text{ 个质量}+\frac{N(N-1)}2\text{ 个角}
+\frac{(N-1)(N-2)}2\text{ 个 CP 相位}.
\] 三代给6质量、3角、1相位。\(\square\)

该计数只涉及 Yukawa 味参数，不把强 CP 参数包含在这10个数中。非零 KM 型
CP 要求至少3代；``至少3''不能推出``恰好3''。有物理 CP
还要求非退化结构和非零重相位不变量，例如
\(J=\operatorname{Im}(V_{ij}V_{kl}V^*_{il}V^*_{kj})\) 的相应组合。

\hypertarget{ux4e2dux5faeux5b50ux4e0dux540cux5b9eux73b0ux4e0dux6df7ux540c}{%
\section{30.3
中微子：不同实现不混同}\label{ux4e2dux5faeux5b50ux4e0dux540cux5b9eux73b0ux4e0dux6df7ux540c}}

增加 \(N_R\) 可产生 Dirac 质量。Majorana EFT 可取明确的规范： \[
\mathcal L_5=-\frac{C_{ij}}{2\Lambda_*}(L_iH)(L_jH)+\mathrm{h.c.},
\qquad m_\nu=\frac{v^2}{2\Lambda_*}C
\] （自旋和 \(SU(2)\)
缩并隐含；\(C=C^T\)）。系数中的2取决于算符定义，必须与质量公式一致。

\protect\hypertarget{p-qnu-01}{}{}

\textbf{P-QNU-01：Type-I seesaw。} 在 \(\|M_D M_R^{-1}\|\ll1\)、\(M_R\)
可逆的重轻分离中， \[
\mathcal M=\begin{pmatrix}0&M_D\\M_D^T&M_R\end{pmatrix},
\quad m_\nu^{\rm eff}=-M_DM_R^{-1}M_D^T
\] 为零动量二次核的 Schur
补；动力学与混合矩阵另有高阶修正。复对称质量的物理质量由 Takagi
分解给出，不能拿一般复矩阵普通本征值代替。\(\square\)

真空超相对论振荡近似为 \[
P_{\alpha\to\beta}(L)=
\left|\sum_i U_{\beta i}U_{\alpha i}^*
e^{-im_i^2L/(2E)}\right|^2,
\]
要求适当制备、相干传播与探测近似。物质效应、波包退相干和非幺正扩展要另外加入。振荡概率并不唯一决定
Dirac/Majorana 性质或高能 seesaw 载体。

\hypertarget{ux5473ux6a21ux578bux4e0eux539fux59cbux5d4cux5957ux76f4ux89c9}{%
\section{30.4
味模型与原始``嵌套''直觉}\label{ux5473ux6a21ux578bux4e0eux539fux59cbux5d4cux5957ux76f4ux89c9}}

离散深度 \(d=1,2,3\) 只有在独立给出 \(d\mapsto Y\)
后才有可检验内容。例如等步指数 \(m(d)=Me^{\kappa d}\) 必须满足
\(m_2^2=m_1m_3\)；这是一条可反驳的关系，不是``深度越大越重''的同义重述。

若写 \(Y_{ij}=c_{ij}\epsilon^{n_{ij}}\)，其结构已是 Froggatt--Nielsen
型。UEOT
要增加物理内容，必须独立选出整数、\(\epsilon\)、系数约束，并预测未用于选择载体的味观测。本文不把可任意拟合的重命名视为质量理论。

\hypertarget{ch-31}{%
\chapter{31. 三代与三色：条件性定理及前提简化}\label{ch-31}}

使用提示：本章保留低能弱渐近自由前提；第36章证明另一套高能渐近自由条件不能自动推出此前提，并给出五色反模型。两章须共同引用。

本章保留可访问主线 v30.11
的有效算术，并给出一个不需要``原始包''的加强版本。它仍不是一般 UEOT Core
的三代定理。

\hypertarget{ux590dux5236ux4e0dux53d8ux6027ux969cux788d}{%
\section{31.1
复制不变性障碍}\label{ux590dux5236ux4e0dux53d8ux6027ux969cux788d}}

\protect\hypertarget{p-qg3-01}{}{}

\textbf{P-QG3-01。} 若公理类 \(\mathcal A\) 含模型 \(T\)，且对某个
\(r>1\) 的家族复制封闭，复制后仍满足全部公理且
\(N_f(R_rT)=rN_f(T)\)，则该类不能唯一选出固定正整数代数。

\textbf{证明。} 若唯一值为 \(N_0>0\)，则同类复制模型取值
\(rN_0\ne N_0\)，矛盾。\(\square\)

每代独立的异常抵消、一般 Born
运动学和对复制封闭的对象证书，因而不能单独解释3。成功选择器必须含真正破坏复制自由的前提；上界、跨代异常或耦合结构都可能承担这个角色。

\hypertarget{ux6a2aux5411ux89c4ux8303ux7236ux7cfbux7edfux7684ux56faux5b9aux8868ux793aux7c7b}{%
\section{31.2
横向规范父系统的固定表示类}\label{ux6a2aux5411ux89c4ux8303ux7236ux7cfbux7edfux7684ux56faux5b9aux8868ux793aux7c7b}}

\textbf{H-QG3-1（HGP 与表示配置）。} 存在动态手征横向/扩展弱规范载体
\(G_H\)，含复表示 \(R\) 且立方异常系数
\(A(R)=a\ne0,A(\bar R)=-a\)。每代轻子是一个无色 \(R\)；夸克带 \(N_c\)
色重数，其中 \(k\) 代取 \(R\)，其余 \(N_f-k\) 代取
\(\bar R\)。本条件不仅是``存在某个父对象''，还明确规定了它的手征表示配置及无其他未计入的该立方异常贡献。

\protect\hypertarget{p-qg3-02}{}{}

\textbf{P-QG3-02：全部整数解。} \[
A_H/a=N_f+N_c(2k-N_f)=0
\iff 2N_ck=(N_c-1)N_f.
\] 令 \(d=\gcd(2N_c,N_c-1)=\gcd(2,N_c-1)\)，则所有正解为 \[
k=t\frac{N_c-1}{d},\qquad N_f=t\frac{2N_c}{d},\qquad
t\in\mathbb Z_{>0}.
\]
证明：先除去最大公因数，剩余两系数互素，整除性迫使两变量取上述共同倍数。奇
\(N_c\) 时 \(d=2\)，偶 \(N_c\) 时 \(d=1\)。\(\square\)

``原始包''定义为 \(t=1\)；一般 Core
原子性并不强迫宇宙只实现一个此包。该异常配置与331模型的经典代数机制有明确先例，不是本书原创：\href{https://arxiv.org/abs/hep-ph/9206242}{Pisano--Pleitez
原始论文}、\href{https://link.aps.org/doi/10.1103/PhysRevD.50.R34}{含右手中微子的扩展}。

\hypertarget{ux4e0dux518dux9700ux8981ux539fux59cbux5305ux524dux63d0ux7684ux6761ux4ef6ux9009ux62e9}{%
\section{31.3
不再需要原始包前提的条件选择}\label{ux4e0dux518dux9700ux8981ux539fux59cbux5305ux524dux63d0ux7684ux6761ux4ef6ux9009ux62e9}}

另外明确：

\begin{itemize}
\tightlist
\item
  \textbf{H-QG3-2：} 色群取 \(SU(N_c)\) 的真正复基本表示，故
  \(N_c\ge3\)；排除伪实 \(SU(2)\) 基本表示是物理类选择。
\item
  \textbf{H-QG3-3：}
  在所声明的全部普通代均活跃、规范结构和物质内容固定的微扰窗口，要求弱
  \(SU(2)_L\) 的一圈系数为渐近自由方向。普通费米子贡献给必要条件 \[
  N_f(N_c+1)<22.
  \] 标量及额外普通带弱荷物质进一步
  screening；改变规范群或引入新的矢量结构则须重算，不能无条件沿用此界。
\end{itemize}

\protect\hypertarget{p-qg3-03}{}{}

\textbf{P-QG3-03：去原始性前提的选择定理。} 在 H-QG3-1、2、3
的交集中，全部正整数异常解唯一为 \[
(N_c,N_f,k)=(3,3,1).
\] 无需另加 \(t=1\)。

\textbf{证明。} 奇数 \(N_c\ge3\) 时 \(N_f=tN_c\)，故
\(tN_c(N_c+1)<22\)。\(N_c\ge5\) 的最小左端为30，不可能；\(N_c=3\) 时
\(12t<22\)，正整数 \(t\) 只能为1，进而 \(N_f=3,k=1\)。偶数 \(N_c\ge4\)
时 \(N_f=2tN_c\)，最小左端40，亦不可能。\(\square\)

这是在固定模型类中删除冗余前提的逻辑优化。它不证明
HGP，不证明该表示类唯一，也不证明自然界所有弱有效理论必须渐近自由。若放宽复色条件，\((2,4,1)\)
通过上述代数界；若去掉弱上界，\(t>1\)
和更大奇色数重新出现。若加入额外横向异常抵消物质，整数方程本身会改变。

对候选点再给一个 Higgs 双重态，完整普通弱系数为 \[
b_2=\frac{22}{3}-\frac{N_f(N_c+1)}3-\frac16
=\frac{19}{6}>0.
\] 这仅核对了一个允许的弱物质配置；不是对所有 HGP
高能完成的充分存在性证明。

\hypertarget{ux5b8cux6574331ux5b58ux5728ux6027ux4e0eux5f02ux5e38ux6e05ux5355}{%
\section{31.4
完整331存在性与异常清单}\label{ux5b8cux6574331ux5b58ux5728ux6027ux4e0eux5f02ux5e38ux6e05ux5355}}

在 \(SU(3)_C\times SU(3)_L\times U(1)_X\) 的一个 RHN
型分支中，全左手表示为：

\begin{longtable}[]{@{}
  >{\raggedleft\arraybackslash}p{(\columnwidth - 2\tabcolsep) * \real{0.3000}}
  >{\raggedright\arraybackslash}p{(\columnwidth - 2\tabcolsep) * \real{0.7000}}@{}}
\toprule\noalign{}
\begin{minipage}[b]{\linewidth}\raggedleft
多重数
\end{minipage} & \begin{minipage}[b]{\linewidth}\raggedright
表示 \((C,L,X)\)
\end{minipage} \\
\midrule\noalign{}
\endhead
\bottomrule\noalign{}
\endlastfoot
3 & \((1,3,-1/3)\) \\
3 & \((1,1,1)\) \\
1 & \((3,3,1/3)\) \\
2 & \((3,\bar3,0)\) \\
4 & \((\bar3,1,-2/3)\) \\
5 & \((\bar3,1,1/3)\) \\
\end{longtable}

计算纯 \(C^3,L^3\)、混合 \(C^2X,L^2X\)、引力混合 \(X\) 与
\(X^3\)，分别给出零。基本/反基本对立方非 Abelian 异常取相反符号，对二次
Dynkin 指数取相同 \(1/2\)；\(U(1)\)
项乘其余表示维数。这些规则足以重现附录的有理数验证。该表说明局域异常自由的载体确实存在，不是自然界选择它的证据，也不完成其全部全局异常和现象学检查。

\hypertarget{ux56faux5b9aux666eux901aux7535ux8377ux7684ux53e6ux4e00ux6761ux6761ux4ef6ux5206ux652f}{%
\section{31.5
固定普通电荷的另一条条件分支}\label{ux56faux5b9aux666eux901aux7535ux8377ux7684ux53e6ux4e00ux6761ux6761ux4ef6ux5206ux652f}}

对一般 \(N_c\)，取 \(Q=T_3+\beta_8T_8+X\)，固定普通电荷
\(Q_\nu=0,Q_e=-1,Q_u=2/3,Q_d=-1/3\)。则 \[
X_L=-\tfrac12-\frac{\beta_8}{2\sqrt3},\quad
X_{Q+}=\tfrac16-\frac{\beta_8}{2\sqrt3},\quad
X_{Q-}=\tfrac16+\frac{\beta_8}{2\sqrt3},
\] 第三分量电荷为 \[
q_E=-\tfrac12-\frac{\sqrt3\beta_8}{2},\quad
q_{J+}=\tfrac16-\frac{\sqrt3\beta_8}{2},\quad
q_{J-}=\tfrac16+\frac{\sqrt3\beta_8}{2}.
\] 为在任意 \(\beta_8\) 下声明完整手征表，应加入全部相应右手场的左手共轭
singlet，\textbf{包括带电的 \(E^c\)}，不能只列 \(e^c,u^c,d^c,J^c\)。当
\(q_E=0\) 时中性 singlet 对这些局域异常无贡献。

\protect\hypertarget{p-qg3-04}{}{}

\textbf{P-QG3-04。} 按此完整表示表并施加 \(L^3\) 异常方程后，直接求和得
\[
\begin{gathered}
A_{L^2X}=\frac{N_f(N_c-3)}{12},\\
A_{X^3}=-\frac{N_f(N_c-3)(1+\beta_8^2)}4,
\\ A_{\mathrm{grav}^2X}=A_{C^2X}=A_{C^3}=0.
\end{gathered}
\] 验证方式：三重态 \(X^3\) 分别为 \(3X_L^3,3N_cX_{Q\pm}^3\)；加上
singlet 的负物理电荷立方，再将 \(k=N_f(N_c-1)/(2N_c)\)
代入并展开，得到上述因式。混合项直接使用 \(T(3)=1/2\)。\(\square\)

于是此类 \(N_f>0\) 的局域抵消迫使 \(N_c=3\)，再得
\(N_f=3k\)；仍需上界或其他复制破缺原则才能选
\(k=1\)。若将普通夸克电荷改为第29章的
\(N_c\)-依赖解，混合异常的强制三色结论消失。因此这不是从零推出普通电荷和三色。

\hypertarget{ux4e0eux5df2ux9a8cux8bc1ux5185ux6838ux7684ux51c6ux786eux5bf9ux5e94ux53caux6c38ux4e45ux53cdux6a21ux578b}{%
\section{31.6
与已验证内核的准确对应及永久反模型}\label{ux4e0eux5df2ux9a8cux8bc1ux5185ux6838ux7684ux51c6ux786eux5bf9ux5e94ux53caux6c38ux4e45ux53cdux6a21ux578b}}

第 31.3 节条件选择器的整数内容已在有限 UEOT-QM 内核中实现为
\texttt{Selector.\allowbreak selects\_\allowbreak three\_\allowbreak without\_\allowbreak atomicity}。其输入明确包括自然数
\(N_c,N_f,k\)、\(N_c\ge3\)、\(N_f>0\)、\(k\le N_f\)、给定立方方程及严格弱预算。它没有从一般
Core 产生这些物理前提。

\texttt{Spectrum.\allowbreak spectrum\_\allowbreak selects\_\allowbreak three}
进一步从有限谱行求和接到该选择器。接口要求轻子朝向同质、非零异常单位、每行普通轻子与夸克双重态在同一弱窗口活跃、非负额外
screening、完整弱系数为正；未列普通行之外的异常补偿总和也须单独等于零。这些都是有实质内容的合同。它没有构造
\(SU(3)_L\)
的量子场、证明额外补偿机制有效，或在改变阈值后继续使用原窗口。

例如放宽同质轻子模板，令 \(l\) 代轻子取 \(R\)、其余取
\(\bar R\)，则立方式变为 \[
(2l-N_f)+N_c(2k-N_f)=0.
\] \((N_c,N_f,l,k)=(3,2,1,1)\) 使其为零，并有
\(N_f(N_c+1)=8<22\)。这是对``谱朝向无关紧要''的具体反例。内核验证了相应的行计数与弱系数；在另行声明的一般化
331 电荷表、\(\beta_8=0\) 等条件下，Python
精确算术还核对了六项局域异常。后者不因此成为 Lean
的完整六异常定理，也不构成已可行的全部粒子模型。

同样，在只计普通费米子的三代点上，加入总数二十个复 Higgs
双重态可把弱系数压到零；所以普通必要预算通过不保证加入任意标量后的完整系数仍为正。若允许额外立方补偿抵消普通行的非零异常，单代等其他点也可能进入更宽模型类。相反例与正例应永久一起保存，防止把物理前提在摘要中删掉。

第 31.4--31.5 节的完整六局域异常求和、第 32 章 framing/指标路线，以及第
36
章正确高能规范群上的分类，是不同的数学对象，仍按各自纸面证明和精确回归登记。一个形式化定理数量不能代替这些具体覆盖关系。

\hypertarget{ch-32}{%
\chapter{32. 其他演进结果、拓扑路线与开放边界}\label{ch-32}}

\hypertarget{framing-completion-ux7684ux7cbeux786eux6761ux4ef6ux5206ux7c7b}{%
\section{32.1 Framing completion
的精确条件分类}\label{framing-completion-ux7684ux7cbeux786eux6761ux4ef6ux5206ux7c7b}}

下面是一条独立于 HGP 的候选路线，不能将两套前提不加说明地相互替代。

\textbf{H-QFR-1（FAD）。} 声明某种具体降维/缺陷实现，其 descendant
必须在无额外 inflow 的自治条件下忘掉指定 framing。再假定每代 completion
的贡献可加，并以 \(\chi=2c_-^{\mathcal C}\in\mathbb R/48\mathbb Z\)
记录；普通每代 \(4N_c+3\) 个 Weyl 的指定 descent
贡献同数单位。则所采用的条件是 \[
N_f(4N_c+3+\chi)=0\quad\text{in }\mathbb R/48\mathbb Z.
\] 这是完整物理假设包中的一条相容式，不是一般四维异常消除公式。

复色 \(N_c\ge3\)、KM 型 CP 的 \(N_f\ge3\)、弱 AF 必要界给七个候选： \[
(3,3),(3,4),(3,5),(4,3),(4,4),(5,3),(6,3).
\]

\protect\hypertarget{p-qfr-01}{}{}

\textbf{P-QFR-01。} 在此七点集内，framing 式的解恰好只有
\((3,3)\)，当且仅当 \(\chi=1\bmod48\)。

\textbf{证明。} \((3,3)\) 通过要求
\(3(15+\chi)=0\bmod48\)，故模48仅可能为1、17、33。取1时条件为
\(12\mid N_f(N_c+1)\)，七点中仅 \((3,3)\) 的积为12；17还允许
\((4,4)\)；33还允许 \((3,4),(3,5)\)。\(\square\)

若 \(\chi=0\)，因为 \(4N_c+3\) 为奇数，要求
\(16\mid N_f\)，与上界矛盾。若只允许每代0或1个 singlet Weyl
completion，可在该狭窄类中选1；不能据此排除耦合 TQFT 或非逐代可加
completion。

相同四维异常类不自动给相同 descendant 的 \(\chi\)。需要具体、自然的
descent 映射及 framing 响应证明；目前该普遍桥并未建立。

\hypertarget{ux81eaux6cbbux6027ux5bf9ux8c61ux5212ux5206ux4e0eux5916ux7f6e-completion}{%
\section{32.2 自治性、对象划分与外置
completion}\label{ux81eaux6cbbux6027ux5bf9ux8c61ux5212ux5206ux4e0eux5916ux7f6e-completion}}

\protect\hypertarget{p-qfr-02}{}{}

\textbf{P-QFR-02。} 若联合过程精确因子化
\(\mathbf T_{FC}=\mathbf T_F\otimes\mathbf T_C\)，合法 cut
仅删除交叉耦合，则 cut 不改变过程，任何操作距离下的该 cut 距离为零。

\textbf{证明。}
原过程没有交叉项可删除，故两个被比较的响应相同。\(\square\)

这只阻止把完全解耦的外置系统凭命名认作因果整合父对象。耦合拓扑扇区、共享规范变量或线/面协议可能有真实交叉结构，不受此反例排除。父系统自治也不推出每个有边界
descendant 自治；边界异常可由体 inflow 抵消。

味空间维数 \(N_f\)
不自动赋予规范无关的逐代子对象基。质量、交互作用和具体可实施协议可选出某些结构，但必须给出证书，不能把任意坐标基当作自然界的分解。

\hypertarget{ux4eceux539fux8bbaux6587ux4fddux7559ux7684ux62d3ux6251ux8defux7ebf}{%
\section{32.3
从原论文保留的拓扑路线}\label{ux4eceux539fux8bbaux6587ux4fddux7559ux7684ux62d3ux6251ux8defux7ebf}}

\textbf{K-QIDX-01。} 对闭 spin 流形与适当扭曲丛， \[
\operatorname{Index}D_E=\int_M\widehat A(TM)\operatorname{ch}(E)
\] 取最高次数分量。具体在 \(S^2\) 上，线丛次数 \(n=(2\pi)^{-1}\int F\)
给 \(\operatorname{Index}D_E=n\)，而 \(b_1(S^2)=0\)。
指数是正负手征零模之差，绝非任意 Betti 数。

因此合法研究链是``独立选择载体 \(M,E,D\) → 计算指数 → 连接净手征零模 →
检查额外成对零模及低能谱''。选一个次数3的丛再计算出3是条件实现，不是三代选择。原始
\(b_1=3\Rightarrow\operatorname{Index}D=3\) 链不进入本书证明体系。

\hypertarget{ux660eux786eux9000ux51faux5df2ux8bc1ux4e3bux7ebfux7684ux4e3bux5f20}{%
\section{32.4
明确退出已证主线的主张}\label{ux660eux786eux9000ux51faux5df2ux8bc1ux4e3bux7ebfux7684ux4e3bux5f20}}

\begin{longtable}[]{@{}
  >{\raggedright\arraybackslash}p{(\columnwidth - 4\tabcolsep) * \real{0.2500}}
  >{\raggedright\arraybackslash}p{(\columnwidth - 4\tabcolsep) * \real{0.3400}}
  >{\raggedright\arraybackslash}p{(\columnwidth - 4\tabcolsep) * \real{0.4100}}@{}}
\toprule\noalign{}
\begin{minipage}[b]{\linewidth}\raggedright
主张
\end{minipage} & \begin{minipage}[b]{\linewidth}\raggedright
本书状态
\end{minipage} & \begin{minipage}[b]{\linewidth}\raggedright
恢复所需材料
\end{minipage} \\
\midrule\noalign{}
\endhead
\bottomrule\noalign{}
\endlastfoot
闭轨相位积分直接推出 CCR & N & 正则操作、表示、连续性与非平凡中心角色 \\
有限信息自动给标准不确定常数 & N & 物理算符、态域、\(\hbar\) 标定 \\
公共记录等于客观唯一坍缩 & O & 可区别于开放系统的动力学及实验 \\
任意两角测试证明全局 Born 规则 & N & 全域量词或已固定有限参数模型 \\
\(b_1\) 直接决定三代 & N & 指数载体和独立选择机制 \\
common-\(q\) 或 CP 容量单独选3 & N & 排除其他一致候选的独立条件 \\
自治父系统强迫自治 descendant & N & 无 inflow 的额外条件 \\
规范异常自动选出完整 SM & N & 规范群、表示、全局结构及参数选择 \\
核心对象数学已证明 HGP/FAD & O & 明确实现定理与反模型排除 \\
\end{longtable}

所谓 QBX/量子基础突破候选只能按``数学模型 → 模拟证据 → 硬件协议 →
真实盲测 →
独立复现''升级。本书没有新的硬件数据，也没有宣布客观坍缩、超量子性能或标准模型以外的新粒子已被发现。

\hypertarget{ch-33}{%
\chapter{33. 全部架构的接口契约}\label{ch-33}}

\hypertarget{ux5206ux5c42ux53caux5176ux6743ux8d23}{%
\section{33.1 分层及其权责}\label{ux5206ux5c42ux53caux5176ux6743ux8d23}}

\begin{longtable}[]{@{}
  >{\raggedright\arraybackslash}p{(\columnwidth - 6\tabcolsep) * \real{0.2100}}
  >{\raggedright\arraybackslash}p{(\columnwidth - 6\tabcolsep) * \real{0.2000}}
  >{\raggedright\arraybackslash}p{(\columnwidth - 6\tabcolsep) * \real{0.2800}}
  >{\raggedright\arraybackslash}p{(\columnwidth - 6\tabcolsep) * \real{0.3100}}@{}}
\toprule\noalign{}
\begin{minipage}[b]{\linewidth}\raggedright
层
\end{minipage} & \begin{minipage}[b]{\linewidth}\raggedright
主要对象
\end{minipage} & \begin{minipage}[b]{\linewidth}\raggedright
输出
\end{minipage} & \begin{minipage}[b]{\linewidth}\raggedright
禁止隐含的升级
\end{minipage} \\
\midrule\noalign{}
\endhead
\bottomrule\noalign{}
\endlastfoot
Core & 协议、历史、响应等价、资源与控制 & 预测因子、载体、持续性证书 &
默认复 Hilbert、Born 或粒子表 \\
Q-Selector & 有限因果 OPT 与附加原则 & 量子实现的条件选择 &
将领域原则称作 Core 定理 \\
Q-Operations & 态、通道、instrument、效应 & 序贯概率、后验、误差界 &
单时通道自动多时无记忆 \\
Q-Objects & 编码/解码、访问、相干父系统 & 有效因子、公共记录、SRN 核 &
协议不可见等同本体不存在 \\
Q-Dynamics & 自伴算符、CPTP 演化、过程张量 & 时间响应、控制与稳态 &
隐藏定义域、初始相关或截断 \\
Q-Fields & 局域代数、场分布、调节、作用量 & 相关函数、振幅、EFT 系数 &
形式级数等于非微扰构造 \\
Q-Particles & 规范群、表示、异常、Higgs/味 & 候选一致性与条件选择 &
一致性等于自然界唯一实现 \\
Q-Evidence & 数据、校准、盲测与统计方案 & 区间、模型比较、失败原因 &
合成验证等于真实实验证据 \\
\end{longtable}

每层必须携带模型身份、单位、适用域、前提和证书。上下游交换的是类型明确的结果，不是相同缩写。

\hypertarget{ux6709ux9650ux91cfux5b50ux64cdux4f5cux6700ux5c0f-api}{%
\section{33.2 有限量子操作最小
API}\label{ux6709ux9650ux91cfux5b50ux64cdux4f5cux6700ux5c0f-api}}

以下保留覆盖全部量子架构的接口规范；其中已实现的有限精确内核在33.8节逐项对应实际名称。规格中的一般
Choi、连续域与过程张量接口不因有限内核完成而自动成为软件功能。矩阵以复数二维数组表达；向量化和复合因子顺序全局固定。

\begin{Verbatim}[breaklines,breakanywhere,fontsize=\scriptsize]
SystemSpec:
  labels: ordered subsystem labels
  dims: positive integers
  basis_id: immutable string
  units: explicit energy/time/action convention

State:
  system: SystemSpec
  rho: Hermitian positive matrix
  trace: 1, or separately tagged subnormalized
  validity: tolerance + minimum eigenvalue + trace residual

Channel:
  input, output: SystemSpec
  choi: input-first, UNNORMALIZED Choi
  validate: J >= 0; Tr_output J = I_input
  apply(rho): Tr_input[(rho^T tensor I_output) J]
  compose(after, before): match system identities, not just dimensions

Instrument:
  maps: outcome -> CP trace-nonincreasing map
  validate: sum of maps is TP
  observe(rho): outcome probabilities + subnormalized branches
  condition(outcome): normalized state only if probability > 0

Protocol:
  initialization
  causal history -> permitted instrument/action
  terminal effects
  resource/time budget
  target and stopping rule

Certificate:
  model_id, protocol_id, target_id
  assumptions, theorem_ids, estimator_spec
  domain, tolerances, confidence, failure_modes
  evidence_kind: proof | exact_arithmetic | synthetic | experiment
\end{Verbatim}

Channel.compose 不允许用 Kronecker 乘积代替串联；tensor
是并行复合。向量化约定为
\(|A\rangle\rangle=\sum_i|i\rangle_{\rm in}\otimes A|i\rangle_{\rm out}\)，
因此 \(J=\sum_\mu|A_\mu\rangle\rangle\langle\langle A_\mu|\)。无归一
Choi 的迹等于输入维数；态归一 Choi 是另一种明确标记的对象。

输入失败应返回
INVALID\_\allowbreak STATE、NOT\_\allowbreak CP、NOT\_\allowbreak TP、BASIS\_\allowbreak MISMATCH、ZERO\_\allowbreak PROBABILITY\_\allowbreak BRANCH
等可解释错误，不得默默把非正矩阵剪裁为有效模型。若请求投影修复，必须同时返回原始残差和所解约束问题。

\hypertarget{ux6709ux8bb0ux5fc6ux8fc7ux7a0bux4e0e-comb}{%
\section{33.3 有记忆过程与
comb}\label{ux6709ux8bb0ux5fc6ux8fc7ux7a0bux4e0e-comb}}

按时间顺序令第 \(j\) 次过程步输入腿为 \(2j-2\)，输出腿为
\(2j-1\)。确定性 \(n\)-comb 的 Choi 算符 \(R^{(n)}\) 满足 \[
R^{(n)}\succeq0,\qquad
\operatorname{Tr}_{2n-1}R^{(n)}
=I_{2n-2}\otimes R^{(n-1)},\quad R^{(0)}=1.
\] 每一递归关系按指定腿置换后比较，不依赖数组默认顺序。外插 instrument
使用相容的 Choi link contraction；共享腿转置约定须与33.2一致。可选结果
comb 的和才满足确定性归一。

过程张量是将初始化、系统自由传播及环境记忆吸收后，对一系列插槽操作给概率/末态的
comb 型对象。API
包含插槽时间、输入/输出腿、可控操作集和初始相关范围。两个过程只在同一
tester 集上比较。CP 可分性、CP-divisibility、Markov
因子化和未来充分性必须分项报告。

\hypertarget{ux5bf9ux8c61ux63a5ux53e3ux4e0eux7b97ux6cd5}{%
\section{33.4
对象接口与算法}\label{ux5bf9ux8c61ux63a5ux53e3ux4e0eux7b97ux6cd5}}

\begin{Verbatim}[breaklines,breakanywhere,fontsize=\scriptsize]
QuantumObjectSpec:
  carrier: physical system or accessible algebra
  protocol_family: allowed future interventions/readouts
  target: predicted observable/history
  encoder, decoder: optional CP maps
  lifecycle: initialization, survival, recovery, termination
  cuts: legal access restrictions or physical interventions
  records: pointer projectors + fragment specification
  tolerance: metric, horizon, confidence

certify_predictive_factor(spec):
  verify admissibility and common protocol identity
  compare future response laws or derive exact intertwining
  distinguish exact equality from finite-test agreement
  return factor certificate, unresolved directions, witness protocols

certify_SRN(channel, projectors):
  check each branch's entire operator algebra, not only diagonal populations
  recover block coefficients K_ab
  test Hermiticity, diagonal=1, PSD and model residual
  return K, generator status if time data available, deviations

discover_access_carriers(spec):
  enumerate permitted access subsets or use a declared search algorithm
  test target sufficiency for each
  output minimal certified subsets relative to this finite universe
  do not claim global minimality beyond searched cuts
\end{Verbatim}

``访问受限''只是减少可用效应；物理切断/擦除是改变过程。两个接口分开，以免把读不出相干错误解释为相干已消失。

\hypertarget{ux8fdeux7eedux53d8ux91cfux7b97ux7b26ux4e0eux6570ux503cux540eux7aef}{%
\section{33.5
连续变量、算符与数值后端}\label{ux8fdeux7eedux53d8ux91cfux7b97ux7b26ux4e0eux6570ux503cux540eux7aef}}

OperatorSpec 至少记录 Hilbert
空间、稠密域、边界条件、对称/自伴状态、谱下界、共同不变核及所采用的扩张。未指定边界的微分表达式不等于一个自伴
Hamiltonian。

CVBackend 记录相空间顺序、\(\hbar\)、辛形式、Weyl 乘子、Gaussian
均值/协方差规范和截断。有限 Fock 截断仅近似低能行为，不可谎报精确
CCR。误差分成尾概率、截断动力学误差、积分误差和统计误差；后两者不能由静态能量尾界自动替代。

\hypertarget{ux573aux8bbaux4e0eux7c92ux5b50ux540eux7aef}{%
\section{33.6
场论与粒子后端}\label{ux573aux8bbaux4e0eux7c92ux5b50ux540eux7aef}}

\begin{Verbatim}[breaklines,breakanywhere,fontsize=\scriptsize]
QFTModel:
  spacetime, metric, units
  local_fields + representations + statistics
  global_gauge_group and permitted bundles
  regulated_action and gauge_fixing
  scheme, scale, renormalization_conditions
  perturbative_order or nonperturbative_definition
  operator_basis, boundary_conditions, infrared_resolution

AmplitudeRequest:
  external states, mass/charge convention
  spin preparation or averaging
  kinematic domain and phase-space cuts
  inclusive/exclusive definition
  requested order, EFT power and uncertainty

AnomalyCertificate:
  all-left-handed representation table
  multiplicities, Dynkin and cubic normalizations
  local anomaly polynomial coefficients
  global/discrete/boundary/inflow checks: done or unresolved
  assumptions under which quotient is defined
\end{Verbatim}

普通场表核对六个局域系数为零时，应返回
LOCAL\_\allowbreak ANOMALIES\_\allowbreak CANCELLED，不能直接返回
FULL\_\allowbreak QUANTUM\_\allowbreak THEORY\_\allowbreak CONSISTENT。微扰可计算性、全局规范问题、稳定真空、幺正性、实际谱和实验排除均是另一些检查。

\hypertarget{ux6570ux636eux7248ux672cux4e0eux8bc1ux636eux63a5ux53e3}{%
\section{33.7
数据、版本与证据接口}\label{ux6570ux636eux7248ux672cux4e0eux8bc1ux636eux63a5ux53e3}}

每条实验数据必须指向原始计数或已说明的充分统计量、制备/测量校准、采集时间、随机化、排除规则和数据冻结点。响应矩阵保留有噪估计，不只保存
PSD 投影后的``物理''矩阵。

版本升级记录：旧命题、变更原因、新前提/删除前提、受影响下游结果、回归结果。更换协议、目标、基、单位或作用量算符基是语义变更；不能只有文件版本号而没有模型身份更新。

\hypertarget{ux5df2ux5b9eux73b0ux7684ux6709ux9650ux5185ux6838ux4e0eux63a5ux53e3ux6620ux5c04}{%
\section{33.8
已实现的有限内核与接口映射}\label{ux5df2ux5b9eux73b0ux7684ux6709ux9650ux5185ux6838ux4e0eux63a5ux53e3ux6620ux5c04}}

\begin{longtable}[]{@{}
  >{\raggedright\arraybackslash}p{(\columnwidth - 4\tabcolsep) * \real{0.2500}}
  >{\raggedright\arraybackslash}p{(\columnwidth - 4\tabcolsep) * \real{0.3400}}
  >{\raggedright\arraybackslash}p{(\columnwidth - 4\tabcolsep) * \real{0.4100}}@{}}
\toprule\noalign{}
\begin{minipage}[b]{\linewidth}\raggedright
研究接口
\end{minipage} & \begin{minipage}[b]{\linewidth}\raggedright
当前实际入口
\end{minipage} & \begin{minipage}[b]{\linewidth}\raggedright
已实现范围
\end{minipage} \\
\midrule\noalign{}
\endhead
\bottomrule\noalign{}
\endlastfoot
精确态与操作 &
\texttt{QComplex}、\allowbreak \texttt{Matrix}、\allowbreak \texttt{Kraus\allowbreak Term}、\allowbreak \texttt{Instrument}
& Gaussian-rational 矩阵、非负有理权重；Python PSD 检查限维数1--4 \\
允许库与有限树 & \texttt{Instrument\allowbreak Library}、\allowbreak \texttt{Protocol\allowbreak Node} &
上下文指纹、完整结果、精确有限叶律 \\
响应表示 & \texttt{Predictive\allowbreak Model}、\allowbreak \texttt{Encoded\allowbreak State} &
词效应闭包搜索、精确坐标、分支与后验更新 \\
可移植表示证书 &
\texttt{export\_\allowbreak representation\_\allowbreak certificate}、\allowbreak \texttt{verify\_\allowbreak representation\_\allowbreak certificate}
& 数值基/生成词/全部局部恒等式与上下文的独立核验 \\
Lean 数值入口 &
\texttt{export\_\allowbreak lean\_\allowbreak representation}、\allowbreak \texttt{Exact\allowbreak Certificate.\allowbreak check} &
生成具体有限模型，再由 Lean 内核检查真值证明及后果 \\
实际效应与误差数学 &
\texttt{Effect}、\allowbreak \texttt{effect\allowbreak Distance}、\allowbreak \texttt{channel\allowbreak Defect} &
全效应上确界、完整通道收缩、正事件后验与轨迹定理；上确界不是运行时优化器 \\
粒子谱入口 &
\texttt{Spectrum}、\allowbreak \texttt{inspect\_\allowbreak spectrum}、\allowbreak \texttt{local\_\allowbreak 331\_\allowbreak anomalies}
& 明确谱行和窗口的精确算术；六种局域异常与 Lean
条件选择的范围分别标记 \\
\end{longtable}

Python 精确算术检查、Lean
的通用定理、具体导出模型编译是三个独立证据层。表示 checker
允许维数1--4、至多 \(d^2\) 个词效应和长度至多 \(d^2-1\)
的证据词；这些限制只约束当前证据格式，不把已获证书的全部未来响应限定为该长度。Lean
exporter 另有4096位输入有理数、65536个补齐标量存储及8
MiB源文件限制，不能将这些产品边界当作物理常数。

最小复查路径为：构造已固定上下文的仪器库；生成表示证书；独立核验局部矩阵条件；导出具体
Lean 数据；实际编译其
\texttt{checked}；读取公理与源哈希回执；最后引用第49--50章相应的全词、后验或误差结论。只完成导出时状态仍为
\texttt{LEAN\_\allowbreak SOURCE\_\allowbreak EXPORTED\_\allowbreak NOT\_\allowbreak COMPILED}，不能提前标为已形式验证。

\hypertarget{ch-34}{%
\chapter{34. 完整应用路径}\label{ch-34}}

\hypertarget{ux5355ux6bd4ux7279ux4e0eux901aux9053ux6821ux51c6}{%
\section{34.1
单比特与通道校准}\label{ux5355ux6bd4ux7279ux4e0eux901aux9053ux6821ux51c6}}

输入：未知制备态、可校准的三组 Pauli 测量以及计数。由
\(\hat r_j=(n_{j,+}-n_{j,-})/n_j\) 线性重建
\(\hat\rho=(I+\sum_j\hat r_j\sigma_j)/2\)。保留原估计与置信区间；若
\(|\hat r|>1\)，可能是采样噪声，也可能是漂移/模型失配，不能单凭剪裁宣布实验符合模型。

通道层析用输入线性基和输出信息完备测量恢复 \(J\)，再检验 CP、TP
及留出输入。SPAM
未独立校准时，门与制备测量的误差不可唯一分解；需要联合估计或明确的 gauge
等价报告。指标可含平均保真度、最坏协议误差、泄漏和漂移，不能仅一个高保真数字代替全部证书。

\hypertarget{ux91cfux5b50ux8bb0ux5f55qdrr-ux4e0e-actualization-ux68c0ux9a8c}{%
\section{34.2 量子记录、QDRR 与 actualization
检验}\label{ux91cfux5b50ux8bb0ux5f55qdrr-ux4e0e-actualization-ux68c0ux9a8c}}

闭环步骤是：相干制备 → 写入指针/环境记录 → 指定等待与公共读取 →
在可控边界内逆写 →
对相干做相位层析。记录访问、不可控环境泄漏和逆操作误差分别记录。

在 SRN 假设内，\(m\) 个分支的 \(m(m-1)/2\) 对，每对实虚两相位，重建
\(K\)。第8章界以每相位 \(n\) 次独立有界结果计算。例
\(m=3,\epsilon=0.12,\alpha=0.05\)
时每设置762次，共6设置4572次（未计校准与对照）。同时必须检验 SRN
模型，而非先假设其精确成立再宣称发现或排除新物理。

关键对照包括不同存储时长、环境访问范围、回波精度、记录冗余、随机化顺序与空协议。若观测额外不可逆损失，先比较已冻结的开放系统模型及设备噪声；新
actualization
模型必须给独立参数、尺度依赖和能被反驳的留出预测。没有真实计数时只称协议设计。

\hypertarget{ux4e8cux76f8ux4f4d-gram-choi-ux68c0ux9a8cux7684ux5e94ux7528ux5206ux5de5}{%
\section{34.3 二相位 Gram / Choi
检验的应用分工}\label{ux4e8cux76f8ux4f4d-gram-choi-ux68c0ux9a8cux7684ux5e94ux7528ux5206ux5de5}}

项目中不同 Paper-A/B/C 或其他内部标签不是证明类别。统一落到三类对象：

\begin{longtable}[]{@{}
  >{\raggedright\arraybackslash}p{(\columnwidth - 4\tabcolsep) * \real{0.2500}}
  >{\raggedright\arraybackslash}p{(\columnwidth - 4\tabcolsep) * \real{0.3400}}
  >{\raggedright\arraybackslash}p{(\columnwidth - 4\tabcolsep) * \real{0.4100}}@{}}
\toprule\noalign{}
\begin{minipage}[b]{\linewidth}\raggedright
应用目标
\end{minipage} & \begin{minipage}[b]{\linewidth}\raggedright
数据与重建
\end{minipage} & \begin{minipage}[b]{\linewidth}\raggedright
合法结论
\end{minipage} \\
\midrule\noalign{}
\endhead
\bottomrule\noalign{}
\endlastfoot
干涉二次结构 & 基态及成对0、\(\pi/2\) 相位响应 & 已声明有限二次模型的
Gram 估计 \\
操作过程 & 输入基、输出完备读出、Choi 约束 & 指定维度/校准下的 CP/TP
相容性 \\
公共记录后的相干 & 写---存---逆写及 \(K_{ab}\) & SRN
参数和对照模型差异 \\
\end{longtable}

先冻结模型和检验统计量，保留 raw
residual，再做必要的约束估计，最后对未用数据检验。不能``投影到量子集合 →
投影结果是量子 → 证明所有概率必须量子''。

\hypertarget{ux91cfux5b50ux5b58ux50a8ux7ea0ux9519ux4e0eux63a7ux5236}{%
\section{34.4
量子存储、纠错与控制}\label{ux91cfux5b50ux5b58ux50a8ux7ea0ux9519ux4e0eux63a7ux5236}}

给逻辑编码 \(V:\mathcal H_L\to\mathcal H_P\)、噪声 \(\mathcal E\)、恢复
\(\mathcal R\)，逻辑通道
\(\mathcal N_L=\mathcal D\mathcal R\mathcal E\mathcal V\)。
Knill--Laflamme 条件给理想可纠错类；近似纠错必须直接评估逻辑误差及泄漏。

例三比特重复码可纠正单个 \(X\)
错误，不纠正任意单比特量子误差；忽略这一点会把经典比特保护误报为完整量子纠错。容错阈值需要具体码、噪声相关、解码器和门调度，本文不给无模型的通用阈值。

闭环控制按第20章计算价值或脉冲梯度。实际控制器只能依赖已获得记录，计算延迟和带宽计入资源。GOA
报告为逻辑存活概率、稳态集合或可恢复域，而不预设一个普适终极态。

\hypertarget{ux591aux4f53ux6decux706bux4e0eux70edux73afux5883}{%
\section{34.5
多体淬火与热环境}\label{ux591aux4f53ux6decux706bux4e0eux70edux73afux5883}}

以有限 TFIM
为例，指定边界、\(J,h\)、初态和淬火时间；计算磁化、相关函数和熵。小系统精确对角化可核对张量网络，小时间可用级数核对积分器。有限体积、bond
dimension、Trotter 步长与观测误差分别收敛。

环境作用存在时，先检验使用 GKSL 的尺度分离；若不同制备历史经同一 causal
break
后仍给不同未来，系统态不是充分预测态，应增加记忆变量或采用过程张量。不能用更精细时间步挽救错误的无记忆模型。

\hypertarget{ux534aux5bfcux4f53ux8d85ux5bfcux91cfux5b50ux5668ux4ef6ux7684ux6a21ux578bux63a5ux53e3}{%
\section{34.6
半导体/超导量子器件的模型接口}\label{ux534aux5bfcux4f53ux8d85ux5bfcux91cfux5b50ux5668ux4ef6ux7684ux6a21ux578bux63a5ux53e3}}

量子点可从有限轨道 Anderson 型 Hamiltonian
出发，声明库仑能、隧穿、温度、偏压与读出。近邻超导纳米线的常见有效 BdG
模型可写为 \[
H_{\rm BdG}(k)=
\left(\frac{k^2}{2m_*}-\mu\right)\tau_z
+\alpha_R k\,\sigma_y\tau_z+V_Z\sigma_x+\Delta\tau_x,
\] 其中 Nambu/自旋基及相位约定固定。此为器件领域的有效模型输入，不是
Core 定理；无序、多带、轨道效应、自能与接触会改变简化预测。理论先例见
\href{https://arxiv.org/abs/1002.4033}{Lutchyn--Sau--Das Sarma} 与
\href{https://arxiv.org/abs/1003.1145}{Oreg--Refael--von Oppen}。

闭合均匀单带模型的能隙闭合条件 \(V_Z^2=\mu^2+\Delta^2\)
只在其假设内成立。局域零偏压峰或 Andreev 束缚态信号不单独证明拓扑
Majorana
模式；需要排除平凡机制并设计更有区分力的协议。这里不声称项目已有这种器件的实验发现。

\hypertarget{ux6563ux5c04eft-ux4e0eux7c92ux5b50ux5019ux9009ux6bd4ux8f83}{%
\section{34.7 散射、EFT
与粒子候选比较}\label{ux6563ux5c04eft-ux4e0eux7c92ux5b50ux5019ux9009ux6bd4ux8f83}}

对 QED 基准，先计算 Born
截面，再逐级加入质量、辐射、接受度和相关不确定度。对 EFT，先明确能区
\(E/\Lambda_*<1\)
和算符截断，匹配系数后在留出过程比较；不得在重极点附近用低阶局域展开硬拟合。

对 HGP/331
候选，异常自由只开启后续任务：标量破缺、谱与混合、额外规范玻色子和
exotic 粒子、味变中性流、精密观测、真空稳定、全局异常及运行窗口。对
framing 候选，则先完成 descent 与 completion
分类。应分别冻结预测、拟合参数和排除条件，不把三代算术通过视为整个模型胜出。

\hypertarget{ch-35}{%
\chapter{35. 量子主干的装配、误差与扩展入口}\label{ch-35}}

\hypertarget{ux5355ux4e00ux88c5ux914dux5f0f}{%
\section{35.1 单一装配式}\label{ux5355ux4e00ux88c5ux914dux5f0f}}

在有限无记忆实现中， \[
p(\mathbf a\mid\pi)=
\operatorname{Tr}\!\left[
\mathcal M_{a_T|\mathbf a_{<T}}\circ\cdots\circ
\mathcal M_{a_1}(\rho_0)\right].
\] 有记忆时以合法 comb 与相应 tester
的配对替换；连续变量必须声明能量/算符域；场论把制备、探测和相空间分辨率接到相关函数/振幅。Core
再以这套概率响应定义对象、未来充分性与控制任务。

因此统一的是``实现---干预---响应---证书''的组合规则，并使各物理层的独立内容可追踪，而不是把所有物理写成同一个不分类型的公式。

\hypertarget{ux53efux7ec4ux5408ux7684ux8befux5deeux8bc1ux4e66}{%
\section{35.2
可组合的误差证书}\label{ux53efux7ec4ux5408ux7684ux8befux5deeux8bc1ux4e66}}

先在数学纸面层给出包含旁系统的合同。一次仪器不能只保留选择分支
\(\mathcal M_o\)，而应写成保留全部标签的 CPTP 通道 \[
\widehat{\mathcal M}(X)=\sum_o|o\rangle\langle o|\otimes\mathcal M_o(X).
\]
自适应控制把此前经典历史加入显式寄存器，按历史块选择下一仪器；停止路径继续保存一个吸收的停止标记。完整性保证每一整步在共同的经典---量子载体上为
CPTP。比较的两个模型必须使用一致的标签、寄存器和允许协议；丢弃结果后才相等的通道可能仍有不同的多时记录律，见第5.3节反例。

对同一协议与目标，若初态（包括被比较的初始记忆）迹距离误差至多
\(\epsilon_0\)，每个完整步骤在任意量子旁系统上有 \[
\tfrac12\|\widehat{\mathcal E}_j-\widehat{\widetilde{\mathcal E}}_j\|_\diamond
\le\epsilon_j,
\] 而末端读出造成的总变差缺陷一致有界为 \(\epsilon_M\)，则 \[
\operatorname{TV}(p,\widetilde p)
\le\min\left\{1,\epsilon_0+\sum_j\epsilon_j+\epsilon_M\right\}.
\] 证明为在实际复合中逐步替换通道，使用 diamond
缺陷对共同旁系统的控制、后续完整通道收缩与三角不等式，再应用测量收缩。这个论证不适用于把未归一选择分支直接当作
TP，或在没有记录寄存器的模型中省略自适应记忆。

第50章给出与此纸面合同相接的\textbf{已由 Lean 检查的有限切片}：实际完整
Kraus 通道对全效应距离 \(d_E\) 收缩，定义的实际输出缺陷 \(\Delta_E\)
给出一步递推，继而接入既有 Core 的误差运输定理。该切片没有形式化 diamond
范数、任意旁系统稳定性或一般历史寄存器装配；这些不随``通道误差''四字被默认获得。

若另有领域模型逼近、截断和估计误差，只有在同一目标概率度量中各自有界后才能相加。能量误差、保真度误差、对数似然和
Wilson
系数误差应先经有证明的转换。后选择则单独使用正事件下界：本轮从实际 Born
分支得到
\(d_E(\rho_o,\sigma_o)\le2d_E(\rho,\sigma)/p_o(\rho)\)；当概率很小时，此界可能没有信息，零事件也没有归一化后验。

多个置信事件相交可用 union
bound，总失败率至多各项之和；独立性不能未经检验就假定。渐近级数的未知高阶估计属于模型不确定性，不是自动有覆盖保证的频率学置信区间。

\hypertarget{ux672cux6b21ux72ecux7acbux68c0ux9a8cux7684ux8303ux56f4}{%
\section{35.3
本次独立检验的范围}\label{ux672cux6b21ux72ecux7acbux68c0ux9a8cux7684ux8303ux56f4}}

附录D包含完整代码和实际运行结果。本版重新执行继承的三套完整程序，共757项有限检查；程序采用
NumPy/SciPy 与有理数实现，覆盖：

\begin{itemize}
\tightlist
\item
  CPTP/Choi、非 CP 反例、rebit 局部盲性、相干记录与逆写；
\item
  SRN 块结构、复 Gram 核、生成元正性及反例、相位恢复；
\item
  截断 CCR、Gaussian 正性、Rabi、自旋、CHSH、控制梯度；
\item
  开放系统、KMS、自由能、因果切断记忆；
\item
  Clifford、QED Ward 树恒等式、散射迹、圈参数积分、EFT 余项；
\item
  Higgs/seesaw 基准、SM 与331有理异常、去原始包的三代搜索、framing
  分类。
\end{itemize}

这些是有限数值与精确算术回归，不是所有定理的形式化验证。尤其没有独立机器证明重建定理、Stone--von
Neumann、GKSL、LSZ、自旋---统计或四维非微扰构造；没有真实实验数据。有限整数搜索用于复核；全整数三代结论依据第31章解析证明，并由
\texttt{Selector} 的完整前提整数定理及 \texttt{Spectrum}
的条件谱桥分别得到形式化支持。第49--50章报告的 Lean
成果只对应其精确定义与定理清单，不覆盖本段所列全部物理计算。

\hypertarget{ux5b8cux6210ux6539ux8fdbux4e0eux5f00ux653e}{%
\section{35.4
完成、改进与开放}\label{ux5b8cux6210ux6539ux8fdbux4e0eux5f00ux653e}}

本书维护一份包含全部声明模块的正式连续主文，并把当前有限量子内核与误差证明整合到第49--50章。实质改进是：保留数学
Core 的领域中立性；给测量的相干父层、公共层与后验层不同接口；把 SRN
的复生成元条件、连续变量定义域、过程记忆、三类 RG
和规范异常责任写清；对三代条件定理删除冗余的原始包前提，并补全一般331表示清单。

尚未完成的是：从 Core
必然选择现实量子物理的无附加前提定理、客观实际化机制、量子场完整非微扰构造、规范群和全部粒子参数选择、HGP/FAD
的独立物理必然性、质量层级与三代的无条件解释。它们保留为开放问题，不以更大的版本号代替证明。

\hypertarget{ch-36}{%
\chapter{36. 三代问题：必须在正确的规范群上计算}\label{ch-36}}

\hypertarget{ux51bbux7ed3ux7684ux8868ux793aux914dux7f6e}{%
\section{36.1
冻结的表示配置}\label{ux51bbux7ed3ux7684ux8868ux793aux914dux7f6e}}

考虑 \[
G=SU(N_c)_C\times SU(3)_L\times U(1)_X,\qquad N_c\ge3.
\] 这里 \(SU(3)_L\) 是物理输入，不由一般 UEOT Core
推出。每代轻子取一个无色 \(3_L\)；\(k\) 代夸克取 \((N_c,3_L)\)，其余
\(N_f-k\) 代取
\((N_c,\bar3_L)\)。每个夸克三重态配三个相应色反基本右手共轭
singlet。为一般电荷嵌入，轻子第三分量的必要 singlet
也要加入。所有费米子均用左手 Weyl 记法。

纯 \(SU(3)_L^3\) 异常为 \[
A_L=N_f+N_c(2k-N_f).
\] 故抵消要求 \[
2N_ck=(N_c-1)N_f.
\] 令 \(d=\gcd(2,N_c-1)\)，全部正整数解为 \[
N_f=\frac{2N_c}{d}t,\qquad
k=\frac{N_c-1}{d}t,\qquad t=1,2,\ldots
\tag{36.1}
\] 证明是互素系数的整除性。奇色数给 \(N_f=tN_c\)，偶色数给
\(N_f=2tN_c\)。不能仅因最小解漂亮就宣称自然界必须 \(t=1\)。

这套表示组织有331研究的既有先例；也存在把代际不普适性放到轻子扇区的
flipped
模型，因此它不是所有扩展规范理论的唯一形式。\href{https://arxiv.org/abs/hep-ph/9206242}{原始331论文}、\href{https://arxiv.org/abs/1606.01109}{flipped
331}

\hypertarget{ux4e3aux4ec0ux4e48ux539fux6765ux7684ux5f31ux4e0aux754cux4e0dux7b49ux4e8eux9ad8ux80fdux8981ux6c42}{%
\section{36.2
为什么原来的弱上界不等于高能要求}\label{ux4e3aux4ec0ux4e48ux539fux6765ux7684ux5f31ux4e0aux754cux4e0dux7b49ux4e8eux9ad8ux80fdux8981ux6c42}}

量子主文的条件定理使用 \[
N_f(N_c+1)<22,
\tag{36.2}
\] 这是特定低能 \(SU(2)_L\) 场内容的一圈渐近自由必要界。在该前提下推出
\((3,3,1)\) 的算术是正确的。

但若在尺度 \(M_B\) 以上，弱群已经扩展成
\(SU(3)_L\)，其矢量自相互作用贡献变了。不能以``整个高能理论应该良好''为理由，把低能
\(b_2>0\)
自动升格为必要条件。有效低能耦合可以在有限窗口增长，然后接到另一高能群；负的有效
beta 系数不独自否定这样的匹配。

因此必须重算，而不是通过加强措辞保住三代。

\hypertarget{ux5b8cux6574ux6d3bux8dc3ux573aux5185ux5bb9ux7684ux4e00ux5708ux7cfbux6570}{%
\section{36.3
完整活跃场内容的一圈系数}\label{ux5b8cux6574ux6d3bux8dc3ux573aux5185ux5bb9ux7684ux4e00ux5708ux7cfbux6570}}

约定 \[
\mu\frac{dg}{d\mu}=-\frac{b\,g^3}{16\pi^2}+O(g^5),\quad
b=\frac{11}{3}C_A-\frac23\sum_{\rm Weyl}T(R)
-\frac13\sum_{\rm complex\ scalar}T(R).
\tag{36.3}
\] 假定全部 exotic 费米子活跃，无额外带色标量，\(n_s\) 个复 \(SU(3)_L\)
基本标量且无其他带 \(L\) 荷物质。

对色群，每代有3个基本 Weyl 加3个反基本 Weyl，等于3个 Dirac 色味，因此 \[
b_C=\frac{11N_c}{3}-2N_f.
\tag{36.4}
\] 对 \(SU(3)_L\)，Weyl 二次指数总和为 \(N_f(N_c+1)/2\)，每个复标量指数
\(1/2\)，所以 \[
b_L=11-\frac{N_f(N_c+1)}3-\frac{n_s}{6}.
\tag{36.5}
\] 这与常见扩展弱群中额外夸克对 QCD running 的作用一致，但这里按声明的
\(N_c,N_f\)
配置重新计数。\href{https://arxiv.org/abs/1205.4483}{扩展规范群原始研究}

\(b_C,b_L>0\) 只是两个非 Abelian
因子的``一圈渐近自由方向''。\(U(1)_X\)、Yukawa、标量势和更高阶仍要审查；不得将两个符号称作完整
UV completion。

\hypertarget{ux9ad8ux80fdux5206ux7c7bux4e0eux4e94ux8272ux53cdux4f8b}{%
\section{36.4
高能分类与五色反例}\label{ux9ad8ux80fdux5206ux7c7bux4e0eux4e94ux8272ux53cdux4f8b}}

\protect\hypertarget{g-01}{}{}

\textbf{G-01：三标量配置的全部整数解。}
在(36.1)、\(N_c\ge3\)、\(n_s=3\)、\(b_C>0,b_L>0\) 下， \[
(N_c,N_f,k)=(3,3,1)\quad\text{或}\quad(5,5,2).
\tag{36.6}
\]

\textbf{证明。} \(b_L>0\) 给 \(N_f(N_c+1)<31.5\)。

偶数 \(N_c\ge4\) 时 \(N_f\ge2N_c\)，乘积至少40，排除。奇数 \(N_c\ge7\)
时乘积至少56，排除。只剩3与5。

\(N_c=3\) 时 \(N_f=3t\)，色系数 \(b_C=11-6t>0\)，迫使 \(t=1\)。
\(N_c=5\) 时 \(N_f=5t\)，\(b_L>0\) 迫使 \(30t<31.5\)，故 \(t=1\)，且
\(b_C=25/3>0\)。两点都通过。\(\square\)

具体 \[
\begin{array}{c|ccc}
(N_c,N_f,k)&b_C&b_L&b_2\text{（低能单双重态）}\\\hline
(3,3,1)&5&13/2&19/6\\
(5,5,2)&25/3&1/2&-17/6
\end{array}
\tag{36.7}
\] 最后一列正好显示：高能非 Abelian 两因子有渐近自由方向，不蕴含低能
\(b_2>0\)。

这不是说自然界有五色。它是对``高能物理一致性已经独立选出三色三代''的反模型；五色分支具有不同于实测夸克的电荷时，不能冒充现实拟合。

\hypertarget{ux4e94ux8272ux53cdux4f8bux7684ux5168ux90e8ux5c40ux57dfux5f02ux5e38ux4e5fux80fdux62b5ux6d88}{%
\section{36.5
五色反例的全部局域异常也能抵消}\label{ux4e94ux8272ux53cdux4f8bux7684ux5168ux90e8ux5c40ux57dfux5f02ux5e38ux4e5fux80fdux62b5ux6d88}}

为避免只检查一个异常，引入 \[
Q_u=\frac{N_c+1}{2N_c},\qquad
Q_d=-\frac{N_c-1}{2N_c},\quad Q_\nu=0,\quad Q_e=-1.
\] 取 \(Q=T_3+\beta_8T_8+X\)，则 \[
X_L=-\frac12-\frac{\beta_8}{2\sqrt3},\quad
X_{Q+}=\frac1{2N_c}-\frac{\beta_8}{2\sqrt3},\quad
X_{Q-}=\frac1{2N_c}+\frac{\beta_8}{2\sqrt3}.
\] 第三分量电荷为 \[
q_E=-\frac12-\frac{\sqrt3\beta_8}{2},\quad
q_{J+}=\frac1{2N_c}-\frac{\sqrt3\beta_8}{2},\quad
q_{J-}=\frac1{2N_c}+\frac{\sqrt3\beta_8}{2}.
\] 加入 \(e^c,E^c,u^c,d^c,J^c\) 及所需重数。

\protect\hypertarget{g-02}{}{}

\textbf{G-02：一般色数局域完成。} 在(36.1)下，六项局域异常
\(C^3,L^3,C^2X,L^2X,\mathrm{grav}^2X,X^3\) 全为零。

\textbf{证明。}
色立方的3基本与3反基本逐代抵消；每代物理电荷总和等于三重态的
\(3X\)，故加入相反电荷 singlet 后引力混合与色混合抵消。\(L^2X\) 为 \[
\tfrac12[N_fX_L+N_c kX_{Q+}+N_c(N_f-k)X_{Q-}]=0
\] 其中常数与 \(\beta_8\) 项分别由所选电荷和(36.1)消掉。

对一组三重态，令其电荷为 \(X+t_i\)，\(\sum_i t_i=0\)。加相反电荷 singlet
后立方贡献 \(3X^3-\sum_i(X+t_i)^3=-3X\sum_i t_i^2-\sum_i t_i^3\)。
基本与反基本的二次和相同，三次和相反；总和的第一项由 \(L^2X=0\)
消掉，第二项由 \(L^3=0\) 消掉。\(\square\)

该构造只证明局域异常与一圈计数相容，不证明全局异常、稳定破缺、全部质量或实验可行性。若直接固定
\(Q_u=2/3,Q_d=-1/3\)，就恢复第31章的三色选择条件；这时实测电荷是输入。

\hypertarget{ux54eaux4e9bux524dux63d0ux771fux6b63ux8fd8ux6ca1ux6709ux89e3ux51b3}{%
\section{36.6
哪些前提真正还没有解决}\label{ux54eaux4e9bux524dux63d0ux771fux6b63ux8fd8ux6ca1ux6709ux89e3ux51b3}}

在上述固定族中，\(0\le n_s\le5\) 时保留两点；\(6\le n_s\le41\) 时只剩
\((3,3,1)\)；\(n_s\ge42\) 无解。证明沿用上一不等式，三色点
\(b_L=7-n_s/6\)，五色点
\(b_L=1-n_s/6\)。不能为了排除5而任意塞入6个标量。

若三色已由实验或另一条独立原理给定，则(36.1)和正确高能色系数已经给
\(N_f=3t,\ b_C=11-6t>0\Rightarrow t=1\)。这是一条物理类型更清楚的条件三代路线，但``为什么有这个扩展弱载体''``为什么其表示清单如此''仍开放。

\hypertarget{ch-37}{%
\chapter{37. 从代际不普适性推出新的可检验味约束}\label{ch-37}}

本章不再要求先解决所有三代起源。问一个更直接的问题：一旦采用上述三代载体，它在低能必须留下什么结构？

\hypertarget{ux56faux5b9aux4e00ux4e2aux4e2dux6027ux77e2ux91cfux5019ux9009}{%
\section{37.1
固定一个中性矢量候选}\label{ux56faux5b9aux4e00ux4e2aux4e2dux6027ux77e2ux91cfux5019ux9009}}

假定一个重中性矢量 \(Z'\) 在弱基底对左手普通夸克双重态的耦合是 \[
G_{\rm weak}=g_0I+\Delta g\,P,\qquad P^2=P=P^\dagger,\quad
\operatorname{rank}P=1,\quad \Delta g\ne0.
\tag{37.1}
\] \(P\) 选出与另外两代不同的一代；\(g_0,\Delta g\)
为实数。这个结构来自特定代际表示差异，并非所有 HGP 或 flipped
模型都满足。

假设普通三代质量旋转为酉矩阵，忽略与 exotic
的显著混合，且比较同一左手双重态上的规范耦合。令 \[
V=U_{uL}^\dagger U_{dL},\quad
P_u=U_{uL}^\dagger PU_{uL}=uu^\dagger,\quad
P_d=V^\dagger P_uV=dd^\dagger,\quad d=V^\dagger u.
\tag{37.2}
\] \(u^\dagger u=d^\dagger d=1\)。一般331的树级 FCNC
早有研究；以下具体代数和下界在这里逐步推导，而不把 FCNC
的存在据为原创。\href{https://arxiv.org/abs/hep-ph/9312312}{Liu
原始分析}

\hypertarget{ux79e9ux4e00ux548cux5f0fux672aux77e5ux53c2ux6570ux4ecdux53d7ux5230ux7ea6ux675f}{%
\section{37.2
秩一和式：未知参数仍受到约束}\label{ux79e9ux4e00ux548cux5f0fux672aux77e5ux53c2ux6570ux4ecdux53d7ux5230ux7ea6ux675f}}

\protect\hypertarget{f-01}{}{}

\textbf{F-01。} 对 \(B=G-g_0I=\Delta g\,uu^\dagger\)， \[
B_{ij}B_{kl}=B_{il}B_{kj},\qquad
|B_{ij}|^2=B_{ii}B_{jj},\qquad
B_{12}B_{23}B_{31}=(\Delta g)^3|u_1u_2u_3|^2\in\mathbb R.
\tag{37.3}
\] 并且 \[
B_d=V^\dagger B_uV.
\tag{37.4}
\]

\textbf{证明。} 逐项代入
\(B_{ij}=\Delta g\,u_i\bar u_j\)，所有等式直接约去公共因子。闭合三角相位不随各夸克场重相位改变。\(\square\)

这给出矩阵秩、幅值和循环相位三种可失败约束。若从多个过程提取的同一重矢量耦合不满足它们，固定的秩一候选被否定；不能逐过程各自重新拟合一个不相容的
\(u\)。

\hypertarget{ux4e0aux4e0bux6247ux533aux4e0dux80fdux540cux65f6ux6d88ux6389ux5168ux90e8-fcnc}{%
\section{37.3 上下扇区不能同时消掉全部
FCNC}\label{ux4e0aux4e0bux6247ux533aux4e0dux80fdux540cux65f6ux6d88ux6389ux5168ux90e8-fcnc}}

\protect\hypertarget{f-02}{}{}

\textbf{F-02。} 在非退化质量基底中，若 \(P_u\) 与 \(P_d\) 都对角，则
\(V\) 至少有一个矩阵元模为1，且其所在行列其他元为零。故对
\(c=\max_{ij}|V_{ij}|<1\)，两扇区至少一个有非零非对角规范耦合。

\textbf{证明。} 对角秩一投影必为 \(|e_i\rangle\langle e_i|\)。若
\(V^\dagger|e_i\rangle\langle e_i|V=|e_j\rangle\langle e_j|\)，则
\(V^\dagger e_i=e^{i\theta}e_j\)，因此
\(|V_{ij}|=1\)。酉性强制该行列其余元为零。\(\square\)

这比``可能有
FCNC''更强：在声明假设内，不能通过把特殊代在两个质量基底同时对齐来完全隐藏它。质量简并或非酉轻扇区需重新处理。

\hypertarget{ux4e0dux4f9dux8d56ux672aux77e5ux53d6ux5411ux7684ux5b9aux91cfux4e0bux754c}{%
\section{37.4
不依赖未知取向的定量下界}\label{ux4e0dux4f9dux8d56ux672aux77e5ux53d6ux5411ux7684ux5b9aux91cfux4e0bux754c}}

记 \(\operatorname{off}(A)\)
为删除对角的矩阵，\(p_i=|u_i|^2,q_j=|(V^\dagger u)_j|^2\)。

\protect\hypertarget{f-03}{}{}

\textbf{F-03。} \[
\|\operatorname{off}G_u\|_F^2+
\|\operatorname{off}G_d\|_F^2
\ge |\Delta g|^2(1-c),\qquad c=\max_{ij}|V_{ij}|.
\tag{37.5}
\]

\textbf{证明。} 秩一投影 Frobenius 范数平方为1，故左端除以
\(|\Delta g|^2\) 为 \(2-\sum_i p_i^2-\sum_jq_j^2\)。 设 \(p_i,q_j\)
取最大值。因为概率归一， \(\sum p_i^2\le p_i,\sum q_j^2\le q_j\)。
两个秩一投影之和 \[
|e_i\rangle\langle e_i|+
V|e_j\rangle\langle e_j|V^\dagger
\] 的最大本征值是
\(1+|V_{ij}|\)（在二向量张成子空间直接求二阶特征方程）。所以
\(p_i+q_j\le1+|V_{ij}|\le1+c\)。整理得到(37.5)。\(\square\)

该界通常保守，不是所有 \(V\) 的精确最小值。\(c\)
接近1时它很小，不能因其严格为正就声称实验容易发现。精确候选比较可再对四个实自由度的射影
\(u\) 优化，但数值最优值不替代已证下界。

\hypertarget{eft-ux5339ux914dux548cux53efux89c2ux6d4bux610fux4e49}{%
\section{37.5 EFT
匹配和可观测意义}\label{eft-ux5339ux914dux548cux53efux89c2ux6d4bux610fux4e49}}

采用
\(\mathcal L\supset \tfrac12M_{Z'}^2Z'_\mu Z'^\mu+Z'_\mu J^\mu\)，低于重质量时消去
\(Z'\)： \[
\mathcal L_{\rm EFT}=-\frac{J_\mu J^\mu}{2M_{Z'}^2}
+O(\partial^2/M_{Z'}^4).
\tag{37.6}
\] 对固定手征算符
\(O_{ij}^{u,d}=(\bar q_{Li}\gamma_\mu q_{Lj}) (\bar q_{Li}\gamma^\mu q_{Lj})\)，
\[
C_{ij}^{u,d}=-\frac{(G_{u,d})_{ij}^{\,2}}{2M_{Z'}^2}.
\tag{37.7}
\] 按有向 \(i\ne j\) 同时记录算符与其共轭，则 \[
\sum_{i\ne j}(|C_{ij}^u|+|C_{ij}^d|)
\ge\frac{|\Delta g|^2}{2M_{Z'}^2}(1-c).
\tag{37.8}
\] 若只列独立无向对，双方再除以2。

这是匹配尺度的系数关系，不是介子混合实测量的直接和式。运行、Hadron
矩阵元、SM干涉、右手耦合和其他重粒子贡献须另算；多矢量可相消。本次没有从这些式子预测
\(M_{Z'}\)
或具体衰变率。它解决的是``这个候选能否完全不留味结构''的问题，而不是已经找到了新粒子。

\hypertarget{ux8f7bux5b50ux666eux9002ux6027ux5fc5ux987bux662fux5b8cux6574ux77e9ux9635ux6761ux4ef6}{%
\section{37.6
轻子普适性必须是完整矩阵条件}\label{ux8f7bux5b50ux666eux9002ux6027ux5fc5ux987bux662fux5b8cux6574ux77e9ux9635ux6761ux4ef6}}

若某个重中性矢量在三代普通轻子上确实具有相同表示耦合，且忽略 exotic
混合后保持酉的轻子质量旋转，那么弱基底耦合应为 \[
G_\ell=g_\ell I.
\] 对任意质量基底旋转
\(U\)，\(U^\dagger G_\ell U=g_\ell I\)，于是该项没有由这种旋转产生的树级味改变中性流。这是一项可用的候选区分条件，不是所有
HGP、flipped 331 或含显著重轻混合模型的结论。

检查各对角元相同还不够。以两维子块为例，取
\(G=\operatorname{diag}(g_1,g_2)\)，用 \[
U=\frac1{\sqrt2}\begin{pmatrix}1&1\\1&-1\end{pmatrix}
\] 旋转后得到 \[
U^\dagger GU=
\begin{pmatrix}
(g_1+g_2)/2&(g_1-g_2)/2\\
(g_1-g_2)/2&(g_1+g_2)/2
\end{pmatrix}.
\] 两个对角元相等，但只要
\(g_1\ne g_2\)，非对角耦合仍存在。三维中嵌入这个子块即可得到同类反例。因此``普适''应检验整个矩阵是标量单位阵，或检验其谱与所有相关矩阵元，不能只对角拟合后宣称味改变消失。低能过程的
running、算符混合与其他介子/重矢量贡献仍按第 37.5 节的合同处理。

\hypertarget{ch-38}{%
\chapter{38. 从相干相位推出不可任意调节的动力学约束}\label{ch-38}}

\hypertarget{srn-ux7684ux6700ux5c0fux7269ux7406ux5b9aux4e49}{%
\section{38.1 SRN
的最小物理定义}\label{srn-ux7684ux6700ux5c0fux7269ux7406ux5b9aux4e49}}

取正交指针分支 \(P_a\)。强记录无扰 SRN
要求通道固定每个分支内部全部态，不仅保持总人口。有限维完全正保迹性等价于
\[
\mathcal A_K(\rho)=\sum_{ab}K_{ab}P_a\rho P_b,\qquad
K\succeq0,\quad K_{aa}=1.
\tag{38.1}
\] 证明：任意 Kraus
在每个被逐点固定的纯态上只能乘标量；线性性迫使同一分支的标量一致。于是
\(L_\mu=\sum_a l_{\mu a}P_a\)，
\(K_{ab}=\sum_\mu l_{\mu a}\bar l_{\mu b}\) 是 Gram 矩阵。反向用 Gram
分解构造 Kraus 即可。

相干父系统经公共读出可完全看不见
\(K_{ab},a\ne b\)。因此新的测量物理必须测相干，而不能只数公共记录。

\hypertarget{ux8fdeux7eedux751fux6210ux5143ux4e0eux5faaux73afux76f8ux4f4d}{%
\section{38.2
连续生成元与循环相位}\label{ux8fdeux7eedux751fux6210ux5143ux4e0eux5faaux73afux76f8ux4f4d}}

时间齐次连续半群有 \[
K_{ab}(t)=e^{t\lambda_{ab}},\qquad
\lambda_{ab}=-\gamma_{ab}+i\omega_{ab}.
\]
\(\gamma_{ab}=\gamma_{ba}\ge0,\omega_{ab}=-\omega_{ba}\)，对角为零。全时间
CP 的条件是 \(\Lambda\) 在零和子空间半正定。

对三个分支定义 \[
\Omega_{abc}=\omega_{ab}+\omega_{bc}+\omega_{ca}.
\tag{38.2}
\] 改变分支能量或采用时间线性的对角旋转参考系，只给 \(\omega_{ab}\)
加一个差 \(h_a-h_b\)，其循环和为零。因此 \(\Omega\) 不能被这样的
Hamiltonian 重标定消掉。它也不依赖 Lindblad 表示中的冗余变换，因为
\(\lambda_{ab}\) 是固定物理通道的系数。

这一障碍与 Fagnola、Gough、Nurdin、Viola 已研究的 Hamiltonian
obstruction 同类；不是 UEOT
首次发现。\href{https://arxiv.org/abs/1811.11784}{原始研究}

\hypertarget{ux4e09ux5206ux652fux7684ux9510-cp-ux754c}{%
\section{38.3 三分支的锐 CP
界}\label{ux4e09ux5206ux652fux7684ux9510-cp-ux754c}}

\protect\hypertarget{m-01}{}{}

\textbf{M-01。} 任意三个分支必须满足 \[
\boxed{\ 
\Omega_{abc}^{\,2}
\le
4\gamma_{ac}\gamma_{bc}
-(\gamma_{ac}+\gamma_{bc}-\gamma_{ab})^2.\ }
\tag{38.3}
\] 对恰好三分支，该条件连同各 \(\gamma\ge0\) 为生成元 CP
的充分必要条件。

\textbf{证明。} 以 \(c\) 为锚，令
\(G_{ij}=\lambda_{ij}-\lambda_{ic}-\lambda_{cj}\)。 其非零 \(2\times2\)
块为 \[
G=
\begin{pmatrix}
2\gamma_{ac}&
\gamma_{ac}+\gamma_{bc}-\gamma_{ab}+i\Omega_{abc}\\
\gamma_{ac}+\gamma_{bc}-\gamma_{ab}-i\Omega_{abc}&
2\gamma_{bc}
\end{pmatrix}.
\]
零和子空间的正性等价于此矩阵半正定。二阶主子式非负正是(38.3)；两个对角非负给其余条件。反向以这个正
Gram 矩阵构造
\(\lambda_{ij}=G_{ij}-(G_{ii}+G_{jj})/2+i(\theta_i-\theta_j)\)，
逐项指数可由 Schur 正性证明为相关矩阵，故所有 \(t\ge0\)
CPTP。\(\square\)

分支更多时，全部三角条件仍只是必要，不替代完整 Gram 矩阵的高阶主子式。

\hypertarget{ux76f8ux4f4dux79efux5fc5ux987bux652fux4ed8ux76f8ux5e72ux635fux5931}{%
\section{38.4
相位积必须支付相干损失}\label{ux76f8ux4f4dux79efux5fc5ux987bux652fux4ed8ux76f8ux5e72ux635fux5931}}

记 \(S=\gamma_{ab}+\gamma_{bc}+\gamma_{ca}\)。右端可写成 \[
S^2-2(\gamma_{ab}^2+\gamma_{bc}^2+\gamma_{ca}^2)\le S^2/3.
\] 因此：

\protect\hypertarget{m-02}{}{}

\textbf{M-02。} \[
S\ge\sqrt3\,|\Omega_{abc}|.
\tag{38.4}
\] 对循环相干积 \(B(t)=K_{ab}(t)K_{bc}(t)K_{ca}(t)\)， \[
|B(t)|=e^{-St}
\le e^{-\sqrt3|\Omega_{abc}|t}
=e^{-\sqrt3|\arg_{\rm unwrap} B(t)|}.
\tag{38.5}
\] 证明由(38.3)和平方和下界直接得到。相位必须连续解缠；不能用模 \(2\pi\)
的主值制造虚假的大时间结论。\(\square\)

若给定有限时刻的三个相干，完全正性还有不需要 Markov 假设的检验： \[
1-|K_{ab}|^2-|K_{bc}|^2-|K_{ca}|^2
+2\operatorname{Re}(K_{ab}K_{bc}K_{ca})\ge0.
\tag{38.6}
\] 这是 \(3\times3\)
相关矩阵行列式。半群界与单时刻界的假设不同，违背前者可能意味着记忆或时变动力学，而不一定违背
QM。

\hypertarget{ux4e3aux4ec0ux4e48ux975eux96f6ux5faaux73afux76f8ux4f4dux4ecdux4e0dux8bc1ux660eux91cfux5b50ux566aux58f0}{%
\section{38.5
为什么非零循环相位仍不证明``量子噪声''}\label{ux4e3aux4ec0ux4e48ux975eux96f6ux5faaux73afux76f8ux4f4dux4ecdux4e0dux8bc1ux660eux91cfux5b50ux566aux58f0}}

对实对角 Hermitian 噪声
\(L_\mu=\operatorname{diag}(\ell_{\mu a}),\ell_{\mu a}\in\mathbb R\)，
\[
\lambda_{ab}=-\tfrac12\sum_\mu(\ell_{\mu a}-\ell_{\mu b})^2
-i(h_a-h_b),
\] 所以 \(\Omega=0\)。反向，若所有循环为零，则 \(\omega_{ab}\)
是一个分支势差；锚 Gram 的剩余部分实正，可分解成实向量，于是存在这种
Gaussian 扩散表示。

但经典 Poisson 跳相位给 \[
\mathcal L(\rho)=r(U\rho U^\dagger-\rho),\qquad
U=\operatorname{diag}(e^{i\theta_a}),
\] \[
\lambda_{ab}=r(e^{i(\theta_a-\theta_b)}-1).
\tag{38.7}
\] 取 \(\theta=(0,2\pi/3,4\pi/3)\)，各 \(\gamma=3r/2\)，
\(|\Omega|=3\sqrt3r/2\)，恰好饱和(38.4)。

因此非零循环相位排除``只有对角 Gaussian 白噪声加
Hamiltonian''的模型，但连锐界饱和都可来自经典 Poisson
噪声。这个反例决定了下一章为什么需要更强的四能级判据。

\hypertarget{ch-39}{%
\chapter{39. 四能级判据：排除整个随机相位噪声类}\label{ch-39}}

\hypertarget{ux95eeux9898ux4e0eux6700ux5c0fux7ef4ux6570}{%
\section{39.1
问题与最小维数}\label{ux95eeux9898ux4e0eux6700ux5c0fux7ef4ux6570}}

任意经典对角相位噪声在一个固定时刻可写 \[
\mathcal R(\rho)=\mathbb E[U_\theta\rho U_\theta^\dagger],
\quad U_\theta=\operatorname{diag}(z_1,\ldots,z_m),\quad |z_a|=1,
\] 所以 \[
K_{\rm cl}=\mathbb E[zz^\dagger].
\tag{39.1}
\] 这里允许任意相关、非 Gaussian、非 Markov 的相位分布，不只指数衰减。

\protect\hypertarget{m-03}{}{}

\textbf{M-03：单时刻的最小分支数为4。} 对 \(m\le3\)，每个相关矩阵
\(K\succeq0,K_{aa}=1\) 都是(39.1)形式。因此只有二或三分支的单时刻完整
SRN 层析也不能排除所有经典随机相位噪声。

\textbf{证明。} 写秩为 \(r\) 的 \(K=VV^\dagger\)。若 \(r^2>m\)，则
\(r\times r\) Hermitian 矩阵空间中存在非零 \(A\)，使
\(v_aAv_a^\dagger=0\) 对全部行成立。足够小 \(\epsilon>0\) 下
\(V(I\pm\epsilon A)V^\dagger\) 为两个不同的相关矩阵，其平均为 \(K\)。故
\(K\) 不极端。于是极端点必要满足 \(r^2\le m\)。当 \(m\le3\) 时只可
\(r=1\)，对角一迫使它为
\(zz^\dagger\)。有限维紧凸集是其极端点的凸包，得结论。\(\square\)

这一维数界与下面的四维矩阵已有明确先例，尤其是
Buscemi、Chiribella、D'Ariano
2005年的研究。本章的工作是给出可直接使用的解析见证、校准界和计数协议，不把已知反例重新命名成新自然定律。\href{https://doi.org/10.1103/PhysRevLett.95.090501}{原论文}

\hypertarget{ux4e00ux4e2aux53eaux9700ux4e8cux7ef4ux73afux5883ux7684ux7269ux7406ux901aux9053}{%
\section{39.2
一个只需二维环境的物理通道}\label{ux4e00ux4e2aux53eaux9700ux4e8cux7ef4ux73afux5883ux7684ux7269ux7406ux901aux9053}}

取 \[
V=\begin{pmatrix}
1&0\\0&1\\1/\sqrt2&1/\sqrt2\\1/\sqrt2&i/\sqrt2
\end{pmatrix},\qquad K_*=VV^\dagger.
\tag{39.2}
\] 各行范数一，故 \(K_*\) 是合法相关矩阵，秩2。Kraus 可直接取 \[
L_1=\operatorname{diag}(1,0,1/\sqrt2,1/\sqrt2),\quad
L_2=\operatorname{diag}(0,1,1/\sqrt2,i/\sqrt2).
\tag{39.3}
\] \(\sum L_\mu^\dagger L_\mu=I\)，所以它是标准 CPTP-SRN。

等距实现为 \[
|a\rangle|0\rangle_E\longmapsto |a\rangle|e_a\rangle_E,
\] 其中 \(|e_a\rangle\) 的两个分量取 \(V\) 的第 \(a\) 行。环境 Gram 满足
\(\langle e_b|e_a\rangle=K_{*,ab}\)。四个受控环境酉操作均可把
\(|0\rangle_E\) 送到相应
\(|e_a\rangle_E\)，所以联合演化可幺正完成。这是四能级系统加一个环境
qubit 的有限模型。

\(K_{*,12}=0\)，它不是连续有限速率时间齐次 SRN
半群在有限时刻的核；本章不依赖 Markov 假设。

\hypertarget{ux89e3ux6790ux5206ux79bbux89c1ux8bc1ux4e0eux4e25ux683cux5e38ux6570}{%
\section{39.3
解析分离见证与严格常数}\label{ux89e3ux6790ux5206ux79bbux89c1ux8bc1ux4e0eux4e25ux683cux5e38ux6570}}

定义 \[
B=\begin{pmatrix}
-1/\sqrt2&-1/\sqrt2&1&0\\
-1/\sqrt2&-i/\sqrt2&0&1
\end{pmatrix},\qquad W=B^\dagger B.
\tag{39.4}
\] 有 \(BV=0\)，所以 \[
\operatorname{Tr}(WK_*)=0.
\tag{39.5}
\]

\protect\hypertarget{m-04}{}{}

\textbf{M-04：锐经典下界。} \[
\boxed{\operatorname{Tr}(WK_{\rm cl})\ge
\kappa:=3-2\sqrt2=0.171572875\ldots}
\tag{39.6}
\] 对整个(39.1)类成立，且常数可达到。

\textbf{证明。} 先考虑任意相位向量 \(z\)。令 \[
r_1=\left|z_3-\frac{z_1+z_2}{\sqrt2}\right|,\qquad
r_2=\left|z_4-\frac{z_1+iz_2}{\sqrt2}\right|.
\] 则 \(z^\dagger Wz=r_1^2+r_2^2\)。记
\(w=z_1\bar z_2\)，\(|w|=1\)。由反三角不等式， \[
|\operatorname{Re}w|
=\left|\left|\frac{z_1+z_2}{\sqrt2}\right|^2-1\right|
\le(\sqrt2+1)r_1,
\] \[
|\operatorname{Im}w|
=\left|\left|\frac{z_1+iz_2}{\sqrt2}\right|^2-1\right|
\le(\sqrt2+1)r_2.
\] 平方相加，利用 \((\operatorname{Re}w)^2+(\operatorname{Im}w)^2=1\)，
得 \(r_1^2+r_2^2\ge1/(\sqrt2+1)^2=3-2\sqrt2\)。
对概率分布取平均即得(39.6)。取 \(z=(1,1,1,(1+i)/\sqrt2)\)，有
\(r_1=\sqrt2-1,r_2=0\)，达到下界。\(\square\)

于是 \(K_*\) 无法由任何经典随机对角相位分布产生。并且 \[
\|W\|_F=3,\qquad \|W\|_{\rm op}=2+1/\sqrt2=:w_{\max}.
\] 由 Cauchy--Schwarz， \[
\inf_{K_{\rm cl}}\|K_*-K_{\rm cl}\|_F\ge\kappa/3.
\tag{39.7}
\] 这是一个证明出来的保守距离界，不是优化器给出的猜测。

\hypertarget{ux53efux627fux53d7ux591aux5c11ux989dux5916ux9000ux76f8ux5e72}{%
\section{39.4
可承受多少额外退相干}\label{ux53efux627fux53d7ux591aux5c11ux989dux5916ux9000ux76f8ux5e72}}

\protect\hypertarget{m-05}{}{}

\textbf{M-05。} 对 \(K_\epsilon=(1-\epsilon)K_*+\epsilon I\)， \[
\operatorname{Tr}(WK_\epsilon)=4\epsilon.
\] 因此本见证在 \[
\epsilon<\frac{3-2\sqrt2}{4}=0.0428932188\ldots
\tag{39.8}
\] 时仍严格排除(39.1)。

这里 \(I\)
表示完全指针去相干核，而不是恒等通道核；后者是全一矩阵。对其他噪声形状需要另算，4.29\%不是普适硬件错误阈值。超过此值只是本见证不再保证检出，不能据此判断通道已经变成随机酉。

若 raw 层析估计满足 \(\|\widehat K-K\|_F\le\eta\)，则 \[
\operatorname{Tr}(WK)\le \operatorname{Tr}(W\widehat K)+3\eta.
\] 上界小于 \(\kappa\) 才是有误差的排除；PSD 投影不能替代该置信责任。

\hypertarget{ux4e00ux4e2aux8f93ux5165ux4e00ux4e2aux6548ux5e94ux7684ux96f6ux4e8bux4ef6ux534fux8bae}{%
\section{39.5
一个输入、一个效应的零事件协议}\label{ux4e00ux4e2aux8f93ux5165ux4e00ux4e2aux6548ux5e94ux7684ux96f6ux4e8bux4ef6ux534fux8bae}}

制备 \(|+\!_4\rangle=(1,1,1,1)^T/2\)。SRN 输出为 \(\rho_K=K/4\)。
取合法二元测量效应 \[
E=\frac{W}{w_{\max}},\qquad 0\preceq E\preceq I.
\tag{39.9}
\] 则 \[
p_E(K_*)=0,\qquad
p_E(K_{\rm cl})\ge
p_{\rm cl}:=\frac{\kappa}{4w_{\max}}
=0.0158446719\ldots
\tag{39.10}
\] 可以对 \(W\) 作谱测量后按本征值 \(\lambda/w_{\max}\)
接受，也可用一个辅助比特实现二元
POVM。复杂性转移到了可信制备与效应校准，并未消失。

\protect\hypertarget{e-01}{}{}

\textbf{E-01：有限样本反证。} 在独立、同分布且模型与测量已冻结时，\(n\)
次实验零接受事件在任意经典假设下的概率至多
\((1-p_{\rm cl})^n\)。因此理想校准下 \(n=188\)
的零事件结果在显著性0.05下排除该整个经典相位类。

\textbf{证明。} Bernoulli 零事件概率对 \(p\) 单调递减；代入最小
\(p_{\rm cl}\)，取 \(n\ge\log(0.05)/\log(1-p_{\rm cl})\)
并向上取整。\(\square\)

若观察到 \(k>0\) 个事件，使用预先固定的 \[
p\text{-value}=\sum_{j=0}^k{n\choose j}p_{\rm cl}^{\,j}(1-p_{\rm cl})^{n-j}.
\tag{39.11}
\] 对保守误差 \(\delta\)，将 \(p_{\rm cl}\) 换成
\(\max(0,p_{\rm cl}-\delta)\)。188不包含校准次数、SRN检验、漂移控制或重复性；不是设备现状的速度承诺。

\hypertarget{ux5fc5ux987bux9632ux6b62ux666eux901aux975eux5bf9ux89d2ux9149ux64cdux4f5cux4f2aux88c5ux6210ux901aux8fc7}{%
\section{39.6
必须防止普通非对角酉操作伪装成通过}\label{ux5fc5ux987bux9632ux6b62ux666eux901aux975eux5bf9ux89d2ux9149ux64cdux4f5cux4f2aux88c5ux6210ux901aux8fc7}}

一个任意酉操作可将该单一输入送到
\(\ker W\)。所以(39.10)不是对所有随机酉通道的无条件检验。需要独立验证``指针态保持''这一物理类。

若随机酉通道精确保持每个
\(|a\rangle\langle a|\)，秩一正项求和迫使其每个酉分量都在该基底对角，才归入(39.1)。

\protect\hypertarget{e-02}{}{}

\textbf{E-02：近似人口保持的稳健界。} 若一个一般随机酉通道
\(\sum_\nu p_\nu U_\nu(\cdot)U_\nu^\dagger\) 对每个基态满足
\(\Pr(a|a)\ge1-\ell\)，则对任意固定输入，它与某个对角随机酉通道输出的迹距离至多
\[
\sqrt{2m\ell}.
\tag{39.12}
\] 因此四能级协议对该更宽假设类的概率下界为
\(p_{\rm cl}-\sqrt{8\ell}\)，再减制备/测量误差。

\textbf{证明。} 为每个 \(U_\nu\) 取对角酉
\(D_{\nu,aa}=U_{\nu,aa}/|U_{\nu,aa}|\)，零元相位任意。 \[
\|U_\nu-D_\nu\|_F^2=2\sum_a(1-|U_{\nu,aa}|)
\le2\sum_a(1-|U_{\nu,aa}|^2).
\] 按 \(\nu\) 平均得不超过
\(2m\ell\)。纯态输出迹距离不超过对应向量距离，再以算符范数小于 Frobenius
范数和 Jensen 不等式求平均；混态由凸性延拓。\(\square\)

这个界保守但诚实：人口只保持99\%远不足以使用理想188次结论。以零额外误差计，甚至需要
\(\ell<p_{\rm cl}^2/8\)
才让这个保守下界为正。可以研究更精细校准界，但不能省略模型验证。

\hypertarget{ux8fd9ux89e3ux51b3ux4e86ux4ec0ux4e48ux6ca1ux6709ux89e3ux51b3ux4ec0ux4e48}{%
\section{39.7
这解决了什么，没有解决什么}\label{ux8fd9ux89e3ux51b3ux4e86ux4ec0ux4e48ux6ca1ux6709ux89e3ux51b3ux4ec0ux4e48}}

解决：单时刻二/三能级 SRN
数据无法完成的``排除所有经典随机相位噪声''，可以在四能级完成，并有显式物理通道、锐见证、噪声半径和有限数据协议。

没有解决：排除所有普通量子环境。\(K_*\)
本身就是标准量子环境可产生的通道；它通过检验说明该经典噪声类不够，而非标准
QM 错了，更非客观坍缩已被发现。这是量子噪声谱学和 QCVV
的候选新工具，不是已经确立的新自然定律。

\hypertarget{ch-40}{%
\chapter{40.
实际化动力学：把``最终单分支''从假设推进成条件定理}\label{ch-40}}

\hypertarget{ux5148ux6392ux9664ux4e00ux79cdux770bux4f3cux5408ux7406ux5374ux53efux4f20ux4fe1ux53f7ux7684ux505aux6cd5}{%
\section{40.1
先排除一种看似合理却可传信号的做法}\label{ux5148ux6392ux9664ux4e00ux79cdux770bux4f3cux5408ux7406ux5374ux53efux4f20ux4fe1ux53f7ux7684ux505aux6cd5}}

设有人根据``叠加越强，实际化越快''提出单比特非线性规则 \[
\dot x=-\gamma(1-z^2)x,\quad
\dot y=-\gamma(1-z^2)y,\quad \dot z=0
\tag{40.1}
\] 并要求它分别作用于每条纯态制备分支，再按通常 Born
权重平均。这不是对线性量子通道的等价改写。

\protect\hypertarget{c-01}{}{}

\textbf{C-01：明确的远端制备矛盾。} 相同密度矩阵 \(\rho=(I+0.6X)/2\)
有两种分解： \[
\mathcal E_1=\{0.8,|+x\rangle;\ 0.2,|-x\rangle\},
\] \[
\mathcal E_2=\{1/2,\mathbf r_+=(0.6,0,0.8);\
1/2,\mathbf r_-=(0.6,0,-0.8)\}.
\] (40.1)逐分支作用给平均 \[
x_1(t)=0.6e^{-\gamma t},\qquad
x_2(t)=0.6e^{-0.36\gamma t}.
\tag{40.2}
\] 两者不同。

\textbf{证明与操作实现。} 共享纯态 \[
|\Psi\rangle=\sqrt{0.8}|0\rangle_A|+x\rangle_B
+\sqrt{0.2}|1\rangle_A|-x\rangle_B.
\] Alice 选择测 \(Z_A\) 或 \(X_A\)，不告知结果，分别制备 Bob
的上述两分解，边缘均为同一 \(\rho\)。在 \(\gamma t=1\) 后 Bob 测
\(X\)，正结果概率差 \[
\Delta p=0.3(e^{-0.36}-e^{-1})=0.0989390655\ldots
\] 非零，因此结合标准远端制备与分支平均规则会允许信号。\(\square\)

这不是``所有数学非线性都必然超光速''的无条件命题；矛盾使用了明确的纯态动力学、复合态和概率解释。标准无信号约束对此已有原始定理。\href{https://arxiv.org/abs/quant-ph/0102125}{Simon--Bužek--Gisin}

直接把 ORC
冗余、态熵或分支权重塞入确定性坍缩速率，也必须过这种分解独立性检验；若不声明复合系统和随机性，单系统轨迹漂亮不构成合法新物理。

\hypertarget{ux4e00ux4e2aux8303ux6570ux4fddux6301ux65e0ux6761ux4ef6ux7ebfux6027ux7684ux968fux673aux5019ux9009}{%
\section{40.2
一个范数保持、无条件线性的随机候选}\label{ux4e00ux4e2aux8303ux6570ux4fddux6301ux65e0ux6761ux4ef6ux7ebfux6027ux7684ux968fux673aux5019ux9009}}

有限维取 Hermitian \(A\)、\(\gamma>0\)，采用 Itō 方程 \[
d|\psi\rangle=
\left[-\frac{i}{\hbar}Hdt
-\frac{\gamma}{2}(A-\bar a)^2dt
+\sqrt\gamma(A-\bar a)dW_t\right]|\psi\rangle,
\quad \bar a=\langle\psi|A|\psi\rangle.
\tag{40.3}
\] \(W_t\) 是相对于物理概率测度的标准 Wiener 过程。\(A\) 无量纲时
\(\gamma\) 单位为时间倒数；一般 \([\gamma]=[A]^{-2}[t]^{-1}\)。

这是标准连续测量/随机约化模型的一类。本次不把方程本身称为原创，而完整推导它在
UEOT
测量问题上能关闭什么、不能关闭什么。\href{https://arxiv.org/abs/1301.3626}{连续量子测量}、\href{https://arxiv.org/abs/quant-ph/0107153}{鞅约化原始研究}

\protect\hypertarget{c-02}{}{}

\textbf{C-02：范数与线性平均。} 归一初态保持归一，且 \[
\bar\rho=\mathbb E|\psi\rangle\langle\psi|
\] 满足 \[
\frac{d\bar\rho}{dt}
=-\frac{i}{\hbar}[H,\bar\rho]
-\frac{\gamma}{2}[A,[A,\bar\rho]].
\tag{40.4}
\]

\textbf{证明。} 对 \(\langle\psi|\psi\rangle\) 使用 Itō 乘积。实漂移为
\(-\gamma\langle(A-\bar a)^2\rangle dt\)，二次变差为相反项；随机项
\(2\sqrt\gamma\langle A-\bar a\rangle dW=0\)，所以范数不变。
对纯态投影作同样计算， \[
d\rho=-i[H,\rho]dt/\hbar-\tfrac\gamma2[A,[A,\rho]]dt
+\sqrt\gamma\{A-\bar a,\rho\}dW.
\]
适应有界噪声积分均值为零，得线性方程(40.4)。有限维系数在单位球面光滑有界，给全局过程；没有通过平均偷偷保留分解依赖项。\(\square\)

局部 \(A\) 的平均通道
CPTP，与普通复合规则相容，故未条件化统计不允许信号。轨迹非线性本身不构成
C-01 那种矛盾。

\hypertarget{ux8bc1ux660eux53d1ux751fux8c31ux5206ux652fux5438ux6536ux800cux975eux5047ux8bbeux5b83ux53d1ux751f}{%
\section{40.3
证明发生谱分支吸收，而非假设它发生}\label{ux8bc1ux660eux53d1ux751fux8c31ux5206ux652fux5438ux6536ux800cux975eux5047ux8bbeux5b83ux53d1ux751f}}

以下增加关键条件 \([H,A]=0\)。令
\(A=\sum_a aP_a\)，不同标签对应不同本征值。\(P_a\)
可以有简并，结论只到谱子空间。

定义 \(p_a(t)=\langle P_a\rangle_t\) 和
\(V(t)=\langle A^2\rangle_t-\bar a_t^2\)。

\protect\hypertarget{c-03}{}{}

\textbf{C-03：单谱分支与 Born 权重。} (40.3)下， \[
dp_a=2\sqrt\gamma(a-\bar a)p_a\,dW_t,
\tag{40.5}
\] 每个 \(p_a\) 是有界鞅，且几乎处处最终趋于一个 one-hot 向量， \[
p_a(t)\longrightarrow 1_{\{A_\infty=a\}},\qquad
\Pr(A_\infty=a)=p_a(0).
\tag{40.6}
\]

\textbf{完整证明。} 由轨迹密度式取 \(P_a\) 迹，交换条件消去 Hamiltonian
与双交换漂移，得到(40.5)。所以 \(p_a\) 有界，鞅收敛定理给几乎处处和
\(L^1\) 收敛。

由(40.5)， \(d\bar a=2\sqrt\gamma VdW\)；再对
\(\langle A^2\rangle-\bar a^2\) 作 Itō 运算得 \[
dV=2\sqrt\gamma\langle(A-\bar a)^3\rangle dW
-4\gamma V^2dt.
\tag{40.7}
\] \(V\) 是有界非负超鞅，故收敛。其均值 \(v(t)=\mathbb EV(t)\) 满足
\(v'\le-4\gamma v^2\)，因此 \[
\mathbb EV(t)\le\frac{V(0)}{1+4\gamma V(0)t}\longrightarrow0.
\tag{40.8}
\] 由有界收敛，极限 \(V_\infty\)
非负且均值零，故几乎处处为零。有限个不同本征值的概率分布方差为零，当且仅当全部权重集中在一个本征值上。鞅极限因而是
one-hot。最后 \(\mathbb E p_a(\infty)=p_a(0)\)，得到(40.6)。\(\square\)

这比第3章``若无偏鞅最终单分支，则有 Born
权重''前进一步：对一个具体动力学，最终单分支也被证明了。仍然没有从 Core
推出应选这个 \(A,\gamma\) 或此随机方程。若 \([H,A]\ne0\)，Hamiltonian
可持续把态带离本征分支，上述吸收证明不能照搬。

这里``谱分支吸收''是 \(t\to\infty\)
的几乎处处极限，不是有限时刻精确投影。事实上对 \(p_a(0)>0\)，令
\(b_a(t)=2\sqrt\gamma[a-\bar a_t]\)，则 \[
p_a(t)=p_a(0)\exp\left[
\int_0^t b_a(s)dW_s-\frac12\int_0^t b_a(s)^2ds\right]>0
\] 在每个有限时间几乎处处成立，因为有限谱使 \(b_a\)
有界。零初权重保持零，而多个非零初权重通常只在无限时极限分离。有限实验需要下面的误差或置信界，不能把
C-03 的渐近陈述读成精确有限时坍缩。

\hypertarget{ux4e00ux4e2aux53efux8ba1ux7b97ux7684ux6709ux9650ux65f6ux95f4ux95edux5408ux7387}{%
\section{40.4
一个可计算的有限时间闭合率}\label{ux4e00ux4e2aux53efux8ba1ux7b97ux7684ux6709ux9650ux65f6ux95f4ux95edux5408ux7387}}

\protect\hypertarget{c-04}{}{}

\textbf{C-04：成对不确定性衰减。} 继续采用第 40.3 节的有限谱与
\([H,A]=0\) 条件，则 \[
\mathbb E\sqrt{p_a(t)p_b(t)}
=\sqrt{p_a(0)p_b(0)}
 e^{-\frac\gamma2(a-b)^2t}.
\tag{40.9}
\] 因而 \[
\mathbb E[1-\max_a p_a(t)]
\le\sum_{a<b}\sqrt{p_a(0)p_b(0)}
 e^{-\frac\gamma2(a-b)^2t}.
\tag{40.10}
\]

\textbf{证明。} 对 \(q_{ab}=\sqrt{p_ap_b}\) 使用 (40.5) 与 Itō
公式。两种乘法噪声的差为 \(2\sqrt\gamma(a-b)\)，由二阶项得到漂移
\(-\frac\gamma2(a-b)^2q_{ab}\)；有界随机积分均值为零。于是其均值解一阶线性方程，得到
(40.9)。正初概率在有限时间保持正；零初概率恒为零，边界可分别处理。对任意概率向量，把最大项记为
\(p_{\max}\)，与该项组成的成对最小值已给 \(1-p_{\max}\)，故 \[
1-p_{\max}\le\sum_{a<b}\min(p_a,p_b)
\le\sum_{a<b}\sqrt{p_ap_b}.
\] 取均值得 (40.10)。\(\square\)

若 \(A\) 有简并，选共同本征基 \(|a,\alpha\rangle\)，使 \[
A|a,\alpha\rangle=a|a,\alpha\rangle,
\qquad H|a,\alpha\rangle=E_{a\alpha}|a,\alpha\rangle.
\] 无条件主方程相应给 \[
\bar\rho_{a\alpha,b\beta}(t)=
 e^{-i(E_{a\alpha}-E_{b\beta})t/\hbar}
 e^{-\frac\gamma2(a-b)^2t}
\bar\rho_{a\alpha,b\beta}(0).
\] 一个 A 的谱标签不必指定唯一 H
能量；同一谱子空间中的相干可以继续酉演化。模型关联的是不同谱分支的确定化与非选择相干损失速率，不能据此宣布简并子空间内也已经选出唯一纯态。

可用无需 Euler 时间步的精确过滤表示模拟这些权重。先按 \(p_a(0)\)
抽取标签 \(A_*\)，在该模拟空间取独立标准 Brownian 过程 \(B_t\)，生成 \[
Y_t=2\sqrt\gamma\,a_{A_*}t+B_t,
\] 并用 Bayes 公式给 \[
p_a(t)=
\frac{p_a(0)e^{2\sqrt\gamma aY_t-2\gamma a^2t}}
{\sum_b p_b(0)e^{2\sqrt\gamma bY_t-2\gamma b^2t}}.
\tag{40.11}
\] 在由观察记录生成的过滤中，创新过程 \[
W_t=Y_t-2\sqrt\gamma\int_0^t\bar a_s\,ds
\] 是物理混合测度下的 Wiener 过程，权重满足 (40.5)。这里的 \(B_t\)
与创新 \(W_t\)
承担不同角色；抽标签只是同一过滤律的模拟表示，不是以``预先选结果''替代
C-03 的动力学吸收证明。

\hypertarget{ux65b0ux7269ux7406ux5019ux9009ux5fc5ux987bux540cux65f6ux652fux4ed8ux80fdux91cfux4ee3ux4ef7}{%
\section{40.5
新物理候选必须同时支付能量代价}\label{ux65b0ux7269ux7406ux5019ux9009ux5fc5ux987bux540cux65f6ux652fux4ed8ux80fdux91cfux4ee3ux4ef7}}

由 (40.4) 及合法的迹循环，额外噪声对能量的贡献为 \[
\frac{d\langle H\rangle}{dt}\bigg|_{\rm noise}
=-\frac\gamma2\langle[A,[A,H]]\rangle.
\tag{40.12}
\] 这条恒等式本身不要求 \([H,A]=0\)。若在一维非相对论延拓中取
\(A=X\)、\(H=P^2/(2m)+V(X)\)，则要另行给无界算符的共同域、状态能量条件及交换迹与生成元的正当性。在这些条件下
\[
[X,[X,H]]=-\hbar^2/m,
\qquad\frac{d\langle H\rangle}{dt}=\frac{\gamma\hbar^2}{2m}.
\tag{40.13}
\] 位置表象的两个点相距 \(d\ne0\)，该附加双交换耗散的相干衰减率为
\(\Gamma(d)=\gamma d^2/2\)，故 \[
\boxed{\frac{d\langle H\rangle}{dt}
=\frac{\hbar^2}{md^2}\Gamma(d).}
\tag{40.14}
\] 这是同一个线性位置噪声模型的双观测约束；不是所有
CSL、耗散或有色噪声模型的普适加热公式。这里没有使用第 \textbf{40.3} 节的
\([H,A]=0\)：自由粒子位置噪声可以持续加热，不能同时被描述成固定位置本征态的永久吸收。有限矩阵的精确全域
CCR 不成立，所以数值位置截断也不能未经误差分析直接继承 (40.13)。

若独立实验给附加加热上界 \(\dot E_{\max}\)，该具体模型必须满足
\(\gamma\le2m\dot E_{\max}/\hbar^2\)，并限制相应的额外干涉衰减。若一个记录触发模型拟合了更大衰减却没有匹配能量效应，它必须明确改变
\(A\)、噪声谱或耗散结构，随后接受新模型的检验；不能把同一噪声参数的另一个后果从合同中删去。

\hypertarget{ux8bb0ux5f55ux4f9dux8d56ux53efux4ee5ux5408ux6cd5ux52a0ux5165ux4f46ux987bux653eux5728ux7269ux7406ux7cfbux7edfux91cc}{%
\section{40.6
记录依赖可以合法加入，但须放在物理系统里}\label{ux8bb0ux5f55ux4f9dux8d56ux53efux4ee5ux5408ux6cd5ux52a0ux5165ux4f46ux987bux653eux5728ux7269ux7406ux7cfbux7edfux91cc}}

若记录 \(R\) 是真实的局部寄存器，有固定投影 \(Q_r\)，可以在联合空间定义
\[
L_r=\sqrt{\gamma_r}\,A\otimes Q_r,\qquad\gamma_r\ge0,
\] 并采用标准 Lindblad 方程。对已经经典条件化到记录值 \(r\)
的分支，系统获得相应速率；在全部记录结果上平均，整体演化仍是线性 CPTP。

这里 \(\gamma_r\) 是固定模型参数，\(Q_r\)
是可实施物理变量；若寄存器尚处于相干叠加，就应继续用联合模型计算，而不是先把记录强行当经典标签。它不等价于把分析者根据未观测全局态算出的熵、冗余或``对象程度''任意代入一个非线性单系统速率。

怎样从资源、自维持、实际可访问记录和目的任务中独立选出
\(A,Q_r,\gamma_r\)，仍是新增物理桥的责任。允许记录反馈的标准量子模型、基本随机场候选以及某个设备的校准控制是不同实现；它们需要共同响应和可区分的实验条件才能比较。

\hypertarget{ch-41}{%
\chapter{41. 为什么仍不能用相干损失单独证明客观坍缩}\label{ch-41}}

\hypertarget{ux7cfbux7edfux5185ux7edfux8ba1ux7684ux53efux8bc6ux522bux6027ux6781ux9650}{%
\section{41.1
系统内统计的可识别性极限}\label{ux7cfbux7edfux5185ux7edfux8ba1ux7684ux53efux8bc6ux522bux6027ux6781ux9650}}

\protect\hypertarget{c-05}{}{}

\textbf{C-05：不受限量子环境的等价障碍。} 每个有限 CPTP 通道
\(\mathcal E(\rho)=\sum_\mu L_\mu\rho L_\mu^\dagger\) 都可写为 \[
\mathcal E(\rho)=\operatorname{Tr}_E[V\rho V^\dagger],
\quad V=\sum_\mu L_\mu\otimes|\mu\rangle,\quad V^\dagger V=I.
\tag{41.1}
\]
因而仅观察系统的协议，不能把该通道唯一归因为基本坍缩或未观察的量子环境。

\textbf{证明。} 展开部分迹即得 Kraus
和；等距由保迹得到，并可在扩充空间中幺正完成。所有系统效应的 Born
配对相同。对有限多时干预，用顺序环境记忆构造因果 comb
的扩张，结论同样成立，但必须匹配整个过程而不只单时态。\(\square\)

该障碍不意味着坍缩模型不可实验检验。它意味着必须限制普通环境假设，提出跨系统、能量、时空或不同实验共用参数的额外关系，才能比较两套不同物理理论。有限所有
CPTP 解释的并集如果无限宽，任何单一通道都在里面。

\hypertarget{ux5b8cux6574ux9006ux6f14ux5316ux7a76ux7adfux6d4bux4ec0ux4e48}{%
\section{41.2
完整逆演化究竟测什么}\label{ux5b8cux6574ux9006ux6f14ux5316ux7a76ux7adfux6d4bux4ec0ux4e48}}

若全部标准环境、写入和动力学已独立表征，并能实施完整
\(U^\dagger\)，则标准模型预测返回初态。无法返回可否定这个完整``模型+控制''包；但也可能是遗漏环境或反演误差。对某新增噪声模型的证据，必须来自它在预注册对照中给出的独特尺度依赖或参数关系。

第39章通过的是更窄且明确的假设检验：拒绝所有对角经典随机相位。它没有跨越
C-05
的环境等价障碍。三分支循环测试、四分支随机酉测试、全环境反演和基本坍缩检验，应当是不同层级的实验标签。

\hypertarget{ch-42}{%
\chapter{42. 完整可执行研究协议与论证状态}\label{ch-42}}

\hypertarget{ux5b9eux9a8cux4e0eux62dfux5408ux5408ux540c}{%
\section{42.1
实验与拟合合同}\label{ux5b9eux9a8cux4e0eux62dfux5408ux5408ux540c}}

\begin{Verbatim}[breaklines,breakanywhere,fontsize=\scriptsize]
ExperimentContract:
  model_family:
    classical_diagonal_phase | approximately_pointer_preserving_random_unitary
    | Gaussian_diffusion | calibrated_quantum_environment | specified_collapse
  system: 4 levels, fixed pointer basis, fixed preparation
  channel: duration, control sequence, environment access
  measurement: E=W/(2+1/sqrt(2)), calibration bound
  assumptions:
    conditional per-shot null lower bound (chapter 46), or IID
    no post-hoc choice of witness or exclusion rules
    pointer survival and SPAM calibration data kept separate
  data:
    raw accept/reject counts
    time order, randomization, rejected-shot reasons
    calibration runs and their confidence allocation
  test:
    predeclared n and alpha
    lower null probability after systematic-error allowance
    lower-tail binomial test, no optional stopping without valid adjustment
  outputs:
    rejected family, unresolved larger families
    parameter intervals, calibration failures, reproducibility record
\end{Verbatim}

数据用于选择 \(W\) 后，若又用同一数据计算未经校正的
p-value，会产生选择偏差。可预先固定本章的
\(W\)，或用训练/留出数据。后选择损失必须计入仪器，不能删去``有事件''的试次。

粒子合同则冻结完整表示表、全局群、活跃质量阈值、标量数、beta
方案、\(Z'\) 混合和右手耦合。对 flavor 数据先匹配 Wilson 系数，再做
running
与矩阵元转换。本文没有使用未经提供的真实谱、介子或器件数据作``验证''。

\hypertarget{ux672cux7ec4ux524dux6cbfux63a8ux5bfcux7684ux72ecux7acbux68c0ux9a8cux4e0eux5f62ux5f0fux5316ux8fb9ux754c}{%
\section{42.2
本组前沿推导的独立检验与形式化边界}\label{ux672cux7ec4ux524dux6cbfux63a8ux5bfcux7684ux72ecux7acbux68c0ux9a8cux4e0eux5f62ux5f0fux5316ux8fb9ux754c}}

附录保留从头编写的数值/有理数验证程序、冻结输入和结果。主要比较相互独立的表达：

\begin{itemize}
\tightlist
\item
  从全左手表示表求六种异常，再与解析因式比较；
\item
  对有限整数域穷举，与正文的全整数分类证明对照；
\item
  由酉质量旋转直接构造矩阵，检验秩一和式、CKM 共轭关系和 FCNC 下界；
\item
  从复 Lindblad 向量构造生成元，再独立计算三角正性及 Poisson 饱和；
\item
  直接构造 Kraus、环境扩张、四维 \(W\) 及实际统计效应；
\item
  检查锐见证的饱和点、随机相位、混合与校准距离界；
\item
  通过同一联合纯态的两种远端测量复核分解依赖反例；
\item
  用精确 QND 过滤采样检查鞅均值、成对衰减和有限时吸收误差上界。
\end{itemize}

历史保存结果与本版重新运行分别登记，以程序
SHA、输入、运行环境和实际退出状态绑定；不能用旧 JSON
的成功标记替代本次执行。整数/有理数等式、有限维浮点残差和 Monte Carlo
覆盖也属于不同证据。

这些高级前沿的全称数学结论依靠正文的解析论证及所引用外部定理，抽样或有限搜索不替代它们。有限
Kraus 内核已经采用 Lean，并有独立 build、公理审计、实例正反控与远端 CI
回执；但不能把这项进展表述成 G-01、F-03、M-04、Itō 吸收、Berry/SSH
或历史依赖统计等全部高级命题都已形式化。具体有源码映射的选择器和谱行合同按第
31.6 节与形式化登记处理。

目前这些程序提供模型回归和算术复核，不是新硬件数据。实际设备必须执行预注册的物理和统计合同；模拟通过不构成经典噪声、基本坍缩、五色或新粒子模型的经验确认。

\hypertarget{ux5bf9ux65b0ux7269ux7406ux548cux65e7ux95eeux9898ux7684ux660eux786eux56deux7b54}{%
\section{42.3
对``新物理''和旧问题的明确回答}\label{ux5bf9ux65b0ux7269ux7406ux548cux65e7ux95eeux9898ux7684ux660eux786eux56deux7b54}}

\begin{longtable}[]{@{}
  >{\raggedright\arraybackslash}p{(\columnwidth - 4\tabcolsep) * \real{0.2500}}
  >{\raggedright\arraybackslash}p{(\columnwidth - 4\tabcolsep) * \real{0.3400}}
  >{\raggedright\arraybackslash}p{(\columnwidth - 4\tabcolsep) * \real{0.4100}}@{}}
\toprule\noalign{}
\begin{minipage}[b]{\linewidth}\raggedright
旧问题
\end{minipage} & \begin{minipage}[b]{\linewidth}\raggedright
本组前沿推导推进
\end{minipage} & \begin{minipage}[b]{\linewidth}\raggedright
尚未解决
\end{minipage} \\
\midrule\noalign{}
\endhead
\bottomrule\noalign{}
\endlastfoot
是否无条件推出三代三色 &
正确高能计数给出完整双候选分类；暴露低/高能条件不能混用 &
HGP、表示选择、排除五色的独立物理原理 \\
候选三代父结构是否有低能后果 & 秩一和式、上下味扇区不能同时对齐、总 FCNC
下界 & 重质量和耦合数值、现实信号及完整竞争模型 \\
相干循环能否随便设定 & 得到锐 CP 相位---衰减界和 Poisson 反例 &
非零循环不等于基本新物理 \\
二角/少能级能否识别所有经典噪声 &
证明单时刻三能级仍不够，构造四能级锐见证与有限样本协议 &
真实硬件执行、优于既有方案的实验比较 \\
Born 实际化只是假设最终吸收 & 对具体随机模型证明吸收及 Born
权重，给有限时间界 & 自然界是否采用这个基本随机机制 \\
额外退相干能否证明坍缩 & 给出无条件环境等价障碍及双观测能量约束 &
用独立校准/多实验共用参数识别具体基本模型 \\
\end{longtable}

本次有可计算的新研究产出，但没有足够证据称``发现了新的自然定律''。在
UEOT
当前基础上，最实在的进展是把不可检验的解释压缩成更少、更清楚的物理假设，然后得到它们不能逃避的定量后果。这也允许推导反过来否定我们喜欢的猜想。

\hypertarget{ux6765ux6e90ux4e0eux539fux521bux6027ux8fb9ux754c}{%
\section{42.4
来源与原创性边界}\label{ux6765ux6e90ux4e0eux539fux521bux6027ux8fb9ux754c}}

本组前沿推导继承原始研究定位、既有完整量子主文与综合 v2 总著，以及其中的
SRN、HGP、QFT 接口。本次维护还核对了可读的 v23--v30.11
历史文本、两份互补 v28.1 来源、正式 Core--Physics 合同及 v31
系列内核报告。材料身份和实际取得范围见附录
C；缓存可读文本、完整原件和已运行验证分别标识，不把其中一种冒充另一种，也不声称掌握未取得的全部历史对话。外部技术依赖保留其原始研究或作者材料的来源。

\begin{itemize}
\tightlist
\item
  331表示、跨代异常、FCNC：Pisano--Pleitez 1992；Liu
  1993；Frampton--Ho--Kephart 2012；Fonseca--Hirsch
  2016。链接见第36、37章。
\item
  非随机酉四维退相干及其先例矩阵：Buscemi--Chiribella--D'Ariano
  2005；几何研究也见
  \href{https://arxiv.org/abs/1508.04027}{Kayser--Luoma--Strunz}。本书未声称这些基础现象原创。
\item
  循环相位/Hamiltonian obstruction：Fagnola
  等，2018预印本、2019修订。网页抓取时出现的动态显示日期不视为论文新近发表日期。
\item
  无信号与线性 CP：Simon--Bužek--Gisin 2001。
\item
  连续测量和随机约化：Adler 等的鞅研究、Barchielli
  的随机方程工作。本文按自己声明的模型推导有限时间式；不是宣称随机约化机制首次提出。
\end{itemize}

对 F-03 的具体保守下界、M-04 的锐见证、E-01/E-02
的实验装配及其组合，本次给出了独立完整推导；当前检索并不构成``文献中从未出现''的证明。若准备投稿，应继续针对这些精确表达做先例核查、独立复核和实验性能对比，不能把``我这次推出来''直接等同于学术原创。

\hypertarget{ch-43}{%
\chapter{43. 全局相位、边界条件与拓扑量子化}\label{ch-43}}

\hypertarget{ux5c40ux90e8-madelung-ux8868ux793aux7f3aux5c11ux4ec0ux4e48}{%
\section{43.1 局部 Madelung
表示缺少什么}\label{ux5c40ux90e8-madelung-ux8868ux793aux7f3aux5c11ux4ec0ux4e48}}

第3章在 \(\rho>0\) 区域使用
\(\psi=\sqrt\rho e^{iS/\hbar}\)。不同局部片上的相位可以相差常数或规范转移函数；遇到节点时
\(S\) 不必是全局单值实函数。一个完整模型应给配置空间 \(X\)、Hermitian
线丛 \(L\to X\)、连接、Hilbert 空间 \(L^2(X,L)\)，以及具体自伴
Hamiltonian 的定义域。局部微分表达式相同的模型仍可因边界或 holonomy
不同而给出不同谱。

固定普通复标量波函数且闭轨 \(C\) 上无节点时，单值性意味着 \[
\oint_C dS=2\pi\hbar n,\qquad n\in\mathbb Z.
\]
它约束给定波函数的相位绕数，不是从整数条件构造所有位置和动量算符的证明。带电时可观测机械动量为
\(dS-qA\)；其积分不是无条件的 \(2\pi\hbar n\)，而含磁通项。

\hypertarget{ux73afux4e0aux5e26ux626dux8f6cux8fb9ux754cux7684ux5b8cux6574ux8c31}{%
\section{43.2
环上带扭转边界的完整谱}\label{ux73afux4e0aux5e26ux626dux8f6cux8fb9ux754cux7684ux5b8cux6574ux8c31}}

取 \(X=[0,L]\) 两端粘合、质量 \(m>0\)、常数切向势 \(A\)。在
\(L^2([0,L])\) 上定义 \[
P_\theta=-i\hbar\partial_x,
\quad D(P_\theta)=\{f\in H^1:f(L)=e^{i\theta}f(0)\}.
\] 边界型为
\(-i\hbar[\bar f g]_0^L\)，上述域使其为零；伴随域由对所有满足边界的
\(f\) 的边界型消失得到同一关系。因此这是自伴域，而非仅形式对称。

\protect\hypertarget{p-ugl-01}{}{}

\textbf{P-UGL-01：环上谱和规范不变量。} \[
f_n(x)=L^{-1/2}e^{i(2\pi n+\theta)x/L},\quad
E_n=\frac{[\hbar(2\pi n+\theta)/L-qA]^2}{2m}.
\] \(H=(P_\theta-qA)^2/(2m)\) 的域是 \(H^2\)
中函数与一阶导数均满足相同扭转边界的元素。基底由乘上 \(e^{i\theta x/L}\)
后的普通 Fourier 基得到，故完备。谱只依赖 \[
\eta=\theta-\frac{qAL}{\hbar}\pmod {2\pi}.
\] 采用 \(A'=A+\partial_x\chi\)、\(f'=e^{iq\chi/\hbar}f\)，有
\(\theta'=\theta+q[\chi(L)-\chi(0)]/\hbar\)；于是 \(\eta\)
不变。\(\theta\mapsto\theta+2\pi\) 只重标整数
\(n\)。若从一开始固定周期域，不能只改变边界标签而忘记同步运输 Hilbert
域。\(\square\)

例 \(\hbar=m=L=q=1\)：\(\theta=0\)、\(A=0\) 时基态能量0；\(A=\pi\)
时最低能量 \(\pi^2/2\)，有两重简并。若同时把 \(\theta\) 改为
\(\pi\)，则回到 \(\eta=0\)
的谱；这是规范运输的例子，不是磁通不产生物理效果。

\hypertarget{ux5bf9ux539fux59cbux76f8ux4f4dux95edux5408ux8defux7ebfux7684ux53cdux5411ux4f18ux5316}{%
\section{43.3
对原始相位闭合路线的反向优化}\label{ux5bf9ux539fux59cbux76f8ux4f4dux95edux5408ux8defux7ebfux7684ux53cdux5411ux4f18ux5316}}

应把``相位闭合''拆为连接的局部曲率、全局
holonomy、波函数丛的转移函数、算符域四项。原始闭轨直觉在这里有可计算落点，但量子作用尺度、复数态空间和边界的物理选择仍是输入。Core
可以接收规范不变谱和允许干涉协议形成响应对象；不能从协议命名重新推出
\(\hbar\) 的数值。

\hypertarget{ch-44}{%
\chapter{44. 拓扑量子对象：Berry 丛与 SSH 边态}\label{ch-44}}

\hypertarget{ux4e8cux80fdux7ea7-berry-ux66f2ux7387ux7684ux5b8cux6574ux8ba1ux7b97}{%
\section{44.1 二能级 Berry
曲率的完整计算}\label{ux4e8cux80fdux7ea7-berry-ux66f2ux7387ux7684ux5b8cux6574ux8ba1ux7b97}}

取
\(H(\mathbf n)=\Delta\mathbf n\cdot\boldsymbol\sigma\)，\(\Delta>0\)，\(\mathbf n\in S^2\)，讨论能量
\(+\Delta\) 的瞬时本征线。球坐标下在北片可取 \[
|u_N\rangle=(\cos(\vartheta/2),e^{i\varphi}\sin(\vartheta/2))^T.
\] 定义连接 \(\mathcal A=i\langle u|du\rangle\)。南片取
\(u_S=e^{-i\varphi}u_N\)，于是 \[
\mathcal A_N=-\frac{1-\cos\vartheta}{2}d\varphi,
\quad \mathcal A_S=\frac{1+\cos\vartheta}{2}d\varphi.
\]

\protect\hypertarget{p-utp-01}{}{}

\textbf{P-UTP-01：此本征线的 Chern 数。} 直接求导给 \[
\mathcal F=-\tfrac12\sin\vartheta\,d\vartheta\wedge d\varphi,
\qquad C_1=\frac1{2\pi}\int_{S^2}\mathcal F=-1.
\] 两个局部连接之差 \(\mathcal A_S-\mathcal A_N=d\varphi\)
与转移函数一致；总通量非零说明不能在整球面选一个处处光滑的本征向量。负能带的符号相反。这里的符号取决于明列的连接、能带和球面定向。\(\square\)

绝热演化的几何相位为
\(\oint\mathcal A\)，还须从总相位中分离动力学相位；有限速度误差须随间隙、路径光滑性与速度另行控制。本例使用
Berry
的既有几何相位结构，本书的计算不是其首次发现。\href{https://doi.org/10.1098/rspa.1984.0023}{Berry
原始论文}

\hypertarget{ssh-ux4f53ux8c31ux548cux7ed5ux6570}{%
\section{44.2 SSH 体谱和绕数}\label{ssh-ux4f53ux8c31ux548cux7ed5ux6570}}

固定晶胞、边界切法以及实正跃迁 \(t_1,t_2\)。单粒子 Hamiltonian 为 \[
H=\sum_{n\ge1}(t_1|n,A\rangle\langle n,B|
+t_2|n+1,A\rangle\langle n,B|+\mathrm{h.c.}).
\] 无限平移系统的 Bloch 矩阵，按约定的 Fourier 变换为 \[
H(k)=\begin{pmatrix}0&t_1+t_2e^{-ik}\\t_1+t_2e^{ik}&0\end{pmatrix},
\quad E_\pm(k)=\pm\sqrt{t_1^2+t_2^2+2t_1t_2\cos k}.
\] 手征算符 \(\Gamma=\operatorname{diag}(1,-1)\) 满足
\(\Gamma H\Gamma=-H\)。取 \(q(k)=t_1+t_2e^{-ik}\)，绕数 \[
\nu=\frac{1}{2\pi i}\int_{-\pi}^{\pi}q^{-1}\partial_kq\,dk
\] 在 \(t_2>t_1\) 时为 \(-1\)，\(t_2<t_1\) 时为0。证明是 \(q(k)\) 描出以
\(t_1\) 为中心、半径 \(t_2\) 的顺时针圆；在 \(t_1=t_2\)
能隙闭合，绕数不定义。若交换 Fourier
或子格约定，绕数符号也随之改变。晶胞重选必须同步改变边界，不能仅变一个数字后宣称物理边态消失。

\hypertarget{ux534aux65e0ux9650ux8fb9ux754cux4e0eux6709ux9650ux94feux8befux5dee}{%
\section{44.3
半无限边界与有限链误差}\label{ux534aux65e0ux9650ux8fb9ux754cux4e0eux6709ux9650ux94feux8befux5dee}}

\protect\hypertarget{p-utp-02}{}{}

\textbf{P-UTP-02：半无限零模。} 对左端终止于 \(A_1\) 的链，\(t_1<t_2\)
时 \[
|\psi_L\rangle=\sqrt{1-r^2}\sum_{n=1}^\infty(-r)^{n-1}|n,A\rangle,
\qquad r=t_1/t_2
\] 是归一化零能本征态。证明：所有 \(B_n\) 分量要求
\(t_1a_n+t_2a_{n+1}=0\)；几何级数当且仅当 \(r<1\)
收敛，归一因子如上。所有 \(A_n\) 方程因 \(b_n=0\) 自动满足。\(\square\)

有限 \(N\) 晶胞上将此向量截断并归一为 \(\psi_{L,N}\)，系数
\(c_N=\sqrt{(1-r^2)/(1-r^{2N})}\)。只有末端 \(B_N\) 有残差，故 \[
\|H_N\psi_{L,N}\|=c_Nt_1r^{N-1}.
\] 谱定理给
\(\operatorname{dist}(0,\operatorname{spec}H_N)\le\|H_N\psi_{L,N}\|\)。不能把有限等子格链的近零模写成一般精确零模；两端耦合会产生小能量劈裂。对子格不平衡或特殊边界须另论。已知
SSH 模型的来源见
\href{https://link.aps.org/doi/10.1103/PhysRevLett.42.1698}{Su--Schrieffer--Heeger
原始论文}。

\hypertarget{ueot-ux63a5ux53e3ux4e0eux53efux505aux7684ux5b9eux9a8c}{%
\section{44.4 UEOT
接口与可做的实验}\label{ueot-ux63a5ux53e3ux4e0eux53efux505aux7684ux5b9eux9a8c}}

以边界子空间作为候选载体，协议包含局部注入、有限时传播、边界存活和扰动响应。证书应同时记录体能隙、手征对称保留程度、边界切法、有限尺寸和寿命；仅有一个局域峰不构成拓扑证书。

可执行路径：先对小链精确对角化核对残差；再改变 \(t_1/t_2\)
观察体隙闭合与边界分布；最后加入已知手征保持或破坏的扰动，比较预注册的谱与存活曲线。由此原论文``持续的结构''成为具体的受保护低能子空间问题。它不直接给出电子世代数、基本粒子或绝对永存对象。

\hypertarget{ch-45}{%
\chapter{45. 非微扰起点：有限格点规范模型与跨尺度接口}\label{ch-45}}

\hypertarget{ux6709ux9650ux7eafux89c4ux8303ux6a21ux578bux7684ux7cbeux786eux5b9aux4e49}{%
\section{45.1
有限纯规范模型的精确定义}\label{ux6709ux9650ux7eafux89c4ux8303ux6a21ux578bux7684ux7cbeux786eux5b9aux4e49}}

取有限四维 Euclidean 超立方格点，定向边上
\(U_{x\mu}\in SU(N)\)，反向边取 \(U^\dagger\)。定向 plaquette \[
U_p=U_{x\mu}U_{x+\hat\mu,\nu}U^\dagger_{x+\hat\nu,\mu}U^\dagger_{x\nu}.
\] 定义 Wilson 作用和归一化测度 \[
\begin{gathered}
S_W=\beta\sum_p\left(1-\frac1N\operatorname{ReTr}U_p\right),
\\ Z=\int\prod_l dU_l\,e^{-S_W},\\ d\mu=Z^{-1}e^{-S_W}\prod_l dU_l,
\\\beta=\frac{2N}{g^2}>0.
\end{gathered}
\] \(dU_l\) 是归一化 Haar 测度，每个无向 plaquette
取一种定向。本节无费米子行列式、化学势或实时间相位因子。

\protect\hypertarget{p-ulg-01}{}{}

\textbf{P-ULG-01：有限模型存在且规范不变。} 有限紧致乘积群上
\(0\le S_W\le2\beta N_p\)，故
\(e^{-2\beta N_p}\le Z\le1\)，测度定义良好。局部变换
\(U_{xy}\mapsto G_xU_{xy}G_y^\dagger\)
使闭环积变为起点共轭；迹不变，Haar 测度左右不变，故所有 Wilson
环期望规范不变。\(\square\)

此存在性不要求选取全局规范。它仅建立有限 Euclidean
统计模型；要重建正能实时量子理论，还需反射正性、连续性和适当极限等条件，不能仅凭
\(e^{-S}>0\) 宣布已证明全部条件。Wilson 原始格点构造及强耦合研究见
\href{https://link.aps.org/doi/10.1103/PhysRevD.10.2445}{Confinement of
quarks}。

\hypertarget{ux8fdeux7eedux89c4ux8303ux4f5cux7528ux7684ux7cfbux6570ux6838ux5bf9}{%
\section{45.2
连续规范作用的系数核对}\label{ux8fdeux7eedux89c4ux8303ux4f5cux7528ux7684ux7cfbux6570ux6838ux5bf9}}

采用第29章约定 \(D_\mu=\partial_\mu+igA_\mu\)，Hermitian 生成元
\(\operatorname{Tr}T^aT^b=\delta^{ab}/2\)，因此 \[
F_{\mu\nu}=\partial_\mu A_\nu-\partial_\nu A_\mu+ig[A_\mu,A_\nu].
\] 对光滑场把边变量视为与上述输运约定相容的路径有序指数，局部展开 \[
U_p=\exp\{iga^2F_{\mu\nu}+O(a^3)\},\quad
1-\frac1N\operatorname{ReTr}U_p
=\frac{g^2a^4}{2N}\operatorname{Tr}F_{\mu\nu}^2+O(a^5).
\] 线性项因 \(F\)
无迹而消失，二次项由指数展开给出。固定体积和光滑场、边界受控时，乘
\(\beta\) 并求和得到 \[
S_W\longrightarrow\int d^4x\sum_{\mu<\nu}\operatorname{Tr}F_{\mu\nu}^2
=\frac14\int d^4x\sum_{\mu,\nu,a}(F_{\mu\nu}^a)^2.
\]
本式证明经典光滑场作用的一致性；量子典型场并不光滑，所以它不证明量子测度的连续极限。经过对称求和可以研究更高阶格距改进，但需声明边界和光滑性，不能把展开阶数直接当作强耦合量子误差界。

\hypertarget{ux4e00ux4e2aux53efux5b8cux5168ux7a77ux4e3eux7684ux89c4ux8303ux95edux5408ux4f8b}{%
\section{45.3
一个可完全穷举的规范闭合例}\label{ux4e00ux4e2aux53efux5b8cux5168ux7a77ux4e3eux7684ux89c4ux8303ux95edux5408ux4f8b}}

为了不把数值平台的缺失藏在抽象符号里，取单方形四条边的 \(\mathbb Z_2\)
子模型，\(u_l=\pm1\)，plaquette \(P=\prod_{l=1}^4u_l\)，作用
\(S=-\beta P\)。用未归一计数测度，有16个配置，其中各8个 \(P=\pm1\)。

\protect\hypertarget{p-ulg-02}{}{}

\textbf{P-ULG-02：精确配分函数与响应。} \[
Z=16\cosh\beta,\qquad \langle P\rangle=\tanh\beta,
\qquad\partial_\beta\langle P\rangle=1-\langle P\rangle^2.
\] 证明由两个8重扇区求和；顶点变换使相邻两边各乘同一 \(g_x=\pm1\)，故
\(P\)
不变。所有单边期望为0，由翻转该边端点并利用测度不变得到。若改为归一 Haar
计数，\(Z\) 除16，归一期望不变。\(\square\)

这给出一种具体 UEOT 读出：允许全部 \(f(P)\) 观测时 \(P\)
是完备不变信息载体。证明其轨道完备性可依次规范固定前三条边为
\(+1\)，最后一条恰为 \(P\)；因此同 \(P\)
的配置属于同一规范轨道。此例只针对该有限群与单环，不把它冒充 \(SU(3)\)
QCD 或四维禁闭证明。

\hypertarget{ux683cux70b9ux7c97ux5316ux4e0e-core-ux7684ux5408ux6cd5ux8fdeux63a5}{%
\section{45.4 格点粗化与 Core
的合法连接}\label{ux683cux70b9ux7c97ux5316ux4e0e-core-ux7684ux5408ux6cd5ux8fdeux63a5}}

定义粗边为一条细路径的有序边积，得到映射
\(B:U_f\mapsto U_c\)。有限模型的粗测度可直接定义为推前
\(\mu_c=B_*\mu_f\)。对每个有界粗观测量， \[
\mathbb E_{\mu_c}f(U_c)=\mathbb E_{\mu_f}f(B(U_f)).
\]
这是推前测度的定义恒等式，具有精确读出一致性。它不意味着粗作用仍只含单
plaquette；多环项与非局域项一般会出现。若截断为有限算符族，应输出被舍弃响应的误差，而不是宣布``同形闭合''。

Euclidean 配置的确定映射是经典统计粗化；它不是自动的实时 CPTP
通道。要接第26章的量子过程接口，需给转移矩阵或受控实时模型、初态与同一干预插槽，再比较响应。由此可分别研究：精确读出推前、近似有效作用和量子干预粗化三种不同的问题。

\hypertarget{ux63a8ux8fdbux5230-qcd-ux7684ux660eux786eux4efbux52a1}{%
\section{45.5 推进到 QCD
的明确任务}\label{ux63a8ux8fdbux5230-qcd-ux7684ux660eux786eux4efbux52a1}}

先在有限 \(SU(2)\)/\(SU(3)\)
格点验证局部更新、规范不变性、热化与自相关；再加入具体费米子离散化和海夸克权重，核对加倍、手征性与反常的处理；最后对同一无量纲观测作体积、格距、质量与统计误差的联合外推。实时响应不能从有噪
Euclidean
数据无条件稳定反演。有限纯规范测度的存在性和局部一圈渐近自由都不足以证明四维连续
Yang--Mills 质量隙。

\hypertarget{ch-46}{%
\chapter{46. 有历史依赖的实验统计与可选停止}\label{ch-46}}

\hypertarget{ux4eceux72ecux7acbux8bd5ux6b21ux6539ux4e3aux6761ux4ef6ux96f6ux5047ux8bbe}{%
\section{46.1
从独立试次改为条件零假设}\label{ux4eceux72ecux7acbux8bd5ux6b21ux6539ux4e3aux6761ux4ef6ux96f6ux5047ux8bbe}}

第39章的四能级见证，固定输入和效应后，经典对角随机相位模型给事件概率下界
\[
p_* =\frac{3-2\sqrt2}{4(2+1/\sqrt2)}.
\] 设第 \(t\) 次事件 \(X_t\in\{0,1\}\)，\(\mathcal F_{t-1}\)
包含此前全部记录和允许反馈。新的零假设明确为 \[
\Pr(X_t=1\mid\mathcal F_{t-1})\ge p_0\quad\text{几乎处处，所有 }t.
\] 其中 \(0<p_0<1\)，且 \(p_0\)
已扣除校准、泄漏和可识别的系统误差。每个条件通道都属于被检验类时，该条件成立；只知道时间平均噪声属于该类并不足够。

\protect\hypertarget{p-ust-01}{}{}

\textbf{P-UST-01：相关试次的固定样本检验。} 在上述条件下，对固定 \(n\)，
\[
\Pr\left(\sum_{t=1}^nX_t\le k\right)
\le\sum_{j=0}^k{n\choose j}p_0^j(1-p_0)^{n-j}.
\] 证明可用顺序耦合：在扩充概率空间取独立
\(U_t\sim\mathrm{Unif}[0,1]\)，以相应历史条件概率 \(p_t\) 实现
\(X_t=1\{U_t\le p_t\}\)。同时令 \(Y_t=1\{U_t\le p_0\}\)，则路径上
\(X_t\ge Y_t\)，而 \(Y_t\) 为独立
Bernoulli，故下尾受其控制。对额外历史变量使用其条件核依次生成，二值试次的边缘联合律不变。\(\square\)

零事件特例亦可直接迭代条件期望得 \((1-p_0)^n\)。因此理想 \(p_0=p_*\)
时188次零事件仍可给固定样本5\%检验，即使存在满足此条件下界的历史依赖。这个结论没有把任意漂移或记忆都排除掉。

\hypertarget{ux4e3aux4ec0ux4e48ux5355ux6b21ux8fb9ux7f18ux4e0bux754cux4e0dux591f}{%
\section{46.2
为什么单次边缘下界不够}\label{ux4e3aux4ec0ux4e48ux5355ux6b21ux8fb9ux7f18ux4e0bux754cux4e0dux591f}}

令所有试次都等于同一个隐藏 Bernoulli 变量
\(Z\)，\(\Pr(Z=1)=p_0\)。每次边缘概率都为
\(p_0\)，但任意多次全零的概率始终为 \(1-p_0\)，不会按 \((1-p_0)^n\)
衰减。首次看到零后，下一次事件条件概率为0，正好违反46.1的条件。这个反例必须随统计证书保存。

\hypertarget{ux5141ux8bb8ux8fdeux7eedux67e5ux770bux6570ux636eux7684ux4e25ux683cux7248ux672c}{%
\section{46.3
允许连续查看数据的严格版本}\label{ux5141ux8bb8ux8fdeux7eedux67e5ux770bux6570ux636eux7684ux4e25ux683cux7248ux672c}}

对预先固定 \(q\in[0,p_0)\)，定义 \[
L_n=\prod_{t=1}^n
\left(\frac q{p_0}\right)^{X_t}
\left(\frac{1-q}{1-p_0}\right)^{1-X_t},\quad L_0=1.
\]

\protect\hypertarget{p-ust-02}{}{}

\textbf{P-UST-02：可选停止证书。} \(L_n\)
是零假设下非负超鞅。因为给定历史后单步因子的期望为 \[
1+\frac{(p_t-p_0)(q-p_0)}{p_0(1-p_0)}\le1.
\] 令 \(\tau=\inf\{n:L_n\ge1/\alpha\}\)。对截断停时 \(\tau\wedge N\)
的非负超鞅使用迭代条件期望，得
\(1\ge E L_{\tau\wedge N}\ge\alpha^{-1}\Pr(\tau\le N)\)；令
\(N\to\infty\) 得
\(\Pr(\tau<\infty)\le\alpha\)。这给出任意时刻查看并按阈值停止的有效检验。\(\square\)

\(q=0\) 时约定无事件因子为
\(1/(1-p_0)\)，一旦有事件整个乘积为0；它最适合专门的零事件目标。为允许少量事件可固定
\(q>0\)，或用事先固定权重的多个 \(q\) 的混合。不能事后取最大 \(L_n(q)\)
而不校正。

\hypertarget{ux6821ux51c6ux6cc4ux6f0fux4e0eux5b9eux9645ux5408ux540c}{%
\section{46.4
校准、泄漏与实际合同}\label{ux6821ux51c6ux6cc4ux6f0fux4e0eux5b9eux9645ux5408ux540c}}

若以概率至少 \(1-\alpha_{\rm cal}\)
的同时校准事件保证所有使用试次的下界成立，则总错误概率至多
\(\alpha_{\rm cal}+\alpha_{\rm test}\)。应证明在该好事件与相应信息条件下所用零假设有效，或采用独立校准冻结的物理一致置信集合；事后筛选的校准事件不能任意当作与测试独立。对整个置信集合取最坏下界更直接。

实际执行应冻结
\(W\)、准备、指针基、随机化、失效率和停止规则，保留时间序列。第39章的近似指针保持界
\(p_0\ge p_*-\sqrt{8\ell}-\delta_{\rm SPAM}\)
很保守；若右端非正，不能声称当前试验仍有该见证的排除能力。新的改进方向是从完整校准置信集合做有证界的最优化，而不是随意减小误差条。

本章在明确条件下取消了独立同分布要求，并补足允许可选停止的检验。仍未执行真实硬件实验，也未证明自然界存在新的退相干机制。

\hypertarget{ch-47}{%
\chapter{47. 研究总规划：从闭合模块到可区分物理}\label{ch-47}}

\hypertarget{ux5de5ux4f5cux4f9dux8d56ux4e0eux4f18ux5148ux7ea7}{%
\section{47.1
工作依赖与优先级}\label{ux5de5ux4f5cux4f9dux8d56ux4e0eux4f18ux5148ux7ea7}}

研究计划按可检验依赖推进，而不是承诺在某个日期解决开放难题。P0
表示任何新物理声明之前必须完成的底座；P1
表示当前已有定理或可计算模型支撑、最可能产生明确产出的方向；P2
表示需要新物理选择原则或非平凡数学突破的长期方向。

\begin{longtable}[]{@{}
  >{\raggedright\arraybackslash}p{(\columnwidth - 6\tabcolsep) * \real{0.2100}}
  >{\raggedright\arraybackslash}p{(\columnwidth - 6\tabcolsep) * \real{0.2000}}
  >{\raggedright\arraybackslash}p{(\columnwidth - 6\tabcolsep) * \real{0.2800}}
  >{\raggedright\arraybackslash}p{(\columnwidth - 6\tabcolsep) * \real{0.3100}}@{}}
\toprule\noalign{}
\begin{minipage}[b]{\linewidth}\raggedright
工作包
\end{minipage} & \begin{minipage}[b]{\linewidth}\raggedright
优先级
\end{minipage} & \begin{minipage}[b]{\linewidth}\raggedright
必须先有
\end{minipage} & \begin{minipage}[b]{\linewidth}\raggedright
可以交付的结果
\end{minipage} \\
\midrule\noalign{}
\endhead
\bottomrule\noalign{}
\endlastfoot
R1 量子选择与最小前提 & P0 & OPT 背景、经典/实量子反模型 &
依赖图、逐假设删除证书、可审查重建链 \\
R2 全局相位与连续域 & P0 & 自伴域、线丛与能量界 &
含边界的模型清单、收敛证书 \\
R3 记录见证与 QCVV & P1 & SRN、见证、校准和统计合同 &
有效置信界、预注册实验与原始数据 \\
R4 实际化与新动力学 & P1/P2 & CP、无信号、吸收/加热约束 &
能被排除的具体模型、共参数多观测预测 \\
R5 记忆过程与对象压缩 & P1 & comb、干预充分性、资源约束 &
预测压缩误差、记忆维数上下界 \\
R6 多体拓扑与控制 & P1 & 有限模型、谱隙、稳定性和噪声 &
可复现边态寿命、逻辑误差和控制收益 \\
R7 QFT/EFT 精确接口 & P0/P1 & 作用量、调节、方案、红外定义 &
带阶数和误差的散射与匹配基准 \\
R8 格点与非微扰 & P1/P2 & 有限规范测度、更新算法 &
有限体积观测、系统外推；连续极限为独立目标 \\
R9 规范与世代选择 & P2 & 全场表、阈值、全局群与反模型 &
排除五色的独立机制或承认多候选 \\
R10 味、质量与中微子 & P1/P2 & 破缺谱、秩一结构、EFT 匹配 &
留出和式、过程相关预测与竞争模型比较 \\
R11 全局异常与几何 & P2 & 具体丛、算符、completion &
可计算异常映射，保持或撤回 framing 条件结论 \\
R12 引力与宇宙学接口 & P2 & 量子场、局域性、能量及几何输入 &
受控背景模型和一致性检查；不预报已完成引力统一 \\
R13 证明与软件复现 & P0 & 本书状态编号和代码 &
类型检查、反例集、形式化小核心和版本追踪 \\
\end{longtable}

第一阶段冻结 R1/R2/R7/R13
的定义与基准；第二阶段优先并行开展不依赖新基本公理的 R3/R5/R6/R8
有限模型；第三阶段才让 R4/R9/R10/R11 的候选与实验、数学反例竞争。R12
先作为接口研究，不能借``统一''的目标绕过前面未闭合的步骤。这里的并行表示研究任务之间的依赖关系；每项成果仍以具体证明、源文件和独立审查为准。

\hypertarget{r1ux91cfux5b50ux9009ux62e9ux4e0eux6700ux5c0fux6027}{%
\section{47.2
R1：量子选择与最小性}\label{r1ux91cfux5b50ux9009ux62e9ux4e0eux6700ux5c0fux6027}}

\textbf{目标。} 把``UEOT 可组织量子理论''和``UEOT
必须选择量子理论''分清，并尝试减少真正独立的量子领域输入。

\textbf{数学入口。} 用完整因果 OPT
的对象、测试、组合、等价、凸闭包和有限维条件构成背景包
\(\mathcal B\)；选择假设集合 \(\mathcal H\)，结论写成
\(\mathcal B\cup\mathcal H\Rightarrow\mathrm{CQM}_{\rm fd}\)。任何更换背景都重新计入最小性成本。第2章的外部重建定理不是通用
Core 定理。

\textbf{步骤。}
首先逐条对照所引用重建原文，建立每个引理实际使用的背景表；其次在经典、实量子、受限效应/变换的量子子理论及其他具体
OPT
中做删除测试；再检查``态空间像量子''是否同时保证全部量子过程与复合规则；最后寻找
UEOT
原始资源/访问条件能否独立蕴含某个领域原则。每次声称蕴含都必须同时尝试构造经典反模型。

\textbf{验收。}
一条没有循环定义、能追到基础量词的证明链；或明确反例证明所尝试的更弱包仍不足够。反例也是有效交付。仅改名纯化为``父对象完备''不构成进展；把目标量子性藏入``可实现''定义也不合格。若原预印本的版本修订改变背景，必须更新
K 依赖，不能静默保留更强结论。

\hypertarget{r2ux5168ux5c40ux76f8ux4f4dux7b97ux7b26ux57dfux4e0eux6570ux503cux6536ux655b}{%
\section{47.3
R2：全局相位、算符域与数值收敛}\label{r2ux5168ux5c40ux76f8ux4f4dux7b97ux7b26ux57dfux4e0eux6570ux503cux6536ux655b}}

\textbf{目标。}
使第3章的局部统计作用量可安全进入真实边界、节点和多连通配置空间。

\textbf{步骤。}
从第43章环模型和二维穿孔平面开始，先固定线丛及连接，再列出 Hamiltonian
的自伴扩张；用环谱、磁通周期和干涉相位做基准。遇到
\(\rho=0\)，使用弱解/波函数域表述，禁止除以零密度。对数值截断同时控制初态尾、生成元近似和时间积分；不能仅凭
\(D(\rho,\rho_{\rm cut})\) 小就宣布长期动力学误差小。

若在共同不变有限域上有 \(\|H-\widetilde H\|\le\delta_H\)，Duhamel 恒等式
\[
U(t)-\widetilde U(t)
=-\frac{i}{\hbar}\int_0^tU(t-s)(H-\widetilde H)\widetilde U(s)\,ds
\] 给
\(\|U-\widetilde U\|\le |t|\delta_H/\hbar\)。无界问题应对所用态的积分范数逐项估计，不能继续使用不存在的全局算符范数。这是动力学误差合同的最小原型。

\textbf{验收。}
边界变换前后谱/响应相同；网格或能量截断的误差在指定时间窗收敛；全局节点处理不再依赖奇异的
Madelung
除法。若不同自伴扩张给不同预测而没有独立物理选择，应保留多个模型，而不是选最符合目标的谱。

\hypertarget{r3srn-ux89c1ux8bc1qdrr-ux4e0eux5b9eux9a8cux7edfux8ba1}{%
\section{47.4 R3：SRN 见证、QDRR
与实验统计}\label{r3srn-ux89c1ux8bc1qdrr-ux4e0eux5b9eux9a8cux7edfux8ba1}}

\textbf{目标。}
把``记录导致什么相干损失''变成能排除明确模型类的实验，而不直接跳到本体坍缩。

\textbf{候选顺序。} 先检验全 SRN；再检验时间齐次半群；再用三分支循环区分
Gaussian
对角扩散与更宽经典跳跃；最后用四分支见证排除整个经典对角随机相位类。每一级都有更宽的未排除类。第39章
\(K_*\) 有零元，因此其精确通道不是有限速率齐次 SRN
半群的有限时刻，实验实现应使用受控环境通道，不把它错误拟合为该半群。

\textbf{步骤。} 以已知四级系统或两比特编码构造
\(|+_4\rangle\)；用环境比特实现指定 Kraus 扩张；测量
\(E=W/(2+1/\sqrt2)\)；独立校准指针存活和 SPAM，冻结
\(p_0\)；按第46章固定样本或超鞅方案采集时间序列；留出不同存储时间和环境控制做模型检查。泄漏事件必须作为输出或概率损失计入，不能删除后提高信号。

\textbf{优化数学。} 令 \(\mathcal C(D_{\rm cal})\)
为有覆盖保证的通道/仪器集合，求 \[
p_{\min}=\inf_{M\in\mathcal C(D_{\rm cal})\cap\mathfrak M_0}p_M(E).
\] 随机相位凸包不自动有精确小规模 SDP
表示；可用有证外逼近得到保守下界。下界大于实测留出上界时才能拒绝该类，优化器的一个局部最小值不足为证。

\textbf{验收与停止。} 在模拟的零假设、漂移和近似 SRN
数据上校准错误率；真实实验同时报告原始计数和最坏误差。若
\(p_{\min}\le0\)
或准备/效应不可独立校准，先改善实验合同，不发布``排除经典噪声''的结论。实际样本量根据校准后的
\(p_0\) 重新计算；188不是通用硬件承诺。

\hypertarget{r4ux5b9eux9645ux5316ux4e0eux65b0ux52a8ux529bux5b66}{%
\section{47.5
R4：实际化与新动力学}\label{r4ux5b9eux9645ux5316ux4e0eux65b0ux52a8ux529bux5b66}}

\textbf{目标。}
从第40章的标准随机候选进一步得到与已校准环境不同的、跨观测共用参数的物理预测。

\textbf{模型先冻结。} 必须给随机概率空间、观测算符 \(A\)、速率
\(\gamma\)、时间与空间关联、系统组合规则、能量预算及记录如何参与动力学。有限维
QND
模型是已闭合起点；它既可解释为连续测量轨迹，也可作为基本随机假设。两者的无条件通道相同不产生实验区分。

\textbf{步骤。} 先证明范数、平均 CP/TP
与局域无信号；再求吸收或稳态，非对易 \([H,A]\ne0\) 时重新分析，不能沿用
QND
定理；随后推导加热、动量扩散、干涉对比度与频谱的共同关系；最后把同一个参数域带入多类实验，限制普通环境模型并留出一组观测。若依赖质量密度场，须给空间平滑尺度以控制
UV，并核对复合体放大规律与粒子置换对称。

\textbf{关键对照。} 第40章位置噪声给
\(\dot E=\gamma\hbar^2/(2m)\)、\(\Gamma(d)=\gamma d^2/2\)，从而
\(\dot E=\hbar^2\Gamma(d)/(md^2)\)。如果某具体候选同时声称该退相干律又没有该能量代价，必须指出模型新增了何种耗散或反馈；不能只保留有利的半条预测。

\textbf{验收。}
同一冻结模型的可重复留出预测，包含允许环境的独立上界。若候选始终可被任意环境
dilation
完全吸收且没有跨观测约束，则只有解释差异，不足以称新物理。相对论协变扩展还须满足局域一致性和能量控制，本书未完成该扩展。

\hypertarget{r5ux975e-markov-ux8bb0ux5fc6ux4e0eux9884ux6d4bux5bf9ux8c61ux538bux7f29}{%
\section{47.6 R5：非 Markov
记忆与预测对象压缩}\label{r5ux975e-markov-ux8bb0ux5fc6ux4e0eux9884ux6d4bux5bf9ux8c61ux538bux7f29}}

\textbf{目标。} 用 UEOT
的未来充分性指导何时一个密度矩阵够用，何时必须保留环境记忆。

\textbf{步骤。} 构建一个有限系统与可控记忆环境，枚举一组层析完备干预；以
causal break 制造相同当前态但不同未来响应；再拟合有限记忆 comb，用独立
tester 检查预测误差。以允许协议族 \(\mathcal U\) 定义
\(d_{\mathcal U}(R,\widetilde R)=\sup_{T\in\mathcal U}\operatorname{TV}(p_R^T,p_{\widetilde R}^T)\)，压缩目标是在给定资源约束内最小化此误差，而非仅最小化单时态重构误差。

\textbf{可推进的数学。}
建立近似充分因子的误差组合界；给不同允许访问集合的记忆维数上下界；用可区分历史数或可证矩阵秩产生压缩下界。过程
Choi
的某个矩阵秩不无条件等于最小物理环境维数，需证明表示与实现之间的对应。

\textbf{验收。} 压缩模型的 Choi
正性与因果递归都通过；未用于训练的自适应协议误差有证界；同样资源下优于仅
\(\rho_t\)
的模型。若只在无干预轨迹上拟合好，不能宣称已取得完整未来充分性。

\hypertarget{r6ux591aux4f53ux62d3ux6251ux70edux529bux5b66ux4e0eux63a7ux5236ux5e94ux7528}{%
\section{47.7
R6：多体、拓扑、热力学与控制应用}\label{r6ux591aux4f53ux62d3ux6251ux70edux529bux5b66ux4e0eux63a7ux5236ux5e94ux7528}}

\textbf{目标。} 从已知具体 Hamiltonian
出发，研究对象持续、编码和控制是否产生可比较的性能收益。

\textbf{步骤。} 以 SSH、TFIM、Hubbard
和小码作为四种基准；先在小体积精确求解，再用
MPS/张量网络或其他声明方法扩展；分别收敛时间步、bond
dimension、截断和体积。拓扑任务同时记录对称性和边界，热任务同时记录浴假设和详细平衡，纠错任务同时记录噪声相关与解码延迟。

\textbf{目标函数。} 例如给存储时间 \(T\)、控制能量代价 \(C[u]\)
和逻辑通道误差，定义
\(J[u]=1-F_e(\mathcal N_L^u)+\lambda C[u]\)。权重与任务由应用指定；不能称所有自然系统必然最小化此
\(J\)。对动力学梯度做有限差分或伴随独立核对，再与同预算基线比较。

\textbf{验收。} SSH
的有限链残差满足第44章公式；多体观测有明确收敛窗口；控制收益在留出噪声、温度和初态仍存在；逻辑错误确实低于未编码基准。若资源提高是唯一原因，应报告收益来自资源而非新的对象原理。无限体积相变、普适容错阈值与所有相互作用拓扑分类仍是独立问题。

\hypertarget{r7qftux89c4ux8303ux56faux5b9aeft-ux4e0eux7ea2ux5916}{%
\section{47.8 R7：QFT、规范固定、EFT
与红外}\label{r7qftux89c4ux8303ux56faux5b9aeft-ux4e0eux7ea2ux5916}}

\textbf{目标。} 让场论模块成为有定义、有阶数、有误差的计算接口。

\textbf{步骤。}
固定作用量、度规、生成元、规范参数、调节与减除条件；从第25/28章的一圈基准扩展到同阶完整振幅；验证
Ward/Slavnov--Taylor
恒等式及物理输出的规范参数独立性；对软/共线发散指定包容观测与分辨率，逐项展示虚实贡献的抵消。带电渐近态的问题不能由一个形式
LSZ 公式略去。

\textbf{EFT 任务。}
固定算符基与冗余消去方式，匹配时保留轻场圈的两侧抵消；在相同方案做系数
running 和阈值转换；把误差拆为圈阶、\(E/\Lambda_*\)
幂次、参数和输入矩阵元。第26章的重场 Schur 补是精确代数测试，但不是所有
EFT 的收敛定理。

\textbf{验收。}
同一物理截面用两个相容表示计算一致；不随任意重整化尺度变化到所保留阶数之外；光学定理和对称恒等式在该精度成立。若改变方案后参数未重新匹配，数值差不能被称作新物理。BRST
局部微扰闭合不自动解决 Gribov/global gauge 问题。

\hypertarget{r8ux975eux5faeux6270ux4e0eux8de8ux5c3aux5ea6}{%
\section{47.9
R8：非微扰与跨尺度}\label{r8ux975eux5faeux6270ux4e0eux8de8ux5c3aux5ea6}}

\textbf{目标。} 将``跨尺度统一''落实为同一指定观测的可证明一致性。

\textbf{第一步有限模型。} 用第45章的精确单环检验代码和归一化；再实现有限
\(SU(N)\) 链变量更新，检验 detailed
balance、热化与自相关，记录有效样本数和误差。只有正确平稳分布而没有充分混合，不代表采样估计可靠。

\textbf{第二步受控极限。} 选择一个明确 Euclidean
可观测，在多个体积、格距和裸参数上固定同一物理尺度；区分统计误差、格距误差、有限体积和参数调谐。研究反射正性和转移矩阵的适用条件，不能把有限测度存在性直接升级为完整连续量子场论。

\textbf{第三步 Core 桥。}
对粗映射先证明读出推前恒等式；若要干预响应，必须另给实时实现和操作插槽。比较有效作用截断时输出未保留多环项对指定观测的误差，避免把
Core 商和 Wilsonian 积分同名化。

\textbf{验收。}
对选定量能给独立算法交叉核对和可追踪极限；若不能建立连续极限，就只发布有限格点结果。四维质量隙、一般手征规范理论与实时解析延拓均不因算出一个有限谱而解决。

\hypertarget{r9ux89c4ux8303ux7fa4ux4e16ux4ee3ux4e0eux9ad8ux80fdux4e00ux81f4ux6027}{%
\section{47.10
R9：规范群、世代与高能一致性}\label{r9ux89c4ux8303ux7fa4ux4e16ux4ee3ux4e0eux9ad8ux80fdux4e00ux81f4ux6027}}

\textbf{目标。}
使第31章条件性三代定理与第36章高能双候选分类成为可比较候选，而非相互混用的结论。

\textbf{冻结表。} 指定全局群 \(G\)、全部左手 Weyl 表示、标量、Yukawa
允许项、质量阈值和规范嵌入。算 beta 时用该能区全部活跃场；\(SU(3)_L\) 的
UV 渐近自由不等于破缺后 \(SU(2)_L\) 在任意假想高能延拓中渐近自由。

\textbf{路线A。}
在三色已由独立观测固定的条件下，研究横向反常约束是否比直接输入三代减少独立信息，并求额外粒子和味预测。这里三色是输入，结论可以是条件的，仍有研究价值。

\textbf{路线B。} 若希望同时选择色数，从不使用目标 \((3,3)\)
的独立物理选择规则出发，检查它能否排除
\((5,5,2)\)；如真空稳定、可实现破缺或某个严格全局异常约束。每条规则都须对完整候选计算，而不能因为五色不合目标就增加标量或改变电荷。

\textbf{路线C。}
扩大表示类别研究反模型：非最小谱、矢量伴侣、多横向因子可能改变算术。原分类只在冻结类成立；扩大类别时重新分类，不保留原唯一性标题。

\textbf{验收。} 完整场表的局域和所需全局异常证书、破缺极小点、质量矩阵和
running
窗口；至少一条未用于选择的低能预测。若所有额外条件都是重新编码已知色/代数，只能称相容性重述。

\hypertarget{r10ux5473ux8d28ux91cfux4e2dux5faeux5b50ux548cux53efux89c2ux6d4bux5339ux914d}{%
\section{47.11
R10：味、质量、中微子和可观测匹配}\label{r10ux5473ux8d28ux91cfux4e2dux5faeux5b50ux548cux53efux89c2ux6d4bux5339ux914d}}

\textbf{目标。}
将第37章秩一耦合限制转为实际过程间的共参数关联，并尝试有物理理由的
Yukawa 降维。

\textbf{步骤。}
先在指定标量与对称性下列出所有可重整化质量项，求质量矩阵、单位变换和重轻混合；然后推导
\(Z'\)
的左右手与混合耦合。第37章只适用于声明的轻扇区幺正、共同双重态和单秩一差异；若
exotic 混合不可忽略，须计算额外非幺正修正，不能原样套用下界。

\textbf{从耦合到数据。} 匹配 \(\Delta F=2\) Wilson 系数，按相同算符基
running
到强子尺度，再使用有来源和协方差的矩阵元。对多个中性介子系统和稀有衰变共同拟合；把其它玻色子、右手电流和相消空间包括在竞争模型里。保守总
FCNC 下界不是某一个单通道振幅的正下界。

\textbf{质量路线。} 对拟议几何/家族结构先计算允许 Yukawa
的自由参数维数，再与可观测维数比较。给足任意复矩阵就不再预测混合。若有真正限制，可用部分谱拟合而留出
CKM/PMNS 和 CP 量；Dirac、Weinberg 与 seesaw
实现分别计数。轻中微子质量矩阵本身不能唯一反演重 seesaw 参数。

\textbf{验收。}
至少一个不能由剩余自由参数随意调整的和式或留出预测；完整误差传播；模型排除条件明列。不能把一个优美拟合点称作质量起源证明。

\hypertarget{r11ux5168ux5c40ux5f02ux5e38framing-ux4e0eux51e0ux4f55ux6307ux6807}{%
\section{47.12 R11：全局异常、framing
与几何指标}\label{r11ux5168ux5c40ux5f02ux5e38framing-ux4e0eux51e0ux4f55ux6307ux6807}}

\textbf{目标。}
判断原论文的相位/拓扑直觉是否能产生真实且自然的粒子选择，而不混同各种拓扑不变量。

\textbf{入口。} 给定带指定切丛结构和规范丛的流形、具体 Dirac
型算符、边界条件和 completion 理论，区分局域异常多项式、全局相位、eta
不变量、手征指标及 framing 响应。不能用一个 Betti 数替代未构造的 Dirac
指数。

\textbf{步骤。} 先在低维可算模型计算谱流或边界响应；再构造明确的
inflow/descent
映射并检验拼接、自然性与背景运输；最后把响应类送入第32章算术条件。必须实际得到
\(\chi\)，而不是选择使三代唯一的 \(\chi=1\) 后再寻找解释。

\textbf{验收。} 给同一 completion
在生成背景上的完整计算，证明所需自然映射，并检验外置理论是否真的不能改变对象自治证书。若不能得到唯一
\(\chi\) 或因不同 completion
有多个类，就保留相应多个候选。局域异常抵消不自动排除全局异常，也不自动证明该拓扑路线存在。

\hypertarget{r12ux5f15ux529bux4e0eux5b87ux5b99ux5b66ux7684ux53d7ux63a7ux5165ux53e3}{%
\section{47.13
R12：引力与宇宙学的受控入口}\label{r12ux5f15ux529bux4e0eux5b87ux5b99ux5b66ux7684ux53d7ux63a7ux5165ux53e3}}

\textbf{范围。}
本书没有量子引力构造。若统一目标包含引力，首先须新增时空几何动力学和量子测量接口，而非把现有
Hilbert 理论换名。

\textbf{第一入口。}
在给定全局双曲背景上研究自由场与局域探测器，明确态、重整化应力能张量和局域观测；比较不同观察者粒子概念时用同一探测器响应，不能把非唯一粒子分解当作直接矛盾。

\textbf{第二入口。} 半经典方程
\(G_{\mu\nu}+\Lambda g_{\mu\nu}=8\pi G_N\langle T_{\mu\nu}\rangle_{\rm ren}\)
是新增候选接口，还需反项、守恒和适用的涨落尺度。它不是量子引力闭合解，也不自动处理测量条件化后的非线性反馈。任何随机实际化耦合几何都须检查
Bianchi 恒等式与能量交换是否相容。

\textbf{第三入口。}
对宇宙学候选明确背景初态、扰动、观测传递函数和参数先验，再谈功率谱等留出预测。暗物质、暗能量、宇宙常数和黑洞信息问题各需新的具体模型和判据；本书不从``对象存在''一句话推导这些现象。

\textbf{验收。}
至少一个受控背景模型在已声明误差下可计算且守恒一致，再考虑升级几何动力学。缺少此起点时，不应把宏大统一图作为已完成理论的一部分。

\hypertarget{r13ux5f62ux5f0fux5316ux8bc1ux660eux4ee3ux7801ux4e0eux6587ux732eux8d23ux4efb}{%
\section{47.14
R13：形式化证明、代码与文献责任}\label{r13ux5f62ux5f0fux5316ux8bc1ux660eux4ee3ux7801ux4e0eux6587ux732eux8d23ux4efb}}

从有限矩阵和整数分类入手：Choi--Kraus、SRN Gram、三角 CP
界、四维见证、秩一 FCNC
和全域整数不等式可作为形式化候选。先编码类型、共轭/转置约定和前提，再证明定理；一项数值随机测试不是替代品。

代码分模型、计算、证书三层。模型层拒绝单位或基不匹配；计算层返回近似值与误差；证书层保留假设和失败状态。遇到非
CP
数据可提供有标记的拟合模型，但保留原残差。证明注册表必须区分标准先例与本项目新组合；投稿前对精确表达做先例查重和独立审稿。

\textbf{已达到的阶段。}
第49--59章的有限操作、响应、表示、证书、误差、混合、平均、条件生存、完整有限复纯化、真实记录记忆和指定非交换动力学继续保留。第60章证明所有合法效应分离完整密度态，并证明原复操作模型对任意联合态满足局部层析。第61章从完整凸极点导出向量表示，构造同辅助载体的实际酉关系，并以明确等距嵌入处理异型辅助空间。第62章把共同固定深度因果控制器的真实历史通道接回原协议，证明初态误差加逐动作记录缺陷的预算。当前主包共35个模块、747个显式定理和985项唯一命名公理审计。语法声明、具体见证与物理规律分别计数。100项运行方法、4个演示及精确证书按实际门记录；三套旧主文程序的757项检查沿用3.0的执行身份及未变源码哈希。

\textbf{下一批可验证义务。}
已补的有限量子性质仍以完整复矩阵、合法效应和所声明过程为前提。要将其用于量子选择，必须逐项审查重建定理的全部背景，并给出UEOT物理实现为何满足选择原则的独立理由。记录预算目前要求共同固定深度因果控制器；变长停止协议须证明补齐语义，相干控制或额外环境须有新的稳定化距离合同。动力学继续研究可验证的有限收缩或本原条件、一般正谱存在与长期条件过程。精确程序细化、真实QDRR数据、完整谱与阈值、连续变量能量桥各有独立验收门。

真实硬件数据必须另标来源和原始计数，不能把合成结果重新命名为实证。只在证明脚本实际检查成功、输入身份与适用域同时明确时，才给对应声明
FORMAL 标签。

\hypertarget{ux7edfux4e00ux5b8cux6210ux6807ux51c6ux4e0eux7814ux7a76ux505cux6b62ux6761ux4ef6}{%
\section{47.15
统一完成标准与研究停止条件}\label{ux7edfux4e00ux5b8cux6210ux6807ux51c6ux4e0eux7814ux7a76ux505cux6b62ux6761ux4ef6}}

每个工作包完成时输出：问题定义、依赖、模型/参数、证明或反例、算法与误差、可观测合同、独立验证、未覆盖域。新结果回写同一主文并更新依赖，不创建一堆互相缺页的``最终版''。

若某方向只增加不能检验的假设、用数据反填选择条件、或反复在未固定模型上拟合，应停止该分支并记录原因。若得到反例使原猜想失败，应保留反例并削弱结论。构造统一理论的实际进展包括严格排除错误路线，而不只包括得到期待的数字。

\hypertarget{ux901aux5411ux5b8cux5907ux7269ux7406ux7406ux8bbaux7684ux4f9dux8d56ux8d26ux672c}{%
\section{47.16
通向完备物理理论的依赖账本}\label{ux901aux5411ux5b8cux5907ux7269ux7406ux7406ux8bbaux7684ux4f9dux8d56ux8d26ux672c}}

用户提出的总体目标是基于 UEOT
构建完备的量子物理理论。为使这个目标能够持续推进，本版以四个层面分别检验完成度：数学定义与逻辑闭合、对物理实现的选择能力、跨模型/尺度的可组合预测、独立实验约束。某一层的进展不替代其他层。

\begin{longtable}[]{@{}
  >{\raggedright\arraybackslash}p{(\columnwidth - 6\tabcolsep) * \real{0.2100}}
  >{\raggedright\arraybackslash}p{(\columnwidth - 6\tabcolsep) * \real{0.2000}}
  >{\raggedright\arraybackslash}p{(\columnwidth - 6\tabcolsep) * \real{0.2800}}
  >{\raggedright\arraybackslash}p{(\columnwidth - 6\tabcolsep) * \real{0.3100}}@{}}
\toprule\noalign{}
\begin{minipage}[b]{\linewidth}\raggedright
当前接口
\end{minipage} & \begin{minipage}[b]{\linewidth}\raggedright
已获得的实质结果
\end{minipage} & \begin{minipage}[b]{\linewidth}\raggedright
仍需关闭的义务
\end{minipage} & \begin{minipage}[b]{\linewidth}\raggedright
下一个可以判成成功或失败的产物
\end{minipage} \\
\midrule\noalign{}
\endhead
\bottomrule\noalign{}
\endlastfoot
Core → 量子扇区 &
三模型、经典纯化障碍、实层析反例；复全态层析、完整纯性与辅助酉关系 &
独立原则来源、全部重构背景与其他模型 &
在明确候选类内核对全部重建前提；给出可检验的UEOT物理桥 \\
量子扇区 → 内核 & 合法操作、全词与有限表示、证书可靠性 &
运行程序通用细化、连续域与无限协议 & 同一输入的通用细化定理及误差界 \\
内核 → 持续与误差 &
真实记录记忆、共同固定深度控制器逐动作预算、指定非交换条件态精确收敛 &
变长停止与额外环境的预算；一般正谱与混合 &
证明明确扩张的控制语义；核验一般动力学充分条件 \\
已知QFT → 独立物理选择 & 作用量、QED/EFT推导、异常与候选谱 &
规范群、现实谱、质量、真空和阈值 &
完整冻结模型的一条未被用于拟合的共参数预测 \\
候选 → 自然界 & 见证、校准与顺序统计方案 & 真实原始数据与竞争模型约束 &
可复查、可拒绝的实验结论 \\
\end{longtable}

这张账本的验收结果是具体命题或反例，不以``更接近终极真理''的形容词替代。总著将新推导写回其所属章节，同时保留反模型和适用条件；开放条目随证据关闭，失败路线也留下可追溯的理由。

\hypertarget{ch-48}{%
\chapter{48. 理论总著、状态接口与持续维护}\label{ch-48}}

\hypertarget{ux552fux4e00ux7406ux8bbaux4e3bux6e90}{%
\section{48.1 唯一理论主源}\label{ux552fux4e00ux7406ux8bbaux4e3bux6e90}}

本书固定名称为《UEOT物理学》，稳定文档身份为
\texttt{UEOT-PHYSICS}，当前正式整合版为3.4。规范正文固定在仓库
\texttt{physics/qm/book/UEOT\_\allowbreak PHYSICS.\allowbreak md}；同目录 PDF
由该正文生成。第0--62章、附录A--F构成连续完整文档，保留全部原有物理主题，并把新证明整合进入相应论证。不要求读者从若干旧版本补齐一个证明。

版本报告、PR、日志和数学模块各有自己的身份，但不再与总著争夺理论正文的权威。源文件清单、章节修订档案及证据回执记录``这句话依据什么、哪一部分被检查过''。历史文件保留沿革责任；它们的旧状态不会覆盖较新的已验证内核，新的内核也不会为历史上未形式化的物理命题追授机器证明。

\hypertarget{ux4e09ux4e2aux72ecux7acbux72b6ux6001ux8f74}{%
\section{48.2
三个独立状态轴}\label{ux4e09ux4e2aux72ecux7acbux72b6ux6001ux8f74}}

正文论证状态仍采用
D/P/K/H/S/E/N/O：定义、手写命题/推导、外部重型定理、附加假设、解释、经验输入、否定结果、开放问题。这一轴回答``文本正在做什么''。物理依赖使用仓库合同
C/B/S/P/E/O/X：Core结果、领域桥、扇区选择、标准物理输入、有效模型假设、可观测后果、开放科学猜想。这一轴记录结论的依赖和物理角色；可观测后果
O 本身不表示已经有实验数据。两套字母在 P、S、E、O
等位置含义不同，记录时必须写明字段名。

3.1版明确修正3.0版在该字段中误用
E=经验、O=开放、X=反例的冲突，并在逐章档案保留迁移理由。经验来源继续由论证状态及数据身份记录；反例使用论证状态
N，不能迁成开放猜想
X；已经有标准数学证明而本项目尚未形式化的事项属于开发义务，也不因此成为
X。上游合同未被改写，标签修订不新增或撤销数学定理。

第三轴 \texttt{verification\_\allowbreak kind}
记录解析推导、精确算术、合成数值、Lean形式检查或真实实验。一个状态为 P
的手工定理不因此已被 Lean 检查；一个由标准物理输入 P
支持的形式化结论也不因此成为 Core
无条件定理。定理内容、适用域和检查证据必须一同提取。

\begin{Verbatim}[breaklines,breakanywhere,fontsize=\scriptsize]
document:
  id: UEOT-PHYSICS
  edition: "3.4"
  canonical_source: physics/qm/book/UEOT_PHYSICS.md
  derivative_pdf: physics/qm/book/UEOT_PHYSICS.pdf
  chapter_range: [0, 56]
  fundamental_unique_derivation_complete: false
  four_dimensional_nonperturbative_qft_complete: false
  quantum_gravity_complete: false
  new_experimental_data_in_this_edition: false
  formal_proof_assistant_used_for_listed_kernel_claims: true
claim_schema:
  required:
    - id
    - argument_status
    - dependency_class
    - statement_and_hypotheses
    - domain_units_protocol
    - proof_or_primary_source
    - countermodels_and_exclusions
    - verification_kind_and_receipt
  argument_status: [D, P, K, H, S, E, N, O]
  dependency_class: [C, B, S, P, E, O, X]
  verification_kind: [analytic, exact_arithmetic, synthetic, lean, experiment]
maintenance:
  preserve_existing_result_ids: true
  update_affected_chapters_and_dossier: true
  bind_verification_to_actual_source_hashes: true
  require_separate_evidence_for_stronger_scope: true
\end{Verbatim}

该模式是文档与证据的组织合同，完整科研知识图谱仍需逐条扩展。书中数字按其实际计数口径给出，不能把一次测试中的多个参数点称为多个普遍物理定理。

\hypertarget{ux5fc5ux987bux643aux5e26ux524dux63d0ux7684ux9ad8ux98ceux9669ux7ed3ux8bba}{%
\section{48.3
必须携带前提的高风险结论}\label{ux5fc5ux987bux643aux5e26ux524dux63d0ux7684ux9ad8ux98ceux9669ux7ed3ux8bba}}

\begin{longtable}[]{@{}
  >{\raggedright\arraybackslash}p{(\columnwidth - 2\tabcolsep) * \real{0.3000}}
  >{\raggedright\arraybackslash}p{(\columnwidth - 2\tabcolsep) * \real{0.7000}}@{}}
\toprule\noalign{}
\begin{minipage}[b]{\linewidth}\raggedright
可提取结论
\end{minipage} & \begin{minipage}[b]{\linewidth}\raggedright
必须同时保留
\end{minipage} \\
\midrule\noalign{}
\endhead
\bottomrule\noalign{}
\endlastfoot
条件三代三色唯一 &
第31章全部低能弱窗口与谱前提；第36章正确高能计算的五色反例 \\
Born吸收 & 有限谱、指定SSE、QND条件；简并只到本征子空间 \\
非零循环相位 & 排除特定Gaussian对角扩散；更宽经典跳跃仍可能解释 \\
四能级见证 &
固定指针、SRN和校准；排除经典对角随机相位类，不排除量子环境 \\
188次零事件 &
理想事件下界、固定样本或合法超鞅；实际校准可能增大样本需求 \\
有限坐标覆盖全部未来 &
实际生成词及所有局部拉回恒等式，固定允许库与上下文 \\
通道误差累积 &
完整Kraus、共同载体、实际缺陷上界；一般记忆/旁系统合同另证 \\
后验有界 & 双侧事件正概率及明确下界；稀有事件可能使界失去信息 \\
格点模型存在 & 指定有限Euclidean规范模型；连续极限和四维质量隙仍独立 \\
形式化内核存在 &
仅列举定理和实例的源码、工具链与实际审计；非整部物理理论已被机器证明 \\
\end{longtable}

\hypertarget{ux5ba1ux67e5ux6539ux5199ux4e0ecoreux56deux9988}{%
\section{48.4
审查、改写与Core回馈}\label{ux5ba1ux67e5ux6539ux5199ux4e0ecoreux56deux9988}}

每次修改先确认命题身份及影响域：是改正一个错误、补出原来遗漏的前提，还是证明了一个更强结论。改正时保留旧表述的问题和最小反例；增强时列新增证明及受影响下游章节。精确全文、构建脚本和证据应在本地完成核查后，再进行必要的集中提交与远端验收。

本版针对原第5.3节加入逐分支 intertwining
条件及反例；第8章拆分写入等距映射与完整酉；第24、27章修复参数分解并补回顶点UV推导；第35章明确完整cq通道；第40章校正范数与简并标签。这些修订与新增第49--50章同属理论建设：排除不成立的推断使内核更可靠。

可回馈通用数学Core的内容，应先抽象成不以特定物理结论为前提的一般定理，例如响应因子、合法递归、误差组合、持续事件或不可识别性。复Hilbert、Born、特定规范群、作用尺度和质量谱属于领域输入，不能反填为对象数学公理后再宣称无条件得到目标。冻结Core的变更继续通过其治理流程；本版只在QM领域完成实现与著作整合。

\hypertarget{ch-49}{%
\chapter{49. UEOT-QM 的已验证有限量子内核}\label{ch-49}}

\hypertarget{ux5185ux6838ux627fux62c5ux7684ux79d1ux5b66ux4efbux52a1}{%
\section{49.1
内核承担的科学任务}\label{ux5185ux6838ux627fux62c5ux7684ux79d1ux5b66ux4efbux52a1}}

本章把第1、4、5、33章的对象与量子操作思想连接到已经实际检查的数学实现。这里的内核是一组共享类型、依赖和可组合定理的有限量子结构：合法态与操作产生记录概率；全部未来记录决定协议相对的响应对象；有限表示证书使该对象可以计算；正概率后验在可实现表示空间内闭合；完整有限适应树保持概率归一。第50章再建立这些实际量子对象的误差合同。

这一结构实现 UEOT
的核心研究方法：先声明允许的干预和未来目标，再问哪些差异能够改变未来响应。它不把物理对象预先等同于某一种粒子，也不把相同的响应坐标误认为完全相同的微观态。它的物理输入是有限复
Hilbert 空间、密度态、Kraus 操作与 Born
配对，属于量子扇区和标准物理实现；第2章的量子重构背景及附加原则承担另一层选择责任。

代码位于仓库
\texttt{formalization/ueot-qm}，可执行的有理数模型及导出器位于
\texttt{physics/qm/kernel}。本书是连续理论说明，代码和逐项日志是相应命题的可复查证明产品。Core
v3 的冻结数学和计数没有因本章新增物理实现而改变。

\hypertarget{ux5b9eux9645ux77e9ux9635ux8bb0ux5f55ux4e0eux54cdux5e94ux5546}{%
\section{49.2
实际矩阵、记录与响应商}\label{ux5b9eux9645ux77e9ux9635ux8bb0ux5f55ux4e0eux54cdux5e94ux5546}}

固定有限指标集 \(\iota\)，记 \(\mathcal H=\mathbb C^{\iota}\)。态为
\(\rho\succeq0\) 且 \(\operatorname{Tr}\rho=1\) 的矩阵。对允许动作
\(a\)，完整仪器给出有限族 \(K_{aok}\)，满足 \[
\sum_{o,k}K_{aok}^{\dagger}K_{aok}=I,
\qquad \Phi_{a,o}(X)=\sum_kK_{aok}XK_{aok}^{\dagger}.
\] \texttt{Density}、\allowbreak \texttt{Instrument}
的合法性字段是这些局部数学条件；没有将``全部未来响应正确''作为字段预先输入。分支正性由
\(X\succeq0\Rightarrow KXK^\dagger\succeq0\)
及有限和证明；总迹保持由迹循环性与完整性证明。特别地 \[
p_{a,o}(\rho)=\operatorname{Tr}\Phi_{a,o}(\rho)\ge0,
\qquad \sum_op_{a,o}(\rho)=1.
\]

词 \(w=((a_1,o_1),\ldots,(a_n,o_n))\) 按实际时间顺序记录操作，定义 \[
R_\rho(\varnothing)=1,\qquad
R_\rho(w)=\operatorname{Tr}\bigl[\Phi_{a_n,o_n}\circ\cdots\circ\Phi_{a_1,o_1}(\rho)\bigr].
\] 协议库固定以后，\(\rho\sim\sigma\) 定义为对\textbf{所有有限词}都有
\(R_\rho(w)=R_\sigma(w)\)。这一量词不限于证书曾经列出的有限深度。对目标因子
\(S\)，若 \(S(\rho)=S(\sigma)\) 蕴含响应相同，响应函数便唯一地经由 \(S\)
的可达像分解；在不可达元素上的任意延伸不属于唯一性结论。这正是第1章的协议相对充分性在该量子实现中的实例。

\protect\hypertarget{p-qkr-01}{}{}

\textbf{P-QKR-01（响应对象的后验闭合）。} 同一动作结果的概率 \(p>0\)
时，实际后验为 \(\rho_{a,o}=\Phi_{a,o}(\rho)/p\)，且 \[
R_{\rho_{a,o}}(w)=\frac{R_\rho((a,o)::w)}{R_\rho([(a,o)])}.
\]
因此响应等价的两态具有相同分支概率；在正概率分支上其后验仍响应等价。更新由响应类本身决定，并留在某个实际密度态的响应像内。

\textbf{证明。} 分支映射和其后所有操作均为线性，将标量 \(1/p\)
从复合与迹中提出即得公式。分子分母由当前响应函数决定，所以更换同一响应类的代表态不改变后验响应。后验矩阵正半定且迹一，故更新输出可实现。若分支概率为零，正半定矩阵零迹蕴含分支矩阵为零；此时没有归一化后验，接口返回
\texttt{none}。\(\square\)

对应 Lean 结果包括
\texttt{response\_\allowbreak posterior}、\allowbreak \texttt{posterior\_\allowbreak response\_\allowbreak eq}、\allowbreak \texttt{next\allowbreak Response\_\allowbreak eq\_\allowbreak none\_\allowbreak iff}、\allowbreak \texttt{exists\allowbreak Unique\_\allowbreak closed\allowbreak Reachable\allowbreak Update}。这比仅证明某个集合论商存在多了一步：实际更新映射不依赖所选代表元，并保持可实现域。

\hypertarget{ux6709ux9650ux751fux6210ux7a7aux95f4ux4e0eux5168ux8bcdux5750ux6807}{%
\section{49.3
有限生成空间与全词坐标}\label{ux6709ux9650ux751fux6210ux7a7aux95f4ux4e0eux5168ux8bcdux5750ux6807}}

定义 Heisenberg 拉回及词效应 \[
\Phi_{a,o}^*(E)=\sum_kK_{aok}^{\dagger}EK_{aok},\qquad
E_{\varnothing}=I,\quad E_{(a,o)::w}=\Phi_{a,o}^*(E_w).
\] 迹循环性逐步给出
\(R_\rho(w)=\operatorname{Re}\operatorname{Tr}(E_w\rho)\)。所有词效应为
Hermitian 矩阵，令 \[
W=\operatorname{span}_{\mathbb R}\{E_w:w\text{为有限词}\}.
\] 它是包含 \(I\)
并对所有允许分支拉回封闭的最小实子空间：包含与闭包使所有词效应通过归纳进入该空间；反过来，词效应的线性张成对拉回封闭。有限矩阵空间保证
\(W\) 有有限个真实词效应生成元。

\protect\hypertarget{p-qkr-02}{}{}

\textbf{P-QKR-02（局部有限表示决定全部未来）。} 取有限生成族
\(B_j=E_{w_j}\)，其中 \(B_u=I\)。若每个分支满足实际矩阵恒等式 \[
\Phi_{a,o}^*(B_j)=\sum_i A_{a,o;ji}B_i,
\] 定义 \(x_j(\rho)=\operatorname{Re}\operatorname{Tr}(B_j\rho)\)，则 \[
x(\rho)=x(\sigma)\ \Longleftrightarrow\ R_\rho=R_\sigma.
\] 未归一分支坐标为 \(y=A_{a,o}x\)，其概率为 \(p=y_u\)；在 \(p>0\)
时，后验坐标为 \(y/p\)，且属于实际态坐标的可达像。

\textbf{证明。} 实际生成词保证所有 \(B_j\) 属于
\(W\)；单位与局部闭包保证所有词效应位于
\(\operatorname{span}\{B_j\}\)，故两空间相同。相同坐标于是给出与每个词效应的相同配对；反向直接对生成词读取响应。对拉回等式与
\(\rho\) 配对，即得分支坐标公式。后验公式由 P-QKR-01
的实际归一化而来。\(\square\)

这里系数的第 \(j\) 行展开被拉回的第 \(j\) 个效应，因而更新是
\(Ax\)。改变行列约定必须同时改变证书、解释器和证明。Lean 的
\texttt{Finite\allowbreak Observable\allowbreak Presentation}
允许冗余生成族；它不要求输入一份尚未证明的最小基。Python 的精确 Gram
独立性检查提供额外信息，但该搜索和消元程序尚没有通用 Lean 细化证明。

在通常 \(d\) 维 Hermitian 矩阵空间中，对角矩阵单位与每对非对角的实、虚
Hermitian 矩阵单位构成 \(d+2\binom d2=d^2\) 个实基，故手工线性代数给出
\(\dim_{\mathbb R}W\le d^2\)。当前 Lean
接口已经证明有限维性和有限词生成存在性，不把这条具体维数计数或 Python
搜索的终止性计入机器证明范围。

单比特实例说明协议相对性的实质：只允许 Z
投影读取时，\(W=\operatorname{span}\{I,Z\}\)；加入 Hadamard 操作后可产生
X 方向；再加入相位门 S 后产生 Y 方向，完整 Hermitian
响应空间有四个实坐标。相同的二能级载体因访问库不同具有不同的预测表示。较少的响应坐标不自动提供较小
Hilbert 空间的物理实现、较少驱动器或较低信息成本。

\hypertarget{ux5b8cux6574ux6709ux9650ux81eaux9002ux5e94ux8bb0ux5f55ux5f8b}{%
\section{49.4
完整有限自适应记录律}\label{ux5b8cux6574ux6709ux9650ux81eaux9002ux5e94ux8bb0ux5f55ux5f8b}}

一个允许的有限确定性适应协议，是在每个非终端节点选择动作、并为该动作的\textbf{每个结果}提供后续子树的有限树。它可以不平衡，不要求所有路径等长；分支选择仅依赖已发生记录。每个叶子的概率是沿其路径复合分支映射后取迹。

\protect\hypertarget{p-qkr-03}{}{}

\textbf{P-QKR-03（全词坐标控制完整有限适应树）。}
对上述同一协议，叶概率非负，总和为一。若初态全词响应相同，或 P-QKR-02
的坐标相同，则每个叶概率相同。

\textbf{证明。}
在叶节点，叶质量就是当前非负矩阵的迹。内节点对各结果分别应用归纳假设，再用完整仪器的总迹恒等式求和，得到输入质量。对密度态，根质量为一。每个叶对应唯一的有限路径词，其概率等于该词响应；响应相同便给出所有对应叶的相同概率。\(\square\)

\texttt{Adaptive.\allowbreak lean}
形式化这一归纳、叶词单射以及叶律的坐标表示。任意删除不利叶、事后改变停时或只保留存活结果，会改变上述归一化合同；新的条件协议须另给事件概率和后选择误差。无限停时的终止与极限也需要另证。第46章的统计超鞅解决另一种条件采样问题，不能仅因都使用``适应''一词而互换定理。

\hypertarget{ux6709ux7406ux52a0ux6743ux4e0eux53efux53d8ux6709ux9650ux7d22ux5f15}{%
\section{49.5
有理加权与可变有限索引}\label{ux6709ux7406ux52a0ux6743ux4e0eux53efux53d8ux6709ux9650ux7d22ux5f15}}

实际精确程序保存 \(q_{aok}\in\mathbb Q_{\ge0}\) 和
\(M_{aok}\in M_d(\mathbb Q[i])\)，以 \[
\Phi_{a,o}(X)=\sum_kq_{aok}M_{aok}XM_{aok}^{\dagger}
\] 定义分支。取复矩阵 \(K_{aok}=\sqrt{q_{aok}}M_{aok}\)，逐项乘法证明
\(KXK^\dagger=qMXM^\dagger\)，且拉回、完整性具有相同转换。平方根的存在在实数域证明，数值证书无需近似该平方根。例如整数
Hadamard 矩阵配权重 \(1/2\) 即可精确表示标准 Hadamard 通道。

不同结果的 Kraus
项数、不同动作的结果数可以不同。\texttt{Ragged\allowbreak Instrument}
先允许结果依赖的有限项类型，再沿明确注入补零到统一类型。有限和沿注入重新索引，新增项为零，因而分支、拉回和
Gram
完整性不变。结果补零也只产生零概率的新结果。该定理要求实际注入与零扩展，不认可含糊的``数组长度兼容''，更不使用取模修补越界词索引。

\hypertarget{ux53efux8ba1ux7b97ux68c0ux67e5ux5668ux7684ux6570ux5b66ux53efux9760ux6027}{%
\section{49.6
可计算检查器的数学可靠性}\label{ux53efux8ba1ux7b97ux68c0ux67e5ux5668ux7684ux6570ux5b66ux53efux9760ux6027}}

复有理矩阵表示为 \(A+iB\)，乘法是 \((AC-BD)+i(AD+BC)\)，共轭转置是
\(A^T-iB^T\)。\texttt{Exact.\allowbreak Matrix\allowbreak Data.\allowbreak to\allowbreak Matrix}
将这份数据解释为真实复矩阵，并已证明其单射性及对上述代数运算、有限和与有理实数缩放的保持性。

数值证书 \(D\) 保存 Kraus 权重与矩阵、生成词 \(w_j\)、基矩阵
\(B_j\)、单位索引 \(u\) 和局部系数 \(A_{a,o;ji}\)。独立 Lean
检查器只核验五类有限条件：权重非负；各动作完整性；单位词为空；每个基矩阵等于其真实词效应；每个局部拉回等于给定有限展开。它没有``对所有未来预测正确''的假设字段。

\protect\hypertarget{p-qkr-04}{}{}

\textbf{P-QKR-04（精确证书可靠性）。} 若实际可计算检查得到
\(\operatorname{check}(D)=\mathrm{true}\)，则解释后的复矩阵形成合法仪器及有限
observable presentation，且 P-QKR-01 至 P-QKR-03
的响应、闭合后验与完整有限树结论成立。

\textbf{证明。} \texttt{check\_\allowbreak sound}
将布尔真值转成五类有限数学条件。代数解释将其提升为真实复矩阵等式；对词递归证明解释保持实际词效应，从而数值生成词与真实基绑定。最后应用
P-QKR-02 的局部闭包定理及 P-QKR-03。\(\square\)

六份已检查实例为 Z、H+Z、H+S+Z、可变结果/Kraus 数的 ragged、\(d=1\)
的复数例及 \(d=4\)
粗投影例。它们覆盖有理权重、非零虚部、零概率结果及非满响应维数。每份数据的
\texttt{checked} 声明由 Lean
内核实际接受，继而导出全词和树的后果。另有四类导出后数值篡改：权重、数值基、局部拉回系数、词顺序。每类既要求原真值证明被拒绝，又要求\textbf{相同坏数据}的假值证明成功，以排除导入或资源故障造成的伪负控。

这是通用数学检查器的可靠性与具体数值实例的接受记录。Python、JSON
parser、搜索器和输出格式的所有执行行为尚未统一形式化；六份实例的接受不能替代那个更强结论。第33章把已实现接口与规格中的未来接口分别标记。

\hypertarget{ux4e0eux4e09ux4ee3ux8c31ux53caux5176ux4ed6ux7269ux7406ux5c42ux7684ux8fdeux63a5}{%
\section{49.7
与三代谱及其他物理层的连接}\label{ux4e0eux4e09ux4ee3ux8c31ux53caux5176ux4ed6ux7269ux7406ux5c42ux7684ux8fdeux63a5}}

\texttt{Spectrum.\allowbreak lean}
由明确有限家族行、额外物质行和有效窗口条件求和；同质轻子模板、非零异常单位及补偿项条件导出选择器的立方方程，非负的额外屏蔽贡献则进入弱预算。随后
\texttt{Selector.\allowbreak lean} 对完整整数前提证明条件唯一的
\((N_c,N_f,k)=(3,3,1)\)。它已经超越仅在有限搜索框中找到一个三元组，但尚未从现实表示论、阈值和真空机制无条件导出那些前提。

第31与36章的区别仍然有效：在低能有效弱预算等前提下的条件唯一性，并不排除正确高能
\(SU(3)_L\)
条件中的五色候选。新增内核没有把独立高能候选偷偷从模型类中删除，也没有把拓扑指标指定为三代。质量、混合、全局异常和现实粒子表的物理选择责任继续由相应章节承担。

连续变量、一般过程张量、相对论量子场及第40章随机实际化保留各自解析推导和域条件。有限矩阵内核能够作为其中受控有限模型的检查基础；从有限截断跨到无界算符、无限体积或连续极限，需要该方向自己的误差与存在性定理。

\hypertarget{ch-50}{%
\chapter{50. 从实际 Born 概率到量子误差传播}\label{ch-50}}

\hypertarget{ux5168ux6548ux5e94ux5bf9ux8c61ux4e0eux6982ux7387ux754c}{%
\section{50.1
全效应对象与概率界}\label{ux5168ux6548ux5e94ux5bf9ux8c61ux4e0eux6982ux7387ux754c}}

本章完成上一版路线 QMU-04 的一个具体数学切片：对同一有限量子载体上的完整
Kraus
通道，先从实际矩阵建立可观测概率距离，再证明一步收缩、模型缺陷、轨迹误差及正事件后选择。所用量均由量子对象定义；收缩结论不是预装在证书字段中的假设。

定义效应 \(E\) 为 \(E\succeq0\) 且 \(I-E\succeq0\) 的矩阵，记 \[
P_E(\rho)=\operatorname{Re}\operatorname{Tr}(E\rho).
\] 这是 \texttt{Effect} 的实际定义，覆盖全部矩阵序区间 \([0,I]\)。

\protect\hypertarget{p-qst-01}{}{}

\textbf{P-QST-01（实际正配对与效应界）。} 对任意 \(X\succeq0\) 和效应
\(E\)，有 \[
0\le\operatorname{Re}\operatorname{Tr}(EX)\le\operatorname{Tr}X.
\] 特别地 \(0\le P_E(\rho)\le1\)，且 \(P_{I-E}(\rho)=1-P_E(\rho)\)。

\textbf{证明。} 由有限复矩阵正半定性取得 \(E=B^\dagger B\)。迹循环给出
\(\operatorname{Tr}(EX)=\operatorname{Tr}(BXB^\dagger)\)，右侧为正半定矩阵之迹，故非负。对
\(I-E\)
同样应用此结论，展开线性配对即得上界。对密度态使用迹一。\(\square\)

Lean 中的 \texttt{pairing\_\allowbreak nonneg} 从已有 PSD 谓词经 Mathlib
矩阵序因子分解取得 \(B\)，没有要求使用者另输入 Gram
因子。\texttt{Effect.\allowbreak probability\_\allowbreak nonneg}、\allowbreak \texttt{probability\_\allowbreak le\_\allowbreak one}
与 \texttt{probability\_\allowbreak complement} 因此适用于任意已有合法密度态。

\hypertarget{ux9010ux7ed3ux679cux62c9ux56deux4e0eux6240ux6709ux6709ux9650ux8bb0ux5f55}{%
\section{50.2
逐结果拉回与所有有限记录}\label{ux9010ux7ed3ux679cux62c9ux56deux4e0eux6240ux6709ux6709ux9650ux8bb0ux5f55}}

对一个完整仪器及每个结果各自的后续效应 \(E_o\)，定义 \[
F=\sum_o\Phi_o^*(E_o)=\sum_{o,k}K_{ok}^{\dagger}E_oK_{ok}.
\]

\protect\hypertarget{p-qst-02}{}{}

\textbf{P-QST-02（续效应的完整拉回）。} \(F\) 仍为效应，且 \[
\operatorname{Re}\operatorname{Tr}(FX)
=\sum_o\operatorname{Re}\operatorname{Tr}(E_o\Phi_o(X)).
\]

\textbf{证明。} 每项为正半定，故 \(F\succeq0\)。完整性给出 \[
I-F=\sum_{o,k}K_{ok}^{\dagger}(I-E_o)K_{ok}\succeq0.
\] 配对恒等式逐项使用迹循环性。\(\square\)

令所有 \(E_o=E\)，得到忘记结果的完整通道 \(\Phi=\sum_o\Phi_o\)
的效应拉回。只对一个选定结果取 \(E\)、其他结果取零，得到分支效应
\(\Phi_o^*(E)\in[0,I]\)。从 \(E_\varnothing=I\)
递归使用这个分支结果，便证明每个已有词效应 \(E_w\) 不仅 PSD，而且
\(I-E_w\) 也 PSD。因此全部有限词响应都是实际效应概率，自动位于
\([0,1]\)。

该结论对应
\texttt{Instrument.\allowbreak pullback\allowbreak Effects}、\allowbreak \texttt{branch\allowbreak Effect}、\allowbreak \texttt{record\allowbreak Test}
及
\texttt{record\allowbreak Effect\_\allowbreak complement\_\allowbreak pos\allowbreak Semidef}，使原有全词响应层与新距离层连接。它适用于结果依赖的后续效应；只证明总通道拉回而忽略结果标签不足以替代这个更完整的接口。

\hypertarget{ux5168ux6548ux5e94ux54cdux5e94ux8dddux79bbux53caux5176ux8fb9ux754c}{%
\section{50.3
全效应响应距离及其边界}\label{ux5168ux6548ux5e94ux54cdux5e94ux8dddux79bbux53caux5176ux8fb9ux754c}}

定义 \[
d_E(\rho,\sigma)=\sup_{0\preceq E\preceq I}|P_E(\rho)-P_E(\sigma)|.
\] 零效应使被取上确界的集合非空，P-QST-01
使其有界。这是真实全部效应的上确界；有限次层析测量或某个训练测试集没有被当作这个上确界。

\protect\hypertarget{p-qst-03}{}{}

\textbf{P-QST-03（可观测距离的基本合同）。} 有
\(0\le d_E\le1\)、对称性、\(d_E(\rho,\rho)=0\) 及三角不等式。并且 \[
d_E(\rho,\sigma)=0
\iff \forall E\in[0,I],\ P_E(\rho)=P_E(\sigma).
\] 任一允许库的任一有限记录满足 \[
|R_\rho(w)-R_\sigma(w)|\le d_E(\rho,\sigma).
\]

\textbf{证明。} 每个概率差的绝对值位于
\([0,1]\)，并受实数绝对值三角不等式控制。逐效应应用后取上确界，得到前三类结论。上确界为零当且仅当全部非负概率差为零。最后使用
P-QST-02 的 \(E_w\in[0,I]\)。\(\square\)

第4章的标准谱分解证明在有限维数学中识别
\(d_E=\tfrac12\|\rho-\sigma\|_1\)；取差矩阵的正谱投影达到正质量和，零迹保证正负质量相等。本轮
Lean 模块\textbf{尚未形式化这项识别，也尚未把零距离升级成密度矩阵相等的
metric
实例}，故其精确已验证接口是全效应概率伪距离及上述零判据。单一受限访问库可能无法区分不同密度态；不得把全效应集合与那种协议受限空间混为一谈。

\hypertarget{ux5b9eux9645ux901aux9053ux6536ux7f29ux4e0eux6a21ux578bux7f3aux9677}{%
\section{50.4
实际通道收缩与模型缺陷}\label{ux5b9eux9645ux901aux9053ux6536ux7f29ux4e0eux6a21ux578bux7f3aux9677}}

\protect\hypertarget{p-qst-04}{}{}

\textbf{P-QST-04（完整 Kraus 通道的一步稳定性）。} 完整仪器所定义的通道
\(\Phi\) 满足 \[
d_E(\Phi(\rho),\Phi(\sigma))\le d_E(\rho,\sigma).
\] 对任意两个同载体的态映射，定义真实最坏输出差 \[
\Delta_E(\Phi,\Psi)=\sup\bigl(\{0\}\cup
\{d_E(\Phi(\tau),\Psi(\tau)):\tau\text{为密度态}\}\bigr).
\] 若 \(\Phi,\Psi\) 为上述完整通道，则 \[
d_E(\Phi(\rho),\Psi(\sigma))
\le d_E(\rho,\sigma)+\Delta_E(\Phi,\Psi).
\]

\textbf{证明。} 任一输出效应通过 \(\Phi^*\) 拉回后仍是输入效应，且实际
Born 配对保持，所以每个输出概率差受输入 \(d_E\)
控制；取上确界得收缩。再在 \(\Phi(\sigma)\)
处插入中间态，使用三角不等式、收缩和 \(\Delta_E\)
的上界性质。定义中插入零，使空密度类型情况下的集合仍非空；对通常正维物理系统不改变值。\(\square\)

\texttt{channel\allowbreak State\_\allowbreak effect\allowbreak Distance\_\allowbreak le} 和
\texttt{channel\allowbreak State\_\allowbreak perturbation\_\allowbreak bound}
证明的正是这些实际对象之间的关系。\(\Delta_E\)
本身为非计算的上确界，不是一个已经实现的高效误差估计算法。第51章进一步给出同库随机化这一具体模型中的有理上界。任何运行时提供的数值上界
\(\widehat\Delta\) 还需独立证明覆盖
\(\Delta_E\)；有限抽样中的最大误差只给下界，不能反过来充当安全上界。

\hypertarget{ux4eceux91cfux5b50ux4e00ux6b65ux5b9aux7406ux63a5ux5165-core-ux7684ux7ec4ux5408ux8bc1ux4e66}{%
\section{50.5 从量子一步定理接入 Core
的组合证书}\label{ux4eceux91cfux5b50ux4e00ux6b65ux5b9aux7406ux63a5ux5165-core-ux7684ux7ec4ux5408ux8bc1ux4e66}}

给定两串完整通道 \(\Phi_j,\Psi_j\)，按实际时间顺序定义 \[
\rho_{j+1}=\Phi_j(\rho_j),\qquad
\sigma_{j+1}=\Psi_j(\sigma_j),\qquad
\delta_j=\Delta_E(\Phi_j,\Psi_j).
\]

\protect\hypertarget{p-qst-05}{}{}

\textbf{P-QST-05（真实量子轨迹的可组合误差）。} 对每个有限 \(n\)， \[
d_E(\rho_n,\sigma_n)
\le d_E(\rho_0,\sigma_0)+\sum_{j=0}^{n-1}\delta_j.
\] 对相同未来仪器库的任一有限词 \(w\)，亦有 \[
|R_{\rho_n}(w)-R_{\sigma_n}(w)|
\le d_E(\rho_0,\sigma_0)+\sum_{j=0}^{n-1}\delta_j.
\] 由于概率差至多一，使用中可将界截在一以内。

\textbf{证明。} 令 \(e_j=d_E(\rho_j,\sigma_j)\)。P-QST-04 给出实际递推
\(e_{j+1}\le e_j+\delta_j\)；概率距离和缺陷均非负。将其作为 \(L_j=1\)
的实例代入既有
\texttt{Transport.\allowbreak scalar\_\allowbreak transport\_\allowbreak certificate}，得到其递归预算上界。对
\(n\) 归纳，预算等于 \(e_0+\sum_{j<n}\delta_j\)。第二式再应用
P-QST-03。\(\square\)

关键进展在于递推来自实际 Kraus
通道的数学定义，已经不再把需要证明的量子一步不等式留给调用者。本轮没有改变
Core 的组合定理，而是完成其 QM 侧前提。Lean 对应
\texttt{channel\allowbreak Trajectory\_\allowbreak transport\_\allowbreak certificate}、\allowbreak \texttt{channel\allowbreak Trajectory\_\allowbreak effect\allowbreak Distance\_\allowbreak le\_\allowbreak sum}、\allowbreak \texttt{channel\allowbreak Trajectory\_\allowbreak response\_\allowbreak sub\_\allowbreak le\_\allowbreak sum}。

这个已形式化的轨迹合同具有共同有限载体；两条通道序列可使用不同的结果与
Kraus
指标类型，但每条序列内部的类型固定，实际矩阵可以随步数改变。若需逐步改变指标，可先显式补零到统一有限类型，再使用第49.5节的保持性定理。它没有自动包括任意外部量子旁系统的最坏输入，因此不能命名为
diamond
距离定理。一般自适应仪器误差需要显式保存经典历史、结果与量子记忆后，构造其共同完整通道；第35.2节给出相应纸面装配合同。仅比较丢弃结果后的通道，会漏掉对控制有用的标签信息。

\hypertarget{ux6b63ux4e8bux4ef6ux6761ux4ef6ux5316ux4e0eux7a00ux6709ux5206ux652f}{%
\section{50.6
正事件条件化与稀有分支}\label{ux6b63ux4e8bux4ef6ux6761ux4ef6ux5316ux4e0eux7a00ux6709ux5206ux652f}}

对同一仪器和同一结果 \(o\)，设
\(p=p_o(\rho)>0\)、\(q=p_o(\sigma)>0\)，令后验分别为
\(\rho_o,\sigma_o\)。

\protect\hypertarget{p-qst-06}{}{}

\textbf{P-QST-06（后选择误差及正支撑保护）。} 有 \[
|p-q|\le d_E(\rho,\sigma),\qquad
d_E(\rho_o,\sigma_o)\le \frac{2d_E(\rho,\sigma)}{p}.
\] 若 \(p\ge\eta>0\)，右侧可换成 \(2d_E(\rho,\sigma)/\eta\)。若已知
\(d_E(\rho,\sigma)<p\)，则 \(q>0\)，所以近似模型也具有合法后验。

\textbf{证明。} 对输出效应 \(E\)，记
\(a=P_{\Phi_o^*(E)}(\rho)\)、\(b=P_{\Phi_o^*(E)}(\sigma)\)。P-QST-02
使该拉回为效应，故 \(|a-b|\le d_E(\rho,\sigma)\)；对 \(E=I\) 得
\(|p-q|\) 的界。实际后验配对为 \(a/p,b/q\)，且 \(0\le b/q\le1\)。恒等式
\[
\frac ap-\frac bq=
\frac{a-b+(q-p)(b/q)}p
\] 与绝对值三角不等式给出 \[
\left|\frac ap-\frac bq\right|
\le\frac{|a-b|+|p-q|}{p}.
\] 逐效应代入并取上确界即得后验界。概率下界改写分母；\(q\ge p-d_E>0\)
给出正支撑保护。\(\square\)

这里并未证明选择性归一化收缩，且常数二不是最优性的声明。第4章从次归一算符迹范数出发的界使用另一种输入误差定义，与本章不可只看分子符号就互换。若一个分支在某模型中概率为零，必须保持``后验未定义''的状态，不能任意指定一个态以延续证明。

举一个纯粹的合同算例：若独立证书给初始误差
\(10^{-3}\)、二十步各有缺陷上界 \(2\times10^{-4}\)，则最终误差不超过
\(0.005\)；再选择概率至少 \(0.2\) 的事件，后验误差界为
\(0.05\)。如果第一模型的事件概率只保证 \(10^{-3}\)，误差上界 \(0.005\)
已不足以推出第二模型事件为正；此时第二后验甚至可能未定义。只有另行确认两个事件都正，才可代入条件化定理；所得上界超过一，只能留下平凡界，说明原证书不足以支持该稀有条件任务。这些数字演示误差合同的使用，不是实验测量或某个现成设备的性能。

\hypertarget{ux8bc1ux660eux6267ux884cux548cux7269ux7406ux4e3bux5f20ux7684ux4e09ux5c42ux95edux5408}{%
\section{50.7
证明、执行和物理主张的三层闭合}\label{ux8bc1ux660eux6267ux884cux548cux7269ux7406ux4e3bux5f20ux7684ux4e09ux5c42ux95edux5408}}

本章的首层成果是解析命题及其 Lean
对应；第二层是锁定工具链上的实际构建、公理检查和既有证书的兼容性回归；第三层才是某个真实器件或领域模型是否满足输入的量子实现与数值误差上界。前两层已经为后续物理研究提供可复用内核，第三层必须由具体制备、控制、噪声、阈值和数据支持。

QMU-04 因而登记为``有限共同载体完整 Kraus
通道及正事件后验切片完成''。全自适应含记忆过程的统一误差、旁系统稳定范数识别、可执行最坏误差上界算法、无界能量控制和连续极限仍分开列入第47章的工作包。下一步应优先从这些有明确接口的义务中选择，不通过给已有定理更强的标题来跳过它们。

\hypertarget{ch-51}{%
\chapter{51. 随机化量子实现与可认证误差预算}\label{ch-51}}

第50章把量子误差写成实际全部效应、全部输入态的上确界。这解决了比较对象的定义，却还没有为每种设备给出可计算上界。本章闭合一个经常出现且能严格计算的切片：实验先用经典随机数选择同一有限库中的某个完整量子操作，模型差异来自选择权重，或者来自可分别控制的库成员误差。这里的随机化是明确的操作机制；不把任意未知噪声都归入这个模型。

\hypertarget{ux4eceux6982ux7387ux6743ux91cdux6784ux9020ux771fux5b9eux91cfux5b50ux5bf9ux8c61}{%
\section{51.1
从概率权重构造真实量子对象}\label{ux4eceux6982ux7387ux6743ux91cdux6784ux9020ux771fux5b9eux91cfux5b50ux5bf9ux8c61}}

固定有限指标集 \(B\)、共同的有限维复 Hilbert 空间，以及概率权重 \[
p_b\ge0,\qquad \sum_{b\in B}p_b=1.
\] 每个库成员是完整仪器 \(\{K^{(b)}_{ok}\}_{o,k}\)，满足
\(\sum_{o,k}K^{(b)\dagger}_{ok}K^{(b)}_{ok}=I\)。本节在一个共同的有限结果集和
Kraus 指标集上书写；指标不齐时可先采用第49.5节已验证的补零适配。

\protect\hypertarget{p-qmx-01}{}{}

\textbf{P-QMX-01（概率混合的实际实现）。} 对任意密度态库
\(\rho_b\)，\(\rho_p=\sum_b p_b\rho_b\) 是密度态。以新指标 \((b,k)\)
定义 \[
L^{(p)}_{o,(b,k)}=\sqrt{p_b}\,K^{(b)}_{ok},
\] 得到完整量子仪器。其分支、Heisenberg 拉回与忘记结果的通道分别为 \[
\begin{aligned}
\Phi^{(p)}_o(X)&=\sum_b p_b\Phi^{(b)}_o(X),\\
\Phi^{(p)*}_o(E)&=\sum_b p_b\Phi^{(b)*}_o(E),\\
\Phi^{(p)}(X)&=\sum_b p_b\Phi^{(b)}(X).
\end{aligned}
\]

\textbf{证明。} 正半定锥对非负加权和封闭，迹线性及权重归一给出
\(\operatorname{Tr}\rho_p=1\)。新仪器满足 \[
\sum_{o,b,k}L^{(p)\dagger}_{o,(b,k)}L^{(p)}_{o,(b,k)}
=\sum_b p_b\sum_{o,k}K^{(b)\dagger}_{ok}K^{(b)}_{ok}
=\Bigl(\sum_b p_b\Bigr)I=I.
\] 逐项展开 \(LXL^\dagger\) 或 \(L^\dagger EL\)
即得三条线性公式。\(\square\)

Lean 的 \texttt{Mixture.\allowbreak Weights}
只保存权重、非负性和归一性。\texttt{density} 与
\texttt{weighted\allowbreak Instrument}
按上述公式构造，完整性由已有仪器推导；\texttt{instrument}
通过已验证的加权到普通 Kraus
转换取得实际平方根。权重为零时平方根亦为零，仍是合法项，不需要删去某个结果或改变协议身份。运行内核保存有理权重与未缩放矩阵，因此
Hadamard 之类可用有理加权表示，不要求程序把 \(\sqrt2\) 假装成有理数。

\hypertarget{ux603bux53d8ux5deeux7684ux4e00ux534aux6765ux81eaux5f52ux4e00ux5316}{%
\section{51.2
总变差的一半来自归一化}\label{ux603bux53d8ux5deeux7684ux4e00ux534aux6765ux81eaux5f52ux4e00ux5316}}

对同一指标集上的两个分布定义 \[
\operatorname{TV}(p,q)=\frac12\sum_b|p_b-q_b|.
\] 有 \(0\le\operatorname{TV}\le1\)。关键估计适用于所有取值于 \([0,1]\)
的函数 \(f_b\)： \[
\left|\sum_b(p_b-q_b)f_b\right|
=\left|\sum_b(p_b-q_b)(f_b-\tfrac12)\right|
\le\frac12\sum_b|p_b-q_b|.
\] 第一步使用 \(\sum_b(p_b-q_b)=0\)；第二步使用
\(|f_b-1/2|\le1/2\)。如果只逐项估计 \(|f_b|\le1\)，会得到较松的
\(\ell^1\) 界；本章实际形式化了归一化带来的一半因子。

\protect\hypertarget{p-qmx-02}{}{}

\textbf{P-QMX-02（同库权重变化的全称上界）。}
对同一密度态库和同一完整通道库，分别有 \[
d_E\left(\sum_b p_b\rho_b,\sum_b q_b\rho_b\right)
\le\operatorname{TV}(p,q),
\] \[
\Delta_E\left(\sum_b p_b\Phi_b,\sum_b q_b\Phi_b\right)
\le\operatorname{TV}(p,q).
\]

\textbf{证明。} 固定任一效应 \(E\)，令
\(f_b=P_E(\rho_b)\)。第50章已证明每个
\(f_b\in[0,1]\)，所以中心化估计给每个效应的概率差上界。取全部效应上确界得第一式。对第二式，任取输入态
\(\tau\)，将第一式用于密度态库
\(\Phi_b(\tau)\)，再对全部输入态取上确界。\(\square\)

这不是有限状态样本上的最大误差。量词次序是``每个合法输入态、每个效应均受同一预算控制''，所以才能安全取两层上确界。其
Lean 入口是
\texttt{expectation\_\allowbreak sub\_\allowbreak le\_\allowbreak tv}、\allowbreak \texttt{density\_\allowbreak weight\_\allowbreak change\_\allowbreak le\_\allowbreak tv}
与 \texttt{channel\allowbreak Defect\_\allowbreak weight\_\allowbreak change\_\allowbreak le\_\allowbreak tv}。

\hypertarget{ux5e93ux6210ux5458ux4e5fux53d8ux5316ux65f6ux7684ux9884ux7b97ux5206ux89e3}{%
\section{51.3
库成员也变化时的预算分解}\label{ux5e93ux6210ux5458ux4e5fux53d8ux5316ux65f6ux7684ux9884ux7b97ux5206ux89e3}}

\protect\hypertarget{p-qmx-03}{}{}

\textbf{P-QMX-03（权重与成员误差的联合上界）。}
给定两个同载体的完整通道库 \(\Phi_b,\Psi_b\)，记
\(\delta_b=\Delta_E(\Phi_b,\Psi_b)\)。则 \[
\Delta_E\left(\sum_b p_b\Phi_b,\sum_b q_b\Psi_b\right)
\le\operatorname{TV}(p,q)+\sum_b q_b\delta_b.
\] 固定权重 \(p=q\) 时，右侧成为
\(\sum_b p_b\delta_b\)。两库可以有不同的结果和 Kraus
指标类型，因为本式比较的是各自忘记结果后的同载体通道。

\textbf{证明。} 固定输入态、效应，在 \(\sum_b q_b\Phi_b\)
插入中间通道。第一段由 P-QMX-02
控制。第二段展开为成员概率差的非负加权和，应用绝对值三角不等式，再逐项使用成员的真实全称缺陷上界。两层上确界保留这个与输入态、效应无关的右侧。\(\square\)

交换两个模型也能得到权重 \(p_b\)
加权的对称版本；程序或应用可择取更小的已证明预算。无论选择哪个方向，成员上界都需要真实证据。相同的权重并不证明两个不同库等价：单成员分布
\(p=q=(1)\) 的 TV 为零，但恒等通道与 Pauli \(X\) 通道对
\(|0\rangle\langle0|\) 的输出可由 \(|0\rangle\langle0|\)
效应以概率差一分开。这个反例直接说明遗漏``同库''前提会使结论错误。

\hypertarget{ux6709ux7406ux9884ux7b97ux4e0eux957fux671fux8befux5deeux88c5ux914d}{%
\section{51.4
有理预算与长期误差装配}\label{ux6709ux7406ux9884ux7b97ux4e0eux957fux671fux8befux5deeux88c5ux914d}}

\protect\hypertarget{p-qmx-04}{}{}

\textbf{P-QMX-04（精确有理预算的解释）。}
设权重为有理数，仍满足非负、归一条件。定义完全由有限有理算术计算的 \[
b_{\rm TV}=\frac12\sum_b|p_b-q_b|\in\mathbb Q.
\] 其到实数的解释恰等于实数分布的 TV。若
\(b_{\rm TV}\le b\)，则同库通道缺陷
\(\Delta_E\le b\)。更一般地，若另有每个成员的已证明有理上界
\(\delta_b\le\varepsilon_b\)，则 \[
b_{\rm joint}=b_{\rm TV}+\sum_b q_b\varepsilon_b
\] 也是合法的实数解释上界。

\textbf{证明。}
有理到实数的嵌入保持有限和、乘法、减法、绝对值及序。逐项转换后使用
P-QMX-02 或 P-QMX-03 即得结论。\(\square\)

对应
\texttt{Rational\allowbreak Weights.\allowbreak to\allowbreak Weights}、\allowbreak \texttt{tv\_\allowbreak to\allowbreak Weights}、\allowbreak \texttt{channel\allowbreak Defect\_\allowbreak weight\_\allowbreak change\_\allowbreak le\_\allowbreak rational\_\allowbreak budget}
与
\texttt{channel\allowbreak Defect\_\allowbreak joint\_\allowbreak change\_\allowbreak le\_\allowbreak rational\_\allowbreak budget}。这些定理证明数学上的有限有理表达式解释；它们没有自动验证
Python 的解析器、循环、序列索引或外部文件传输。

若第 \(j\) 步各自具有这种预算 \(b_j\)，第50.5节的既有轨迹证书给出 \[
d_E(\rho_n,\sigma_n)\le d_E(\rho_0,\sigma_0)+\sum_{j<n}b_j.
\] 这为``量子实际递推---可计算局部预算---Core
组合证书''接上了一个具体入口。时间依赖的库可逐步声明；本结论仍要求每一步的两模型共享该步所声明的组件对应关系。现有
Lean 轨迹接口在每一串演化内固定有限结果与 Kraus
指标类型，左右两串可不同，并共享同一矩阵载体；各步索引改变时，有限时域须先提供明确的补零或重索引适配。

\hypertarget{ux5c16ux9510ux89c1ux8bc1ux7a0bux5e8fux8d23ux4efbux4e0eux4e0bux4e00ux6b65}{%
\section{51.5
尖锐见证、程序责任与下一步}\label{ux5c16ux9510ux89c1ux8bc1ux7a0bux5e8fux8d23ux4efbux4e0eux4e0bux4e00ux6b65}}

取
\(\Phi_0=\mathrm{id}\)、\(\Phi_1(\rho)=X\rho X\)，\(p=(3/4,1/4)\)、\(q=(1/4,3/4)\)。输入
\(\rho=|0\rangle\langle0|\) 时，输出分别为
\(\operatorname{diag}(3/4,1/4)\) 与
\(\operatorname{diag}(1/4,3/4)\)。测量 \(E=|0\rangle\langle0|\) 得概率差
\(1/2\)，恰好达到 TV 界。若两个组件都改成恒等通道，同样权重仍有预算
\(1/2\)，而实际通道缺陷为零。因此该预算是可证上界，通常不是精确距离算法。

运行入口 \texttt{Mixture\allowbreak Weight\allowbreak Budget}
将两个权重表绑定到同一组已验证完整仪器，预算由输入精确计算，不接受随意填写的更小数字。它记录实际
Kraus
表示指纹；不同但物理等价的表示仍可能有不同指纹。测试同时覆盖负权重、未归一、浮点输入、维数或结果标签失配，以及零权重组件的上下文检查。运行测试提供独立有限见证，通用量词由
Lean 定理负责。

本章依赖明确的量子实现和随机化模型，属于标准物理输入 P、有效模型 E
与领域桥 B 所支持的条件结论；其中可测试的概率约束属依赖分类
O。它不选择复数量子理论，不给任意未知通道的通用计算算法，也不自动控制保留结果的仪器距离或任意旁系统上的
diamond
距离。后续优先把随机分支和控制历史作为实际经典寄存器保留下来，再对完整联合通道使用同类预算。

\hypertarget{ch-52}{%
\chapter{52. 长期平均、紧性与量子平稳态}\label{ch-52}}

UEOT
的``对象持续''需要一个可以核验的动力学含义。给定完整量子通道，最基本的问题是：在状态可能持续变化时，是否存在不变的平均结构？本章在实际有限复矩阵拓扑中证明答案，并把它与逐时收敛、唯一性、条件生存和目标最优性分开。

\hypertarget{ux5b9eux9645ux8fedux4ee3ux548cux5e73ux5747ux6001}{%
\section{52.1
实际迭代和平均态}\label{ux5b9eux9645ux8fedux4ee3ux548cux5e73ux5747ux6001}}

固定同一个完整 Kraus 通道 \(\Phi\) 与初始密度态 \(\rho_0\)，定义 \[
\rho_{n+1}=\Phi(\rho_n),\qquad
\bar\rho_N=\frac1N\sum_{n=0}^{N-1}\rho_n,\qquad N>0.
\] 每一步是合法密度态，平均也由正性、迹线性和 \(N>0\)
得到密度态。平均包含初态，终点为
\(N-1\)；这使下面望远镜恒等式的端点没有歧义。

\protect\hypertarget{p-qce-01}{}{}

\textbf{P-QCE-01（平均态的精确残差与全效应预算）。} 有 \[
\Phi(\bar\rho_N)-\bar\rho_N=\frac{\rho_N-\rho_0}{N},
\qquad
d_E\bigl(\Phi(\bar\rho_N),\bar\rho_N\bigr)\le\frac1N.
\]

\textbf{证明。} 由通道线性， \[
\Phi\left(\sum_{n=0}^{N-1}\rho_n\right)-\sum_{n=0}^{N-1}\rho_n
=\sum_{n=0}^{N-1}(\rho_{n+1}-\rho_n)=\rho_N-\rho_0.
\] 除以 \(N\) 得第一式。任意效应 \(E\) 的两端概率均在
\([0,1]\)，因而其概率残差绝对值至多
\(1/N\)。取全部效应上确界即得第二式。\(\square\)

这不是对某个候选平稳态的距离界；它控制的是``平均态再演化一步''与``平均态自身''的差。没有稳定性常数或谱隙时，小残差不能直接转换为到指定平稳态的距离。\texttt{matrix\_\allowbreak residual}、\allowbreak \texttt{probability\_\allowbreak residual}
与 \texttt{effect\allowbreak Distance\_\allowbreak residual\_\allowbreak le} 形式化的是这个精确合同。

\hypertarget{ux5bc6ux5ea6ux6001ux96c6ux5408ux4e3aux4ec0ux4e48ux7d27}{%
\section{52.2
密度态集合为什么紧}\label{ux5bc6ux5ea6ux6001ux96c6ux5408ux4e3aux4ec0ux4e48ux7d27}}

在 \(d\times d\) 复矩阵空间的通常有限维拓扑中，令 \[
\mathcal D_d=\{\rho:\rho\succeq0,\ \operatorname{Tr}\rho=1\}.
\] 这里 \(\|\cdot\|_{\rm op}\) 表示由复 Euclidean
向量范数诱导的算符范数，不是迹范数。Lean 采用 Mathlib 的
\texttt{Matrix.\allowbreak Norms.\allowbreak L2\allowbreak Operator}
作用域；该结构保持通常矩阵拓扑，不能为了方便而给矩阵人为换成使任意序列收敛的拓扑。

\protect\hypertarget{p-qce-02}{}{}

\textbf{P-QCE-02（实际密度集合的紧性）。} \(\mathcal D_d\)
是闭且有界的集合，故在有限维中紧。每个 \(\rho\in\mathcal D_d\) 满足
\(\|\rho\|_{\rm op}\le1\)，并有 \[
\|\Phi(\bar\rho_N)-\bar\rho_N\|_{\rm op}\le\frac2N.
\]

\textbf{证明。} 正半定锥是闭集；等价地，对任意向量 \(v\)，连续函数
\(\rho\mapsto v^\dagger\rho v\)
的非负条件在极限中保持。迹是连续线性函数，迹一集合也闭，两者交集为闭。有限维
Hermitian 谱定理给非负特征值 \(\lambda_i\)，且
\(\sum_i\lambda_i=1\)；所以每个
\(\lambda_i\le1\)，酉共轭不改变算符范数，得到范数界。有限维闭且有界蕴含紧。最后将
P-QCE-01 的矩阵残差取范数，使用 \(\|\rho_N-\rho_0\|_{\rm op}\le2\)
即得。\(\square\)

这条证明中的紧性、范数界和连续性都在实际矩阵对象上建立，没有把它们写成需要调用者先相信的``持续性证书字段''。Lean
对应
\texttt{density\_\allowbreak operator\allowbreak Norm\_\allowbreak le\_\allowbreak one}、\allowbreak \texttt{density\allowbreak Matrices\_\allowbreak closed}、\allowbreak \texttt{density\allowbreak Matrices\_\allowbreak compact}、\allowbreak \texttt{continuous\_\allowbreak channel}
和 \texttt{matrix\_\allowbreak residual\_\allowbreak norm\_\allowbreak le}。Mathlib 提供谱理论、C*
序与有限维拓扑的通用定理，本项目证明其输入确实由密度态和 Kraus
通道满足。

\hypertarget{ux4e0dux53d8ux5b50ux6781ux9650ux4e0eux5e73ux7a33ux6001ux5b58ux5728}{%
\section{52.3
不变子极限与平稳态存在}\label{ux4e0dux53d8ux5b50ux6781ux9650ux4e0eux5e73ux7a33ux6001ux5b58ux5728}}

\protect\hypertarget{p-qce-03}{}{}

\textbf{P-QCE-03（有限完整通道的平稳密度态）。} 任给初态
\(\rho_0\)，存在密度态 \(\sigma\) 及严格递增的自然数子列
\(N_k\to\infty\)，使 \[
\bar\rho_{N_k}\longrightarrow\sigma,\qquad \Phi(\sigma)=\sigma.
\] 任何收敛的这种平均子列，其极限均为平稳密度态。特别地，当 \(d>0\)
时，每个完整有限 Kraus 通道都至少有一个平稳密度态。

\textbf{证明。} 平均态全在紧集 \(\mathcal D_d\)
内，故存在收敛子列。闭性保证极限正半定且迹一。线性映射 \(\Phi\)
在有限维连续，因而 \[
\Phi(\sigma)-\sigma
=\lim_k\bigl(\Phi(\bar\rho_{N_k})-\bar\rho_{N_k}\bigr)=0,
\] 末项由 P-QCE-02 的 \(2/N_k\)
界趋于零。这个论证适用于任一收敛且指标趋于无穷的子列。若只给通道而没有初态，正维条件提供
\(I/d\) 作为种子，再应用上述结果。\(\square\)

Lean 入口 \texttt{exists\_\allowbreak stationary\_\allowbreak subsequence}
同时返回密度态、平稳性和严格子列；\texttt{subsequence\_\allowbreak limit\_\allowbreak density}
与 \texttt{subsequence\_\allowbreak limit\_\allowbreak invariant}
分别承担极限的物理合法性和不变性；\texttt{exists\_\allowbreak stationary}
用非空有限载体构造归一单位阵。该形式化结果是存在性，不是已实现的平稳态求解程序，也没有给找到足够长子列的通用计算复杂度。

Core 的 P-PER-02/M-OI
与相关持续性章节提供了``平均---残差---不变结构''的通用研究接口。本章在
QM
层对密度集合和实际通道独立完成相应前提与有限矩阵证明；源码没有直接调用
Core 的占据测度定理，因此不把思想上的对应写成代码依赖。第53章的 QSD
桥则会直接消费冻结 Core
定理，两种连接方式在来源上分别记录。{[}S2，§7--8；S3，持续性端口{]}

\hypertarget{ux5468ux671fux8f68ux9053ux4e0eux591aux5e73ux7a33ux6001ux7684ux53cdux4f8b}{%
\section{52.4
周期轨道与多平稳态的反例}\label{ux5468ux671fux8f68ux9053ux4e0eux591aux5e73ux7a33ux6001ux7684ux53cdux4f8b}}

\protect\hypertarget{n-qce-01}{}{}

\textbf{N-QCE-01（平均不变不等于轨道收敛或唯一性）。} P-QCE-01
至03的前提不能推出原始态序列 \(\rho_n\) 收敛，也不能推出平稳态唯一。

\textbf{证明与精确见证。} 取
\(\Phi(\rho)=X\rho X\)、\(\rho_0=|0\rangle\langle0|\)。偶数步为
\(\rho_0\)，奇数步为 \(\rho_1=|1\rangle\langle1|\)，两者由 \(P_0\)
效应以概率差一分开，故原始轨道不收敛。但 \[
\bar\rho_N=\operatorname{diag}\left(\frac{\lceil N/2\rceil}{N},
\frac{\lfloor N/2\rfloor}{N}\right)\longrightarrow I/2.
\] 奇数 \(N\) 时 \(P_0\) 的一步残差恰为 \(1/N\)，偶数 \(N\)
时为零。另取恒等通道，每个密度态都平稳；至少 \(\rho_0\) 与 \(\rho_1\)
不相同，所以唯一性不能成立。\(\square\)

这两类反例有正文代数证明和有理运行见证，不能将其运行检查计入通用 Lean
theorem 数量。实际代码拒绝 \(N=0\)
的平均、非整数步数和非法初态；其平均的起止点与 P-QCE-01 完全一致。

\hypertarget{ux957fux671fux5e73ux5747ux5728ux5b8cux6574ux7269ux7406ux7406ux8bbaux4e2dux7684ux4f4dux7f6e}{%
\section{52.5
长期平均在完整物理理论中的位置}\label{ux957fux671fux5e73ux5747ux5728ux5b8cux6574ux7269ux7406ux7406ux8bbaux4e2dux7684ux4f4dux7f6e}}

通道平稳态、唯一平稳态、所有初态的逐时混合、量子轨迹的几乎处处经验收敛，是不同命题。有限维量子通道文献对遍历性与混合性有专门区分，本章使用与之相容的有限存在性范围；更强的谱结构及均值遍历结果没有由本轮的子列证明自动获得。\href{https://arxiv.org/abs/1210.5625}{Burgarth等原始研究}

本章已经给 UEOT
对象持续提供一种具体且无循环定义的证书：在已声明的物理实现内，存在平均意义不变的合法量子态。这个事实不指定什么对象应该持续，不选择一个价值函数，也不推出控制最优性。要讨论
GOD/GOA，仍须声明动作、代价、约束和评价；要讨论寿命，则需指定``仍然存活''的事件，进入下一章。

下一步可选择实际有限通道，证明其平稳态唯一及额外谱条件，再推导到平稳态的速率；或先形式化整个
Cesàro
序列的投影极限。后一命题在有限维线性理论中有更强的标准工具，本版只将已经实际完成的子列和存在性登记为机器成果。若进入无限维，密度态集合通常不再具有这里使用的紧性，必须加入能量紧性、尾部控制和定义域合同。

\hypertarget{ch-53}{%
\chapter{53. 条件生存、对角量子实现与 Core 谱桥}\label{ch-53}}

``对象仍然存在''只有在指定事件后才能成为概率命题。一个完整量子通道总是保持总概率一，却可以同时包含``继续存活''和``已离开任务域''两个结果。前者的概率一般衰减，按它条件化得到的态也未必等于完整通道的无条件态。本章给出一个可逐项核对的有限实现，把冻结
Core 的次随机与 QSD 定理接到真实量子仪器和记录概率。

\hypertarget{ux7ea6ux5b9aux7ed3ux679cux4e0eux5931ux8d25ux8865ux5168}{%
\section{53.1
约定、结果与失败补全}\label{ux7ea6ux5b9aux7ed3ux679cux4e0eux5931ux8d25ux8865ux5168}}

固定有限状态指标 \(i=1,\ldots,d\)，采用行分布约定 \[
(pQ)_j=\sum_i p_iQ_{ij}.
\] 给定非负次随机矩阵 \[
Q_{ij}\ge0,\qquad s_i:=\sum_jQ_{ij}\le1.
\] 对角嵌入为
\(D(p)=\operatorname{diag}(p_1,\ldots,p_d)\)；非负且和为一的 \(p\)
给出实际复密度态。矩阵单位 \(U_{ji}=|j\rangle\langle i|\) 表示从源 \(i\)
测量并制备到目的 \(j\)，不能将两个指标交换而继续沿用同一 \(Q\)。

\protect\hypertarget{p-qsv-01}{}{}

\textbf{P-QSV-01（次随机核的完整量子仪器）。} 取存活 Kraus 矩阵和失败
Kraus 矩阵 \[
A_{ji}=\sqrt{Q_{ij}}\,|j\rangle\langle i|,
\qquad B_i=\sqrt{1-s_i}\,|i\rangle\langle i|.
\] 二者构成完整仪器，其存活分支为 \[
\mathcal K_Q(X)=\sum_{i,j}Q_{ij}U_{ji}XU_{ji}^{\dagger}.
\]

\textbf{证明。}
\(U_{ji}^{\dagger}U_{ji}=|i\rangle\langle i|\)，所以存活项的 Gram 和为
\(\operatorname{diag}(s_i)\)，失败项的 Gram 和为
\(\operatorname{diag}(1-s_i)\)，二者之和为
\(I\)。非负性及行和上界保证全部权重合法，故两个分支完全正且总通道保持迹。\(\square\)

\texttt{Survival.\allowbreak weighted\allowbreak Instrument} 和 \texttt{instrument}
构造了这组实际仪器。Lean 的两个结果为 \texttt{true} 与
\texttt{false}，运行内核分别命名为 \texttt{survive} 与
\texttt{lost}。Lean 为失败分支也使用完整 \((i,j)\)
指标，仅对角权重非零；Python
省略那些为零的非对角失败项，二者的操作公式相同，表示指纹则不被要求相同。

失败后仍使用同一矩阵载体只是一种完整结果编码；``已死亡''由经典结果记录标记，而不是由
Hilbert
空间中一个未声明的吸收态定义。本章的寿命任务在第一次失败时终止，因此不把失败之后再次调用存活操作解释成复活过程。

\hypertarget{ux540cux4e00ux4e2aux5b58ux6d3bux8bb0ux5f55ux7684ux91cfux5b50ux4e0eux7ecfux5178ux89e3ux91ca}{%
\section{53.2
同一个存活记录的量子与经典解释}\label{ux540cux4e00ux4e2aux5b58ux6d3bux8bb0ux5f55ux7684ux91cfux5b50ux4e0eux7ecfux5178ux89e3ux91ca}}

\protect\hypertarget{p-qsv-02}{}{}

\textbf{P-QSV-02（对角动力学与有限成功词的交织）。} 对任意复算符 \(X\)，
\[
\mathcal K_Q(X)=\operatorname{diag}\left(\sum_iQ_{ij}X_{ii}\right)_j.
\] 因此对任意实向量 \(p\)，任意 \(n\ge0\)， \[
\mathcal K_Q(D(p))=D(pQ),\qquad
\mathcal K_Q^n(D(p))=D(pQ^n).
\] 若 \(p\) 是概率分布，则连续 \(n\) 次记录为存活的概率为
\(pQ^n\mathbf1\)；当该质量为正时，条件态为 \[
D\left(\frac{pQ^n}{pQ^n\mathbf1}\right).
\]

\textbf{证明。} 直接计算
\(U_{ji}XU_{ji}^{\dagger}=X_{ii}|j\rangle\langle j|\)，求和得第一式。代入
\(D(p)\) 得一次交织，随后按记录实际时间顺序归纳，并使用
\((pQ)Q^n=pQ^{n+1}\)。对角态的迹等于分量和，给出记录概率；正质量下按实际归一公式得到条件态。\(\square\)

这里的 \(\mathcal K_Q^n\) 对应既有 \texttt{record\allowbreak Operator} 的长度 \(n\)
全成功词，\texttt{survival\allowbreak Mass\_\allowbreak eq\_\allowbreak response} 将它明确识别为同一操作库的
\texttt{response}，所以不是另造一个与量子记录无关的标量过程。Lean 对所有
\(p,n\) 给出 \texttt{killed\allowbreak Iterate\_\allowbreak diagonal} 和
\texttt{survival\allowbreak Mass\_\allowbreak diagonal}；单步归一的显式矩阵恒等式为
\texttt{success\_\allowbreak posterior\_\allowbreak matrix}。一般多步条件公式按同一归一构造解析推出，后面的
QSD 特例另有完整命名形式定理。

第一式也显示该模型的物理限制：任意输入相干在一次存活测量制备后消失。它是合法的量子实现和经典过程的嵌入，未描述所有非交换存活通道。不能仅因仪器使用复矩阵，就声称已经证明一般量子
Perron 理论。

\hypertarget{qsd-ux7684ux751fux5b58ux5f8bux4e0eux5bffux547d}{%
\section{53.3 QSD
的生存律与寿命}\label{qsd-ux7684ux751fux5b58ux5f8bux4e0eux5bffux547d}}

给定非负归一分布 \(q\)、正数 \(r\)，并明确要求左特征关系 \[
qQ=rq,\qquad q_i\ge0,\quad \sum_iq_i=1,\quad r>0.
\] 这些谱数据是本定理输入；本章没有先证明每个 \(Q\)
都具有所需正谱数据，也没有假设任意初态最终收敛到 \(q\)。此处 \(r\)
是满足该关系的存活特征因子；其全局谱半径身份和主特征值唯一性不属于以下定理的输出。

\protect\hypertarget{p-qsv-03}{}{}

\textbf{P-QSV-03（实际全成功记录的几何生存律）。}
上述前提下，\(0<r\le1\)，且 \[
\mathcal K_Q^n(D(q))=r^nD(q),\qquad
\Pr_q(\text{前 }n\text{ 步全部存活})=r^n.
\] 每个有限 \(n\) 的存活概率为正，按该事件条件化后的实际量子态仍为
\(D(q)\)。

\textbf{证明。} 冻结 Core 的 \texttt{left\allowbreak Perron\_\allowbreak pow} 由左特征关系给
\(qQ^n=r^nq\)。将它经 P-QSV-02
的实际交织送入量子记录，即得第一式。归一化 \(\sum q_i=1\)
后取迹，得到第二式，这一步实际调用 Core 的
\texttt{qsd\_\allowbreak survival\_\allowbreak mass\_\allowbreak pow}。\(r>0\) 保证各有限步质量正；Core 的
\texttt{qsd\_\allowbreak all\_\allowbreak steps} 给归一分布仍为
\(q\)，再经对角嵌入及密度态外延相等得到量子条件态相等。最后 \[
r=\sum_j(qQ)_j=\sum_iq_is_i\le\sum_iq_i=1,
\] 这由 Core 的 \texttt{perron\_\allowbreak root\_\allowbreak le\_\allowbreak one\_\allowbreak of\_\allowbreak substochastic}
正式消费。\(\square\)

本轮 QM 入口为
\texttt{qsd\_\allowbreak killed\allowbreak Iterate}、\allowbreak \texttt{qsd\_\allowbreak survival\_\allowbreak pow}、\allowbreak \texttt{qsd\_\allowbreak survival\_\allowbreak pos}、\allowbreak \texttt{qsd\_\allowbreak conditional\allowbreak State}、\allowbreak \texttt{qsd\_\allowbreak rate\_\allowbreak le\_\allowbreak one}。依赖链先完成量子构造，再实际调用冻结
Core，最后回到真实记录和条件态；量子---经典相等式没有作为结构字段预先给出。

一个直接的解析推论是寿命公式。令 \(\tau\in\{1,2,\ldots\}\cup\{\infty\}\)
表示第一次失败的步数，永不失败时记为无穷，则 \(\Pr(\tau>n)=r^n\)。当
\(r<1\) 时，\(\tau\) 几乎处处有限，并有 \[
\Pr(\tau=m)=(1-r)r^{m-1},\qquad
\mathbb E\tau=\sum_{n=0}^{\infty}\Pr(\tau>n)=\frac1{1-r}.
\] 若每步代表已校准时间 \(\Delta t\)，平均物理寿命为
\(\Delta t/(1-r)\)。本版的 Lean 结论止于有限记录的 \(r^n\)
与条件态；无穷尾和期望及时间标定在这里分别作为解析推论和物理输入，未冒称新增机器证明。\(r=1\)
的理想模型则无有限步失败。

\hypertarget{doob-ux5b9eux73b0ux53caux5176ux4e0dux540cux7684ux4e0dux53d8ux5206ux5e03}{%
\section{53.4 Doob
实现及其不同的不变分布}\label{doob-ux5b9eux73b0ux53caux5176ux4e0dux540cux7684ux4e0dux53d8ux5206ux5e03}}

进一步给严格正的右特征向量 \(h\)，满足 \[
Qh=rh,\qquad h_i>0,\qquad \sum_iq_ih_i=1.
\] 定义 \[
\widetilde Q_{ij}=\frac{Q_{ij}h_j}{rh_i},\qquad \pi_i=q_ih_i.
\]

\protect\hypertarget{p-qsv-04}{}{}

\textbf{P-QSV-04（Doob 核的实际完整实现与不变成功记录）。}
\(\widetilde Q\) 非负且每行和为一，\(\pi\) 是概率分布并满足
\(\pi\widetilde Q=\pi\)。将 \(\widetilde Q\) 按 P-QSV-01
实现，得到实际量子仪器：任一归一对角初态的一步成功概率为一；以
\(D(\pi)\) 为初态，成功分支及任意有限全成功词均保持
\(D(\pi)\)，其记录概率为一，条件态亦为 \(D(\pi)\)。

\textbf{证明。} 分母严格正给非负性，右特征关系给 \[
\sum_j\widetilde Q_{ij}=\frac{\sum_jQ_{ij}h_j}{rh_i}=1.
\] 归一条件保证 \(\sum_i\pi_i=1\)。左特征关系给 \[
(\pi\widetilde Q)_j
=\sum_iq_ih_i\frac{Q_{ij}h_j}{rh_i}
=\frac{h_j}{r}(qQ)_j=q_jh_j=\pi_j.
\] 对这些 Core 已证的行和、不变分布归一及平稳关系，分别通过 P-QSV-01
和02得到实际仪器、概率与成功分支恒等式；按记录长度归纳得到所有有限成功词，再对质量一归一即可。\(\square\)

Lean 实际消费
\texttt{doob\allowbreak Transform\_\allowbreak row\_\allowbreak sum}、\allowbreak \texttt{doob\allowbreak Invariant\_\allowbreak sum} 和
\texttt{doob\allowbreak Invariant\_\allowbreak stationary}，对应本层
\texttt{doob\allowbreak Instrument}、\allowbreak \texttt{doob\allowbreak Density}、\allowbreak \texttt{doob\_\allowbreak success\_\allowbreak probability}、\allowbreak \texttt{doob\_\allowbreak success\_\allowbreak branch}、\allowbreak \texttt{doob\_\allowbreak killed\allowbreak Iterate}、\allowbreak \texttt{doob\_\allowbreak survival\allowbreak Mass}
与
\texttt{doob\_\allowbreak conditional\allowbreak State}。这些名称准确覆盖分支、记录和条件态；本版没有另外登记忘记结果后的
\texttt{channel\allowbreak State} 平稳推论为新的命名机器定理。

Doob
变换常用于研究条件轨迹或稀有事件的替代动力学；一般量子版本还涉及非交换谱数据和适当的轨迹极限。\href{https://arxiv.org/abs/1711.10951}{Carollo等原始研究}
本章只证明给定谱前提下的有限构造和上述不变关系，尚未建立``先条件存活到遥远未来再取极限''与本变换相等的过程定理。尤其不能把
\(q\) 的有限时存活条件律与 \(\pi=qh\) 的 Doob 平稳律直接互换。

\hypertarget{ux533aux5206-qsddoob-ux5f8bux4e0eux521dux6001ux524dux63d0ux7684ux7cbeux786eux6a21ux578b}{%
\section{53.5 区分 QSD、Doob
律与初态前提的精确模型}\label{ux533aux5206-qsddoob-ux5f8bux4e0eux521dux6001ux524dux63d0ux7684ux7cbeux786eux6a21ux578b}}

取 \[
Q=\begin{pmatrix}1/4&1/2\\1/8&1/4\end{pmatrix},\quad
q=(1/3,2/3),\quad r=1/2,\quad h=(3/2,3/4).
\] 两行存活概率为 \(3/4\) 与 \(3/8\)，因此损失不均匀。直接算得 \[
qQ=(1/6,1/3)=rq,\qquad
Qh=(3/4,3/8)=rh,\qquad \sum_iq_ih_i=1.
\] 所以从 \(D(q)\) 开始，第 \(n\) 次全成功记录概率为
\(2^{-n}\)，每次条件态仍为
\(D(q)\)；每步无条件态则还包括失败分支的贡献。

\protect\hypertarget{n-qsv-01}{}{}

\textbf{N-QSV-01（QSD、Doob 不变律与任意初态不可混同）。} 在本例中， \[
\widetilde Q=\begin{pmatrix}1/2&1/2\\1/2&1/2\end{pmatrix},\qquad
\pi=(1/2,1/2)\ne q.
\] 从均匀初态 \(p=(1/2,1/2)\) 出发的第一次存活概率为
\(9/16\ne r\)；若错误地将 \(Q\) 转置而保持行分布解释，从 \(q\)
出发的第一次存活概率成为 \(5/8\ne r\)。

\textbf{证明。} 逐项代入 Doob 公式得到上述两行；乘积 \(q_ih_i\) 各为
\(1/2\)。另有 \(pQ=(3/16,3/8)\)，其和为
\(9/16\)；\(qQ^T=(5/12,5/24)\)，其和为
\(5/8\)。因此三种替换都改变了待证命题或实际概率。\(\square\)

运行见证按有理矩阵逐次计算前六个成功词，得到 \(1/2,1/4,\ldots,1/64\)
和相同条件态；同时检查 Doob 的 \(\pi\) 固定、\(q\)
不固定。该模型还给平均失败步数二的解析值，但运行的六步检查不被当作无穷期望的证明。负控包括行和超过一、负项、非方阵和不精确输入；删除失败分支会破坏完整性。取
\(Q=0\) 时存活概率为零，运行入口返回后验未定义，不能人为指定 \(q\)。

\hypertarget{ux5bf9-ueot-ux6838ux5fc3ux8ffdux6c42ux7684ux5b9eux9645ux8d21ux732e}{%
\section{53.6 对 UEOT
核心追求的实际贡献}\label{ux5bf9-ueot-ux6838ux5fc3ux8ffdux6c42ux7684ux5b9eux9645ux8d21ux732e}}

本章把``持续''分成三个可以各自核验的问题：完整操作的概率守恒；指定存活事件的概率及条件态；给定左右谱数据构造的替代动力学不变结构。它让
UEOT 对象持续的抽象语言具有实际 quantum record
的可计算内容，且能指出缺少哪个前提时结论失败。

分类上，冻结 QSD 代数属于 C，复密度态和 Kraus 仪器属于标准物理输入
P，选择对角测量制备和终止事件属于有效模型 E，实际交织属于桥
B，寿命与记录概率约束属于 O。一般非交换 Perron
存在性、不可约/本原条件、条件极限和谱隙在本项目仍有标准数学、领域桥或具体模型义务；``尚未形式化''本身不使它们成为科学猜想
X。只有额外提出且未证的 UEOT 机制必然性或物理选择关系，才另按 X
登记。原始双驱动哲学可指导选择哪些持续机制值得研究，却不从这些有限定理直接推出目标函数或自然最优化。{[}S1；S2，QSD/Perron
与 GOA 各节；S4，持续性和评价的分型{]}

第56章现已给出保留相干的有限 CP 存活映射、明确左右谱对象、实际非交换
Doob
实现与任意初态的有限条件化连接。下一步继续分别证明适用条件下的谱存在、唯一、条件极限和收敛率。与第51章的扰动预算结合后，才讨论当制备、转移或存活判定具有误差时，寿命与条件预测能否稳定转移。这种推进能积累跨模型可复用的物理合同，同时保持``UEOT
何以选择量子结构''这一根本问题的独立研究位置。

\hypertarget{ch-54}{%
\chapter{54. 同一 Core 过程上的经典纯化边界}\label{ch-54}}

第1、2章区分了对象数学与量子扇区选择。一个通用对象接口能够被量子系统实现，并不等于只有量子系统能够实现它。本章把这项责任落实到冻结
Core
的实际有限随机核、终端稳定商、控制商和无限轨迹。所用经典解释公开给出；纯化的缺失由完整状态空间中的凸极点定义证明，且发生在同一轨迹的实际商边际上。

\hypertarget{ux5b8cux6574ux7ecfux5178ux72b6ux6001ux7a7aux95f4ux53caux5176ux64cdux4f5cux89e3ux91ca}{%
\section{54.1
完整经典状态空间及其操作解释}\label{ux5b8cux6574ux7ecfux5178ux72b6ux6001ux7a7aux95f4ux53caux5176ux64cdux4f5cux89e3ux91ca}}

对有限载体 \(Z\)，定义 \[
\Delta(Z)=\{p:Z\to\mathbb R\mid p_z\ge0,\ \sum_zp_z=1\}.
\] 允许态是整个单纯形，效应是所有 \(e:Z\to[0,1]\)，概率配对为 \[
P_e(p)=\sum_zp_ze_z.
\] 复合系统采用 \(\Delta(Z\times B)\)，丢弃右系统为实际求和边际 \[
(\operatorname{marg}_Z\omega)_z=\sum_b\omega_{zb}.
\] 这些选择构成一个明确的经典领域解释 B。本章列明的 H1--H6
接口组没有替我们选择或排除这个复合规则，不能将它们写成 Core
的必然物理结论。

\protect\hypertarget{p-qcl-01}{}{}

\textbf{P-QCL-01（完整单纯形的状态、响应与运输）。} 上述配对在 \([0,1]\)
内，全部效应分离状态。Dirac
制备、合法凸组合、任意有限函数的推送以及上述边际均保持归一态；推送与二态等权混合相容，边际就是投影推送。

\textbf{证明。} 非负性逐项给 \(P_e(p)\ge0\)；由 \(e_z\le1\) 得
\(P_e(p)\le\sum_zp_z=1\)。若两态对全部效应响应相同，取点指示效应
\(e^{(x)}_z=\mathbf1_{z=x}\) 即得每个坐标相等。对函数 \(f:Z\to W\)， \[
(f_*p)_w=\sum_{z:f(z)=w}p_z.
\]
有限和换序给归一性，每个源点恰在一个纤维中。相同换序给所用的等权混合相容性与边际公式。\(\square\)

Lean 的 \texttt{Classical.\allowbreak State} 复用 \texttt{Mixture.\allowbreak Weights}
的三个纯标量字段，没有把复密度矩阵或 Born
结构作为新经典定理的假设。其导入闭包含量子模块，是软件依赖事实；本章的新状态定义及证明前提仍是有限非负实数和。\texttt{all\_\allowbreak effects\_\allowbreak separate}
量化整个 \texttt{Effect}
类型，不能缩减成几次有限样本测量就沿用同一结论。

对 Core 的任意实际有限随机核 \(K\)、任意 \(p\in\Delta(Z)\) 和固定动作
\(a\)，\texttt{step\allowbreak State} 还定义完整单纯形上的真实随机作用 \[
(T_ap)_y=\sum_xp_xK(x,a,y).
\]
非负性逐项成立，归一性由核的每行和为一及有限和换序给出。\texttt{to\allowbreak PMF}
采用 Core 已有的实概率行构造；\texttt{to\allowbreak PMF\_\allowbreak step\allowbreak State} 证明 \[
\operatorname{toPMF}(T_ap)
=\operatorname{toPMF}(p)\mathbin{\mathrm{bind}}
\bigl(x\mapsto\operatorname{toPMF}(K(x,a,\cdot))\bigr).
\] 因此经典状态作用与实际 PMF
转移在任意输入态上相符。后文的路径连接再为同一模型选择明确策略；这一固定动作公式本身不承担任意历史反馈的形式化责任。

称 \(p\) 为纯态，当且仅当它是整个允许态空间的凸极点：对任何
\(q,r\in\Delta(Z)\) 以及 \(0<t<1\)，\(p=tq+(1-t)r\) 蕴涵
\(q=r=p\)。\texttt{Is\allowbreak Pure} 正是这个量词定义。Dirac
态的纯性另行证明：在其零坐标处，两个非负加权项之和为零，故两项都为零；归一性迫使剩余坐标都为一。没有先把``纯''定义成
Dirac，再把同义改写当作分类成果。

\hypertarget{ux5bf9ux4efbux610fux6709ux9650ux8f85ux52a9ux7cfbux7edfux7684ux7eafux5316ux5426ux5b9a}{%
\section{54.2
对任意有限辅助系统的纯化否定}\label{ux5bf9ux4efbux610fux6709ux9650ux8f85ux52a9ux7cfbux7edfux7684ux7eafux5316ux5426ux5b9a}}

\protect\hypertarget{n-qcl-01}{}{}

\textbf{N-QCL-01（二点混态没有有限经典纯化）。} 设
\(z_0\ne z_1\)，\(p=\tfrac12\delta_{z_0}+\tfrac12\delta_{z_1}\)。对任意有限类型
\(B\)，不存在 \(\omega\in\Delta(Z\times B)\) 同时满足凸极点纯性及
\(\operatorname{marg}_Z\omega=p\)。

\textbf{证明。} 假定这样的联合态存在。边际在 \(z_0,z_1\) 各有质量
\(1/2\)，其他行质量为零。非负性使其他行的每个条目为零。定义 \[
\omega_i(z,b)=2\mathbf1_{z=z_i}\omega(z,b),\quad i=0,1.
\] 两者非负；第 \(z_i\) 行的质量 \(1/2\) 给
\(\sum_{z,b}\omega_i(z,b)=1\)。逐坐标有 \[
\omega=\tfrac12\omega_0+\tfrac12\omega_1,
\qquad \operatorname{marg}_Z\omega_i=\delta_{z_i}.
\] 极点纯性迫使 \(\omega_0=\omega\)，取边际在 \(z_1\) 的坐标则得到
\(0=1/2\)，矛盾。\(\square\)

这项证明不需要 \(B\) 非空；如果 \(B\) 为空，本来就不存在归一联合态。Lean
的 \texttt{no\_\allowbreak pure\_\allowbreak extension\_\allowbreak half\allowbreak Mix}
对任意有限辅助类型量化，没有只测试一个预选维数。否定的是纯化的存在部分，因此无需先形式化纯化的可逆唯一性。

删除纯性要求会得到完全不同的结论。给定辅助点 \(b\)，\(p\otimes\delta_b\)
是任意 \(p\) 的普通联合扩展；对 Dirac 制备，\(\delta_{(z,b)}\)
则是真正纯的联合扩展。\texttt{marginal\allowbreak Left\_\allowbreak attach} 与
\texttt{dirac\_\allowbreak has\_\allowbreak pure\_\allowbreak extension}
分别验证这两件事。父载体存在、记录扩展存在和纯联合态存在具有不同的数学要求。

\hypertarget{ux4fddux7559ux64cdux4f5cux6df7ux6001ux7684ux5b9eux9645ux7ec8ux7aefux5546}{%
\section{54.3
保留操作混态的实际终端商}\label{ux4fddux7559ux64cdux4f5cux6df7ux6001ux7684ux5b9eux9645ux7ec8ux7aefux5546}}

采用冻结 Core 的同一核 \texttt{survival\allowbreak Micro\allowbreak Kernel}。状态为
\(x=(s,h)\in\mathrm{Bool}^2\)，动作为 \(a\in\mathrm{Bool}\)， \[
K(x,a,z)=\begin{cases}1/2,&z_1=a,\\0,&z_1\ne a.\end{cases}
\] 记
\(x_0=(\mathrm{true},\mathrm{false})\)、\(x_1=(\mathrm{true},\mathrm{true})\)。第一位作为注册安全谓词。新实例
\texttt{fully\allowbreak Observed\allowbreak Survival\allowbreak Model} 明确注册完整读出
\(x\mapsto x\)，同时注册安全位、平凡程序／资源字段和零奖励。完整读出是本模型的
B 解释，尚不代表自然界中任意标签都能被测量。

\protect\hypertarget{p-qcl-02}{}{}

\textbf{P-QCL-02（实际终端商保留二点支持）。} 该模型的实际
\texttt{terminal\allowbreak Setoid} 满足 \(x\sim y\iff x=y\)。由
\texttt{organizational\allowbreak Terminal\allowbreak Control} 实际构造的控制商映射 \(f\)
因而单射，原核的第一行 \[
p_* = K((\mathrm{true},\mathrm{false}),\mathrm{true},\cdot)
=\tfrac12\delta_{x_0}+\tfrac12\delta_{x_1}
\] 被推送为
\(\bar p=\tfrac12\delta_{f(x_0)}+\tfrac12\delta_{f(x_1)}\)，其中
\(f(x_0)\ne f(x_1)\)。

\textbf{证明。}
终端等价保持已注册输出；其第一项是完整源态，所以等价蕴涵源态相等。反向由等价关系自反性成立。若实际控制商的两个像相等，取其商类分量，使用商相等的消去规则得到终端等价，进而得到源态相等。原核行的二点公式由四个条目直接计算，再用
P-QCL-01 的等权混合相容性和 Dirac 推送公式。\(\square\)

这里没有手写恒等商代替算法。商载体是该实际稳定关系的
\texttt{Reachable\allowbreak State}，映射是实际
\texttt{terminal\allowbreak Control.\allowbreak f}。完整观察使本例没有严格的状态数压缩；这仍是算法的合法输出。压缩率和操作上是否保留相关混合是两个不同问题。

\hypertarget{ux540cux4e00ux63a7ux5236ux5546ux771fux5b9eux8defux5f84ux4e0eux7b2cux4e00ux65f6ux523bux5546ux8fb9ux9645}{%
\section{54.4
同一控制商、真实路径与第一时刻商边际}\label{ux540cux4e00ux63a7ux5236ux5546ux771fux5b9eux8defux5f84ux4e0eux7b2cux4e00ux65f6ux523bux5546ux8fb9ux9645}}

固定初态 \(\delta_{x_0}\) 及明确的常真策略
\(\pi_*(x)=\mathrm{true}\)。转移 PMF 直接来自同一个控制商的
\texttt{micro.\allowbreak transition\allowbreak PMF}。令 \(\mathbb P_*\) 为其实际
\texttt{stationary\allowbreak Traj\allowbreak Measure}，没有另造一条恰好满足期望结论的路径。

\protect\hypertarget{p-qcl-03}{}{}

\textbf{P-QCL-03（同源安全轨迹与纯化分离证书）。}
本模型满足下表列明的有限 Core 接口组；并有 \[
\mathbb P_*\{\omega:\forall n\in\mathbb N,\ (\omega_n)_1=\mathrm{true}\}=1,
\] \[
(\omega\mapsto f(\omega_1))_*\mathbb P_*
=\operatorname{Measure}(\operatorname{toPMF}\bar p).
\] 其右侧状态 \(\bar p\) 对任意有限经典辅助系统均无纯化。

\begin{longtable}[]{@{}
  >{\raggedright\arraybackslash}p{(\columnwidth - 2\tabcolsep) * \real{0.3000}}
  >{\raggedright\arraybackslash}p{(\columnwidth - 2\tabcolsep) * \real{0.7000}}@{}}
\toprule\noalign{}
\begin{minipage}[b]{\linewidth}\raggedright
接口
\end{minipage} & \begin{minipage}[b]{\linewidth}\raggedright
本例的实际责任
\end{minipage} \\
\midrule\noalign{}
\endhead
\bottomrule\noalign{}
\endlastfoot
H1 正规化受控过程 & 同一 \(K\) 的每行非负且和为一 \\
H2 注册接口 & transition 恰为 \(K\)，输出恰为完整读出与已声明属性 \\
H3 终端稳定商 & 实际 \texttt{terminal\allowbreak Setoid} 稳定，且终端等价恰为相等 \\
H4 精确控制 & 实际控制商、折扣 \(1/2\)、零奖励；微观与宏观最优值经 \(f\)
运输 \\
H5 非空安全性 & \(x_0\) 在每个有限时域的实际 viability 集内，初始 PMF
支持满足相应条件 \\
H6 实际安全路径 & 明确常真策略、实际 PMF
源等式、全部离散时刻安全与上述真实测度边际 \\
\end{longtable}

\textbf{证明。}
核的正规化来自冻结定义。注册保持、终端稳定与最优值运输使用冻结接口的实际定理。对
viability 按时域归纳：选动作 true 后所有正概率目标的第一位都为
true，故递归安全。对于路径，初始分布在安全集合内，常真策略每次都把其支持送回该集合。实际坐标分布定理给每个时刻安全事件的概率一，随后
\texttt{all\_\allowbreak times\_\allowbreak safe\_\allowbreak of\_\allowbreak coordinate\_\allowbreak one}
给可数全部时刻的交集概率一。

第一时刻边际不是从``存在某个安全策略''推出的。先使用实际坐标定理，把
\(\omega_1\) 的分布化为从 \(\delta_{x_0}\) 经常真策略转移一次的 PMF，即
\(K(x_0,\mathrm{true},\cdot)\)。再以可测 \(f\)
推送，利用两次测度推送的合成和
\texttt{to\allowbreak PMF\_\allowbreak pushforward}，得到显示的测度等式。最后以 P-QCL-02
的两个不同商点实例化 N-QCL-01。\(\square\)

\texttt{Finite\allowbreak Core\allowbreak Slice}
把上述具体事实写成合取命题，\texttt{finite\allowbreak Core\allowbreak Slice\_\allowbreak holds}
从实际定理逐项证明它。\texttt{classical\_\allowbreak purification\_\allowbreak separation}
再连接真实路径边际与任意有限辅助空间上的否定量词。证书没有把纯化失败或边际相等作为构造对象时的输入。

零奖励是合法控制实例，最优值运输在这里不选择非平凡偏好。动作仍有真实安全差异：同一
\(K\) 在动作 false 下将第一位置为
false，形式化见证明确给一个不安全目标的概率
\(1/2\)。全部时刻安全表示这个注册谓词的概率保持，不等于已证明物质修复、自生产、谱系持续或内生目标选择。{[}S2、S4；QM316{]}

\hypertarget{ux5220ux9664ux8bfbux51faux540eux53cdux4f8bux4e3aux4ec0ux4e48ux5fc5ux987bux5931ux6548}{%
\section{54.5
删除读出后，反例为什么必须失效}\label{ux5220ux9664ux8bfbux51faux540eux53cdux4f8bux4e3aux4ec0ux4e48ux5fc5ux987bux5931ux6548}}

\protect\hypertarget{n-qcl-02}{}{}

\textbf{N-QCL-02（微观混合可被实际操作商消去）。} 保留完全同一个
\(K\)，改用既有只读第一位的
\texttt{registered\allowbreak Survival\allowbreak Model}，则其实际终端控制商合并
\(x_0,x_1\)。因此 \(p_*\) 的实际商推送为一个 Dirac
态，且该商态由凸极点定义确为纯态。

\textbf{证明。} 冻结定理
\texttt{registered\_\allowbreak survival\_\allowbreak two\_\allowbreak hidden\_\allowbreak states\_\allowbreak merge}
给两个点的实际终端等价，商引入规则给其实际控制商像相等。于是 \[
f_*p_* = \tfrac12\delta_{f(x_0)}+\tfrac12\delta_{f(x_0)}
=\delta_{f(x_0)}.
\] 再用 Dirac 纯性的独立证明。\(\square\)

这项删除见证表明，操作对象的身份会改变可见混合。微观第二位仍随机，但被指定读出与实际商完全识别；若继续把它称为操作混态，纯化论证的对象就错了。反之，主模型通过完整注册读出明确保留该差别，并承担这个领域解释的责任。受限读出下这一特定混态反例失效，并不证明受限模型的一切状态都满足普遍纯化。

\hypertarget{ux5df2ux95edux5408ux7684ux8d23ux4efbux4e0eux4e0bux4e00ux9879ux91cfux5b50ux9009ux62e9ux4efbux52a1}{%
\section{54.6
已闭合的责任与下一项量子选择任务}\label{ux5df2ux95edux5408ux7684ux8d23ux4efbux4e0eux4e0bux4e00ux9879ux91cfux5b50ux9009ux62e9ux4efbux52a1}}

\begin{longtable}[]{@{}
  >{\raggedright\arraybackslash}p{(\columnwidth - 2\tabcolsep) * \real{0.3000}}
  >{\raggedright\arraybackslash}p{(\columnwidth - 2\tabcolsep) * \real{0.7000}}@{}}
\toprule\noalign{}
\begin{minipage}[b]{\linewidth}\raggedright
正文结论
\end{minipage} & \begin{minipage}[b]{\linewidth}\raggedright
主要 Lean 入口
\end{minipage} \\
\midrule\noalign{}
\endhead
\bottomrule\noalign{}
\endlastfoot
全部效应、推送与边际 &
\texttt{all\_\allowbreak effects\_\allowbreak separate}、\allowbreak \texttt{pushforward\_\allowbreak half\allowbreak Mix}、\allowbreak \texttt{marginal\allowbreak Left\_\allowbreak eq\_\allowbreak pushforward} \\
任意态的同源随机作用 &
\texttt{step\allowbreak State}、\allowbreak \texttt{to\allowbreak PMF\_\allowbreak step\allowbreak State}、\allowbreak \texttt{to\allowbreak PMF\_\allowbreak pushforward} \\
极点与任意有限辅助的否定 &
\texttt{dirac\_\allowbreak is\allowbreak Pure}、\allowbreak \texttt{no\_\allowbreak pure\_\allowbreak extension\_\allowbreak half\allowbreak Mix} \\
实际终端商与非混同 &
\texttt{terminal\_\allowbreak eq\_\allowbreak iff}、\allowbreak \texttt{terminal\allowbreak Control\_\allowbreak injective}、\allowbreak \texttt{first\allowbreak Quotient\allowbreak State\_\allowbreak half\allowbreak Mix} \\
同源真实路径 &
\texttt{actual\allowbreak Path\_\allowbreak all\_\allowbreak times\_\allowbreak safe}、\allowbreak \texttt{actual\allowbreak Path\_\allowbreak first\_\allowbreak quotient\_\allowbreak marginal} \\
接口合取与分离 &
\texttt{finite\allowbreak Core\allowbreak Slice\_\allowbreak holds}、\allowbreak \texttt{classical\_\allowbreak purification\_\allowbreak separation} \\
删除读出 &
\texttt{restricted\allowbreak First\allowbreak Quotient\allowbreak State\_\allowbreak dirac}、\allowbreak \texttt{restricted\allowbreak First\allowbreak Quotient\allowbreak State\_\allowbreak pure} \\
\end{longtable}

本章收紧的结论是：表中这组有限 Core
接口与公开的完整经典解释相容，而且同一可达操作混态违反纯化存在性。因此，若要从这些接口进一步得到普遍纯化，必须增加能够排除此实例的独立原则或证明额外的真实物理桥。第1章的整个重构问题没有被一个有限证书替代；本章不构成全部106条及所有物理解释的模型，也不否定更丰富的
UEOT 理论最终重构 QM 的可能。

第57章现在把候选原则的系统类型、完整状态集合、复合和丢弃写成实际接口，并证明经典反例经适配后仍违反纯化存在。复量子模型的任意有限态纯化也在完整密度空间中得到证明。第55章实扇区的排除责任则仍属于局部层析；纯化存在与本质唯一性分别维护。继续推进时，不能把预先装入的复Hilbert、Born或Kraus结构误称为已经从Core推出的起源。

\hypertarget{ch-55}{%
\chapter{55. 实量子扇区的局域层析分离}\label{ch-55}}

纯化与局域层析承担不同的选择责任。第54章否定一个完整经典解释的纯化存在性；本章固定通常的实态与实效应限制，直接构造两个联合态：任何实局域乘积测试都看不出差别，一个合法实联合测试却能完全区分它们。这里使用具体矩阵与全部效应量词，不把一次测量失败误称为所有局域测量失败。

两 rebit 的实组合结构与 \(Y\otimes Y\) 额外方向已有明确研究先例，见
Caves、Fuchs、Rungta
原文第4--5页及式(1.6)--(1.12)。本章的工作是把所用态对、全部效应量词和删除见证接入当前形式化内核，不把这个标准数学现象登记为新的自然规律。\href{https://arxiv.org/abs/quant-ph/0009063}{原文}

\hypertarget{ux5b9eux6027ux662fux989dux5916ux9650ux5236ux8054ux5408ux7cfbux7edfux5fc5ux987bux72ecux7acbux58f0ux660e}{%
\section{55.1
实性是额外限制，联合系统必须独立声明}\label{ux5b9eux6027ux662fux989dux5916ux9650ux5236ux8054ux5408ux7cfbux7edfux5fc5ux987bux72ecux7acbux58f0ux660e}}

在既有有限复矩阵实现中，称矩阵实，指其每个条目的虚部为零。局域密度态与效应仍须满足原有半正定、迹或单位上界条件，再加实性。实厄米矩阵就是实对称矩阵。两局域系统以四维实
Hilbert 空间为联合载体，用实对称 \(4\times4\) 密度矩阵表示联合态。

固定 \[
Y=\begin{pmatrix}0&-i\\i&0\end{pmatrix},\qquad H=Y\otimes Y.
\] 按 \((00,01,10,11)\) 的基顺序， \[
H=\begin{pmatrix}
0&0&0&-1\\
0&0&1&0\\
0&1&0&0\\
-1&0&0&0
\end{pmatrix}.
\] 单个 \(Y\) 厄米但不实，而 \(Y\otimes Y\) 是实对称矩阵。不能因为局域
\(Y\)
不允许，就把它的张量积自动从完整实联合系统中删掉；这正是需明确声明组合规则的地方。

定义 \[
P_\pm=\frac{I\pm H}{2},\qquad
\rho_\pm=\frac{P_\pm}{2}=\frac{I\pm H}{4}.
\]

\protect\hypertarget{p-qrt-01}{}{}

\textbf{P-QRT-01（两个合法且不同的实联合态）。} \(\rho_+,\rho_-\)
都是实、厄米、半正定且迹为一的密度矩阵，并且 \(\rho_+\ne\rho_-\)。

\textbf{证明。} \(Y^\dagger=Y\)、\(Y^2=I\)，故
\(H^\dagger=H\)、\(H^2=I\)。于是 \[
P_\pm^2=P_\pm,\quad P_+P_-=0,\quad P_++P_-=I.
\] 厄米幂等矩阵满足 \(P=P^\dagger P\)，因而半正定。由
\(\operatorname{Tr}H=0\)、联合维数为4，得到
\(\operatorname{Tr}P_\pm=2\)，再除以2即有迹一。实性由上面的具体 \(H\)
条目成立。两态的差为 \(H/2\)，其 \((00,11)\) 条目为
\(-1/2\)，所以不同。\(\square\)

Lean 的 \texttt{involution\allowbreak Effect}
从一个厄米对合矩阵构造实际效应；正性通过真实 Gram
表达式证明，单位上界通过互补投影证明。\texttt{rho\allowbreak Plus} 与
\texttt{rho\allowbreak Minus} 是实际 \texttt{Density}
对象，而非只有预期矩阵、尚未证明合法性的符号标签。

\hypertarget{ux5bf9ux6240ux6709ux5b9eux5c40ux57dfux4e58ux79efux6548ux5e94ux7edfux8ba1ux5747ux76f8ux540c}{%
\section{55.2
对所有实局域乘积效应，统计均相同}\label{ux5bf9ux6240ux6709ux5b9eux5c40ux57dfux4e58ux79efux6548ux5e94ux7edfux8ba1ux5747ux76f8ux540c}}

任意局域效应 \(E,F\) 满足
\(0\le E,F\le I\)。其张量积也是合法联合效应：张量积保持半正定，且 \[
I-E\otimes F=(I-E)\otimes I+E\otimes(I-F)\ge0.
\] \texttt{product\allowbreak Effect}
实际证明了这个上界，未把局域产品的合法性当作额外未核对输入。

实对称局域矩阵具有 \[
E=\begin{pmatrix}a&b\\b&d\end{pmatrix},
\qquad \operatorname{Tr}(EY)=ib-ib=0.
\] 这个等式对所有实厄米局域矩阵成立，不仅对某几个常用投影成立。

\protect\hypertarget{n-qrt-01}{}{}

\textbf{N-QRT-01（实局域乘积测试不能确定联合态）。} 对任意两个实局域效应
\(E,F\)， \[
\operatorname{Tr}[(E\otimes F)\rho_+]
=\operatorname{Tr}[(E\otimes F)\rho_-].
\] 尽管如此，P-QRT-01 给出的两个联合态不同。

\textbf{证明。} 有限 Kronecker 积的乘法与迹公式给 \[
\begin{aligned}
\operatorname{Tr}[(E\otimes F)(\rho_+-\rho_-)]
&=\tfrac12\operatorname{Tr}[(E\otimes F)(Y\otimes Y)]\\
&=\tfrac12\operatorname{Tr}(EY)\operatorname{Tr}(FY)=0.
\end{aligned}
\] 因此两态对这些测试有相同 Born 响应。共同值解析上可写为
\(\tfrac14\operatorname{Tr}E\operatorname{Tr}F\)。\(\square\)

\texttt{all\_\allowbreak real\_\allowbreak product\_\allowbreak probabilities\_\allowbreak equal} 的量词覆盖全部实际
\texttt{Effect}，再以 \texttt{Is\allowbreak Real}
限定允许的局域集合。\texttt{real\_\allowbreak local\_\allowbreak tomography\_\allowbreak fails}
将全称相等与态不等合并。这里的``局域层析''指联合态由所有局域乘积响应确定的性质；本版没有把它扩张成任意未建模的通信、适应协议或外部辅助资源的分类定理。

一个解析的维数解释是：单个二维实对称矩阵空间有维数3，两个局域空间的乘积线性张成有维数9，而四维实对称联合矩阵空间有维数10。额外方向可以取
\(Y\otimes Y\)。本版的主要 Lean
结论直接证明实际全称响应等式，不以未形式化的维数计数代替它。

\hypertarget{ux5408ux6cd5ux5b9eux8054ux5408ux6d4bux8bd5ux80fdux591fux5b8cux5168ux533aux5206}{%
\section{55.3
合法实联合测试能够完全区分}\label{ux5408ux6cd5ux5b9eux8054ux5408ux6d4bux8bd5ux80fdux591fux5b8cux5168ux533aux5206}}

\protect\hypertarget{p-qrt-02}{}{}

\textbf{P-QRT-02（联合效应的确定性区分）。} \(P_+\)
是合法的实联合效应，且 \[
\operatorname{Tr}(P_+\rho_+)=1,\qquad
\operatorname{Tr}(P_+\rho_-)=0.
\]

\textbf{证明。} P-QRT-01 的投影正性给 \(P_+\ge0\)，而
\(I-P_+=P_-\ge0\)，所以它满足完整效应条件。幂等性给
\(P_+\rho_+=P_+/2=\rho_+\)，迹为1；正交性给
\(P_+\rho_-=P_+P_-/2=0\)。\(\square\)

因而反例并非``两个数学标签其实代表同一个操作态''。它们在联合允许效应下确有可区分统计，只是局域乘积效应集合不足以重建该差别。\texttt{joint\allowbreak Test}
及其两个概率定理落实了这项要求。

\hypertarget{ux5220ux9664ux5c40ux57dfux5b9eux6027ux9650ux5236ux540eux5deeux522bux91cdux65b0ux53efux89c1}{%
\section{55.4
删除局域实性限制后，差别重新可见}\label{ux5220ux9664ux5c40ux57dfux5b9eux6027ux9650ux5236ux540eux5deeux522bux91cdux65b0ux53efux89c1}}

\protect\hypertarget{n-qrt-02}{}{}

\textbf{N-QRT-02（允许复局域效应后的区分）。} 取复局域效应 \[
E_Y=\frac{I+Y}{2}.
\] 它是合法厄米投影，虽然不满足实性。其乘积测试满足 \[
\operatorname{Tr}[(E_Y\otimes E_Y)\rho_+]=\tfrac12,
\qquad \operatorname{Tr}[(E_Y\otimes E_Y)\rho_-]=0.
\]

\textbf{证明。} \(Y\) 为厄米对合，故 \(E_Y\) 与 \(I-E_Y\)
为正交半正定投影。由
\(\operatorname{Tr}E_Y=1\)、\(\operatorname{Tr}(E_YY)=1\)，代入
\(\rho_\pm\) 的张量迹表达式得 \((1\pm1)/4\)。\(\square\)

该删除对照区分的是这里的具体态对，不等于一般复量子局域层析的完整定理。这是针对关键限制的删除见证。联合态没有被修改；改变的是允许的局域测试集合。它说明实局域限制在本例中实际承担了统计不可区分性，也防止把
N-QRT-01 错写为复量子理论自身缺少局域层析。

\hypertarget{ux7ebfux6027ux540eux5904ux7406ux4e0dux80fdux8865ux51faux7f3aux5931ux8054ux5408ux6548ux5e94}{%
\section{55.5
线性后处理不能补出缺失联合效应}\label{ux7ebfux6027ux540eux5904ux7406ux4e0dux80fdux8865ux51faux7f3aux5931ux8054ux5408ux6548ux5e94}}

\protect\hypertarget{n-qrt-03}{}{}

\textbf{N-QRT-03（区分效应不在有限实产品效应的线性张成中）。}
不存在有限指标集 \(J\)、实系数 \(c_j\) 以及实局域效应 \(E_j,F_j\)，使 \[
P_+=\sum_{j\in J}c_j(E_j\otimes F_j).
\]

\textbf{证明。} 如果存在这样的表达式，迹配对的实线性及 N-QRT-01 使右侧对
\(\rho_+,\rho_-\) 的响应相同。但 P-QRT-02
使左侧的响应分别为1与0，矛盾。\(\square\)

系数允许为负，故这个结论比``不能写成产品效应的凸组合''更强。它针对所声明的有限线性张成；没有借此证明全部物理实现或无限操作协议的排除。专用
\texttt{Real\allowbreak Tomography\allowbreak Two\allowbreak Rebit}
见证将这一线性推论与主模块的态、效应及全称等式连接。

\hypertarget{ux5728-ueot-ux91cfux5b50ux9009ux62e9ux94feux4e2dux7684ux51c6ux786eux4f4dux7f6e}{%
\section{55.6 在 UEOT
量子选择链中的准确位置}\label{ux5728-ueot-ux91cfux5b50ux9009ux62e9ux94feux4e2dux7684ux51c6ux786eux4f4dux7f6e}}

\begin{longtable}[]{@{}
  >{\raggedright\arraybackslash}p{(\columnwidth - 4\tabcolsep) * \real{0.2500}}
  >{\raggedright\arraybackslash}p{(\columnwidth - 4\tabcolsep) * \real{0.3400}}
  >{\raggedright\arraybackslash}p{(\columnwidth - 4\tabcolsep) * \real{0.4100}}@{}}
\toprule\noalign{}
\begin{minipage}[b]{\linewidth}\raggedright
选择责任
\end{minipage} & \begin{minipage}[b]{\linewidth}\raggedright
本轮实际结论
\end{minipage} & \begin{minipage}[b]{\linewidth}\raggedright
仍需承担的部分
\end{minipage} \\
\midrule\noalign{}
\endhead
\bottomrule\noalign{}
\endlastfoot
排除完整经典解释 & 第54章在指定同源 Core 接口上证明纯化失败 &
为什么自然界必须满足相应纯化原则 \\
排除这里的实组合解释 & 本章证明局域层析失败且联合态可区分 & 为什么 UEOT
的物理实现必须满足局域层析 \\
唯一选出复量子结构 & 本章只提供一个必要分离见证 &
候选理论论域、其他重构前提及完整重构定理 \\
对应自然界 & 本章给出精确概率分离 & 实际制备、测量能力及经校准的数据 \\
\end{longtable}

本章的机器证明使用复矩阵中的实条目子空间表示这个全称分离。第57章进一步给出全部有限实态、实效应、实Kraus过程、复合和偏迹的实际操作呈现，并证明七项基本规律；这仍不等于所有OPT重建背景、范畴相干性或纯化本质唯一性已形式化。本例也没有被改写为第54章H1--H6的另一个完整Core证书。

分类上，有限态与效应的标准结构属于 P，实性限制及所选组合解释承担
B／E，若要求物理系统必须局域层析则是额外的扇区选择原则
S，其显示的响应差别属于 O。只有进一步声称某个尚未证明的 UEOT
机制必然产生这些选择原则时，才另登记相应的
X。下一轮需要比较候选原则的独立性和物理意义，使``为什么选出量子''逐步成为有明确排除责任的数学问题。{[}S2、S3；QM316{]}

\hypertarget{ch-56}{%
\chapter{56. 保留相干的量子条件演化与 Doob 通道}\label{ch-56}}

第53章完成了经典次随机过程的对角量子实现，但该实现每次成功后都消去输入相干。本章直接从实际量子仪器的一个成功分支出发，在明确的正谱和可逆过滤前提下构造完整量子
Doob
通道。关键结论不只涉及一个特选稳态：任意初态的真实有限时条件态，经过指定归一化变换后，都与这个通道的实际迭代相连。

量子 Doob 变换已有开放量子系统研究先例；Carollo 等讨论了稀有动力学与另一
Markov
动力学的关系。本章独立证明以下有限离散仪器命题，不援引该论文来跳过所需有限证明，也不据此登记无限时条件极限。\href{https://arxiv.org/abs/1711.10951}{原文}

\hypertarget{ux6210ux529fux5206ux652fux771fux5b9eux8bb0ux5f55ux4e0eux5de6ux53f3ux8c31ux8d23ux4efb}{%
\section{56.1
成功分支、真实记录与左右谱责任}\label{ux6210ux529fux5206ux652fux771fux5b9eux8bb0ux5f55ux4e0eux5de6ux53f3ux8c31ux8d23ux4efb}}

固定有限维复矩阵载体和一个完整仪器，选定成功结果 \(o\)。其分支及
Heisenberg 伴随为 \[
\Phi(X)=\sum_kK_kXK_k^\dagger,
\qquad \Phi^*(A)=\sum_kK_k^\dagger AK_k.
\]
完整仪器还包含其他结果，因而本成功分支通常只保持或减少正算符的迹。\texttt{success\allowbreak Operator}
使用已有 \texttt{record\allowbreak Operator}，具体记录是同一成功结果连续出现 \(n\)
次；它不是另造的迭代代理。定义 \[
m_n(\rho)=\operatorname{Tr}\Phi^n(\rho),\qquad
\rho_n^{\mathrm{cond}}=\frac{\Phi^n(\rho)}{m_n(\rho)}
\] 时，第二式要求真实质量严格正。\texttt{success\allowbreak Mass\_\allowbreak eq\_\allowbreak response}
将第一式连接到同一仪器库的实际记录概率。

右特征密度条件是 \(\Phi(\rho_*)=r\rho_*\)。左特征条件则作用于可观测矩阵
\(H\)：\(\Phi^*(H)=rH\)。两者承担不同任务。前者给准稳态的几何存活律；后者给
Doob
完备性和任意初态的有限条件化连接。不能把一个谱条件默认为另一个已经存在的证明。

\hypertarget{ux53efux9006ux5171ux8f6dux4e0eux5f52ux4e00ux5316ux6001ux53d8ux6362}{%
\section{56.2
可逆共轭与归一化态变换}\label{ux53efux9006ux5171ux8f6dux4e0eux5f52ux4e00ux5316ux6001ux53d8ux6362}}

给定互逆矩阵 \(S,T\)，满足 \(TS=ST=I\)。记 \[
F(X)=SXS^\dagger,\quad F^{-1}(X)=TXT^\dagger,
\quad H=S^\dagger S.
\] 对密度态 \(\rho\)，令 \[
c_\rho=\operatorname{Tr}F(\rho),\qquad
\mathcal F(\rho)=F(\rho)/c_\rho.
\]

\protect\hypertarget{p-qqd-01}{}{}

\textbf{P-QQD-01（过滤归一化的合法性）。} \(F\) 与 \(F^{-1}\)
互逆并保持半正定性。对每个密度态，\(c_\rho>0\)，故 \(\mathcal F(\rho)\)
是合法密度态；而且 \(c_\rho=\operatorname{Tr}(H\rho)\)。

\textbf{证明。} 共轭保持半正定，代入 \(TS=ST=I\) 给两个逆等式。若正算符
\(F(\rho)\) 的迹为零，则由半正定零迹判据它必须为零；再施以 \(F^{-1}\) 得
\(\rho=0\)，与迹一矛盾。最后用有限矩阵迹的循环性得到
\(\operatorname{Tr}(S\rho S^\dagger)=\operatorname{Tr}(S^\dagger S\rho)\)。\(\square\)

\texttt{Invertible\allowbreak Filter}
只输入两个矩阵及真实互逆等式；\texttt{mass\_\allowbreak apply\_\allowbreak pos}
从这些条件推出正质量，没有把正质量本身塞进输入结构。此处不要求 \(S\) 是
\(H\) 的某个预先选定正平方根，因而避免附带未证明的规范选择。

必须区分这个数学变换和物理操作：\(F\)
是线性的完全正共轭，但不必保持或减少迹；\(\mathcal F\)
因归一化通常是非线性的。\(c_\rho\)
是归一化因子，可能大于一，所以不能直接称为一次筛选成功概率。本章真正证明完整性的量子通道是下一节的
\(\widetilde\Phi\)。

\hypertarget{ux4eceux5de6ux7279ux5f81ux6761ux4ef6ux8bc1ux660eux5b9eux9645-kraus-ux5b8cux5907ux6027}{%
\section{56.3 从左特征条件证明实际 Kraus
完备性}\label{ux4eceux5de6ux7279ux5f81ux6761ux4ef6ux8bc1ux660eux5b9eux9645-kraus-ux5b8cux5907ux6027}}

明确要求 \[
r>0,\qquad \Phi^*(H)=rH.
\]

\protect\hypertarget{p-qqd-02}{}{}

\textbf{P-QQD-02（实际量子 Doob 仪器）。} 定义 \[
\widetilde K_k=r^{-1/2}SK_kT.
\] 则这些矩阵构成完整单出口仪器，其通道满足 \[
\widetilde\Phi(X)=r^{-1}F\bigl(\Phi(F^{-1}(X))\bigr),
\qquad \widetilde\Phi(F(X))=r^{-1}F(\Phi(X)).
\]

\textbf{证明。} 直接求 Gram 和： \[
\begin{aligned}
\sum_k\widetilde K_k^\dagger\widetilde K_k
&=r^{-1}T^\dagger\Bigl(\sum_kK_k^\dagger S^\dagger SK_k\Bigr)T\\
&=r^{-1}T^\dagger\Phi^*(H)T
=T^\dagger HT=(ST)^\dagger(ST)=I.
\end{aligned}
\] 逐项展开 \(\widetilde K_kX\widetilde K_k^\dagger\) 得通道公式，再用
\(F^{-1}F=\mathrm{id}\) 得一次交织。\(\square\)

Lean 先用显式实权重 \(r^{-1}\) 与矩阵 \(SK_kT\) 构造
\texttt{doob\allowbreak Weighted}，再经既有正确转换得到实际平方根 Kraus 的
\texttt{doob\allowbreak Instrument}。\texttt{doob\_\allowbreak complete}
是推导出的完备性定理。该仪器的可见结果类型为 \texttt{Unit}，Kraus
指标在分支内求和；本版没有把这些指标改成可见监测结果，也没有证明任意细粒度字串的过程等价。

\hypertarget{ux4efbux610fux521dux6001ux7684ux6709ux9650ux65f6ux6761ux4ef6ux6f14ux5316}{%
\section{56.4
任意初态的有限时条件演化}\label{ux4efbux610fux521dux6001ux7684ux6709ux9650ux65f6ux6761ux4ef6ux6f14ux5316}}

\protect\hypertarget{p-qqd-03}{}{}

\textbf{P-QQD-03（全算符有限交织与真实条件态连接）。} 在 P-QQD-02
的前提下，对任意复算符 \(X\) 与任意有限 \(n\)， \[
\widetilde\Phi^n(F(X))=r^{-n}F(\Phi^n(X)),
\qquad \operatorname{Re}\operatorname{Tr}F(\Phi^n(X))=r^n\operatorname{Re}\operatorname{Tr}F(X).
\] 对任意密度初态 \(\rho\)，每个有限 \(n\) 都有 \(m_n(\rho)>0\)，且 \[
\mathcal F\bigl(\rho_n^{\mathrm{cond}}\bigr)
=\widetilde\Phi^n\bigl(\mathcal F(\rho)\bigr).
\]

\textbf{证明。} 对一次交织按 \(n\)
归纳即可得到第一式。左特征关系及伴随配对恒等式给 \[
\operatorname{Tr}F(\Phi(X))
=\operatorname{Tr}(H\Phi(X))
=\operatorname{Tr}(\Phi^*(H)X)
=r\operatorname{Tr}F(X),
\]
迭代得到第二式的实部迹关系；上述循环迹计算还给出相应的全复迹解析恒等式，本版命名
Lean 定理的类型是实
\texttt{mass}。对密度态，相关算符均半正定，迹为实数，其右侧为
\(r^nc_\rho>0\)。若 \(m_n(\rho)=0\)，半正定的 \(\Phi^n(\rho)\)
将由零迹判据等于零，使左侧也为零，矛盾。因此所有有限条件态都定义良好。

最后把实际归一化逐项代入： \[
\begin{aligned}
\mathcal F\!\left(\frac{\Phi^n(\rho)}{m_n(\rho)}\right)
&=\frac{F(\Phi^n(\rho))}{\operatorname{Tr}F(\Phi^n(\rho))}\\
&=\frac{r^{-n}}{c_\rho}F(\Phi^n(\rho))
=\widetilde\Phi^n\!\left(\frac{F(\rho)}{c_\rho}\right).
\end{aligned}
\] 第一步的 \(m_n\) 消去依赖其严格正性，不能用于零概率事件。\(\square\)

Lean 对应
\texttt{doob\_\allowbreak success\allowbreak Operator\_\allowbreak apply}、\allowbreak \texttt{eigenmetric\_\allowbreak success\allowbreak Mass}、\allowbreak \texttt{eigenmetric\_\allowbreak success\allowbreak Mass\_\allowbreak pos}
与
\texttt{doob\_\allowbreak finite\_\allowbreak conditioning}。前两项在全部算符上成立，条件态结论则明确限制为实际
\texttt{Density}。归一化证明使用原先定义的真实成功记录，而非在另一条演化上比较相似符号。

这个等式中的左侧仍含
\(\mathcal F\)。它不声称原条件态直接等于未变换初态的 Doob
轨道，更不声称所有初态的存活率都为 \(r^n\)。解析地反解第一式还可写成 \[
m_n(\rho)=r^nc_\rho\operatorname{Tr}\!\left[
T\,\widetilde\Phi^n(\mathcal F(\rho))T^\dagger\right].
\]
该式显示一般初态的真实存活质量还含一个态依赖因子；本版没有把这个书面推论另计为命名
Lean 定理。

\hypertarget{ux53f3ux7279ux5f81ux5bc6ux5ea6ux624dux7ed9ux51c6ux7a33ux6001ux4e0e-doob-ux5e73ux7a33ux6001}{%
\section{56.5 右特征密度才给准稳态与 Doob
平稳态}\label{ux53f3ux7279ux5f81ux5bc6ux5ea6ux624dux7ed9ux51c6ux7a33ux6001ux4e0e-doob-ux5e73ux7a33ux6001}}

\protect\hypertarget{p-qqd-04}{}{}

\textbf{P-QQD-04（准稳态、几何存活与实际通道平稳）。} 设
\(r>0\)，若另外给定密度态 \(\rho_*\) 满足 \(\Phi(\rho_*)=r\rho_*\)，则
\[
\Phi^n(\rho_*)=r^n\rho_*,\quad m_n(\rho_*)=r^n,
\quad \rho_{*,n}^{\mathrm{cond}}=\rho_*.
\] 该实际仪器的概率界给 \(r\le1\)；若也满足前述左特征条件，则 \[
\widetilde\Phi(\mathcal F(\rho_*))=\mathcal F(\rho_*).
\] 因此 \(\mathcal F(\rho_*)\) 是实际 \texttt{channel\allowbreak State}
的平稳密度态，全部有限重复单出口记录均保持该态且质量为一。

\textbf{证明。} 右特征式按次数归纳给第一式；取迹及使用
\(\operatorname{Tr}\rho_*=1\) 给几何质量；严格正的 \(r^n\)
使归一条件态还原为 \(\rho_*\)。一次成功概率就是
\(r\)，完整仪器给其不超过一。将右特征式代入一次交织，再按 \(c_{\rho_*}\)
归一化，得到 Doob 固定点。单出口分支就是其全部通道，因此此结论连接到实际
\texttt{channel\allowbreak State}。\(\square\)

Doob
的单出口记录质量一有迹保持的直接背景；其科学内容在于这个真正完整的通道与原成功分支的全算符交织、任意初态的条件化等式以及明确的平稳矩阵。当前没有证明所有
CP
分支都存在所需正可逆左谱矩阵，也未证明谱半径身份、简单主特征值或长期吸引。

\hypertarget{ux4e00ux4e2aux7cbeux786eux6709ux7406ux4fddux7559ux8f93ux5165ux76f8ux5e72ux7684ux5b9eux4f8b}{%
\section{56.6
一个精确有理、保留输入相干的实例}\label{ux4e00ux4e2aux7cbeux786eux6709ux7406ux4fddux7559ux8f93ux5165ux76f8ux5e72ux7684ux5b9eux4f8b}}

取二维源仪器，成功 Kraus 为 \[
K_0=\tfrac14I,\quad
K_1=\begin{pmatrix}1/4&0\\1/2&0\end{pmatrix},\quad
K_2=\begin{pmatrix}0&1/8\\0&1/4\end{pmatrix},\quad K_3=0.
\] 失败 Kraus 为 \[
L_0=\operatorname{diag}(3/4,7/8),\quad
L_1=\operatorname{diag}(1/4,1/4),\quad
L_2=L_3=\operatorname{diag}(0,1/8).
\] 直接计算给 \[
\sum_kK_k^\dagger K_k=\operatorname{diag}(3/8,9/64),
\quad \sum_kL_k^\dagger L_k=\operatorname{diag}(5/8,55/64),
\] 其和为单位矩阵。成功分支对任意复矩阵 \(X\) 的实际公式是 \[
\Phi(X)=\begin{pmatrix}
X_{00}/8+X_{11}/64&X_{01}/16+X_{00}/8+X_{11}/32\\
X_{10}/16+X_{00}/8+X_{11}/32&X_{00}/4+X_{11}/8
\end{pmatrix}.
\] 因此非对角输入 \(|0\rangle\langle1|\) 被送到其 \(1/16\)
倍，而非直接消去；同时
\(K_1K_2\ne K_2K_1\)。这两个精确见证说明本例超出了第53章的一步全对角测量制备模型，不据此额外宣称已经分类所有非经典性资源。

取 \[
\begin{aligned}
\rho_* &=\tfrac15\begin{pmatrix}1&2\\2&4\end{pmatrix},\qquad r=3/16,\\
S&=\operatorname{diag}(2,1),\quad T=\operatorname{diag}(1/2,1),\quad H=\operatorname{diag}(4,1).
\end{aligned}
\] 逐条精确核验 \[
\Phi(\rho_*)=r\rho_*,\qquad \Phi^*(H)=rH,\qquad
c_{\rho_*}=8/5,
\] \[
\sigma_*:=\mathcal F(\rho_*)=\tfrac12\begin{pmatrix}1&1\\1&1\end{pmatrix}
=|+\rangle\langle+|.
\] 归一化因子 \(8/5\)
正好展示了第56.2节的解释边界：它大于一，故不能被当作源仪器某个结果的概率。

\protect\hypertarget{p-qqd-05}{}{}

\textbf{P-QQD-05（全矩阵可计算的非对角 Doob 见证）。} 这个实例的实际
Doob 通道对每个复算符 \(X\) 满足 \[
\widetilde\Phi(X)=\tfrac13X+\tfrac23\operatorname{Tr}(X)\,\sigma_*.
\] 源准稳态连续成功 \(n\) 次的真实概率为 \((3/16)^n\)，条件态仍为
\(\rho_*\)；变换后的实际通道保持 \(\sigma_*\)。P-QQD-03
的任意初态有限条件化结论在本实例同样成立。

\textbf{证明。} 代入 \(S,T,r\)
和源分支全矩阵公式，逐条整理四个矩阵条目，得到所示线性算子恒等式。后续结论分别实例化
P-QQD-03、04。\(\square\)

专用 \texttt{Quantum\allowbreak Doob\allowbreak Examples} 对任意复 \(X\) 证明
\texttt{source\_\allowbreak branch\_\allowbreak formula} 与
\texttt{transformed\_\allowbreak branch\_\allowbreak formula}。后者明确使用复数迹；不能把只对厄米输入等于迹的实
\texttt{mass}
误用于全部复算符。各矩阵和系数均为精确有理输入，平方根只由已证明的普通仪器转换产生，浮点近似不参与证明。

\protect\hypertarget{n-qqd-01}{}{}

\textbf{N-QQD-01（必须保留的三项区分）。} 本例中，删去失败 Kraus
会破坏源仪器完备性；原准稳态 \(\rho_*\) 与过滤平稳态 \(\sigma_*\)
不同；从均匀初态 \(I/2\) 出发的一次成功概率为 \(33/128\)，并不等于
\(r=3/16\)。

\textbf{证明。} 成功 Gram 和的对角条目分别为
\(3/8,9/64\)，都不等于1。两态的非对角条目分别为 \(2/5\) 与 \(1/2\)。最后
\[
\operatorname{Tr}\Phi(I/2)=\tfrac12(3/8+9/64)=33/128\ne3/16.
\] 三个结论均有专用 Lean 见证。\(\square\)

\hypertarget{ux5411ux52a8ux529bux5b66ux5185ux6838ux63a8ux8fdbux7684ux4e0bux4e00ux9053ux95e8}{%
\section{56.7
向动力学内核推进的下一道门}\label{ux5411ux52a8ux529bux5b66ux5185ux6838ux63a8ux8fdbux7684ux4e0bux4e00ux9053ux95e8}}

这条线把第53章的对角成功嵌入推进到一般有限 Kraus
分支的条件构造，并在非对角实例上完成全矩阵计算。它与第54、55章互补：前两章检查量子选择原则的排除责任，本章在已给量子结构的条件下补强真正可使用的动力学。三者不能相互替代。

本章主要属于标准量子结构 P 和实际操作连接 B
下的条件结论；指定成功任务与所选源仪器为 E，记录概率为 O。没有直接调用
Core 来选择复量子结构，因此不能仅因主题与 Core 的
QSD／持续思想相关就把全部结果登记为
C。尚待完成的谱存在、不可约性、唯一性和条件极限是相应标准数学或领域桥责任；额外的
UEOT 物理必然性主张才另按 X 登记。

第59章已对本章同一个有理非对角模型证明实际精确迭代、唯一Doob不变态和真实条件态的几何收敛率，并以合法恒等Doob通道证明正谱度量不足以保证混合。接下来应固定可验证的有限不可约、本原或收缩条件，研究一般正谱对象及长期条件过程。原始双驱动思想可提示哪些持续机制值得研究，但成功事件、初态、控制代价与真实观测仍须分别定义。{[}S1、S2、S4；QM316；QM317{]}

\hypertarget{ch-57}{%
\chapter{57. 完整状态空间中的纯化与操作原则}\label{ch-57}}

第54章已经说明，一组真实的有限 Core
过程接口可以有完整经典解释，而其中的混态没有任何有限经典纯扩展。这个结果提出了一个正向责任：如果把纯化作为量子选择候选，必须在同一套操作语言中说清楚系统、全部允许态、复合与丢弃，再证明复量子模型确实满足所声明的原则。本章完成有限维复密度态的纯化存在构造，并明确它与量子重建之间仍然存在的距离。

\hypertarget{ux5148ux7ed9ux5b8cux6574ux64cdux4f5cux5448ux73b0ux518dux8ba8ux8bbaux9009ux62e9ux539fux5219}{%
\section{57.1
先给完整操作呈现，再讨论选择原则}\label{ux5148ux7ed9ux5b8cux6574ux64cdux4f5cux5448ux73b0ux518dux8ba8ux8bbaux9009ux62e9ux539fux5219}}

一个系统的操作呈现至少包含状态集合、效应集合、效应在状态上的概率、凸混合、带输入输出类型的过程及其状态作用。复合系统还需要产品态、产品效应与实际丢弃映射。单独的规律包包括七项：概率界、概率对混合的仿射性、产品概率乘法、产品态的边际等式、过程的仿射性、恒等过程的状态作用及顺序复合的状态作用。它们必须在具体模型中证明，不能因为一个结构字段具有所需类型，就把模型的所有物理性质视为已证。

系统类型必须真实反映复合。有限载体为 \(\iota\) 与 \(\alpha\)
时，本章的联合矩阵载体为
\(\iota\times\alpha\)。这避免把维数相同当作类型相同，也使偏迹的求和指标与被丢弃系统直接对应。输入、输出不同的过程采用第58章的矩形
Kraus 通道，其 Kraus 矩阵是从输入向量空间到输出向量空间的实际线性映射。

对有限复量子系统，状态空间始终是

\[
\mathcal D(\iota)=\{\rho\in M_\iota(\mathbb C):\rho\succeq0,\ \operatorname{Tr}\rho=1\}.
\]

效应是全部 \(0\preceq E\preceq I\) 的矩阵，概率为
\(\operatorname{Re}\operatorname{Tr}(E\rho)\)。矩阵乘积 \(E\rho\)
一般并非 Hermitian；其迹的实性由 PSD Gram
分解、迹循环和正矩阵的实迹推出，不能以乘积仍为正来代替证明。实部约定使概率函数在类型上直接取实值。对于
\(0\le t\le1\)，混合是实际矩阵
\(t\rho+(1-t)\sigma\)；正性和迹一由各自的线性关系推出。

纯性采用这个完整凸集的极点定义：

\[
\begin{aligned}
\operatorname{Pure}(\rho)\iff{}&
\forall\sigma,\tau\in\mathcal D(\iota),\ \forall 0<t<1,\\
&\rho=t\sigma+(1-t)\tau
\ \Longrightarrow\ \sigma=\rho\ \land\ \tau=\rho.
\end{aligned}
\]

这是一个对所有合法密度态量化的性质。只把秩一矩阵命名为``纯态''，或者只在某个挑选后的状态子集中检查极点，都不能替代本章的证明责任。

\hypertarget{ux5355ux4f4dux5411ux91cfux6001ux4e3aux4f55ux662fux5b8cux6574ux51f8ux96c6ux4e2dux7684ux6781ux70b9}{%
\section{57.2
单位向量态为何是完整凸集中的极点}\label{ux5355ux4f4dux5411ux91cfux6001ux4e3aux4f55ux662fux5b8cux6574ux51f8ux96c6ux4e2dux7684ux6781ux70b9}}

\protect\hypertarget{p-pur-01}{}{}

\textbf{P-PUR-01（完整密度态空间中的向量纯态）。} 令
\(P=|v\rangle\langle v|\) 且 \(\operatorname{Tr}P=1\)。则 \(P\) 在全部
\(\mathcal D(\iota)\) 中是凸极点。

\textbf{证明。} 归一条件给出 \(P^2=P\)；令 \(Q=I-P\)，则 \(Q=Q^\dagger\)
且 \(QP=0\)。假设 \(P=t\sigma+(1-t)\tau\)，其中
\(0<t<1\)，\(\sigma,\tau\) 都是任意密度态。左右压缩得到

\[
0=QPQ=tQ\sigma Q+(1-t)Q\tau Q.
\]

右边两项均为正半定矩阵。正锥的反对称性与严格正的两个系数迫使
\(Q\sigma Q=Q\tau Q=0\)。

还需要从``正压缩为零''取得整个支撑结论。对任意正半定 \(X\)，由有限矩阵的
\(C^*\) 正性表示取 \(X=B^\dagger B\)。若 \(QXQ=0\)，则

\[
\operatorname{Tr}\bigl((BQ)^\dagger(BQ)\bigr)=0,
\]

故 \(BQ=0\)，从而 \(XQ=B^\dagger BQ=0\)。Hermitian 性再给 \(QX=0\)。因此
\(XP=PX=X\)。对秩一 \(P\)，直接矩阵乘法有

\[
PXP=\operatorname{Tr}(XP)P.
\]

代入 \(X=\sigma\) 并使用 \(\operatorname{Tr}\sigma=1\)，得到
\(\sigma=P\)；对 \(\tau\) 同理。整个论证中的 \(\sigma,\tau\)
没有被限制为秩一、对角态或有限样本。\(\square\)

Lean 中的 \texttt{Is\allowbreak Pure} 对实际 \texttt{Density}
类型量化；\texttt{positive\_\allowbreak compression\_\allowbreak zero}、\allowbreak \texttt{density\_\allowbreak eq\_\allowbreak vector\_\allowbreak of\_\allowbreak compression\_\allowbreak zero}
与 \texttt{vector\allowbreak State\_\allowbreak is\allowbreak Pure}
分别承担压缩支撑、迹一识别和完整极点证明。调用的矩阵正性定理属于明确的数学依赖，不是新增物理选择原则。{[}QM317{]}

\hypertarget{ux5b9eux9645ux504fux8ff9ux4e0eux4efbux610fux6709ux9650ux6001ux7684-gram-ux7eafux5316}{%
\section{57.3 实际偏迹与任意有限态的 Gram
纯化}\label{ux5b9eux9645ux504fux8ff9ux4e0eux4efbux610fux6709ux9650ux6001ux7684-gram-ux7eafux5316}}

被丢弃的辅助系统具有载体 \(\alpha\)。对每个 \(a\in\alpha\)，定义矩形
Kraus 矩阵

\[
L_a(i,(j,b))=\delta_{ij}\delta_{ab}.
\]

直接求和给
\(\sum_aL_a^\dagger L_a=I_{\iota\times\alpha}\)，故这些矩阵定义一个实际通道。它对任意复矩阵的作用是

\[
\operatorname{Tr}_{\alpha}(X)_{ij}
=\sum_aX_{(i,a),(j,a)}.
\]

因此偏迹保持正性和复迹；其任意有限辅助放大也保持所有联合正输入的正性。这里没有把检查限制为产品态。对于产品密度态，进一步由
\(\operatorname{Tr}\sigma=1\) 得到
\(\operatorname{Tr}_{\alpha}(\rho\otimes\sigma)=\rho\)。

\protect\hypertarget{p-pur-02}{}{}

\textbf{P-PUR-02（任意有限复密度态的纯化存在）。} 对任意
\(\rho\in\mathcal D(\iota)\)，存在
\(\Psi\in\mathcal D(\iota\times\iota)\)，使 \(\Psi\)
在完整联合密度态空间中为纯态，且 \(\operatorname{Tr}_{2}\Psi=\rho\)。

\textbf{证明。} 从 \(\rho\succeq0\) 导出一个矩阵 \(A\) 满足
\(\rho=AA^\dagger\)，并定义联合向量 \(\psi(i,a)=A_{ia}\)。其外积
\(\Psi=|\psi\rangle\langle\psi|\) 为正，且

\[
\operatorname{Tr}\Psi
=\sum_{i,a}|A_{ia}|^2
=\operatorname{Tr}(AA^\dagger)=1.
\]

P-PUR-01 证明 \(\Psi\) 对完整联合状态空间确为极点。实际偏迹逐项给出

\[
(\operatorname{Tr}_{2}\Psi)_{ij}
=\sum_a A_{ia}\overline{A_{ja}}
=(AA^\dagger)_{ij}=\rho_{ij}.
\]

因子 \(A\)
的存在来自原密度态的正性，不作为调用者必须另行提供的假设。\(\square\)

形式构造 \texttt{purification} 选择一个由 PSD 存在性给出的 Gram
因子；\texttt{purification\_\allowbreak is\allowbreak Pure} 与 \texttt{purification\_\allowbreak marginal}
分别证明纯性和边际，\texttt{exists\_\allowbreak purification}
组合二者。这里的选择使用已列入基础审计的
\texttt{Classical.\allowbreak choice}。它证明存在并给出 Lean
中的非计算选择，并未附带一个已经认证的数值矩阵平方根算法。对于空载体，迹一密度态根本不存在；定理不靠隐藏的非空假设制造状态。{[}QM317{]}

Bell 向量 \(\psi=(|00\rangle+|11\rangle)/\sqrt2\)
给出熟悉的二维正控，其边际为 \(I/2\)。具体 Lean 见证在两个非零位置使用
\((1+i)/2\)，保持有理复系数；这与标准向量仅差整体相位，外积为同一个 Bell
投影。一般定理的量词覆盖任意有限维度与任意合法初态；这个小例子用于核对索引与边际，不能代替一般证明。

\hypertarget{ux4e09ux79cdux5b9eux9645ux6a21ux578bux4e0eux76f8ux4e92ux72ecux7acbux7684ux5019ux9009ux539fux5219}{%
\section{57.4
三种实际模型与相互独立的候选原则}\label{ux4e09ux79cdux5b9eux9645ux6a21ux578bux4e0eux76f8ux4e92ux72ecux7acbux7684ux5019ux9009ux539fux5219}}

\protect\hypertarget{p-opt-01}{}{}

\textbf{P-OPT-01（同一操作语言中的三个具体呈现）。}
全经典单纯形、全实密度矩阵与全复密度矩阵分别具有实际状态、效应、凸混合、产品、丢弃和所声明过程；上述七项基本规律均已分别由具体数据证明，对应
\texttt{classical\_\allowbreak laws}、\allowbreak \texttt{real\_\allowbreak laws}、\allowbreak \texttt{complex\_\allowbreak laws}。纯化存在和局部层析另定义为候选性质，不作为呈现的必备字段。

经典模型的状态是全部有限概率分布，效应是全部 \([0,1]\)
值函数；复合是笛卡尔积上的全部分布，丢弃是求和边际，过程采用实际随机核。实模型使用全部实对称
PSD 迹一矩阵、实效应和实 Kraus
过程；复模型使用上述复矩阵对象。产品概率的乘法来自经典乘积求和或矩阵迹的
Kronecker
乘法，过程仿射性来自实际线性作用。实矩阵嵌入复矩阵时，必须保存取值实性与所用联合载体，不能把整个复状态集合误当作实模型的状态空间。

在这种语言中，纯化存在是

\[
\forall A\ \forall\rho\in\operatorname{State}(A),\
\exists B\ \exists\Psi\in\operatorname{State}(A\otimes B),\
\operatorname{Pure}(\Psi)\ \land\ \operatorname{discard}_{B}\Psi=\rho.
\]

局部层析则要求：若两个联合态在所有产品效应上的概率相同，它们就是同一个联合态。前者讨论状态的纯扩展，后者讨论联合状态能否被局部测试统计识别；逻辑责任不同。

\hypertarget{ux53cdux6a21ux578bux5982ux4f55ux7ea6ux675fux91cfux5b50ux9009ux62e9ux800cux4e0dux9884ux8bbeux7b54ux6848}{%
\section{57.5
反模型如何约束量子选择，而不预设答案}\label{ux53cdux6a21ux578bux5982ux4f55ux7ea6ux675fux91cfux5b50ux9009ux62e9ux800cux4e0dux9884ux8bbeux7b54ux6848}}

\protect\hypertarget{n-opt-01}{}{}

\textbf{N-OPT-01（基本操作规律不自动给出量子选择）。}
完整经典模型满足所声明基本规律，但第54章的非平凡半混合不具有任意有限纯扩展；实量子模型的第55章两态在所有实局部产品效应下同概率，却能被合法实联合效应区分。把两个证据接到同一操作呈现后，前者否定该经典模型的纯化存在候选，后者否定该实模型的局部层析候选。

第54章的经典混态仍有普通扩展，例如与一个确定辅助态取产品。失败的是``扩展必须纯''这一附加要求。另一方面，若故意把允许态限制为只含某个混态的单元素集合，该态在这个缩小集合中会成为极点。这个负控表明，证明纯性时的完整状态量词本身属于科学内容，不能悄悄缩减。

实量子反例也不能改写成``实量子不允许纯化''。本章使用它排除的是局部层析：第55章的两个实联合态具有相同的全部实产品测试统计。允许其具体复局部效应后，原有区分重新可见。这个删除限制的对照检验准确定位了反例所依赖的局部实性。

由此可以得到一条严格的研究路径：对某组独立提出的选择原则，逐模型检查满足与失败，并追查原则是否有
UEOT
内部或经验上的独立理由。不能从``在三个检查过的模型中复模型表现合适''推出``所有替代理论已经排除''。本章也没有把完整106条
Core 命题的所有解释穷尽为这三个模型。

\hypertarget{ux7eafux5316ux5b58ux5728ux4e0eux5b8cux6574ux91cfux5b50ux91cdux5efaux7684ux4e0bux4e00ux5c42ux8d23ux4efb}{%
\section{57.6
纯化存在与完整量子重建的下一层责任}\label{ux7eafux5316ux5b58ux5728ux4e0eux5b8cux6574ux91cfux5b50ux91cdux5efaux7684ux4e0bux4e00ux5c42ux8d23ux4efb}}

量子重建文献中的纯化原则通常包含辅助系统上可逆变换意义下的本质唯一性。Chiribella、D'Ariano
与 Perinotti
的重建还使用因果性、完美可区分性、理想压缩、局部可区分性和纯条件化等背景原则；不能只截取``纯化存在''和``局部层析''两个名称，就引用其结论为本项目的已证推论。{[}CDP2010；CDP2011{]}

本章给出有限复矩阵实现中的存在构造与向量态凸极点责任。3.3版在这里保留的反向纯性与唯一性缺口，已由第61章补上：完整凸纯态必为归一向量态；同辅助载体的任意两个纯化由实际辅助酉连接；不同辅助载体单列等距嵌入和共同扩展。第60章证明当前复操作呈现对任意联合态满足局部层析。这些结果绑定新的实际源码与验收，不追溯改写本章原始证明的范围。完整重建背景和UEOT为何必须满足相应原则，仍是独立桥接责任。

对原始思想的推进在于，``父对象能支持子对象的混合描述''现在有了具体构造和明确边界：纯父态的实际偏迹给出任意有限复密度态；同一句思想在完整经典模型里并不普遍成立。只有补充原则的来源也受到证明或经验约束，才可能把这种差别继续推进为物理选择。

\hypertarget{ch-58}{%
\chapter{58. 联合记录、量子记忆与有限自适应协议}\label{ch-58}}

仪器把一个输入态映到多个带结果标签的未归一化分支。把分支相加得到无条件量子通道，但这一步已经丢弃了结果信息。第49章的完整记录律和第50章的效应距离因而需要一个共同载体：把真实经典记录与每个记录对应的量子记忆保存为一个联合状态。本章从实际
Kraus 矩阵构造这个载体，并把它推广到已有的任意有限自适应协议。

\hypertarget{ux8f93ux5165ux8f93ux51faux5f02ux578bux7684-kraus-ux901aux9053}{%
\section{58.1 输入输出异型的 Kraus
通道}\label{ux8f93ux5165ux8f93ux51faux5f02ux578bux7684-kraus-ux901aux9053}}

输入载体为 \(\iota\)，输出载体为 \(\omega\)。一个有限矩形 Kraus 通道给出
\(K_k\in M_{\omega\times\iota}(\mathbb C)\)，并满足输入空间上的完整性

\[
\sum_kK_k^\dagger K_k=I_\iota.
\]

它的实际算符映射为
\(\mathcal C(X)=\sum_kK_kXK_k^\dagger\)。输入与输出维数可以不同；迹恒等式仍在各自空间中有意义。

\protect\hypertarget{p-cq-01}{}{}

\textbf{P-CQ-01（矩形通道的完全正性、迹保持与效应收缩）。}
上述映射对任意复矩阵保持复迹。对任意有限辅助载体 \(\nu\)，由
\(K_k\otimes I_\nu\) 定义的实际放大保持所有 \(X\succeq0\)
的正性。对任意输出效应 \(E\)，其拉回
\(\mathcal C^*(E)=\sum_kK_k^\dagger E K_k\) 是合法输入效应，且

\[
\operatorname{Tr}(E\mathcal C(X))=\operatorname{Tr}(\mathcal C^*(E)X),\qquad
d_{\rm eff}(\mathcal C(\rho),\mathcal C(\sigma))\le d_{\rm eff}(\rho,\sigma).
\]

\textbf{证明。} 每项 \(K_kXK_k^\dagger\)
对正输入为正；迹循环与完整性给出

\[
\operatorname{Tr}\mathcal C(X)
=\operatorname{Tr}\!\left(\left(\sum_kK_k^\dagger K_k\right)X\right)
=\operatorname{Tr}X.
\]

放大的 Kraus 完整性为
\(\sum_k(K_k^\dagger K_k)\otimes I=I\)，故同一正性证明对任意联合正输入成立。逐项展开还识别其作用为真正的恒等放大，而不是另一个只在产品态上相符的映射。

对于效应，\(\mathcal C^*(E)\succeq0\)，并且
\(I-\mathcal C^*(E)=\mathcal C^*(I-E)\succeq0\)。迹配对关系让每个输出测试成为一个输入测试。对所有输出效应取上确界即得收缩。恒等过程和顺序组合也由实际
Kraus 给出：先 \(\mathcal C\) 后 \(\mathcal D\) 的 Kraus 为
\(D_lC_k\)。\(\square\)

本章的完全正性包括任意有限辅助系统上的正输入；误差度量仍是明示输入空间上的全效应距离。这个正性结论本身不把后面的缺陷上确界升级为
diamond 距离。{[}QM317{]}

\hypertarget{ux4e00ux6b65ux4eeaux5668ux7684ux8bb0ux5f55ux5bc4ux5b58ux5668ux4e0eux4e24ux4e2aux8fb9ux9645}{%
\section{58.2
一步仪器的记录寄存器与两个边际}\label{ux4e00ux6b65ux4eeaux5668ux7684ux8bb0ux5f55ux5bc4ux5b58ux5668ux4e0eux4e24ux4e2aux8fb9ux9645}}

给定完整仪器 \(I\) 的 Kraus \(K_{o,k}\)，记录载体为结果集合 \(O\)。定义

\[
V_{o,k}((p,i),j)=\delta_{po}K_{o,k}(i,j).
\]

这些是从 \(\iota\) 到 \(O\times\iota\) 的实际矩形矩阵。对标签求和有
\(V_{o,k}^\dagger V_{o,k}=K_{o,k}^\dagger K_{o,k}\)，因此完整性直接来自原仪器。

\protect\hypertarget{p-cq-02}{}{}

\textbf{P-CQ-02（真实仪器的联合块状态与边际）。} 上述 Kraus 定义通道

\[
\mathcal R_I(X)=\bigoplus_{o\in O}\Phi_{I,o}(X),\qquad
\Phi_{I,o}(X)=\sum_kK_{o,k}XK_{o,k}^\dagger.
\]

块等式对任意复输入 \(X\) 成立。对于密度态
\(\rho\)，\(\Omega_{I,\rho}=\mathcal R_I(\rho)\)
是合法联合密度态，其量子边际和记录边际分别为

\[
\operatorname{Tr}_{O}\Omega_{I,\rho}=\sum_o\Phi_{I,o}(\rho),\qquad
\operatorname{Tr}_{\iota}\Omega_{I,\rho}
=\operatorname{diag}\bigl(p_I(o\mid\rho)\bigr).
\]

\textbf{证明。} 逐项计算输出 \((p,i),(q,j)\) 元素。两个标签约束使
\(p\ne q\) 的块为零，\(p=q=o\) 的块恰为实际分支
\(\Phi_{I,o}(X)\)。两个偏迹分别对标签或记忆指标求和，给出上述等式；每个记录块的迹就是已有仪器概率。\(\square\)

若 \(p_I(o\mid\rho)>0\)，把该块除以其迹，得到已有的
\texttt{posterior}，没有重新指定另一种条件更新。若该概率为零，则该正块整体为零：正半定矩阵的零迹迫使矩阵为零。因此零概率事件没有在此被赋予任意的``实际后验态''。

这是一种操作上明确的测量记录表示。它证明已声明仪器与公开记录之间的一致性，没有从纯粹丢弃或块对角化推导客观坍缩机制。

\hypertarget{ux4efbux610fux6709ux9650ux81eaux9002ux5e94ux6811ux4e2dux7684ux771fux5b9eux53f6ux5757}{%
\section{58.3
任意有限自适应树中的真实叶块}\label{ux4efbux610fux6709ux9650ux81eaux9002ux5e94ux6811ux4e2dux7684ux771fux5b9eux53f6ux5757}}

已有协议类型递归为停止节点或动作节点
\(\operatorname{node}(a,\pi_o)\)。每个叶 \(\ell\) 有既定的时间顺序词
\(\operatorname{word}_\pi(\ell)\)；不同结果可以进入深度不同的子树。联合输出空间采用实际叶类型
\(\operatorname{Leaf}(\pi)\times\iota\)，保留完整记录与终态记忆。

对每个叶定义 Kraus 路径类型。停止叶只有一个空路径，其矩阵是 \(I\)；在叶
\((o,\ell')\) 上，一条路径由本步 Kraus 下标 \(k\)
与子树路径组成。路径矩阵按时间顺序为

\[
K_{\text{path}}=K_{\text{tail}}K_{a,o,k}.
\]

总 Kraus 下标是
\(\sum_{\ell:\operatorname{Leaf}(\pi)}\operatorname{Path}(\pi,\ell)\)。将每条路径矩阵嵌入其真实叶标签块，就得到协议的矩形通道。这个依赖索引同时保存叶标签、路径长度与所选子树，不要求所有叶等深。

\protect\hypertarget{p-cq-03}{}{}

\textbf{P-CQ-03（全部有限协议的记录与记忆实现）。}
对任意当前有限完整协议 \(\pi\)、实际仪器族 \(I\) 和密度态
\(\rho\)，构造得到的联合密度态满足

\[
\Omega_{\pi,\rho}[\ell,\ell]
=\operatorname{recordOperator}\bigl(I,\operatorname{word}_\pi(\ell),\rho\bigr),
\qquad
\Omega_{\pi,\rho}[\ell,m]=0\quad(\ell\ne m).
\]

因此该块的迹等于已有的
\texttt{leaf\allowbreak Probability}，记录边际是完整叶分布，量子边际是所有真实终端未归一化记忆的和。

\textbf{证明。} 对协议树归纳。停止情形是单位 Kraus
与原输入态。动作节点处，对每个结果先使用子树的完整性，再对本步 Kraus
求和，得到

\[
\sum_{o,k}K_{a,o,k}^\dagger
\left(\sum_{\ell',\text{tail}}K_{\text{tail}}^\dagger K_{\text{tail}}\right)
K_{a,o,k}
=\sum_{o,k}K_{a,o,k}^\dagger K_{a,o,k}=I.
\]

对固定叶，子树路径共轭和等于其记录算符；再代入本步分支，即与已有按时间顺序定义的
\texttt{record\allowbreak Operator}
一致。标签嵌入给出联合块等式，最后取偏迹。这个叶块关系是由路径乘积和树归纳证明的结论，没有作为通道构造的输入字段。\(\square\)

由于保存的是整个未归一化分支矩阵，本结果比只保存叶概率更强。后续效应可以依赖完整叶记录并作用于相应量子记忆；模型仍限定在明示的有限协议和记忆载体内，没有涵盖一个未建模的任意外部环境。

\hypertarget{ux5b9eux9645ux8bb0ux5f55ux7f3aux9677ux4e8bux4ef6ux6548ux5e94ux4e0eux7cbeux786e-tv-ux5e38ux6570}{%
\section{58.4 实际记录缺陷、事件效应与精确 TV
常数}\label{ux5b9eux9645ux8bb0ux5f55ux7f3aux9677ux4e8bux4ef6ux6548ux5e94ux4e0eux7cbeux786e-tv-ux5e38ux6570}}

对于共同输入、共同输出载体的矩形通道，定义实际缺陷

\[
\delta(\mathcal C,\mathcal D)=
\sup\left(\{0\}\cup\{d_{\rm eff}(\mathcal C(\tau),\mathcal D(\tau)):
\tau\in\mathcal D(\iota)\}\right).
\]

并入 \(0\) 使空输入状态集合的边界也有明确含义；所有实际距离均在
\([0,1]\)，故上确界有界。对记录通道取
\(\delta_{\rm rec}(I,J)=\delta(\mathcal R_I,\mathcal R_J)\)。它由实际通道定义，不把用户给定的某个容差数字命名为``缺陷''。

对结果事件 \(S\subseteq O\)，联合效应为

\[
E_S=\operatorname{diag}(\mathbf1_S)\otimes I_\iota.
\]

它及其补效应都为正，其对联合块状态的概率恰为事件 \(S\) 的真实记录概率。

\protect\hypertarget{p-cq-04}{}{}

\textbf{P-CQ-04（记录分布的联合距离与扰动控制）。} 设 \(p_{I,\rho}\)
为实际记录分布，则

\[
\operatorname{TV}(p_{I,\rho},p_{J,\sigma})
\le d_{\rm eff}(\mathcal R_I(\rho),\mathcal R_J(\sigma))
\le d_{\rm eff}(\rho,\sigma)+\delta_{\rm rec}(I,J).
\]

同样的结论适用于第58.3节构造的整个有限协议联合通道及其真实叶分布；此时缺陷指该整体协议通道的实际缺陷。

\textbf{证明。} 第二个不等式先在 \(\mathcal R_I(\sigma)\)
处分解三角不等式，再使用 P-CQ-01
的输入距离收缩和缺陷的上确界性质。第一个不等式对每个事件使用
\(E_S\)，得到

\[
|p(S)-q(S)|\le d_{\rm eff}(\Omega_p,\Omega_q).
\]

对有限归一化分布，令 \(S_+=\{o:p(o)\ge q(o)\}\)。由于
\(\sum_o(p(o)-q(o))=0\)，正差与负差总量相同，故

\[
p(S_+)-q(S_+)=\frac12\sum_o|p(o)-q(o)|=\operatorname{TV}(p,q).
\]

这个实际事件达到 TV。于是``全部事件差不超过
\(\varepsilon\)''与``完整分布 TV 不超过 \(\varepsilon\)''等价，取
\(S_+\) 即得第一式。\(\square\)

\texttt{Record\allowbreak Probability}
单独证明了上述事件等价；联合量子模块则构造实际效应并证明概率相符。这保留了界中的常数。逐单点估计再相加只能直接得到另一个可能随结果数增长的界，不能未经证明保留这里的常数。{[}QM317{]}

\hypertarget{ux4e22ux5f03ux6807ux7b7eux4f1aux62b9ux53bbux771fux5b9eux53efux533aux5206ux6027}{%
\section{58.5
丢弃标签会抹去真实可区分性}\label{ux4e22ux5f03ux6807ux7b7eux4f1aux62b9ux53bbux771fux5b9eux53efux533aux5206ux6027}}

\protect\hypertarget{n-cq-01}{}{}

\textbf{N-CQ-01（相同无条件通道不控制记录误差）。}
取相同的二维量子记忆，仪器 \(I\) 确定输出标签 \(0\) 并保持输入态，仪器
\(J\) 确定输出标签 \(1\) 并保持输入态。两者都由完整 Kraus
家族实现。它们丢弃标签后的通道同为恒等，旧无标签通道缺陷为零；其实际记录分布
TV 为 \(1\)，联合记录缺陷为 \(1\)。

\textbf{证明。} 对 \(I\)，标签 \(0\) 分支取单位 Kraus，标签 \(1\)
分支取零矩阵；对 \(J\) 交换标签。完整性均为
\(I^\dagger I=I\)，无条件作用均为输入矩阵本身。任意相同初态下，两种记录律分别是两个不同
Dirac 分布。效应 \(E_{\{0\}}\) 的概率分别为 \(1\) 和
\(0\)，故联合距离至少为 \(1\)；全效应距离至多为 \(1\)，因此等于
\(1\)。再对初态取上确界给出缺陷。\(\square\)

若共同控制器根据第一步的标签选择第二步动作，记录差异就会进入后续真实协议。例如读到标签
\(o\) 后选择确定输出 \(o\)
的动作，第二步输出再次区分两种仪器族。这个两步对照已通过实际协议叶词、后续事件效应与完整记录概率核对：第二次结果为零的事件概率分别为1和0，完整记录TV距离为1。量子记忆在这个删除控制中保持相同，已证明的差异位于实际记录及其后续选择。

\hypertarget{ux5df2ux95edux5408ux7684ux63a5ux53e3ux4e0eux4e0bux4e00ux9879ux8befux5deeux8d23ux4efb}{%
\section{58.6
已闭合的接口与下一项误差责任}\label{ux5df2ux95edux5408ux7684ux63a5ux53e3ux4e0eux4e0bux4e00ux9879ux8befux5deeux8d23ux4efb}}

本章把原有单步仪器、完整有限记录律和量子态误差接在一个实际联合系统上。已有信息足以讨论每个有限协议的整个输出与全事件误差；其实际缺陷仍可能难以计算。

第62章完成了这里提出的固定深度共同因果控制器任务：逐步构造真实历史标签加量子记忆通道，接通原Protocol的实际叶与词顺序，并从共同输入的PSD块、归一化权重、零块及合法效应限制推出逐动作记录预算。新证明承担了从整体缺陷到可组合预算的责任；本章的整体协议缺陷没有被重新命名为逐动作结论。变长停止树、相干控制器和任意额外环境仍需要各自的扩展证明。

这也明确了 UEOT
预测对象的一项结构要求：若允许的未来动作依赖记录，丢弃该记录通常会改变可用的未来测试和对象等价。第54章的受限读出反例与本章的标签反例从两个方向说明，选择接口的范围会改变可区分性。完整物理理论必须把允许观察和允许控制作为真实模型内容维护。

\hypertarget{ch-59}{%
\chapter{59. 非交换条件持续的精确动力学}\label{ch-59}}

第56章已经从给定的正左特征度量构造实际 Doob
通道，并证明有限次条件演化的交织关系。一般构造并不保证唯一平稳态或吸引性。本章固定第56章同一个有理、非对角且
Kraus
不交换的模型，证明实际迭代、实际记录概率和条件态的闭式及收敛率，再用一个合法恒等
Doob 模型检验不能推广的结论。

\hypertarget{ux56faux5b9aux6a21ux578bux4e0eux4e24ux79cdux72b6ux6001ux5750ux6807}{%
\section{59.1
固定模型与两种状态坐标}\label{ux56faux5b9aux6a21ux578bux4e0eux4e24ux79cdux72b6ux6001ux5750ux6807}}

沿用

\[
r=\frac3{16},\qquad S=\operatorname{diag}(2,1),\qquad T=S^{-1},\qquad
H=S^\dagger S=\operatorname{diag}(4,1),
\]

\[
\rho_* =\frac15\begin{pmatrix}1&2\\2&4\end{pmatrix},\qquad
\sigma_+=\frac12\begin{pmatrix}1&1\\1&1\end{pmatrix}.
\]

源的实际成功分支记为 \(\Phi\)，其已验证右特征密度为
\(\rho_*\)，\(\Phi(\rho_*)=r\rho_*\)；左谱关系是
\(\Phi^*(H)=rH\)。完整源仪器还具有真实失败分支。本章不改变其
Kraus，不把这个实例替换为一个更易处理的对角经典模型。

Doob 通道是

\[
D(X)=\frac13X+\frac23\operatorname{Tr}(X)\sigma_+.
\]

这个公式对所有复矩阵成立。对非 Hermitian
输入，迹可以是复数，不能提前只保留其实部。对密度态，归一化坐标变换为

\[
F(\tau)=\frac{S\tau S^\dagger}{\operatorname{Tr}(H\tau)}.
\]

\(F\) 的分母依赖输入，所以它一般不是线性通道。\(\Phi\) 的成功概率、\(F\)
的归一因子和 \(D\) 的迹保持各自有不同含义。

公共模块 \texttt{Quantum\allowbreak Doob\allowbreak Model}
从上一版具体见证中提取同一模型，使通用动力学源码不再依赖见证脚本。提取的原有定义与证明仍按来源登记；搬入公共模块不计为另一个新的物理结论。{[}QM317{]}

\hypertarget{doob-ux5168ux8f68ux9053ux7684ux95edux5f0fux552fux4e00ux6027ux4e0eux7cbeux786eux6536ux7f29}{%
\section{59.2 Doob
全轨道的闭式、唯一性与精确收缩}\label{doob-ux5168ux8f68ux9053ux7684ux95edux5f0fux552fux4e00ux6027ux4e0eux7cbeux786eux6536ux7f29}}

令 \(a_n=3^{-n}\)。

\protect\hypertarget{p-qdd-01}{}{}

\textbf{P-QDD-01（当前 Doob 通道的精确动力学）。} 对任意复矩阵 \(X\)
与任意 \(n\ge0\)，

\[
D^n(X)=a_nX+(1-a_n)\operatorname{Tr}(X)\sigma_+.
\]

\(\sigma_+\) 是唯一不变密度态。对任意密度态 \(\tau,\upsilon\)，

\[
d_{\rm eff}(D^n\tau,D^n\upsilon)
=a_n d_{\rm eff}(\tau,\upsilon),
\qquad D^n\tau\longrightarrow\sigma_+.
\]

最后的极限是实际有限矩阵空间通常拓扑中的整条轨道极限。

\textbf{证明。}
由迹保持，\(\operatorname{Tr}D(X)=\operatorname{Tr}X\)。\(n=0\)
时公式成立；把公式应用于 \(D(X)\)，用 \(a_{n+1}=a_n/3\)
合并系数，归纳即得全复算符闭式。若 \(D(\tau)=\tau\)，则
\(\tau=\frac13\tau+\frac23\sigma_+\)，逐矩阵元素移项得
\(\tau=\sigma_+\)。这个唯一性证明直接逐矩阵元素移项；第60章另行补上全效应距离的分离定理，不改变此处原证明。

两个密度态轨道相减得到
\(D^n\tau-D^n\upsilon=a_n(\tau-\upsilon)\)。对每个效应，实际概率差的绝对值缩放
\(a_n\)
倍；当系数严格正时，除以该系数给出上确界反向不等式，系数为零时直接处理。因此全效应上确界也精确缩放，而不只是有一个上界。\(a_n\to0\)
和矩阵加法、实数乘法的连续性给出整轨道极限。\(\square\)

这种全效应缩放证明只使用概率配对与上确界。它没有把当前效应距离与另一种矩阵范数的同一性作为隐藏前提，也没有将
Cesàro 子列收敛改称整轨道收敛。

\hypertarget{ux4eceux4ea4ux7ec7ux5173ux7cfbux56deux5230ux771fux5b9eux6210ux529fux8bb0ux5f55}{%
\section{59.3
从交织关系回到真实成功记录}\label{ux4eceux4ea4ux7ec7ux5173ux7cfbux56deux5230ux771fux5b9eux6210ux529fux8bb0ux5f55}}

对任意输入密度态 \(\tau\)，定义

\[
c_\tau=\operatorname{Tr}(H\tau),\qquad
b_\tau=\frac58c_\tau,\qquad
z_n(\tau)=a_n+(1-a_n)b_\tau.
\]

因为 \(\tau_{00},\tau_{11}\ge0\) 且 \(\tau_{00}+\tau_{11}=1\)，有

\[
c_\tau=4\tau_{00}+\tau_{11},\qquad
1\le c_\tau\le4,\qquad
\frac58\le b_\tau\le\frac52.
\]

又因为 \(0<a_n\le1\)，\(z_n\) 是 \(1\) 与 \(b_\tau\) 的凸组合，故

\[
\frac58\le z_n(\tau)\le\frac52.
\]

这些界来自完整密度态的正性和迹一；没有只对角态或有限样本输入作检查。\(c_\tau\)
是滤波归一因子，可能大于 \(1\)，不是物理成功概率。

\protect\hypertarget{p-qdd-02}{}{}

\textbf{P-QDD-02（实际源成功算符与质量闭式）。} 对上述同一个源模型，

\[
\Phi^n(\tau)
=r^n\left[a_n\tau+(1-a_n)b_\tau\rho_*\right],
\]

\[
m_n(\tau)=\operatorname{Tr}\Phi^n(\tau)
=r^nz_n(\tau).
\]

这里 \(\Phi^n\) 是已有 \texttt{success\allowbreak Operator}，\(m_n\) 是已有
\texttt{success\allowbreak Mass}，也等于实际仪器族重复成功词的 \texttt{response}。

\textbf{证明。} 第56章的交织关系在任意复矩阵上给出

\[
D^n(SXS^\dagger)=r^{-n}S\Phi^n(X)S^\dagger.
\]

左右乘 \(T,T^\dagger\)，代入 P-QDD-01，并使用

\[
T\sigma_+T^\dagger=\frac58\rho_*.
\]

即得任意复 \(X\) 的回拉公式。令 \(X=\tau\)，用
\(\operatorname{Tr}(S\tau S^\dagger)=c_\tau\)，得到密度态版本。取迹并用两态的迹一，取得实际质量等式。\(\square\)

这个证明先连接实际算符迭代，再取迹；没有另行定义一个相同外观的数列并把它视为原记录概率。成功质量处处严格正，而且

\[
\frac58r^n\le m_n(\tau)\le\frac52r^n.
\]

上下界是统一控制，未声称一般初态有 \(m_n=r^n\)。

\hypertarget{ux771fux5b9eux6761ux4ef6ux6001ux7684ux7cbeux786eux51e0ux4f55ux8d8bux8fd1}{%
\section{59.4
真实条件态的精确几何趋近}\label{ux771fux5b9eux6761ux4ef6ux6001ux7684ux7cbeux786eux51e0ux4f55ux8d8bux8fd1}}

由于每个有限 \(n\) 的成功质量严格正，已有 \texttt{conditional\allowbreak State}
对全部密度输入都定义良好。

\protect\hypertarget{p-qdd-03}{}{}

\textbf{P-QDD-03（实际条件态及其定量收敛）。} 记
\(\tau_{\mid n}=\Phi^n(\tau)/m_n(\tau)\)，则

\[
\tau_{\mid n}
=\frac{a_n}{z_n(\tau)}\tau
+\frac{(1-a_n)b_\tau}{z_n(\tau)}\rho_*.
\]

两个系数非负且和为 \(1\)。进一步有

\[
d_{\rm eff}(\tau_{\mid n},\rho_*)
=\frac{a_n}{z_n(\tau)}d_{\rm eff}(\tau,\rho_*)
\le\frac85\,3^{-n}d_{\rm eff}(\tau,\rho_*).
\]

实际条件态矩阵收敛到 \(\rho_*\)，其全效应距离也趋于零。

\textbf{证明。} 用 P-QDD-02 除以严格正的 \(r^nz_n\)，消去公共因子
\(r^n\) 得到条件态公式。归一因子的定义使第二个系数等于
\(1-a_n/z_n\)，因此

\[
\tau_{\mid n}-\rho_*
=\frac{a_n}{z_n(\tau)}(\tau-\rho_*).
\]

重复 P-QDD-01 的全效应上确界缩放证明得到精确等式。\(z_n\ge5/8\)
给出系数上界 \(a_n/z_n\le(8/5)a_n\)。由于 \(z_n\to b_\tau>0\)，系数
\(a_n/z_n\to0\)，从而得到矩阵与效应距离的两个真实极限。\(\square\)

在这个指定模型中，归一化的源条件演化具有非线性分母，Doob
轨道则由线性迹保持通道产生。交织关系让两者可比较；这不意味着两种过程的物理概率或平稳态可以互换。

\hypertarget{ux6761ux4ef6ux7a33ux6001ux4e0eux65e0ux6761ux4ef6ux6210ux529fux8d28ux91cfux7684ux4e0dux540cux6781ux9650}{%
\section{59.5
条件稳态与无条件成功质量的不同极限}\label{ux6761ux4ef6ux7a33ux6001ux4e0eux65e0ux6761ux4ef6ux6210ux529fux8d28ux91cfux7684ux4e0dux540cux6781ux9650}}

\protect\hypertarget{p-qdd-04}{}{}

\textbf{P-QDD-04（有限成功概率衰减与初态责任）。} 对任意密度态
\(\tau\)，每个有限 \(n\) 均有 \(m_n(\tau)>0\)，但

\[
m_n(\tau)\longrightarrow0.
\]

对 \(\tau=\rho_*\)，\(c_{\rho_*}=8/5\)、\(b_{\rho_*}=1\)，故
\(m_n(\rho_*)=r^n\)。对均匀初态 \(\tau_0=I/2\)，

\[
c_{\tau_0}=\frac52,\qquad b_{\tau_0}=\frac{25}{16},\qquad
m_1(\tau_0)=\frac{33}{128}\ne\frac3{16}.
\]

\textbf{证明。} \(z_n(\tau)\) 收敛到有限的正数 \(b_\tau\)，而
\(0<r<1\)，故
\(r^nz_n\to0\)。特定初态的等式由实际矩阵元素与通用质量公式代入得到。\(\square\)

因此一个条件分布可以逐渐稳定，同时满足这个条件的有限记录越来越稀少。这里严格证明的是有限事件概率序列的极限；本章未构造无限路径测度或无限时间条件事件，也没有把条件极限称为无条件永久存在。

对原始 UEOT
的持续性追求，这个区分是必要的：一个组织结构在存活样本中越来越稳定，并不独自证明总体存活概率保持正下界。要讨论何种控制能提高寿命或改变成功率，还必须声明可用动作、成功事件、约束和代价。当前无控制可解模型并未选择自然界必须优化的目标函数。{[}S1；S2；S4；QM317{]}

\hypertarget{ux6b63ux8c31ux5ea6ux91cfux4e0dux80fdux4fddux8bc1ux552fux4e00ux5438ux5f15ux6001}{%
\section{59.6
正谱度量不能保证唯一吸引态}\label{ux6b63ux8c31ux5ea6ux91cfux4e0dux80fdux4fddux8bc1ux552fux4e00ux5438ux5f15ux6001}}

\protect\hypertarget{n-qdd-01}{}{}

\textbf{N-QDD-01（合法恒等 Doob 负控）。}
存在实际完整源仪器，其成功分支对非零 PSD
输入严格缩小迹，并具有正左特征度量，但其 Doob
通道为恒等；所有密度态都平稳，不存在共同的唯一吸引态。归一密度输入的成功概率恒为
\(1/4\)。

\textbf{构造与证明。} 在二维空间取一个成功 Kraus \(I/2\)、三个失败 Kraus
\(I/2\)，必要的其余下标补零。四个非零矩阵的 \(K^\dagger K\) 之和为
\(I\)，因此这确实是完整仪器。其成功分支为 \(\Phi(X)=X/4\)。取
\(H=I\)、\(S=T=I\) 和 \(r=1/4\)，左谱关系立即成立，按第56章实际 Doob
构造得到 \(D(X)=X\)。

于是每个密度态的整个 Doob
轨道都保持原态。选两个不同密度态，即得两个不同平稳态；若所有初态都收敛到同一个矩阵，常值轨道极限的唯一性会迫使这两个初态相等，矛盾。源的成功质量仍为
\((1/4)^n\)，实际有限条件态始终等于输入态。\(\square\)

这个反例保留了成功分支、完整仪器、严格衰减、正特征度量和实际 Doob
构造。它准确否定从这些条件单独推出唯一性或混合的推理。要得到一般收敛定理，还需明确的不可约性、本原性、收缩或谱条件及其证明，不能把原始
GOA 的唯一固定点图景作为所有对象的普遍结论。

\hypertarget{ux4eceux5177ux4f53ux52a8ux529bux5b66ux5185ux6838ux5230ux66f4ux4e00ux822cux7684ux7269ux7406ux6865}{%
\section{59.7
从具体动力学内核到更一般的物理桥}\label{ux4eceux5177ux4f53ux52a8ux529bux5b66ux5185ux6838ux5230ux66f4ux4e00ux822cux7684ux7269ux7406ux6865}}

本章完成的是一个指定非交换模型重复指定成功分支的有限动力学，以及相应
Doob
轨道和条件态的长时极限。一般正谱对象存在、谱半径识别、一般准稳态唯一性与无限条件过程仍有独立责任；它们不能从这个可解模型的公式外推。有限模型的物理实现、成功事件对应的实验读出及参数校准也仍需单独验证。

与前两章合看，当前内核获得了三种相互配合的能力：用实际纯态与偏迹表示任意有限复混态，用联合记录寄存器保存协议历史与量子记忆，用实际谱交织定量解释指定模型中的条件持续。这些都在已经声明的量子实现内成立。下一步应优先补齐选择原则之间的严格逻辑关系、可组合的记录误差预算和有明确一般条件的动力学推广，再把所得可观测量接到可实施实验。这样，UEOT
的解释目标会不断受到模型反例、形式证明和经验责任的共同约束。

\hypertarget{ch-60}{%
\chapter{60. 全效应分离与复模型的普遍局部层析}\label{ch-60}}

第50章从实际 Born
概率构造了全效应距离，并证明概率界、对称性、三角不等式和通道收缩；其中距离为零的结论准确停在``所有效应给出同样概率''。第55章则证明，在完整实量子模型中，即使知道全部局域产品效应的统计，也可能无法确定联合态。第57章给出了同一操作语言中的完整复模型，留下了它是否对所有联合态满足局部层析的正向责任。

本章在原有完整有限复密度空间中闭合这两项责任。关键不是直接把矩阵元素当作可测量量，而是先把任意
Hermitian
测试写成合法效应的实线性组合，再从真实效应概率恢复矩阵。联合情形对任意密度态成立，包括并非产品态的输入；所用信息是两侧效应的联合概率，包含相关性，并非只给两个边际。

\hypertarget{ux5b8cux6574ux72b6ux6001ux5b8cux6574ux6548ux5e94ux4e0eux5b9eux503cux914dux5bf9}{%
\section{60.1
完整状态、完整效应与实值配对}\label{ux5b8cux6574ux72b6ux6001ux5b8cux6574ux6548ux5e94ux4e0eux5b9eux503cux914dux5bf9}}

固定任意有限载体 \(\iota\)。本章的矩阵、状态和效应分别为

\[
\begin{aligned}
M_\iota&=M_\iota(\mathbb C),\\
\mathcal D(\iota)&=\{\rho\in M_\iota:\rho\succeq0,\ \operatorname{Tr}\rho=1\},\\
\mathcal E(\iota)&=\{E\in M_\iota:E\succeq0,\ I-E\succeq0\}.
\end{aligned}
\]

正性采用完整复矩阵的正半定锥。没有把态限制成对角态、向量态、选定的 Gram
因子输入或有限样本，也没有把效应限制为某个仪器库恰好能实现的测试。Lean
的 \texttt{Density} 与 \texttt{Effect}
是此前已经使用的实际类型；新增证明不改动它们的定义。

对任意矩阵 \(H,X\)，记实值迹配对

\[
\langle H,X\rangle_{\mathbb R}
=\operatorname{Re}\operatorname{Tr}(HX).
\]

效应概率就是
\(p_\rho(E)=\langle E,\rho\rangle_{\mathbb R}\)。这个配对对两边都具有实线性。把复矩阵乘以复数之后再取实部，不会得到一个以复数为值域的复线性概率函数，因此以下全部概率扩张只使用实系数。

正算符次序 \(H\preceq K\) 表示 \(K-H\) 正半定。对 Hermitian
\(H\)，有限矩阵的算符范数给出

\[
-\|H\|_{\rm op}I\preceq H\preceq\|H\|_{\rm op}I.
\]

这里的范数是复欧氏空间上的矩阵算符范数，Lean 在相应证明段显式使用
\texttt{Matrix.\allowbreak Norms.\allowbreak L2\allowbreak Operator}。上述次序界来自已锁定的有限矩阵 \(C^*\)
数学依赖；它不是额外的物理谱假设，也没有预先要求调用者提供分离效应。{[}QM318{]}

\hypertarget{ux6bcfux4e2a-hermitian-ux7b97ux7b26ux90fdux7531ux5408ux6cd5ux6548ux5e94ux5b9eux7ebfux6027ux5f20ux6210}{%
\section{60.2 每个 Hermitian
算符都由合法效应实线性张成}\label{ux6bcfux4e2a-hermitian-ux7b97ux7b26ux90fdux7531ux5408ux6cd5ux6548ux5e94ux5b9eux7ebfux6027ux5f20ux6210}}

\protect\hypertarget{p-sep-01}{}{}

\textbf{P-SEP-01（Hermitian 算符的合法效应表示）。} 对任意有限载体
\(\iota\) 和任意 Hermitian 矩阵 \(H\)，存在实数 \(r>0\) 与实际效应
\(E\in\mathcal E(\iota)\)，使

\[
H=2rE-rI.
\]

因此全部 Hermitian 矩阵都属于全部效应矩阵在实数域上的线性张成。

\textbf{证明。} 取 \(r=\|H\|_{\rm op}+1\)。由范数非负知
\(r>0\)；由上一节两个次序界，分别加上正半定的 \(I\)，得到

\[
H+rI\succeq0,\qquad rI-H\succeq0.
\]

定义

\[
E=\frac{H+rI}{2r}.
\]

系数 \((2r)^{-1}\) 严格为正，故
\(E\succeq0\)。还须验证效应的上界，而不能只验证正性。实际矩阵恒等式给出

\[
I-E=\frac{rI-H}{2r}\succeq0.
\]

这同时建立了 \texttt{Effect} 所要求的两个正性字段。由于
\(2r\ne0\)，重新整理定义即得 \(H=2rE-rI\)。单位矩阵本身是合法效应，因此
\(H\) 是两个效应矩阵的实线性组合。\(\square\)

这里``张成''不是说每个 Hermitian 矩阵都是效应的凸混合。系数 \(-r\)
可以为负，重建公式允许对不同效应的概率做有符号的实数运算；它并未把这个运算本身定义为一个新的随机混合实验。

形式化中的 \texttt{bounded\allowbreak Hermitian\allowbreak Effect}
先处理任意给定的双侧次序界；\texttt{bounded\allowbreak Hermitian\allowbreak Effect\_\allowbreak affine}
证明实际矩阵恒等式；\texttt{exists\_\allowbreak effect\_\allowbreak affine\_\allowbreak of\_\allowbreak hermitian} 从
\(H\) 本身导出所需界；\texttt{hermitian\_\allowbreak mem\_\allowbreak real\_\allowbreak span\_\allowbreak effects}
则把结论明确写进
\texttt{Submodule.\allowbreak span\ ℝ}。实系数性质是声明的一部分。{[}QM318{]}

这个构造在空载体上同样有意义：唯一矩阵为零，单位矩阵也为零，仍可取
\(r=1\)。后面关于密度态的前提则会自动排除空载体上的虚构归一态，见第60.6节。

\hypertarget{ux5168ux6548ux5e94ux540cux6982ux7387ux8febux4f7fux5b9eux9645ux5bc6ux5ea6ux6001ux76f8ux7b49}{%
\section{60.3
全效应同概率迫使实际密度态相等}\label{ux5168ux6548ux5e94ux540cux6982ux7387ux8febux4f7fux5b9eux9645ux5bc6ux5ea6ux6001ux76f8ux7b49}}

先从 P-SEP-01 得到一个可复用的中间步骤。若任意两个矩阵 \(X,Y\) 满足

\[
\langle E,X\rangle_{\mathbb R}
=\langle E,Y\rangle_{\mathbb R}
\quad\text{对全部 }E\in\mathcal E(\iota),
\]

则它们对任意 Hermitian \(H\) 的配对也相同。事实上，取 \(H=2rE-rI\)
后，实线性分别给出

\[
\langle H,X\rangle_{\mathbb R}
=2r\langle E,X\rangle_{\mathbb R}
-r\langle I,X\rangle_{\mathbb R},
\]

对 \(Y\) 同理。假设涵盖实际效应 \(E\) 和单位效应
\(I\)，故两式相等。这个中间步骤不要求 \(X,Y\) 为密度态，真正用到
Hermitian 性的是测试 \(H\) 的合法效应表示。

\protect\hypertarget{p-sep-02}{}{}

\textbf{P-SEP-02（全效应分离与距离零的矩阵含义）。} 对任意
\(\rho,\sigma\in\mathcal D(\iota)\)，

\[
\rho=\sigma
\quad\Longleftrightarrow\quad
\forall E\in\mathcal E(\iota),\ p_\rho(E)=p_\sigma(E).
\]

对第50章实际定义的全效应距离

\[
d_{\rm eff}(\rho,\sigma)
=\sup_{E\in\mathcal E(\iota)}
|p_\rho(E)-p_\sigma(E)|,
\]

相应有 \(d_{\rm eff}(\rho,\sigma)=0\) 当且仅当 \(\rho=\sigma\)。

\textbf{证明。} 相等的密度态当然给出相等概率。反过来，令
\(D=\rho-\sigma\)。由两态的正性知它们都是 Hermitian，因此 \(D\) 也是
Hermitian。以上实线性扩张可用 \(H=D\)，给出

\[
\langle D,\rho\rangle_{\mathbb R}
=\langle D,\sigma\rangle_{\mathbb R}.
\]

再对配对的右侧做减法，得到

\[
\operatorname{Re}\operatorname{Tr}(D^\dagger D)
=\operatorname{Re}\operatorname{Tr}(DD)=0.
\]

\(D^\dagger D\)
正半定，其迹的虚部为零。于是完整复数迹也为零。有限矩阵恒等式

\[
\operatorname{Tr}(D^\dagger D)
=\sum_{i,j}|D_{ij}|^2
\]

说明每个矩阵元素都为零，故 \(D=0\)，即 \(\rho=\sigma\)。形式证明使用
\texttt{Matrix.\allowbreak trace\_\allowbreak conj\allowbreak Transpose\_\allowbreak mul\_\allowbreak self\_\allowbreak eq\_\allowbreak zero\_\allowbreak iff}
完成这一步，之前已单独补足实部为零到完整迹为零的理由。

距离部分使用第50章已经证明的上确界性质：距离为零当且仅当每个效应的概率相同。接上刚才的分离结论即得双向等价。\(\square\)

新公开号
\texttt{Quantum\allowbreak Separation.\allowbreak effect\allowbreak Distance\_\allowbreak eq\_\allowbreak zero\_\allowbreak iff\_\allowbreak density\_\allowbreak eq}
与旧的 \texttt{UEOTQM.\allowbreak effect\allowbreak Distance\_\allowbreak eq\_\allowbreak zero\_\allowbreak iff}
分开命名。旧定理仍准确表达全效应同概率，新定理给出它与实际密度相等之间此前缺失的连接。这不需要改写旧的收缩、条件化或扰动证明，也不附带另一个矩阵范数的同一性声明。{[}QM318{]}

\hypertarget{ux4eceux56dbux4e2aux5408ux6cd5ux4ea7ux54c1ux6d4bux8bd5ux6269ux5f20ux5230-hermitian-ux4ea7ux54c1}{%
\section{60.4 从四个合法产品测试扩张到 Hermitian
产品}\label{ux4eceux56dbux4e2aux5408ux6cd5ux4ea7ux54c1ux6d4bux8bd5ux6269ux5f20ux5230-hermitian-ux4ea7ux54c1}}

现在固定两个任意有限载体 \(\iota,\nu\)，联合系统实际采用
\(\iota\times\nu\)。对局部效应 \(E,F\)，原有
\texttt{Real\allowbreak Local\allowbreak Tomography.\allowbreak product\allowbreak Effect} 的矩阵就是
\(E\otimes F\)。这个历史模块名并不限制其输入为实矩阵：定义本身接受任意复效应，实模型的限制由另一个实性字段单独承担。

产品效应的合法性来自 Kronecker 正性及恒等式

\[
I_{\iota\times\nu}-E\otimes F
=(I_\iota-E)\otimes I_\nu
+E\otimes(I_\nu-F).
\]

右边两项均正半定。因此这里的联合测试是一个具有完整正性和上界证明的实际效应。

设两个任意联合密度态 \(\rho,\sigma\)
对全部局部产品效应给出相同概率。给定任意局部 Hermitian 矩阵
\(H,K\)，分别用 P-SEP-01 写成

\[
H=2rE-rI_\iota,\qquad K=2sF-sI_\nu,
\qquad r,s>0.
\]

实际 Kronecker 乘法分配律给出四项实展开：

\[
\begin{aligned}
H\otimes K={}&4rs\,E\otimes F
-2rs\,E\otimes I_\nu\\
&-2rs\,I_\iota\otimes F
+rs\,I_\iota\otimes I_\nu.
\end{aligned}
\]

假设分别涵盖
\((E,F)\)、\((E,I_\nu)\)、\((I_\iota,F)\)、\((I_\iota,I_\nu)\)
四个真实局部测试组合。对四项使用实值配对的实线性，即得

\[
\langle H\otimes K,\rho\rangle_{\mathbb R}
=\langle H\otimes K,\sigma\rangle_{\mathbb R}.
\]

这就是 \texttt{hermitian\_\allowbreak product\_\allowbreak pairing\_\allowbreak eq} 的实际结论。四个测试随
\(H,K\) 的表示确定；证明没有假设 \(\rho\) 或 \(\sigma\)
自身可以分解为产品态的概率混合。展开的是测试算符，未对输入态做任何可分离假设。

\hypertarget{ux5b9eux90e8ux865aux90e8ux4e0eux5168ux90e8ux8054ux5408ux77e9ux9635ux5143ux7d20ux7684ux6062ux590d}{%
\section{60.5
实部、虚部与全部联合矩阵元素的恢复}\label{ux5b9eux90e8ux865aux90e8ux4e0eux5168ux90e8ux8054ux5408ux77e9ux9635ux5143ux7d20ux7684ux6062ux590d}}

令 \(e_{ij}\) 表示第 \((i,j)\) 项为一、其余项为零的矩阵单位。定义两个
Hermitian 矩阵

\[
S_{ij}=e_{ij}+e_{ji},\qquad
T_{ij}=\mathrm i(e_{ij}-e_{ji}).
\]

这里 \(\mathrm i\) 为虚数单位。转置共轭分别给出
\(S_{ij}^\dagger=S_{ij}\) 和 \(T_{ij}^\dagger=T_{ij}\)。对任意 Hermitian
\(X\)，设 \(X_{ij}=u+\mathrm i v\)，则
\(X_{ji}=u-\mathrm i v\)，从矩阵单位的迹配对可得

\[
\begin{aligned}
\langle S_{ij},X\rangle_{\mathbb R}
&=2\operatorname{Re}X_{ij},\\
\langle T_{ij},X\rangle_{\mathbb R}
&=2\operatorname{Im}X_{ij}.
\end{aligned}
\]

第二式的符号是正号：\(\operatorname{Tr}(T_{ij}X)=\mathrm i(X_{ji}-X_{ij})=2v\)。若改用
\(\mathrm i(e_{ji}-e_{ij})\)，恢复式的符号就必须反转。声明
\texttt{pairing\_\allowbreak imaginary\allowbreak Unit} 固定了本章方向，避免把两种约定混用。

两个局部载体上的上述矩阵单位还满足

\[
S_{(i,a),(j,b)}
=\frac12\bigl(S_{ij}\otimes S_{ab}
-T_{ij}\otimes T_{ab}\bigr),
\]

\[
T_{(i,a),(j,b)}
=\frac12\bigl(T_{ij}\otimes S_{ab}
+S_{ij}\otimes T_{ab}\bigr).
\]

逐项展开即可验证：第一式中交叉的 \(e_{ij}\otimes e_{ba}\) 与
\(e_{ji}\otimes e_{ab}\)
抵消，留下联合矩阵单位及其转置；第二式留下它们之差的 \(\mathrm i\)
倍。所有外部系数仍为实数。

这些恒等式不要求 \(i\ne j\) 或 \(a\ne b\)。例如
\(S_{ii}=2e_{ii}\)、\(T_{ii}=0\)；两个指标都重复时第一式仍恰好等于两倍联合对角矩阵单位，第二式两边都为零。对
Hermitian \(X\)，对角元的虚部本来也为零。这些边界包含在实际通用声明中。

\protect\hypertarget{p-clt-01}{}{}

\textbf{P-CLT-01（完整有限复模型的普遍局部层析）。} 对任意有限载体
\(\iota,\nu\) 和任意 \(\rho,\sigma\in\mathcal D(\iota\times\nu)\)，若

\[
\forall E\in\mathcal E(\iota),\ 
\forall F\in\mathcal E(\nu),\qquad
p_\rho(E\otimes F)=p_\sigma(E\otimes F),
\]

则 \(\rho=\sigma\)。因此第57章已经构造的
\texttt{Operational\allowbreak Systems.\allowbreak complex\allowbreak Model} 满足原接口中的完整
\texttt{Local\allowbreak Tomography} 原则。

\textbf{证明。} 第60.4节把产品效应的同概率扩张为所有局部 Hermitian
产品的同配对。给定任意联合行列 \((i,a),(j,b)\)，用 \(S\)
的产品恒等式与配对实线性，得到

\[
\langle S_{(i,a),(j,b)},\rho\rangle_{\mathbb R}
=\langle S_{(i,a),(j,b)},\sigma\rangle_{\mathbb R}.
\]

两态均 Hermitian，故矩阵单位恢复式推出相应联合矩阵元素的实部相等。再使用
\(T\)
的产品恒等式，以同样方式得到虚部相等。于是每个复矩阵元素相等，两个实际
\texttt{Density} 相等。

原 \texttt{complex\allowbreak Model}
的联合状态类型就是这里的完整联合密度空间，局部效应与产品概率也正是本章所用定义。将上述定理代入原有
\texttt{Local\allowbreak Tomography}
的全称量词，便得到所声明的模型原则，未另建一个较小状态空间的替代呈现。\(\square\)

这条结果说明，在该完整复矩阵模型中，全部产品效应的联合统计足以识别整个联合态。它没有要求先掌握一个全局纠缠效应，也没有把``只知道各侧边际''加强成不成立的层析前提。为特定设备设计有限测试表、估计抽样误差、控制病态重建或证明测量库完备性，仍是另有输入的任务。

\hypertarget{ux7a7aux8f7dux4f53ux5355ux5143ux7d20ux8f7dux4f53ux4e0eux4e0dux540cux5c40ux90e8ux7ef4ux6570}{%
\section{60.6
空载体、单元素载体与不同局部维数}\label{ux7a7aux8f7dux4f53ux5355ux5143ux7d20ux8f7dux4f53ux4e0eux4e0dux540cux5c40ux90e8ux7ef4ux6570}}

有限类型可以为空。空载体上迹的求和为零，因此不存在迹一密度态；具体见证
\texttt{empty\_\allowbreak carrier\_\allowbreak has\_\allowbreak no\_\allowbreak density}
保留了这个事实。通用定理接受任意有限载体，关于密度输入的结论在空态域上按逻辑正确成立，不需要添加一个隐藏的
\texttt{Nonempty} 前提来制造状态。

在单元素载体上，唯一矩阵元素就是迹。任意密度态的迹一条件因此迫使这个元素为一，任意两个密度态相等。见证
\texttt{singleton\_\allowbreak densities\_\allowbreak equal}
从实际迹一证明，而不是用一个未声明的非平凡二维测试。

两个局部载体的维数可以不同。本章的指标从一开始分成 \(\iota\) 与
\(\nu\)，产品载体为
\(\iota\times\nu\)，所有分配律和矩阵单位恒等式分别保持两组指标。\texttt{unequal\_\allowbreak local\_\allowbreak dimensions\_\allowbreak tomography}
将通用结论实际实例化到
\(\operatorname{Fin}2\times\operatorname{Fin}3\)。这个见证检验类型边界；通用结论仍对任意两个有限维数成立。

这些见证与本章通用定理分开计数。空态排除继承既有结果，单元素计算与
\(2\times3\) 实例分别承担边界核查，它们不是三条新的自然定律。{[}QM318{]}

\hypertarget{ux5bf9ux89d2ux6548ux5e94ux8d1fux63a7ux4e0eux5b8cux6574ux5b9eux6a21ux578bux7684ux7ee7ux627fux969cux788d}{%
\section{60.7
对角效应负控与完整实模型的继承障碍}\label{ux5bf9ux89d2ux6548ux5e94ux8d1fux63a7ux4e0eux5b8cux6574ux5b9eux6a21ux578bux7684ux7ee7ux627fux969cux788d}}

全效应分离的量词不能删成``某些常用效应''。为同时检查实、虚两类相干信息，取二维
Hermitian 对合

\[
X=\begin{pmatrix}0&1\\1&0\end{pmatrix},\qquad
Y=\begin{pmatrix}0&-\mathrm i\\\mathrm i&0\end{pmatrix}.
\]

它们满足
\(X^2=Y^2=I\)、\(\operatorname{Tr}X=\operatorname{Tr}Y=0\)。本章
\(T_{01}\) 的约定为 \(-Y\)，所以这里使用标准 \(Y\)
不改变第60.5节已经固定的恢复符号。

\protect\hypertarget{n-sep-01}{}{}

\textbf{N-SEP-01（全部对角效应仍不能识别相干态）。} 定义

\[
\rho_{x,\pm}=\frac12(I\pm X),\qquad
\rho_{y,\pm}=\frac12(I\pm Y).
\]

这四个矩阵都是实际密度态。对每一个合法对角效应 \(E\)，分别有

\[
p_{\rho_{x,+}}(E)=p_{\rho_{x,-}}(E),\qquad
p_{\rho_{y,+}}(E)=p_{\rho_{y,-}}(E).
\]

然而两对密度态各自不同；存在实际效应使每一对的概率分别为一与零，而且

\[
d_{\rm eff}(\rho_{x,+},\rho_{x,-})
=d_{\rm eff}(\rho_{y,+},\rho_{y,-})=1.
\]

\textbf{构造与证明。} 对任意 Hermitian 对合 \(A\)，\(P_\pm=(I\pm A)/2\)
是互补正交投影：直接乘法给出 \(P_\pm^2=P_\pm\)、\(P_+P_-=0\)，Hermitian
性与 \(P_\pm=P_\pm^\dagger P_\pm\) 给出正性。对
\(A=X,Y\)，迹零和二维单位矩阵的迹二进一步给出
\(\operatorname{Tr}P_\pm=1\)，故得到上述四个密度态。

每个态的两个对角元都为 \(1/2\)。任意对角效应写成
\(E=\operatorname{diag}(e_0,e_1)\)，其概率只由对角元决定：

\[
\operatorname{Re}\operatorname{Tr}(E\rho)
=\frac12(e_0+e_1).
\]

合法效应的 Hermitian 性保证 \(e_0,e_1\) 为实数，正性和补正性给出
\(0\le e_0,e_1\le1\)。所以等概率结论量化的是全部合法对角效应，而不是挑选出的两个例子。

另一方面取
\(E_x=(I+X)/2\)、\(E_y=(I+Y)/2\)。上面的互补正投影性质同时证明它们是合法效应。由
\(P_+^2=P_+\) 和 \(P_+P_-=0\)
得到各自的概率为一与零，直接证明两对态不同。每个全效应距离至少是这个合法测试的概率差一，而第50章已经证明它至多为一，故距离恰等于一。\(\square\)

对角测量家族在其声明的范围内完全合法，但它的统计等价类可以包含不同的完整密度态。这个反例同时保留了实非对角项与虚非对角项的丢失，避免仅检查一种相干方向就宣称测量库完备。

第55章的完整实模型障碍仍然成立，范围也没有缩成对角库。其两个合法实联合态为

\[
\omega_\pm=\frac14(I\pm Y\otimes Y).
\]

\(Y\otimes Y\) 虽由两个含虚部的局部矩阵组成，联合矩阵本身却是实的
Hermitian 对合。对任意局部实效应 \(E,F\)，实 Hermitian
性给出对称性，进而
\(\operatorname{Tr}(EY)=\operatorname{Tr}(FY)=0\)。因此

\[
p_{\omega_+}(E\otimes F)
=p_{\omega_-}(E\otimes F)
=\frac14\operatorname{Tr}E\operatorname{Tr}F.
\]

两个态仍由合法联合效应 \((I+Y\otimes Y)/2\)
以概率一与零区分；允许复局部效应 \((I+Y)/2\) 后，其两侧产品的概率也变为
\(1/2\)
与零。这些是第55章已编号、已证明的结果，在本章只作继承交叉核对。\texttt{full\_\allowbreak real\_\allowbreak model\_\allowbreak not\_\allowbreak local\allowbreak Tomography}
直接调用原
\texttt{Operational\allowbreak Systems.\allowbreak real\_\allowbreak not\_\allowbreak local\allowbreak Tomography}，未修改实模型字段，也不重新计为本章新增的编号反例。{[}QM316；QM317；QM318{]}

这解释了复局部层析证明中虚部测试的具体角色。第60.5节的局部 \(T\) 类
Hermitian
测试可以在复模型内写成合法效应的实组合；实模型不允许据此悄然加入含虚矩阵元的局部效应。两个模型共享真实的组合和迹运算，允许测试的范围仍有实际差别。

\hypertarget{ux516cux5f00ux53f7ux5b9eux9645ux8bc1ux636eux4e0eux4ecdux5f85ux5efaux7acbux7684ux7269ux7406ux6865}{%
\section{60.8
公开号、实际证据与仍待建立的物理桥}\label{ux516cux5f00ux53f7ux5b9eux9645ux8bc1ux636eux4e0eux4ecdux5f85ux5efaux7acbux7684ux7269ux7406ux6865}}

本章四个编号结果与实际代码按以下关系对应。表内前两组声明位于
\texttt{UEOTQM.\allowbreak Quantum\allowbreak Separation}，第三组位于
\texttt{UEOTQM.\allowbreak Complex\allowbreak Local\allowbreak Tomography}，最后一组见证位于
\texttt{UEOTQM.\allowbreak Quantum\allowbreak Separation\allowbreak Examples}；因此命名空间与局部公开号共同确定稳定名称。

\begin{longtable}[]{@{}
  >{\raggedright\arraybackslash}p{(\columnwidth - 2\tabcolsep) * \real{0.3000}}
  >{\raggedright\arraybackslash}p{(\columnwidth - 2\tabcolsep) * \real{0.7000}}@{}}
\toprule\noalign{}
\begin{minipage}[b]{\linewidth}\raggedright
编号结果或证明环节
\end{minipage} & \begin{minipage}[b]{\linewidth}\raggedright
公共声明与实际责任
\end{minipage} \\
\midrule\noalign{}
\endhead
\bottomrule\noalign{}
\endlastfoot
P-SEP-01：双侧正性与效应表示 &
\texttt{bounded\allowbreak Hermitian\allowbreak Effect}；\texttt{bounded\allowbreak Hermitian\allowbreak Effect\_\allowbreak matrix}；\texttt{bounded\allowbreak Hermitian\allowbreak Effect\_\allowbreak affine}；\texttt{exists\_\allowbreak effect\_\allowbreak affine\_\allowbreak of\_\allowbreak hermitian} \\
P-SEP-01：实数域张成 & \texttt{hermitian\_\allowbreak mem\_\allowbreak real\_\allowbreak span\_\allowbreak effects} \\
P-SEP-02：实扩张和密度分离 &
\texttt{pairing\_\allowbreak eq\_\allowbreak of\_\allowbreak all\_\allowbreak effect\_\allowbreak pairings\_\allowbreak eq}；\texttt{pairing\_\allowbreak sub\_\allowbreak right}；\texttt{density\_\allowbreak eq\_\allowbreak of\_\allowbreak all\_\allowbreak effect\_\allowbreak probabilities\_\allowbreak eq}；\texttt{density\_\allowbreak eq\_\allowbreak iff\_\allowbreak all\_\allowbreak effect\_\allowbreak probabilities\_\allowbreak eq} \\
P-SEP-02：距离零的实际含义 &
\texttt{effect\allowbreak Distance\_\allowbreak eq\_\allowbreak zero\_\allowbreak iff\_\allowbreak density\_\allowbreak eq} \\
P-CLT-01：局部矩阵单位 &
\texttt{symmetric\allowbreak Unit}；\texttt{imaginary\allowbreak Unit}；两者的
\texttt{\_\allowbreak hermitian}
定理；\texttt{pairing\_\allowbreak symmetric\allowbreak Unit}；\texttt{pairing\_\allowbreak imaginary\allowbreak Unit} \\
P-CLT-01：实产品展开 &
\texttt{sub\_\allowbreak kronecker}；\texttt{kronecker\_\allowbreak sub}；\texttt{symmetric\allowbreak Unit\_\allowbreak product}；\texttt{imaginary\allowbreak Unit\_\allowbreak product}；\texttt{hermitian\_\allowbreak product\_\allowbreak pairing\_\allowbreak eq} \\
P-CLT-01：任意联合态与原模型 &
\texttt{density\_\allowbreak eq\_\allowbreak of\_\allowbreak all\_\allowbreak product\_\allowbreak effect\_\allowbreak probabilities\_\allowbreak eq}；\texttt{complex\_\allowbreak local\allowbreak Tomography} \\
N-SEP-01：全部对角效应同统计 &
\texttt{diagonal\_\allowbreak effect\_\allowbreak pairing\_\allowbreak eq}；\texttt{all\_\allowbreak diagonal\_\allowbreak effects\_\allowbreak agree\_\allowbreak x}；\texttt{all\_\allowbreak diagonal\_\allowbreak effects\_\allowbreak agree\_\allowbreak y} \\
N-SEP-01：实际区分和距离 &
\texttt{x\allowbreak Test\_\allowbreak probability\_\allowbreak plus}、\allowbreak \texttt{x\allowbreak Test\_\allowbreak probability\_\allowbreak minus}；对应的
\texttt{y\allowbreak Test}
概率定理；\texttt{x\allowbreak Plus\_\allowbreak ne\_\allowbreak x\allowbreak Minus}、\allowbreak \texttt{y\allowbreak Plus\_\allowbreak ne\_\allowbreak y\allowbreak Minus}；\texttt{x\_\allowbreak effect\allowbreak Distance\_\allowbreak exact}、\allowbreak \texttt{y\_\allowbreak effect\allowbreak Distance\_\allowbreak exact} \\
维度边界与继承原则交叉核对 &
\texttt{empty\_\allowbreak carrier\_\allowbreak has\_\allowbreak no\_\allowbreak density}；\texttt{singleton\_\allowbreak densities\_\allowbreak equal}；\texttt{unequal\_\allowbreak local\_\allowbreak dimensions\_\allowbreak tomography}；\texttt{full\_\allowbreak complex\_\allowbreak model\_\allowbreak local\allowbreak Tomography}；\texttt{full\_\allowbreak real\_\allowbreak model\_\allowbreak not\_\allowbreak local\allowbreak Tomography} \\
\end{longtable}

两份新增公共模块含20条显式 theorem 与3项定义；具体见证单列19条
theorem、8项定义和1项缩写。阶段性的实际聚焦回执分别完整核对23个公共声明与28个见证声明的逐名公理输出，并绑定成功编译的源码字节。允许基础只有
\texttt{propext}、\allowbreak \texttt{Classical.\allowbreak choice}、\allowbreak \texttt{Quot.\allowbreak sound}。实际编译保留普通
linter 警告。相应记录位于
\texttt{physics/qm/kernel/evidence/v31.\allowbreak 8/A}；本轮统一验收、公开源码和远端发布状态由
v31.8
的统一回执另行承担。聚焦通过、统一验收与真实远端身份是不同的证据事实，定义、见证和编号结果也不是同一种计数。{[}QM318{]}

科学上，本章补足了已给复量子模型内部的可识别性：完整效应区分完整密度态，局部产品的完整联合统计区分完整联合态。对原有预测对象语言，这使``全效应响应相同''能够在明确的复矩阵实现中识别同一个密度对象；对受限仪器库，只能按其实际可用响应讨论相应商，N-SEP-01
保留了严格更粗的可能性。

本章的有限复向量空间、完整正半定态锥、完整效应锥、张量复合和迹概率属于已经声明的数学与量子扇区输入。其结果主要属于标准量子结构
P 及实际操作连接 B，具体概率承担
O；若要求自然界必须满足局部层析，则还须把这项选择原则及其背景按 S
登记并论证。冻结 Core
的106项与四生成器没有因此增加``必须为复数''的公理。本章也未证明这些候选原则唯一排除所有其他操作模型，未建立无限维或连续极限，未产生新的实验数据。

后续的物理桥应说明哪些 UEOT
机制选择这些状态、效应与复合规则，哪些真实设备实现所需测试，以及候选理论之间有哪些可区别的经验预测。本章为这项工作提供了完整且可检验的有限实现；它不把有限实现的完备性改称自然理论的最终完成。

\hypertarget{ch-61}{%
\chapter{61. 完整纯性与有限纯化的酉等价}\label{ch-61}}

第57章已经从任意有限复密度态的正性导出 Gram
因子，并用实际偏迹证明纯化存在。然而，只证明某些向量态是纯态，还不能比较所有按完整凸极点定义给出的纯化；只构造一个
Gram
纯化，也没有说明同一边际的其他纯化与它有什么关系。本章补上这两个责任，并分别处理相同辅助载体与不同辅助载体。

结论建立在已经明确的有限复密度模型中：完整凸极点恰为归一向量的外积；相同辅助载体上的任意两个纯化由实际辅助酉变换连接；不同载体可先经实际等距通道嵌入共同直和，再在这个共同载体上比较。最后两种说法的矩阵类型不同，不能合并成一个没有维数条件的``矩形酉矩阵''断言。{[}QM318{]}

\hypertarget{ux5b8cux6574ux72b6ux6001ux91cfux8bcdux4e0eux5f52ux4e00ux7ea6ux5b9a}{%
\section{61.1
完整状态量词与归一约定}\label{ux5b8cux6574ux72b6ux6001ux91cfux8bcdux4e0eux5f52ux4e00ux7ea6ux5b9a}}

继续使用

\[
\mathcal D(\iota)=\{\rho\in M_\iota(\mathbb C):\rho\succeq0,\ \operatorname{Tr}\rho=1\}.
\]

其中 \(\iota\) 是任意有限载体。\texttt{Density}
的两个条件是实际矩阵的正半定性和复迹一。\texttt{Is\allowbreak Pure}
沿用第57章对全部密度态、全部 \(0<t<1\)
的凸极点定义，没有缩小允许态集，也没有改为把向量表示作为定义字段。

对向量 \(v:\iota\to\mathbb C\)，记

\[
P_v=|v\rangle\langle v|,
\qquad (P_v)_{ij}=v_i\overline{v_j}.
\]

本章 Lean 声明中的归一条件写成 \(\operatorname{Tr}P_v=1\)。实际展开为
\(\sum_i|v_i|^2=1\)，即有限复欧几里得空间中的单位向量条件。这里的向量范数是二范数；不能把函数类型上某个不同的默认范数暗中代入。

对正矩阵 \(X\)，质量仍记作

\[
m(X)=\operatorname{Re}\operatorname{Tr}X\ge0.
\]

继承的正矩阵零质量定理给出 \(m(X)=0\iff X=0\)。只有在 \(m(X)>0\)
时，才使用合法归一密度态
\(\widehat X=X/m(X)\)。这一条件在下面的凸分解中实际检查，零余项不被赋予一个虚构的归一态。{[}QM317；QM318{]}

\hypertarget{ux6781ux70b9ux7684ux6bcfux4e2aux6b63ux8d28ux91cfux6b63ux5206ux91cfux90fdux6307ux5411ux540cux4e00ux6001}{%
\section{61.2
极点的每个正质量正分量都指向同一态}\label{ux6781ux70b9ux7684ux6bcfux4e2aux6b63ux8d28ux91cfux6b63ux5206ux91cfux90fdux6307ux5411ux540cux4e00ux6001}}

先证明一个直接反映完整极点量词的中间结论。设 \(\rho\) 为纯密度态，且

\[
\rho=X+Y,\qquad X,Y\succeq0,\qquad m(X)>0.
\]

则
\(\widehat X=\rho\)。由迹一，\(m(X)+m(Y)=1\)。证明分为两个实际可区分的情形。

如果 \(m(Y)=0\)，正性迫使 \(Y=0\)，因此
\(m(X)=1\)、\(X=\rho\)，归一不改变 \(X\)。如果 \(m(Y)>0\)，则
\(0<m(X)<1\)，而

\[
\rho=m(X)\widehat X+m(Y)\widehat Y
=m(X)\widehat X+(1-m(X))\widehat Y.
\]

两份归一态都属于完整 \(\mathcal D(\iota)\)，所以可直接把它们代入原有
\texttt{Is\allowbreak Pure} 定义，得到 \(\widehat X=\rho\) 和
\(\widehat Y=\rho\)。这里没有额外假定 \(X\)、\(Y\)
的秩、支撑或对角形式。Lean 定理
\texttt{Is\allowbreak Pure.\allowbreak normalize\_\allowbreak eq\_\allowbreak of\_\allowbreak positive\_\allowbreak summand}
逐项完成正性、严格正质量、归一重组和零余项分支。{[}QM318{]}

这个结论的作用是把极点性转成一个可以用于 Gram
列分解的约束。它不是把所需秩一结论放进假设；分量是否来自向量，还要在下一步由真实的正矩阵分解得到。

\hypertarget{ux4eceux4efbux610fux51f8ux6781ux70b9ux5bfcux51faux5f52ux4e00ux5411ux91cf}{%
\section{61.3
从任意凸极点导出归一向量}\label{ux4eceux4efbux610fux51f8ux6781ux70b9ux5bfcux51faux5f52ux4e00ux5411ux91cf}}

\protect\hypertarget{p-puc-01}{}{}

\textbf{P-PUC-01（完整有限复密度空间的纯性刻画）。} 对任意
\(\rho\in\mathcal D(\iota)\)，

\[
\operatorname{Pure}(\rho)
\iff
\exists v:\iota\to\mathbb C,\quad
\operatorname{Tr}P_v=1\ \land\ \rho=P_v.
\]

\textbf{证明。} 反向，即归一向量态为完整凸极点，已经由 P-PUR-01
证明。现在设 \(\rho\) 是按原定义给出的任意极点。从正性取得实际矩阵
\(A\)，使 \(\rho=AA^\dagger\)。对每一列 \(a\)，记
\(u_a(i)=A_{ia}\)，则逐项相乘给出

\[
\rho=AA^\dagger=\sum_a |u_a\rangle\langle u_a|
=\sum_a P_a,\qquad P_a\succeq0.
\]

取迹得到 \(\sum_a m(P_a)=1\)，所以至少有一列满足
\(m(P_a)>0\)。若所有列质量都不为正，其有限和至多为零，与迹一矛盾。对选中的一列，定义实际余项

\[
R=\sum_{b\ne a}P_b\succeq0.
\]

有限求和的删除重组给出 \(\rho=P_a+R\)。上一节的两分支结论于是推出
\(\rho=P_a/m(P_a)\)。令

\[
v=\frac{u_a}{\sqrt{m(P_a)}}.
\]

因为 \(m(P_a)>0\)，实平方根的平方等于
\(m(P_a)\)，且其实数倍在外积两侧各贡献一次系数。因此

\[
P_v=\frac1{m(P_a)}P_a=\rho,
\qquad\operatorname{Tr}P_v=1.
\]

这就从任意完整凸极点导出了归一向量，而非把向量作为调用者提供的见证。\(\square\)

Lean 中 \texttt{gram\_\allowbreak eq\_\allowbreak sum\_\allowbreak outer} 证明真实 Gram
矩阵的列外积和，\texttt{outer\_\allowbreak real\_\allowbreak smul}
证明实数缩放的平方律，\texttt{Is\allowbreak Pure.\allowbreak exists\_\allowbreak vector}
完成上述构造，\texttt{is\allowbreak Pure\_\allowbreak iff\_\allowbreak exists\_\allowbreak vector}
再接回既有正向定理。选择正质量列和 Gram
因子的存在过程使用已允许的经典基础，但不引入新的秩一公理或物理谱假设。

空载体也没有被隐藏地排除。若存在迹一密度态，上述总质量为一的论证本身就能推出一个正质量列；空载体上这不可能发生，因此该情形根本没有合法
\(\rho\)
可作为输入。定理无需额外添加一个用来掩盖此边界的非空字段。{[}QM318{]}

\hypertarget{ux76f8ux540c-gram-ux6570ux636eux600eux6837ux4ea7ux751fux771fux6b63ux7684ux9149ux53d8ux6362}{%
\section{61.4 相同 Gram
数据怎样产生真正的酉变换}\label{ux76f8ux540c-gram-ux6570ux636eux600eux6837ux4ea7ux751fux771fux6b63ux7684ux9149ux53d8ux6362}}

下一步先处理一个纯线性代数问题。设
\(F,G:\mathbb C^{\iota}\to\mathbb C^{\alpha}\) 是任意有限复矩阵，满足

\[
F^\dagger F=G^\dagger G.
\]

对任意输入 \(x\)，欧几里得平方范数由实际 Gram 二次型给出，因而

\[
\|Fx\|_2^2=x^\dagger F^\dagger Fx
=x^\dagger G^\dagger Gx=\|Gx\|_2^2.
\]

所以 \(Fx=0\) 蕴含 \(Gx=0\)。这使

\[
L:\operatorname{ran}F\longrightarrow\mathbb C^\alpha,
\qquad L(Fx)=Gx
\]

成为良定义的复线性映射：同一个 \(Fx\) 的两个表示之差在 \(\ker F\)
内，也就在 \(\ker G\) 内。范数等式进一步说明 \(L\)
是其实际生成支撑上的线性等距映射。

Lean 构造并未把 \(L\) 的良定义当作注释跳过。它先经 \(\ker F\)
商空间定义诱导映射，再使用商核到像的线性等价，把诱导映射转回
\(\operatorname{ran}F\)；随后证明每个像向量的范数被保持。\texttt{Linear\allowbreak Isometry.\allowbreak extend}
将这个子空间等距映射扩展为整个相同有限辅助空间上的线性等距映射。有限维自映射的单射性给出满射性，于是得到线性等距等价。

用标准正交基取其实际矩阵 \(U\)，得到

\[
U^\dagger U=I_\alpha,\qquad UU^\dagger=I_\alpha,
\qquad UF=G.
\]

两条酉恒等式都进入最终结论。这里不假定 \(F\)、\(G\)
满秩，也不要求它们占满辅助空间；未被生成支撑占据的方向由有限维等距扩张处理。\texttt{exists\_\allowbreak unitary\_\allowbreak of\_\allowbreak column\_\allowbreak gram}
的输入只有实际 Gram
等式，输入中没有所求酉矩阵、可扩张性字段或酉等价结论。{[}QM318{]}

这是一项数学存在构造。它使用所固定 Mathlib
中已证明的商空间和有限等距扩张定理，并由精确命名公理审计检查其基础；本章没有另外声称得到一个已经认证、带数值稳定性界的矩阵分解算法。

\hypertarget{ux884c-gramux8f6cux7f6eux548cux8f85ux52a9ux4f5cux7528ux7684ux65b9ux5411}{%
\section{61.5 行
Gram、转置和辅助作用的方向}\label{ux884c-gramux8f6cux7f6eux548cux8f85ux52a9ux4f5cux7528ux7684ux65b9ux5411}}

纯化向量沿用固定的向量化约定

\[
\psi(i,a)=A_{ia},\qquad\phi(i,a)=B_{ia}.
\]

实际系统边际分别为 \(AA^\dagger\) 和
\(BB^\dagger\)。如果两边际相同，得到的是\textbf{行 Gram} 等式

\[
AA^\dagger=BB^\dagger.
\]

本轮实际源码把这条等式转置，取
\(F=A^\mathsf T\)、\(G=B^\mathsf T\)。因为

\[
(A^\mathsf T)^\dagger A^\mathsf T
=\overline A A^\mathsf T=(AA^\dagger)^\mathsf T,
\]

上一节直接给出一个辅助酉矩阵，满足
\(UA^\mathsf T=B^\mathsf T\)。逐项读出便是

\[
B_{ia}=\sum_bU_{ab}A_{ib},
\qquad
\phi(i,a)=\sum_bU_{ab}\psi(i,b).
\]

因此真正作用在第二个向量指标上的矩阵就是 \(U\)，而矩阵因子的右乘写法为
\(B=AU^\mathsf T\)。相应联合向量关系是

\[
\phi=(I_\iota\otimes U)\psi.
\]

\texttt{exists\_\allowbreak auxiliary\_\allowbreak unitary\_\allowbreak of\_\allowbreak gram} 从转置后的真实 Gram
等式导出所需 \(U\)；\texttt{auxiliary\_\allowbreak action\_\allowbreak apply}
直接计算张量矩阵作用的每个条目；\texttt{auxiliary\_\allowbreak action\_\allowbreak vectorize}
明确证明右乘因子中的转置。这样，\(U\)、\(U^\dagger\)、\(U^\mathsf T\) 与
\(\overline U\) 不会在不同表示间被悄悄互换。

如果另从 \(F=A^\dagger\)
开始，也可以走一条等价的数学路线，但从支撑等距矩阵转回向量辅助作用时必须再跟踪共轭。本章实际实现采用上面的转置路线；准备阶段讨论过的其他路线不作为这份源码已经执行的证明记录。{[}QM318{]}

\hypertarget{ux540cux8f85ux52a9ux8f7dux4f53ux4e0aux7684ux4efbux610fux7eafux5316}{%
\section{61.6
同辅助载体上的任意纯化}\label{ux540cux8f85ux52a9ux8f7dux4f53ux4e0aux7684ux4efbux610fux7eafux5316}}

\protect\hypertarget{p-pun-01}{}{}

\textbf{P-PUN-01（同辅助载体的完整纯化酉等价）。} 设
\(\Psi,\Phi\in\mathcal D(\iota\times\alpha)\)
都是完整联合密度空间中的纯态，且实际偏迹满足

\[
\operatorname{Tr}_\alpha\Psi=\operatorname{Tr}_\alpha\Phi.
\]

则存在实际矩阵 \(U\in M_\alpha(\mathbb C)\)，使

\[
U^\dagger U=UU^\dagger=I_\alpha,
\qquad
\Phi=(I_\iota\otimes U)\Psi(I_\iota\otimes U)^\dagger.
\]

\textbf{证明。} 先对 \(\Psi\)、\(\Phi\) 分别使用 P-PUC-01，取得归一向量
\(\psi\)、\(\phi\)，使两联合密度为其真实外积。把原偏迹等式代入这些外积，得到对应矩阵因子的行
Gram 等式。上一节产生 \(U\)，使
\(\phi=(I\otimes U)\psi\)。最后用实际外积恒等式

\[
|K\psi\rangle\langle K\psi|
=K|\psi\rangle\langle\psi|K^\dagger
\]

得到联合密度的酉共轭关系。\(\square\)

这里的向量中间定理
\texttt{exists\_\allowbreak auxiliary\_\allowbreak unitary\_\allowbreak of\_\allowbreak equal\_\allowbreak marginal}
甚至不需要归一假设：任意两个同载体联合向量，只要实际外积偏迹相同，就具有上述精确向量作用关系。最终
\texttt{pure\_\allowbreak states\_\allowbreak unitarily\_\allowbreak equivalent} 则通过 P-PUC-01
处理所有按完整凸极点定义给出的纯密度，因而没有把``任意纯化''退化成``调用者已提供向量表示的纯化''。{[}QM318{]}

``唯一性''在这里是辅助酉作用下的等价性。它没有说
\(\Psi=\Phi\)，也没有说连接二者的 \(U\)
是唯一矩阵。未使用的辅助方向仍可容纳不同酉扩张；向量代表的整体相位也可吸收到辅助酉中。第57章的纯化存在与本节的等价结论合起来，才描述了同一有限复模型中``有纯化，且所有同辅助载体纯化彼此可比较''的完整逻辑链。

\hypertarget{ux4e0dux540cux8f7dux4f53ux5148ux4f7fux7528ux771fux6b63ux7684ux77e9ux5f62ux7b49ux8dddux901aux9053}{%
\section{61.7
不同载体先使用真正的矩形等距通道}\label{ux4e0dux540cux8f7dux4f53ux5148ux4f7fux7528ux771fux6b63ux7684ux77e9ux5f62ux7b49ux8dddux901aux9053}}

\protect\hypertarget{p-aux-01}{}{}

\textbf{P-AUX-01（辅助等距嵌入的实际通道、边际和纯性）。} 设
\(J:\mathbb C^\alpha\to\mathbb C^\beta\) 是实际矩形矩阵，满足
\(J^\dagger J=I_\alpha\)。则

\[
\mathcal E_J(X)=(I_\iota\otimes J)X(I_\iota\otimes J)^\dagger
\]

是从 \(\iota\times\alpha\) 到 \(\iota\times\beta\) 的实际单 Kraus CPTP
通道，并且对\textbf{任意复输入矩阵} \(X\)，

\[
\operatorname{Tr}_\beta\mathcal E_J(X)
=\operatorname{Tr}_\alpha X.
\]

如果输入是完整联合状态空间中的纯密度态，输出也是完整放大后状态空间中的纯密度态。

\textbf{证明。} 令 \(K=I_\iota\otimes J\)。张量乘法与共轭转置给出

\[
K^\dagger K=I_\iota\otimes(J^\dagger J)
=I_{\iota\times\alpha},
\]

所以这个实际 Kraus
矩阵满足矩形通道完备性。第58章的矩形通道理论因此适用，包括任意有限辅助放大上的完全正性。

固定系统指标 \(i,j\)，令
\(X^{ij}_{cd}=X_{(i,c),(j,d)}\)，这是辅助载体上的实际复矩阵。直接计算输出相应块，得到

\[
(\mathcal E_JX)^{ij}=JX^{ij}J^\dagger.
\]

实际偏迹取这个块的迹，利用迹循环，有

\[
\operatorname{Tr}(JX^{ij}J^\dagger)
=\operatorname{Tr}(J^\dagger JX^{ij})
=\operatorname{Tr}X^{ij}.
\]

逐个 \(i,j\) 即得边际恒等式。这个计算没有对 \(X\)
施加正性、迹一、纯态或产品形式限制。

最后，若 \(\Psi\) 为完整凸纯态，P-PUC-01 给
\(\Psi=P_\psi\)。通道输出等于
\(P_{K\psi}\)，归一由实际通道的迹保持得到，再由 P-PUR-01 或 P-PUC-01
的反向得知输出对全部放大后密度态仍为极点。\(\square\)

\texttt{auxiliary\allowbreak Embedding\allowbreak Channel}
真正构造了这个带输入、输出类型的通道；\texttt{auxiliary\allowbreak Embedding\allowbreak Channel\_\allowbreak operator\_\allowbreak apply}
给出块公式；\texttt{auxiliary\allowbreak Embedding\allowbreak Channel\_\allowbreak discard\_\allowbreak operator}
量化全部复矩阵；\texttt{auxiliary\allowbreak Embedding\allowbreak Channel\_\allowbreak discard\_\allowbreak state} 与
\texttt{auxiliary\allowbreak Embedding\allowbreak Channel\_\allowbreak preserves\_\allowbreak purity}
分别把它接回实际密度和完整纯性。其前提只要求
\(J^\dagger J=I_\alpha\)。当输出空间更大时，\(JJ^\dagger=I_\beta\)
一般不成立，不能把它加进结论。{[}QM318{]}

\hypertarget{ux4e0dux540cux8f85ux52a9ux7ef4ux6570ux7ecfux5171ux540cux76f4ux548cux540eux6bd4ux8f83}{%
\section{61.8
不同辅助维数经共同直和后比较}\label{ux4e0dux540cux8f85ux52a9ux7ef4ux6570ux7ecfux5171ux540cux76f4ux548cux540eux6bd4ux8f83}}

令共同辅助载体为不交和
\(\alpha\sqcup\beta\)；其复向量空间是两辅助空间的直和。定义实际矩形矩阵

\[
J_\alpha=\begin{bmatrix}I_\alpha\\0\end{bmatrix},
\qquad
J_\beta=\begin{bmatrix}0\\I_\beta\end{bmatrix}.
\]

两者分别满足 \(J_\alpha^\dagger J_\alpha=I_\alpha\) 和
\(J_\beta^\dagger J_\beta=I_\beta\)。注意两个等式右侧属于不同源载体。由它们分别构造

\[
\widetilde\Psi=\mathcal E_{J_\alpha}(\Psi),
\qquad
\widetilde\Phi=\mathcal E_{J_\beta}(\Phi).
\]

\protect\hypertarget{p-pun-02}{}{}

\textbf{P-PUN-02（不同有限辅助载体的共同扩展酉等价）。} 若
\(\Psi\in\mathcal D(\iota\times\alpha)\)、\(\Phi\in\mathcal D(\iota\times\beta)\)
都在各自完整联合密度空间中为纯态，且实际系统边际相同，则存在
\(U\in M_{\alpha\sqcup\beta}(\mathbb C)\)，满足两条酉恒等式，并有

\[
\widetilde\Phi
=(I_\iota\otimes U)\widetilde\Psi(I_\iota\otimes U)^\dagger.
\]

\textbf{证明。} 两个规范嵌入的等距恒等式由分块矩阵乘法直接得到。P-AUX-01
保证 \(\widetilde\Psi\)、\(\widetilde\Phi\)
是同一完整共同状态空间中的纯密度，并保持各自的原系统边际。原边际相同，因此可在共同辅助载体
\(\alpha\sqcup\beta\) 上应用 P-PUN-01。\(\square\)

形式定义 \texttt{auxiliary\allowbreak Inl}、\allowbreak \texttt{auxiliary\allowbreak Inr}
给出真实矩形嵌入；\texttt{extend\allowbreak Auxiliary\allowbreak Left}、\allowbreak \texttt{extend\allowbreak Auxiliary\allowbreak Right}
给出真实通道输出。四个纯性与边际定理先独立证明必要前提，\texttt{pure\_\allowbreak states\_\allowbreak common\_\allowbreak extension\_\allowbreak unitary}
才调用同载体酉等价定理。

本结论无需比较 \(|\alpha|\) 与
\(|\beta|\)，也不把它们当作同一类型。代价是共同载体维数采用
\(|\alpha|+|\beta|\)，没有声称它最小。它没有直接给出一个从任意指定一侧到另一侧的满空间等距映射，更没有给出维数不同时两侧都满足逆关系的``酉矩阵''。若要登记最小辅助维数、在额外维数不等式下的直接等距延拓，或某个量子重建体系要求的特定补齐形式，仍须另写相应类型和证明。{[}QM318{]}

\hypertarget{ux76f8ux540cux8fb9ux9645ux800cux8054ux5408ux6001ux4e0dux540cux7684ux7cbeux786eux6b63ux63a7}{%
\section{61.9
相同边际而联合态不同的精确正控}\label{ux76f8ux540cux8fb9ux9645ux800cux8054ux5408ux6001ux4e0dux540cux7684ux7cbeux786eux6b63ux63a7}}

本轮选取的辅助矩阵含有实、虚混合条目，以便检查共轭与转置的方向：

\[
U=\begin{bmatrix}3/5&4i/5\\4/5&-3i/5\end{bmatrix}.
\]

它的两列分别为 \((3,4)^\mathsf T/5\) 与
\(i(4,-3)^\mathsf T/5\)，每列平方范数为一，列间内积为零。实际见证同时验证
\(U^\dagger U=I\) 和 \(UU^\dagger=I\)。

取

\[
\psi_1=\tfrac35|00\rangle+\tfrac45|11\rangle,
\]

并按已经固定的第二指标作用计算

\[
\begin{aligned}
\psi_2=(I\otimes U)\psi_1
={}&\tfrac9{25}|00\rangle+\tfrac{12}{25}|01\rangle\\
&+\tfrac{16i}{25}|10\rangle-\tfrac{12i}{25}|11\rangle.
\end{aligned}
\]

对应向量化矩阵是

\[
A=\begin{bmatrix}3/5&0\\0&4/5\end{bmatrix},
\qquad
B=\begin{bmatrix}9/25&12/25\\16i/25&-12i/25\end{bmatrix}.
\]

精确相乘给出

\[
AA^\dagger=BB^\dagger
=\begin{bmatrix}9/25&0\\0&16/25\end{bmatrix}=\rho.
\]

例如 \(BB^\dagger\) 的非对角条目为
\(-144i/625+144i/625=0\)。两向量的总平方范数都为一，所以对应两联合密度均为完整凸纯态，其实际偏迹都为
\(\rho\)。

两联合向量和两联合密度确实不同：其 \(|00\rangle\langle00|\) 条目分别是
\(9/25\) 与
\(81/625\)。这正好说明边际相同不意味着联合密度相等；正确结论是存在辅助酉变换连接二者。\texttt{exact\_\allowbreak auxiliary\_\allowbreak vector\_\allowbreak action}
检查的是精确向量等式，\texttt{same\_\allowbreak actual\_\allowbreak marginal}
检查的是原有实际丢弃通道，\texttt{general\_\allowbreak uniqueness\_\allowbreak instantiated}
还实际调用一般完整纯态唯一性定理，避免只验证专门算出来的小矩阵。{[}QM318{]}

\hypertarget{ux4e09ux79cdux9519ux8befux4f5cux7528ux65b9ux5411ux7684ux8054ux5408ux5bc6ux5ea6ux8d1fux63a7}{%
\section{61.10
三种错误作用方向的联合密度负控}\label{ux4e09ux79cdux9519ux8befux4f5cux7528ux65b9ux5411ux7684ux8054ux5408ux5bc6ux5ea6ux8d1fux63a7}}

\protect\hypertarget{n-pun-01}{}{}

\textbf{N-PUN-01（伴随、转置和逐项共轭不能替代实际辅助作用）。}
对上一节固定的 \(U\) 与 \(\psi_1\)，若把正确辅助矩阵分别错换为
\(U^\dagger\)、\(U^\mathsf T\) 或
\(\overline U\)，得到的联合密度都不等于 \(P_{\psi_2}\)。

\textbf{证明。} 定义
\(\chi_T=(I\otimes T)\psi_1\)。以下三个实际复矩阵条目已经分别精确计算：

\begin{longtable}[]{@{}
  >{\raggedright\arraybackslash}p{(\columnwidth - 6\tabcolsep) * \real{0.2100}}
  >{\raggedright\arraybackslash}p{(\columnwidth - 6\tabcolsep) * \real{0.2000}}
  >{\raggedright\arraybackslash}p{(\columnwidth - 6\tabcolsep) * \real{0.2800}}
  >{\raggedright\arraybackslash}p{(\columnwidth - 6\tabcolsep) * \real{0.3100}}@{}}
\toprule\noalign{}
\begin{minipage}[b]{\linewidth}\raggedright
所用辅助矩阵 \(T\)
\end{minipage} & \begin{minipage}[b]{\linewidth}\raggedright
所比较的联合密度条目
\end{minipage} & \begin{minipage}[b]{\linewidth}\raggedright
\(P_{\chi_T}\) 的值
\end{minipage} & \begin{minipage}[b]{\linewidth}\raggedright
正确 \(P_{\psi_2}\) 的值
\end{minipage} \\
\midrule\noalign{}
\endhead
\bottomrule\noalign{}
\endlastfoot
\(U^\dagger\) & \(((0,0),(0,1))\) & \(108i/625\) & \(108/625\) \\
\(U^\mathsf T\) & \(((0,0),(0,1))\) & \(-108i/625\) & \(108/625\) \\
\(\overline U\) & \(((0,0),(1,0))\) & \(144i/625\) & \(-144i/625\) \\
\end{longtable}

每行的虚部均不同，因此各自的实际联合密度不相等。\(\square\)

\texttt{wrong\_\allowbreak adjoint\_\allowbreak joint\_\allowbreak density}、\allowbreak \texttt{wrong\_\allowbreak transpose\_\allowbreak joint\_\allowbreak density}
和 \texttt{wrong\_\allowbreak conjugate\_\allowbreak joint\_\allowbreak density}
是三条分别接受公理审计的负控定理，其条目值各有对应的精确计算定理。它们比较联合密度，因此即使某个错误向量只差整体相位，也不会被一个仅看向量表达式的表面差异误报为物理区分；本例的实际密度不等排除了这种遮蔽。

仅检查系统边际不足以发现这些方向错误。由 \(U\)
的酉恒等式及转置、共轭关系可知，这三种替代矩阵也都是辅助酉矩阵，因此其辅助作用仍保持系统边际。负控所要检验的是``是否得到指定的同一个联合纯化''，所以证据必须落在精确联合向量或联合密度上。它不声称边际测试可以识别被丢弃辅助系统上的酉选择。{[}QM318{]}

\hypertarget{ux4e8cux7ef4ux5230ux4e09ux7ef4ux7684ux5b9eux9645ux5d4cux5165ux4e0eux975eux9149ux8d1fux63a7}{%
\section{61.11
二维到三维的实际嵌入与非酉负控}\label{ux4e8cux7ef4ux5230ux4e09ux7ef4ux7684ux5b9eux9645ux5d4cux5165ux4e0eux975eux9149ux8d1fux63a7}}

取 \(\alpha=\mathrm{Fin}\,2\)、\(\beta=\mathrm{Unit}\)，则共同载体
\(\alpha\sqcup\beta\) 有三维。左规范嵌入的实际矩阵是

\[
J=\begin{bmatrix}1&0\\0&1\\0&0\end{bmatrix}.
\]

\protect\hypertarget{n-aux-01}{}{}

\textbf{N-AUX-01（严格辅助扩维是等距嵌入，不是双侧酉等价）。} 上述实际
\(3\times2\) 矩阵满足 \(J^\dagger J=I_2\)，但 \(JJ^\dagger\ne I_3\)。

\textbf{证明。} 前一等式由列正交直接得到。后一乘积是
\(\operatorname{diag}(1,1,0)\)；第三个对角条目为零，而 \(I_3\)
的相应条目为一。\(\square\)

见证 \texttt{strict\_\allowbreak embedding\_\allowbreak isometry} 与
\texttt{strict\_\allowbreak embedding\_\allowbreak not\_\allowbreak coisometry}
分别证明两项，\texttt{auxiliary\_\allowbreak dimensions\_\allowbreak differ}
证明这两个辅助载体的基数确实不同。把第61.9节的第二联合纯态沿此 \(J\)
嵌入，得到
\texttt{larger\allowbreak State}：它在六维联合空间、三维辅助载体上仍是完整凸纯态，并保持原系统边际。它与原二维辅助载体上的
\texttt{first\allowbreak State} 因而构成一对真正不同辅助维数的纯化。

\texttt{different\_\allowbreak auxiliary\_\allowbreak common\_\allowbreak unitary} 对这一对实际密度调用
P-PUN-02。其共同辅助载体为
\(\mathrm{Fin}\,2\sqcup(\mathrm{Fin}\,2\sqcup\mathrm{Unit})\)，维数是五；一般定理给出这个共同载体上满足两条酉恒等式的矩阵。这里的五维矩阵通过已经证明的存在构造得到，见证没有额外列出一张经过数值求解的五维酉矩阵，也没有将两个原始辅助维数混写为相同维数。{[}QM318{]}

\hypertarget{ux6e90ux7801ux9010ux58f0ux660eux5ba1ux8ba1ux4e0eux8bc1ux636eux8303ux56f4}{%
\section{61.12
源码、逐声明审计与证据范围}\label{ux6e90ux7801ux9010ux58f0ux660eux5ba1ux8ba1ux4e0eux8bc1ux636eux8303ux56f4}}

本章新增公共源码为 \texttt{UEOTQM/Pure\allowbreak State\allowbreak Characterization.\allowbreak lean} 与
\texttt{UEOTQM/Purification\allowbreak Uniqueness.\allowbreak lean}，具体见证为
\texttt{scripts/witnesses/Purification\allowbreak Uniqueness\allowbreak Examples.\allowbreak lean}。它们分别包含：

\begin{longtable}[]{@{}
  >{\raggedright\arraybackslash}p{(\columnwidth - 8\tabcolsep) * \real{0.2000}}
  >{\raggedleft\arraybackslash}p{(\columnwidth - 8\tabcolsep) * \real{0.2000}}
  >{\raggedleft\arraybackslash}p{(\columnwidth - 8\tabcolsep) * \real{0.2000}}
  >{\raggedleft\arraybackslash}p{(\columnwidth - 8\tabcolsep) * \real{0.2000}}
  >{\raggedleft\arraybackslash}p{(\columnwidth - 8\tabcolsep) * \real{0.2000}}@{}}
\toprule\noalign{}
\begin{minipage}[b]{\linewidth}\raggedright
源文件
\end{minipage} & \begin{minipage}[b]{\linewidth}\raggedleft
显式 theorem
\end{minipage} & \begin{minipage}[b]{\linewidth}\raggedleft
定义 def
\end{minipage} & \begin{minipage}[b]{\linewidth}\raggedleft
缩写 abbrev
\end{minipage} & \begin{minipage}[b]{\linewidth}\raggedleft
精确命名公理审计
\end{minipage} \\
\midrule\noalign{}
\endhead
\bottomrule\noalign{}
\endlastfoot
\texttt{Pure\allowbreak State\allowbreak Characterization.\allowbreak lean} & 5 & 0 & 0 & 5 \\
\texttt{Purification\allowbreak Uniqueness.\allowbreak lean} & 22 & 5 & 0 & 27 \\
\texttt{Purification\allowbreak Uniqueness\allowbreak Examples.\allowbreak lean} & 28 & 10 & 0 & 38 \\
\end{longtable}

两个公共模块合计27条显式定理、5个定义；见证另有28条定理、10个定义。公共与见证共70个不同稳定命名声明，在本路实际聚焦审计中全部只依赖
\texttt{propext}、\allowbreak \texttt{Classical.\allowbreak choice}、\allowbreak \texttt{Quot.\allowbreak sound}。见证的每个
theorem、def、abbrev 按源码顺序打印公理依赖，没有只抽查最后一个总定理。

实际接受日志、最终源码摘要、源码与 olean 绑定及逐名库存分别登记在
\texttt{physics/qm/kernel/evidence/v31.\allowbreak 8/B/focused-verification.\allowbreak json} 和
\texttt{declaration\_\allowbreak inventory.\allowbreak json}。聚焦模块验证对应
\texttt{pure-005}、\allowbreak \texttt{uniqueness-005}，最终具体见证对应
\texttt{witness-003}，公共精确审计对应
\texttt{public-audit-001}。这些都是实际退出零的调用；早先有意义的失败及其可取得源码快照被另行保留，不计入接受结果。

这份本路回执的层级是聚焦证明验收。版本统一构建、全量公理审计、继承见证及远端
CI
的完成与身份由主版本回执另行登记，不能用本地结果替代尚未实际回读的远端运行。普通
linter
提示与逻辑验收状态也应分开记录，不能把允许的普通警告改写为``零警告''。这里没有新增真实实验；有理复矩阵正负控是形式化见证，不是自然界的实验确证。{[}QM318{]}

\hypertarget{ux5df2ux95edux5408ux7684ux6570ux5b66ux8d23ux4efbux4e0eux4ecdux5f00ux653eux7684ux7406ux8bbaux9009ux62e9}{%
\section{61.13
已闭合的数学责任与仍开放的理论选择}\label{ux5df2ux95edux5408ux7684ux6570ux5b66ux8d23ux4efbux4e0eux4ecdux5f00ux653eux7684ux7406ux8bbaux9009ux62e9}}

第57章的存在构造、本章的完整纯性反向、同辅助载体酉等价与不同载体共同扩展，现已在同一个有限复密度模型中接通。纯态定义、Gram
因子、联合向量、实际偏迹和实际矩形通道各有清晰接口，结论不再依赖``所有纯态已经被假定为向量态''或``所求辅助酉由调用者给出''这类未完成责任。

与此同时，辅助酉矩阵本身的唯一性、共同扩展的最小维数、额外维数条件下的直接等距延拓及认证数值构造仍与已完成结论不同。无限维系统中的闭像、非有限秩与可扩张性问题也不能直接沿用本章的有限维满射论证。若某个量子重建原则对辅助补齐、允许的可逆过程或系统同构有更具体要求，必须把那个原则逐字段对接到当前已证明的类型，不能仅凭``本质唯一性''这一名称视为已经满足全部文献合同。

这些结果主要属于在标准有限复量子结构 P 与真实操作连接 B
下的数学结论。它们加强了这个候选模型的内部完备性和可检验接口，但从冻结的106项
Core
与四生成器唯一选择复量子结构的桥仍未建立。完整经典模型的纯化障碍与实模型的局部层析反例继续约束候选原则的选择责任；排除这两个已检查模型，不等于排除所有可能自然理论。

三代结论继续依赖原有谱、表示、阈值与有效窗口等前提；本章既不修改那些前提，也不产生新的经验数据。UEOT
的目标仍是把理论思想、实际推导、反例与经验责任连成可审查的整体。本章完成的是其中一段有限量子数学责任，下一步的自然理论主张仍须提出独立前提、实际证明及可检验的物理对应。{[}QM318{]}

\hypertarget{ch-62}{%
\chapter{62. 保留历史的逐动作记录误差预算}\label{ch-62}}

第58章给出了真实仪器、全部有限自适应叶记录与量子记忆的联合通道，并用该通道的整体缺陷控制完整记录
TV。整体缺陷说明怎样度量两个协议的差异，但通常仍须进一步求出或估计它。本章在固定有限深度和同一个因果控制器下，把误差追到每一步实际选择的仪器：先构造读取历史的矩形
Kraus 通道，再从共同输入的 PSD
主块、真实块权重和合法块效应导出动作预算。

这种控制器实际读取并去相干历史寄存器。形式化通道可以作用于任意历史---记忆密度矩阵，证明也确实使用所有输入主块；这不等于把经典历史控制改成相干量子控制。原有有限树解释器继续保持其更一般的变长叶语义，本章通过一个明确的固定深度子类与它连接。{[}QM318{]}

\hypertarget{ux6709ux9650ux65f6ux95f4ux771fux5b9eux5386ux53f2ux4e0eux540cux4e00ux4e2aux63a7ux5236ux5668}{%
\section{62.1
有限时间、真实历史与同一个控制器}\label{ux6709ux9650ux65f6ux95f4ux771fux5b9eux5386ux53f2ux4e0eux540cux4e00ux4e2aux63a7ux5236ux5668}}

固定有限结果载体 \(O\)、有限复量子记忆载体 \(\iota\) 和有限深度
\(N\in\mathbb N\)。第 \(t\) 步之前的完整结果历史为

\[
\mathsf H_t=\{h:\operatorname{Fin}(t)\to O\}.
\]

这里的下标就是已记录结果数；\(\mathsf H_0\) 只有空历史。动作载体记为
\(A\)。有限深度只要求控制器在相应时间和历史上给出动作，不要求 \(A\)
本身有限。实际类型为

\[
C:\prod_{t:\operatorname{Fin}(N)}(\mathsf H_t\to A).
\]

两种实现分别使用 \(I_a\) 与 \(J_a\)，但都把同一个 \(C_t(h)\)
当作该历史上的动作。共同结果字母表并不足以蕴含共同控制器；这里两者使用的是同一个函数值。每个仪器都有完整
Kraus 家族，例如

\[
\Phi^I_{a,o}(X)=\sum_{k:\kappa}K^I_{a,o,k}X(K^I_{a,o,k})^\dagger,
\qquad
\sum_{o,k}(K^I_{a,o,k})^\dagger K^I_{a,o,k}=I_\iota.
\]

\(J\) 的 Kraus 下标可以来自另一个有限载体
\(\eta\)。本章比较真实通道，不要求两种实现具有相同的 Kraus
表示。保留结果的单动作联合通道仍是第58章的

\[
\mathcal R^I_a(X)=\sum_o|o\rangle\langle o|\otimes\Phi^I_{a,o}(X).
\]

记 \(d_{\rm eff}\) 为全效应距离，并定义实际动作记录缺陷

\[
\delta_{\rm rec}(I_a,J_a)
=\sup\left(\{0\}\cup
 \{d_{\rm eff}(\mathcal R^I_a(\rho),\mathcal R^J_a(\rho)):
   \rho\in\mathcal D(\iota)\}\right).
\]

上确界的同一输入量词不能删去，记录标签也不能在定义缺陷前丢弃。插入零保证空输入态类型下定义仍有明确非空集合；给历史通道缺陷设统一上界时，源码显式保留该上界非负的条件。

\hypertarget{psd-ux5386ux53f2ux4e3bux5757ux771fux5b9eux6743ux91cdux4e0eux5408ux6cd5ux6548ux5e94}{%
\section{62.2 PSD
历史主块、真实权重与合法效应}\label{psd-ux5386ux53f2ux4e3bux5757ux771fux5b9eux6743ux91cdux4e0eux5408ux6cd5ux6548ux5e94}}

对任意有限历史载体 \(H\) 和任意
\(X\in M_{H\times\iota}(\mathbb C)\)，定义主块

\[
X_h(i,j)=X((h,i),(h,j)).
\]

它不是取出某些任意矩阵元后另行宣称为状态。将嵌入
\(V_h:e_i\mapsto e_{(h,i)}\) 写出，便有
\(X_h=V_h^\dagger XV_h\)。同一个主子矩阵构造也用于输出效应。

\protect\hypertarget{p-arb-01}{}{}

\textbf{P-ARB-01（历史块、归一权重与合法块效应）。} 若
\(\zeta\in\mathcal D(H\times\iota)\)，则每个主块 \(\zeta_h\) 都是
PSD，且

\[
p_h=\operatorname{Tr}\zeta_h\ge0,
\qquad \sum_h p_h=1,
\qquad p_h=0\ \Longleftrightarrow\ \zeta_h=0.
\]

若 \(E\in\mathcal E(H\times\nu)\) 是任意合法输出效应，则主块 \(E_h\)
也是 \(\mathcal E(\nu)\) 中的合法效应。对任意块族
\(B_h\)，还有真实迹配对等式

\[
\left\langle E,\operatorname{diag}_hB_h\right\rangle_{\mathbb R}
=\sum_h\langle E_h,B_h\rangle_{\mathbb R}.
\]

\textbf{证明。} PSD 在主子矩阵下保持，因此
\(\zeta_h\succeq0\)。总迹按积指标求和，给出
\(\sum_h\operatorname{Tr}\zeta_h=\operatorname{Tr}\zeta=1\)。对 PSD
矩阵，迹为零当且仅当矩阵为零；这一已有定理直接用于每个主块，未给零权重历史强行指定一个后验态。

对效应，\(E\succeq0\) 与 \(I-E\succeq0\) 分别限制到历史 \(h\)。由于
\(i\mapsto(h,i)\) 是单射，单位矩阵的主块就是 \(I_\nu\)，故
\(E_h\succeq0\) 且
\(I_\nu-E_h\succeq0\)。最后展开矩阵乘法和迹，输出块的非对角历史元素为零，剩余项恰为主块配对之和。\(\square\)

这里 \(\zeta\)
可以有不同历史间的非零相干项。主块正性与迹归一性不要求输入事先块对角。下一节构造的通道会实际消去这些历史相干项；证明范围由该明确动作决定。

\hypertarget{ux4eceux771fux5b9eux8bfbux53d6-kraus-ux5230ux52a0ux6743ux52a8ux4f5cux754c}{%
\section{62.3 从真实读取 Kraus
到加权动作界}\label{ux4eceux771fux5b9eux8bfbux53d6-kraus-ux5230ux52a0ux6743ux52a8ux4f5cux754c}}

先考虑一般的历史依赖矩形通道 \(\mathcal C_h:M_\iota\to M_\nu\)。若其
Kraus 矩阵为 \(L_{h,k}\)，在历史与记忆联合空间中定义

\[
\widetilde L_{h,k}((g,a),(g',i))
=\mathbf 1_{g=h}\mathbf 1_{g'=h}\,L_{h,k}(a,i).
\]

这个矩阵读取输入历史 \(h\)、对记忆执行 \(\mathcal C_h\)，并把同一个
\(h\) 写入输出。对不同 \(h\) 的 Kraus 求和是通道的非相干 Kraus
求和，因而是历史去相干的真实实现。

\protect\hypertarget{p-arb-02}{}{}

\textbf{P-ARB-02（实际历史读取通道的动作加权界）。} 上述 Kraus
家族构成迹保持完全正通道
\(\mathcal S_{\mathcal C}\)，对每个实际矩阵输入都有

\[
\mathcal S_{\mathcal C}(X)
=\operatorname{diag}_h\mathcal C_h(X_h).
\]

若 \(\mathcal D_h\) 是另一组具有相同输入输出载体的通道，则任意
\(\zeta\in\mathcal D(H\times\iota)\) 满足

\[
d_{\rm eff}(\mathcal S_{\mathcal C}(\zeta),
             \mathcal S_{\mathcal D}(\zeta))
\le \sum_h p_h\,\delta(\mathcal C_h,\mathcal D_h),
\qquad p_h=\operatorname{Tr}\zeta_h.
\]

特别地，若所有 \(h\) 上的实际缺陷都不超过
\(\varepsilon\)，则该同输入距离不超过 \(\varepsilon\)。若还显式要求
\(\varepsilon\ge0\)，历史通道的实际缺陷也不超过
\(\varepsilon\)，包括空输入载体的上确界边界。

\textbf{证明。} 对固定输入指标 \((g,i),(g',j)\)，矩阵积
\(\widetilde L_{h,k}^\dagger\widetilde L_{h,k}\) 只有 \(g=g'=h\)
时非零。先对 \(k\) 使用各 \(\mathcal C_h\) 的完整性，再对 \(h\)
求和，得到联合单位矩阵。矩形通道的完全正性由已证的有限辅助放大定理给出。类似展开
\(\widetilde L_{h,k}X\widetilde L_{h,k}^\dagger\)，只有 \(X_h\)
进入对应输出块，给出作用公式。

对 PSD 块 \(X_h\)，分两种情况。若 \(p_h=0\)，P-ARB-01 给出
\(X_h=0\)，两个线性通道输出均为零。若 \(p_h>0\)，令
\(\rho_h=X_h/p_h\)，它具有实际 PSD 与迹一证明。对任意合法效应
\(F\)，由通道实线性和缺陷定义，

\[
\begin{aligned}
&\left|\langle F,\mathcal C_h(X_h)\rangle_{\mathbb R}
 -\langle F,\mathcal D_h(X_h)\rangle_{\mathbb R}\right|\\
&\quad=p_h\left|p_{\mathcal C_h(\rho_h)}(F)
                 -p_{\mathcal D_h(\rho_h)}(F)\right|\\
&\quad\le p_h\delta(\mathcal C_h,\mathcal D_h).
\end{aligned}
\]

再取任意联合输出效应 \(E\)。P-ARB-01 保证每个 \(E_h\)
合法，并把整个概率差展开成主块概率差的和。绝对值三角不等式给出所述加权界；最后对全部效应取上确界。统一界使用
\(\sum_hp_h=1\)，历史通道缺陷界再对所有输入态取上确界并检查所插入的零。\(\square\)

关键权重来自同一个
\(\zeta\)。没有把两个不同演化轨迹的历史分布假定相同；本结果为下一节的三角分解提供一个可合法使用的共同输入。

\hypertarget{ux8ffdux52a0ux7ed3ux679cux7684ux771fux5b9eux65f6ux95f4ux987aux5e8fux4e0eux539fux534fux8bae}{%
\section{62.4
追加结果的真实时间顺序与原协议}\label{ux8ffdux52a0ux7ed3ux679cux7684ux771fux5b9eux65f6ux95f4ux987aux5e8fux4e0eux539fux534fux8bae}}

对长度 \(t\) 的历史 \(h\) 和新结果 \(o\)，记 \(h\mathbin{\|}o\) 为把
\(o\) 追加到末尾的长度 \(t+1\) 历史。真实载体双射为

\[
\mathsf H_t\times(O\times\iota)
\longleftrightarrow\mathsf H_{t+1}\times\iota,
\qquad(h,(o,i))\longmapsto(h\mathbin{\|}o,i).
\]

逆映射取前 \(t\) 项和最后一项。这一双射由一个矩形置换 Kraus
矩阵实现，其逆双射也给出真实通道。因此重标记保持完整态，并保持全效应距离，而不只是保持迹。

第 \(t\) 个实际步骤 \(S^I_t\) 先读取历史，使用
\(\mathcal C_h=\mathcal R^I_{C_t(h)}\)，再经上述置换追加结果。于是

\[
S^I_t(X)_{h\|o,h\|o}
=\Phi^I_{C_t(h),o}(X_h),
\]

不同完整历史之间的输出块为零。初态通过
\(\iota\simeq\mathsf H_0\times\iota\)
的真实置换嵌入；依次执行这些步骤得到联合密度态
\(\zeta^I_{C,\rho}\)。零深度没有动作调用，不要求动作载体非空。

\protect\hypertarget{p-arb-03}{}{}

\textbf{P-ARB-03（有限历史演化与原完整协议的矩阵一致性）。} 任意固定
\(N\)、控制器 \(C\)、实际仪器族 \(I\) 和初态 \(\rho\)
所构造的真实输出满足

\[
\zeta^I_{C,\rho}
=\operatorname{diag}_{h\in\mathsf H_N}
 \operatorname{recordOperator}(I,\operatorname{word}_C(h),\rho).
\]

从同一个 \(C\) 构造原类型中的协议 \(\pi_C\)，则存在实际叶双射
\(e_C:\operatorname{Leaf}(\pi_C)\simeq\mathsf H_N\)，并有

\[
\operatorname{word}_{\pi_C}(\ell)
=\operatorname{word}_C(e_C\ell).
\]

经过 \(e_C\times\operatorname{id}_\iota\)
的真实置换，两份联合密度态相等；相应叶概率、完整记录 TV
和终端量子边际全部一致。量子边际具体为

\[
\operatorname{Tr}_{\mathsf H_N}\zeta^I_{C,\rho}
=\sum_{\ell:\operatorname{Leaf}(\pi_C)}
 \operatorname{recordOperator}(I,\operatorname{word}_{\pi_C}(\ell),\rho).
\]

\textbf{证明。} 原协议从词的首项递归：停止对应空历史，动作节点对全部
\(o\in O\)
给出后继控制器。叶到历史的映射将当前结果放到首项并递归处理尾部；逆映射读取首项，再递归还原尾叶。两个复合分别等于恒等，由此得到真实双射。

通道执行采用逐步末尾追加，因此还须证明两种递归顺序相符。源码先证明

\[
\operatorname{word}_{C}(h\|o)
=\operatorname{word}_{C_{<t}}(h)
 \mathbin{+\!+}[(C_t(h),o)],
\]

其中 \(C_{<t}\) 是去掉最后动作时刻的同一控制器。另对原解释器证明

\[
\operatorname{recordOperator}(I,w\mathbin{+\!+}v,X)
=\operatorname{recordOperator}
  (I,v,\operatorname{recordOperator}(I,w,X)).
\]

零深度时置换嵌入保留原矩阵。归纳步把真实 \(S^I_t\)
的主块公式、归纳假设和上述两条词顺序等式连接起来，就得到完整输出矩阵，而非仅有概率相等。最后用叶双射换序有限和，得到原协议全态置换等式及量子边际。\(\square\)

原 \texttt{Protocol} 允许不同叶具有不同长度，本章的 \(\pi_C\)
则是由同一个固定 \(N\)
构造的完整子类。零概率结果仍保留其合法后继动作。这里没有补齐任意停止树的语义证明，也没有借由叶类型转换自动取得那项更广范围。

\hypertarget{ux9010ux52a8ux4f5cux9012ux63a8ux6709ux9650ux603bux548cux4e0eux5b8cux6574ux8bb0ux5f55-tv}{%
\section{62.5 逐动作递推、有限总和与完整记录
TV}\label{ux9010ux52a8ux4f5cux9012ux63a8ux6709ux9650ux603bux548cux4e0eux5b8cux6574ux8bb0ux5f55-tv}}

对于同一 \(C\)，把两种初态分别记为 \(\rho,\sigma\)。令 \(\zeta^I_t\)
是从 \(\rho\) 经 \(I\) 和前 \(t\)
步控制器所得的实际联合态，\(\zeta^J_t\) 则从 \(\sigma\) 经 \(J\)
得到。定义

\[
\begin{aligned}
q^J_t(h)&=\operatorname{Tr}(\zeta^J_t)_h,\\
a_t&=\sum_{h\in\mathsf H_t}q^J_t(h)
 \delta_{\rm rec}(I_{C_t(h)},J_{C_t(h)}).
\end{aligned}
\]

这就是源码的 \texttt{action\allowbreak Cost}。其中前缀态由
\texttt{run\ J\ (C.\allowbreak truncate\ t)\ σ} 实际构造；\(a_t\)
非负，零质量历史按零块定理贡献零。它不是重新命名整个协议的缺陷。

\protect\hypertarget{p-arb-04}{}{}

\textbf{P-ARB-04（实际逐动作和预算与完整记录 TV）。}
在上述固定有限深度、共同控制器与共同记录/记忆载体条件下，对任意初态
\(\rho,\sigma\) 有

\[
\begin{aligned}
d_{\rm eff}(\zeta^I_N,\zeta^J_N)
&\le d_{\rm eff}(\rho,\sigma)+\sum_{t=0}^{N-1}a_t,\\
\operatorname{TV}(p^I_N,p^J_N)
&\le d_{\rm eff}(\rho,\sigma)+\sum_{t=0}^{N-1}a_t.
\end{aligned}
\]

若对每个时间 \(t\)、每个长度 \(t\) 历史 \(h\) 已实际证明

\[
\delta_{\rm rec}(I_{C_t(h)},J_{C_t(h)})\le\varepsilon_t,
\]

则两式右端都可换成
\(d_{\rm eff}(\rho,\sigma)+\sum_{t<N}\varepsilon_t\)。\(p^I_N,p^J_N\)
是完整历史分布；经 P-ARB-03 的叶双射后，同样的界适用于原协议的完整叶律。

\textbf{证明。} 设 \(D_t=d_{\rm eff}(\zeta^I_t,\zeta^J_t)\)。在同一个
\(S^I_t(\zeta^J_t)\) 处使用三角不等式，得到

\[
\begin{aligned}
D_{t+1}
&\le d_{\rm eff}(S^I_t(\zeta^I_t),S^I_t(\zeta^J_t))\\
&\quad+d_{\rm eff}(S^I_t(\zeta^J_t),S^J_t(\zeta^J_t))\\
&\le D_t+a_t.
\end{aligned}
\]

第一项由真实 CPTP 效应拉回给出收缩。第二项的两个仪器实现在同一个
\(\zeta^J_t\) 上比较，P-ARB-02 给出以 \(q^J_t(h)\)
为共同权重的动作界；结果追加置换保持距离。归纳和有限求和给
\(D_N\le D_0+\sum_ta_t\)，零历史的初态置换给
\(D_0=d_{\rm eff}(\rho,\sigma)\)。若每个动作缺陷不超过
\(\varepsilon_t\)，P-ARB-01 的权重归一性给
\(a_t\le\varepsilon_t\)，再求和即可。

对任意有限记录事件 \(S\subseteq\mathsf H_N\)，构造真实效应

\[
E_S=\sum_{h\in S}|h\rangle\langle h|\otimes I_\iota.
\]

它与补效应均 PSD，且在任意联合态上的概率恰为主块质量在 \(S\)
上之和。因此
\(|p^I_N(S)-p^J_N(S)|\le D_N\)。有限归一化分布的全部事件刻画给
\(\operatorname{TV}(p^I_N,p^J_N)\le D_N\)；这一步保持系数一，没有按单点误差再乘历史总数。最后使用
P-ARB-03 已证的 TV 重标记等式，传回原协议。\(\square\)

\(J\) 在此只是选取三角分解共同输入的一个明确方向。公式没有要求
\(\zeta^I_t=\zeta^J_t\)，也没有要求两者的历史权重相同。实际 \(a_t\)
可以反映 \(J\) 前缀中很少访问的历史；统一 \(\varepsilon_t\)
则只需核对所有被控制器选择的底层动作，使用更保守但较易外部估计的上界。后者仍需真实仪器误差证明、校准或其他独立误差输入；本章不把一个定义为上确界的数自动称为已数值算出。

\hypertarget{ux7cbeux786eux6709ux7406ux7684ux4e24ux6b65ux9884ux7b97ux6b63ux63a7}{%
\section{62.6
精确有理的两步预算正控}\label{ux7cbeux786eux6709ux7406ux7684ux4e24ux6b65ux9884ux7b97ux6b63ux63a7}}

以下例子固定 \(O=\iota=\operatorname{Fin}(2)\)，每个仪器使用四个 Kraus
下标。令 \(\mathsf D_b\) 确定输出 \(b\in\{0,1\}\)，唯一非零 Kraus
是该标签的第零个单位矩阵，因此每个矩阵输入均满足

\[
\Phi_{\mathsf D_b,o}(X)=\mathbf1_{o=b}X.
\]

再定义有偏仪器 \(\mathsf B\)：标签零下的 Kraus 四元组为
\((I/2,I/2,I/2,0)\)，标签一下为 \((0,0,0,I/2)\)。四个非零项的 Gram
之和是 \(I\)，所以它是实际完整仪器，且

\[
\Phi_{\mathsf B,0}(X)=\tfrac34X,
\qquad \Phi_{\mathsf B,1}(X)=\tfrac14X.
\]

所有矩阵系数都是精确有理数，没有浮点近似或通过估算平方根设置 Kraus。

控制器在第一步选择动作 \(\mathrm{none}\)，读到首个结果 \(o\)
后，在第二步选择动作 \(\mathrm{some}(o)\)。参考实现令
\(I_{\mathrm{none}}=\mathsf D_0\)；近似实现令
\(J_{\mathrm{none}}=\mathsf B\)；两种实现均令
\(I_{\mathrm{some}(o)}=J_{\mathrm{some}(o)}=\mathsf D_o\)。因此控制器的确读取前一步记录并影响后一步仪器。

\protect\hypertarget{p-arb-05}{}{}

\textbf{P-ARB-05（精确有理逐步预算的饱和见证）。}
对任意二维输入密度态，第一动作的真实记录缺陷为
\(1/4\)，第二动作的对应缺陷为零。实际加权成本分别为
\(a_0=1/4\)、\(a_1=0\)，完整两步记录 TV 精确等于 \(1/4\)。不同初态时则有

\[
d_{\rm eff}(\zeta^I_{2,\rho},\zeta^J_{2,\sigma})
\le d_{\rm eff}(\rho,\sigma)+\tfrac14.
\]

相同输入态 \(\rho\) 时，两种实现的终端量子边际都恰为
\(\rho\)。原协议叶表示的完整记录 TV 也精确等于 \(1/4\)。

\textbf{证明。} 对任意合法联合效应 \(E\)，令
\(x=p_\rho(E_0)\)、\(y=p_\rho(E_1)\)。P-ARB-01 给出
\(0\le x,y\le1\)。两种第一步输出的概率分别为

\[
x,\qquad \tfrac34x+\tfrac14y.
\]

它们之差的绝对值是
\(|x-y|/4\le1/4\)，所以这是所有输入态、所有合法联合效应上的上界。记录事件``输出零''在两种输出中具有概率
\(1\) 和 \(3/4\)，给出匹配下界；对输入态取上确界得到实际动作缺陷
\(1/4\)，没有把单一被测试输入当作上确界证明。

第二步两方在共同历史 \(h\) 上选取的仪器同为
\(\mathsf D_{h(0)}\)，故该动作缺陷为零。前缀权重归一后得到
\(a_0=1/4,a_1=0\)，P-ARB-04 给出全程上界。完整历史 \((0,0)\)
的实际概率在两实现中分别为 \(1\) 与 \(3/4\)，其事件差给 TV 下界
\(1/4\)，与上界相等。

终端量子边际直接由实际 \texttt{run\_\allowbreak matrix}
对四个长度二历史求和：参考实现总和为 \(\rho\)，近似实现总和为
\(3\rho/4+\rho/4=\rho\)。最后使用叶 TV 等式得到原协议结论。\(\square\)

该例检验了``当下结果加入历史、控制器读取既有历史、下一步动作依赖它''的实际顺序，同时检验权重、总和与
TV 系数。精确计算属于有限形式见证，不是实验确证。{[}QM318{]}

\hypertarget{ux540cux4e00ux4e2aux53cdux9988ux4e0bux7684ux5220ux9664ux6807ux7b7eux53cdux4f8b}{%
\section{62.7
同一个反馈下的删除标签反例}\label{ux540cux4e00ux4e2aux53cdux9988ux4e0bux7684ux5220ux9664ux6807ux7b7eux53cdux4f8b}}

保留上一节的控制器和参考实现 \(I\)，另取
\(\widehat J_{\mathrm{none}}=\mathsf D_1\)，并令
\(\widehat J_{\mathrm{some}(o)}=\mathsf D_o\)。对于每一个动作，这两种实现丢弃标签之后的实际量子通道都相同：都是恒等通道。

\protect\hypertarget{n-arb-01}{}{}

\textbf{N-ARB-01（无标签动作零缺陷仍可给出最大完整记录误差）。}
上述两个实际仪器族的每个无标签动作缺陷均为零，但相同任意初态下的两步完整记录
TV 为 \(1\)。第二次结果为零的事件概率分别为 \(1\) 和
\(0\)。两种实现的终端量子密度态均等于输入 \(\rho\)，因而彼此相同。

\textbf{证明。} \(\mathsf D_0\) 与 \(\mathsf D_1\)
的无标签通道都把每个输入矩阵 \(X\) 送到
\(X\)，后续动作两方更是相同仪器。因此无标签通道作为函数相等，已有
\texttt{channel\allowbreak Defect\_\allowbreak self} 给出零缺陷。

参考实现必定得到历史 \((0,0)\)；另一实现必定得到
\((1,1)\)。这些概率由原时间顺序词及真实分支公式求得。历史事件
\(\{(0,0)\}\) 的概率差为 \(1\)，给出完整 TV 下界；有限概率分布 TV 至多为
\(1\)，故正好为 \(1\)。对实际事件 \(\{h:h(1)=0\}\)
展开四个历史，得到第二次结果的概率 \(1\) 与
\(0\)。两边终端量子边际分别对实际记录块求和，均为 \(\rho\)。原协议叶 TV
由实际叶双射保持。\(\square\)

本反例把第58章的标签责任具体接到了新逐步实现与旧协议之间，未把它重新解释成终端量子态不同。若误用零的无标签动作缺陷替换
P-ARB-04 中的
\(\delta_{\rm rec}\)，同初态预算会给出错误的零记录误差；真正的带标签动作缺陷不允许这一步替换。

\hypertarget{ux516cux5171ux63a5ux53e3ux5b9eux9645ux8bc1ux636eux4e0eux7269ux7406ux8d23ux4efb}{%
\section{62.8
公共接口、实际证据与物理责任}\label{ux516cux5171ux63a5ux53e3ux5b9eux9645ux8bc1ux636eux4e0eux7269ux7406ux8d23ux4efb}}

本章新增公共模块为 \texttt{UEOTQM.\allowbreak History\allowbreak Channels} 与
\texttt{UEOTQM.\allowbreak Action\allowbreak Record\allowbreak Budget}。下表给出稳定责任入口；具体正负控位于
\texttt{UEOTQM.\allowbreak Action\allowbreak Record\allowbreak Budget\allowbreak Examples}。{[}QM318{]}

\begin{longtable}[]{@{}
  >{\raggedright\arraybackslash}p{(\columnwidth - 2\tabcolsep) * \real{0.3000}}
  >{\raggedright\arraybackslash}p{(\columnwidth - 2\tabcolsep) * \real{0.7000}}@{}}
\toprule\noalign{}
\begin{minipage}[b]{\linewidth}\raggedright
编号结果或接口
\end{minipage} & \begin{minipage}[b]{\linewidth}\raggedright
实际源码责任
\end{minipage} \\
\midrule\noalign{}
\endhead
\bottomrule\noalign{}
\endlastfoot
P-ARB-01：PSD、质量和零块 &
\texttt{principal\_\allowbreak pos\allowbreak Semidef}、\allowbreak \texttt{sum\_\allowbreak mass\_\allowbreak principal}、\allowbreak \texttt{weight\_\allowbreak nonneg}、\allowbreak \texttt{sum\_\allowbreak weight}、\allowbreak \texttt{zero\_\allowbreak weight\_\allowbreak iff\_\allowbreak zero\_\allowbreak principal} \\
P-ARB-01：合法块效应与真实配对 &
\texttt{restrict\allowbreak Effect}、\allowbreak \texttt{restrict\allowbreak Effect\_\allowbreak matrix}、\allowbreak \texttt{pairing\_\allowbreak block} \\
P-ARB-02：读取 Kraus 与实际块作用 &
\texttt{read\allowbreak Kraus}、\allowbreak \texttt{read\allowbreak Kraus\_\allowbreak gram}、\allowbreak \texttt{read\allowbreak Kraus\_\allowbreak conjugation}、\allowbreak \texttt{controlled}、\allowbreak \texttt{controlled\_\allowbreak operator\_\allowbreak eq\_\allowbreak block} \\
P-ARB-02：从 PSD 块到动作界 &
\texttt{pairing\_\allowbreak operator\_\allowbreak difference\_\allowbreak le\_\allowbreak mass\_\allowbreak mul\_\allowbreak defect}、\allowbreak \texttt{controlled\_\allowbreak state\_\allowbreak distance\_\allowbreak le\_\allowbreak weighted\_\allowbreak defect}、\allowbreak \texttt{controlled\_\allowbreak defect\_\allowbreak le\_\allowbreak uniform} \\
P-ARB-03：有限控制器和时间顺序 &
\texttt{Controller.\allowbreak leaf\allowbreak Equiv}、\allowbreak \texttt{Controller.\allowbreak protocol\_\allowbreak word\_\allowbreak eq}、\allowbreak \texttt{Controller.\allowbreak word\_\allowbreak snoc}、\allowbreak \texttt{record\allowbreak Operator\_\allowbreak append} \\
P-ARB-03：真实步骤与原完整状态 &
\texttt{step}、\allowbreak \texttt{step\_\allowbreak operator\_\allowbreak eq\_\allowbreak block}、\allowbreak \texttt{run\_\allowbreak matrix}、\allowbreak \texttt{run\_\allowbreak eq\_\allowbreak protocol\_\allowbreak state}、\allowbreak \texttt{terminal\_\allowbreak memory\_\allowbreak matrix}、\allowbreak \texttt{weight\_\allowbreak run\_\allowbreak leaf} \\
P-ARB-04：实际逐动作和 &
\texttt{action\allowbreak Cost}、\allowbreak \texttt{run\_\allowbreak succ\_\allowbreak distance\_\allowbreak le}、\allowbreak \texttt{run\_\allowbreak distance\_\allowbreak le\_\allowbreak action\allowbreak Cost\_\allowbreak sum}、\allowbreak \texttt{action\allowbreak Cost\_\allowbreak le}、\allowbreak \texttt{run\_\allowbreak distance\_\allowbreak le\_\allowbreak budget\_\allowbreak sum} \\
P-ARB-04：完整记录与原叶律 &
\texttt{record\_\allowbreak tv\_\allowbreak le\_\allowbreak action\allowbreak Cost\_\allowbreak sum}、\allowbreak \texttt{record\_\allowbreak tv\_\allowbreak le\_\allowbreak budget\_\allowbreak sum}、\allowbreak \texttt{protocol\_\allowbreak record\_\allowbreak tv\_\allowbreak eq}、\allowbreak \texttt{protocol\_\allowbreak record\_\allowbreak tv\_\allowbreak le\_\allowbreak budget\_\allowbreak sum} \\
P-ARB-05：精确有理正控 &
\texttt{biased\_\allowbreak record\_\allowbreak defect\_\allowbreak exact}、\allowbreak \texttt{first\_\allowbreak action\allowbreak Cost\_\allowbreak exact}、\allowbreak \texttt{second\_\allowbreak action\allowbreak Cost\_\allowbreak exact}、\allowbreak \texttt{two\_\allowbreak step\_\allowbreak record\_\allowbreak tv\_\allowbreak exact}、\allowbreak \texttt{original\_\allowbreak protocol\_\allowbreak record\_\allowbreak tv\_\allowbreak exact} \\
N-ARB-01：删除标签与量子边际 &
\texttt{negative\_\allowbreak unlabelled\_\allowbreak defect\_\allowbreak zero}、\allowbreak \texttt{negative\_\allowbreak two\_\allowbreak step\_\allowbreak record\_\allowbreak tv}、\allowbreak \texttt{exact\_\allowbreak second\_\allowbreak record\_\allowbreak zero}、\allowbreak \texttt{flipped\_\allowbreak second\_\allowbreak record\_\allowbreak zero}、\allowbreak \texttt{negative\_\allowbreak terminal\_\allowbreak memories\_\allowbreak equal} \\
\end{longtable}

两份公共模块含54条显式
theorem、20项定义和1项缩写，共75个实际逐名公理审计目标。具体见证另含36条
theorem、14项定义和2项缩写，共52个目标。三份源码均有实际聚焦编译成功回执；两个完整审计序列与实际源码顺序对应，基础只包含
\texttt{propext}、\allowbreak \texttt{Classical.\allowbreak choice}、\allowbreak \texttt{Quot.\allowbreak sound}
或其子集。开发中遇到的失败与最终成功日志均保留，最终日志有普通 linter
警告。独立 B→C 审查分别覆盖公共证明与完整具体见证。实际证据位于
\texttt{physics/qm/kernel/evidence/v31.\allowbreak 8/C}，本轮统一验收和远端发布另由
root 的完整回执承担；这里的局部通过不替代那两项身份。

已闭合的数学责任是：给定明确的有限复量子仪器和共同因果动作规则，预测误差可以通过真实历史通道分解到实际动作记录缺陷，并传到完整记录分布。这让一个难以整体评估的协议对象有了可逐项审查的误差入口。量子记忆维数、每个动作的完整性、记录范围、实际控制器以及输入态误差都保留在前提中。

变长停止树的逐动作补齐、相干控制器、未纳入当前载体的任意环境记忆、diamond
距离和连续时间极限各需新的定义与证明。这里的全效应距离不自动具有这些放大后量词。把动作缺陷变成现实设备校准值、处理建模外泄漏和验证实验适用条件，也有独立经验责任。本章没有新增实验，没有从有限误差内核推出完整自然理论，更没有改变三代结论所依赖的谱、表示、阈值和有效窗口条件。冻结
Core 的106项与四生成器保持原身份。

\hypertarget{appendix-a}{%
\chapter{附录 A. 符号、单位与标准定理责任表}\label{appendix-a}}

\hypertarget{a.1-ux5e38ux7528ux7b26ux53f7}{%
\section{A.1 常用符号}\label{a.1-ux5e38ux7528ux7b26ux53f7}}

\begin{longtable}[]{@{}
  >{\raggedright\arraybackslash}p{(\columnwidth - 2\tabcolsep) * \real{0.3000}}
  >{\raggedright\arraybackslash}p{(\columnwidth - 2\tabcolsep) * \real{0.7000}}@{}}
\toprule\noalign{}
\begin{minipage}[b]{\linewidth}\raggedright
符号
\end{minipage} & \begin{minipage}[b]{\linewidth}\raggedright
含义与约定
\end{minipage} \\
\midrule\noalign{}
\endhead
\bottomrule\noalign{}
\endlastfoot
\(\rho,\sigma\) & 密度算符；允许次归一时显式标记 \\
\(\mathcal E,\mathcal M_a\) & CPTP 通道、instrument 分支 CP 映射 \\
\(E_a,P_a\) & POVM 效应、正交分支投影；通常不相同 \\
\(J_{\mathcal E}\) & 输入在前、无归一 Choi，迹为输入维数 \\
\(D(\rho,\sigma)\) & 使用时依正文定义为迹距离；相对熵写
\(D(\rho\|\sigma)\) \\
\(\delta_\diamond\) & 半 diamond 范数差，除非明确另定 \\
\(C(h)\) & 固定未来目标和协议下的规范响应 \\
FFIPS & 未来/干预预测等价及其规范充分表示 \\
ORC & Objective Record Certificate；操作公共记录证书 \\
SRN & strong record nondisturbance，分支内所有态保持 \\
QDRR & 项目记录写入---逆写/恢复类协议；本书按第8章严格定义使用 \\
\(K,\Lambda\) & SRN 相关矩阵、逐项指数生成元；不是一般 Kraus 算符 \\
\(\mathsf G\) 或 \(G\) & 按上下文为 Gram 矩阵、Gaussian
二次核或群；不可跨章混算 \\
\(\beta_T,\beta_C,\beta_8\) &
逆温、控制折扣、331电荷嵌入参数，三者完全不同 \\
\(\hbar\) & 作用单位标定；场论章自然单位取1；随机动力学等章显式恢复 \\
\(\sigma,\Omega\) & 相空间辛双线性型、矩阵形式 \\
\(\Phi,\Pi\) & Core 结构/资源符号时须给具体功能；不等于固定物理熵 \\
\(\Pi_{\rm cut}\) & 能量截断投影，避免与资源代价混同 \\
\(\Gamma\) & 场论有效作用或顶点；不是 Core 的未定义通用复杂度 \\
\(\lambda,\lambda_H\) & 标量四次耦合、Higgs 四次耦合 \\
\(\mu,\Lambda_*\) & 按上下文为重整化尺度、EFT 高能尺度；化学势另标 \\
\(N_c,N_f,k,t\) & 色数、代数、横向 \(R\) 夸克代数、异常解倍数 \\
\(\chi\) & framing completion 响应类，不是普通四维手征指标 \\
HGP, FAD & 第31、32章各自独立的物理前提包 \\
GOIT & 规范轨道身份运输的 UEOT 解释接口 \\
D/P/K/H/S/E/N/O & 第0章的状态；P 不等于原创或机器证明 \\
\end{longtable}

\hypertarget{a.2-ux5355ux4f4dux89c4ux5219}{%
\section{A.2 单位规则}\label{a.2-ux5355ux4f4dux89c4ux5219}}

第1--20章及实际化、全局相位章节保留 \(\hbar\)，Hamiltonian
具有能量单位，Lindblad 生成元具有时间倒数单位；写
\(\mathcal L\rho=-i[H,\rho]/\hbar+\sum L\rho L^\dagger-\frac12\{L^\dagger L,\rho\}\)
时 \(L\) 具有时间的负二分之一次方单位。Fisher
位置信息具有长度的负二次方单位，作用量具有作用单位。

自然单位下四维 \([\phi]=[A_\mu]=1,[\psi]=3/2,[\partial_\mu]=1\)，
\([\mathcal L]=4\)，规范和四次耦合维数0。截面维数为负2。 EFT
中算符维数、Wilson 系数与显式高能尺度的乘积必须总计为4。

熵用自然对数时是 nat；如改 bits，应除以 \(\log2\)。含温度物理熵时是否乘
\(k_B\) 必须注明；本文 \(\beta_T=(k_BT)^{-1}\)，信息熵默认无量纲。

\hypertarget{a.3-ux8c03ux7528ux800cux672aux5b8cux6574ux91cdux8bc1ux7684ux4e3bux8981ux7ed3ux679c}{%
\section{A.3
调用而未完整重证的主要结果}\label{a.3-ux8c03ux7528ux800cux672aux5b8cux6574ux91cdux8bc1ux7684ux4e3bux8981ux7ed3ux679c}}

\begin{longtable}[]{@{}
  >{\raggedright\arraybackslash}p{(\columnwidth - 4\tabcolsep) * \real{0.2500}}
  >{\raggedright\arraybackslash}p{(\columnwidth - 4\tabcolsep) * \real{0.3400}}
  >{\raggedright\arraybackslash}p{(\columnwidth - 4\tabcolsep) * \real{0.4100}}@{}}
\toprule\noalign{}
\begin{minipage}[b]{\linewidth}\raggedright
标准工具
\end{minipage} & \begin{minipage}[b]{\linewidth}\raggedright
必要条件/本书用途
\end{minipage} & \begin{minipage}[b]{\linewidth}\raggedright
不在本文证明范围
\end{minipage} \\
\midrule\noalign{}
\endhead
\bottomrule\noalign{}
\endlastfoot
有限 OPT 两原则重建 & 使用原预印本的完整标准背景，LE/CIT 与 ESP &
不将摘要当成全部公理清单或已公认最终定理 \\
复极化/二次型刻画 & 全域齐次与平行四边形关系，连续性 &
有限采样无法替代全域量词 \\
相对熵数据处理 & CPTP、适当支撑；有限维主用 & 任意非 CP
近似或无定义熵差 \\
Knill--Laflamme & 固定码与噪声线性跨度、完美恢复 &
未指定噪声的通用硬件容错 \\
谱定理、Stone & 自伴性、强连续性、合适算符域 &
任意形式微分算符自动自伴 \\
Kato--Rellich/Friedrichs & 对称性、相对界、半有界/闭型等条件 &
没有域条件的 Hamiltonian 合法性 \\
Stone--von Neumann & 有限自由度、正则、不可约、非零中心特征 & QFT
无限自由度的全局唯一表示 \\
Williamson 与 Gaussian CP & 有限模、正定协方差、给定辛约定 & 任意非
Gaussian 操作由二阶矩决定 \\
Wigner & 射线概率保持双射与连续性要求 & 任意有噪通道是可逆对称 \\
Lieb--Robinson & 合适局域有界相互作用及空间求和界 &
所有长程、无界场无条件同一光锥 \\
GKSL & 有限维范数连续 CPTP 半群等版本条件 & 任意约化动力学为半群 \\
Davies 弱耦合极限 & 浴平衡、相关衰减、尺度极限等 &
强耦合或长记忆系统自动热 Lindblad \\
过程张量/comb & 完全正性、因果递归、指定插槽 &
只测无干预相关就获得完整过程 \\
自旋---统计、LSZ & 相对论局域/谱/正性或稳定粒子渐近态条件 &
任意非相对论模型或简单带电红外极限 \\
微扰重整化、Ward/ST & 指定调节、方案、无相关规范异常 &
非微扰存在性和全局规范固定 \\
Atiyah--Singer/Dai--Freed & 指定流形、丛、算符、边界及异常类型 & Betti
数代替指数或任意 framing 映射 \\
\end{longtable}

重型定理的条件调用是理论依赖，不是删去证明责任。本书明列可计算实现，读者无需先阅读这些原文才能使用基础公式；如要把其结论推广出声明域，必须回到相应严格定理。

\hypertarget{a.4-ux672cux7248ux65b0ux589eux7b26ux53f7}{%
\section{A.4
本版新增符号}\label{a.4-ux672cux7248ux65b0ux589eux7b26ux53f7}}

\(\theta\) 为环边界扭转，\(\eta\)
为规范不变扭转；\(\mathcal A,\mathcal F\) 在第44章为 Berry
连接和曲率；\(t_1,t_2\) 为 SSH 跃迁；\(a\) 在格点章为格距；\(U_l\)
为群值边变量；\(p_0\) 为校准后的条件事件下界；\(L_n\) 为统计超鞅，不是
Lindblad 算符或线丛。所有重用符号以本章类型为准。

\hypertarget{a.5-ux5df2ux9a8cux8bc1ux5185ux6838ux7684ux65b0ux589eux7b26ux53f7}{%
\section{A.5
已验证内核的新增符号}\label{a.5-ux5df2ux9a8cux8bc1ux5185ux6838ux7684ux65b0ux589eux7b26ux53f7}}

本书第49--50章以 \(R_\rho(w)\) 表示全部有限词响应，以 \(W\)
表示实际词效应的实张成，以 \(x(\rho)\) 表示有限响应坐标。\(d_E\)
是对全部效应 \(0\preceq E\preceq I\) 的概率差上确界；\(\Delta_E\)
是实际通道输出差上确界；它们不是有限样本最大误差。\(p,q\)
是两个输入态对同一结果的 Born 概率，\(\eta\)
是声明的正概率下界。书中阐述的 trace/diamond 范数数学与当前 Lean
对象的已证明识别关系按第50.3、50.5节区分。

\hypertarget{appendix-b}{%
\chapter{附录 B. 编号证明结果注册表}\label{appendix-b}}

本版共147个稳定编号结果，完整保留3.3版全部131个编号，第60--62章新增16项命题或反例。既有131项包括2.0的78项、第49--50章10项、第51--53章13项、第54--56章16项和第57--59章14项。编号登记正文解析结果，不表示原创性、Lean声明或实验确认的计数；完整前提、证明和开放责任见相应章节。

\begin{longtable}[]{@{}
  >{\raggedright\arraybackslash}p{(\columnwidth - 4\tabcolsep) * \real{0.2400}}
  >{\raggedleft\arraybackslash}p{(\columnwidth - 4\tabcolsep) * \real{0.0700}}
  >{\raggedright\arraybackslash}p{(\columnwidth - 4\tabcolsep) * \real{0.6900}}@{}}
\toprule\noalign{}
\begin{minipage}[b]{\linewidth}\raggedright
编号
\end{minipage} & \begin{minipage}[b]{\linewidth}\raggedleft
章
\end{minipage} & \begin{minipage}[b]{\linewidth}\raggedright
结果
\end{minipage} \\
\midrule\noalign{}
\endhead
\bottomrule\noalign{}
\endlastfoot
\protect\hyperlink{p-qbr-01}{P-QBR-01} & 1 &
P-QBR-01：经典实现反模型。 \\
\protect\hyperlink{p-qsel-01}{P-QSEL-01} & 2 &
P-QSEL-01：经典与实量子删除见证。 \\
\protect\hyperlink{p-qact-01}{P-QACT-01} & 3 &
P-QACT-01：Madelung--Schrödinger 条件等价。 \\
\protect\hyperlink{p-qact-02}{P-QACT-02} & 3 & P-QACT-02。 \\
\protect\hyperlink{p-qborn-02}{P-QBORN-02} & 3 &
P-QBORN-02：无偏最终分支的概率。 \\
\protect\hyperlink{p-qop-01}{P-QOP-01} & 4 & P-QOP-01：有限 Choi--Kraus
等价。 \\
\protect\hyperlink{p-qop-02}{P-QOP-02} & 4 & P-QOP-02：无信号。 \\
\protect\hyperlink{p-qop-03}{P-QOP-03} & 4 &
P-QOP-03：操作误差与传播。 \\
\protect\hyperlink{p-qff-01}{P-QFF-01} & 5 & P-QFF-01：有限量子
FFIPS。 \\
\protect\hyperlink{p-qff-02}{P-QFF-02} & 5 &
P-QFF-02：量子商的路径保持。 \\
\protect\hyperlink{p-qmeas-01}{P-QMEAS-01} & 6 &
P-QMEAS-01：父系统完成。 \\
\protect\hyperlink{p-qmeas-02}{P-QMEAS-02} & 6 & P-QMEAS-02。 \\
\protect\hyperlink{p-qmeas-03}{P-QMEAS-03} & 6 &
P-QMEAS-03：公共盲性。 \\
\protect\hyperlink{p-qmeas-04}{P-QMEAS-04} & 6 & P-QMEAS-04。 \\
\protect\hyperlink{p-qsrn-01}{P-QSRN-01} & 7 & P-QSRN-01：SRN
完备分类。 \\
\protect\hyperlink{p-qsrn-02}{P-QSRN-02} & 7 & P-QSRN-02。 \\
\protect\hyperlink{p-qsrn-03}{P-QSRN-03} & 7 &
P-QSRN-03：生成元锥的必要充分条件。 \\
\protect\hyperlink{p-qqdr-01}{P-QQDR-01} & 8 & P-QQDR-01：相干拉回。 \\
\protect\hyperlink{p-qqdr-02}{P-QQDR-02} & 8 & P-QQDR-02。 \\
\protect\hyperlink{p-qqdr-03}{P-QQDR-03} & 8 & P-QQDR-03：同时误差。 \\
\protect\hyperlink{p-qinfo-01}{P-QINFO-01} & 9 & P-QINFO-01。 \\
\protect\hyperlink{p-qinfo-02}{P-QINFO-02} & 9 & P-QINFO-02。 \\
\protect\hyperlink{p-qinf-01}{P-QINF-01} & 10 &
P-QINF-01：正常态的迹范数连续性。 \\
\protect\hyperlink{p-qinf-02}{P-QINF-02} & 10 &
P-QINF-02：能量截断误差。 \\
\protect\hyperlink{p-qccr-01}{P-QCCR-01} & 11 &
P-QCCR-01：辛协变的单尺度。 \\
\protect\hyperlink{p-qccr-02}{P-QCCR-02} & 11 & P-QCCR-02。 \\
\protect\hyperlink{p-qccr-03}{P-QCCR-03} & 11 & P-QCCR-03。 \\
\protect\hyperlink{p-qgauss-01}{P-QGAUSS-01} & 13 &
P-QGAUSS-01：协方差必要条件。 \\
\protect\hyperlink{p-qmem-01}{P-QMEM-01} & 18 &
P-QMEM-01：遗漏记忆的不可消除误差。 \\
\protect\hyperlink{p-qtherm-01}{P-QTHERM-01} & 19 &
P-QTHERM-01：自由能与相对熵。 \\
\protect\hyperlink{p-qtherm-02}{P-QTHERM-02} & 19 & P-QTHERM-02：有限
Gibbs 的 KMS。 \\
\protect\hyperlink{p-qtherm-03}{P-QTHERM-03} & 19 & P-QTHERM-03。 \\
\protect\hyperlink{p-qctl-01}{P-QCTL-01} & 20 & P-QCTL-01。 \\
\protect\hyperlink{p-qkg-01}{P-QKG-01} & 22 &
P-QKG-01：正则关系与传播子。 \\
\protect\hyperlink{p-qkg-02}{P-QKG-02} & 22 & P-QKG-02：类空因果性。 \\
\protect\hyperlink{p-qdir-01}{P-QDIR-01} & 23 & P-QDIR-01。 \\
\protect\hyperlink{p-qunit-01}{P-QUNIT-01} & 24 &
P-QUNIT-01：光学定理。 \\
\protect\hyperlink{p-qren-01}{P-QREN-01} & 25 & P-QREN-01。 \\
\protect\hyperlink{p-qren-02}{P-QREN-02} & 25 & P-QREN-02：泡图减除。 \\
\protect\hyperlink{p-qren-03}{P-QREN-03} & 25 & P-QREN-03：二粒子割。 \\
\protect\hyperlink{p-qeft-01}{P-QEFT-01} & 26 & P-QEFT-01：二次 Schur
补。 \\
\protect\hyperlink{p-qeft-02}{P-QEFT-02} & 26 &
P-QEFT-02：受限的实时桥。 \\
\protect\hyperlink{p-qqed-01}{P-QQED-01} & 28 & P-QQED-01。 \\
\protect\hyperlink{p-qqed-02}{P-QQED-02} & 28 & P-QQED-02：阈值
beta。 \\
\protect\hyperlink{p-qsm-01}{P-QSM-01} & 29 &
P-QSM-01：异常解与非唯一性。 \\
\protect\hyperlink{p-qhig-01}{P-QHIG-01} & 30 & P-QHIG-01：质量矩阵。 \\
\protect\hyperlink{p-qflav-01}{P-QFLAV-01} & 30 & P-QFLAV-01：一般非退化
\(N\) 代夸克计数。 \\
\protect\hyperlink{p-qnu-01}{P-QNU-01} & 30 & P-QNU-01：Type-I
seesaw。 \\
\protect\hyperlink{p-qg3-01}{P-QG3-01} & 31 & P-QG3-01。 \\
\protect\hyperlink{p-qg3-02}{P-QG3-02} & 31 & P-QG3-02：全部整数解。 \\
\protect\hyperlink{p-qg3-03}{P-QG3-03} & 31 &
P-QG3-03：去原始性前提的选择定理。 \\
\protect\hyperlink{p-qg3-04}{P-QG3-04} & 31 & P-QG3-04。 \\
\protect\hyperlink{p-qfr-01}{P-QFR-01} & 32 & P-QFR-01。 \\
\protect\hyperlink{p-qfr-02}{P-QFR-02} & 32 & P-QFR-02。 \\
\protect\hyperlink{g-01}{G-01} & 36 & G-01：三标量配置的全部整数解。 \\
\protect\hyperlink{g-02}{G-02} & 36 & G-02：一般色数局域完成。 \\
\protect\hyperlink{f-01}{F-01} & 37 & F-01。 \\
\protect\hyperlink{f-02}{F-02} & 37 & F-02。 \\
\protect\hyperlink{f-03}{F-03} & 37 & F-03。 \\
\protect\hyperlink{m-01}{M-01} & 38 & M-01。 \\
\protect\hyperlink{m-02}{M-02} & 38 & M-02。 \\
\protect\hyperlink{m-03}{M-03} & 39 & M-03：单时刻的最小分支数为4。 \\
\protect\hyperlink{m-04}{M-04} & 39 & M-04：锐经典下界。 \\
\protect\hyperlink{m-05}{M-05} & 39 & M-05。 \\
\protect\hyperlink{e-01}{E-01} & 39 & E-01：有限样本反证。 \\
\protect\hyperlink{e-02}{E-02} & 39 & E-02：近似人口保持的稳健界。 \\
\protect\hyperlink{c-01}{C-01} & 40 & C-01：明确的远端制备矛盾。 \\
\protect\hyperlink{c-02}{C-02} & 40 & C-02：范数与线性平均。 \\
\protect\hyperlink{c-03}{C-03} & 40 & C-03：单谱分支与 Born 权重。 \\
\protect\hyperlink{c-04}{C-04} & 40 & C-04：成对不确定性衰减。 \\
\protect\hyperlink{c-05}{C-05} & 41 &
C-05：不受限量子环境的等价障碍。 \\
\protect\hyperlink{p-ugl-01}{P-UGL-01} & 43 &
P-UGL-01：环上谱和规范不变量。 \\
\protect\hyperlink{p-utp-01}{P-UTP-01} & 44 & P-UTP-01：此本征线的 Chern
数。 \\
\protect\hyperlink{p-utp-02}{P-UTP-02} & 44 & P-UTP-02：半无限零模。 \\
\protect\hyperlink{p-ulg-01}{P-ULG-01} & 45 &
P-ULG-01：有限模型存在且规范不变。 \\
\protect\hyperlink{p-ulg-02}{P-ULG-02} & 45 &
P-ULG-02：精确配分函数与响应。 \\
\protect\hyperlink{p-ust-01}{P-UST-01} & 46 &
P-UST-01：相关试次的固定样本检验。 \\
\protect\hyperlink{p-ust-02}{P-UST-02} & 46 &
P-UST-02：可选停止证书。 \\
\protect\hyperlink{p-qkr-01}{P-QKR-01} & 49 &
P-QKR-01（响应对象的后验闭合）。 \\
\protect\hyperlink{p-qkr-02}{P-QKR-02} & 49 &
P-QKR-02（局部有限表示决定全部未来）。 \\
\protect\hyperlink{p-qkr-03}{P-QKR-03} & 49 &
P-QKR-03（全词坐标控制完整有限适应树）。 \\
\protect\hyperlink{p-qkr-04}{P-QKR-04} & 49 &
P-QKR-04（精确证书可靠性）。 \\
\protect\hyperlink{p-qst-01}{P-QST-01} & 50 &
P-QST-01（实际正配对与效应界）。 \\
\protect\hyperlink{p-qst-02}{P-QST-02} & 50 &
P-QST-02（续效应的完整拉回）。 \\
\protect\hyperlink{p-qst-03}{P-QST-03} & 50 &
P-QST-03（可观测距离的基本合同）。 \\
\protect\hyperlink{p-qst-04}{P-QST-04} & 50 & P-QST-04（完整 Kraus
通道的一步稳定性）。 \\
\protect\hyperlink{p-qst-05}{P-QST-05} & 50 &
P-QST-05（真实量子轨迹的可组合误差）。 \\
\protect\hyperlink{p-qst-06}{P-QST-06} & 50 &
P-QST-06（后选择误差及正支撑保护）。 \\
\protect\hyperlink{p-qmx-01}{P-QMX-01} & 51 &
P-QMX-01（概率混合的实际实现）。 \\
\protect\hyperlink{p-qmx-02}{P-QMX-02} & 51 &
P-QMX-02（同库权重变化的全称上界）。 \\
\protect\hyperlink{p-qmx-03}{P-QMX-03} & 51 &
P-QMX-03（权重与成员误差的联合上界）。 \\
\protect\hyperlink{p-qmx-04}{P-QMX-04} & 51 &
P-QMX-04（精确有理预算的解释）。 \\
\protect\hyperlink{p-qce-01}{P-QCE-01} & 52 &
P-QCE-01（平均态的精确残差与全效应预算）。 \\
\protect\hyperlink{p-qce-02}{P-QCE-02} & 52 &
P-QCE-02（实际密度集合的紧性）。 \\
\protect\hyperlink{p-qce-03}{P-QCE-03} & 52 &
P-QCE-03（有限完整通道的平稳密度态）。 \\
\protect\hyperlink{n-qce-01}{N-QCE-01} & 52 &
N-QCE-01（平均不变不等于轨道收敛或唯一性）。 \\
\protect\hyperlink{p-qsv-01}{P-QSV-01} & 53 &
P-QSV-01（次随机核的完整量子仪器）。 \\
\protect\hyperlink{p-qsv-02}{P-QSV-02} & 53 &
P-QSV-02（对角动力学与有限成功词的交织）。 \\
\protect\hyperlink{p-qsv-03}{P-QSV-03} & 53 &
P-QSV-03（实际全成功记录的几何生存律）。 \\
\protect\hyperlink{p-qsv-04}{P-QSV-04} & 53 & P-QSV-04（Doob
核的实际完整实现与不变成功记录）。 \\
\protect\hyperlink{n-qsv-01}{N-QSV-01} & 53 & N-QSV-01（QSD、Doob
不变律与任意初态不可混同）。 \\
\protect\hyperlink{p-qcl-01}{P-QCL-01} & 54 &
P-QCL-01（完整单纯形的状态、响应与运输）。 \\
\protect\hyperlink{n-qcl-01}{N-QCL-01} & 54 &
N-QCL-01（二点混态没有有限经典纯化）。 \\
\protect\hyperlink{p-qcl-02}{P-QCL-02} & 54 &
P-QCL-02（实际终端商保留二点支持）。 \\
\protect\hyperlink{p-qcl-03}{P-QCL-03} & 54 &
P-QCL-03（同源安全轨迹与纯化分离证书）。 \\
\protect\hyperlink{n-qcl-02}{N-QCL-02} & 54 &
N-QCL-02（微观混合可被实际操作商消去）。 \\
\protect\hyperlink{p-qrt-01}{P-QRT-01} & 55 &
P-QRT-01（两个合法且不同的实联合态）。 \\
\protect\hyperlink{n-qrt-01}{N-QRT-01} & 55 &
N-QRT-01（实局域乘积测试不能确定联合态）。 \\
\protect\hyperlink{p-qrt-02}{P-QRT-02} & 55 &
P-QRT-02（联合效应的确定性区分）。 \\
\protect\hyperlink{n-qrt-02}{N-QRT-02} & 55 &
N-QRT-02（允许复局域效应后的区分）。 \\
\protect\hyperlink{n-qrt-03}{N-QRT-03} & 55 &
N-QRT-03（区分效应不在有限实产品效应的线性张成中）。 \\
\protect\hyperlink{p-qqd-01}{P-QQD-01} & 56 &
P-QQD-01（过滤归一化的合法性）。 \\
\protect\hyperlink{p-qqd-02}{P-QQD-02} & 56 & P-QQD-02（实际量子 Doob
仪器）。 \\
\protect\hyperlink{p-qqd-03}{P-QQD-03} & 56 &
P-QQD-03（全算符有限交织与真实条件态连接）。 \\
\protect\hyperlink{p-qqd-04}{P-QQD-04} & 56 &
P-QQD-04（准稳态、几何存活与实际通道平稳）。 \\
\protect\hyperlink{p-qqd-05}{P-QQD-05} & 56 &
P-QQD-05（全矩阵可计算的非对角 Doob 见证）。 \\
\protect\hyperlink{n-qqd-01}{N-QQD-01} & 56 &
N-QQD-01（必须保留的三项区分）。 \\
\protect\hyperlink{p-pur-01}{P-PUR-01} & 57 &
P-PUR-01（完整密度态空间中的向量纯态）。 \\
\protect\hyperlink{p-pur-02}{P-PUR-02} & 57 &
P-PUR-02（任意有限复密度态的纯化存在）。 \\
\protect\hyperlink{p-opt-01}{P-OPT-01} & 57 &
P-OPT-01（同一操作语言中的三个具体呈现）。 \\
\protect\hyperlink{n-opt-01}{N-OPT-01} & 57 &
N-OPT-01（基本操作规律不自动给出量子选择）。 \\
\protect\hyperlink{p-cq-01}{P-CQ-01} & 58 &
P-CQ-01（矩形通道的完全正性、迹保持与效应收缩）。 \\
\protect\hyperlink{p-cq-02}{P-CQ-02} & 58 &
P-CQ-02（真实仪器的联合块状态与边际）。 \\
\protect\hyperlink{p-cq-03}{P-CQ-03} & 58 &
P-CQ-03（全部有限协议的记录与记忆实现）。 \\
\protect\hyperlink{p-cq-04}{P-CQ-04} & 58 &
P-CQ-04（记录分布的联合距离与扰动控制）。 \\
\protect\hyperlink{n-cq-01}{N-CQ-01} & 58 &
N-CQ-01（相同无条件通道不控制记录误差）。 \\
\protect\hyperlink{p-qdd-01}{P-QDD-01} & 59 & P-QDD-01（当前 Doob
通道的精确动力学）。 \\
\protect\hyperlink{p-qdd-02}{P-QDD-02} & 59 &
P-QDD-02（实际源成功算符与质量闭式）。 \\
\protect\hyperlink{p-qdd-03}{P-QDD-03} & 59 &
P-QDD-03（实际条件态及其定量收敛）。 \\
\protect\hyperlink{p-qdd-04}{P-QDD-04} & 59 &
P-QDD-04（有限成功概率衰减与初态责任）。 \\
\protect\hyperlink{n-qdd-01}{N-QDD-01} & 59 & N-QDD-01（合法恒等 Doob
负控）。 \\
\protect\hyperlink{p-sep-01}{P-SEP-01} & 60 & P-SEP-01（Hermitian
算符的合法效应表示）。 \\
\protect\hyperlink{p-sep-02}{P-SEP-02} & 60 &
P-SEP-02（全效应分离与距离零的矩阵含义）。 \\
\protect\hyperlink{p-clt-01}{P-CLT-01} & 60 &
P-CLT-01（完整有限复模型的普遍局部层析）。 \\
\protect\hyperlink{n-sep-01}{N-SEP-01} & 60 &
N-SEP-01（全部对角效应仍不能识别相干态）。 \\
\protect\hyperlink{p-puc-01}{P-PUC-01} & 61 &
P-PUC-01（完整有限复密度空间的纯性刻画）。 \\
\protect\hyperlink{p-pun-01}{P-PUN-01} & 61 &
P-PUN-01（同辅助载体的完整纯化酉等价）。 \\
\protect\hyperlink{p-aux-01}{P-AUX-01} & 61 &
P-AUX-01（辅助等距嵌入的实际通道、边际和纯性）。 \\
\protect\hyperlink{p-pun-02}{P-PUN-02} & 61 &
P-PUN-02（不同有限辅助载体的共同扩展酉等价）。 \\
\protect\hyperlink{n-pun-01}{N-PUN-01} & 61 &
N-PUN-01（伴随、转置和逐项共轭不能替代实际辅助作用）。 \\
\protect\hyperlink{n-aux-01}{N-AUX-01} & 61 &
N-AUX-01（严格辅助扩维是等距嵌入，不是双侧酉等价）。 \\
\protect\hyperlink{p-arb-01}{P-ARB-01} & 62 &
P-ARB-01（历史块、归一权重与合法块效应）。 \\
\protect\hyperlink{p-arb-02}{P-ARB-02} & 62 &
P-ARB-02（实际历史读取通道的动作加权界）。 \\
\protect\hyperlink{p-arb-03}{P-ARB-03} & 62 &
P-ARB-03（有限历史演化与原完整协议的矩阵一致性）。 \\
\protect\hyperlink{p-arb-04}{P-ARB-04} & 62 &
P-ARB-04（实际逐动作和预算与完整记录 TV）。 \\
\protect\hyperlink{p-arb-05}{P-ARB-05} & 62 &
P-ARB-05（精确有理逐步预算的饱和见证）。 \\
\protect\hyperlink{n-arb-01}{N-ARB-01} & 62 &
N-ARB-01（无标签动作零缺陷仍可给出最大完整记录误差）。 \\
\end{longtable}

\hypertarget{appendix-c}{%
\chapter{附录 C. 来源、继承关系与研究诚实性}\label{appendix-c}}

\hypertarget{c.1-ux672cux7248ux5b9eux9645ux6765ux6e90ux8fb9ux754c}{%
\section{C.1
本版实际来源边界}\label{c.1-ux672cux7248ux5b9eux9645ux6765ux6e90ux8fb9ux754c}}

本版直接核对五份原始/Core资料：原始UEOT PDF、Core
v3完整数学主文、开放端口矩阵、post-Core科学闭合说明、理论完成路线。另恢复并核对已形成的统一主文2.0
Markdown、LaTeX、PDF及完整源/验证包；Markdown和LaTeX与包内同名成员逐字节相同。独立PDF与包内PDF字节不同，分别保留各自身份，不声称二者是同一渲染产物。

现行数学与运行代码以v31.8源清单的逐文件SHA-256及实际验证为准。本轮精确基线为已完成v31.7源码和后续远端验收元数据的主线7d05d0369b6793e28f0e39efaa77c0c9887479d1。该元数据主线与v31.7实际受验源码和主线分别具有身份，详见历史远端回执；不将其改称本轮CI。新模块在独立工作树开发，旧正式正文、原始来源登记和历史证据按版本保留。两份同名v28.1来源继续保持独立身份；本轮未声称重新阅读所有原始附件或重跑未变旧程序。

\hypertarget{c.2-ux9879ux76eeux6765ux6e90ux6620ux5c04}{%
\section{C.2
项目来源映射}\label{c.2-ux9879ux76eeux6765ux6e90ux6620ux5c04}}

\begin{longtable}[]{@{}
  >{\raggedright\arraybackslash}p{(\columnwidth - 2\tabcolsep) * \real{0.3000}}
  >{\raggedright\arraybackslash}p{(\columnwidth - 2\tabcolsep) * \real{0.7000}}@{}}
\toprule\noalign{}
\begin{minipage}[b]{\linewidth}\raggedright
来源
\end{minipage} & \begin{minipage}[b]{\linewidth}\raggedright
主文使用方式
\end{minipage} \\
\midrule\noalign{}
\endhead
\bottomrule\noalign{}
\endlastfoot
UEOT\_\allowbreak v3\_\allowbreak top.pdf &
原始论文的对象、相位、资源与物理统一动机；物理主张须经后续严格化 \\
UEOT\_\allowbreak v1.0\_\allowbreak 核心理论正式版.md & 原始定义、Core 与物理研究方向 \\
UEOT\_\allowbreak 数学基础\_\allowbreak 统一正式版\_\allowbreak v1.1(1).md &
早期数学支撑；与后续严谨边界比较 \\
UEOT\_\allowbreak Core\_\allowbreak Mathematics\_\allowbreak v3.0\_\allowbreak Complete.md &
通用协议/预测/载体/控制桥接；不把量子假设回填为一般公理 \\
UEOT-QOD\_\allowbreak Interface\_\allowbreak v1.0；Core v2.2 双主线 &
物理实现与通用对象数学的接口分离 \\
UEOT-QM v8.1、v9--20 演进 & 纯化、预测对象、访问载体、过程
RG、盲理论竞争的历史来源；通用部分在数学主文展开 \\
UEOT-QM v21、v22、v23.0 & 相干测量父层、公共盲性、SRN
分类/生成元、QDRR；第4--9章 \\
P2\_\allowbreak Quantum\_\allowbreak Selection\_\allowbreak Minimality\_\allowbreak Frozen.md & 固定 OPT
背景下的两原则选择与相对独立性；第2章 \\
UEOT-QM v24.4 P1-CrossStage-Harmonized & 自伴域、Weyl、UCA、有限
canonical 自由度及共同作用尺度条件；第10--13章 \\
UEOT-QM v25.2 P3-CrossStage-Harmonized &
对称、自旋、相同粒子和跨阶段边界；第14--15章 \\
UEOT-QM v26.3 P4-Final-Closure &
多体、开放、过程记忆与统计物理；第16--20章 \\
UEOT-QM v27.1 Relativistic-Free-Fields &
相对论假设、自由场、局域性；第21--23章 \\
UEOT-QM v28.1 Interacting-QFT-RG-EFT & 相互作用、圈图、重整化和三类 RG
桥；第24--26章 \\
UEOT-QM v29.2 Gauge-QED-Final-Closure & QED、Ward、UV
极点、散射、阈值与磁矩；第27--28章 \\
UEOT-QM v30.0 P8A、v30.2 P8C、v30.3 P8D & 非
Abelian/SM候选、异常、Higgs/味、旧选择器边界；第29--30章 \\
UEOT-QM v30.4--v30.9 演进及 v30.10 归纳 & CP、common-q、framing、自治与
completion 路线的取舍；第32章 \\
UEOT-QM v30.11 Gauge-Orbit-Horizontal-Anomaly-Atomicity &
GOIT、复制障碍、HGP、331和条件代数；第29、31章 \\
\end{longtable}

表中的早期历史材料有些仅作为后续归纳和接口来源，不冒称本轮逐字重读了其全部历史正文。主文没有把旧版本声称的
PASS 数直接搬成本次验证结果。

\hypertarget{c.3-ux5916ux90e8ux539fux59cbux7814ux7a76ux53caux4f5cux8005ux8d44ux6599}{%
\section{C.3
外部原始研究及作者资料}\label{c.3-ux5916ux90e8ux539fux59cbux7814ux7a76ux53caux4f5cux8005ux8d44ux6599}}

以下链接用于定位数学/物理先例，不把标准理论据为 UEOT 原创：

\begin{enumerate}
\def\labelenumi{\arabic{enumi}.}
\tightlist
\item
  Kenji Nakahira，\emph{Two Operational Principles Single Out Quantum
  Theory}，2026预印本。第2章依赖其完整 OPT
  定理背景；预印本状态与已完成实验验证不同。\href{https://arxiv.org/abs/2605.23217}{原文}
\item
  John Watrous，\emph{The Theory of Quantum
  Information}，作者提供的2018版。第4、9章的通道、范数、信息基本工具可与此核对。\href{https://cs.uwaterloo.ca/~watrous/TQI/}{作者页面}
\item
  Pollock 等，\emph{Operational Markov Condition for Quantum
  Processes}。第18章的因果干预判据属于既有过程理论。\href{https://link.aps.org/doi/10.1103/PhysRevLett.120.040405}{原论文}
\item
  Pollock 等，\emph{Operational Quantum Theory of Non-Markovian
  Processes}。用于定位多时过程张量的原始框架。\href{https://arxiv.org/abs/1512.00589}{原文}
\item
  David Tong，\emph{Quantum Field
  Theory}，作者讲义。作为自由场、散射及规范场计算的标准背景；本书推导和数值基准另在正文给出。\href{https://www.damtp.cam.ac.uk/user/tong/qft.htm}{作者页面}
\item
  Kazuya Yonekura，\emph{Dai-Freed theorem and topological phases of
  matter}。第29章 determinant-line
  解释的数学物理来源。\href{https://arxiv.org/abs/1607.01873}{原文}
\item
  I. García-Etxebarria、M. Montero，\emph{Dai-Freed anomalies in
  particle
  physics}。区分局域与全局异常检查。\href{https://arxiv.org/abs/1808.00009}{原文}
\item
  F. Pisano、V. Pleitez，\emph{SU(3)⊗U(1) Model for Electroweak
  Interactions}，1992。331表示组织有既有先例。\href{https://arxiv.org/abs/hep-ph/9206242}{原文}
\item
  R. Foot 等，\emph{SU(3)L⊗U(1)N and SU(4)L⊗U(1)N gauge models with
  right-handed
  neutrinos}，1994。用于331RHN类的先例定位。\href{https://link.aps.org/doi/10.1103/PhysRevD.50.R34}{原论文}
\item
  R. M. Lutchyn、J. D. Sau、S. Das Sarma，\emph{Majorana Fermions and a
  Topological Phase Transition in Semiconductor-Superconductor
  Heterostructures}。第34章器件有效模型的理论来源，非新实验声明。\href{https://arxiv.org/abs/1002.4033}{原文}
\item
  Y. Oreg、G. Refael、F. von Oppen，\emph{Helical liquids and Majorana
  bound states in quantum
  wires}。同一应用方向的独立理论先例。\href{https://arxiv.org/abs/1003.1145}{原文}
\item
  D. Burgarth、G. Chiribella、V. Giovannetti、P. Perinotti、K.
  Yuasa，\emph{Ergodic and Mixing Quantum Channels in Finite
  Dimensions}。第52章的平稳、遍历和混合层次参照；本轮机器证明的实际范围单列。\href{https://arxiv.org/abs/1210.5625}{原文}
\item
  F. Carollo、J. P. Garrahan、I. Lesanovsky、C.
  Pérez-Espigares，\emph{Making rare events typical in Markovian open
  quantum
  systems}。第53章之外的一般量子条件动力学背景；不作为本项目已实现无限条件极限的证明。\href{https://arxiv.org/abs/1711.10951}{原文}
\item
  {[}CDP2010{]} G. Chiribella、G. M. D'Ariano、P.
  Perinotti，\emph{Probabilistic theories with purification}，Physical
  Review A 81, 062348
  (2010)。第57章用于区分纯化存在与含本质唯一性的完整原则。\href{https://arxiv.org/abs/0908.1583}{原文}
\item
  {[}CDP2011{]} 同上作者，\emph{Informational derivation of Quantum
  Theory}，Physical Review A 84, 012311
  (2011)。其完整重建使用明示背景公理，未由本书有限纯化存在结果替代。\href{https://arxiv.org/abs/1011.6451}{原文}
\end{enumerate}

\hypertarget{c.4-ux672cux4e66ux65b0ux589eux6574ux7406ux4e0eux903bux8f91ux4f18ux5316ux7684ux8eabux4efd}{%
\section{C.4
本书新增整理与逻辑优化的身份}\label{c.4-ux672cux4e66ux65b0ux589eux6574ux7406ux4e0eux903bux8f91ux4f18ux5316ux7684ux8eabux4efd}}

本书不是简单粘贴历史文章。它重编了单位、Choi
腿序、过程接口、物理/公共层区分、条件证明、失败出口和应用误差责任；第31章明确证明原始包前提在所声明弱上界下冗余；一般331写法显式补齐带电
\(E^c\)，防止遗漏异常贡献。

这些工作属于本次独立组织、核算与条件推导，不等于声称发现了新的331算术或新粒子物理。要主张学术原创，需要独立先例检索、精确新定理范围、同行核查和可区别的物理后果。文档自足性也不意味着拒绝标准定理依赖。

\hypertarget{c.5-ux7ee7ux627fux7684ux9ad8ux7ea7ux6a21ux5757ux4e0eux672cux7248ux65b0ux589eux6765ux6e90}{%
\section{C.5
继承的高级模块与本版新增来源}\label{c.5-ux7ee7ux627fux7684ux9ad8ux7ea7ux6a21ux5757ux4e0eux672cux7248ux65b0ux589eux6765ux6e90}}

量子基础主文与前沿推导的全部编号结论整合于正文；新增全局环模型、Berry/SSH、有限格点及条件试次检验。第44--45章的标准物理来源分别为
Berry 1984、Su--Schrieffer--Heeger 1979、Wilson
1974，链接紧邻相应推导。新写的解析证明和组合用途不构成``文献首次''的证明。现有标准文献与项目材料承担不同的来源责任。

本版未获取新的完整历史对话或新增真实硬件数据。原始附件、可访问演进稿以及此前两份完整主文提供研究沿革；不可访问内容不作为隐藏证明依据。

\hypertarget{c.6-ux7b2cux4e09ux7248ux7684ux4ee3ux7801ux4e0eux539fux59cbux5b66ux672fux6838ux5bf9}{%
\section{C.6
第三版的代码与原始学术核对}\label{c.6-ux7b2cux4e09ux7248ux7684ux4ee3ux7801ux4e0eux539fux59cbux5b66ux672fux6838ux5bf9}}

第49--59章分别保留Kernel1.1--1.7的有限内核、稳定性、混合平均、生存、选择分离、纯化、真实协议与条件动力学。第60--62章接入Kernel1.8的全效应分离和复局部层析、完整纯性和辅助空间比较、共同控制器逐动作记录预算。当前代码与证据入口为
\texttt{formalization/ueot-qm} 和
\texttt{physics/qm/kernel/evidence/v31.\allowbreak 8}；历史回执保留原始输入与执行范围。冻结Core与领域接口决定分类，不将量子域内性质回填为通用公理。

3.0版修订时曾直接核对 Nakahira 预印本的 Theorem1 及 Appendix
A.2--A.3：完整OPT背景是调用条件，结论不要求每个自然数维数都必须有物理系统。正文第2章保留外部标准定理K的身份，未将其计入本项目通用Lean定理。\href{https://arxiv.org/html/2605.23217v1}{原始HTML}。通道与范数的标准数学可参照\href{https://cs.uwaterloo.ca/~watrous/TQI/}{Watrous作者资料}，本版新Lean证明另从实际矩阵构造，未复制其书稿或把其全部理论标为本项目机器重证。

原始论文第30--33页的
Fisher→不确定关系、闭作用→CCR、Betti数→指标→三代和Massey积→规范群的强推断，需要本书已经分别补出的桥、领域输入或开放义务。保留研究思想不等于恢复已撤回的推断。现代严格版本中的具体前提、反例和可复查证明承担结论责任。

来源标签QM316指向经典扇区、Core纯化边界、实局部层析与一般量子Doob；QM317指向完整态纯化、三操作呈现、矩形通道、联合记录记忆和指定条件动力学；其历史执行分别由对应版本回执承担。QM318指向
\texttt{Quantum\allowbreak Separation}、\allowbreak \texttt{Complex\allowbreak Local\allowbreak Tomography}、\allowbreak \texttt{Pure\allowbreak State\allowbreak Characterization}、\allowbreak \texttt{Purification\allowbreak Uniqueness}、\allowbreak \texttt{History\allowbreak Channels}、\allowbreak \texttt{Action\allowbreak Record\allowbreak Budget}
及三份新见证、交审与统一实际验收。第60--62章使用的Mathlib辅助定理按锁定源码注明，未将库中重定理写成本项目重新原创的物理结果。既有外部文献的引用范围继续以原登记为准。

\hypertarget{appendix-d}{%
\chapter{附录 D. 全部继承验证代码与版本回执}\label{appendix-d}}

\hypertarget{d.1-ux8fd0ux884cux5408ux540c}{%
\section{D.1 运行合同}\label{d.1-ux8fd0ux884cux5408ux540c}}

3.0版在2026-10-19T23:13:31.030592+00:00实际重新执行三套继承程序，共通过757项检查（237、423、97）。3.1、3.2、3.3与3.4版逐字节保留这三套程序并核对哈希，因此沿用这份版本回执，没有将它重命名为本轮执行。它们是有限精确算术、确定性数值回归和带固定种子的合成抽样，不能计作757个普遍定理或实验结果。下面保留全部代码与该次实际环境/输出；逐项结果保存在本书evidence/v3.0目录。

第49--53章的Lean形式化、100项运行测试、4个演示及具体证明见证另见附录F，不与这些检查混计。下列脚本不调用网络、不读取实验数据，依赖Python3、NumPy与SciPy，可分别保存并直接运行。

\hypertarget{d.2-verify_ueot_quantum.py}{%
\section{D.2
verify\_\allowbreak ueot\_\allowbreak quantum.py}\label{d.2-verify_ueot_quantum.py}}

代码SHA-256：79dc78035c239957\allowbreak a8ac4f5b2f6a9d1c\allowbreak 1a6b2a84d2cee094\allowbreak f831bad531f983c0

\begin{Verbatim}[breaklines,breakanywhere,fontsize=\scriptsize]
"""UEOT quantum master: deterministic finite/numerical regression, not formal proof."""
import json
import math
import platform
from fractions import Fraction as F
import numpy as np
import scipy
from scipy.linalg import expm, logm
from scipy.integrate import quad, dblquad
from numpy.polynomial.legendre import leggauss

rng = np.random.default_rng(31011)
checks = []
metrics = {}
I = np.eye(2, dtype=complex)
X = np.array([[0, 1], [1, 0]], complex)
Y = np.array([[0, -1j], [1j, 0]], complex)
Z = np.diag([1, -1]).astype(complex)
paulis = [I, X, Y, Z]

def check(name, value, detail=None):
    ok = bool(value)
    checks.append({"name": name, "passed": ok})
    if not ok:
        raise AssertionError((name, detail))

def close(name, a, b, atol=1e-10, rtol=1e-9):
    check(name, np.allclose(a, b, atol=atol, rtol=rtol),
          float(np.max(np.abs(np.asarray(a)-np.asarray(b)))))

def state(d):
    a = rng.normal(size=(d,d)) + 1j*rng.normal(size=(d,d))
    a = a@a.conj().T
    return a/np.trace(a)

def pure(v):
    return np.outer(v, np.conj(v))

def apply(ks, rho):
    return sum(k@rho@k.conj().T for k in ks)

def trace_out_second(rho, da, db):
    return np.einsum("abcb->ac",rho.reshape(da,db,da,db))

def td(a,b):
    return float(np.sum(np.abs(np.linalg.eigvalsh((a-b+(a-b).conj().T)/2)))/2)

def choi(ks):
    d = ks[0].shape[1]
    blocks = []
    for i in range(d):
        row = []
        for j in range(d):
            e = np.zeros((d,d),complex); e[i,j] = 1
            row.append(apply(ks,e))
        blocks.append(row)
    return np.block(blocks)

def herm_min(a):
    return float(np.linalg.eigvalsh((a+a.conj().T)/2).min())

# Complex channel construction and input-first Choi reconstruction.
for rep in range(12):
    z = rng.normal(size=(6,2))+1j*rng.normal(size=(6,2))
    v,_ = np.linalg.qr(z)
    ks = [v[i*2:(i+1)*2] for i in range(3)]
    j = choi(ks); rho = state(2)
    close(f"channel_TP_{rep}",sum(k.conj().T@k for k in ks),I)
    check(f"channel_Choi_PSD_{rep}",herm_min(j)>-1e-12)
    close(f"channel_Choi_partial_trace_{rep}",trace_out_second(j,2,2),I)
    product = np.kron(rho.T,I)@j
    recovered = np.einsum("iaib->ab",product.reshape(2,2,2,2))
    close(f"channel_reconstruction_{rep}",recovered,apply(ks,rho))
    sig=state(2)
    check(f"channel_trace_contraction_{rep}",
          td(apply(ks,rho),apply(ks,sig)) <= td(rho,sig)+1e-12)

bell=np.array([1,0,0,1],complex)/np.sqrt(2)
bellrho=pure(bell)
pt=bellrho.reshape(2,2,2,2).transpose(0,3,2,1).reshape(4,4)
close("transpose_nonCP_eigenvalue",np.linalg.eigvalsh(pt).min(),-0.5)
rho=state(2)
r=[np.trace(rho@s).real for s in [X,Y,Z]]
close("Pauli_state_tomography",rho,(I+sum(a*s for a,s in zip(r,[X,Y,Z])))/2)

# Rebit deletion witness.
rplus=(np.eye(4)+np.kron(Y,Y))/4
rminus=(np.eye(4)-np.kron(Y,Y))/4
check("rebit_positive",min(herm_min(rplus),herm_min(rminus))>-1e-12)
check("rebit_distinct_real",td(rplus,rminus)>0.99 and np.max(abs(rplus.imag))<1e-12)
for aidx,a in enumerate([I,X,Z]):
    for bidx,b in enumerate([I,X,Z]):
        close(f"rebit_local_blind_{aidx}_{bidx}",np.trace((rplus-rminus)@np.kron(a,b)),0)

# Coherent write / public records / inverse write.
cnot=np.array([[1,0,0,0],[0,1,0,0],[0,0,0,1],[0,0,1,0]],complex)
initial=pure(np.array([1,0,1,0],complex)/np.sqrt(2))
written=cnot@initial@cnot
public=sum(p@written@p for p in [np.diag([1,1,0,0]),np.diag([0,0,1,1])])
close("write_isometry",written,bellrho)
close("public_record_probabilities",np.diag(written),np.diag(public))
close("coherent_parent_distinguishability",td(written,public),0.5)
close("inverse_write_exact",cnot@written@cnot,initial)
close("inverse_write_dephased_X",np.trace((cnot@public@cnot)@np.kron(X,I)),0)

# SRN Gram kernel, block-scalar Kraus, phase-sensitive tomography.
m=3
env=rng.normal(size=(m,4))+1j*rng.normal(size=(m,4))
env/=np.linalg.norm(env,axis=1)[:,None]
K=env@env.conj().T
dims=[2,1,2]; labels=np.repeat(np.arange(m),dims)
ks=[np.diag(env[labels,mu]) for mu in range(4)]
rho=state(5)
expected=K[labels[:,None],labels[None,:]]*rho
close("SRN_block_Kraus_Schur",apply(ks,rho),expected)
for a in range(m):
    idx=np.where(labels==a)[0]
    rr=np.zeros((5,5),complex)
    rr[np.ix_(idx,idx)]=state(len(idx))
    close(f"SRN_fixes_branch_internal_{a}",apply(ks,rr),rr)
close("SRN_composition",apply(ks,apply(ks,rho)),
      (K*K)[labels[:,None],labels[None,:]]*rho)
for a in range(m):
    for b in range(a+1,m):
        for phase in [0,np.pi/2]:
            v=np.zeros(m,complex);v[a]=1/np.sqrt(2);v[b]=np.exp(1j*phase)/np.sqrt(2)
            out=K*pure(v)
            xx=np.zeros((m,m),complex);xx[a,b]=xx[b,a]=1
            close(f"QDRR_phase_{a}_{b}_{phase:.3f}",
                  np.trace(xx@out).real,(np.exp(-1j*phase)*K[a,b]).real)
n=math.ceil(2*math.log(2*m*(m-1)/0.05)/0.12**2)
metrics["QDRR_shots_per_pair_phase"]=n
check("QDRR_sample_count",n==762)

for rep in range(12):
    u=rng.normal(size=(4,3))+1j*rng.normal(size=(4,3))
    G=u@u.conj().T; diag=G.diagonal().real
    theta=rng.normal(size=4)
    lam=G-(diag[:,None]+diag[None,:])/2+1j*(theta[:,None]-theta[None,:])
    C=np.eye(4)-np.ones((4,4))/4
    check(f"SRN_generator_CPSD_{rep}",herm_min(C@lam@C)>-1e-10)
    kt=np.exp(0.4*lam)
    check(f"SRN_generator_exponential_PSD_{rep}",herm_min(kt)>-1e-10)
    close(f"SRN_generator_diagonal_{rep}",kt.diagonal(),np.ones(4))
bad=np.array([[0,-1,-1],[-1,0,-10],[-1,-10,0]],float)
z=np.array([2,-1,-1])
close("invalid_generator_quadratic",z@bad@z,-12)
check("invalid_generator_fails_CP",herm_min(np.exp(0.001*bad))<0)

# Finite CV boundary, Gaussian, spin, Bell and control.
N=9
a=np.diag(np.sqrt(np.arange(1,N)),1)
close("truncated_CCR_boundary",a@a.T-a.T@a,
      np.eye(N)-N*np.diag([0]*(N-1)+[1]))
q=(a+a.T)/np.sqrt(2);p=(a-a.T)/(1j*np.sqrt(2))
close("oscillator_vacuum_uncertainty",(q@q)[0,0]*(p@p)[0,0],0.25)
Omega=np.array([[0,1],[-1,0]],float)
check("Gaussian_vacuum_covariance",herm_min(np.eye(2)/2+1j*Omega/2)>-1e-12)
check("Gaussian_subvacuum_rejected",herm_min(np.eye(2)*0.2+1j*Omega/2)<0)
eta=.7
check("Gaussian_loss_CP",herm_min((1-eta)*np.eye(2)/2+
      1j*(Omega-eta*Omega)/2)>-1e-12)
for t in [.1,.7,1.5]:
    H=(.8*X+.6*Z)/2
    u=expm(-1j*H*t)
    close(f"Rabi_{t}",abs(u[1,0])**2,.8**2*np.sin(t/2)**2)
close("spin_commutator",(X@Y-Y@X)/4,1j*Z/2)
close("spin_time_reversal_square",(1j*Y)@(1j*Y).conj(),-I)
B=np.kron(Z,(Z+X)/np.sqrt(2))+np.kron(Z,(Z-X)/np.sqrt(2))+np.kron(X,(Z+X)/np.sqrt(2))-np.kron(X,(Z-X)/np.sqrt(2))
close("CHSH_Tsirelson",np.vdot(bell,B@bell),2*np.sqrt(2))
Hs=(np.kron(X,X)+np.kron(Y,Y)+np.kron(Z,Z))/4
close("Heisenberg_spectrum",np.linalg.eigvalsh(Hs),[-.75,.25,.25,.25])
H=.3*X+.4*Z; O=Z; psi=np.array([1,0],complex); T=.7
U=expm(-1j*H*T); psit=U@psi
chi=U.conj().T@O@psit
# integrated derivative for constant pulse, compared to central finite difference
def grad_integrand(t):
    pt=expm(-1j*H*t)@psi
    ct=expm(-1j*H*(T-t)).conj().T@O@psit
    return 2*np.imag(np.vdot(ct,X@pt))
g=quad(grad_integrand,0,T)[0]
def objective(e):
    v=expm(-1j*(H+e*X)*T)@psi
    return np.vdot(v,O@v).real
close("control_gradient_sign",g,(objective(1e-6)-objective(-1e-6))/2e-6,atol=1e-8)

# Open dynamics, KMS, causal breaks, thermodynamics.
eta=np.exp(-.7)
damp=[np.diag([1,np.sqrt(eta)]),np.array([[0,np.sqrt(1-eta)],[0,0]])]
close("amplitude_damping_TP",sum(k.conj().T@k for k in damp),I)
close("amplitude_damping_population",apply(damp,np.diag([0,1]))[1,1],eta)
beta=.8; H=np.diag([0,1.3]); th=expm(-beta*H);th/=np.trace(th)
A=X+.2*Z; BB=Y+.3*X; t=.4
def tau(t):
    return expm(1j*H*t)@BB@expm(-1j*H*t)
close("Gibbs_KMS_complex_boundary",np.trace(th@A@tau(t+1j*beta)),
      np.trace(th@tau(t)@A))
def relent(r,s):
    return np.trace(r@(logm(r)-logm(s))).real
def entropy(r):
    ev=np.linalg.eigvalsh(r);ev=ev[ev>1e-14]
    return -np.sum(ev*np.log(ev))
rho=state(2)
free=lambda r:np.trace(r@H).real-entropy(r)/beta
close("Gibbs_free_energy_identity",free(rho)-free(th),relent(rho,th)/beta)
relax=lambda r:.63*r+.37*th*np.trace(r)
check("relative_entropy_stationary_contraction",relent(relax(rho),th)<=relent(rho,th)+1e-12)
# causal break resets S, leaves previous record in E; then E controls S
reverse=np.array([[1,0,0,0],[0,0,0,1],[0,0,1,0],[0,1,0,0]],complex)
outs=[]
for bit in [0,1]:
    inp=np.zeros(4);inp[2*bit]=1
    joint=cnot@pure(inp)@cnot
    env=np.einsum("abad->bd",joint.reshape(2,2,2,2))
    reset=np.kron(np.diag([1,0]),env)
    final=reverse@reset@reverse
    outs.append(trace_out_second(final,2,2))
close("causal_break_memory_witness",td(*outs),1)

# Clifford and Dirac mass shell.
zero=np.zeros((2,2),complex)
gam=[np.block([[I,zero],[zero,-I]])]
gam += [np.block([[zero,s],[-s,zero]]) for s in [X,Y,Z]]
metric=[1,-1,-1,-1]
for mu in range(4):
    for nu in range(4):
        close(f"Clifford_{mu}_{nu}",gam[mu]@gam[nu]+gam[nu]@gam[mu],
              2*metric[mu]*np.eye(4) if mu==nu else np.zeros((4,4)))
def slash(v):
    return sum(metric[i]*v[i]*gam[i] for i in range(4))
pp=np.array([np.sqrt(1+.2**2+.3**2),.2,.3,0])
close("Dirac_mass_shell_factorization",(slash(pp)-np.eye(4))@(slash(pp)+np.eye(4)),np.zeros((4,4)))
qq=np.array([.4,.2,-.1,.3])
close("QED_tree_Ward",slash(qq),(slash(pp+qq)-np.eye(4))-(slash(pp)-np.eye(4)))

# Independent numerical spin traces for e+e- -> mu+mu-.
for costh in [-.7,0,.4,.9]:
    E=2.0; st=np.sqrt(1-costh**2)
    p1=np.array([E,0,0,E]);p2=np.array([E,0,0,-E])
    k1=np.array([E,E*st,0,E*costh]);k2=np.array([E,-E*st,0,-E*costh])
    traces=0j
    for mu in range(4):
        for nu in range(4):
            L=np.trace(slash(p2)@gam[mu]@slash(p1)@gam[nu])/4
            R=np.trace(slash(k1)@gam[mu]@slash(k2)@gam[nu])
            traces+=metric[mu]*metric[nu]*L*R
    close(f"QED_scattering_trace_{costh}",traces/(4*E*E)**2,1+costh**2)
alpha=1/137.035999084
sigma=quad(lambda c:2*np.pi*alpha**2/(4*100)*(1+c*c),-1,1)[0]
close("QED_integrated_cross_section",sigma,4*np.pi*alpha**2/300)
metrics["QED_tree_cross_section_inverse_energy_squared"]=sigma

# phi4 renormalization and EFT.
mass=1.3; cutoff=5.0
numeric=quad(lambda q:q**3/(q*q+mass*mass),0,cutoff)[0]/(8*np.pi**2)
closed=(cutoff**2-mass**2*np.log1p(cutoff**2/mass**2))/(16*np.pi**2)
close("phi4_tadpole_radial_integral",numeric,closed)
f=lambda x:np.log((1+x*(1-x)*4)/(1+x*(1-x)*25))/(16*np.pi**2)
bubble=quad(f,0,1,epsabs=1e-13)[0]
gx,gw=leggauss(100)
close("phi4_bubble_two_quadratures",bubble,np.dot(gw,[f((x+1)/2) for x in gx])/2)
close("phi4_bubble_benchmark",bubble,-.006683040905266517)
metrics["phi4_one_loop_F"]=.4-1.5*.4**2*bubble
s=5.;lam=.2;vel=np.sqrt(1-4/s)
im=lam**2*vel/(32*np.pi)
phase=quad(lambda c:vel/(16*np.pi),-1,1)[0]
close("phi4_optical_cut",2*im,lam**2*phase/2)
mass=.7;M=10.;kap=3.;p2=4.
kernel=np.array([[p2-mass**2,-kap],[-kap,p2-M*M]])
schur=kernel[0,0]-kernel[0,1]*kernel[1,0]/kernel[1,1]
close("EFT_Schur_inverse",1/schur,np.linalg.inv(kernel)[0,0])
errors=[]
for order in range(3):
    approx=kap**2/M**2*sum((p2/M**2)**j for j in range(order+1))
    err=abs(kap**2/(M*M-p2)-approx)
    bound=kap**2/M**2*(p2/M**2)**(order+1)/(1-p2/M**2)
    close(f"EFT_remainder_{order}",err,bound)
    errors.append(err)
metrics["EFT_errors_LO_NLO_NNLO"]=errors
def pi_hat(aq):
    return 2*alpha/np.pi*quad(lambda x:x*(1-x)*np.log1p(aq*x*(1-x)),0,1,epsabs=1e-13)[0]
aq=1e-5
close("QED_polarization_IR",pi_hat(aq)/(alpha*aq/(15*np.pi)),1,atol=2e-6)
aq=1e7
close("QED_polarization_UV",pi_hat(aq),alpha/(3*np.pi)*(np.log(aq)-5/3),atol=1e-8)
def threshold(aq):
    return 4*alpha**2/np.pi*quad(lambda x:aq*x*x*(1-x)**2/(1+aq*x*(1-x)),0,1)[0]
close("QED_threshold_UV",threshold(1e7)/(2*alpha**2/(3*np.pi)),1,atol=2e-6)
aq=3.;step=1e-4
close("QED_threshold_derivative",threshold(aq),
      alpha*(pi_hat(aq*np.exp(2*step))-pi_hat(aq*np.exp(-2*step)))/(2*step),
      atol=1e-11)
schwinger=dblquad(lambda y,z:z/(1-z),0,1,lambda z:0,lambda z:1-z)[0]*alpha/np.pi
close("Schwinger_parameter_integral",schwinger,alpha/(2*np.pi))
metrics["Schwinger_one_loop"]=schwinger

# Higgs and seesaw.
g=.65;gp=.36;v=246.
mat=v*v/4*np.array([[g*g,-g*gp],[-g*gp,gp*gp]])
close("Higgs_neutral_eigenvalues",np.linalg.eigvalsh(mat),[0,v*v*(g*g+gp*gp)/4],atol=1e-9)
md=.1;mr=20.;ev=np.linalg.eigvalsh([[0,md],[md,mr]])
check("seesaw_controlled_error",abs(ev[0]+md*md/mr)<md**4/mr**3*1.01)

# Exact local anomaly functions; rows = multiplicity, C dimension/sign, L dimension/sign, X.
def anomalies(rows):
    out=[F(0) for _ in range(6)]
    for mult,dc,sc,dl,sl,x in rows:
        x=F(x)
        out[0]+=mult*dl*sc
        out[1]+=mult*dc*sl
        out[2]+=mult*dl*(F(1,2) if dc>1 else 0)*x
        out[3]+=mult*dc*(F(1,2) if dl>1 else 0)*x
        out[4]+=mult*dc*dl*x
        out[5]+=mult*dc*dl*x**3
    return out
rows331=[(3,1,0,3,1,F(-1,3)),(3,1,0,1,0,1),
         (1,3,1,3,1,F(1,3)),(2,3,1,3,-1,0),
         (4,3,-1,1,0,F(-2,3)),(5,3,-1,1,0,F(1,3))]
check("331RHN_all_six_local_anomalies",anomalies(rows331)==[0]*6)
metrics["331RHN_local_anomalies"]=[str(x) for x in anomalies(rows331)]
for nc in range(3,16):
    q=F(1,2*nc); u=-F(nc+1,2*nc);d=F(nc-1,2*nc);l=F(-1,2);e=F(1)
    vals=[2*q+u+d,nc*q+l,2*nc*q+nc*u+nc*d+2*l+e,
          2*nc*q**3+nc*u**3+nc*d**3+2*l**3+e**3]
    check(f"generalized_SM_local_anomalies_{nc}",vals==[0]*4)
# Generic beta8=sqrt(3)*b, exact rational grid, including E^c.
for nc in range(3,9):
    nf=2*nc; k=nc-1
    for b in [F(-1),F(-1,3),F(0),F(2,3),F(1)]:
        xl=-F(1,2)-b/2;xp=F(1,6)-b/2;xm=F(1,6)+b/2
        qe=-F(1,2)-3*b/2;jp=F(1,6)-3*b/2;jm=F(1,6)+3*b/2
        rows=[(nf,1,0,3,1,xl),(nf,1,0,1,0,1),(nf,1,0,1,0,-qe),
              (k,nc,1,3,1,xp),(nf-k,nc,1,3,-1,xm),
              (nf,nc,-1,1,0,F(-2,3)),(nf,nc,-1,1,0,F(1,3)),
              (k,nc,-1,1,0,-jp),(nf-k,nc,-1,1,0,-jm)]
        pred=[0,0,0,F(nf*(nc-3),12),0,-F(nf*(nc-3),4)*(1+3*b*b)]
        check(f"generalized_331_full_anomaly_{nc}_{b}",anomalies(rows)==pred)
survivors=[];loose=[]
for nc in range(2,31):
    for nf in range(1,101):
        for k in range(nf+1):
            if 2*nc*k==(nc-1)*nf and nf*(nc+1)<22:
                loose.append((nc,nf,k))
                if nc>=3:survivors.append((nc,nf,k))
check("G3_no_primitive_filter",survivors==[(3,3,1)])
check("G3_complex_color_needed",(2,4,1) in loose)
check("G3_AF_needed",2*3*2==(3-1)*6)
metrics["G3_survivors_without_atomicity"]=survivors
candidates=[(nc,nf) for nc in range(3,23) for nf in range(3,23) if nf*(nc+1)<22]
close("framing_candidate_count",len(candidates),7)
framing={str(chi):[(nc,nf) for nc,nf in candidates if nf*(4*nc+3+chi)%48==0] for chi in [0,1,17,33]}
check("framing_zero_empty",framing["0"]==[])
check("framing_one_unique",framing["1"]==[(3,3)])
check("framing_seventeen_competitor",(4,4) in framing["17"])
check("framing_thirtythree_competitors",all(x in framing["33"] for x in [(3,3),(3,4),(3,5)]))
metrics["framing_survivors"]=framing

result={"summary":{"status":"all_passed","checks":len(checks),"seed":31011,
                   "formal_proof_assistant_checked":False,"experimental_data_used":False},
        "environment":{"python":platform.python_version(),"numpy":np.__version__,"scipy":scipy.__version__},
        "metrics":metrics,"checks":checks}
print(json.dumps(result,ensure_ascii=False,indent=2))
\end{Verbatim}

3.0版实际结果摘要（代码在3.1--3.4版保持不变）：

\begin{Verbatim}[breaklines,breakanywhere,fontsize=\scriptsize]
{
  "summary": {
    "status": "all_passed",
    "checks": 237,
    "seed": 31011,
    "formal_proof_assistant_checked": false,
    "experimental_data_used": false
  },
  "environment": {
    "python": "3.12.14",
    "numpy": "2.3.5",
    "scipy": "1.17.0"
  },
  "metrics": {
    "QDRR_shots_per_pair_phase": 762,
    "QED_tree_cross_section_inverse_energy_squared": 2.230587522067968e-06,
    "phi4_one_loop_F": 0.40160392981726395,
    "EFT_errors_LO_NLO_NNLO": [
      0.0037500000000000033,
      0.00014999999999999736,
      5.999999999992123e-06
    ],
    "Schwinger_one_loop": 0.0011614097328858596,
    "331RHN_local_anomalies": [
      "0",
      "0",
      "0",
      "0",
      "0",
      "0"
    ],
    "G3_survivors_without_atomicity": [
      [
        3,
        3,
        1
      ]
    ],
    "framing_survivors": {
      "0": [],
      "1": [
        [
          3,
          3
        ]
      ],
      "17": [
        [
          3,
          3
        ],
        [
          4,
          4
        ]
      ],
      "33": [
        [
          3,
          3
        ],
        [
          3,
          4
        ],
        [
          3,
          5
        ]
      ]
    }
  }
}
\end{Verbatim}

\hypertarget{d.3-verify_frontier.py}{%
\section{D.3 verify\_\allowbreak frontier.py}\label{d.3-verify_frontier.py}}

代码SHA-256：f0eb63e7d743bb9e\allowbreak 807526a47708fc36\allowbreak b73e6762ab4fbd27\allowbreak 510927f3876094d2

\begin{Verbatim}[breaklines,breakanywhere,fontsize=\scriptsize]
"""Independent checks for UEOT physics frontier. No experimental data."""
import json, math, platform
from fractions import Fraction as F
import numpy as np
import scipy
from scipy.linalg import expm
from scipy.special import logsumexp
from scipy.stats import binom

rng=np.random.default_rng(9062026)
checks=[];metrics={}
def test(name,truth,detail=None):
    if not bool(truth):raise AssertionError((name,detail))
    checks.append({"name":name,"passed":True})
def close(name,a,b,tol=1e-10):
    test(name,np.allclose(a,b,atol=tol,rtol=tol),float(np.max(abs(np.asarray(a)-np.asarray(b)))))
def pure(v):return np.outer(v,v.conj())
def td(a,b):return np.linalg.eigvalsh((a-b+(a-b).conj().T)/2).__abs__().sum()/2
def off(a):return a-np.diag(np.diag(a))
def haar(n):
    z=rng.normal(size=(n,n))+1j*rng.normal(size=(n,n))
    q,r=np.linalg.qr(z)
    return q@np.diag(np.diag(r)/abs(np.diag(r)))
def psd(a):return np.linalg.eigvalsh((a+a.conj().T)/2).min()
def anomalies(rows):
    val=[F(0) for _ in range(6)]
    for n,dc,sc,dl,sl,x in rows:
        x=F(x)
        val[0]+=n*dl*sc
        val[1]+=n*dc*sl
        val[2]+=n*dl*(F(1,2) if dc>1 else 0)*x
        val[3]+=n*dc*(F(1,2) if dl>1 else 0)*x
        val[4]+=n*dc*dl*x
        val[5]+=n*dc*dl*x**3
    return val
def spectrum(nc,nf,b):
    k=F((nc-1)*nf,2*nc)
    assert k.denominator==1
    k=int(k)
    qu=F(nc+1,2*nc);qd=-F(nc-1,2*nc)
    xl=-F(1,2)-b/2;xp=F(1,2*nc)-b/2;xm=F(1,2*nc)+b/2
    qe=-F(1,2)-3*b/2;jp=F(1,2*nc)-3*b/2;jm=F(1,2*nc)+3*b/2
    return [(nf,1,0,3,1,xl),(nf,1,0,1,0,1),(nf,1,0,1,0,-qe),
            (k,nc,1,3,1,xp),(nf-k,nc,1,3,-1,xm),
            (nf,nc,-1,1,0,-qu),(nf,nc,-1,1,0,-qd),
            (k,nc,-1,1,0,-jp),(nf-k,nc,-1,1,0,-jm)]
def survivors(ns):
    ans=[]
    for nc in range(3,31):
        for nf in range(1,121):
            if (nc-1)*nf%(2*nc):continue
            bc=F(11*nc,3)-2*nf
            bl=11-F(nf*(nc+1),3)-F(ns,6)
            if bc>0 and bl>0:ans.append((nc,nf,(nc-1)*nf//(2*nc)))
    return ans
for ns in range(46):
    expected=[(3,3,1),(5,5,2)] if ns<=5 else ([(3,3,1)] if ns<=41 else [])
    test(f"UV_integer_classification_ns_{ns}",survivors(ns)==expected)
metrics["UV_three_scalar_survivors"]=survivors(3)
for nc in range(3,11):
    nf=2*nc
    for b in [F(-1),F(-1,3),F(0),F(1,2)]:
        test(f"generalized_complete_anomalies_{nc}_{b}",anomalies(spectrum(nc,nf,b))==[0]*6)
test("five_color_five_family_all_local_anomalies",anomalies(spectrum(5,5,F(-1,3)))==[0]*6)
metrics["five_color_coefficients"]={"bC":str(F(25,3)),"bL":str(F(1,2)),"low_b2":str(F(-17,6))}
# The zero-U(1) anomaly identities fail if generic charged exotic E^c is omitted.
rows=spectrum(5,5,F(0))
test("omitted_charged_Ec_detected",anomalies(rows[:2]+rows[3:])!=[0]*6)

# Flavor, computed through independent mass-basis rotations.
for rep in range(30):
    uu=haar(3);ud=haar(3);P=np.diag([1,0,0]);V=uu.conj().T@ud
    Pu=uu.conj().T@P@uu;Pd=ud.conj().T@P@ud
    close(f"flavor_CKM_transport_{rep}",Pd,V.conj().T@Pu@V)
    close(f"flavor_rank_one_minor_{rep}",Pu[0,1]*Pu[1,2],Pu[0,2]*Pu[1,1])
    close(f"flavor_cycle_real_{rep}",np.imag(Pu[0,1]*Pu[1,2]*Pu[2,0]),0)
    strength=np.linalg.norm(off(Pu))**2+np.linalg.norm(off(Pd))**2
    test(f"flavor_universal_lower_bound_{rep}",strength>=1-abs(V).max()-1e-12)
    C_u=-Pu**2/2;C_d=-Pd**2/2
    test(f"flavor_EFT_coefficient_bound_{rep}",
         np.sum(abs(off(C_u)))+np.sum(abs(off(C_d)))>=(1-abs(V).max())/2-1e-12)
# Illustrative, frozen angles; these are not a fresh fit to data.
t12,t23,t13,delta=.227,.041,.0035,1.2
c12,c23,c13=np.cos([t12,t23,t13]);s12,s23,s13=np.sin([t12,t23,t13])
V=np.array([[c12*c13,s12*c13,s13*np.exp(-1j*delta)],
            [-s12*c23-c12*s23*s13*np.exp(1j*delta),
             c12*c23-s12*s23*s13*np.exp(1j*delta),s23*c13],
            [s12*s23-c12*c23*s13*np.exp(1j*delta),
             -c12*s23-s12*c23*s13*np.exp(1j*delta),c23*c13]])
close("illustrative_CKM_unitarity",V.conj().T@V,np.eye(3))
metrics["illustrative_flavor_lower_bound_over_delta_g_squared"]=float(1-abs(V).max())

# Three-cycle complete positivity, phase loss, and the classical Poisson counterexample.
for rep in range(24):
    L=rng.normal(size=(3,2))+1j*rng.normal(size=(3,2))
    G=L@L.conj().T;di=np.diag(G).real;h=rng.normal(size=3)
    lam=G-(di[:,None]+di[None,:])/2-1j*(h[:,None]-h[None,:])
    rates=-lam.real
    omega=(lam[0,1]+lam[1,2]+lam[2,0]).imag
    bound=4*rates[0,2]*rates[1,2]-(rates[0,2]+rates[1,2]-rates[0,1])**2
    test(f"cycle_sharp_CP_{rep}",omega**2<=bound+1e-10)
    total=rates[0,1]+rates[1,2]+rates[2,0]
    test(f"cycle_phase_loss_{rep}",total>=np.sqrt(3)*abs(omega)-1e-10)
    K=np.exp(.3*lam)
    test(f"cycle_finite_CP_{rep}",psd(K)>-1e-10)
    determinant=1-abs(K[0,1])**2-abs(K[1,2])**2-abs(K[2,0])**2+2*np.real(K[0,1]*K[1,2]*K[2,0])
    close(f"cycle_determinant_identity_{rep}",determinant,np.linalg.det(K))
theta=np.array([0,2*np.pi/3,4*np.pi/3])
lam=np.exp(1j*(theta[:,None]-theta[None,:]))-1
omega=(lam[0,1]+lam[1,2]+lam[2,0]).imag
close("Poisson_saturates_phase_loss",sum(-lam.real[i,(i+1)%3] for i in range(3)),np.sqrt(3)*abs(omega))
test("Poisson_nonzero_cycle",abs(omega)>2)
metrics["classical_Poisson_cycle"]=float(omega)
bad=np.array([[0,-1+2j,-1],[-1-2j,0,-1],[-1,-1,0]])
test("excess_cycle_not_CP",psd(np.exp(.001*bad))<-1e-6)

# Four-level channel and witness.
V=np.array([[1,0],[0,1],[1,1],[1,1j]],complex);V[2:]/=np.sqrt(2)
K=V@V.conj().T
B=np.array([[-1/np.sqrt(2),-1/np.sqrt(2),1,0],
            [-1/np.sqrt(2),-1j/np.sqrt(2),0,1]],complex)
W=B.conj().T@B;kappa=3-2*np.sqrt(2);wmax=2+1/np.sqrt(2);E=W/wmax
close("four_level_nullspace",B@V,np.zeros((2,2)))
close("four_level_witness_zero",np.trace(W@K),0)
close("witness_frobenius",np.linalg.norm(W),3)
close("witness_operator_norm",np.linalg.eigvalsh(W).max(),wmax)
test("effect_valid",psd(E)>-1e-12 and psd(np.eye(4)-E)>-1e-12)
ks=[np.diag(V[:,mu]) for mu in range(2)]
close("four_level_Kraus_TP",sum(a.conj().T@a for a in ks),np.eye(4))
psi=np.ones(4)/2;rho=pure(psi)
out=sum(a@rho@a.conj().T for a in ks)
close("four_level_output",out,K/4)
joint=np.zeros((8,4),complex)
for a in range(4):joint[a*2:(a+1)*2,a]=V[a]
close("joint_isometry",joint.conj().T@joint,np.eye(4))
jointstate=joint@rho@joint.conj().T
reduced=np.einsum("abcb->ac",jointstate.reshape(4,2,4,2))
close("environment_dilation",reduced,out)
zsat=np.array([1,1,1,(1+1j)/np.sqrt(2)])
close("classical_bound_attained",np.vdot(zsat,W@zsat),kappa)
min_random=1e9
for rep in range(20):
    zz=np.exp(1j*rng.uniform(-np.pi,np.pi,size=(5000,4)))
    vals=np.einsum("ni,ij,nj->n",zz.conj(),W,zz).real
    min_random=min(min_random,float(vals.min()))
    test(f"classical_witness_random_batch_{rep}",vals.min()>=kappa-1e-12)
for eps in [0,.01,.04,kappa/4,.08,1]:
    Ke=(1-eps)*K+eps*np.eye(4)
    close(f"mixture_witness_{eps}",np.trace(W@Ke),4*eps)
pcl=kappa/(4*wmax)
n=math.ceil(math.log(.05)/math.log(1-pcl))
test("zero_event_exact_n",n==188)
test("zero_event_size",binom.cdf(0,n,pcl)<=.05)
test("one_less_shot_insufficient",binom.cdf(0,n-1,pcl)>.05)
metrics["four_level"]={"kappa":kappa,"p_classical_min":pcl,"ideal_zero_event_n95":n,
                      "noise_mixture_threshold":kappa/4,"random_min_witness":min_random,
                      "classical_pvalue_zero_188":float(binom.cdf(0,n,pcl)),
                      "population_leakage_positive_bound_limit":pcl**2/8}
# Verify the approximate population-preserving random-unitary bound on genuine unitary samples.
for rep in range(12):
    us=[]
    for j in range(4):
        h=rng.normal(size=(4,4))+1j*rng.normal(size=(4,4));h=(h+h.conj().T)/2
        us.append(expm(-1j*.002*h))
    weights=rng.dirichlet(np.ones(4))
    surv=sum(p*abs(np.diag(u))**2 for p,u in zip(weights,us))
    ell=1-surv.min()
    rout=sum(p*u@rho@u.conj().T for p,u in zip(weights,us))
    ds=[np.diag(np.diag(u)/abs(np.diag(u))) for u in us]
    rdiag=sum(p*d@rho@d.conj().T for p,d in zip(weights,ds))
    test(f"approximate_RU_distance_{rep}",td(rout,rdiag)<=np.sqrt(8*ell)+1e-12)
    test(f"approximate_RU_effect_{rep}",np.trace(E@rout).real>=pcl-np.sqrt(8*ell)-1e-12)

# Nonlinear ensemble inconsistency.
x1=.8*np.exp(-1)+.2*(-np.exp(-1))
x2=.5*.6*np.exp(-.36)+.5*.6*np.exp(-.36)
close("initial_ensemble_same_bloch",.8-.2,.6)
close("signalling_probability_difference",(x2-x1)/2,.0989390654698766)
metrics["nonlinear_signal_probability_difference"]=(x2-x1)/2

# Exact physical filtering sampler, no Euler step or negative probabilities.
p0=np.array([.2,.3,.5]);aval=np.array([-1.,.2,1.3]);gamma=.7;nsamp=120000
v0=np.dot(p0,aval**2)-np.dot(p0,aval)**2
filter_metrics=[]
for t in [.2,1.,4.]:
    labels=rng.choice(3,size=nsamp,p=p0)
    yt=2*np.sqrt(gamma)*aval[labels]*t+np.sqrt(t)*rng.normal(size=nsamp)
    logs=np.log(p0)[None,:]+2*np.sqrt(gamma)*yt[:,None]*aval-2*gamma*aval**2*t
    pp=np.exp(logs-logsumexp(logs,axis=1)[:,None])
    close(f"filter_normalization_{t}",pp.sum(axis=1),np.ones(nsamp))
    for a in range(3):
        stderr=pp[:,a].std(ddof=1)/np.sqrt(nsamp)
        test(f"filter_Born_martingale_{t}_{a}",abs(pp[:,a].mean()-p0[a])<=6*stderr+1e-4)
    bound=0.
    for a in range(3):
        for b in range(a+1,3):
            samples=np.sqrt(pp[:,a]*pp[:,b])
            theory=np.sqrt(p0[a]*p0[b])*np.exp(-gamma*(aval[a]-aval[b])**2*t/2)
            stderr=samples.std(ddof=1)/np.sqrt(nsamp)
            test(f"filter_pair_decay_{t}_{a}_{b}",abs(samples.mean()-theory)<=6*stderr+1e-4)
            bound+=theory
    variance=pp@(aval**2)-(pp@aval)**2
    test(f"filter_variance_bound_{t}",variance.mean()<=v0/(1+4*gamma*v0*t)+.004)
    err=1-pp.max(axis=1)
    test(f"filter_localization_bound_{t}",err.mean()<=bound+.002)
    filter_metrics.append({"t":t,"mean_p":pp.mean(axis=0).tolist(),
        "mean_variance":float(variance.mean()),"mean_unresolved":float(err.mean()),
        "proved_unresolved_bound":float(bound)})
metrics["exact_filtering"]=filter_metrics
# Heating commutator checked on an interior state, not as false full finite CCR.
a=np.diag(np.sqrt(np.arange(1,12)),1);q=(a+a.T)/np.sqrt(2);p=(a-a.T)/(1j*np.sqrt(2))
H=(p@p+q@q)/2
comm=lambda x,y:x@y-y@x
close("position_noise_heating_interior",-comm(q,comm(q,H))[0,0]/2,.5)
result={"summary":{"status":"all_passed","checks":len(checks),"seed":9062026,
        "experimental_data":False,"formal_proof_assistant":False},
        "environment":{"python":platform.python_version(),"numpy":np.__version__,"scipy":scipy.__version__},
        "metrics":metrics,"checks":checks}
print(json.dumps(result,ensure_ascii=False,indent=2))
\end{Verbatim}

3.0版实际结果摘要（代码在3.1--3.4版保持不变）：

\begin{Verbatim}[breaklines,breakanywhere,fontsize=\scriptsize]
{
  "summary": {
    "status": "all_passed",
    "checks": 423,
    "seed": 9062026,
    "experimental_data": false,
    "formal_proof_assistant": false
  },
  "environment": {
    "python": "3.12.14",
    "numpy": "2.3.5",
    "scipy": "1.17.0"
  },
  "metrics": {
    "UV_three_scalar_survivors": [
      [
        3,
        3,
        1
      ],
      [
        5,
        5,
        2
      ]
    ],
    "five_color_coefficients": {
      "bC": "25/3",
      "bL": "1/2",
      "low_b2": "-17/6"
    },
    "illustrative_flavor_lower_bound_over_delta_g_squared": 0.0008465021129667738,
    "classical_Poisson_cycle": -2.598076211353316,
    "four_level": {
      "kappa": 0.1715728752538097,
      "p_classical_min": 0.0158446719248555,
      "ideal_zero_event_n95": 188,
      "noise_mixture_threshold": 0.04289321881345243,
      "random_min_witness": 0.17513824113339804,
      "classical_pvalue_zero_188": 0.04965532503989115,
      "population_leakage_positive_bound_limit": 3.1381703550788005e-05
    },
    "nonlinear_signal_probability_difference": 0.0989390654698766,
    "exact_filtering": [
      {
        "t": 0.2,
        "mean_p": [
          0.19907099901956785,
          0.29966622147876876,
          0.5012627795016538
        ],
        "mean_variance": 0.5373589168712367,
        "mean_unresolved": 0.40514124557832,
        "proved_unresolved_bound": 0.7956706280194334
      },
      {
        "t": 1.0,
        "mean_p": [
          0.19923581748426195,
          0.30010875188535324,
          0.5006554306303944
        ],
        "mean_variance": 0.2081865339186897,
        "mean_unresolved": 0.21236481946950928,
        "proved_unresolved_bound": 0.4512082432205973
      },
      {
        "t": 4.0,
        "mean_p": [
          0.19889313068851436,
          0.2987103279415504,
          0.5023965413698956
        ],
        "mean_variance": 0.034615822280499135,
        "mean_unresolved": 0.03598012354677628,
        "proved_unresolved_bound": 0.10399501336004843
      }
    ]
  }
}
\end{Verbatim}

\hypertarget{d.4-verify_completion.py}{%
\section{D.4 verify\_\allowbreak completion.py}\label{d.4-verify_completion.py}}

代码SHA-256：c2107ac235f0d87f\allowbreak aee246412c6478a0\allowbreak 405113cb0113b647\allowbreak 027820106d95ad32

\begin{Verbatim}[breaklines,breakanywhere,fontsize=\scriptsize]
"""Independent finite checks for the unified additions; no experimental data."""
import itertools, json, platform
import numpy as np
import scipy
from scipy.integrate import quad
from scipy.linalg import expm
from scipy.stats import binom

rng=np.random.default_rng(9062048)
checks=[];metrics={}
def test(name,ok):
    if not bool(ok): raise AssertionError(name)
    checks.append(name)
def close(name,a,b,tol=1e-10):
    test(name,np.max(np.abs(np.asarray(a)-np.asarray(b)))<tol)

# Spectral gauge transport and flux periodicity, not finite exact CCR.
n=np.arange(-8,9);hbar=1.3;L=2.1;m=.7;q=.8
for j in range(5):
    theta,A,chi=rng.normal(size=3)
    energy=lambda th,aa: (hbar*(2*np.pi*n+th)/L-q*aa)**2/(2*m)
    close(f"ring_gauge_transport_{j}",energy(theta,A),energy(theta+q*chi*L/hbar,A+chi))
    close(f"ring_spectral_relabel_{j}",energy(theta+2*np.pi,A)[:-1],energy(theta,A)[1:])
close("Berry_Chern_integral",quad(lambda t:-.5*np.sin(t),0,np.pi)[0],-1)
for theta in [.2,.9,2.5]:
    close(f"Berry_patch_difference_{theta}",.5*(1+np.cos(theta))+.5*(1-np.cos(theta)),1)

for ratio in [.2,.7,1.4,2.0]:
    k=np.linspace(-np.pi,np.pi,10001)
    z=ratio+np.exp(-1j*k)
    winding=(np.unwrap(np.angle(z))[-1]-np.unwrap(np.angle(z))[0])/(2*np.pi)
    close(f"SSH_winding_{ratio}",winding,-1 if ratio<1 else 0)
for N,r in [(8,.4),(12,.7),(24,.8)]:
    H=np.zeros((2*N,2*N))
    for j in range(N):
        H[2*j,2*j+1]=H[2*j+1,2*j]=r
        if j+1<N:H[2*j+1,2*j+2]=H[2*j+2,2*j+1]=1.
    norm=np.sqrt((1-r*r)/(1-r**(2*N)))
    psi=np.zeros(2*N);psi[::2]=norm*(-r)**np.arange(N)
    residual=norm*r**N
    close(f"SSH_normalization_{N}",np.linalg.norm(psi),1)
    close(f"SSH_finite_residual_{N}",np.linalg.norm(H@psi),residual)
    test(f"SSH_near_zero_spectrum_{N}",min(abs(np.linalg.eigvalsh(H)))<=residual+1e-12)

# Direct noncommutative plaquette gauge covariance using SU(2) matrices.
pauli=[np.array([[0,1],[1,0]]),np.array([[0,-1j],[1j,0]]),np.diag([1,-1])]
def su2():
    v=rng.normal(size=3)
    return expm(1j*sum(v[i]*pauli[i] for i in range(3)))
for j in range(12):
    U=[su2() for _ in range(4)];G=[su2() for _ in range(4)]
    transformed=[G[k]@U[k]@G[(k+1)%4].conj().T for k in range(4)]
    prod=lambda a:a[0]@a[1]@a[2]@a[3]
    close(f"plaquette_gauge_covariance_{j}",prod(transformed),G[0]@prod(U)@G[0].conj().T)
    close(f"plaquette_trace_invariance_{j}",np.trace(prod(transformed)),np.trace(prod(U)))
F=(.3*pauli[0]+.7*pauli[2])/2;g=.8
for a in [.1,.05,.025]:
    exact=(1-np.trace(expm(1j*g*a*a*F)).real/2)*4/g**2
    predicted=a**4*np.trace(F@F).real
    test(f"Wilson_continuum_coefficient_{a}",abs(exact-predicted)/predicted<2e-6)

configs=np.array(list(itertools.product([-1,1],repeat=4)))
P=configs.prod(axis=1)
for beta in [0,.2,.7,1.4,3.]:
    weights=np.exp(beta*P);Z=weights.sum()
    close(f"Z2_partition_{beta}",Z,16*np.cosh(beta))
    close(f"Z2_plaquette_{beta}",weights@P/Z,np.tanh(beta))
    close(f"Z2_link_expectation_{beta}",weights@configs/Z,np.zeros(4))
for vertex in range(4):
    gauged=configs.copy();gauged[:,vertex]*=-1;gauged[:,(vertex-1)%4]*=-1
    close(f"Z2_vertex_gauge_{vertex}",gauged.prod(axis=1),P)

# Enumerate all histories in a genuinely history-dependent binary model.
p0=.13;ntrial=10;qalt=.02
histories=[((),1.)]
for t in range(ntrial):
    new=[]
    for hist,prob in histories:
        pt=p0+.3*((sum(hist)+t)%3)/2
        new += [(hist+(0,),prob*(1-pt)),(hist+(1,),prob*pt)]
    histories=new
distribution=np.zeros(ntrial+1)
for hist,prob in histories:distribution[sum(hist)]+=prob
close("history_tree_normalization",distribution.sum(),1)
for k in range(ntrial+1):
    test(f"dependent_binomial_domination_{k}",distribution[:k+1].sum()<=binom.cdf(k,ntrial,p0)+1e-12)
for pt in [p0,.3,.7,1.]:
    factor=pt*qalt/p0+(1-pt)*(1-qalt)/(1-p0)
    close(f"supermartingale_conditional_identity_{pt}",factor,1+(pt-p0)*(qalt-p0)/(p0*(1-p0)))
    test(f"supermartingale_factor_{pt}",factor<=1+1e-12)
likelihood=lambda k:(qalt/p0)**k*((1-qalt)/(1-p0))**(ntrial-k)
test("dependent_likelihood_mean",sum(distribution[k]*likelihood(k) for k in range(ntrial+1))<=1+1e-12)
test("marginal_bound_counterexample",1-p0>(1-p0)**ntrial)
pstar=(3-2*np.sqrt(2))/(4*(2+1/np.sqrt(2)))
test("ideal_zero_188_valid",(1-pstar)**188<.05)
test("ideal_zero_187_insufficient",(1-pstar)**187>.05)
metrics.update({"ideal_pstar":pstar,"zero_event_188_pvalue":(1-pstar)**188,
    "Berry_upper_band_Chern":-1,"history_model_zero_probability":distribution[0],
    "history_null_zero_bound":(1-p0)**ntrial})
result={"summary":{"status":"all_passed","checks":len(checks),"seed":9062048,
    "experimental_data":False,"formal_proof_assistant":False},
    "environment":{"python":platform.python_version(),"numpy":np.__version__,"scipy":scipy.__version__},
    "metrics":metrics,"checks":checks}
print(json.dumps(result,ensure_ascii=False,indent=2))
\end{Verbatim}

3.0版实际结果摘要（代码在3.1--3.4版保持不变）：

\begin{Verbatim}[breaklines,breakanywhere,fontsize=\scriptsize]
{
  "summary": {
    "status": "all_passed",
    "checks": 97,
    "seed": 9062048,
    "experimental_data": false,
    "formal_proof_assistant": false
  },
  "environment": {
    "python": "3.12.14",
    "numpy": "2.3.5",
    "scipy": "1.17.0"
  },
  "metrics": {
    "ideal_pstar": 0.0158446719248555,
    "zero_event_188_pvalue": 0.04965532503989112,
    "Berry_upper_band_Chern": -1,
    "history_model_zero_probability": 0.039600353867614305,
    "history_null_zero_bound": 0.24842341419143568
  }
}
\end{Verbatim}

\hypertarget{appendix-e}{%
\chapter{附录 E. 正式阅读版、编译与修订合同}\label{appendix-e}}

\hypertarget{e.1-ux5b8cux6574ux8bedux4e49ux4e3bux6e90ux4e0eux6d3eux751fux9605ux8bfbux7248}{%
\section{E.1
完整语义主源与派生阅读版}\label{e.1-ux5b8cux6574ux8bedux4e49ux4e3bux6e90ux4e0eux6d3eux751fux9605ux8bfbux7248}}

稳定主文件 \texttt{UEOT\_\allowbreak PHYSICS.\allowbreak md}
包含全部63章和6个附录。\texttt{UEOT\_\allowbreak PHYSICS.\allowbreak pdf}
是该文件实际编译的阅读版，包含数学排版、全文目录、稳定结果编号、原始推导和代码。第三版在完整第二版基础上进行有记录的实质修订，不是需要读者拼接的补丁稿。版本号属于版本身份，已解决/未解决的物理问题按正文和证据分别报告。

\hypertarget{e.2-ux53efux590dux73b0ux6784ux5efa}{%
\section{E.2 可复现构建}\label{e.2-ux53efux590dux73b0ux6784ux5efa}}

仓库中 \texttt{physics/qm/book/build/build\_\allowbreak pdf.\allowbreak py}
仅接受完整当前主源，不会运行旧版装配器或静默拉取其他稿件。需要Pandoc、XeLaTeX及常规TeX包；用
\texttt{-\/-font} 指定一个本地CJK字体，或安装Noto Sans CJK
SC。编译不启用shell-escape。实际阅读版沿用已恢复源包中的Droid Sans
Fallback字体，字体字节身份在PDF构建回执中记录；字体文件不是理论来源。

\begin{Verbatim}[breaklines,breakanywhere,fontsize=\scriptsize]
python3 physics/qm/book/build/build_pdf.py \
  --output-dir /tmp/ueot-physics-build \
  --font /path/to/UEOT_CJK.ttf
\end{Verbatim}

构建器先检查连续章号、稳定锚点、内部链接和未完成占位，再产生打印源与LaTeX，执行三遍XeLaTeX以稳定目录。源码、字体、TeX、PDF及日志分别记录哈希。最后必须检查实际页面的中文、公式、表格与代码换行，不能仅凭编译退出码宣布排版完成。

\hypertarget{e.3-ux7ef4ux62a4ux89c4ux5219}{%
\section{E.3 维护规则}\label{e.3-ux7ef4ux62a4ux89c4ux5219}}

新理论写回相关章节和编号注册表；代码变更回写形式化范围与实际回执；跨章节依赖变更记录在CHAPTER\_\allowbreak DOSSIER.json和审查说明。文档正式性来自连续内容、可追溯依据和完成的校核，不要求把尚未解决的物理问题写成已解结论。

编译输出的PDF哈希由外部回执记录，不写入PDF正文形成循环。远端源码提交、PR和主线工作流也各有实际回执，不把本地构建记录改称远端执行。本文正文所述数学与本地验证可先完整核查，随后用必要的集中提交同步相同字节。

\hypertarget{appendix-f}{%
\chapter{附录 F. 实际形式化范围与完整命名审计}\label{appendix-f}}

\hypertarget{f.1-ux672cux7248ux5b9eux9645ux68c0ux67e5ux8eabux4efdux4e0eux8ba1ux6570}{%
\section{F.1
本版实际检查身份与计数}\label{f.1-ux672cux7248ux5b9eux9645ux68c0ux67e5ux8eabux4efdux4e0eux8ba1ux6570}}

本版沿用Lean4.33.1（编译器提交819816b2e0a3bf405af45ae5c7af2491d8f5bee6）及Mathlib0df444a360eaa60ab8c11dca51a86af692955474；九项Git依赖实际身份与锁文件一致。35模块统一构建退出0，实际8799
jobs；747个显式定理与985项唯一完整命名公理输出得到核对。允许基础仅为
\texttt{propext}、\allowbreak \texttt{Classical.\allowbreak choice}、\allowbreak \texttt{Quot.\allowbreak sound}。实际构建保留普通linter警告，不声称零警告。

新增六模块共101个显式定理和130个审计目标。新增显式声明、命名非局部实例与结构构造器均逐项列出；继承855个请求原样保留。总数985由747个显式定理与238个列举非定理目标组成，不表示遍历每一个编译器生成内部常量，也不是新增物理规律的条数。

\begin{longtable}[]{@{}
  >{\raggedright\arraybackslash}p{(\columnwidth - 4\tabcolsep) * \real{0.3000}}
  >{\raggedleft\arraybackslash}p{(\columnwidth - 4\tabcolsep) * \real{0.1200}}
  >{\raggedright\arraybackslash}p{(\columnwidth - 4\tabcolsep) * \real{0.5800}}@{}}
\toprule\noalign{}
\begin{minipage}[b]{\linewidth}\raggedright
模块
\end{minipage} & \begin{minipage}[b]{\linewidth}\raggedleft
显式定理
\end{minipage} & \begin{minipage}[b]{\linewidth}\raggedright
精确责任
\end{minipage} \\
\midrule\noalign{}
\endhead
\bottomrule\noalign{}
\endlastfoot
Basic & 7 & 密度态、正迹和归一化 \\
Instrument & 14 & 实际Kraus分支、概率和后验 \\
Operational & 20 & 全部记录响应、响应商与可达更新 \\
Observables & 25 & 真实词效应、最小闭包和有限生成 \\
FiniteCoordinates & 21 & 局部表示到全词坐标和闭合后验 \\
Adaptive & 11 & 完整有限适应树、叶律与坐标 \\
Weighted & 16 & 非负权重到普通复Kraus操作 \\
Ragged & 19 & 结果依赖有限指标及补零保持 \\
Exact & 24 & 复有理解释的注入与代数保持 \\
ExactCertificate & 9 & 有限可计算检查器的可靠性与实例接口 \\
Transport & 4 & 既有Core误差运输的标量实例 \\
Selector & 11 & 全整数条件三代与删除前提反例 \\
Spectrum & 15 & 实际有限谱行到异常/窗口预算 \\
Effects & 30 & 全部实际效应、分支/通道拉回与词效应上界 \\
Stability & 30 & 全效应伪距离、实际通道缺陷/轨迹及后验 \\
Mixture & 23 & 实际混合、精确TV和有理预算 \\
Cesaro & 23 & 矩阵平均残差、紧性、连续性与平稳子列 \\
Survival & 32 & 实际对角存活记录及Core QSD/Doob消费 \\
ClassicalSector & 24 & 完整经典单纯形、全效应与任意有限辅助纯化障碍 \\
CorePurificationBoundary & 32 &
同一实际Core核、终端商、安全路径与纯化分离 \\
RealLocalTomography & 37 &
全部实局域效应不可区分与联合／复局域删除见证 \\
QuantumDoob & 32 & 实际Doob完备性、有限条件化与额外谱条件下平稳 \\
RecordProbability & 9 & 实际有限记录事件达到TV及全事件等价 \\
RectangularChannels & 29 &
异型Kraus通道、真实任意辅助放大、偏迹与效应收缩 \\
QuantumPurification & 25 &
完整凸极点、PSD导出Gram与任意有限密度态纯化 \\
OperationalSystems & 25 & 三种完整状态呈现、七项规律和独立候选原则 \\
ClassicalQuantumMemory & 39 &
实际联合记录、任意有限树记忆、全事件TV与标签缺陷 \\
QuantumDoobModel & 17 &
先前非对角模型提取及明确公共态；非新增独立物理 \\
QuantumDoobDynamics & 43 &
指定成功记录闭式、真实条件态到固定态的精确距离衰减与概率极限 \\
QuantumSeparation & 9 &
合法效应的Hermitian实张成、完整密度分离及距离零判据 \\
ComplexLocalTomography & 11 &
实际产品效应、实虚矩阵单位及任意联合态复局部层析 \\
PureStateCharacterization & 5 & PSD正列分解导出的完整凸纯态向量刻画 \\
PurificationUniqueness & 22 &
等Gram支撑等距、同辅助酉、异型等距嵌入与真实张量作用 \\
HistoryChannels & 23 &
读取历史的真实通道、共同PSD块权重、零块及合法效应界 \\
ActionRecordBudget & 31 &
原固定深度协议真实语义、逐步与逐动作误差和及完整记录TV \\
\end{longtable}

\hypertarget{f.2-ux8fd0ux884cux4e0eux5177ux4f53ux8bc1ux660eux4ea7ux54c1}{%
\section{F.2
运行与具体证明产品}\label{f.2-ux8fd0ux884cux4e0eux5177ux4f53ux8bc1ux660eux4ea7ux54c1}}

100项内核运行方法和4份演示与通用Lean数学分别计数。六份精确证书共42项审计；四类篡改同时满足原真值证明被拒绝和同数据假值证明通过，假值另有4项审计。本轮因公共数学接口扩充而在冻结的新源集上实际执行统一门；没有为v31.7纯元数据补录重跑旧科学证明。

持久化有限见证门包括七份继承见证和三份新见证：全效应与相干负控、纯性与相位敏感辅助酉、实际逐动作记录预算。十组共193个具体定理、71个定义、5个缩写及269个逐名审计，未计入主包747/985。真实物理实验、精确有限计算与通用形式证明分别维护。

首次统一验收的实际输入为89项，源集SHA-256为
\texttt{a34337cd75cd70b0\allowbreak f2e1125899c4dcd0\allowbreak 9aa7b132642adaad\allowbreak 2c17ef8a69f12f6e}。实际92项QM/Core源与olean绑定；56项实际递归Core源码与已闭环v31.7基线和上一版身份相符。局部编译、公理审计、统一验收与远端提交各自有回执。

发布前发现，现有CI原始文本预检查会把纯化唯一性模块一条英文文档注释中的``admit''误认为未完成证明。该自然语言单词仅改为``have''；原统一构建、清单和日志完整保留。\texttt{publication-source-reconciliation.\allowbreak json}
记录唯一差异、脱注释后代码与声明库存相同、其余88项原字节相同以及最终原始文本预检查通过；未为这个注释更改重跑数学。最终89项发布源码的集合SHA-256为
\texttt{054f2e01c50f25d2\allowbreak 519a2200eeda5ccd\allowbreak 9849c27cfc67141e\allowbreak eaaf552dd4d995d1}。8799
jobs仍绑定首次实际输入，不冒称是在最终原始字节上重新执行；实际远端CI另以最终提交与日志承担发布验收。

\hypertarget{f.3-985ux9879ux5b8cux6574ux547dux540dux516cux7406ux8bf7ux6c42}{%
\section{F.3
985项完整命名公理请求}\label{f.3-985ux9879ux5b8cux6574ux547dux540dux516cux7406ux8bf7ux6c42}}

下列逐项列出实际 \texttt{Axiom\allowbreak Audit.\allowbreak lean}
的完整请求。每项必须具有唯一同名实际输出；缺项、重复、未经审查的公理或构建错误都不算通过。科学范围以真实声明类型和第49--62章所列前提为准。

\hypertarget{f.3.1-ux8bf7ux6c42160}{%
\subsection{F.3.1 请求1--60}\label{f.3.1-ux8bf7ux6c42160}}

\begin{Verbatim}[breaklines,breakanywhere,fontsize=\scriptsize]
UEOTQM.Density.ext
UEOTQM.mass_nonneg
UEOTQM.trace_im_eq_zero
UEOTQM.mass_smul
UEOTQM.Density.mass_one
UEOTQM.mass_eq_zero_iff
UEOTQM.normalize_matrix
UEOTQM.krausMap_posSemidef
UEOTQM.krausMap_smul
UEOTQM.trace_krausMap
UEOTQM.Instrument.branch_posSemidef
UEOTQM.Instrument.branch_smul
UEOTQM.Instrument.sum_trace_branch
UEOTQM.Instrument.sum_mass_branch
UEOTQM.Instrument.probability_nonneg
UEOTQM.Instrument.sum_probability
UEOTQM.Instrument.probability_le_one
UEOTQM.Instrument.posterior_matrix
UEOTQM.Instrument.zero_probability_iff_zero_branch
UEOTQM.ancillaKrausMap_posSemidef
UEOTQM.ancillaKrausMap_entry
UEOTQM.recordOperator_posSemidef
UEOTQM.recordOperator_smul
UEOTQM.response_nonneg
UEOTQM.response_nil
UEOTQM.response_singleton
UEOTQM.totalRecordMass_eq_mass
UEOTQM.totalRecordMass_density
UEOTQM.responseClass_surjective
UEOTQM.existsUnique_responseDecoder
UEOTQM.probability_eq_of_response_eq
UEOTQM.response_posterior
UEOTQM.posterior_response_eq
UEOTQM.nextResponse_eq_none_iff
UEOTQM.nextResponse_eq_of_response_eq
UEOTQM.nextResponse_inputFiberCompatible
UEOTQM.existsUnique_reachableResponseUpdate
UEOTQM.Selector.selector_finite_bounds
UEOTQM.Selector.selects_three_without_atomicity
UEOTQM.Selector.dropping_color_bound_admits_four_families
UEOTQM.Selector.dropping_weak_budget_admits_six_families
UEOTQM.Selector.dropping_cubic_equation_admits_one_family
UEOTQM.Selector.dropping_nonempty_family_count_admits_empty_spectrum
UEOTQM.Selector.homogeneous_leptons_recover_cubic_equation
UEOTQM.Selector.mixed_lepton_template_admits_two_families
UEOTQM.Selector.cubic_cancellation_replicates
UEOTQM.Selector.replicated_weak_screening
UEOTQM.Selector.two_or_more_three_family_packets_fail_weak_budget
UEOTQM.Transport.scalarBudget_zero
UEOTQM.Transport.scalarBudget_succ
UEOTQM.Transport.scalarBudget_nonneg
UEOTQM.Transport.scalar_transport_certificate
UEOTQM.normalize
UEOTQM.Instrument.posterior
UEOTQM.nextReachableResponse_map_val
UEOTQM.nextReachableResponse_eq_of_response_eq
UEOTQM.nextReachableResponse_inputFiberCompatible
UEOTQM.existsUnique_closedReachableUpdate
UEOTQM.nextReachableResponse
UEOTQM.pairing_one
\end{Verbatim}

\hypertarget{f.3.2-ux8bf7ux6c4261120}{%
\subsection{F.3.2 请求61--120}\label{f.3.2-ux8bf7ux6c4261120}}

\begin{Verbatim}[breaklines,breakanywhere,fontsize=\scriptsize]
UEOTQM.pairing_zero_left
UEOTQM.pairing_add_left
UEOTQM.pairing_smul_left
UEOTQM.pairing_smul_right
UEOTQM.pairing_sum_left
UEOTQM.heisenberg_zero
UEOTQM.heisenberg_add
UEOTQM.heisenberg_smul
UEOTQM.trace_heisenberg
UEOTQM.pairing_heisenberg
UEOTQM.heisenberg_posSemidef
UEOTQM.recordEffect_posSemidef
UEOTQM.recordEffect_isHermitian
UEOTQM.mass_recordOperator_eq_pairing
UEOTQM.response_eq_pairing
UEOTQM.recordEffect_mem
UEOTQM.one_mem_observableSpace
UEOTQM.recordEffect_mem_of_closed
UEOTQM.observableSpace_le_of_closed
UEOTQM.observableSpace_pullbackClosed
UEOTQM.isHermitian_of_mem_observableSpace
UEOTQM.response_eq_iff_pairing_observableSpace
UEOTQM.observableSpace_finiteDimensional
UEOTQM.exists_finite_wordEffect_generators
UEOTQM.FiniteObservablePresentation.observable_mem_span
UEOTQM.FiniteObservablePresentation.observable_unit
UEOTQM.FiniteObservablePresentation.one_mem_span
UEOTQM.FiniteObservablePresentation.pullback_observable
UEOTQM.FiniteObservablePresentation.span_pullbackClosed
UEOTQM.FiniteObservablePresentation.span_eq_observableSpace
UEOTQM.FiniteObservablePresentation.recordEffect_mem_span
UEOTQM.FiniteObservablePresentation.operatorCoordinates_unit
UEOTQM.FiniteObservablePresentation.coordinates_unit
UEOTQM.FiniteObservablePresentation.coordinates_apply
UEOTQM.FiniteObservablePresentation.coordinates_eq_iff_response_eq
UEOTQM.FiniteObservablePresentation.operatorCoordinates_branch
UEOTQM.FiniteObservablePresentation.operatorCoordinates_record
UEOTQM.FiniteObservablePresentation.response_eq_recordCoordinates_unit
UEOTQM.FiniteObservablePresentation.probability_eq_applyBranch_unit
UEOTQM.FiniteObservablePresentation.coordinates_posterior
UEOTQM.FiniteObservablePresentation.coordinates_posterior_eq_normalizedBranch
UEOTQM.FiniteObservablePresentation.normalizedBranch_mem_range
UEOTQM.FiniteObservablePresentation.closedCoordinateUpdate_map_val
UEOTQM.FiniteObservablePresentation.coordinateUpdate_eq_posterior
UEOTQM.FiniteObservablePresentation.coordinateUpdate_eq_none_iff
UEOTQM.Adaptive.Protocol.word_injective
UEOTQM.Adaptive.sum_leafWeight
UEOTQM.Adaptive.leafWeight_nonneg
UEOTQM.Adaptive.leafProbability_eq_response
UEOTQM.Adaptive.leafProbability_nonneg
UEOTQM.Adaptive.sum_leafProbability
UEOTQM.Adaptive.leafProbability_le_one
UEOTQM.Adaptive.leafProbability_eq_of_response_eq
UEOTQM.Adaptive.leafProbability_eq_recordCoordinates_unit
UEOTQM.Adaptive.leafProbability_eq_of_coordinates_eq
UEOTQM.Adaptive.leafProbability_posterior
UEOTQM.Spectrum.quarkCount_le_length
UEOTQM.Spectrum.familyCubic_eq_counts
UEOTQM.Spectrum.leptonCount_eq_length_of_homogeneous
UEOTQM.Spectrum.ExtraMatter.screening_nonneg
\end{Verbatim}

\hypertarget{f.3.3-ux8bf7ux6c42121180}{%
\subsection{F.3.3 请求121--180}\label{f.3.3-ux8bf7ux6c42121180}}

\begin{Verbatim}[breaklines,breakanywhere,fontsize=\scriptsize]
UEOTQM.Spectrum.extraScreening_nonneg
UEOTQM.Spectrum.familyScreening_eq_of_active
UEOTQM.Spectrum.selector_anomaly_of_spectrum
UEOTQM.Spectrum.selector_budget_of_spectrum
UEOTQM.Spectrum.spectrum_selects_three
UEOTQM.Spectrum.canonical_rows_and_higgs
UEOTQM.Spectrum.mixed_two_family_rows
UEOTQM.Spectrum.replicated_six_family_rows
UEOTQM.Spectrum.compensated_one_family_rows
UEOTQM.Spectrum.zero_anomaly_unit_one_family_rows
UEOTQM.Spectrum.extra_screening_closes_positive_window
UEOTQM.pairingLinear
UEOTQM.heisenbergLinear
UEOTQM.FiniteObservablePresentation.closedCoordinateUpdate
UEOTQM.Adaptive.Protocol.leafFintype
UEOTQM.Spectrum.oneHiggs
UEOTQM.Spectrum.twentyHiggs
UEOTQM.sqrt_smul_gram
UEOTQM.sqrt_smul_krausTerm
UEOTQM.sqrt_smul_heisenbergTerm
UEOTQM.krausMap_sqrtWeightedKraus
UEOTQM.heisenberg_sqrtWeightedKraus
UEOTQM.weightedKrausMap_posSemidef
UEOTQM.weightedHeisenberg_posSemidef
UEOTQM.pairing_weightedKrausMap
UEOTQM.WeightedInstrument.toInstrument
UEOTQM.WeightedInstrument.toInstrument_kraus
UEOTQM.WeightedInstrument.toInstrument_branch
UEOTQM.WeightedInstrument.toInstrument_heisenberg
UEOTQM.WeightedInstrument.toInstrument_complete
UEOTQM.WeightedInstrument.branch_posSemidef
UEOTQM.WeightedInstrument.toInstrument_probability
UEOTQM.WeightedInstrument.sum_mass_branch
UEOTQM.WeightedInstrument.toInstrument_posterior_matrix
UEOTQM.zeroExtend_apply
UEOTQM.zeroExtend_of_not_mem_range
UEOTQM.sum_zeroExtend
UEOTQM.zeroExtend_nonneg
UEOTQM.weightedKrausMap_zeroExtend
UEOTQM.weightedHeisenberg_zeroExtend
UEOTQM.weightedGram_zeroExtend
UEOTQM.RaggedInstrument.toWeighted
UEOTQM.RaggedInstrument.toWeighted_weight_apply
UEOTQM.RaggedInstrument.toWeighted_kraus_apply
UEOTQM.RaggedInstrument.toWeighted_branch
UEOTQM.RaggedInstrument.toWeighted_toInstrument_branch
UEOTQM.RaggedInstrument.toWeighted_heisenberg
UEOTQM.RaggedInstrument.toWeighted_toInstrument_kraus_of_not_mem_range
UEOTQM.WeightedInstrument.padOutcomes
UEOTQM.WeightedInstrument.padOutcomes_branch_apply
UEOTQM.WeightedInstrument.padOutcomes_branch_of_not_mem_range
UEOTQM.WeightedInstrument.padOutcomes_heisenberg_apply
UEOTQM.WeightedInstrument.padOutcomes_heisenberg_of_not_mem_range
UEOTQM.WeightedInstrument.padOutcomes_probability_apply
UEOTQM.WeightedInstrument.padOutcomes_probability_of_not_mem_range
UEOTQM.Exact.MatrixData.ext
UEOTQM.Exact.MatrixData.toMatrix_re
UEOTQM.Exact.MatrixData.toMatrix_im
UEOTQM.Exact.MatrixData.toMatrix_injective
UEOTQM.Exact.MatrixData.toMatrix_eq_iff
\end{Verbatim}

\hypertarget{f.3.4-ux8bf7ux6c42181240}{%
\subsection{F.3.4 请求181--240}\label{f.3.4-ux8bf7ux6c42181240}}

\begin{Verbatim}[breaklines,breakanywhere,fontsize=\scriptsize]
UEOTQM.Exact.MatrixData.toMatrix_zero
UEOTQM.Exact.MatrixData.toMatrix_one
UEOTQM.Exact.MatrixData.toMatrix_add
UEOTQM.Exact.MatrixData.toMatrix_neg
UEOTQM.Exact.MatrixData.toMatrix_sub
UEOTQM.Exact.MatrixData.toMatrix_mul
UEOTQM.Exact.MatrixData.toMatrix_adjoint
UEOTQM.Exact.MatrixData.toMatrix_scale
UEOTQM.Exact.MatrixData.toMatrix_sum
UEOTQM.Exact.MatrixData.instDecidableEq
UEOTQM.Exact.toMatrix_weightedBranch
UEOTQM.Exact.toMatrix_weightedPullback
UEOTQM.Exact.toMatrix_weightedCompleteness
UEOTQM.Exact.toMatrix_weightedBranch_univ
UEOTQM.Exact.toMatrix_weightedPullback_univ
UEOTQM.Exact.Examples.imaginaryUnit_square_exact
UEOTQM.Exact.Examples.imaginaryUnit_adjoint_exact
UEOTQM.Exact.Examples.imaginaryUnit_square_complex
UEOTQM.Exact.Examples.hadamard_complete_exact
UEOTQM.Exact.Examples.hadamard_complete_complex
UEOTQM.ExactCertificate.check_eq_true_iff
UEOTQM.ExactCertificate.check_sound
UEOTQM.ExactCertificate.branch_eq_exact
UEOTQM.ExactCertificate.heisenberg_eq_exact
UEOTQM.ExactCertificate.exactRecordEffect_toMatrix
UEOTQM.ExactCertificate.basis_toMatrix
UEOTQM.ExactCertificate.coordinates_eq_iff_response_eq
UEOTQM.ExactCertificate.response_eq_recordCoordinates_unit
UEOTQM.ExactCertificate.adaptive_leafProbability_eq
UEOTQM.ExactCertificate.conditionsDecidable
UEOTQM.ExactCertificate.weightedInstrument
UEOTQM.ExactCertificate.instrumentFamily
UEOTQM.ExactCertificate.presentation
UEOTQM.pairing_sub_left
UEOTQM.pairing_add_right
UEOTQM.pairing_sum_right
UEOTQM.pairing_nonneg
UEOTQM.Effect.ext
UEOTQM.Effect.zero
UEOTQM.Effect.one
UEOTQM.Effect.complement
UEOTQM.Effect.zero_matrix
UEOTQM.Effect.one_matrix
UEOTQM.Effect.complement_matrix
UEOTQM.Effect.pairing_nonneg
UEOTQM.Effect.pairing_le_mass
UEOTQM.Effect.probability
UEOTQM.Effect.probability_nonneg
UEOTQM.Effect.probability_le_one
UEOTQM.Effect.probability_zero
UEOTQM.Effect.probability_one
UEOTQM.Effect.probability_complement
UEOTQM.Instrument.channel
UEOTQM.Instrument.channel_posSemidef
UEOTQM.Instrument.trace_channel
UEOTQM.Instrument.channelState
UEOTQM.Instrument.channelState_matrix
UEOTQM.Instrument.pullbackEffects
UEOTQM.Instrument.pullbackEffects_matrix
UEOTQM.Instrument.pairing_pullbackEffects
\end{Verbatim}

\hypertarget{f.3.5-ux8bf7ux6c42241300}{%
\subsection{F.3.5 请求241--300}\label{f.3.5-ux8bf7ux6c42241300}}

\begin{Verbatim}[breaklines,breakanywhere,fontsize=\scriptsize]
UEOTQM.Instrument.pullbackEffect
UEOTQM.Instrument.pairing_pullbackEffect
UEOTQM.Instrument.probability_channelState
UEOTQM.Instrument.branchEffect
UEOTQM.Instrument.branchEffect_matrix
UEOTQM.Instrument.pairing_branchEffect
UEOTQM.Instrument.probability_branchEffect_one
UEOTQM.Instrument.probability_posterior
UEOTQM.recordTest
UEOTQM.recordTest_matrix
UEOTQM.recordEffect_complement_posSemidef
UEOTQM.response_eq_effect_probability
UEOTQM.response_le_one
UEOTQM.effectGapSet
UEOTQM.Effect.probability_abs_sub_le_one
UEOTQM.effectGapSet_nonempty
UEOTQM.effectGapSet_bddAbove
UEOTQM.effectDistance
UEOTQM.Effect.probability_sub_le_effectDistance
UEOTQM.effectDistance_nonneg
UEOTQM.effectDistance_le_iff
UEOTQM.effectDistance_le_one
UEOTQM.effectDistance_self
UEOTQM.effectDistance_comm
UEOTQM.effectDistance_triangle
UEOTQM.effectDistance_eq_zero_iff
UEOTQM.response_sub_le_effectDistance
UEOTQM.channelDefect
UEOTQM.channelDefect_bddAbove
UEOTQM.channelDefect_nonneg
UEOTQM.channelDefect_le_one
UEOTQM.effectDistance_le_channelDefect
UEOTQM.channelDefect_self
UEOTQM.Instrument.channelState_effectDistance_le
UEOTQM.Instrument.probability_sub_le_effectDistance
UEOTQM.Instrument.probability_pos_of_effectDistance_lt
UEOTQM.Instrument.channelState_perturbation_bound
UEOTQM.normalized_probability_difference_le
UEOTQM.Instrument.posterior_effectDistance_le
UEOTQM.Instrument.posterior_effectDistance_le_of_probability_lower_bound
UEOTQM.channelTrajectory
UEOTQM.channelTrajectory_zero
UEOTQM.channelTrajectory_succ
UEOTQM.channelTrajectory_transport_certificate
UEOTQM.scalarBudget_one_eq_sum
UEOTQM.channelTrajectory_effectDistance_le_sum
UEOTQM.channelTrajectory_response_sub_le_sum
UEOTQM.Mixture.tv
UEOTQM.Mixture.tv_nonneg
UEOTQM.Mixture.tv_self
UEOTQM.Mixture.tv_comm
UEOTQM.Mixture.tv_le_one
UEOTQM.Mixture.expectation_sub_le_tv
UEOTQM.Mixture.density
UEOTQM.Mixture.density_matrix
UEOTQM.Mixture.probability_density
UEOTQM.Mixture.density_weight_change_le_tv
UEOTQM.Mixture.density_fixed_weights_le_sum
UEOTQM.Mixture.density_joint_change_le
UEOTQM.Mixture.weightedInstrument
\end{Verbatim}

\hypertarget{f.3.6-ux8bf7ux6c42301360}{%
\subsection{F.3.6 请求301--360}\label{f.3.6-ux8bf7ux6c42301360}}

\begin{Verbatim}[breaklines,breakanywhere,fontsize=\scriptsize]
UEOTQM.Mixture.instrument
UEOTQM.Mixture.instrument_kraus
UEOTQM.Mixture.instrument_branch
UEOTQM.Mixture.instrument_heisenberg
UEOTQM.Mixture.instrument_channel
UEOTQM.Mixture.instrument_channelState
UEOTQM.Mixture.channelDefect_weight_change_le_tv
UEOTQM.Mixture.channelDefect_fixed_weights_le_sum
UEOTQM.Mixture.channelDefect_joint_change_le
UEOTQM.Mixture.RationalWeights.toWeights
UEOTQM.Mixture.rationalTV
UEOTQM.Mixture.tv_toWeights
UEOTQM.Mixture.density_weight_change_le_rational_budget
UEOTQM.Mixture.channelDefect_weight_change_le_rational_budget
UEOTQM.Mixture.channelDefect_fixed_weights_le_rational_sum
UEOTQM.Mixture.rationalJointBudget
UEOTQM.Mixture.channelDefect_joint_change_le_rational_budget
UEOTQM.Cesaro.iterateState
UEOTQM.Cesaro.iterateState_zero
UEOTQM.Cesaro.iterateState_succ
UEOTQM.Cesaro.channelLinear
UEOTQM.Cesaro.channelLinear_apply
UEOTQM.Cesaro.orbitSum
UEOTQM.Cesaro.orbitSum_positive
UEOTQM.Cesaro.orbitSum_trace
UEOTQM.Cesaro.averageState
UEOTQM.Cesaro.averageState_matrix
UEOTQM.Cesaro.orbitSum_telescope
UEOTQM.Cesaro.matrix_residual
UEOTQM.Cesaro.pairing_sub_right
UEOTQM.Cesaro.probability_residual
UEOTQM.Cesaro.probability_residual_le
UEOTQM.Cesaro.effectDistance_residual_le
UEOTQM.Cesaro.density_operatorNorm_le_one
UEOTQM.Cesaro.densityMatrices
UEOTQM.Cesaro.densityMatrices_closed
UEOTQM.Cesaro.densityMatrices_compact
UEOTQM.Cesaro.continuous_channel
UEOTQM.Cesaro.matrix_residual_norm_le
UEOTQM.Cesaro.averageSequence
UEOTQM.Cesaro.matrix_residual_tendsto_zero
UEOTQM.Cesaro.subsequence_limit_invariant
UEOTQM.Cesaro.subsequence_limit_density
UEOTQM.Cesaro.exists_stationary_subsequence
UEOTQM.Cesaro.exists_stationary_of_density
UEOTQM.Cesaro.normalizedIdentity
UEOTQM.Cesaro.exists_stationary
UEOTQM.Survival.diagonalOperator
UEOTQM.Survival.diagonalOperator_posSemidef
UEOTQM.Survival.mass_diagonalOperator
UEOTQM.Survival.diagonalOperator_injective
UEOTQM.Survival.diagonalOperator_smul
UEOTQM.Survival.diagonalDensity
UEOTQM.Survival.diagonalDensity_matrix
UEOTQM.Survival.matrixUnit
UEOTQM.Survival.matrixUnit_gram
UEOTQM.Survival.matrixUnit_sandwich
UEOTQM.Survival.measurePrepare
UEOTQM.Survival.measurePrepare_apply
UEOTQM.Survival.measurePrepare_diagonal
\end{Verbatim}

\hypertarget{f.3.7-ux8bf7ux6c42361420}{%
\subsection{F.3.7 请求361--420}\label{f.3.7-ux8bf7ux6c42361420}}

\begin{Verbatim}[breaklines,breakanywhere,fontsize=\scriptsize]
UEOTQM.Survival.measurePrepare_gram
UEOTQM.Survival.failureKernel
UEOTQM.Survival.failureKernel_nonneg
UEOTQM.Survival.outcomeKernel
UEOTQM.Survival.outcomeKernel_complete
UEOTQM.Survival.weightedInstrument
UEOTQM.Survival.instrument
UEOTQM.Survival.instrument_branch
UEOTQM.Survival.success_branch
UEOTQM.Survival.rowApply_nonneg
UEOTQM.Survival.success_probability
UEOTQM.Survival.success_posterior_matrix
UEOTQM.Survival.killedIterate
UEOTQM.Survival.killedIterate_posSemidef
UEOTQM.Survival.killedIterate_diagonal
UEOTQM.Survival.survivalMass
UEOTQM.Survival.survivalMass_eq_response
UEOTQM.Survival.survivalMass_diagonal
UEOTQM.Survival.conditionalState
UEOTQM.Survival.qsd_killedIterate
UEOTQM.Survival.qsd_survival_pow
UEOTQM.Survival.qsd_survival_pos
UEOTQM.Survival.qsd_conditionalState
UEOTQM.Survival.qsd_rate_le_one
UEOTQM.Survival.doobTransform_nonneg
UEOTQM.Survival.doobInstrument
UEOTQM.Survival.doobDensity
UEOTQM.Survival.doob_success_branch
UEOTQM.Survival.doob_success_probability
UEOTQM.Survival.doob_killedIterate
UEOTQM.Survival.doob_survivalMass
UEOTQM.Survival.doob_conditionalState
UEOTQM.Classical.State
UEOTQM.Classical.Effect
UEOTQM.Classical.Effect.mk
UEOTQM.Classical.state_ext
UEOTQM.Classical.evaluate
UEOTQM.Classical.evaluate_nonneg
UEOTQM.Classical.evaluate_le_one
UEOTQM.Classical.indicator
UEOTQM.Classical.evaluate_indicator
UEOTQM.Classical.all_effects_separate
UEOTQM.Classical.dirac
UEOTQM.Classical.dirac_weight
UEOTQM.Classical.evaluate_dirac
UEOTQM.Classical.mix
UEOTQM.Classical.mix_weight
UEOTQM.Classical.halfMix
UEOTQM.Classical.halfMix_weight
UEOTQM.Classical.pushforward
UEOTQM.Classical.pushforward_weight
UEOTQM.Classical.pushforward_dirac
UEOTQM.Classical.pushforward_halfMix
UEOTQM.Classical.marginalLeft
UEOTQM.Classical.marginalLeft_weight
UEOTQM.Classical.marginalLeft_eq_pushforward
UEOTQM.Classical.marginalLeft_dirac
UEOTQM.Classical.marginalLeft_halfMix
UEOTQM.Classical.IsPure
UEOTQM.Classical.dirac_isPure
\end{Verbatim}

\hypertarget{f.3.8-ux8bf7ux6c42421480}{%
\subsection{F.3.8 请求421--480}\label{f.3.8-ux8bf7ux6c42421480}}

\begin{Verbatim}[breaklines,breakanywhere,fontsize=\scriptsize]
UEOTQM.Classical.dirac_ne
UEOTQM.Classical.halfMix_not_pure
UEOTQM.Classical.halfRow
UEOTQM.Classical.halfRow_weight
UEOTQM.Classical.marginalLeft_halfRow
UEOTQM.Classical.no_pure_extension_halfMix
UEOTQM.Classical.attach
UEOTQM.Classical.marginalLeft_attach
UEOTQM.Classical.dirac_has_pure_extension
UEOTQM.Classical.toPMF
UEOTQM.Classical.toPMF_apply
UEOTQM.Classical.toPMF_apply_toReal
UEOTQM.Classical.toPMF_dirac
UEOTQM.Classical.toPMF_pushforward
UEOTQM.Classical.rowState
UEOTQM.Classical.rowState_weight
UEOTQM.Classical.stepState
UEOTQM.Classical.stepState_dirac
UEOTQM.Classical.toPMF_stepState
UEOTQM.CorePurificationBoundary.Micro
UEOTQM.CorePurificationBoundary.fullyObservedSurvivalModel
UEOTQM.CorePurificationBoundary.terminal_eq_iff
UEOTQM.CorePurificationBoundary.Terminal
UEOTQM.CorePurificationBoundary.terminalFintype
UEOTQM.CorePurificationBoundary.terminalMeasurableSpace
UEOTQM.CorePurificationBoundary.terminalControl
UEOTQM.CorePurificationBoundary.terminalGood
UEOTQM.CorePurificationBoundary.terminalControl_source
UEOTQM.CorePurificationBoundary.terminalControl_reward
UEOTQM.CorePurificationBoundary.terminalControl_optimalValue
UEOTQM.CorePurificationBoundary.terminalControl_injective
UEOTQM.CorePurificationBoundary.terminalGood_source
UEOTQM.CorePurificationBoundary.viable_all_horizons
UEOTQM.CorePurificationBoundary.initial_support_viable
UEOTQM.CorePurificationBoundary.firstRow
UEOTQM.CorePurificationBoundary.firstRow_halfMix
UEOTQM.CorePurificationBoundary.firstQuotientState
UEOTQM.CorePurificationBoundary.quotient_support_distinct
UEOTQM.CorePurificationBoundary.firstQuotientState_halfMix
UEOTQM.CorePurificationBoundary.firstQuotientState_not_pure
UEOTQM.CorePurificationBoundary.firstQuotientState_no_pure_extension
UEOTQM.CorePurificationBoundary.sourcePMF
UEOTQM.CorePurificationBoundary.piStar
UEOTQM.CorePurificationBoundary.actualPathLaw
UEOTQM.CorePurificationBoundary.sourcePMF_toReal
UEOTQM.CorePurificationBoundary.sourcePMF_firstRow
UEOTQM.CorePurificationBoundary.sourcePMF_rowState
UEOTQM.CorePurificationBoundary.piStar_preserves
UEOTQM.CorePurificationBoundary.actualPath_all_times_safe
UEOTQM.CorePurificationBoundary.actualPath_first_quotient_marginal
UEOTQM.CorePurificationBoundary.false_action_destructive
UEOTQM.CorePurificationBoundary.FiniteCoreSlice
UEOTQM.CorePurificationBoundary.finiteCoreSlice_holds
UEOTQM.CorePurificationBoundary.classical_purification_separation
UEOTQM.CorePurificationBoundary.RestrictedTerminal
UEOTQM.CorePurificationBoundary.restrictedTerminalFintype
UEOTQM.CorePurificationBoundary.restrictedFirstQuotientState
UEOTQM.CorePurificationBoundary.restricted_support_merge
UEOTQM.CorePurificationBoundary.restrictedFirstQuotientState_dirac
UEOTQM.CorePurificationBoundary.restrictedFirstQuotientState_pure
\end{Verbatim}

\hypertarget{f.3.9-ux8bf7ux6c42481540}{%
\subsection{F.3.9 请求481--540}\label{f.3.9-ux8bf7ux6c42481540}}

\begin{Verbatim}[breaklines,breakanywhere,fontsize=\scriptsize]
UEOTQM.RealLocalTomography.IsReal
UEOTQM.RealLocalTomography.symmetric_of_real_hermitian
UEOTQM.RealLocalTomography.productEffect
UEOTQM.RealLocalTomography.productEffect_matrix
UEOTQM.RealLocalTomography.plusProjector
UEOTQM.RealLocalTomography.minusProjector
UEOTQM.RealLocalTomography.plusProjector_hermitian
UEOTQM.RealLocalTomography.minusProjector_hermitian
UEOTQM.RealLocalTomography.plusProjector_square
UEOTQM.RealLocalTomography.minusProjector_square
UEOTQM.RealLocalTomography.plusProjector_mul_minusProjector
UEOTQM.RealLocalTomography.one_sub_plusProjector
UEOTQM.RealLocalTomography.plusProjector_posSemidef
UEOTQM.RealLocalTomography.minusProjector_posSemidef
UEOTQM.RealLocalTomography.involutionEffect
UEOTQM.RealLocalTomography.Qubit
UEOTQM.RealLocalTomography.TwoQubit
UEOTQM.RealLocalTomography.pauliY
UEOTQM.RealLocalTomography.parity
UEOTQM.RealLocalTomography.pauliY_hermitian
UEOTQM.RealLocalTomography.pauliY_square
UEOTQM.RealLocalTomography.parity_hermitian
UEOTQM.RealLocalTomography.parity_square
UEOTQM.RealLocalTomography.pauliY_trace
UEOTQM.RealLocalTomography.parity_trace
UEOTQM.RealLocalTomography.parity_isReal
UEOTQM.RealLocalTomography.rhoPlus
UEOTQM.RealLocalTomography.rhoMinus
UEOTQM.RealLocalTomography.rhoPlus_matrix
UEOTQM.RealLocalTomography.rhoMinus_matrix
UEOTQM.RealLocalTomography.rhoPlus_isReal
UEOTQM.RealLocalTomography.rhoMinus_isReal
UEOTQM.RealLocalTomography.rhoPlus_symmetric
UEOTQM.RealLocalTomography.rhoMinus_symmetric
UEOTQM.RealLocalTomography.density_difference
UEOTQM.RealLocalTomography.trace_mul_pauliY_zero
UEOTQM.RealLocalTomography.trace_product_density_difference_zero
UEOTQM.RealLocalTomography.all_real_product_probabilities_equal
UEOTQM.RealLocalTomography.jointTest
UEOTQM.RealLocalTomography.jointTest_matrix
UEOTQM.RealLocalTomography.jointTest_isReal
UEOTQM.RealLocalTomography.jointTest_probability_plus
UEOTQM.RealLocalTomography.jointTest_probability_minus
UEOTQM.RealLocalTomography.rhoPlus_ne_rhoMinus
UEOTQM.RealLocalTomography.real_local_tomography_fails
UEOTQM.RealLocalTomography.complexLocalTest
UEOTQM.RealLocalTomography.complexLocalTest_not_real
UEOTQM.RealLocalTomography.complex_product_probability_plus
UEOTQM.RealLocalTomography.complex_product_probability_minus
UEOTQM.RealLocalTomography.complex_product_distinguishes
UEOTQM.QuantumDoob.successOperator
UEOTQM.QuantumDoob.successOperator_zero
UEOTQM.QuantumDoob.successOperator_succ
UEOTQM.QuantumDoob.successOperator_smul
UEOTQM.QuantumDoob.successOperator_posSemidef
UEOTQM.QuantumDoob.eigen_successOperator
UEOTQM.QuantumDoob.successMass
UEOTQM.QuantumDoob.successMass_eq_response
UEOTQM.QuantumDoob.eigen_successMass
UEOTQM.QuantumDoob.eigen_successMass_pos
\end{Verbatim}

\hypertarget{f.3.10-ux8bf7ux6c42541600}{%
\subsection{F.3.10 请求541--600}\label{f.3.10-ux8bf7ux6c42541600}}

\begin{Verbatim}[breaklines,breakanywhere,fontsize=\scriptsize]
UEOTQM.QuantumDoob.conditionalState
UEOTQM.QuantumDoob.eigen_conditionalState
UEOTQM.QuantumDoob.eigen_rate_nonneg
UEOTQM.QuantumDoob.eigen_rate_le_one
UEOTQM.QuantumDoob.InvertibleFilter
UEOTQM.QuantumDoob.InvertibleFilter.mk
UEOTQM.QuantumDoob.InvertibleFilter.apply
UEOTQM.QuantumDoob.InvertibleFilter.unapply
UEOTQM.QuantumDoob.InvertibleFilter.metric
UEOTQM.QuantumDoob.InvertibleFilter.apply_smul
UEOTQM.QuantumDoob.InvertibleFilter.unapply_apply
UEOTQM.QuantumDoob.InvertibleFilter.apply_unapply
UEOTQM.QuantumDoob.InvertibleFilter.apply_posSemidef
UEOTQM.QuantumDoob.InvertibleFilter.mass_apply_pos
UEOTQM.QuantumDoob.InvertibleFilter.state
UEOTQM.QuantumDoob.InvertibleFilter.state_matrix
UEOTQM.QuantumDoob.InvertibleFilter.mass_apply_eq_pairing
UEOTQM.QuantumDoob.InvertibleFilter.metric_conjugate
UEOTQM.QuantumDoob.transformed_gram
UEOTQM.QuantumDoob.transformed_map
UEOTQM.QuantumDoob.doob_complete
UEOTQM.QuantumDoob.doobWeighted
UEOTQM.QuantumDoob.doobInstrument
UEOTQM.QuantumDoob.doob_branch
UEOTQM.QuantumDoob.doob_branch_apply
UEOTQM.QuantumDoob.doob_branch_state
UEOTQM.QuantumDoob.doob_channelState
UEOTQM.QuantumDoob.doob_successOperator_apply
UEOTQM.QuantumDoob.eigenmetric_successMass
UEOTQM.QuantumDoob.eigenmetric_successMass_pos
UEOTQM.QuantumDoob.doob_finite_conditioning
UEOTQM.QuantumDoob.doob_successOperator_state
UEOTQM.QuantumDoob.doob_successMass_state
UEOTQM.RecordProbability.eventProbability
UEOTQM.RecordProbability.positiveEvent
UEOTQM.RecordProbability.eventProbability_empty
UEOTQM.RecordProbability.eventProbability_univ
UEOTQM.RecordProbability.eventProbability_nonneg
UEOTQM.RecordProbability.eventProbability_le_one
UEOTQM.RecordProbability.eventProbability_eq_indicator
UEOTQM.RecordProbability.eventProbability_abs_sub_le_tv
UEOTQM.RecordProbability.positiveEvent_sub_eq_tv
UEOTQM.RecordProbability.positiveEvent_abs_sub_eq_tv
UEOTQM.RecordProbability.tv_le_iff_events
UEOTQM.RectangularChannel
UEOTQM.RectangularChannel.operator
UEOTQM.RectangularChannel.operator_posSemidef
UEOTQM.RectangularChannel.operator_zero
UEOTQM.RectangularChannel.operator_add
UEOTQM.RectangularChannel.operator_smul
UEOTQM.RectangularChannel.operator_real_smul
UEOTQM.RectangularChannel.operatorLinear
UEOTQM.RectangularChannel.operator_sum
UEOTQM.RectangularChannel.trace_operator
UEOTQM.RectangularChannel.mass_operator
UEOTQM.RectangularChannel.state
UEOTQM.RectangularChannel.state_matrix
UEOTQM.RectangularChannel.adjoint
UEOTQM.RectangularChannel.adjoint_posSemidef
UEOTQM.RectangularChannel.adjoint_one
\end{Verbatim}

\hypertarget{f.3.11-ux8bf7ux6c42601660}{%
\subsection{F.3.11 请求601--660}\label{f.3.11-ux8bf7ux6c42601660}}

\begin{Verbatim}[breaklines,breakanywhere,fontsize=\scriptsize]
UEOTQM.RectangularChannel.adjoint_sub
UEOTQM.RectangularChannel.trace_adjoint
UEOTQM.RectangularChannel.pullbackEffect
UEOTQM.RectangularChannel.pullbackEffect_matrix
UEOTQM.RectangularChannel.probability_state
UEOTQM.RectangularChannel.state_effectDistance_le
UEOTQM.RectangularChannel.identity
UEOTQM.RectangularChannel.identity_operator
UEOTQM.RectangularChannel.identity_state
UEOTQM.RectangularChannel.comp
UEOTQM.RectangularChannel.comp_operator
UEOTQM.RectangularChannel.comp_state
UEOTQM.RectangularChannel.amplify
UEOTQM.RectangularChannel.amplify_posSemidef
UEOTQM.RectangularChannel.amplify_operator_apply
UEOTQM.RectangularChannel.discardRightKraus
UEOTQM.RectangularChannel.discardRight
UEOTQM.RectangularChannel.discardRight_operator_apply
UEOTQM.RectangularChannel.discardLeftKraus
UEOTQM.RectangularChannel.discardLeft
UEOTQM.RectangularChannel.discardLeft_operator_apply
UEOTQM.RectangularChannel.defect
UEOTQM.RectangularChannel.defect_bddAbove
UEOTQM.RectangularChannel.defect_nonneg
UEOTQM.RectangularChannel.defect_le_one
UEOTQM.RectangularChannel.state_distance_le_defect
UEOTQM.RectangularChannel.perturbation_bound
UEOTQM.QuantumPurification.mix
UEOTQM.QuantumPurification.mix_matrix
UEOTQM.QuantumPurification.IsPure
UEOTQM.QuantumPurification.outer
UEOTQM.QuantumPurification.outer_apply
UEOTQM.QuantumPurification.outer_positive
UEOTQM.QuantumPurification.outer_hermitian
UEOTQM.QuantumPurification.outer_trace
UEOTQM.QuantumPurification.outer_square
UEOTQM.QuantumPurification.outer_sandwich
UEOTQM.QuantumPurification.vectorState
UEOTQM.QuantumPurification.vectorState_matrix
UEOTQM.QuantumPurification.exists_gram
UEOTQM.QuantumPurification.positive_compression_zero
UEOTQM.QuantumPurification.density_eq_vector_of_compression_zero
UEOTQM.QuantumPurification.vectorState_isPure
UEOTQM.QuantumPurification.vectorize
UEOTQM.QuantumPurification.vectorize_outer_trace
UEOTQM.QuantumPurification.discard_vectorize_outer
UEOTQM.QuantumPurification.gram
UEOTQM.QuantumPurification.gram_spec
UEOTQM.QuantumPurification.purificationVector
UEOTQM.QuantumPurification.purificationVector_normalized
UEOTQM.QuantumPurification.purification
UEOTQM.QuantumPurification.purification_matrix
UEOTQM.QuantumPurification.purification_isPure
UEOTQM.QuantumPurification.purification_marginal
UEOTQM.QuantumPurification.exists_purification
UEOTQM.QuantumPurification.density_empty_elim
UEOTQM.QuantumPurification.product
UEOTQM.QuantumPurification.product_matrix
UEOTQM.QuantumPurification.discard_product
UEOTQM.QuantumPurification.probability_mix
\end{Verbatim}

\hypertarget{f.3.12-ux8bf7ux6c42661720}{%
\subsection{F.3.12 请求661--720}\label{f.3.12-ux8bf7ux6c42661720}}

\begin{Verbatim}[breaklines,breakanywhere,fontsize=\scriptsize]
UEOTQM.QuantumPurification.state_mix
UEOTQM.OperationalSystems.FiniteSystem
UEOTQM.OperationalSystems.FiniteSystem.of
UEOTQM.OperationalSystems.FiniteSystem.product
UEOTQM.OperationalSystems.Presentation
UEOTQM.OperationalSystems.Presentation.discardRight
UEOTQM.OperationalSystems.Laws
UEOTQM.OperationalSystems.IsPure
UEOTQM.OperationalSystems.Purifies
UEOTQM.OperationalSystems.PurificationExistence
UEOTQM.OperationalSystems.LocalTomography
UEOTQM.OperationalSystems.ClassicalModel.product
UEOTQM.OperationalSystems.ClassicalModel.productEffect
UEOTQM.OperationalSystems.ClassicalModel.act
UEOTQM.OperationalSystems.ClassicalModel.probability_mix
UEOTQM.OperationalSystems.ClassicalModel.product_probability
UEOTQM.OperationalSystems.ClassicalModel.marginal_product
UEOTQM.OperationalSystems.ClassicalModel.act_discard
UEOTQM.OperationalSystems.ClassicalModel.act_mix
UEOTQM.OperationalSystems.ClassicalModel.act_identity
UEOTQM.OperationalSystems.ClassicalModel.act_comp
UEOTQM.OperationalSystems.classicalModel
UEOTQM.OperationalSystems.classical_laws
UEOTQM.OperationalSystems.classical_pure_iff
UEOTQM.OperationalSystems.classical_discard_eq
UEOTQM.OperationalSystems.classical_not_purification
UEOTQM.OperationalSystems.coreSystem
UEOTQM.OperationalSystems.core_first_marginal_no_purifies
UEOTQM.OperationalSystems.QuantumProcess
UEOTQM.OperationalSystems.QuantumModel.trace_product_im_zero
UEOTQM.OperationalSystems.QuantumModel.product_probability
UEOTQM.OperationalSystems.complexModel
UEOTQM.OperationalSystems.complex_laws
UEOTQM.OperationalSystems.complex_pure_iff
UEOTQM.OperationalSystems.complex_purification
UEOTQM.OperationalSystems.RealModel.IsRealMatrix
UEOTQM.OperationalSystems.RealModel.real_add
UEOTQM.OperationalSystems.RealModel.real_smul
UEOTQM.OperationalSystems.RealModel.real_mul
UEOTQM.OperationalSystems.RealModel.real_conjTranspose
UEOTQM.OperationalSystems.RealModel.real_kronecker
UEOTQM.OperationalSystems.RealModel.State
UEOTQM.OperationalSystems.RealModel.Effect
UEOTQM.OperationalSystems.RealModel.mix
UEOTQM.OperationalSystems.RealModel.product
UEOTQM.OperationalSystems.RealModel.productEffect
UEOTQM.OperationalSystems.RealModel.Process
UEOTQM.OperationalSystems.RealModel.operator_real
UEOTQM.OperationalSystems.RealModel.act
UEOTQM.OperationalSystems.RealModel.identity
UEOTQM.OperationalSystems.RealModel.compose
UEOTQM.OperationalSystems.RealModel.discardProcess
UEOTQM.OperationalSystems.realModel
UEOTQM.OperationalSystems.real_laws
UEOTQM.OperationalSystems.real_not_localTomography
UEOTQM.ClassicalQuantumMemory.taggedKraus
UEOTQM.ClassicalQuantumMemory.taggedKraus_gram
UEOTQM.ClassicalQuantumMemory.taggedKraus_conjugation
UEOTQM.ClassicalQuantumMemory.block
UEOTQM.ClassicalQuantumMemory.block_apply_same
\end{Verbatim}

\hypertarget{f.3.13-ux8bf7ux6c42721780}{%
\subsection{F.3.13 请求721--780}\label{f.3.13-ux8bf7ux6c42721780}}

\begin{Verbatim}[breaklines,breakanywhere,fontsize=\scriptsize]
UEOTQM.ClassicalQuantumMemory.block_apply_ne
UEOTQM.ClassicalQuantumMemory.trace_block
UEOTQM.ClassicalQuantumMemory.channel
UEOTQM.ClassicalQuantumMemory.channel_operator_apply
UEOTQM.ClassicalQuantumMemory.channel_amplify_operator_apply
UEOTQM.ClassicalQuantumMemory.channel_operator_eq_block
UEOTQM.ClassicalQuantumMemory.state
UEOTQM.ClassicalQuantumMemory.state_matrix
UEOTQM.ClassicalQuantumMemory.discard_record_operator
UEOTQM.ClassicalQuantumMemory.discard_record_state
UEOTQM.ClassicalQuantumMemory.discard_memory_operator
UEOTQM.ClassicalQuantumMemory.discard_memory_state_matrix
UEOTQM.ClassicalQuantumMemory.state_block_mass
UEOTQM.ClassicalQuantumMemory.positive_block_posterior
UEOTQM.ClassicalQuantumMemory.zero_probability_iff_zero_block
UEOTQM.ClassicalQuantumMemory.eventEffect
UEOTQM.ClassicalQuantumMemory.eventEffect_matrix
UEOTQM.ClassicalQuantumMemory.eventEffect_pairing_block
UEOTQM.ClassicalQuantumMemory.weights
UEOTQM.ClassicalQuantumMemory.weights_weight
UEOTQM.ClassicalQuantumMemory.event_probability
UEOTQM.ClassicalQuantumMemory.record_tv_le_joint_distance
UEOTQM.ClassicalQuantumMemory.defect
UEOTQM.ClassicalQuantumMemory.perturbation_bound
UEOTQM.ClassicalQuantumMemory.record_tv_perturbation_bound
UEOTQM.ClassicalQuantumMemory.ProtocolMemory.Leaf
UEOTQM.ClassicalQuantumMemory.ProtocolMemory.KrausPath
UEOTQM.ClassicalQuantumMemory.ProtocolMemory.pathKraus
UEOTQM.ClassicalQuantumMemory.ProtocolMemory.path_complete
UEOTQM.ClassicalQuantumMemory.ProtocolMemory.path_branch_eq_recordOperator
UEOTQM.ClassicalQuantumMemory.ProtocolMemory.channel
UEOTQM.ClassicalQuantumMemory.ProtocolMemory.channel_operator_apply
UEOTQM.ClassicalQuantumMemory.ProtocolMemory.channel_amplify_operator_apply
UEOTQM.ClassicalQuantumMemory.ProtocolMemory.channel_operator_eq_block
UEOTQM.ClassicalQuantumMemory.ProtocolMemory.state
UEOTQM.ClassicalQuantumMemory.ProtocolMemory.state_matrix
UEOTQM.ClassicalQuantumMemory.ProtocolMemory.discard_record_operator
UEOTQM.ClassicalQuantumMemory.ProtocolMemory.discard_memory_operator
UEOTQM.ClassicalQuantumMemory.ProtocolMemory.discard_memory_state_matrix
UEOTQM.ClassicalQuantumMemory.ProtocolMemory.state_block_mass
UEOTQM.ClassicalQuantumMemory.ProtocolMemory.zero_probability_iff_zero_block
UEOTQM.ClassicalQuantumMemory.ProtocolMemory.weights
UEOTQM.ClassicalQuantumMemory.ProtocolMemory.weights_weight
UEOTQM.ClassicalQuantumMemory.ProtocolMemory.event_probability
UEOTQM.ClassicalQuantumMemory.ProtocolMemory.record_tv_le_joint_distance
UEOTQM.ClassicalQuantumMemory.ProtocolMemory.defect
UEOTQM.ClassicalQuantumMemory.ProtocolMemory.perturbation_bound
UEOTQM.ClassicalQuantumMemory.ProtocolMemory.record_tv_perturbation_bound
UEOTQM.QuantumDoobModel.successK
UEOTQM.QuantumDoobModel.failureK
UEOTQM.QuantumDoobModel.success_gram
UEOTQM.QuantumDoobModel.failure_gram
UEOTQM.QuantumDoobModel.sourceK
UEOTQM.QuantumDoobModel.source_complete
UEOTQM.QuantumDoobModel.sourceInstrument
UEOTQM.QuantumDoobModel.source_branch_formula
UEOTQM.QuantumDoobModel.rhoMatrix
UEOTQM.QuantumDoobModel.rhoMatrix_positive
UEOTQM.QuantumDoobModel.rhoMatrix_trace
UEOTQM.QuantumDoobModel.rho
\end{Verbatim}

\hypertarget{f.3.14-ux8bf7ux6c42781840}{%
\subsection{F.3.14 请求781--840}\label{f.3.14-ux8bf7ux6c42781840}}

\begin{Verbatim}[breaklines,breakanywhere,fontsize=\scriptsize]
UEOTQM.QuantumDoobModel.S
UEOTQM.QuantumDoobModel.T
UEOTQM.QuantumDoobModel.inverse_mul
UEOTQM.QuantumDoobModel.mul_inverse
UEOTQM.QuantumDoobModel.filter
UEOTQM.QuantumDoobModel.metric_exact
UEOTQM.QuantumDoobModel.source_right_eigen
UEOTQM.QuantumDoobModel.source_left_eigen
UEOTQM.QuantumDoobModel.filtered_mass_exact
UEOTQM.QuantumDoobModel.plusMatrix
UEOTQM.QuantumDoobModel.filtered_state_exact
UEOTQM.QuantumDoobModel.plus
UEOTQM.QuantumDoobModel.plus_matrix
UEOTQM.QuantumDoobModel.transformedInstrument
UEOTQM.QuantumDoobModel.transformed_branch_formula
UEOTQM.QuantumDoobModel.success_kraus_noncommuting
UEOTQM.QuantumDoobModel.input_coherence_retained
UEOTQM.QuantumDoobDynamics.decay
UEOTQM.QuantumDoobDynamics.decay_zero
UEOTQM.QuantumDoobDynamics.decay_pos
UEOTQM.QuantumDoobDynamics.decay_le_one
UEOTQM.QuantumDoobDynamics.decay_succ
UEOTQM.QuantumDoobDynamics.decay_tendsto
UEOTQM.QuantumDoobDynamics.effectDistance_eq_of_matrix_sub_smul
UEOTQM.QuantumDoobDynamics.doob_branch_formula
UEOTQM.QuantumDoobDynamics.doob_branch_trace
UEOTQM.QuantumDoobDynamics.doob_successOperator_formula
UEOTQM.QuantumDoobDynamics.doob_iterateState_matrix
UEOTQM.QuantumDoobDynamics.doob_stationary_plus
UEOTQM.QuantumDoobDynamics.doob_stationary_unique
UEOTQM.QuantumDoobDynamics.doob_effectDistance_exact
UEOTQM.QuantumDoobDynamics.doob_matrix_tendsto
UEOTQM.QuantumDoobDynamics.filterMass
UEOTQM.QuantumDoobDynamics.pullbackWeight
UEOTQM.QuantumDoobDynamics.normalizer
UEOTQM.QuantumDoobDynamics.filterMass_diag
UEOTQM.QuantumDoobDynamics.filterMass_trace
UEOTQM.QuantumDoobDynamics.filterMass_metric_trace
UEOTQM.QuantumDoobDynamics.filterMass_bounds
UEOTQM.QuantumDoobDynamics.pullbackWeight_bounds
UEOTQM.QuantumDoobDynamics.normalizer_bounds
UEOTQM.QuantumDoobDynamics.normalizer_pos
UEOTQM.QuantumDoobDynamics.unapply_plus
UEOTQM.QuantumDoobDynamics.source_successOperator_formula
UEOTQM.QuantumDoobDynamics.source_successOperator_density
UEOTQM.QuantumDoobDynamics.source_successMass_formula
UEOTQM.QuantumDoobDynamics.source_response_formula
UEOTQM.QuantumDoobDynamics.source_successMass_pos
UEOTQM.QuantumDoobDynamics.source_successMass_bounds
UEOTQM.QuantumDoobDynamics.source_conditionalState_matrix
UEOTQM.QuantumDoobDynamics.conditioned
UEOTQM.QuantumDoobDynamics.conditioned_matrix
UEOTQM.QuantumDoobDynamics.conditioned_matrix_affine
UEOTQM.QuantumDoobDynamics.conditioned_matrix_sub
UEOTQM.QuantumDoobDynamics.conditioned_effectDistance_exact
UEOTQM.QuantumDoobDynamics.conditioned_effectDistance_le
UEOTQM.QuantumDoobDynamics.normalizer_tendsto
UEOTQM.QuantumDoobDynamics.conditionalWeight_tendsto
UEOTQM.QuantumDoobDynamics.conditioned_matrix_tendsto
UEOTQM.QuantumDoobDynamics.conditioned_effectDistance_tendsto
\end{Verbatim}

\hypertarget{f.3.15-ux8bf7ux6c42841900}{%
\subsection{F.3.15 请求841--900}\label{f.3.15-ux8bf7ux6c42841900}}

\begin{Verbatim}[breaklines,breakanywhere,fontsize=\scriptsize]
UEOTQM.QuantumDoobDynamics.source_successMass_tendsto_zero
UEOTQM.QuantumDoobDynamics.filterMass_rho
UEOTQM.QuantumDoobDynamics.pullbackWeight_rho
UEOTQM.QuantumDoobDynamics.normalizer_rho
UEOTQM.QuantumDoobDynamics.source_successMass_rho
UEOTQM.OperationalSystems.finiteSystemCoeSort
UEOTQM.ClassicalQuantumMemory.ProtocolMemory.krausPathFintype
UEOTQM.OperationalSystems.FiniteSystem.finite
UEOTQM.OperationalSystems.FiniteSystem.decidableEq
UEOTQM.RectangularChannel.mk
UEOTQM.OperationalSystems.FiniteSystem.mk
UEOTQM.OperationalSystems.Presentation.mk
UEOTQM.OperationalSystems.Laws.mk
UEOTQM.OperationalSystems.QuantumProcess.mk
UEOTQM.OperationalSystems.RealModel.Process.mk
UEOTQM.QuantumSeparation.boundedHermitianEffect
UEOTQM.QuantumSeparation.boundedHermitianEffect_matrix
UEOTQM.QuantumSeparation.boundedHermitianEffect_affine
UEOTQM.QuantumSeparation.exists_effect_affine_of_hermitian
UEOTQM.QuantumSeparation.hermitian_mem_real_span_effects
UEOTQM.QuantumSeparation.pairing_eq_of_all_effect_pairings_eq
UEOTQM.QuantumSeparation.pairing_sub_right
UEOTQM.QuantumSeparation.density_eq_of_all_effect_probabilities_eq
UEOTQM.QuantumSeparation.density_eq_iff_all_effect_probabilities_eq
UEOTQM.QuantumSeparation.effectDistance_eq_zero_iff_density_eq
UEOTQM.ComplexLocalTomography.symmetricUnit
UEOTQM.ComplexLocalTomography.imaginaryUnit
UEOTQM.ComplexLocalTomography.symmetricUnit_hermitian
UEOTQM.ComplexLocalTomography.imaginaryUnit_hermitian
UEOTQM.ComplexLocalTomography.sub_kronecker
UEOTQM.ComplexLocalTomography.kronecker_sub
UEOTQM.ComplexLocalTomography.pairing_symmetricUnit
UEOTQM.ComplexLocalTomography.pairing_imaginaryUnit
UEOTQM.ComplexLocalTomography.symmetricUnit_product
UEOTQM.ComplexLocalTomography.imaginaryUnit_product
UEOTQM.ComplexLocalTomography.hermitian_product_pairing_eq
UEOTQM.ComplexLocalTomography.density_eq_of_all_product_effect_probabilities_eq
UEOTQM.ComplexLocalTomography.complex_localTomography
UEOTQM.QuantumPurification.outer_real_smul
UEOTQM.QuantumPurification.gram_eq_sum_outer
UEOTQM.QuantumPurification.IsPure.normalize_eq_of_positive_summand
UEOTQM.QuantumPurification.IsPure.exists_vector
UEOTQM.QuantumPurification.isPure_iff_exists_vector
UEOTQM.QuantumPurification.inner_matrix_image
UEOTQM.QuantumPurification.norm_matrix_image_eq_of_gram
UEOTQM.QuantumPurification.exists_unitary_of_column_gram
UEOTQM.QuantumPurification.exists_auxiliary_unitary_of_gram
UEOTQM.QuantumPurification.auxiliary_action_apply
UEOTQM.QuantumPurification.auxiliary_action_vectorize
UEOTQM.QuantumPurification.outer_mulVec
UEOTQM.QuantumPurification.exists_auxiliary_unitary_of_equal_marginal
UEOTQM.QuantumPurification.pure_states_unitarily_equivalent
UEOTQM.QuantumPurification.auxiliary_tensor_isometry
UEOTQM.QuantumPurification.auxiliaryEmbeddingChannel
UEOTQM.QuantumPurification.auxiliaryEmbeddingChannel_operator
UEOTQM.QuantumPurification.auxiliaryEmbeddingChannel_operator_apply
UEOTQM.QuantumPurification.auxiliaryEmbeddingChannel_discard_operator
UEOTQM.QuantumPurification.auxiliaryEmbeddingChannel_discard_state
UEOTQM.QuantumPurification.auxiliaryEmbeddingChannel_preserves_purity
UEOTQM.QuantumPurification.auxiliaryInl
\end{Verbatim}

\hypertarget{f.3.16-ux8bf7ux6c42901960}{%
\subsection{F.3.16 请求901--960}\label{f.3.16-ux8bf7ux6c42901960}}

\begin{Verbatim}[breaklines,breakanywhere,fontsize=\scriptsize]
UEOTQM.QuantumPurification.auxiliaryInr
UEOTQM.QuantumPurification.auxiliaryInl_isometry
UEOTQM.QuantumPurification.auxiliaryInr_isometry
UEOTQM.QuantumPurification.extendAuxiliaryLeft
UEOTQM.QuantumPurification.extendAuxiliaryRight
UEOTQM.QuantumPurification.extendAuxiliaryLeft_pure
UEOTQM.QuantumPurification.extendAuxiliaryRight_pure
UEOTQM.QuantumPurification.extendAuxiliaryLeft_marginal
UEOTQM.QuantumPurification.extendAuxiliaryRight_marginal
UEOTQM.QuantumPurification.pure_states_common_extension_unitary
UEOTQM.HistoryChannels.principal
UEOTQM.HistoryChannels.principal_apply
UEOTQM.HistoryChannels.principal_posSemidef
UEOTQM.HistoryChannels.sum_mass_principal
UEOTQM.HistoryChannels.weight
UEOTQM.HistoryChannels.weight_nonneg
UEOTQM.HistoryChannels.sum_weight
UEOTQM.HistoryChannels.zero_weight_iff_zero_principal
UEOTQM.HistoryChannels.restrictEffect
UEOTQM.HistoryChannels.restrictEffect_matrix
UEOTQM.HistoryChannels.pairing_block
UEOTQM.HistoryChannels.readKraus
UEOTQM.HistoryChannels.readKraus_gram
UEOTQM.HistoryChannels.readKraus_conjugation
UEOTQM.HistoryChannels.controlled
UEOTQM.HistoryChannels.controlled_operator_apply
UEOTQM.HistoryChannels.controlled_operator_eq_block
UEOTQM.HistoryChannels.pairing_operator_difference_le_mass_mul_defect
UEOTQM.HistoryChannels.controlled_state_distance_le_weighted_defect
UEOTQM.HistoryChannels.controlled_state_distance_le_uniform
UEOTQM.HistoryChannels.controlled_defect_le_uniform
UEOTQM.HistoryChannels.weights
UEOTQM.HistoryChannels.weights_weight
UEOTQM.HistoryChannels.event_probability
UEOTQM.HistoryChannels.record_tv_le_joint_distance
UEOTQM.HistoryChannels.relabelKraus
UEOTQM.HistoryChannels.relabelKraus_gram
UEOTQM.HistoryChannels.relabel
UEOTQM.HistoryChannels.relabel_operator_apply
UEOTQM.HistoryChannels.relabel_state_matrix
UEOTQM.HistoryChannels.relabel_state_distance
UEOTQM.ActionRecordBudget.History
UEOTQM.ActionRecordBudget.Controller
UEOTQM.ActionRecordBudget.Controller.dropLast
UEOTQM.ActionRecordBudget.Controller.shift
UEOTQM.ActionRecordBudget.Controller.dropLast_shift
UEOTQM.ActionRecordBudget.Controller.word
UEOTQM.ActionRecordBudget.Controller.word_length
UEOTQM.ActionRecordBudget.Controller.protocol
UEOTQM.ActionRecordBudget.Controller.leafEquiv
UEOTQM.ActionRecordBudget.Controller.protocol_word_eq
UEOTQM.ActionRecordBudget.Controller.snoc_head
UEOTQM.ActionRecordBudget.Controller.tail_snoc
UEOTQM.ActionRecordBudget.Controller.word_snoc
UEOTQM.ActionRecordBudget.Controller.truncate
UEOTQM.ActionRecordBudget.Controller.truncate_last
UEOTQM.ActionRecordBudget.Controller.truncate_dropLast
UEOTQM.ActionRecordBudget.initialEquiv
UEOTQM.ActionRecordBudget.extendEquiv
UEOTQM.ActionRecordBudget.step
\end{Verbatim}

\hypertarget{f.3.17-ux8bf7ux6c42961985}{%
\subsection{F.3.17 请求961--985}\label{f.3.17-ux8bf7ux6c42961985}}

\begin{Verbatim}[breaklines,breakanywhere,fontsize=\scriptsize]
UEOTQM.ActionRecordBudget.step_defect_le_uniform
UEOTQM.ActionRecordBudget.record_defect_self
UEOTQM.ActionRecordBudget.step_state_distance_le_weighted
UEOTQM.ActionRecordBudget.step_state_distance_le_uniform
UEOTQM.ActionRecordBudget.step_operator_apply
UEOTQM.ActionRecordBudget.step_operator_eq_block
UEOTQM.ActionRecordBudget.run
UEOTQM.ActionRecordBudget.recordOperator_append
UEOTQM.ActionRecordBudget.run_matrix
UEOTQM.ActionRecordBudget.run_eq_protocol_state
UEOTQM.ActionRecordBudget.terminal_memory_matrix
UEOTQM.ActionRecordBudget.weight_run
UEOTQM.ActionRecordBudget.weight_run_leaf
UEOTQM.ActionRecordBudget.actionCost
UEOTQM.ActionRecordBudget.actionCost_nonneg
UEOTQM.ActionRecordBudget.actionCost_castSucc
UEOTQM.ActionRecordBudget.actionCost_last
UEOTQM.ActionRecordBudget.run_succ_distance_le
UEOTQM.ActionRecordBudget.run_distance_le_actionCost_sum
UEOTQM.ActionRecordBudget.actionCost_le
UEOTQM.ActionRecordBudget.run_distance_le_budget_sum
UEOTQM.ActionRecordBudget.record_tv_le_actionCost_sum
UEOTQM.ActionRecordBudget.record_tv_le_budget_sum
UEOTQM.ActionRecordBudget.protocol_record_tv_eq
UEOTQM.ActionRecordBudget.protocol_record_tv_le_budget_sum
\end{Verbatim}

\backmatter
\end{document}
